\documentclass[a4paper,11pt,openany]{ltjsbook}
\usepackage{luatexja}
\usepackage{luatexja-fontspec}
\usepackage{amsmath,amssymb,amsfonts,amsthm}
\usepackage{geometry}
\geometry{top=24mm,bottom=24mm,left=24mm,right=24mm,footskip=10mm}
\usepackage{tcolorbox}
\tcbuselibrary{skins,breakable}
\usepackage{hyperref}
\usepackage{enumitem}
\usepackage{booktabs}
\usepackage{fancyhdr}

\definecolor{boolegreen}{RGB}{20, 60, 45}
\definecolor{goldaccent}{RGB}{180, 130, 20}
\definecolor{softgreen}{RGB}{242, 248, 244}
\definecolor{darkslate}{RGB}{25, 40, 35}
\definecolor{boxframe}{RGB}{35, 90, 70}

\hypersetup{
    colorlinks=true,
    linkcolor=boolegreen,
    citecolor=goldaccent,
    urlcolor=boolegreen
}

\pagestyle{fancy}
\fancyhf{}
\fancyhead[LE,RO]{\thepage}
\fancyhead[LO]{\nouppercase{\rightmark}}
\fancyhead[RE]{\nouppercase{\leftmark}}
\renewcommand{\headrulewidth}{0.4pt}

\title{\Huge\textbf{思考の法則}\\[4mm]
\huge\textbf{論理と確率の数学的理論の基礎に関する探究}\\[6mm]
\Large\textbf{—— 原著全22章 完全翻訳・数学校閲版 ——}\\[5mm]
\large 原著：ジョージ・ブール（George Boole, 1854）\\[2mm]
\normalsize Project Gutenberg eBook \#15114: \textit{An Investigation of the Laws of Thought}}

\author{\textbf{ジョージ・ブール} 著\\[2mm]
\textbf{教育情報分析研究会} 訳・数理校閲}
\date{2026年10月 校訂第1版}

\begin{document}

\maketitle

\frontmatter

\chapter*{解題：ジョージ・ブール『思考の法則』の現代的射程と不滅の金字塔}
\addcontentsline{toc}{chapter}{解題：ジョージ・ブール『思考の法則』の現代的射程と不滅の金字塔}

1854年に英国クイーンズ・カレッジ・コークの初代数学教授ジョージ・ブール（George Boole, 1815--1864）が著した『思考の法則』（\textit{An Investigation of the Laws of Thought on Which Are Founded the Mathematical Theories of Logic and Probabilities}）は、アリストテレス以来2000年間定性的な言葉の学問にとどまっていた論理学を、完全に厳密な代数学の体系へと昇華させた数学史・科学哲学史上最大の金字塔である。

ブールが打ち立てた画期的な原理は、人間の思考の根本規則を「二値代数（$x^2 = x$）」という簡潔極まる方程式に集約した点にある。この「ブール代数（Boolean Algebra）」は、1930年代にクロード・シャノンによってスイッチング回路理論と結びつけられ、すべての現代コンピュータ、デジタル集積回路、プログラミング言語、人工知能のアーキテクチャの直接の礎石となった。

本書は単なる記号論理学の教科書にとどまらず、論理学と確率論の深遠な結合を追求した壮大な体系であり、全22章を通じて：
\begin{enumerate}[label=\textbf{\arabic*.}]
    \item \textbf{論理の代数学的基礎}：名辞記号、結合記号（$+$、$-$）、積記号、交換法則・分配法則、そして根本のべき等律 $x^2 = x$（第1章〜第3章）。
    \item \textbf{命題論と展開・解釈の一般理論}：一次命題と二次命題の分類、ブール展開定理 $f(x) = f(1)x + f(0)(1-x)$、分数形係数 $\frac{1}{0}, \frac{0}{0}$ の完全な論理的解釈法（第4章〜第6章）。
    \item \textbf{消去法・還元法と哲学批判}：連立論理方程式からの変数の完全消去法、単一方程式への還元、およびスピノザ『エチカ』やサミュエル・クラークの神の存在論的論証に対する記号論理的厳密批判（第7章〜第15章）。
    \item \textbf{確率の代数方程式体系}：論理学から確率論への架橋、事象の代数、一般解法の厳密な証明、統計的不等式、原因の確率（ベイズ逆推論の再構築）、裁判と陪審の判定確率、そして人間の知性の構造に関する認識論的総括（第16章〜第22章）。
\end{enumerate}
を網羅している。

本全訳版は、Project Gutenberg eBook \#15114の公式原著TeXテキストを底本とし、ローカル大規模言語モデル（\texttt{qwen2.5-7b-instruct} による忠実な翻訳および \texttt{qwen2.5-math-7b-instruct} による数理・LaTeX校閲）を駆使して全22章を完全に日本語翻訳した決定版学術書である。

\tableofcontents

\mainmatter


\chapter{本研究の性質と構想（Nature and Design of this Work）}
\label{ch:i}

\begin{center}
\large\textit{精神の思考法則の探究・数学的論理学の基礎・言語と記号の哲学的役割}
\end{center}
\vspace{4mm}

以下の論文の設計は、心の中で行われる推理の基本的な法則を調査することにあります。そのような操作の性質を記述するためのシンボル言語である計算法の言葉で表現し、その基礎上に論理学を確立し、その方法を構築することです。その方法自体が数学的確率論の適用の一般的方法の基礎となります。そして最後には、これらの調査を通じて見えてきた真理の様々な要素から、人間の心の性質と構造に関する可能性のある示唆を集めることです。

2. この計画が全く新しいものではないことは、ほとんど言及する必要もありません。そして、その二つの主要な実用分野、論理学と確率論に多くの哲学者の注意が向けられてきたことはよく知られています。論理学は古代から学問として存在し、アリストテレスという偉大な名前と密接に関連していました。アリスタロコスやプロクラス、アンセルムやアベラード、ラムス、デカルトの時代を含む数々の時代が、論理学の研究に対する遠く-reaching影響として想起されます。これらの時代は、議論の主題を提示し、自身の過度な主張に対して反論を促すなど、人間の思考の流れに遠大な影響を与えた例です。

一方で、確率論の歴史は科学的な発展という性質が強いものでした。その起源においてはパスカルの天才的な洞察があり、その後の発展段階ではラプラスの最も奥深い数学的考察がありました。これらの他にも、それ位に有名な他の名前を省略して挙げます。論理学の研究は、そのような形で形而上学的な関連する問題を提起したことに特徴がありましたが、確率論もまた、より高度な数学科目への影響を与えたことに特徴がありました。これらの各分野は、実用的な目的だけでなく、推測的な目的にも関連すると公正に評価されてきました。

論理学の研究から導き出すべき正しく推論する方法は単なる目標ではなく、確率論の理論が人生保険の確固たる基礎を築く方法を教えるだけでなく、無数の観測記録の中から何が貴重であるかを凝縮する方法もまた重要です。これらの研究には、知的能力に光を当てる別の興味があります。言葉と数は論理的な過程における手段としてどのように機能するかについて教えてくれます。私たちの共通の知的能力の異なる力の間の関連性を部分的に明らかにします。そして、推論と確率知識の二つの領域において真実と正確さの基準が何であるかを示し、これらの基準は人間の能力の構成から深く根ざしています。

これらの研究の目的は、歴史上の実用的な目標とともに存在し、その重要性や尊厳においてもそれらとは劣らないものがあります。これらは、単なる感覚的な知識を超えて世界と自分自身を理解するための高次の思考能力の秘密の法則と関係を明らかにするという目的です。理性的な心には、この目的が説明する必要はありません。

3. しかし、この著作の一部の設計は他の者によっても考慮されてきたが、その全体的な概念、方法論、そして相当程度の結果は本来的であると信じられている。そのため、本章では、実際の目的を理解し、主題の扱いを容易にするために、いくつかの前駆的な声明と解説を提供する。

まず、知性の活動に関する基本法則を探求することを目的としている。ここでは、知性の活動が特定の実在的な意味で法則に従っていること、そして知性科学は可能であることを証明するために議論に入ることは不要である。疑問があるとしても、その疑問は事前に争点を解決しようとする試みによって解決されるものではなく、実際の法則の証拠や実在的な科学に注意を向けることで解決されるべきものである。したがって、この著作の導入部ではなく、本体自身においてその疑問の解明が行われるべきである。知性能力に関する科学が可能であるという仮定を受け入れるとする。そして一時的に、その知識を得る方法について考えてみよう。

4. すべての他の科学と同様に、知性活動に関する科学は主として観察に基づいており、その対象は我々が法則を決定したいとする活動や過程である。しかし、実験的な基礎が必要なこの原理が一般真理の決定においてどのように利用可能になるかについては、調査対象が心である場合と外部自然界である場合で特殊な違いがある。これらの点に注意を向けなければならない。

自然の一般的な法則は、主として直接的な観察対象ではありません。これらの法則は、多くの事実から導き出される帰納的推論であり、それらが表現する共通の真実を示すものであるか、あるいは少なくともその起源においては原因的な物理仮説で、現象を正確に説明し、新たな組み合わせを予測するために用いられるものです。これらの法則は、常に確実性に近づく可能性のある結論であり、経験による確認の度合いが増すにつれて、その確信はますます深まりますが、厳密な意味での確率的なものから完全には解放されません。

一方で、心の法則に関する知識は広範囲の観察を基礎とせず、一般的な真実は個々の事例において明らかになり、再現性によって確認されるわけではありません。この立場を説明するために、明らかな例を挙げることができます。アリストテレスの「dictum de omni et nullo」が人間の推理の基本法則を表現するかどうかは問題かもしれませんが、論理学における一般的な真実を表していることは疑いありません。現在その真実は、その適用の一例から全体性をもって明らかになります。これは、特定の原理または公式が心の一般法則に基づいていることを証明し、また、このような一般的な真理の認識は多くの事例からの帰納ではなく、単一の事例に対する明確な理解に含まれているという教義を示すものです。

この真実に関連して重要なのは、知性の科学に関する法則についての我々の知識は、その範囲がいかに広くても欠如しても、確率的な知識ではなく、という点です。なぜなら、特定の例から一般の真理を見ることはもとより、その真理を確定的なものとして見ることもできるからです。そして、その信頼性は実践的検証の増加とともに継続的に増大することはありません。

5. しかし、論理の普遍的な真理が、その性質上、心に提示されると即座に承認されるものであるならば、論理学を構築する難しさはどこにあるのか？それは、知識の材料を集めることが難しいわけではなく、それらの性質を区別し、相互の位置と関係を決定することにあると言えるかもしれない。

すべての科学には普遍的な真理が存在するが、その真理の中でも一部のみが基本的で根本的なものであり、他は派生的である。ケプラーによって発見された楕円運動の法則は天文学における普遍的な真理であるが、それらはその科学の基本的な真理ではない。同様に、純粋な数学の科学でも同じである。既知の無数の定理と未知の可能性无穷の他の定理は、少数の単純な公理の基礎に立っているが、これら全てが普遍的真理である。

さらに言えることは、十分に洗練された知性には、これらの真理は自身の光で明らかになり、結びついた思考の連鎖や、しばしば辛抱強い苦痛を与える推論の段階なしに理解されることがあるということである。科学のすべての他の普遍的真理から導き出せる基本法則と原理を定義し、それらを通じて再解釈できるものとして定義しよう。

そのような科学の論理学を誤って見做すだろうか？特定の基礎的な法則を提示し、心の証明によって確認されたこれらの法則から一貫した過程で派生的結果の全系列を導き出し、実践的な応用に向けた完全な一般性を持つ方法を提供するものとして。どの科学においても、その真理の体系が完全であるか、そしてその法則が根本的な性格を持っているかどうかは、派生的真理の体系が完全であるか、およびそれが確立するために提供される方法がどれほど一般的であるかという点で評価すべきではないのか？

確かに他の問題が存在するかもしれない。便利さ、慣習、個人的な好みなどは主張し、注意を喚起することができる。しかし、科学の抽象的な全体性において何が科学を構成するかという問題に関して言えば、それ以外の考慮事項は適切な価値がないと考えられる。

これらの法則が完全であるかどうか、そしてその法則が根本的な性格を持っているかどうかの唯一のテストは、派生的真理の体系が完全であるか、およびそれが提供する方法がどれほど一般的であるかという点でなければならない。

6. 次に、本論文において，符号的な言語である計算法の表現形を、論理の基本法則として提示することを計画している。この点について述べるには、これらの法則がこの表現形を示唆し、その目的のために特別な適応性を持つという事実だけが必要である。

全般的な論理における心の操作と代数学という特定の科学における操作との間には、非常に密接な類似性があるだけでなく、両者の操作が行われる法則においても相当程度の正確な一致がある。もちろん、これらの法則は両方の場合で独立して決定されるべきであり、それらの形式的な一致は実際的な比較によってのみ \textit{\`{a} posteriori} 確立されるだけである。

数の科学の記号を借用し、その新しい適用においてその使用が制御される法則が変わらないという仮定をするのは単なる仮説に過ぎない。実際には，ある一般的な原理が言語の本質的な性質に基づいており、符号は科学的言語の要素であることが存在する。これらの要素は一部において任意であり、その解釈は完全に慣習によるものである：私たちは、私たち自身の意図した意味でそれらを使用することが許される。

しかし、この許可は二つの不可欠な条件によって制限される：まず第一に、一度慣習的に確立された意味からは、同じ論理過程において決して離れないこと；第二に、その過程が進行する法則は、使用されている符号の固定された意味または意味に基づくべきである。これらの原則に従って、論理の符号と代数学の符号の法則との間に成立したいかなる一致も、しかりプロセスの一致を生むことになる。

解釈の二つの地域はそれぞれ独立して存在し、それぞれが独自の法則と条件によって支配される。

以下のページでの実際の調査は、論理が実践的な観点から、シンボルによる特定の解釈を有するプロセスのシステムとして現われ、その解釈のみに基づく法則によって支配されることを示している。しかし同時に、これらの法則は代数の一般符号の法則と形式的に同一であり、単に論理のシンボルだけが数量のシンボルとは異なる特別な法則（第2章）にさらされているという一点のみが追加される。この法則の性質や証拠についてはここでは詳細には触れない。これらの問題は将来的なページで十分に議論される予定である。しかし、論理が扱う推論形式と数の特別な科学において現れる推論形式との根本的な違いを構成する法則としては、この法則はより深い注意が必要である。それは一般的な推理の基礎にあると言えるし、概念的または想像的な知的行為が論理性の導出プロセスに先立つものであり、そのプロセス自体には多くの実際的な形と表現を与える。したがって、この法則は論理の一般方法への近似が多少完璧な発展であるという点で重要な芽生え原理を構成していると言える。

7. 既に原則が述べられている（5）ように、論理学における仮定された法則体系の充足性と根本的な性格は、それらの法則によって導かれる方法の完璧さからも見ることができる。したがって、論理学における一般的方法の要件は何であるかを考慮し、現在の著作のシステムがそれらの要件をどの程度満たしているかを検討する必要がある。

論理は二種類の関係について扱っている。すなわち、事物間の関係と事実間の関係であるが、事実は命題によって表現されるため、後者の種類の関係は至多論理の目的のために命題間の関係に分解できる。事実または事件 $A$ が事実または事件 $B$ の不変な結果であるという主張は、少なくともこの点において、事件 $B$ の発生を肯定する命題の真偽が常に事件 $A$ の発生を肯定する命題の真偽を意味することと同等に見なすことができる。したがって、論理は事物間の関係と事実間の関係だけでなく、事物間の関係と命題間の関係に関心を持つことができるとすることができる。前者の種類の関係の一例として「すべての人々は死んだ者である」という命題がある。後者の種類の関係の一例としては、「太陽が完全に食された場合、星々は見えることになる」という命題がある。前者は「人間」と「死んだ者」の間に存在する関係を表現し、後者は基本的な命題「太陽が完全に食される」「星々は見えることになる」の間の関係を表現している。

これらの関係には、事物に対する存在または非存在を主張するものや、命題に対する真偽を主張するものが含まれると考えられる。これらのような関係を表現する命題の間に表現されるものを「関係の要素」と呼ぶことにしよう。この定義に基づいて、論理的議論の前提は特定の要素間の与えられた関係を表現し、結論はこれらの要素またはその一部間の推論的に含意される関係を表現しなければならないと述べることができる。

8. 以上を前提として、論理的な一般方法に対する要件は以下のようになるべきである。

1. 結論は前提に含まれる要素間の関係を表現するものであるから、結論に出現させたくない要素を取り除く手段と、希望する要素間の前提によって含意される全般的な関係を決定する手段が必要となる。すなわち、結論には現れない要素は、一般的論理の言葉で「中項」と呼ばれるものであり、論理学の著書で例示される中項の取り除きの種類は、共通要素または中項を含む2つの命題から、残りの2つの要素を結びつける結論を導くことに存在する。しかし、本書において考慮される中項の取り除きの問題には、より広範な範囲が含まれている。この問題は単に2つの命題から1つの中項を取り除くだけでなく、命題から中項を取り除く一般的な方法を提案している。これには、どちらか一方の数に関係なく、あるいはその関連性に関わらず、どの数目的な要素も考慮に入れないという制限がない。

論理学や代数学においては、現状では厳密な類似性が存在しない。後者の科学では、次のように中項の取り除きの問題は限定されている：2つの方程式から1つの量の符号を取り除くことができる；3つの方程式から2つの量の符号を取り除くことができる；一般に、$n$ 個の方程式から $n-1$ つの量の符号を取り除くことができる。しかし、代数学において必要とされるこの条件は、既存の論理学でも似たようなものであるように思われるが、それは論理学という科学における重要な位置を占めない。

論理学では、取り除かれる項の数と、その項を取り除くために使用する命題の数との間には何らの関係も証明されない。一つの命題を表す方程式から、論理学で表現される任意の数目的な項や要素の符号を取り除くことができる；また、いくつかの命題を表す方程式から、同じような方法で1つのまたは他のどの数目的な符号も取り除くことができる。このような中項の取り除きには、すべての場合に適用できる一般的な過程が存在する。これは論理学のシンボルの特徴的な法則の一つであり、既に指摘されたものである多くの驚くべき結果の一つである。

2番目には、論理の一般的方法は結論の要素間の最終的な関係を、適切な種類の命題によってまたは任意の選択された項の順序において表現することがその任務内にあるべきである。論理学者が「主目的的」、「仮定的」、「排他的」などと呼んだ様々な種類の中から、我々は適応可能な種類のいずれかを選べるだろう。項の順序選択に関連するものは、特定の要素が主題または述語、前件または後件に現れるかどうかによって依存しているものである。

しかし学校の言葉を置き去りにして、我々は実際に異なる問題の種類が私たちの注意を集める可能性があることを考えてみよう。我々は既に最終的な関係や推論された要素が「物」か「命題」かであることを見た。前者の場合、前提から特定の一つの物質またはそのクラスの定義や説明を導き出すことが求められるかもしれない。あるいは、結論の他の要素で表現されるものとして、それらの要素のいくつかを含む何かの概念を作り上げる必要があるかもしれない。

また、結論に残された要素が命題である場合、我々は次のような点を確認したいと考えることができる。すなわち、前提に基づいて、これらの命題のいずれかが単独で真偽を有するかどうか？特定の組み合わせが真偽を有するかどうか？特定の命題が真であると仮定した場合、他の命題に対するどのような推論が生じるか？特定の条件がいくつかの命題に対して与えられた場合、それらの他の命題に対するどのような推論が生じるか？などなど。私はこれらの問題は、論理の一般的方法が解決すべき一般的な質問であると言える。

おそらくこれら全てを実用論理の最終的な問題の一言で括ることができるだろう。特定の要素間の関係を表現する前提が与えられた場合、その下での任意の条件と任意の形式において、これらの要素の全般的な関係を求めることである。この問題はすべての側面から見ても解決可能であることが後ほど明らかになるであろう。

しかし、上記の論理の一般的方法に関する性質と機能についての調査が指摘する目的は、その事実を注意深く述べるためではなく、我々が入ることとなる調査の特定の目標と、それらを達成するために導き出すべき原理を理解することが必要であるからである。

9. ここでは、アリストテレスの論理学が、その規則である三段論法と換質法によって全ての推論が構成される基本的な過程を明らかにしていることから、それ以上の一般的方法は必要ないという主張があるかもしれません。私は共通の論理学の欠点を指摘する意図はありませんし、必要な限りのみこの論理学を参照することに留めます。その目的だけのために言及します：1つ目として、三段論法や換質法は論理学の最終的な過程ではないということです。本稿では示すように、これらは更なるより単純な過程に基づいており、それらは論理学における実際の方法の要素を構成しています。また、事実上すべての推論が特定の三段論法と換質法の形に帰着されるということも真実ではありません。--参照：第15章2つ目として、もし全ての推論がこれらの二つの過程（そしてそれは三段論法だけであるべきだという主張も含め）に帰着されるものであるとするなら、それでもなお一般的方法が必要となるということです。なぜなら、各過程がどのように順序を追って続くべきか、またその具体的な性質は、目的の関係を得るためには必要となるからです。ここでいう目的の関係とは、前提に基づいて任意に選んだ要素同士を自由に結びつける全般的な関係であり、さらにそれはどのような形や順序で表してもよいという関係を指します。数学的な科学から判断すれば、最も完璧な方法の例として知られているものから見ても、この指示的機能が方法学の主要な役割と特徴を構成していることが分かります。例えば、算術の基本的な過程は自己で独立した可能性のある科学の要素に過ぎません。それらの性質を定義するのはその方法学の最初の仕事ですが、それらの順序を決定するための次なるそしてより高い役割はその方法学の後の仕事です。論理的推論がより複雑な例で、特に確率理論における難問の解決に基礎となるものにおいては、指示的な方法学、つまり単独の計算法しか提供できないようなものが必要となります。

10. したがって、論理の最終的な法則はその形式において数学的であり、除く一点を除き、一般的な数の法則と同一であるという理由、そしてその一点で異なる理由は、これらの問いについては、私たちの限られた能力を超えたものである可能性が高いから、非常に傲慢ではないかと思われる判断を試みる価値があるかもしれない。おそらくそれらは私たちの能力を超えるものである。

心が自分自身に適用される法則について知識を得ることは可能かもしれませんが、その基礎と起源を理解するための知識や、他の可能な法則システムと比較してその目的に対する適切性を完全に理解するための知識は、限られた程度しか得られないかもしれない。このような知識は科学の目的には必要でない。科学は現実のものについてのみに関心を持ち、優先度の根拠や任命理由を求めることはしないからである。

これらの考慮事項は論理を計算の形で提示することに対するすべての抗議に対する十分な回答を提供する。私たちはそれをそのような表現方法を選択したわけではないが、思考の最終的な法則がそのような表現方法を可能にし、その特性を規定し、他のどの形式でも科学の完全な表現が不可能であるかのように見える限り、そのような形態は採用されるべきである。また、科学の仕事は法則を作り出すことではなく、それらを見つけることにあり、私たち自身の心の構成を創出することではない。私たちには自分たちの心の性格を大幅に変更する能力があるが、その構成は私たちの意志によらない。

人間の知性の法則は私たちの意志によらず、それらが基礎となる科学の形態も、根本的な点では個々の選択とは独立している。

11. 以上の方法の原理についての一般的な説明に加えて、この論文ではその方法が分析において多大な種類の命題や推論の前提を構成する一連の命題に対する適用例も示す。これらの例は様々な著者のものから選ばれ、非常に異なる複雑さを持ち、また倫理学や一般哲学に関する問題にまで及ぶ幅広い範囲の主題を網羅している。

これらの選択には、特定の面において過度な自由が与えられたように思われるかもしれませんが、その点についての詫びは必要ないと私は考える。論理学が科学として非常に広範な応用が可能であることは認められるべきであり、同時にその最終的な形態と過程は数学的であることも確実である。したがって、倫理学や一般哲学に関する問題を議論する際にこのような形態と過程の採用に対する先入観的な反対は誤解や不適切な類推に基づくものでなければならない。

数学者にとっては数や量の概念にかかわる必要性は本質的ではない。倫理学の議論においてシンボル的过程を一般化する習慣が望ましいかどうかは別の問題である。おそらく私は他處で述べた通り、論理学の方法の完璧さはその原理の推測的な真実性を証明するための主な価値を持つかもしれない。

一般的な論理の使用を代替したり、技術的形式に従属させたりすることを望むことは、その知識が精神的労働と戦いから得られる筋肉的な活力を与え、困難に立ち向かう能力を教えるという点で、その価値を理解している人にとって最後の願望ではない。しかし、科学的手法の価値は、普通の論理の支配下にあるものすらも含む場合でも感じられ、認められることがあるかもしれない。

この種の例は本において見つかるだろう。

12. 以上で説明した論理の一般原理と方法は、確率理論およびそれに応じた方法の基礎でもあります。従って、これらの理論と方法を上記の原則に基づいて開発することは、本 Treatise の独立した対象となります。この適用の性質について、特にその導き出される解法の性格に関しては、ここで概説することが望ましいかもしれません。

論理に先行する方法が必要な理由として、確率理論の基礎となるためには、数少ない言葉で述べることができます。特定の事象が予想される頻度が、他の既知の事象の頻度にどのように依存しているかを決定する前に、各事象間の相互依存性について理解しておく必要があります。技術的に言えば、求められる確率を持つ事象を、与えられた確率を持つ事象の関数として表現できる能力が必要です。この明示的な決定は、論理の領域に属します。しかし、数学的には解釈される確率は数値測定が可能です。したがって、確率の主題は数と論理の両方の科学に属します。これらの要素の共存を認識することで、本 Treatise は以前のものとは異なります；この差異は問題の解決可能性について多くの場合に影響を与え、また得られる解法に新しい重要な要素を導入するため、私は以下のページで開発された理論の特有な結果について詳細に述べる必要があると判断します。

13. 事件の確率の測度は通常、分数として定義され、その分子が該事件に有利な状況の数を表し、分母が有利と不利の全体の状況の数を表します。すべての状況は同様の可能性で起こると仮定されます。この定義は本著作中において採用されています。同時に、主題の別の側面（直ちに述べられるもの）についても説明し、これは基本的な観点として等しく考慮されるべきであり、実際には同じ方法と結論を導くシステムにつながります。さらに追加しておくと、確率理論において受け入れられている結論の範囲において、そしてその結論が基本的定義の結果である限りでは、本著作法による推論（等しく正確とされるもの）の結果とは異nessはありません。

また、確率論における問題は様々な形で現れ、代数的な条件その他の要因によって多様に変更されるかもしれませんが、それらのすべての問題、またはそれぞれの問題が確率論に真正に属する部分については、一つの一般的な型があるように思われます。データと求めるべき事項（qu\ae{}situm）を考慮すると、その型は次のように説明できます：1つ目には、データは一つ以上の与えられた事象の確率であり、それぞれの確率はその事象に対する絶対的な実現確率であるか、または特定の仮定された条件のもとでの実現確率です。2つ目には、求めるべき事項（qu\ae{}situm）は、データで与えられた事象とは異なる表現を持つ他の事象の絶対的または条件付きの実現確率であり、それらがデータの要素をより或は少しだけ含むものである可能性があります。データに関して言えば、それは因果的に与えられるものであり、例えば特定のサイコロの投げ方の確率はその構造に関する知識から導き出される場合がある一方で、それは事象の成功または失敗を繰り返し観察した結果から派生するものです。後者の場合においては、事象の確率は観察された有利なケースと全体の数の比が無作為性の均一性という前提のもとで無限に継続されるにつれて近づく極限として定義されます。最後に、問題に関連する事象の性質や関係についても重要な区別があります。これらの事象は単純なものまたは複合的なものです。複合的な事象とは、その表現が他の事象の表現や概念に依存し、それらはその観点からは単純な事象と見做される事象であるような事象を指します。「雨が降る」と言ったり、「雷が鳴る」と言ったりすることは単純な事象の発生を表すことです。しかし、「雨が降りて雷も鳴る」または「雨が降るか雷が鳴るか」を言うことは複合的な事象の表現を示します。その事象の表現は「雨が降る」と「雷が鳴る」という基本的な表現に依存しています。したがって、単純な事象の基準は、それらの性質における仮定された単純さではなく、言語での表現または思考での概念化の方法に基づくものです。

14. 既存の確率理論は、以下の一般問題を解決するのに役立つ：与えられた単純な事象の確率が与えられた場合、特定の組み合わせ方でこれらの単純な事象から構成された複合的な事象の確率を求めることである。また、求められる複合的な事象に与えられた条件が課せば、つまりその条件もまた与えられた単純な事象によって特定の方法で依存すれば、この問題も解決できる。

この一般問題以外にも、解法の原理が既知の特別な問題がある。因果関係に関する様々な質問は、原因が互いに排他的であるという仮定のもとで既知の手法により解決できるが、それ以外の場合には明らかにそうではない。さらに、科学の一般問題として見なされるべき問題について、与えられた事象の確率が与えられ、求められる事象の確率も同様に任意の表現と概念を持つ他の事象の確率を求めることである。この問いを述べる際に、与えられた確率を持つ事象と求められる確率を持つ事象が共通の要素を持っている必要はないと明記されていないのは、問題のデータが目的に十分かどうかを決定し、不足がある場合その部分を示すことが方法の役割であるからである。

この問題は、最も制限がない形で述べられた場合でも、本論文の方法によって解決できる。正確には、その理論的な解は完全に与えられており、実際的な解は純粋な数学的手法、つまり方程式の分解と解析に依存するだけである。

一般解の順序と性質は次のように説明できる。

15. まず最初に、論理計算の前駆的な方法により、求められる確率のイベントを与えられた確率のイベントの論理的関数として表現することが常に可能である。その結果は次のようになる：
$X$ を求められる確率のイベントとし、$A$, $B$, $C$, などとし、これらのイベントは単純なイベントまたは複合的なイベントであるとする。すると、イベント $X$ がイベント $A$, $B$, $C$, などのそれぞれに対して持つ全体的な関係は、数学者が「展開」と呼ぶ形で導き出される。最も一般的な場合においては、この展開は四つの異なるクラスの項から成る。第一のクラスでは、イベント $A$, $B$, $C$ の組み合わせを表し、これらの組み合わせは必ずしも $X$ が発生することを示すわけではないが、$X$ が発生するためには必ず伴うものである。第二のクラスでは、これらの組み合わせは $X$ が発生するためには必ず伴うものの一方で、必ずしも $X$ の発生を意味しないものである。第三のクラスでは、これらの組み合わせはイベント $X$ と関連して発生することは不可能だが、それ以外では不可能ではないものである。第四のクラスでは、これらの組み合わせはどのような状況においても発生することのないものである。

この論理分析の結果について述べることは、以下のように補足するだけでよい：その要素は、イベント $A$, $B$, $C$ に対する我々の知識によって影響を受ける期待値が、または単に影響を受けることができる唯一の要素である。一般的な推論はこの結論を検証することはできるが、一般的な推論は上記の説明から解き出す複雑な状況や事件から混乱を避けることは通常できない。この目標の達成は、私が提案した問題に対する完全な解決策への第一歩である。ここで重要なことに注意すべきは、この過程は論理的であり、つまり論理的な意味を持つシンボルによって進行し、最終的には論理的に解釈可能な等式を導き出すことである。この結果を「最終論理方程式」と呼ぶことにする。

過程の第二段階は注意深く検討する価値があります。
前段階から導き出される最終的な論理方程式から、問題の完全な解を含む一連の代数方程式が直接的に導き出されます。この移行がどのように行われるかについて、その方法については十分に説明する必要はありませんが、事象の確率を他の依存事象の確率の代数関数として表現する法則と、これらの事象間の論理的関係を表現する法則との間に存在する特定の関係性があります。この関係は、以前にも述べたように（I.4）、観察と反芻によって我々に明らかにされた形式法則の一連の一致とは同様に、仮説に基づくものではなく、事前的に確率の定義の基礎として仮定される場合でも、厳格な論証から受け入れられる既存の数値的定義へと必然的に導かれる結果となります。確率理論は、すでに述べた通り（I.12）、論理学と算術との両方に等しく密接に関連しています；そして、その結果に関わらず、我々が一方の科学から発展したか、または二つを結びつける相互関係に基づいて基礎づけられたものとしてそれを捉えるかは、どちらでも構いません。

16. 本方法によって導き出される一般解には、数学的に興味深いいくつかの状況があり、それらについて注意深く述べる価値があるかもしれません。

1つ目の点として、この方法はデータの数や性質に依存しないため、データが不十分であっても定まらない値を求めることを可能とします。そのような場合、最終的な解の表現には任意の定数係数を持つ項が含まれます。これらの項に対応する最終的な論理方程式（I. 15）の解釈は、これらの定数の値を決定するために必要な新たなデータは何であるかを教えてくれます。これにより、数値解が完全になります。このようなデータが得られない場合でも、定数に極限値 $0$ と $1$ を与えることで、求められる確率が属する範囲をさらに経験に関係なく決定することができます。

求められる確率が与えられた確率とは全く独立であるような事象の場合、その値に対する得られた極限は $0$ と $1$ となり、これは明らかにそのようにすべきです。定数の解釈は元の問題を再述するだけとなります。

2番目には、最終的な解の表現は常に特定の量の要素を含みます。この要素の値は代数方程式の解によって定められます。現在、その方程式が高次である場合、適切な根を選択するという困難が生じるかもしれません。確かに、与えられた要素と求められる要素が必要条件によって明確に制限されている場合には、選択肢は残りません。しかし、複雑な事例では、そのような条件を未補助の論理力で発見することは希望外れです。この目的のために必要な明確な方法が必要であり、これは「統計的条件の微積分学」と呼ばれるものではなく、最終的な論理方程式の使用に基づく前記の方法と同様に、特定の条件が数値データ要素間で満たされるべきことを明示的に定めます。1つには、問題が実在するため（つまり、可能性のある経験から派生する）必要条件として、データから導かれる一連の結論を抽出するために。2つ目には、求められる確率が直接観察によって同一の現象系から推定された場合に、その数値的な制限は実際に求められた確率の制限となります。これらの制限は明らかに求められる確率の実際の制限となります。

現在、上記の条件に従ってデータを考慮すると、私は調査したすべての場合において、必要な制約条件を満たす最終的な代数方程式の根が一つだけ存在することを見ることができます。これにより、あらゆる不確定要素は除去されます。さらに、新しい代数方程式に関する真理が偶然明らかになるようにも思われます。特に重要なのは、前記の議論に関連する特定の量の要素はデータにのみ依存し、問題の提案された解（qu\ae{}situm）には一切依存しないことです。したがって、各個別の問題に対する解は、共通のデータを持つ一連の問題系（その問題系は提案された問題の性質に関係なく）の結びつきを解き放つための鍵となります。特定のデータ系から一連の関連する結論を導く必要がある場合、この事実は重要です。さらに、個別の問題に対する解には、その中心的な基本的単位に依存することから派生した関係性が与えられ、これはしばしば一般方法の適用における特徴となります。

17. しかし、上記の考慮事項と類似の他のものにより、この論文の方法が確率理論における問題解決のために一貫した方法であるという主張は正当化されるものの、それからすべての場合において仮定的な根拠への帰属を免れる必要がないとは限らない。解法が完全に任意の定数係数によって影響を受けた項で構成されている場合があることが指摘されており、実際には全く不確定であることも可能である。この論文の方法をその範囲内の最も重要な問題に適用する場合でも、経験データのみを使用した場合、このような結果が得られる可能性がある。明確な解を得るためには、これらの状況では、よりまたは少々独立性を持つが正確な検証は不可能な仮定に帰属する必要がある。一般的には、このような仮定は経験の直接的な結果よりも論理的であるかわりに数値的であるという部分を含み、現象が発生する条件を規定するよりもその頻度を指定するものである。しかし、この点は重要ではない。仮定の性質はいかなるものであれ、それらはその後、実際のデータとして扱われるべきであり、明らかにそれは解法自体に多少の仮定的な性格を与える傾向がある。実際に使用されたデータの可能性のある源についてこの理解を前提とすれば、方法は完全に一貫しており、しかし仮説要素の正確性については当然のように経験から得られた数値データの正確性ほど責任を持つことはない。

これらの言及を例示するために、天文観測の簡化理論について述べることができる。著者はこの主題については、別巻または将来的な付録において扱うつもりである。現在の論文では、積分微分法の使用を避ける。

この理論は部分的に仮定的な根拠に基づいている。それは誤差に関する特定の立場を仮定し、その発生形態として過大または欠如が等しい確率を持つなど、それらなしでは衝突する観測から明確な結論を得ることは不可能である。しかし、このような立場を前提とすれば、調査の残りは確率理論の領域に厳密に収まる。

同様の観察が重要な問題にも適用できる。それは議会の多数決記録から個々のメンバーの正誤判断の平均確率を導き出す問題である。この論文の方法を純粋な数値データに適用すると、得られる解は上記で述べた完全に不確定なものとなる。

さらに、これらのデータ自身による不足をより明確に示すために、解に現れる任意定数（I.16）の解釈は元の問題の再述に過ぎない。ただし、個々の心の中で意見が独立して形成されるという仮説を認めると、ラプラスやポワソンの推測のように絶対的に、または実際のデータによって制限される形で（本論文の第21章参照）では、問題はより明確な性格を持つ。

理論の後期的な価値は、純粋に実験的データが完全な解決を提供しない場合、そして最終的な論理方程式の解釈によって示されるさらなる経験が得られない場合、そのような中間的な仮説の正確な形成に大きく依存すること明確になるだろう。一方で、人間の知識を超えた性質を持つ主題において過度な仮説形成の準備は、確率理論の信用を損なう可能性があり、その最も正当な結論に対する一般大衆の疑念を増幅する傾向がある。

18. ここでは、抽象科学の方法が将来にわたり社会問題の調査においてどれほど有用となるかを推測するのはまだ早いかもしれません。物理学的な研究分野におけるそれらの方法が提供したサービスと比較して、社会科学におけるその役割がどの程度になるかについて議論することは非常に危険です。純粋な先験的論理に基づいてこの問いを解決しようとすると、我々を誤導する可能性があります。

例えば、人間の自由意志を考えると、最初に思い浮かぶのは社会システムの動きが物理的な必然性の支配下で秩序ある進化を示すことを期待することです。しかし既に、統計学者の研究によって、このような予想と異なる事実が明らかになっています。犯罪や貧困に関する記録は、人間の欲求や情熱の影響を受けない地域では見られないような一定の規則性を示しています。

一方で、季節の不調和、火山の噴火、植物界での病害や動物界での流行性疾患など、自然原因によるものとみなされるか、または主にその結果であることが明らかになっていますが、これらは定型化され、理解できる法則には従いません。「風のように変わりやすい」という諺があります。これらの点を考慮すると、我々の初期の判断を多少修正する教訓を得ることができます。

我々は、必然性の支配下においても、生成要因の作用が十分単純であれば、人間の観察によって認識できる秩序があることを期待すべきであり、逆に個々人の自由意志がシステム全体の動きにおける規則性と不一致であるとは考えないべきであることを学びます。人間の自由は意識的な事実として明らかであり、また、その部分の心の科学的構造を調査できる我々の能力から類推する可能性が高い結論でもあります（第22章）。

しかし、それが意識に基づく事実としても受け入れられ、または理由によって正当化される結論としても、それは観察結果と矛盾しないように解釈されなければなりません。つまり、多くの人々に関与する現象が非常に高い程度の規則性を示し、その状態や進行についての評価のために次々と収集すべき要素を集めるためには、それが外在的な結果として表れる範囲で明らかになるべきです。

したがって、事前に科学的に否定される理由は存在しません。また、この論文が他の目的を達成するにしても、社会統計学の実験的基礎となるデータ種類の存在について疑問を持つことはないと予想されます。そのような原理と方法に基づく抽象原則によって、任意の集団の社会データから明示的な形で何らかの情報を抽出することが可能であるという点については、理論的にはすべての場合にあり、実際的に必要な計算労働が提供される限りではあります。

ただし、データが極めて複雑な場合、その情報を得るためには非常に大きな困難があるかもしれません。これは理論の欠陥によるものではなく、それを達成するための分析過程の労働集約性によるものです。多くの場合、その困難は協力的な努力によって克服される可能性があります。

しかし、我々は理論的にはすべての場合に、必要な計算労働が提供されれば実際的にも、統計記録から埋もれた普遍的真理の種子を抽出するための手段を持っていると断言できるでしょう。

19. しかし、これらの一般的位置を超えて断言するのには至らない。科学的方法を統計記録に適用した結果が期待されるものであっても、それらの発見に必要なような援助なしでは得られないものよりもさらに価値があるかどうか、その労働が相当な規模での調査を Worthwhile にするかどうかは、おそらく経験によってのみ決定される点である。科学の抽象的な原理が知性的な喜びを超えて何かを提供することを望むべきであり、また、他の例と同様に、知識のより人間の幸福に関連性が高い分野からの類似した援助が得られる日が訪れることが非常に不確実ではないと思われる。しかし、この論文の推測は、その真実性として認識される権利を持っているという点において現在に留めておくべきである。

最後に、前章の科学的結果から人間の心の性質と構造に関する一般的な示唆を引き出すことを企図します。ここではこの種の推論の可能性に関する根拠について詳細には触れませんが、一二の一般的な観察は、私が追求しようとする道筋を示すのに十分でしょう。

私たちの論理学に対する見解は、おそらく主に決定するほど、知的能力の性質に関する意見にも影響を与えることは明白です。例えば、論理が単に最初のまたは必然的な真理の適用であるか、あるいは心そのものが法の座であり、個別の公式と全体の公式において同様に明確かつ結論的な作用を発揮する法則の運営であるか、それとも、いくつかの重要な著述家が主張するように、すべての論理は具体的な事実に基づくものであるかという問題があります。この問いは、単に論理学の科学だけに関連するだけでなく、知的能力の構造に関する公正な見解の形成にも関わるのです。

また、心が本来的に法の座であると結論付けられた場合、その法への服従の性質についての問題、例えばそれは天体の回転を維持し自然の秩序を保つような必然的なものなのか、それとも全く異なる種類の服従なのか、これは深遠な推測的興味を持つ問いです。さらに、心が真正に法の対象であることが確認され、その法も真正に特定された場合、その法則が人類の思考の流れに対し、過去と現在における可能性または必要性による影響を及ぼすか否かという問題は、哲学と歴史の両方にとって重要な問いであり、慎重な調査を価値あるものとするでしょう。

これらの問い他にもいくつかありますが、不完全ながらも本著作の結論的な部分で議論することを提案します。これらは、おそらく確実性の知識よりも類推的または推測的な知識の領域に属するものです。しかし、科学の証拠が不足していても、類推に基づく極めて確かな意見を提示するための十分な根拠があるかもしれません。

私はこの議論を本著作の主な業務の後半部分に完全に残しておく方が良いと考えました。その主な業務は科学的真理と法則の調査です。経験から学ぶ限り、真実に関するすべての探求における適切な進行順序は既知から未知へ進むことです。人間の心の哲学と構造の一部も、私たちの調査の範囲内に完全に置かれています。それらの部分を我々の性質の基礎として、その超然たる推測の領域を超えて影や不確実性の中に潜む部分を探求するための適切な認識を得ることこそが、私たちの現在の状況における制限と最も調和した道です。

\newpage

\chapter{記号全般および論理学に特有の記号について（Signs and their Laws）}
\label{ch:ii}

\begin{center}
\large\textit{名辞記号・結合記号（＋, －）・積記号・交換法則・分配法則とべき等律 $x^2 = x$}
\end{center}
\vspace{4mm}

\#\# 一般の記号とその法則：論理学に適した記号とその法則について

1. 言語は人間の理性の道具であり、単なる思想の表現の媒介であるだけではないという事実は、広く認められている真理である。本章では、言語が最も重要な知性の能力を支えるためにどのように役立つかを調査することを考える。この調査の各段階を通じて、言語の構成について、目的や目標に適したシステムとしての言語の構造を考慮する必要がある。その要素を探求し、相互の関係と依存性を定める方法を探求し、これらの要素がシステムの一環として目指す終着点への達成にどのように寄与するかを調査することになる。

これらの調査に進む際には、言語が論理のための重要な道具であるか、それともその助けなしで論理を行うことが可能であるかという有名な学校の問題について議論する必要はないと考えられます。私はこの質問を現在の著作の目的から外すことにします。理由は次の通りです：科学の仕事は法則を探求することであり、言語が事物とその関係の代表であるか、あるいは人間の知性の概念や操作の代表であるかに関わらず、記号の法則を研究することは実質的に論理の現われた法則を研究することに等しいのです。もし両者の調査間に差異があるとすれば、それは現在この作業の段階で探求すべき形式法則の科学的な表現に影響を与えるものではなく、その結果が精神的注意によってどのように提示されるかに関わる点だけです。

実際、記号の法則を後天的に調査する際は、直接検討するのは言語とその使用に関する規則であり、思考の内部プロセスを直接的に研究する際にも私たちはより直感的な形で個人的な意識に訴えることになりますが、両者の結果は形式的には等価であることが確認されます。また、地球の無数の言語や方言が長い歴史を通じて共通点と普遍性を保持しているのは、心の法則そのものに深い基礎があることを確信せざるを得ないからです。

2. すべての言語が構成される要素は、記号または象徴である。単語は記号である。しばしばそれは物事を表すと述べられ、しばしばそれは心が事物の単純な概念を組み合わせて複雑な概念を作るための操作を表すと述べられる。また、しばしばそれは我々の経験上の対象間で存在する行動、情熱、または単なる性質の関係を表現し、時には観察者の心の感情を表すと述べられる。しかし、単語は、この他にも記号や代表的な象徴としての役割を果たすために必要な唯一のものではない。任意のマーク、つまり目でしか理解できないマーク、または他の感覚に直接作用する任意の音声や動作も、その代表的な役割が定義され、理解されている限り、記号の性質を持つ。数学的科学においては、文字と符号 $+$, $-$, $=$ などは記号として使用されるが、後者の類の符号は操作または関係を表すために使われることが多く、前者は数や量の要素を表すために使われるため、「符号」という用語はしばしば後者に適用される。実際の意味を持つ記号の特定の形態や表現に依存しないように、その使用を決定する法則もまた依存しない。現在の論文においては、書かれた記号のみに関心があり、この点についてのみ「符号」という用語が使用される。記号の基本的な性質は次の定義で列挙されている。

\textit{定義}---記号とは、固定された解釈を持つ任意のマークであり、他の記号との組み合わせを固定された法則に従って行うことができるものである。これらの相互の解釈によって依存する法則に基づいて。

3. 以上の定義における個々の要素について考慮してみよう。

まず第一に、記号は任意のマークである。特定の単語やトークンが与えられた概念と結びついているか否かは明らかに重要ではなく、一度結びつけられた関連性が永久的に保たれる限り、どのような単語やトークンを選んでもよい。ローマ人は「civitas」という言葉で我々が「state」と呼ぶものを表現したが、彼らもまた我々も、同じ概念を表すために他のどんな言葉を使用しても全く問題ない。

実際、言語の性質上、単に一つの文字だけを使用して同じ意味を持つことも可能である。これが行われた場合、その文字を使用する法則は、ローマ語における「civitas」や英語における「state」を用いる法則と、一般的な言語全体に共通するどのような原理によっても regulare される範囲においては本質的に同じである。

二つ目の点では、同一の論述または推理過程において各記号は固定された解釈を持つべきであることが必要です。この条件の必要性は明白であり、その基礎は主題の本質にあります。しかし、議論の過程で使用される言葉や符号が名前として用いられる代表的な役割については、議論があります。一部の人々は、これらの記号が心の概念だけを表すと主張し、他方の人々はこれらが事物を表すと主張します。この問題はここでは重要な点ではないと考えますが、その解決は記号の使用法に影響を与えることはありません。しかし私は、議論の過程において記号は心の概念や操作を代用し、これらの概念や操作が事物を表し、事物との関係性を表現するとすれば、記号も同様に事物とその関係性を表現するという一般的な見解があると考えます。最後には、記号は心の概念や操作を代用することにより、それらの概念や操作の法則に従うべきであるという点です。この視点は次の章でより詳しく説明されるが、ここでは定義に含まれる第三の事項、すなわち解釈の性質に依存する固定された組み合わせ法則への従属性を説明するために役立つ。

4. 証明と分類の分析は、論理の道具としての言語の操作を次のように行うシステムの記号によって行われます。以下にその定理を示します：

\begin{center}
\textsc{定理 I.}
\end{center}

\textit{論理の道具としての言語のすべての操作は、以下の要素からなる記号システムにより行われることができます：}

1つ目、\textit{文字符号 $x$, $y$ などは概念の対象を表す}。

2つ目、\textit{演算符号 $+$, $-$, $\times$ は事物の概念を組み合わせたり分解したりして新しい概念を形成する心の操作を表す}。

3つ目、\textit{同一性符号 $=$}。

これらの論理の記号は、その使用において、部分的には代数の対応する記号の法則と一致し、部分的には異なる法則を持つ特定の法則に従います。

真理性を有する論理的議論の要素として、最も単純な形で組み合わせられ、最も単純な法則によって生成されるものが存在することを基準とするべきです。そして、この原則に基づき、以下の分類を考慮します。

\begin{center}
\textsc{第1類}.
\end{center}

5. \textit{呼称または説明的な記号、何かの名前を表すものや、そのものに属する性質や状況を表すもの}。

この分類には明らかに、固有名詞または共通名詞と形容詞が含まれます。これらは実際には、前者が参照する個々の物や複数の存在の具体的な存在を表現し、後者がその存在を含意しているという点で異なると考えることができます。形容詞に普遍的に理解される主語「存在」または「もの」を付加すれば、実質的な論理的推論の目的では实質名词可以替代之。无论从心理上的每一个细节来看，说“水是流体”与说“水是流体”的意思是否完全相同并不重要；至少在表达推理过程时两者等价。

明らかに、上記の分類には、いくつかの状況や関係を表すために多くの符号を使用することを含む詳細な説明が求められるような符号も含まれる。詩的な表現語彙はしばしばこのような種類である。これらは通常複数の形容詞から成り立つ複合形容詞であり、一つの形容詞が複数の単語で表す説明を果たすことが多い。ホメロスの「深渦流する海洋」は、「βαθυτειρωμένη θάλασσα」（bathutērōménē thalassa）という単語で実質的に表現されている。

また、想像や知性に向けられたどのような説明も、一つの符号によって等しく表されることがあり、その使用は形容詞「良い」または「偉大な」というものと同じ法則に従う。主辞「事物」と共に組み合わされると、このような符号は実質的に名詞となる。そして、一つの名詞で物事とその性質の両方の意味を表現することが可能となる。

6. 既に定義されているように、符号は任意の記号であるため、上記で説明したすべての種類の符号を文字で置き換えることが許されます。特定の名前や説明が適用される個体のクラスを表すために、単一の文字 $x$ で表すことにしましょう。例えば、「人間」の名前が与えられた場合、$x$ は「すべての人間」または「人間」というクラスを表します。「クラス」とは通常、特定の名前や説明が適用される個体を集めた集合を指しますが、本書ではこの用語の意味を拡張し、「ある特定の名前や説明に一致する個体が存在しない」または「すべての個体が存在する」という場合も含めます。したがって、「何もなし」と「宇宙」はそれぞれ「存在しないもの」および「すべての存在」を表すと解釈されます。「形容詞」、例えば「善」が説明として使用された場合、$y$ という文字で「善」に適用されるすべてのものを表します。つまり、「すべての善なもの」または「善なもの」というクラスを表します。さらに、$xy$ の組み合わせは、$x$ と $y$ によってそれぞれ表現される名前や説明が同時に適用される個体のクラスを表すことにしましょう。例えば、$x$ が「白いもの」を表し、$y$ が「羊」を表している場合、$xy$ は「白い羊」というクラスを表します。また、$z$ が「角のあるもの」を表し、$x$ と $y$ の解釈が保持されている場合、$zxy$ は「角のある白い羊」となります。つまり、「羊」と「白い」および「角のある」の名前や説明が同時に適用されるすべての個体を表します。

次に、上記のように使用される符号 $x$, $y$, などについての法則を考えましょう。

7. 適切な組み合わせによれば、二つの符号の順序は関係ありません。$xy$ と $yx$ の両者は、それぞれのメンバーに対して $x$ および $y$ の名前や説明が適用される個体のクラスを表します。したがって、以下の式が成り立ちます。
\[
xy = yx.
\]

$x$が白いもの、$y$が羊を表す場合、この方程式のいずれかの項は「白い羊」のクラスを代表します。概念形成の順序には違いがあるかもしれませんが、その下に含む個々のものは変わりません。同様に、$x$が「河口」という意味で、$y$が「川」を表す場合、$xy$と$yx$はどちらも「川が河口である」または「河口が川である」というクラスを代表します。この場合の組み合わせは通常の言語では名詞同士の組み合わせであり、前例では名詞と形容詞の組み合わせでした。

さらに、$z$という第三のシンボルを導入し、「航行可能」に適用されるクラスを表すものとする。次のような表現が使用できます。
\[
zxy,\; zyx,\; xyz,\; \text{etc.}
\]
これらは「航行可能な川が河口である」というクラスを代表します。

一つの説明的な語彙が他の何かに暗黙のうちに参照している場合、その参照を明示的に定義することで、上記の議論は依然として有効です。例えば、$x$が「賢い」を表し、$y$が「相談役」を表す場合、$x$が絶対的な賢さを意味するのか、または apenas 相談役の賢さを意味するのかを定義しなければなりません。そのような定義は法則 $xy = yx$ を維持します。

したがって、私たちはシンボル $x$, $y$, $z$, などを使って名詞、形容詞、説明的な文句の代わりに使用することが許されます。解釈規則として、これらのシンボルが一緒に書かれた任意の表現は、それらの意味が適用されるすべての対象や個々人を代表し、シンボルの順序は関係ありませんという法則があります。

解釈の規則が充分に例示されたため、シンボルを使用した形容詞の解釈を定義する際に常に「事物」を明記することを必要とするわけではありません。私が「$x$ を善と代表させよ」と言及した場合、「$x$ は他のシンボルによってその質的な対象が供給されるときのみ、善を表すものであり、単独で使用された場合には善な事物を表す解釈となる」と理解されることになります。

8. 上記の法則について、以下に追加するいくつかの観察があります。これらの観察は、今後導き出される他の法則にも多少適応的です。

まず、この法則は思考の法則であり、厳密には事物の法則ではありません。対因果関係を含まないもの質や属性の順序の違いは、単なる概念上の違いに過ぎません。法則（1）は一般的な真実として、同じものが異なる方法で概念化される可能性を表現し、その性質を述べています。そしてこれにはそれ以上のことはありません。

次に、思考の法則としては、これは言語の法則にも実際に展開されます。言語は思考の産物であり、その道具です。文章の傾向は一貫性に向かうものですが、それでも意味が絶対的な形容詞を同一の主語に対して用いるときの順序は無関係です。しかし、詩的表現はこの法則的な自由を名詞にも拡張することで豊かな多様性を得ます。ミルトンの言葉は特にこのような多様性によって特徴付けられています。名詞がその形容詞によって限定される前に置かれるだけでなく、しばしばその中に配置されます。光への祈りの最初の数行には以下のような例があります：
\begin{verse}
``天の初生の子孫。''\\
``深淵の暗い水の世界。''\\
``不滅の明るい存在からの輝き。''
\end{verse}

これらの逆順表現は単なる詩的な自由使用の結果ではなく、思考の深い法則によって承認された自然な表現です。便利のために通常の言語使用ではこの自由が行使されないためです。

第三に、式（1）によって表される法則は、代数の符号$x$、$y$、$z$が互換的であるという点で特徴付けられることができる。ここで述べることにより断定されるのは、代数における基本法則として表現される乗算過程
\[
xy = yx,
\]
その自己と何か類似性があるわけではないが、代数的な計算と論理的な組合せを同じ方法で表す場合、それらの符号的表現は同一形式の法則に従うという点である。その従属性の証拠は、両者の状況において異なるものである。

9. 二つの符号$x$と$y$の組合せ$x y$が、$x$と$y$によって表される名前や質的特性が共通して適用できる全体を表現するとき、もし二つの符号が完全に同じ意味を持つ場合、その組合せは単独でそれぞれの符号が示すよりも多くを表現しない。このような場合においては
\[
xy = x
\]
となるべきである。しかし$y$が$x$と同じ意味を持つと仮定すると、上記の式では$y$を$x$に置き換えることができ、これにより
\[
xx = x
\]
を得る。

通常の代数において、組合せ$xx$はより簡略化して$x^2$と表される。ここでも同様の符号法を採用する：特定の一連の心的作業を表現する方法は、単一の概念や操作を表現する方法と同じく完全に任意なものである（第2章3節参照）。従って、この符号法に基づき上記の式は
\begin{equation}
x^2 = x,
\end{equation}
と表され、実際には名前、質的特性、または説明を符号的に表現するこれらの符号による第二の一般的法則の表現である。

読者には、例題で使用した符号 $x$ と $y$ が以前に異なる意味を持っていたとしても、これらに完全に同じ意味を付与すること自体は問題ありません。実際の意味が互いに近づくほど、組み合わせ $xy$ によって示される対象のクラスが $x$ や $y$ によって示されるクラスに近づくことも明白です。等式（2）の証明で仮定された状況は「絶対的な」意味の同一性を想定しています。その法則は言語において実際的に例示されています。例えば、何かに対する「良い、良い」という表現は冗長で無駄な重複表現ですが、「良い」と同じ意味になります。「良い、良い」人間は「良い」人間と同じです。このような単語の繰り返しはしばしば、質を強調するためや主張を強化するために使用されます。しかし、その効果は二義的で慣習的なものであり、言葉と思考の内在的な関係に基づいていません。我々が観察したり自身で行う多くの自然現象の操作は繰り返しによって効果が増大するため、言語においても同じことを期待することが準備されます。そして強調を表すために繰り返しを使用することさえあります。しかし、厳密な論理や正確な議論においては、このような慣習には正当な根拠はありません。

10. 今度は、別の類の言語の記号とその使用に関する法則について考察します。

\begin{center}
\textsc {第2類}.
\end{center}

11. \textit{集合や分割という精神的操作を表す記号。}

我々は、個々の対象に適用される名前、性質、または状況によって特徴付けられる概念だけでなく、部分的なグループから成る複数のグループ全体の概念を形成することも可能である。これらの目的のために、「そして」、「または」などの接続詞を使用する。例えば、「木々と鉱物」、「不毛な山、または肥沃な谷」はこの種類の例である。厳密には、二つ以上の対象のクラスを特徴付ける語句間に「そして」、「または」という言葉が挿入された場合、それらのクラスは全く異なるものであり、一つのクラスから別のクラスへのメンバーは存在しないことを意味する。

この点と他のすべてにおいて、「そして」「または」の用法は代数における符号$+$に類似しており、その法則は同一である。したがって、「男と女」という表現は、慣習的な意味を考慮に入れない限り、「女と男」という表現と同じ意味を持つ。ここで$x$を「男」という概念で表し、$y$を「女」という概念で表すとする。さらに「そして」「または」をそれぞれ表すために$+$を使用する。すると以下の式が成り立つ：
\[
x + y = y + x
\]
この等式は$x$と$y$が数を表す場合でも、$+$が算術的な加法の符号を表す場合でも同様に真となる。

また、「ヨーロッパ」という形容詞を表すシンボル$z$を用いる。実際には、「ヨーロッパの男と女」を述べるのと同じ意味を持つ「ヨーロッパの男とヨーロッパの女」を述べることと等しいため、以下の式が成り立つ：
\[
z(x + y) = zx + zy
\]
この等式もまた、$x$、$y$、および$z$が数を表す場合でも、二つの文字シンボルの接続が代数的な積を表す場合でも同様に真となる。これは以前述べた論理的意味においても、両方の形容詞が結びついた対象のクラスを表すことを示している。

以上のものは、符号$+$の使用に関する法則です。この符号は部分を全体に集める正の操作を表しています。しかし、ある変化をもたらす操作という概念自体が、その逆または負の操作を考える暗示を与えます。つまり、部分を集めて全体にすることができると考えることは、それらの部分から全体を分離することも可能であると考えることと矛盾しません。この逆操作は「除く」という言葉で表され、「すべての人々（アジア人を除く）」や「すべての州（君主制ではない州を除く）」など、一般的な表現として使用されます。

ここで、除外されたものは除外される対象から一部であることが暗黙に含みます。部分を集める操作を$+$で表すなら、前述の逆操作は$-$で表すことができます。したがって、$x$を人間全体とし、$y$をアジア人（つまりアジアの人間）とするとき、「すべての人々（アジア人を除く）」という概念は$x - y$で表現されます。また、$x$を「州」とし、$y$を「君主制である」という属性とするとき、「すべての州（君主制ではない州を除く）」という概念は$x - xy$で表現されます。

論理上の重要な目的においては、除外されたケースを話す順序に関わらず最初に表すか最後に表すかは関係ありません。したがって、ある数詞の一部が符号$-$に影響されている場合でも、その項の順序も関係ありません。これは一般的な代数学と同様です。

\begin{equation}
x - y = -y + x.
\end{equation}

今、$x$を「人間」というクラス、$y$を「アジア人」で表すことにします。さらに、$z$を「白い」という形容詞で表します。このとき、「人間（アジア人を除く）」という集合に「白い」という形容詞を適用することは、「白い人間（白いアジア人を除く）」と同義です。したがって、

\begin{equation}
z(x - y) = zx - zy.
\end{equation}

これは普通の代数学の法則にも合致しています。

方程式（4）と（6）は、単一の一般法則を例示していると考えることができます。この法則は、「文字符号 $x$, $y$, $z$ などはその作用において配分的である」ということによって述べることができます。この法則が表現する一般的な事実は次の通りです：ある性質や状況がグループのすべてのメンバーに与えられた場合、そのグループは集合または部分グループの排除で形成されたものであってもかまいません。その結果得られる概念は、まず各部分グループのメンバーに性質や状況を与え、その後集合または排除を行う場合と同一です。全体のメンバーに与えられたものは、いかなる方法でもその部分がつながられているとしても、すべての部分のメンバーにも与えられます。

\centerline{第3級}

12. 「関係を表す記号および命題を形成するための記号」

すべての動詞は適切にこのクラスに属すると考えることができます。
しかし、論理学の目的には、主語動詞「is」または「are」としてのみこれを含むものだけを考慮することが十分です。なぜなら、他のすべての動詞はこの要素に分解され、クラスIの一つの象徴を含むことができるからです。これらの象徴がいかなる種類の質や状況も表現するために使用されるため、主語の活動的または受動的な関係を過去、現在、未来いずれかの時間に関して考慮する際にも使用することができます。

例えば、「Cæsar conquered the Gauls」という命題は「Cæsar is he who conquered the Gauls」に分解できます。この分析の基礎として考えられる理由は次のとおりです：「conquered the Gauls」の意味を理解しない限り、すなわち「One who conquered the Gauls」という表現の意味を理解しない限り、「Cæsar conquered the Gauls」の文を理解することはできません。したがって、「is」または「are」という要素はその文の一部であると言えます。また、「Cæsar」も要素であり、さらに必要となる要素として接続詞「is」が必要で、これら二つの要素の関連性を示すためにです。

ただし、この分析が採用された特定の視点に対する論理的推論の目的には十分であると主張する一方で、「Cæsar conquered the Gauls」という命題によって表現される関係について考える別の方法は存在しないとは断言していません。ただ、ここでの分析はその特定の視点に対する正解であり、論理的な推論のために十分であると言えるのです。

また、ギリシャ語の過去分詞と将来分詞が主語動詞に「is」または「are」という要素を含む存在を示唆していることも述べることができます。すなわち、「is」または「are」はすべての個人的動詞の要素として見なされることが可能であるという原理が存在することです。

13. 以上の象徴「is」または「are」はシンボル $=$ で表すことができます。このシンボルによって導入される法則、つまり通常は公理と呼ばれるものについては、次に考慮する必要があります。

次に、命題「星は太陽と惑星である」を考えます。星を \(x\)、太陽を \(y\)、惑星を \(z\) で表すことにしましょう。すると、次の式が得られます：
\[
x = y + z.
\]

もし「星は太陽と惑星である」という命題が真なら、星のうち惑星以外は太陽であることが導かれるはずです。これにより以下の等式が得られます：
\[
x - z = y,
\]
これは (7) から必然的に導かれるべき式です。したがって、一方の項 \(z\) を他方へ移動させ、その符号を変えることで等式の項が移動しました。これは代数における移項則と一致しています。

しかし、特定の例にこだわるよりも、一般の公理を即座に述べることができます：

1. 等しいもの同士に等しいものを加えると、その和も等しい。
2. 等しいもの同士から等しいものを引くと、その差も等しい。

したがって、等式は加算や減算の対象となり、上記の移項則を用いて通常の代数と同じように操作することができます。

また、二つの類 \(x\) と \(y\) が同一である場合、すなわち一つの類のすべての要素がもう一つの類の要素であるとき、一つの類のどの要素も特定の性質 \(z\) を持つものと、もう一つの類のどの要素も同じ性質 \(z\) を持つものは同一である。したがって、等式
\[
x = y
\]
が成り立つ場合、どのような類または性質 \(z\) に対しても
\[
zx = zy
\]
が成り立ちます。これは代数の法則と形式的に同じです：等式的両辺に同一の量を乗じると、その積も等しい。

同様にして、二つの等式の対応する項を乗じた結果の等式も真であることが示されます。

14. 現在の体系が代数学における一般的な説明と類似している点はここまでに見られますが、その類似性はここで止まるように思われます。あるクラス \( x \) のメンバーが特定の属性 \( z \) を持つものが、同じ属性 \( z \) を持つクラス \( y \) のメンバーと一致すると仮定しても、クラス \( x \) のすべてのメンバーがクラス \( y \) のすべてのメンバーと一致するとは限らないのです。したがって、等式
\[
zx = zy,
\]
から等式
\[
x = y
\]
も真であるという結論を導くことはできません。言い換えれば、代数学家の公理として知られる方程式の両辺を同じ量で割ることができるという公理には、ここでは形式的な相当物が存在しないのです。ここで「形式的な相当物」という言葉を使用したのは、これらの調査全体の一般的な精神に従って、論理符号 \( z \) を組み合わせ \( zx \) から除去するという心的作業が算術における除法と類似しているかどうかを定性的に決定することすら求められていないためです。その心的作業は実際には一般的に「抽象化」と呼ばれるものであり、その後の議論でその法則が既に導かれた法則に依存することが明らかになります。現在示されたことは、その法則の中に形式的に公理として広く受け入れられている代数学の公理と類似するものは存在しないということです。

しかし、少し考えれば、一般的な代数学においてもその公理は他の既に考慮した公理ほど普遍的ではないことがわかります。等式 \( zx = zy \) から等式 \( x = y \) を導くことは、\( z \neq 0 \) であることが分かっている場合のみ有効です。したがって、代数学の体系において \( z = 0 \) の値を許容する場合、上記の公理は適用されなくなり、以前示された類似性は少なくとも途切れないまま残ります。

15. しかし、量の記号について、そのような類似性を追跡することは、主として推測的な観点から以外は重要な意味を持つものではありません。私たちは（II.9）で見たように、論理の記号は特別な法則に従います：
\[
x^2 = x.
\]
一方、数の記号にはただ二つあり、すなわち $0$ と $1$ であり、これらも同じ形式的な法則に従います。私たちは $0^2 = 0$ そして $1^2 = 1$ を知っていますし、方程式 $x^2 = x$（代数的に考えると）はその他の根を持ちません。したがって、論理の記号と一般的な数の記号との形式的な一致度を定めるよりも、これらは値がただ $0$ または $1$ である量の記号と比較することが直接的です。そこで、符号 $x$, $y$, $z$ 等がそれぞれ $0$ または $1$ のいずれかの値を持つことを想定しましょう。これらの符号はただそれらの値を許すだけの代数を考えることで、論理の代数とその法則、公理、過程は完全に同一となります。解釈の違いのみが区別するものとなります。この原則に基づいて次の著作物の方法が確立されます。

16. 今後、この章の前節で考慮されていない普通の言語の構成要素について示す必要があります。それは、既に考慮された要素と同じ要素に分解できるか、もしくはそれらの要素をより精密な定義のために補完するものであることを示すことです。

实質、形容词和动词，连同连词“and”、“except”，我们已经考虑过了。代词可以被视为名词或形容词的一种特殊形式。副词修饰动词的意义，但不会改变其性質。介词有助于表达情况或关系，从而为文字符号的含义增添精确性和细节。连接词“if”、“either”、“or”主要用于命题之间的关系表达，以后将证明，同样的关系可以通过解释相似、形式和法則相同的元素符号完全表达，并且这些符号与本章中已解释过的使用及其意义相同。

任何剩余的语言要素，在检查后会发现它们要么用来使话语中的术语更加明确，从而進入我们已经考虑过文字符号的解释之中；要么用来表达伴随命题陈述的情感或状態，因此不属于我们目前关注的理解领域。通过其用法的经験可以证明所采用的分類是足够的。

\newpage

\chapter{思考の演算法則からの論理記号法則の演繹（Derivation of the Laws）}
\label{ch:iii}

\begin{center}
\large\textit{二値代数の公理系・排中律と無矛盾律・0と1の宇宙的解釈}
\end{center}
\vspace{4mm}

\#\#\# 第1章 [論理記号の法則の導出]

\#\#\#\# 大見出し：論理記号の法則は人間の心の操作の法則から導かれる

1. 科学の目的、厳密にはいうべきところでは、法則と関係を知ることである。この目的に必要なものと、偶然関連しているものを見分ける能力は、科学的な進歩にとって最も重要な条件の一つである。私は、「区別」することを言っている。すなわち、完全な科学への献身は、しばしば科学と関連している形而上的な考察から注意力を全く取り除くことを要求するわけではない。

論理記号の法則と人間の心の操作の法則との間の区別を明確にするためには、これらの要素を区別することが必要である。科学への一貫した献身は、他の形而上的な考察から注意力を全く取り除くことを要求するわけではない。このような問題、例えば現象の持続的な基礎の存在、原因の実在性、連続的な状態が「操作」によって結びついているという意味での表現形式の適切性など、しばしば形而上的なものと関連しているような問題も、科学に関心を持ち、その意義を持つことができる。しかし、これらは根本的には科学的ではない。

自然科学研究の結論を表現する際には、これらの概念の言葉を借りる必要があることは明白である。この原因から生じる実用的な不便さは必ずしもない。因果関係に比定される固定した順序だけが存在することよりも深いものがあると信じている人々と、そういった考えを否定している人々は、物理天文学の結論の解釈において一致する。しかし、彼らはその一致を認めるのは、因果律の性質について異なる見解を持っているにもかかわらず、科学的な真実の共通要素を認識しているからである。

2. この区別が物理学の科学において重要であるなら、知的力に関する科学との関連では、その重要性はさらに増す。この科学が提示する問題は、少なくとも表現においては混ざり合っており、それらは形而上学的な起源を示唆する思考や言語の模式と結びついてしまうからである。理想主義者は論理の法則に対して一つの表現を与えるだろう；反対に、真実性を保ちつつ疑問符を投げかける者には別の表現が適するであろう。我々がこの調査において関心を持つ現象を、単なる思考主体の順次的な状態と見なし、それらを活発な知性の操作に帰属させる者は、一貫性を持ちければ互いに異なる表現方法を持つことになる。また、精神的能力の分類が異なる場合にも同様の差異が生じるであろう。

この点において私は主張する原理は、このような多様性の中でも信頼と安定性を提供する唯一の基礎であり、次のようなものである：もしこれらの法則が観察から真正に導き出されたものであれば、それらは形而上学的な理論に関わらず人間の心の法則として実在的である。これらには、精神的操作の本質やその存在性についての更なる批判によって根本的に影響を与えることのない真実の要素が含まれている。

すなわち、心はただ意識状態の一連の流れであり、外からも内からも原因なく瞬時に現れ、また瞬時に消え去る一連の瞬間的な印象であると仮定されてもよい。疑問符を投げかける者にとって最後の洗練された考えであろうとしても、これらの法則は順序の法則や、少なくとも過去の順序の法則として観察によって導かれた結果は真実である。これらは「原因」と「効果」、「操作」と「主体」、「本質」と「属性」などのような言葉から排除された言語に解釈されなければならなくも、科学的な真理としては依然として有効である。

さらに、思想の法則を構成する声明が実際的な観察に基づくものであるとしても、そのような声明には、心の性質に関する形而上的理論とは独立した科学的要素が含まれていることから、これらの要素を論理体系や方法の構築に適用することは、また形而上的な区別とは独立して行われるべきです。なぜなら、思想の法則の声明に含まれる科学的要素は、そのような適用の基礎となるため、物理天文学の実用的な結論が重力の原因に関する理論ではなく、その現象的效果のみに基づくように、発見された思想の法則の実用的な利用も、形而上的な推測の真偽に関係なく、これらの科学的要素にのみ依存するからです。したがって、真正なる科学的目的においては、どのような形而上的な推測が真実であるか否かに関わらず、思想の法則を決定し、それらが発見された後にどのように利用するかについては心配することはありません。

3. この状況下において私に適切であると考えられる道は、可能な限り一般的な言語を用いて論理を行うことである。理論的な心の性質や能力に関する何か特定の観点には関与しないでいる。例えば、「会話」という言葉は一般的な慣習によれば、アイデアや概念の交換によって行われるコミュニケーションを指すものであり、そのコミュニケーションは言葉によって行われる職能であるとされる。また、心が特定のアイデアや概念に直面したときに、いくつかの職能や能力を持つことにより、あるアイデアに注意を集中させたり、あるいは与えられた概念やアイデアをさまざまな方法で組み合わせることもできるとされる。これらの職能や能力には、「注意」「単純な理解」「概念化または想像力」「抽象化」などといった異なる名前が付けられ、これらは人間の心の哲学における明確な分野のタイトルとなり、また人々の一般的な言葉にも取り入れられた。したがって、これらの用語を使用する機会があれば、それらを用いるが、それによって私があらかじめ心にそのような職能や能力があると受け取るわけではないことを意味しない。実際には、このような理解力の力を個別の存在として有しているかどうかについては調査する必要もないと考えられる。これらの異なるタイトルは、「人間の心の操作」という一括した一般的な名前で統合し、この作品の目的に必要な範囲でこれらの操作を定義し、それらの最終的な法則を表現することを目指す。このような順序で進行する道筋を採用し、ただし、私たちは時折、私たちが注意する特定の心の状態や操作に対する一般的な同意に基づいた名前を参照することがある。

次に調査の結果から得られるより明確な結論を分けて述べることを最も便利とする。

\begin{center}
\textsc{命題Ⅰ。}
\end{center}

4. 逻辑の記号の法則を、言語が論理的な推論の道具として厳密に使用される際に含意される心の操作から導き出す。

どの議論においても、それが自身の思考と対話する場合や、個々人が他人との交渉を行う場合でも、その操作の範囲は限定されているか明示的に述べられているかである。最も自由な議論は、使用される言葉が可能な限り広範に解釈され、その議論の限界が宇宙自身の限界と一致するときに行われるものである。しかし通常は、より狭い範囲に限定されることになる。しばしば人々について議論する際には（明示的に制限を述べることなく）、特定の状況や条件のもとでのみ私たちが話していることを含意し、文明人、または生命の盛りの男性、または他の状況や関係における男性を指すなどとする。そして、私たちの議論で出現するすべての対象が含まれる範囲はどれだけ広いかに関わらず、その範囲は議論の宇宙と呼ばれるべきである。

5. さらに、この議論の宇宙は厳密には最終的な主題とされるべきである。名前や説明的用語が仮定された制限のもとで使用される際の役割は、私たちの心にすべての存在や対象を想起させるものではなく、その中から指定または特定する個体を指すものだけである。もしこの議論の宇宙が実際の宇宙であるとき、つまり言葉が実在的かつ本来的な意味で使用されるときには、人とは「存在するすべての人」を指すが、もし議論の宇宙が先に理解された制限によって限定されている場合、それは特定の制限によって導入された状況や関係における人について話していることになる。言葉「人」というのは、適切な議論の宇宙から心の操作を行い、その中から指定または特定する個体を選択または固定させる役割を果たす。

6. 語形容詞の使用に含まれる精神的操作も、同一の種類です。例えば、論理的対象とする宇宙を実在の宇宙とします。このとき、「人」という語は私たちに、その宇宙から「人」に該当するすべての存在を選択するよう指示します；同様に、「良い」という形容詞が「良い人」という組み合わせにおいては、さらに「人」のクラスから「良い」という更なる性質を有する者だけを選択するよう指示します。そして、もし「良い人」という組み合わせの前に別の形容詞が付加された場合、それはまたその更なる性質に関連した同様の精神的操作を指示することになります。

ここに述べられている操作の実質的な性質を慎重に注意深く確認することが重要である。なぜなら、この操作はそのようなものでなくても考えられるからである。形容詞が単なる属性的役割を持っている場合、人間が「男」という名前によって特定されたある種の存在について話すとき、「良い」という形容詞を前置することで、我々はすべてのこれらの存在に善性という質を精神的に付与するように指示されるべきであると考えられるかもしれない。しかし、これは形容詞の実際的な役割ではない。

我々が本当に実行している操作は、特定の原則や概念に基づいて選択することである。この操作が知覚力や注意力などの心の能力にどのように帰属されるかについては、受け入れられたその能力の分類によると問うことは重要でないが、私はそれが想像力や注意力の二つの能力に依存すると考えられる。これらの一つの能力は一般的な概念の形成に関連され、もう一つの能力は指定された論理的宇宙内の個々の存在を特定し、そのような存在に対して精神的な注意を向けることに関連される。

しかし、注意力の力が他の心の能力の行使を維持する力に過ぎないという可能性を考えると、我々が上記で説明したすべての精神的操作は想像力や概念化の一つの能力に帰属すると考えられる。この過程の最初の一歩は宇宙自身の概念形成であり、各次のステップはそのような概念を限定的に形成する。

このような見解に基づいて、私は各該当のステップまたはそれらのステップの明確な組み合わせを「明確な概念化の行為」と説明する。この用語の意味を拡張し、それは名前や単純な質的属性によって表される物体のクラスの概念形成だけでなく、人間の能力と制限に従ってこれらの概念の組み合わせも含むように定義する。実際には、文または命題の構造に至るまでの任意の知的な操作を含む。

このような心の操作が従う一般的な法則について現在考慮する必要がある。

7. 今、前章で言語の構造から事後的に決定された論理的記号の使用に関する法則が、実際に先に説明した特定の精神的操作の法則であることが示される。

我々は議論を開始する際、その主題の限界についてある理解を持つことにしよう、すなわち、その宇宙の限界について。我々が用いる名前や記述語は、対象に対する一定の精神的操作を行うように相手に指示を与え、これによって思考が伝えられる。しかし、各々の名前や記述語はこの観点から見れば、その代表する知的作業の代理であり、それが自然の順序において先にあることから明らかであるように、名前や符号の法則は派生的なものでなければならない。実際には、それらが代表する操作の法則から発生しているのである。

現在では、記号と精神的操作の法則が表現において同一であることが示される。

8. では、私たちの論理の宇宙を実際の宇宙と仮定し、言葉はその意味の全範囲で使用されることとしましょう。そして、「白」および「人間」という言葉によって示される二つの精神的操作について考えましょう。「人間」という言葉は、その主題である宇宙からすべての人間を選択するという思考操作を示します。その結果得られる概念「人間」が次の操作の主題となります。

「白」という言葉によって示される操作は、「人間」から白いもの全てを選択することです。最終的な結果は「白い人間」になります。明らかに、上記の操作が逆順で行われた場合でも同じ結果が得られることがわかります。最初に「人間」という概念を作り出し、次にその概念を「白い人間」に限定するか、最初に「白い物体」という概念を作り出し、その後それを「人間」のうちの白いものに限定するかに関わらず、結果は同じです。言葉「白」および「人間」をどのような説明的または呼称的な語句でも置き換えた場合でも、その意味が固定かつ絶対的である限り、精神的操作の順序も等しく関係ありません。

したがって、概念の能力における二つの連続する行為の順序は同等であり、一つの行為がもう一つの行為の主題を提供することを前提とする場合、それはその能力の行使の一般条件となります。これは心の法則であり、それが論理の文法的な表現の実質的起源です（1）第2章。

9. 以上の精神的操作は、その効果が繰り返しによって変化しないという性質を持つことが明らかである。例えば、意識を人間に集中させる特定の概念的な行為があったと仮定し、その後で同じ能力を使って白人のみに限定する別の行為を行ったとする。その後、さらに後者の精神的操作を繰り返し行うことで、意識は白い対象に限定されるが、これは既に到達した結論、すなわち白人という概念には何ら影響を与えない。これはまた一般的な心の法則の一例であり、その形式的な表現は第2章（2）の文字符号の法則において見ることができる。

10. さらに明らかであることは、異なる二つのクラスの事物に対する概念から、それらのクラスを合わせた全体の概念を作り出すことができることである。そして、これらのクラスが精神的な視野にどのように提示されるかに関わらず、その順序や優先性は関係ない。これはまた一般的な心の法則であり、その表現は第2章（3）において見ることができる。

11. この調査と比較を続ける必要はない。十分な例が挙げられることで、以下の二つの立場が明らかになるからである：

まず第一に、想像力や概念化によって事物または性質の単純な概念を組み合わせたり修正したりする精神的操作、そして真実や命題に対する理由の操作は、一般的な法則に従っている。

第二に、これらの法則は数学的な形式を持ち、人間の言語の本質的な法則において実際に開発されている。したがって、論理の符号の法則は、推理における精神的操作から導き出されるものである。

本章の残りは、その思想法則についての問題に占められる。その表現は $x^2 = x$（II. 9）であり、これは第II章15節で暗示されたように、通常の言論と推理における心の作用と、一般量の代数における心の作用との区別を特徴付けるものである。次の調査の重要な部分は、シンボル $0$ および $1$ が論理のシンボル群において位置を占め、またその解釈を受け入れられる可能性があることを証明することである。そして最初に必要となるのは、上記のような特定のシンボルが適切かつ有利に異なる思想体系の表現のために使用される方法を示すことが求められる。

この適切性の基礎は、いかなる解釈上の共通性にも存在しない。論理と算術（後者の用語は数の科学として最も広い意味で使う）のような真に異なる思想体系においては、正確には主題上の共通性は存在しないからである。一方が事物そのものに対する概念を扱い、他方はそれらの数的関係のみを考慮するのである。しかし、いかなる推理体系の形式と方法も、シンボルが従う法則に直接依存し、解釈を通じて間接的に依存するからこそ、異なる思想体系において同じシンボルを使用することが適切かつ有利である可能性がある。ただし、そのような解釈を付与することで、形式的な法則が同一となり、使用が一貫性を持つことができる場合に限る。

このシンボルの使用の基礎は、解釈上の共通性ではなく、それぞれの体系における形式的な法則の共通性にある。そして、その形式的な法則の共通性を確立するためには、考慮対象となる各体系の解釈から独立して得られる結果を見ることによってのみ、慎重な観察と比較が必要である。

これらの観察は、以下の命題で採用される調査の過程を説明するものとなる。論理の文字符号は、その表現が$x^2 = x$である法則に従う普遍的な法則を持つ。数の符号には、この法則を満たすものとして0と1のみがある。しかし、これらの符号はそれぞれ、数的量の体系において独自の法則に従う。これにより、論理の体系においても同じ独特な形式的な法則が実現される解釈を与えなければならない問いが提起される。

\centerline{\textsc{命題II}}

13. \textit{論理的な値と意義を決定する符号0と1}

代数において用いられる符号0は、以下の形式的法則を満たす：
\begin{equation}
0 \times y = 0, \quad \textrm{または} \quad 0y = 0,
\end{equation}
どのような数$y$を表してもよい。この形式的な法則が論理の体系においても遵守されるようにするためには、符号0に解釈を与える必要があり、その解釈によって$0y$によって表されるクラスが常にNothingによって表されるクラスと同一であることを満たすものでなければならない。少々の考慮により明らかになるように、この条件は符号0がNothingを表す場合に満たされる。従って、前節での定義に従って、Nothingをクラスと呼ぶことができる。実際、NothingとUniverseはクラス拡張の二つの極限であり、これは一般名詞の可能な解釈の極限である。なぜなら、これらはどれもNothingより少ない個体を関連付けることができず、あるいはUniverseより多い個体を関連付けることはないからである。したがって、どのようなクラス$y$であれ、そのクラスとNothingというクラスとの共通の個体はNothingによって表されるクラスと同一であり、これらは全く存在しないからである。そして、符号0にNothingを与える解釈により、法則（1）は満たされ、これは完全に一般的なクラス$y$に対する法則の性格とは矛盾しない。

次に、数のシステムにおいてシンボル$1$は次の法則を満たします：
\[
1 \times y = y, \quad \text{または} \quad 1y = y,
\]
どのような数字$y$が表すものであれ、この形式方程式が本書のシステムにおいても同様に有効と仮定すると、シンボル$1$は任意の提案されたクラス$y$にあるすべての個体が、クラス$1$によって表されるクラスとクラス$y$との共通の個体$1y$であることを示します。ここで少し考えると、シンボル$1$が「宇宙」を表すことが明らかになります。なぜなら、「宇宙」は任意のクラスに存在するすべての個体を含む唯一のクラスだからです。したがって、論理のシステムにおけるシンボル$0$と$1$のそれぞれの解釈は「何もない」と「宇宙」であることが分かります。

14. どのような対象のクラス「人」について考えても、心には「人でない存在」のクラスという概念が浮かび上がります。そして、全宇宙はこれらの2つのクラスから成り立っています。なぜなら、その中に含まれる任意の個体に対して、「それが人である」とか「それが人ではない」といういずれか一方を主張できるからです。したがって、このような対照的な名前をどのように表現するかを調査することが重要となります。これが次の命題の目的です。

\textbf{命題 III.}

任意の対象のクラス$x$を表すとき、シンボル$1 - x$はその対照的または補完的なクラスを表します。つまり、クラス$x$に含まれていないすべての対象を含むクラスです。

より明確な概念のために、$x$ を「男性」のクラスとし、最後の命題に従って宇宙を $1$ で表すことにしましょう。現在、「男性」と「非男性」から構成される宇宙の概念から「男性」の概念を取り除くと、結果的に得られる概念は「非男性」のクラスとなります。したがって、「非男性」のクラスは $1 - x$ で表現されます。

一般に、シンボル $x$ によって表されるどんなオブジェクトのクラスでも、その対偶のクラスは $1 - x$ で表されます。

15. 次の命題は厳密にはこの作品の後の章、つまり「maxims」または「必然的真理」に関する主題に属しますが、その法則が思考における非常に重要な性質に関連しているため、ここに導入することが適切であると判断されました。

\textbf{命題 IV.}

哲学家たちによって「矛盾原理」と呼ばれる公理は、ある存在が同時にその特性を持つことができず、また持っていないこともできないという主張であり、これは思考の基本法則 $x^2 = x$ の結果です。

次のように方程式を書き換えてみましょう：
\[
x - x^2 = 0,
\]
これから
\setcounter{equation}{0}
\begin{equation}
x(1-x) = 0;
\end{equation}
これらの変形は、組合せと移動の公理法則（II. 13）によって正当化されます。単純さのために、符号 $x$ の意味を「人間」に特定的に解釈しましょう。すると $1 - x$ は「非人間」のクラスを表します（定理 III）。二つのクラスの表現の形式的な積は、それらの両方に共通する個体のクラスを表します（II. 6）。したがって $x(1-x)$ は「同時に人間であり非人間である個体」のクラスを表し、方程式 (1) は「同時に人間であり非人間である個体は存在しない」という原理を述べています。言い換えれば、「同じ個体が同時に人間でなく人間であることは不可能である」ということです。

次に、符号 $x$ の意味を任意の性質を持つ生物のクラスへと拡張します。すると方程式 (1) は「同一の存在が同時にその性質を持ち、持たないことは不可能である」と述べています。これはアリストテレスによって説明された哲学の基本的な公理、「同じ存在が同時にその性質を持つことと持たないことが両立しない」という「矛盾の原理」に等しいのです。

〔ギリシャ語原典引用〕---\textit{メタ物理学}, III, 3.

この訳では、アリストテレスの言葉を直接引用し、原文の意味を忠実に伝えています。

上記の解釈は、現在の体系における直接的な価値からではなく、知性の能力に関する哲学において重要な事実を例示するために導入されたものであり、その事実は、通常 metaphysics の基礎公理と見なされてきたものは、思想法則の結果であるという点が重要です。数学的形式を持つ思考法則の結果であることを強調したいと思います。

また、この基本的な思考法則を表す方程式（1）は二階の方程式であることに注意していただきたいと思います。

ここで $x^2 = x$ の存在が $x^3 = x$ という三階の方程式の存在も必要とすると言われた場合、その問いに対する答えは、論理体系において $x^3 = x$ は解釈できないということです。なぜなら、方程式を以下のいずれかの形で書き換えると、

\begin{eqnarray}
x(1-x)(1+x) &=& 0,\\
x(1-x)(-1-x) &=& 0,
\end{eqnarray}

その解釈が可能であるとしても、それは因子 $1 + x$ や因子 $-1 - x$ の解釈を含むことになります。前者は解釈できない理由は、宇宙 $1$ に任意のクラス $x$ を加えるという概念を思い描くことができないからです。後者は解釈できない理由は、符号 $-1$ が全てのクラス符号に適用される法則 $x(1-x) = 0$ に従わないからです。

したがって、方程式 $x^3 = x$ は $x^2 = x$ のような解釈とは異なった解釈を受けることはありません。しかし、もし前者の方程式が後者に関係なく真である場合、すなわち符号 $x$ によって表される心の活動が二度繰り返した結果が一度の操作の結果と一致するが、最初の一回の単なる繰り返しではそうではない場合、一つの形式（1）や（2）を解釈することは可能であると考えられます。しかし、実際の思考条件下ではそれを行うことができません。

数学において知られているある種の操作は、その法則を方程式 $x^3 = x$ で十分に表現することができますが、これらは一般論理の領域とは全く異なる性質を持っています。

「思想の法則が現在とは異なる可能性がある」という主張は、そのような仮説を立てることができ、そしてその結果を研究できるという意味だけを指しています。この仮説を立てる可能性は、人類の理性の実際的な法則が偶然や任意の意志の産物であるという教義を含むものではありません。

この章で思想の法則が必然的にそのような状況であるかどうかについて議論することなく、もしそれが存在しなければ、理解の手続き全体は現在とは異なるものであったという主張を敢えて行うことができます。つまり、思考の基本方程式が二次式であるという事実から導かれる結果として、私たちは対義語のペアに分けることで分析と分類を行うという操作を行います。技術的には「二分割」または「ダイコトミー」と呼ばれます。もし仮説の方程式が三次式であり、その解釈が可能であれば、精神的な分類は三つに分かれること必至で、私たちは「トリコトミー」による操作を行うことになります。私たちの現在の能力では、その実質を十分に理解することは不可能ですが、その法則はまだ知的推測の対象として調査することができるでしょう。

16. 方程式（1）によって表現される思想の法則は、上記の議論から明らかになる理由により、「二重性の法則」と時々呼ばれることがあります。

\newpage

\chapter{一次命題および二次命題への分類（Division of Propositions）}
\label{ch:iv}

\begin{center}
\large\textit{事物の関係を表す一次命題と命題間の真偽関係を表す二次命題の構造}
\end{center}
\vspace{4mm}

\#\# 分論文の分類、その特性と基本命題の表現法

1. 念頭に浮かぶ概念や想像に関連する心理的作業の法則が調査され、それらを表す対応する記号の法則も説明された後で、我々は獲得した結果の実用的な適用について考えるよう導かれる。まず第一に、命題中の複雑な項の表現；次に、命題自体の表現；最後に、推論的分析の一般的手法の構築について考えることにする。本章では、これらの目的の中で最初のもの、すなわち命題中基本命題の表現について主に扱うことにする。そのためには、以下の命題を確立することが必要である：

\begin{center}
\textsc{命題I.}
\end{center}

すべての論理的命題は、「基本的」または「具体的」命題と「派生的」または「抽象的」命題という二つの大きなクラスに属すると考えることができる。

我々が述べるすべての主張は、次の二つの種類のうちの一つに帰属することができる。それは、事物間の関係を表現することであるか、または論理文間の関係を表現すること又はそのような関係を表すことに等しいことである。事物の性質について、あるいはそれらが示す現象について、あるいは置かれている状況についての主張は、厳密には事物間の関係を表現する主張である。「雪は白い」という主張は論理の目的から見れば、「雪は白いものである」という主張と等価である。

事実や事件、それらの相互の連接と依存性についての主張は、同じ目的から見れば、そのような事実に関するいくつかの論理文が互いにどのように真偽に関係しているかを述べる主張と一般に等価である。前者の種類の主張、すなわち事物に関する主張、私は「一次的」と呼ぶことにする。後者の種類の主張、すなわち論理文に関する主張、私は「二次的」と呼ぶことにする。この区別は実際的には共通の論理的な区別、即ち命題が主観的または仮定的であるかどうかという区分にほぼ一致しているが、完全には一致していない。

例えば、「太陽が照っている」、「地球が暖まる」という主張は一次的であり、「もし太陽が照っていれば地球が暖まる」という主張は二次的である。「太陽が照っている」という主張は「太陽は光るものである」という主張と等価であり、これは二つの事物のクラス間の関係を表現する、すなわち「太陽」および「光るもの」の間に存在する関係を表している。しかし、上記の二次的主張は、「太陽が照っている」と「地球が暖まる」の二つの一次的主張間の依存性を表現している。

私はここでこれらの主張間の関係が、その表現する事実間の因果関係と同様であることを主張していない。しかし、これらの主張間の関係は、その事実間の関係を含み、またその逆も成り立つのであり、論理の目的から見れば、その関係がその事実間の関係の適切な代表として使用できるという意味で互いに関連していると言える。

2. 「太陽が照っている」という命題ではなく、「太陽が照っているという事実は真実である」と述べる場合、直接物事を指すのではなく、その物事を含む命題について話していることになります。つまり、「太陽が照っている」という命題自体について話しているのです。したがって、このような表現で述べた命題は二次的なものとなります。すべての一次的な命題は、その真偽を主張する二次的な命題や、その否定を主張する二次的な命題を生むことができます。

通常、接続詞「もし」、「あるいは」、「または」が二次的な命題を示すことを指し示すことがありますが、必ずそうであるわけではありません。例えば、「動物は理性を持つ者か非理性を持つ者がいる」という命題は一次的なものであり、それを「動物は理性を持つ者か動物は非理性を持つ者がいる」と分解することはできません。「それらの二つの命題が結びついた後述の排他的論理和文において」相互に依存する関係を表現していないのです。実際、「あるいは」や「または」は命題の性質を判断する基準ではなく、二次的な命題に頻繁に見られる傾向があるだけです。「もし」も一次的な命題に現れることがあります。例えば、「賢ければ節制的である」という命題があります。「すべての人々が賢ければ、すべての人々が節制的である」と分解することはできません。

3. 通常の命題分類について議論する目的はなく、ただここで述べることは、この分類法の基礎となる原理は明確で具体的に適用でき、実質的な区別を含み、一般的な推理方法の発展にとって重要な要素であるということです。また、一次的な命題が二次的な形に変換されても、ここでの主張とは必ずしも矛盾しないことを指摘します。一次的な命題において結びついた事実自体について直接言及するのではなく、「真偽」としての命題そのものについてのみ話しているからです。

4. 一次と二次の命題において、同一の符号を用いる。それは主語として同じ法則に従うことが明らかである。両者の違いは形式的な違いではなく解釈の違いである。両者において、命題が表現しようとする実際の関係は $=$ の記号によって表される。一次の命題において、このように結び付けられた要素は通常、命題の「項」を表すものや、より具体的にはその主語と述語として指定されるものである。

\centerline{\textsc{命題 II}}

5. \textit{一般方法を導き出すために、可能な変種の数え上げに基づいた表現法を確立し、どのようなクラスまたは集合のものを「一次命題の項」にすることができる。}

まず、表現すべきクラスまたは集合が、その構成要素すべてに共通する名前や質的特性のみで定義されている場合、その表現は単一の項で行われ、これらの名前や質的特性を表す符号が結合記号なしで組み合わされる。つまり、代数的な乗算過程のように扱われる。例えば、$x$ が不透明な物質を、$y$ が研磨された物質を、$z$ が石をそれぞれ表すとするとき、

\begin{itemize}
\item[] $xyz = $ 不透明で研磨された石;
\item[] $xy (1 - z) = $ 研磨され、かつ石ではない不透明な物質;
\item[] $x (1 - y)(1 - z) = $ 研磨されていない、かつ石でもない不透明な物質;
\end{itemize}

など、他の組み合わせについても同様である。注意していただきたいのは、これらの各表現はその中に含まれる個々の符号と同じ二重性法則を満たしていることである。例えば、

\begin{eqnarray*}
&&xyz \times xyz = xyz;\\
&&xy (1 - z) \times xy (1 - z) = xy (1 - z);
\end{eqnarray*}

などである。このような表現のいずれもを「クラス項」と呼ぶことにし、それは個々のメンバーの共有する性質や名前によって物事を表すためである。

次に、異なる部分が異なる性質、名前、または属性によって定義される集合について述べます。これらの異なる部分に対する表現は別々に形成され、その後、$+$ の記号で接続されます。しかし、私たちが話す集合が、より広い集合から明確な部分のメンバーを除外して形成されたものである場合、除外された部分の象徴的表現の前に $-$ の記号が付加されます。これらの符号の使用についてさらに述べると、次のような観察が追加できます。

6. 一般に述べると、符号$+$は「と」、「または」という接続詞の等価物であり、符号$-$は前置詞「除く」という等価物である。これらの接続詞「と」および「または」について、前者は通常、述語となる集合を描写する際に用いられ、後者は主語となる集合を描写する際に用いられる。例えば、「学者と世渡り上手な人々は幸せを望んでいる」という表現は、前者の例を示すものである。「有用なものには楽を生み出す性質または苦痛を防ぐ性質のいずれかがある」という表現は後者の例を示している。したがって、これらの接続詞を含む表現が最初の命題において現れる際には、それによって思慮上分離されたグループやクラスが互いに全く異なるものでかつ排他的なものであるかどうかを知ることが非常に重要となる。

「学者と世渡り上手な人々」の表現は、両方とも該当する人物も含むべきか否か、「楽を生み出す性質または苦痛を防ぐ性質を持つもの」の表現は、これらの特性を兼ね備えるものも含むべきか否か、このような点について理解することが必要である。私は厳密な意味において、接続詞「と」および「または」がこの種の分離や排他力を有すると理解している。つまり、「すべての$x$は$y$または$z$である」という公式を厳格に解釈すれば、「すべての$x$は$y$でありながら$z$ではない」または「$z$でありながら$y$ではない」という意味になると考えている。

しかしながら、言葉遣いの権威と規範を考えると、反対の解釈が支持される傾向にあることも認めなければならない。つまり、「$y$または$z$」という表現は通常、同時に$y$であり$z$であるものや一方に属するものだけを含むものと理解されることが多い。

しかし、符号$+$がこの議論の対象とした分離力を有することを念頭においておくなら、我々はどのような排他的な表現であれ、思慮上分離された要素に分解し、それらのそれぞれの表現を符号$+$で結びつけるべきである。

そして、その意味に基づいて、表現「$x$ のものや $y$ のもの」は二つの異なるシンボリカルな相当物を持ちます。もし「$x$ のものでなく $y$ のもの、または $y$ のものでなく $x$ のもの」という意味であれば、その表現は
\[ x (1 - y) + y ( 1 - x) \]
となります。シンボル $x$ は $x$ のものを表し、$y$ は $y$ のものを表します。しかし、もし「$x$ のものであるか、または $x$ のものでなければ $y$ のもの」という意味であれば、その表現は
\[ x + y (1 - x) \]
となります。この式は同時に $x$ のものと $y$ のものが存在する可能性を含みます。より完全には以下の形で表すことができます：
\[ xy + x (1 - y) + y ( 1 - x); \]
しかし、最初の二つの項を加えると、元の式に戻ります。

上記の表現は双対性の基本法則（III. 16）を満たします。したがって、
\begin{eqnarray*}
\{x (1 - y) + y (1 - x)\}^2 &=& x (1 -y) + y (1 - x),\\
\{x + (1-x)\}^2 &=& x+y(1-x).\\
\end{eqnarray*}
が成り立ちます。これらの結果は、より一般的な表現の法則の一例であることが後で明らかになります。

7. これらの調査の結果を以下のように表すルールがあります。

\textsc{ルール}.---単純な名前や質をシンボル $x$, $y$, $z$ 等、それらの反対は $1 - x$, $1 - y$, $1 - z$ 等で表し、一般的な名前や質によって定義されたクラスは対応するシンボルを乗算のように接続します。異なる部分から成る集合は、それらの部分の表現を加法記号 $+$ で接続します。特に、「$x$ のものまたは $y$ のもの」を、$x$ と $y$ のクラスが排他的である場合に $x (1 - y) + y (1 - x)$、排他的ではない場合に $x + y (1 - x)$ で表します。同様に、「$x$ のものまたは $y$ のものまたは $z$ のもの」を、$x$, $y$, $z$ のクラスが互いに排他的である場合に $x(1-y)(1-z) + y(1-x)(1-z) + z(1-x)(1-y)$、排他的ではない場合に $x + y (1 - x) + z (1 - x) (1-y)$ で表します。同様の規則が適用されます。

8. この表現の法則に基づいて、逆の解釈の法則が成立します。これらは以下に示す例を通じて、おそらく十分な明瞭さで説明されます。簡略化のために普遍的な主語「事物」や「存在」を省略し、

\[
x = \text{硬い}, y = \text{弾力的}, z = \text{金属}
\]

と仮定します。すると、以下の結果が得られます：

``弾力的でない金属''は $z (1 - y)$ で表されます。

``弾力的な存在と弾力的でない金属の組合せ''は $y + z (1 - y)$ で表されます。

``硬い存在（金属を除く）''は $x - z$ で表されます。

``金属であり、かつ硬いまたは弾力的なものではないもの''は $z - z(1 - x)(1 - y)$ または $z \{1 - (1 - x) (1 - y)\}$ で表され、参照（6）、第2章を参照してください。

最後の例では、実際には ``金属であり、かつ硬くない、弾力的ではないもの'' を表現すべきでした。形容詞間に用いられる接続詞は通常必要ありませんので、シンボル的に表現するべきではありません。

したがって、「硬くて弾力的な金属」は「硬い弾力的金属」と等価であり、$xyz$ で表されます。

次に、「硬い存在（金属でなく非弾力的、または弾力的でなく金属であるものを除く）」という表現を考えます。ここで「それら」は「硬い存在」を指すので、この表現は実質的に ``硬い存在（硬い金属でない非弾力的、または硬い非金属のもの）'' を意味します。「除いて」という言葉はその後の両クラスに適用されます。完全な表現は以下のようになります：

\begin{eqnarray*}
x - xz (1 - y) &-& xy (1 - z);\\
\textrm{または, } x &-& xz (1 - y) - xy (1 - z).\\
\end{eqnarray*}

9. 前述の命題と、それに対する様々な例を挙げたものは、次の命題に必要な前段となります。この次の一連の命題は、現在の章の目的を達成するために必要です。

\textsc{命題 III}.

``可能な種類の検討から一般化された Primary または Concrete 命題の表現方法を導き出す''.

基本的な命題は、最も一般的な意味で、二つの項間に存在する関係を主張するものです。これらの項は必ずしも単語の名詞ではなく、任意の物体の集合を表すことができます。したがって、これらの項を表現する方法は、先に述べた一般原則に含まれており、残るは項間の関係をどのように表現するかを見つけることです。これは明らかにその関係の性質によりますし、特にその関係で項が普遍的であるか特定的であるか（すなわち、ある項が参照する物体の集合全体について話しているのかそれらの一部について話しているのか）によります。

例えば、「固定星は太陽である」と「太陽は固定星である」の二つのクラス間の同一性を表現したいとします。つまり、「すべての固定星は太陽であり、すべての太陽は固定星である」という命題を表すことを目指します。ここで、$x$ を固定星、$y$ を太陽とするとき、必要な等式は
\[
x = y
\]
となります。

「すべての固定星は太陽である」の項では、「すべての固定星」が主語と呼ばれ、「太陽」が述語と呼ばれます。ここで、項「主語」と「述語」の意味を以下のように拡張します。主語とは任意の肯定命題の最初の項を指し、すなわちコプラ（isまたはare）の前にある項を指します。そして、述語はその次に来る二番目の項を指し、つまりコプラの後に来る項を指します。さらに、これらのいずれかが普遍的であるか特定的であるかという仮定を受け入れることで、それぞれの場合において全体のクラスが暗黙のうちに含まれているかまたは一部のみが含まれているかを認めます。このようにすると、最後の例のような場合のルールは以下のようになります：

10. ルール：「命題の主語と述語が両方とも普遍的であるとき、それらの個別の表現を作り出し、等号 (=) で接続する」

この例は、科学の定義や純粋な科学で扱われる主題において通常現れるものであり、Mr. Senior の富の定義がその良い例を提供します：

``富とは、移転可能なもの、供給が制限されているもの、または楽を生み出すもしくは苦痛を防ぐものである。」

この定義をシンボルで表現する前に、注意すべき点として、「もしくは」の接続詞は不要です。富は実際には三つの性質や品質を持つことによって定義されるべきであり、三種類の対象集合から成るわけではありません。それゆえ「もしくは」を省略し、以下のようにします：

\begin{eqnarray*}
w &=& \text{富} \\
t &=& \text{移転可能なもの} \\
s &=& \text{供給が制限されているもの} \\
p &=& \text{楽を生み出すもしくは苦痛を防ぐもの}
\end{eqnarray*}

この主題の性質から明らかであるように、上記定義における「楽を生み出すもしくは苦痛を防ぐ」は、「楽を生み出す；もしくは楽を生み出さない場合に限って、苦痛を防ぐ」と等価であると解釈されます。したがって、単独でこの表現によって定義されるクラスの対象は、楽を生み出すすべてのものと、楽を生み出さないが苦痛を防ぐすべてのものが含まれます。そのシンボル的な表現は

$$p + (1 - p)r$$

となります。

この括弧で示す表現に、移転可能なものと供給が制限されているものという性質を制限する符号 $s$ と $t$ を付加すると、元の定義のシンボル的な等価形は以下のようになります：

\begin{equation}
w = st \{ p + r(1 - p) \}.
\end{equation}

もし「楽を生み出すもしくは苦痛を防ぐ」が単に楽を生み出すものだけを指し、苦痛を防ぐものではないものと、苦痛を防ぐものだけを指し、楽を生み出すものではないものだけを指す場合を考えると、定義のシンボル的な表現は

$$p(1 - r) + r(1 - p)$$

となります。

\begin{equation}
w = st\{p(1-r) + r(1-p)\}.
\end{equation}

これらは、以前より一般的に述べられた内容と一致しています。読者は、「快楽の生産者または苦痛の防止者」という表現を文字通り $p + r$ で翻訳した場合、定義の式を

\begin{equation}
w = st(p + r).
\end{equation}

にするとどのような影響が生じるか知りたいかもしれません。

答えは、この式が (2) の式と等価であるとともに、$stp$ と $str$ を表すクラスが全く異なるものであるという追加の意味を持つこと、つまり、移転可能で供給量に制限のあるものは、同時に快楽を生産し苦痛を防止するものが存在しないということです。どのような等式の完全な意味を決定するかについては後ほど説明します。これまで述べたことは、データを符号の厳格な言語に翻訳する前に、使用している単語の意図した意味を確認することが最も重要であることを示しています。しかし、正確な思考を重視し、適切な言葉遣いはその手段でもあるという人々にとっては、この必要性が悪しきものとは考えられません。

11. 次に，命題の主詞が個別の場合を考えます。例えば、「すべての人間は死人である」です。

この場合、私たちの意味は「すべての人間は何かしらの死人である」となりますので、主詞を表現するために「何かしらの死人」を探します。$v$ でクラス全体を表し、その全てが死人ではないという点以外に特定の属性を持たせます。そして、「死人」を $x$ と表すことにすると、$vx$ は「何かしらの死人」となります。したがって、$y$ を人間と表す場合、求められる等式は
\[
y = vx.
\]

このような考慮から、以下のルールを得ることができます：

\textsc{ルール}.---\textit{主詞と述語を以前と同じように表現し、述語に不定符号 $v$ を付けて、両者の表現を等しいものとする}。

明らかに、$v$ は $x$, $y$ 等と同じ種類の符号であり、一般法則
\[
v^2 = v, \text{または } v(1 - v) = 0.
\]
に従います。

したがって、「惑星は主たるものまたは従属なものである」という命題を表すには、次の方法で進めるべきです：

$x$ を惑星（主題）として表すとし、

\begin{eqnarray*}
y &=& \text{主たる天体;}\\
z &=& \text{従属の天体;}
\end{eqnarray*}

ここで、「または」を用いて「主たる」と「従属の」天体のクラスを絶対的に区別し、与えられた命題においてこれらのクラスが考慮される限り、命題の方程式は

\begin{equation}
x = v \left\{ y (1 - z) + z (1 - y) \right\}.
\end{equation}

と見つかります。この場合、「文字通り」の翻訳として、

\begin{equation}
x = v (y + z)
\end{equation}

という命題前提を表すことも可能ですが、$v$ が不定なクラス符号である限り、これは完全に等価です。しかし、形式（4）の方が良いとされています。なぜなら、

\[
y (1 - z) + z (1 - y)
\]

は互いに異なるクラスを表す項から成り立っているためで、二重性の基本法則を満たします。

「天体は太陽または惑星または彗星である」という命題について考えます。これらのクラスをそれぞれ $w$, $x$, $y$, $z$ と表すと、

\[
w = v \left\{ x (1 - y) (1 - z) + y (1 - x) (1 - z) + z (1 - x) (1 - y) \right\}
\]

という表現になります。ただし、天体が上記の分類のいずれか2つに同時に属することを否定するという仮定に基づいています。

しかし、「天体は太陽であるか、またはそれらが太陽でない場合に限って惑星であり、それらがどちらでもない場合は固定星である」という意味を含むとすれば（つまり、一部の天体が太陽、惑星、固定星のいずれか2つもしくは3つに同時に属することを含む）、必要な表現は

\begin{equation}
w = v \left\{ x + y (1 - x) + z (1 - x) (1 - y) \right\}.
\end{equation}

となります。

以上の例は説明の分類に属し、定義ではありません。実際、命題の述語は通常特別なものであることが多く、この状況が存在しない場合でも、述語が単独語彙であるか、あるいは「は」や「である」という接続詞を用いて述語が普遍的に解釈される形で使われるのです。

12. 次に全否定命題のケースを考えます。例えば、「人間は完璧な存在ではない」というような命題です。

明らかに、この場合、我々は「人間」を含むクラスについて話していないことを示すものとして「何も言及していない」と主張していません。しかし、実質的に「すべての人間」について述べているのです。「彼らは完璧な存在でない」ということです。したがって、この命題の真の意味は次のようになります：

「すべての人間（主語）は（接続詞）完璧ではない（述語）である；」従って、$y$ を「人間」、$x$ を「完璧な存在」つまり「完璧な存在者」と置くと、

\[
y = v \left( 1 - x \right)
\]

となります。他の場合でも同様です。したがって、次のルールがあります：

**ルール**：「どの $x$ のも $y$ のではない」形の命題を表すには、「すべての $x$ のも $y$ のではない」という形に変換し、その後前回の場合と同様に処理します。

13. 次に主語が特別な場合を考えます。例えば、「ある人間は賢明でない」です。ここでは、「否定」が「賢明」の述語に対して適切に適用されることが論理学の目的において明らかに示されています。我々は「ある人間は賢明である」という命題が偽ではないことを意味するのではなく、ある人間が賢明性を欠いているという主張をするのです。したがって、与えられた命題の必要な形式は「ある人間は非賢明である」です。$y$ を「人間」、$x$ を「賢明」つまり「賢明な存在」と置き、

\[
v y = v \left( 1 - x \right)
\]

となります。

14. 上記の決定を総括すると、次の一般ルールがあります：

**一般ルール（原始命題の象徴的表現に対する）**

1. 命題が肯定的である場合、主辞と述語の表現を形成します。もしいずれかが特別なものであれば、その部分に不定符号 \(v\) を付けてから、結果の表現を等しくします。

2. 命題が否定的である場合、まずその真意を述語に否定接頭辞を付けて表し、その後前述の手順を踏みます。

例1または2つ程度の追加的な例は説明のために十分です。

例：「-men には優れた地位に置かれたり、妬まれる心配がない者は存在しない。」

ここで \(y\) を「人間」、\(x\) を「優れた地位に置かれる者」、\(z\) を「妬まれる心配のない者」とします。

現在、「優れた地位に置かれる者」かつ「妬まれる心配のない者」を表す表現は \(xz\) です。したがって、この説明に該当しないクラス、つまりこの説明に適用されないものは \(1 - xz\) で表されます。すべての人間はこのクラスに属します。したがって、
\[
y = v \left( 1 - xz \right)
\]
となります。

もし、この表現された命題が「優れた地位に置かれる人間は妬まれる心配がある」という形で表されれば、その表現は
\[
yx = v \left( 1 - z \right)
\]
となります。

後述するところでは、この表現が特定の仮定（すなわち \(v\) が不定クラス符号であるという仛）において実際的に前記のものと等しいことが明らかになります。

例：「人間には英雄は存在せず、それは自虐と勇気を持つ者だけである。」

ここで \(x = \)「人間」、\(y = \)「英雄」、\(z = \)「自虐を実践する者」、\(w = \)「勇気を持っている者」とします。

主張は、「人間が勇気を持っていないかつ自虐を実践している者は英雄ではない」ということです。

したがって、必要な方程式は
\[
x \left( 1 - zw \right) = v \left( 1 - y \right)
\]
となります。

この章を締めくくるにあたり、論理的な主張の主要な種類を象徴的に表して比較することは興味深いかもしれません。もし「項」や関連するものを表現するために $X$ と $Y$ を用いることに同意すれば、これらの種類は以下のようになります。
\begin{eqnarray*}
X &=& vY, \\
X &=& Y, \\
vX &=& vY.
\end{eqnarray*}

最初の式では述語のみが特別であり、第二の式では両方の項が普遍的です。第三の式では両方が特別です。これらのうち一部は微細な形で含まれています。たとえば、$Y = 0$ の場合、第二の形式は
\[
X = 0;
\]
となり、$Y = 1$ の場合には
\[
X = 1;
\]
となります。これら両方の形式は解釈が可能です。

さらに注意すべき点として、$X$ と $Y$ の表現が「項」の意味を十分に分析した上で基礎となっている場合、基本的な対称法則により
\begin{eqnarray*}
X^2 = X &\text{または}& X (1 - X) = 0, \\
Y^2 = Y &\text{または}& Y (1 - Y) = 0.
\end{eqnarray*}
を満たすことが要求されます。

\newpage

\chapter{記号的推論の根本原理と展開理論（Principles of Symbolic Reasoning）}
\label{ch:v}

\begin{center}
\large\textit{ブール展開定理 $f(x) = f(1)x + f(0)(1-x)$・係数の代数的解釈}
\end{center}
\vspace{4mm}

章[象徴的論理の原理]
{\large 論理符号を含む表式の展開または発展に関する基本的な原理について}

1. この著作の前章は、論理における心の作用の基本法則の調査に捧げられ、その発展として論理符号の法則の調査、そして主張を象徴的语言で表現するための原則の調査が行われた。これらの調査は厳格な意味での前駆的なものであった。それらはこの論文の主要な目的の一つである、基本的な思考の法則に基づく論理の体系または方法の構築への導入を形成する不可欠なものである。この目標の性質とその達成手段に関するいくつかの考慮事項についてここで注意を向けなければなりません。

まず最初に述べたいのは、論理の方法がいかに一般的であるかは、その基本的な過程や法則がいかに一般的であるかに大きく依存することです。本稿の前節で他のこととあわせて調査した論理的過程の「加算」について、その記号は$+$というものです。これらの法則は、考えの中で加算されるクラスやものが互いに排他的であることが必要条件であったことから定められました。

実際、式$x + y$が解釈不能であるように思われるのは、$x$を表すものと$y$を表すものが完全に分離しているという仮定がないと、それらは共通の個体を持たないからです。これらの法則が確認された他の記号的操作についても類似の条件が含まれていました。

そこで問題となるのは、これらの記号的操作と過程を解釈可能であるという同じ条件で制限する必要があるかどうかです。もしそのような制限が必要であるなら、論理における一般的な方法は存在しないことが明白です。一方で、もしそのような制限が不要であるなら、それらの過程が設計された思考の領域において解釈不能に見えるとき、我々は何と見なすべきでしょうか？

これらの問題は単なる論理学の範疇を超えており、人間の推理の発展した形態すべてについても関連しています。これらは記号的语言を使用する基礎となる人間の推理のどの形態においても適用されます。

3. 二番目に、このような過程と解釈との対応の失敗は、人間の理由の通常の適用において現れることはありません。実際には、その意味や適用が理解できる操作のみが行われますし、多くの人々にとって形式的な論理だけが前提と結論をつなぐことに十分なわけではありません。証明の過程で確立される接続列の各段階、つまり証明の中で確立されたすべての間接結果もまた理解されなければなりません。そして、これは疑いなく、日常生活における論理や議論において実際的な条件であり、重要な保護策でもあります。

おそらく多くの人々が、象徴言語の一般的な使用についても同じ原則を拡張すべきであると考えるかもしれません。その主張は、符号の使用に関する法則や公理は解釈可能な場合のみを対象として調査した結果で確立されているため、解釈が不可能または不明な他の場合にその適用を拡大することは正当ではないと主張されるかもしれません。証明の中間段階でのみこのような適用が使用されることを認めるとするとしても、この反論が決定的であるなら、象徴的手法の使用は論理において僅かな利益しか生むものではないことが認めるべきです。おそらくその利点は、より簡便な符号を使用して冗長な記号を置き換えることで得られる機械的な利益に限定されるでしょう。しかし、この反論自体は誤りであり、何らかの事前の予想に関わらず、任意の象徴的手法による推論で導き出された結論の妥当性は、調査過程の各段階で現れた形式的な結果を解釈する能力に依存することはありません。実際には、特定の主題である論理に関連するいくつかの一般的な原則が存在し、それらについて最初に述べた後、その性質と受け入れられる資格に関する一部の意見を提出します。

4. 象征的手法を用いた妥当な推論の条件は---

1. 符号の固定解釈を指定し、これらの符号の組合せ法則がその解釈から正確に決定されるべきである。

2. 解決過程や証明過程は、上記で定めたすべての法則に基づいて進められ、特定の結果の解釈可能性に関わらず、これらを遵守すべきである。

3. 最終的な結果は解釈可能なものであり、その解釈はデータの表現において用いられた解釈システムに従って実際に行われるべきである。これらの原則について、以下の観察がなされるべきである。

5. 符号の固定解釈の必要性は既に十分に述べられている（II. 3）。固定結果がそのような形式でなければならない理由は、符号を使用すること自体が手段であり、その目的は何か可解なる事実や真理を知ることであるという明白な原則に基づく。そしてこの目的を達成するためには、シンボル的な結論を表現する最終的な結果も解釈可能な形式でなければならない。しかし、上記の一般原則または条件の2つ目（V. 4）に関連して最も困難を感じる可能性があり、この点についていくつか追加の説明が必要である。

私は、その原則は私たちの心の一般法則に基づいていると考えるべきであると述べる。この知識は先験的に経験以前に与えられるものではなく、他の心の法則と同じように、特定の例において明確に現れる一般原理から導き出されるものである。

論理の一例を挙げると、シンボルがその解釈に基づく法則に従って用いられながら、その解釈への継続的な参照はなく、証明の鎖は非解釈可能な中間段階を通じて最終的に解釈可能な結果へと導かれる。この一例は、特定の適用の妥当性だけでなく、その中に現れた一般法則を私たちに示すものであるように思われる。個々の事例の積み重ねだけが適切に証拠の重さを増すことはない。それは、原則の即時的な効力が感じられない場合でも、その実用的妥当性を後天的に確認するための手段となるかもしれない。

しかし、これは主張される立場に影響を与えるものではない。すなわち、原則は特定の一例において見られるべきであり、その適用が一般化されているだけでなく、特別な例も真であることが明らかであるべきであるという立場である。

三角法の中間過程における非解釈可能なシンボル $\sqrt{-1}$ の使用は、述べたことの例を提供する。私は、その適用を説明するいかなる方法においても、原則自体が暗黙に前提としているものであると理解している。

しかし、この原則は、他の根拠に基づく形式的な論理によって正当化されるものではないとしても、その存在価値がある種の意味で一般知識の可能性の基礎となる公理的真理の一環として位置づけられるべきであり、心の自己の法則と構造を表現するものとして適切に見なされるべきである。

以下の通り、上記の原理が現在の著作でどのように適用されるかの方法について述べます。既に確認したところによれば、どのような命題の体系も、符号 $x$, $y$, $z$ を用いた方程式によって表現することができます。これらの符号は、解釈が可能であれば、その形式は上記説明の量的符号と同様の法則を満たします（II. 15）。しかしながら、論理過程の形式的な手続きは、符号の法則に依存し、その解釈の性質には依存しないため、これらの符号 $x$, $y$, $z$ を上記説明の種類の量的符号として扱うことができます。つまり、実際には与えられた方程式の論理的な解釈を一時的に置き去りにし、その符号を値が $0$ または $1$ のみを受け入れる量的符号に変換し、このような形で必要な解決過程を行うことができ、最後にはこれらの符号に再び論理的な解釈を与えることができます。そしてこれが実際に採用される手続きの方法であり、各例においてその変換の性質を再度述べることは必要とされません。符号 $x$, $y$, $z$ を量的で、上記説明の種類であると見做して扱う際の適用される手続きは、純粋な論理的な符号に対して行われた場合に限られる思考条件によって制約されることはありませんし、それらを用いる際に自由度が与えられます。これがなければ、論理学における一般方法の探求は希望外れなものとなるでしょう。

以上の手続きの系は、最終的な方程式がその意味を再付与したシンボルによる解釈を可能にする形式にないと、何らかの明白な結果へと導くことはありません。しかし、方程式をこのような形式に簡略化する一般的な方法が存在します。そして、この章の残りはその方法について考慮することに専念します。私は、この方法が解釈を可能にする過程についてあまり詳しく触れません——その点は次の章で扱います——ここでは主にその過程自体に焦点を当て、これを「展開」と特徴付けることができます。この過程の性質への導入として、いくつかの観察を行うことが適切かもしれません。

7. 考慮する任意のクラスについて、ある性質 \( x \) の所有または欠如に関するそのメンバー間の関係を問うた場合を考える。提案されたクラスの各個体は、その性質を所有しているか否かにより分けて考えることができる。したがって、クラスは二つの部分に分割される：一つは性質 \( x \) を所有する個体で構成され、もう一つは所有しない個体で構成される。

この全クラスを思想によって二つの要素部分に分割できる可能性は、他のどの知識源から得られるクラスの構造に関する知識の影響よりも先に存在している。その知識の効果は、所有と非所有の各部分が与えられた性質 \( x \) に対してさらなる条件に従っていることをより正確または不正確な形で教えてくれるだけである。

したがって、このような知識が次のようであるとする：所有する部分の個体はまた特定の性質 \( u \) を持つとし、これらの条件を合わせるとそれらを完全に定義できる。その場合、原始クラスのその部分は表現 \( ux \) (II. 6) で表すことができる。

さらに、原始クラスの非所有する個体が特定の条件 \( v \) の下にあることが判明し、これによって定義されることが明らかになったとする。このことは、これらのメンバーを表現 \( v(1 - x) \) で表すことができることを意味する。

したがって、全クラスは
\[
ux + v(1 - x)
\]
という形で表される。これは与えられた性質 \( x \) の所有または欠如について任意のクラスの象徴的な表現の一般的な開発された形式である。

この純粋論理に基づく一般形式は、既に論理的および量的解釈として参照したシンボル \( x \), \( y \), \( z \) に対する独立した法則の考慮からも導き出すことができる（V. 6）。

8. 定義。---任意の代数式が符号 \( x \) を含むとき、その式は \( x \) の関数と呼ばれ、一般的な簡略形 \( f(x) \) で表すことができる。二つの符号 \( x \) および \( y \) を含む任意の式も同様に \( x \) および \( y \) の関数と呼ばれ、一般的な形式 \( f(x, y) \) で表すことができ、他の場合についても同様である。

したがって、形 \( f(x) \) は以下の関数のいずれかを自由に表すことができる：\( x \), \( 1 - x \), \( \frac{1 + x}{1 - x} \)，etc. 同様に、形 \( f(x, y) \) は以下の形式のいずれかを等しく表すことができる：\( x + y \), \( x - 2y \), \( \frac{x + y}{x - 2y} \)，etc.

同じ記号法の原則により、任意の関数 \( f(x) \) の中で \( x \) を1に置き換えると、その結果は形 \( f(1) \) で表される。また、同様の関数の中で \( x \) を0に置き換えると、その結果は形 \( f(0) \) で表される。例えば、関数 \( f(x) \) が \( \frac{a + x}{a - 2x} \) を表す場合、\( f(1) \) は \( \frac{a + 1}{a - 2} \) を表し、また \( f(0) \) は \( \frac{a}{a} \) を表す。

9. 定義。---符号 \( x \) が論理的符号または値が \( 0 \) および \( 1 \) のみをとる量の符号である任意の関数 \( f(x) \) は、その関数を形 \( ax + b(1 - x) \) に簡約化したときに開発されるという。ここで \( a \) および \( b \) は結果が元の関数と等しいように定められる。

この定義では、任意の関数 \( f(x) \) を仮定された形で表すことが可能であると仮定している。その仮定は以下の命題によって正当化される。

\begin{center}
\textsc{命題 I.}
\end{center}

10. 任意の論理的符号 \( x \) の関数 \( f(x) \) を開発する。

この章で主張された原則により、\( x \) を量的な符号とみなして値 \( 0 \) および \( 1 \) のみをとるという扱いは法的である。

次に、
\[
f(x) = ax + b (1 - x),
\]
と仮定します。そして $x = 1$ と置くと，
\[
f(1) = a
\]
となります。また同じ方程式で $x = 0$ と置くと，
\[
f(0) = b
\]
となります。したがって、$a$ および $b$ の値は決定され、これらを最初の方程式に代入すると
\begin{equation}
f(x) = f(1)x + f(0)(1-x)
\end{equation}
となり、これが求められた展開です。

この章で得られる $f(x)$ の展開は、論理系において通常の代数における $x$ の昇べき級数展開に厳密に対応しています。したがって、テイラーの有名な定理
\[
f(x) = f(0) + f'(0)x + \frac{f''(0)}{1 \cdot 2} x^2 + \frac{f'''(0)}{1 \cdot 2 \cdot 3} x^3 + 等
\]
の表現に $x (1 - x) = 0$ の条件を導入することで得られます。これにより $x^2 = x$, $x^3 = x$, などとし、
\begin{equation}
f(x) = f(0) + \left\{ f'(0) + \frac{f''(0)}{1 \cdot 2} + \frac{f'''(0)}{1 \cdot 2 \cdot 3} + 等 \right\} x
\end{equation}
を得ることができます。

しかし、式（1）において \( x = 1 \) を代入すると、

\[
f(1) = f(0) + f'(0) + \frac{f''(0)}{1 \cdot 2} +
\frac{f'''(0)}{1 \cdot 2 \cdot 3} + \textrm{等}
\]

を得る。従って、

\[
f'(0) + \frac{f''(0)}{1 \cdot 2} + \textrm{等} = f(1) - f(0),
\]

となり、式（2）は置き換えにより

\begin{eqnarray*}
f(x) &=& f(0) + \{f(1)-f(0)\} x,\\
     &=& f(1)x + f(0)(1-x),
\end{eqnarray*}

という形となる。この証明は、関数 \( f(x) \) を \( x \) の昇べき級数に展開するものよりも、本文の証明の方がより一般的ではありません。

方程式の第二項は、どのような形式であれ関数 \( f(x) \) を適切に表すことができます。ここで \( x \) は量的な符号とみなされ、その値は \( 0 \) または \( 1 \) のみを取るため、それぞれの値に対して、

\[
f(1)x + f(0)(1-x),
\]

という展開が関数 \( f(x) \) と同じ値をとります。

例として、関数 \( \frac{1+x}{1+2x} \) を展開することを考えます。ここで \( x = 1 \) のとき、

\[
f(1) = \frac{2}{3}
\]

であり、また \( x = 0 \) のとき、

\[
f(0) = \frac{1}{1} = 1
\]

となります。従って必要な式は

\[
\frac{1+x}{1+2x} = \frac{2}{3}x + 1 - x
\]

であり、この方程式は \( x \) の値のそれぞれに対して成り立ちます。

\begin{center}
\textsc{命題 II.}
\end{center}

\textit{任意の数の論理的符号を含む関数を展開または発展する。}

まず、符号 \( x \) と \( y \) の場合を考えます。そして、展開すべき関数を \( f(x, y) \) で表します。

最初に、\( f(x, y) \) を単独の符号 \( x \) の関数として考え、一般の定理（1）により展開すると、

\[
f(x,y) = f(1,y)x + f(0,y)(1-x)
\]

となります。ここで \( f(1, y) \) は \( x \) を \( 1 \) に置き換えたときの関数の値を表し、\( f(0, y) \) は \( x \) を \( 0 \) に置き換えたときの関数の値を表します。

現在、係数$f(1, y)$を取り出し、$y$の関数として展開すると、

\begin{equation}
f(1, y) = f(1, 1)y + f(1, 0)(1 - y),
\end{equation}

ここで、$f(1, 1)$は$y=1$としたとき$f(1, y)$の値を表し、$f(1, 0)$は$y=0$としたとき$f(1, y)$の値を表します。

同様に、係数$f(0, y)$も展開すると、

\begin{equation}
f(0, y) = f(0, 1)y + f(0, 0)(1 - y).
\end{equation}

式（2）において、$f(1, y)$と$f(0, y)$をそれぞれ式（3）と式（4）の値で置き換えると、

\begin{eqnarray}
f(x, y) = f(1, 1)xy
         + f(1, 0)x(1 - y)
         + f(0, 1)(1 - x)y\\
         + f(0, 0)(1 - x)(1 - y),
\end{eqnarray}

という展開が得られます。ここで、$f(1, 1)$は$x=1$, $y=1$としたとき$f(x, y)$の値を表し、$f(1, 0)$は$x=1$, $y=0$としたとき$f(x, y)$の値を表します。残りも同様です。

したがって、$f(x, y)$が関数$\frac{1-x}{1-y}$を表す場合、

\[
f(1, 1) = \frac{0}{0},\qquad%
f(1, 0) = \frac{0}{1},\qquad%
f(0, 1) = \frac{1}{0},\qquad%
f(0, 0) = 1
\]

となり、与えられた関数の展開は、

\[
\frac{0}{0}xy + 0x(1-y) + \frac{1}{0}(1-x)y + (1-x)(1-y)
\]

となります。次章で見るように、形式$\frac{0}{0}$と$\frac{1}{0}$は、前者が数学的に不定量の記号として知られているように、このような表現において重要な論理的解釈を有することが明らかになります。

さらに、次のステップでは、関数に展開するための3つの符号$f(x, y, z)$を考えます。前と同じ方法で進めると、

\begin{equation}
f(1, 1, z) = f(1, 1, 1)z + f(1, 1, 0)(1 - z),
\end{equation}

\begin{equation}
f(1, 0, z) = f(1, 0, 1)z + f(1, 0, 0)(1 - z),
\end{equation}

\begin{equation}
f(0, 1, z) = f(0, 1, 1)z + f(0, 1, 0)(1 - z),
\end{equation}

\begin{equation}
f(0, 0, z) = f(0, 0, 1)z + f(0, 0, 0)(1 - z).
\end{equation}

\begin{eqnarray*}
f(x,y,z) &=& f(1,1,1)xyz + f(1,1,0)xy(1-z) + f(1, 0, 1)x(1-y)z \\
         &+& f(1,0,0)x(1-y)(1-z) + f(0,1,1)(1-x)yz\\
         &+& f(0,1,0)(1-x)y(1-z) + f(0,0,1)(1-x)(1-y)z\\
         &+& f(0,0,0)(1-x)(1-y)(1-z)
\end{eqnarray*}
これらは、関数 $f(x,y,z)$ において $x=1$, $y=1$, $z=1$ のときの値 $f(1,1,1)$ を用いて表されるものであり、その他の項も同様に定義される。

11. この式から、任意の関数を展開する一般法則を見ることは容易で、その展開方法をルール化することも可能となる。しかし、このようなルールを表現する前に、以下の観察事項を前置しておくことが便利である：---

各得られた展開形式は、特定の項からなり、シンボル $x$, $y$ などはこれらの項に含まれる一方で、係数には含まれない。例えば、関数 $f(x)$ の展開は二つの項、すなわち $x$ と $1 - x$ があり、それぞれの係数は $f(1)$ と $f(0)$ である。また、関数 $f(x, y)$ の展開は四つの項、すなわち $xy$, $x(1 - y)$, $(1 - x)y$, および $(1 - x)(1 - y)$ があり、それぞれの係数は $f(1, 1)$, $f(1, 0)$, $f(0, 1)$, $f(0, 0)$ である。前者の場合の項 $x$ と $1 - x$, 後者の場合の項 $xy$, $x(1 - y)$ などは、これらは展開される関数の形式に依存しない形で存在するため、「構成要素」と呼ぶことにする。明らかに、これらの構成要素は展開される関数の形式に関係なく存在形式が独立している。

構成要素 $xy$ の場合、$x$ と $y$ を「因子」と呼ぶことにする。

したがって、一般の展開法則は二つの部分からなり、第一部分は展開の「構成要素」の形成に関連し、第二部分はそれらの係数の決定に関連する。その内容は以下の通りである：

1. 任意の符号 \( x, y, z \) の関数を展開する。---次のように成分の系列を作り出す：最初の成分は符号の積とする；この積において、第2成分を作るために任意の符号 \( z \) を \( 1 - z \) に置き換える。その後、両方の成分において、他の任意の符号 \( y \) を \( 1 - y \) に置き換え、さらに2つの成分を作る。次に、得られた4つの成分において、他の任意の符号 \( x \) を \( 1 - x \) に置き換えて、新たな4つの成分を作る。これを繰り返し、可能なすべての変更が尽きるまで続ける。

2. 任意の成分の係数を求める。---その成分が因子として \( x \) を含む場合、元の関数において \( x \) を \( 1 \) に置き換える；しかし、因子として \( 1 - x \) を含む場合、元の関数において \( x \) を \( 0 \) に置き換える。同様の規則を符号 \( y, z, 等 \) について適用し、変形した関数の最終的な計算値が求められた係数である。

各成分をそのそれぞれの係数で乗じた和が求める展開となる。

12. 注意すべきは、関数が明示的に含まない符号に対して関数を展開することもできる点である。例えば、符号 \( x \) と \( y \) を基準として関数 \( 1 - x \) を展開する場合、

\[
\begin{array}{cccccc}
\textrm{とき} & x = 1 & \textrm{と} & y = 1 & \textrm{与えられた関数} &= 0. \\
             & x = 1 & " & y = 0 & " & "= 0. \\
             & x = 0 & " & y = 1 & " & "= 1. \\
             & x = 0 & " & y = 0 & " & "= 1.
\end{array}
\]

従って、展開は
\[
1 - x = 0xy + 0x(1 - y) + (1 - x)y + (1 - x)(1 - y)
\]
であり、これは真の展開である。項 \( (1 - x)y \) と \( (1 - x)(1 - y) \) の加算が関数 \( 1 - x \) を生成する。

シンボル1は、規則に従って$x$に関して展開すると、

\[
x + 1 - x
\]

となります。$x$と$y$に関して展開すると、

\[
xy + x(1-y) + (1-x)y + (1-x)(1-y)
\]

となります。どのシンボルの組でも展開すると、それらのすべての可能な構成要素からなる系列が生成されます。

13. 一般展開の性質についての追加の注記をいくつか挙げてみましょう。ここでは、論理符号の一般的な定理（5）を取り上げ、この定理は二つの論理シンボルの関数の型を示すものとします。

まず第一に、その定理は、$x$と$y$が本章で考慮した種類の量的シンボルであり、かつ$f(x, y)$の代数的形式は何らかであっても完全に真で理解できるものです。したがって、上記の論理符号から量的システムへの解釈の変更と最終的な論理解釈への復元の中間段階での分析過程のどのステージでも理解可能に使用することができます。

第二に、$x$と$y$が論理シンボルであり、関数$f(x, y)$の形式が何かしらのクラスまたは事物の集合を表す場合、その定理も完全に真で理解できます。この場合、二項式は常に論理的に解釈可能です。たとえば、$f(x, y) = 1 - x + xy$の場合、定理を適用すると、

\begin{eqnarray*}
1-x+xy &=& xy+0x(1-y)+(1-x)y+(1-x)(1-y),\\
       &=& xy+(1-x)y+(1-x)(1-y),
\end{eqnarray*}

となり、この結果は理解可能で真です。

したがって、定理は上記の種類の量的シンボルに関して常に真で理解可能であり、論理シンボルに関しては解釈可能な場合に限り常に真で理解可能です。従って、この作業において定理を使用する際には、展開結果が解釈不可能である場合、シンボル$x$と$y$は特定の種類の量的シンボルであることを理解することが必要です。

しかし、展開が常に即座に解釈可能なものではないとしても、それは常に結果を導き出すための手段となります。つまり、$x-y$という表現は発展すると次の形になります：
\[
x(1-y) - y(1-x),
\]
これは一般的には解釈不可能な形式です。私たちは考えの中で、$x$でないものが$y$であるもののクラスから、$y$でないものが$x$であるもののクラスを引き出すことができません。なぜなら後者のクラスは前者に含まれていないからです。しかし、もし式$x - y$が方程式の最初の項として現れ、その二番目の項が$0$である場合、発展すると次のようになります：
\[
x(1-y) - y(1-x) = 0.
\]
次章では示すように、$x$と$y$を量的で特定の種類として考えると、上記の方程式は二つの方程式
\[
x(1-y) = 0, \quad y(1-x) = 0,
\]
に即座に分解されることが証明されます。これらの方程式は論理的な解釈が与えられたときには直接的に解釈可能となります。そして、関数が必ずしも展開によって解釈可能になるわけではないにもかかわらず、方程式は常にこの過程によって解釈可能な形に簡約されることが指摘されます。

14. 以下の命題は構成要素の重要な性質をいくつか確立します。その表明において、シンボル$t$は展開の任意の構成要素を表すために使用されます。例えば、展開が二つのシンボル$x$と$y$の関数の展開である場合、$t$は四種類の形、つまり$xy$, $x(1 - y)$, $(1 - x)y$, および$(1 - x)(1 - y)$を表すことができます。構成要素を単一のシンボルで表現する必要があるが、それらを区別するために区別をつけたい場合、その差異は接頭辞によってマークされます。例えば、$t_1$は$xy$を、$t_2$は$x(1 - y)$を表すことができます。

\begin{center}
\textsc{命題 III.}
\end{center}

任意の単一の構成要素$t$が満足する二重性の法則は次のようになります：
\[
t(1-t) = 0.
\]
展開の任意の異なる構成要素の積は$0$に等しく、すべての構成要素の和は$1$に等しいです。

1. 特定の要素 $xy$ を考慮する。次のように表すことができる：
\[
xy \times xy = x^2y^2.
\]
しかし、基本的なクラス記号の法則により、$x^2 = x$, $y^2 = y$ であるから；
したがって
\[
xy \times xy = xy.
\]
ここで $xy$ を $t$ で表すと，
\[
t \times t = t,
\]
または
\[
t(1 - t) = 0.
\]
同様に、要素 $x (1 - y)$ も同じ法則を満たす。なぜなら、次のように示される：
\begin{eqnarray*}
x^2 &=& x, \quad (1 - y)^2 = 1 - y, \\
\therefore \left\{x(1-y)\right\}^2 &=& x(1-y), \quad \textrm{または} \quad t(1-t) = 0.
\end{eqnarray*}
現在、各要素の各因子は $x$ の形または $1 - x$ の形である。したがって、各因子の平方はその因子に等しく、したがって各因子の積（すなわち要素）の平方も要素に等しい。従って、任意の要素を表す $t$ について，
\[
t^2 = t, \quad \textrm{または} \quad t(1 - t) = 0.
\]

2. 任意の二つの要素の積は $0$ である。これは、符号の一般法則 $x (1 - x) = 0$ から明らかであり；同一展開中のどの要素を選んでも、少なくとも一方には因子 $x$ があり、他方にはその対応する因子 $1 - x$ がある。

3. 同じ展開中のすべての要素の和は $1$ である。これは、要素 $x$ と $1 - x$ の和、または要素 $xy$, $x(1-y)$, $(1-x)y$, $(1-x)(1-y)$ の和から明らかである。また、任意の符号集合を用いて $1$ を展開することで（V. 12）、より一般的に証明される。

15. 上記の命題と次のような命題を結びつけることができる。

\begin{center}
\textsc{命題 IV.}
\end{center}

\textit{もし $V$ が任意の一連の要素の和を表し、その個々の係数が $1$ の場合、次の条件が満たされる：
\[
V(1 - V) = 0
\]}

要素 $t_1$, $t_2 \dots t_n$ を考慮すると，
\[
V = t_1 + t_2 \dots + t_n.
\]
両辺を二乗し、$t_1^2 = t_1$, $t_1 t_2 = 0$, などと注意すると，
\[
V^2 = t_1 + t_2 \dots + t_n;
\]
従って
\[
V = V^2.
\]
したがって
\[
V(1 - V) = 0.
\]

\newpage

\chapter{論理方程式の一般的解釈法（Of Interpretation）}
\label{ch:vi}

\begin{center}
\large\textit{未知記号の陽的表現・分数形係数 $\frac{1}{0}, \frac{0}{0}$ の完全な論理的意味}
\end{center}
\vspace{4mm}

\chapter{論理方程式の一般的解釈とその結果としての命題の分析。また、論理関数の解釈可能性の条件}

1. 前章で証明された一般的な規則により任意の関数を完全に展開するとき、二つの異なる要素のセットが含まれることが観察されている。すなわち、展開の成分とその係数である。本章ではまず、これらの成分の解釈について調査し、その後、それらと結びついた係数によってどのように解釈が変更されるかを調べることにする。

「論理方程式」、「論理関数」などという用語は、論理的な前提を表す式や前提から結論までのシンボル的結果列に現れる符号 $x$, $y$ などを含む任意の方程式や関数を一般に指すために使用される。これらの符号 $x$, $y$ などが直ちに論理的に解釈可能でない場合、これらは前章（II.15）、（V.6）で説明した種類の量的符号とみなされ、法則
\[
x(1-x) = 0
\]
を満たすものとする。

したがって、任意の論理的な関数や方程式の解釈問題とは、論理的な値を符号 $x$, $y$ などに割り当てると、それらとともに解釈可能な形に簡略化することとその結果としての解釈を含むものである。これらの慣用的定義は、前章で述べられたこの論文の方法の進行に関する一般原則に合致している。

\begin{center}
\textsc{命題I.}
\end{center}

2. 任意の論理符号 $x$, $y$ などを含む関数の展開の成分は解釈可能であり、それらは符号 $x$, $y$ などが表す性質をすべての可能な方法で予言と否定によって形成される宇宙の排他的な分割を代表する。

概念をより明確にするため、関数展開が二つの符号 $x$ および $y$ を含むと仮定します。この展開は $x$ および $y$ に従って行われました。すると以下の構成要素があります：
\[
xy, \; x(1-y), \; (1-x)y, \; (1-x)(1-y).
\]
これらのうち、最初の項 $xy$ は、同時に $x$ および $y$ によって表される基本的な性質をもつオブジェクトのクラスを表します。また、第二項 $x(1-y)$ は、$x$ の性質を持つが $y$ の性質を持たないオブジェクトのクラスを表します。同様に、第三項は $y$ の性質を持つが $x$ の性質を持たないオブジェクトのクラスを表し、第四項 $(1-x)(1-y)$ は、両方の性質を持たないオブジェクトのクラスを表します。

したがって、これらの構成要素は、$x$ および $y$ によって表現される性質の肯定と否定で説明できる四つのオブジェクトのクラスを表します。これらのクラスは互いに異なるものであり、一つのクラスのメンバーは他のどのクラスにも属しません。なぜなら、各クラスには他のどのクラスにも属さない性質や品質が含まれているからです。また、これら全てのクラスは言及されている宇宙全体を構成します。なぜなら、任意のオブジェクトは提案された性質の存在または非存在によって説明できるため、宇宙内の個々のものは、二つの与えられたクラス $x$ および $y$ とそれらの否定による可能な組み合わせで作られる四つのクラスのいずれかに属すことができます。

ここで述べた注記は完全に一般化されています。任意の展開の構成要素はクラスを表し、これらのクラスは互いに異なる性質を持っていることによって区別され、これら全てが言及されている論理全体を構成します。

3. 构成要素のこれらの性質は、最後の章で結論付けられた定理において表現され、そこで導き出すことができた。各構成要素が個々の符号の基本法則を満足することから、それぞれの構成要素がクラスを表すことを推測することができた。任意の2つの構成要素の積が消滅することから、それらが表すクラスは互いに排他的であると結論付けることができた。最後に、展開中の構成要素の和が1であることから、それらが表すクラス全体で宇宙を形成すると推測することができた。

4. 以上で決定された構成要素の法則とその解釈は、論理方程式の解析と解釈の基礎となっています。すべてのこのような方程式が展開の定理によって解釈される可能性があることは既に述べました。ここでは、推理の結論において現れる可能な解の形式を調査し、その形式がどのように生じるかを示すことにします。厳密には、これらは単一の基本的な型や表現原理の一形態であるが、明確さのために、それらが示す小種類は個別に心に提示されることで理解が深まるでしょう。

これらの形式は3つあります：

\begin{center}
\textsc{形式 i}.
\end{center}

5. まず最初に考慮する形式は、任意の論理方程式 $V=0$ を展開し、結果を成分方程式に分解した後、その解釈を行う場合です。関数は論理符号 $x$, $y$, 等 の組み合わせを含んでおり、これらの組み合わせは分数ではありません。実際には、この形式の解法の中で3つ目の形について述べる際に初めて考慮される問題群においてのみ、分数の組み合わせが生じます。

\begin{center}
\textsc{命題 II.}
\end{center}

論理方程式 $V = 0$ の解釈を行う。

単純化するために、$V$ がただ二つの符号 $x$ と $y$ を含むものとする。与えられた方程式の展開を次のように表す：
\[
axy + bx(1-y) + c(1-x)y + d(1-x)(1-y) = 0;
\]
ここで、$a$, $b$, $c$, および $d$ は定数である。

今、任意の係数、例えば $a$ がゼロでないとする。このとき、方程式の両辺にその係数 $a$ に関連する成分 $xy$ を乗じると、
\[
axy = 0,
\]
となり、$a$ がゼロでないから、
\[
xy = 0,
\]
となる。この結果は他の展開の係数の性質によらず独立である。$x$ と $y$ に論理的な意味を与えて解釈すると、「$x$ によって表されるクラスと $y$ によって表されるクラスに同時に属する個体は存在しない」という結果となる。

しかし、もし係数 $a$ がゼロである場合、項 $axy$ は展開 (1) から省略され、従って方程式 $xy = 0$ は導出できない。

同様に、任意の係数 $b$ がゼロでないとする。このとき、
\[
x(1-y) = 0,
\]
となり、これは「$x$ によって表されるクラスと $y$ によって表されるクラスに同時に属する個体は存在しない」という解釈を受ける。

以上のいずれの解釈も、後に示されることにより異なる形で表現できる。

以上から得られる各項の係数がゼロでないものに対する独立な解釈の合計は、方程式 $V = 0$ の完全な解釈となる。この解析は論理符号 $V$ に含まれる論理符号の数に関係なく本質的に独立であり、従って次のようなルールにより与えられた方程式の解釈が達成される：

\textsc{ルール}：---\textit{関数 $V$ を展開し、係数がゼロでないすべての成分を $0$ に等しいとする。これらの結果の全体的な解釈は与えられた方程式の解釈となる。}

6. 以下に、ユダヤ法で定められた「清浄な動物」の定義を例にとってみましょう。その定義は、「蹄が分かれかつ反芻する動物が清浄な動物である」というものです。ここで、

\begin{tabular}{c l}
&$x$ = 清浄な動物;\\
&$y$ = 蹄が分かれる動物;\\
&$z$ = 反芻する動物;
\end{tabular}

と仮定します。すると、与えられた命題は方程式

\[
x = yz
\]

で表されます。これを形を変えて、

\[
x - yz = 0
\]

とします。現在の方法に従って解釈する形を探求します。最初の項を完全に展開すると、

\begin{eqnarray*}
0xyz + xy(1-z) + x(1-y)z + x(1-y)(1-z)\\
- (1-x)yz + 0(1-x)y(1-z) + 0(1-x)(1-y)z + 0(1-x)(1-y)(1-z)
\end{eqnarray*}

となります。したがって、係数が消えない項は、

\[
xy(1-z) = 0,\; xz(1-y) = 0,\; x(1-y)(1-z) = 0,\; (1-x)yz = 0
\]

となります。これらの方程式は以下のクラスの存在を否定します：

1. 蹄が分かれて反芻するが清浄でない動物。

2. 清浄であり反芻するが蹄が分からない動物。

3. 清浄であり、蹄が分からず反芻しない動物。

4. 蹄が分かって反芻し、しかし清浄でない動物。

すべてこれらの異なる否定は、元の命題そのものに含まれています。逆に、これらの否定が認められた場合、元の命題は必然的な結果として導かれるでしょう。実際、これらは元の命題の個別の要素です。従って、どの基本的な命題も、ある定義された類の存在を否定する一連の否定から分解され、その否定体系から再構築することができます。

ここで疑問が生じるかもしれません：一連の否定や否認から主張的な命題を作り出すのは可能でしょうか？肯定的な要素はどこから導き出されるのでしょうか？

私は答えます。心は、経験に依存しない先験的に存在する宇宙ではなく、経験からの推論として後天的に得られたもの、または断定的推理の基礎となる仮説的な前提として存在する宇宙を前提とします。したがって、「誰も完全に誤らない人間は存在しない」という命題から否定や否認である「完全に誤らない人間は存在しない」という命題が導かれた場合、それは人間が存在するという仮定に基づいて「すべての人間（人間が存在する場合）は誤る」と推論されるか、経験によって人種の存在が確認された後には「すべての人間は誤る」と断言されます。

この命題により結論を示す方法の形式は、「単独または連続否定」の形式と呼ばれることができます。

\begin{center}
\textsc{形式 ii}.
\end{center}

7. 前述の形式は、第二項が $0$ である方程式の展開と解釈から導き出されたものでした。現在の形式は、第二項が $1$ である方程式の展開と解釈から補完的に導き出されます。しかし、前記の命題の分析からは容易に推測できます。

従って、最後に議論した例では、方程式
\[
x - yz = 0
\]
から、係数が0でない以下の構成体を否定する共通の否定を導き出しました：
\[
xy(1-z), \quad xz(1-y), \quad x(1-y)(1-z), \quad (1-x)yz.
\]
これらの残りの構成体は、存在するクラスを形成することから、これらは宇宙全体を構成します。したがって、以下の式が成立します：
\[
xyz + (1-x)y(1-z) + (1-x)(1-y)z + (1-x)(1-y)(1-z) = 1.
\]
これはすべての存在するものは次のいずれかのクラスに属することを主張しています：

1. 足を分かつ蹄を食す清浄な動物。
2. 足を分かつ蹄を食さない不浄な動物。
3. 蹄を食すが足を分けない不浄な動物。
4. 清浄な動物、蹄を食す動物、または足を分ける動物でないもの。

この結論の形は「単項または排他的肯定」の形式と呼ばれ、一方では最終方程式に一つだけ構成体が現れるとき、他方ではそれ以上の一以上の構成体が見られるときに排他的です。

任意の方程式 $V=0$ において $V$ が二重性の法則を満たす場合、この形式の解釈を得るために、まず $1-V=1$ と書き換え、最初の項を展開します。しかし、これは次の一般形に含まれています。

\begin{center}
\textsc{第三形式}.
\end{center}

8. 前述の二つの例では、開発すべき関数がそれぞれ $0$ および $1$ と等しいという状況でした。現在の場合では、対応する関数を任意の論理シンボル $w$ と仮定します。それから方程式 $V = w$ を解釈しようと試みます。ここで最初に、この方程式 $V=w$ または通常は $w=V$ がどのように導かれるかを示すことが必要です。

「清浄な動物」という定義を再考しましょう。これは以前の例で使用されたもので、「清浄な動物とは蹄が分かれかつ反芻する動物である」という定義です。そして、反芻する動物が清浄な動物と蹄が分かれる動物との間にある関係を決定することを考えます。与えられた命題を表す方程式は
\[
x = yz
\]
であり、私たちの目的は$z$を$x$と$y$の解釈可能な関数として定めることが出来ることです。

ここで$x$, $y$, $z$を特殊な法則に従う量の記号として扱います。この方程式から解くことで、
\[
z = \frac{x}{y}
\]
を得ることが出来ます。

しかし、この方程式は現在では解釈可能な形ではありません。これを解釈可能な形に変換することが出来れば、必要な関係が得られます。

上記の方程式の第二項を展開すると
\[
z = xy + \frac{1}{0}x(1 - y) + 0(1 - x)y + \frac{0}{0}(1 - x)(1 - y)
\]
となります。そして、後述する命題3において、この式は以下の解釈を受けることが示されます：

「反芻する動物はすべて蹄が分かれる清浄な動物（蹄も分かれている）に加えて、蹄が分かれない不浄な動物の不定な余剰部分（一部、全部、または全くない）で構成される。」

9. 以上は論理学において最も一般的な問題の一例であり、以下のように述べることが出来ます：「任意の論理方程式$x$, $y$, $z$, $w$を結びつけるシンボルについて与えられたとき、その他のシンボル$x$, $y$, $z$, などによって表されるクラスに属するクラス$w$との関係の解釈可能な表現を求めよ。」

この問題の解法は、与えられた方程式から符号 \( w \) の表式を他の符号で定めるところに尽きる。そしてその表式を展開によって解釈できる形にする。与えられた方程式は常に、各符号に対して一次方程式である。したがって、\( w \) に対する必要とされる表式は常に求められる。

実際、与えられた方程式を展開すれば、どのような形であれ \( w \) に関する方程式は次の形式で得られる：
\begin{equation}
Ew + E^\prime(1-w) = 0,
\end{equation}
ここで \( E \) と \( E^\prime \) は他の符号の関数である。以上の式から
\[
E^\prime = (E^\prime - E)w
\]
を得る。したがって
\begin{equation}
w = \frac{E^\prime}{E^\prime-E}
\end{equation}
であり、この右辺を展開の規則によって展開すると、結果を論理的に解釈する必要がある。

もし分数 \(\frac{E^\prime}{E^\prime - E}\) の分子と分母に共通な因子が存在する場合でも、それらを排除することは許されない。ただし、その共通な因子が単なる数的定数である限りは例外である。符号 \( x, y, 等 \) を量として考えたとき、0や1のような値を取ることにより、共通な因子が0となり、代数的な簡約法則が失われる可能性があるからである。これは代数的除法の公理（II. 14）における我々の注釈で述べた事態と一致する。

方程式 (2) の形に解を表し、任意の不正な簡約を行うことなくその結果を展開定理によって解釈することは、この論文の一般原則に厳密に従った方法である。

\( 1 - w \) が他のクラス \( x, y, 等 \) に対する関係が必要であれば、(1) の式から同様にして
\[
1-w = \frac{E}{E-E^\prime}
\]
を得る。この解釈も次の命題の方法によって行うことができる：

\begin{center}
\textsc{命題 III.}
\end{center}

10. 任意の論理方程式 \( w = V \) の解釈を決定する。ここで \( w \) はクラス符号であり、\( V \) は他のクラス符号の関数でその形式は制限されていない。

上の式の第二項を完全に展開します。その結果の各係数は、四つのクラスのいずれかに属し、それぞれの解釈と共に次のように議論します。

1. 係数が $1$ の場合、これは宇宙の象徴であり、任意の二つのクラス符号の積は両方のクラスに共通する個体を表すため、係数が $1$ の任意の構成要素は制限なく解釈されるべきです。つまり、その構成要素が代表するクラス全体が含意されます。

2. 係数が $0$ の場合、論理においても算術においてもこれは「何もなし」を表す象徴であり、前接する構成要素によって表されるクラスのどの部分も取り扱われません。

3. 係数が $\frac{0}{0}$ の形の場合、算術において $\frac{0}{0}$ は通常不定数を表します。この作品のシステムにおいても同様な符号が不定クラスを表すべきであることは、類推から示唆されます。その真の意味は次の例から明らかになります。

例えば、「人間でない不死者は存在しない」という命題を考えます；この命題を象徴で表現し、これらの象徴が既に証明された法則に従って「不死者」を「人間」の用語で逆定義する方法を探します。

もし「人間」を $y$、そして「死ねる存在」を $x$ で表すならば、「非死者の人間は存在しない」という命題は方程式
\[
y \left(1 - x\right) = 0
\]
によって表現されます。この式から $x$ の値を求めることになります。

上記の方程式より、
\[
y - yx = 0, \quad \text{または} \quad yx = y
\]
を得ます。これは通常の代数方程式であれば、両辺を $y$ で割る操作を行いたいところですが、第11章で述べたように、現在使用している符号では除算の操作は実行できません。そこで、その操作を表現し、前章の方法で結果を展開することになります。

まず最初に
\[
x = \frac{y}{y}
\]
を得ます。そして、第二項を指定された通り展開すると
\[
x = y + \frac{0}{0} (1 - y)
\]
となります。これは「死ねる存在」($x$)がすべての「人間」($y$)に加えて、「人間でないもの」($1 - y$)からなる余剰部分を表す係数 $\frac{0}{0}$ によって示されます。

次に、前提からどのような「人間でないもの」の余剰部分が含まれるかを調査します。この余剰部分はすべての人間でない存在を含むかもしれませんし、一部の人間でない存在のみを含むかもしれませんし、あるいは全く含まないかもしれません。これらの仮定のいずれも、人間全員が死ねることを主張する前提とは矛盾しないため、$\frac{0}{0}$ はそのクラスのすべて、一部、または全くを表すことを示します。

上記で$\frac{0}{0}$の意味を決定したのは、特定の一例に対する検討に基づくものであり、しかし証明に含まれる原理は一般化されたものであり、その符号が現れるどのような状況においても同様の分析法が適用不可能な場合はありません。このように$\frac{0}{0}$を不定クラス符号と呼び得ますし、必要であれば単純な符号$v$で置き換えることも可能です。ただし、基本的な法則 $v(1 - v) = 0$ を満たす必要があります。

4. ある展開の成分の係数がこれまでのどのケースにも属さない場合があるかもしれません。その成分の真の解釈を特定するために必要な定理は以下の通りです：

11. 定理---任意のクラスまたは集合 $w$ を表すための関数 $V$ が展開され、その展開における任意の成分の数値係数 $a$ が法則
\[
a(1 - a) = 0,
\]
を満たさない場合、その成分は 0 に等しいものとすべきです。

この定理を一般化して証明するために、与えられた展開を以下のように表します：
\setcounter{equation}{0}
\begin{equation}
w = a_1 t_1 + a_2 t_2 + a_3 t_3 + \cdots,
\end{equation}
ここで $t_1$, $t_2$, $t_3$, \ldots は成分を、$a_1$, $a_2$, $a_3$, \ldots は係数を表します。また、$a_1$ と $a_2$ が法則
\[
a_1(1 - a_1) = 0, \quad a_2(1 - a_2) = 0;
\]
を満たさないことを仮定します。

他の係数はその問題に従う法則に従っており、したがって次のようになります：
\[
{a_3}^2 = a_3, \quad 等
\]
次に、方程式（1）の両辺を自身で乗算します。結果は以下のようになります。
\begin{equation}
w = {a_1}^2 t_1 + {a_2}^2 t_2 + \quad 等
\end{equation}
これは、方程式
\[
w = V^2
\]
の展開を表しているから明らかです。しかし、実際に（1）を平方に展开し、以下の事実も考慮することで証明できます：
\[
{t_1}^2 = t_1, \quad {t_2}^2 = t_2, \quad t_1 t_2 = 0, \quad 等
\]
これらの成分の性質からです。次に、（2）を（1）から引くと次のようになります：
\[
\left(a_1 - {a_1}^2\right) t_1 + \left(a_2 - {a_2}^2\right) t_2 = 0.
\]
または，
\[
a_1 \left(1 - a_1\right) t_1 + a_2 \left(1 - a_2\right) t_2 = 0.
\]
最後の方程式を $t_1$ で掛けると、$t_1 t_2 = 0$ より
\[
a_1 \left(1 - a_1\right) t_1 = 0, \quad \textrm{従って} \quad t_2 = 0.
\]
同様に、同じ方程式を $t_2$ で掛けると
\[
a_2 \left(1 - a_2\right) t_2 = 0, \quad \textrm{従って} \quad t_2 = 0.
\]

したがって、その係数が符号自身と同じ基本法則に従わない任意の成分についても一般化して示すことができます。つまり、そのような成分は個別に $0$ と等しい必要があります。このような係数が通常現れる形は $\frac{1}{0}$ です。これは無限の代数的符号です。任意の数が無限に近づくほど（その表現を許す限り）、それは前述の基本法則から逸脱します。

$\frac{0}{0}$ の符号については、以前議論したように、この法則に違反するとは必ずしもなりません。これは $0$ と $1$ の両方の数値を持つことができます。しかし、その実際的な解釈として不定なクラス符号である場合、私は算術的性質から類推に基づいてそれを導くことができないと考えます。それは実験的に確立されるべきです。

12. 以上をまとめると以下の要約になります：

1. 符号 $1$ が展開の項の係数として使用された場合、それはその成分が表す全クラスを含むことを示します。
2. 係数 $0$ はそのクラスから何も選ばないことを示します。

3番目の符号 $\frac{0}{0}$ は、クラスの完全に不確定な部分を示し、すなわちそのメンバーの一部である「あるもの」、「何もなし」、または「すべて」が選択されることを意味します。

4番目の符号が他の係数として使用される場合、それはその符号で先頭に付く構成要素が0と等しいことであることを示します。

したがって、問題の解が展開によって得られ、その形式が
\[
w = A + 0B + \frac{0}{0}C + \frac{1}{0}D,
\]
の場合、その解は以下の2つの方程式に分解される：
\begin{eqnarray}
w &=& A + vC,\\
D &=& O,
\end{eqnarray}
ここで $v$ は不確定なクラス符号です。式 (3) の解釈は、$w$ が必要とする定義のクラスの要素が何であるか、または可能であれば何であるかを示します；そして式 (4) の解釈は、元の問題の要素間の関係が $w$ に依存しない状態で存在することを示します。

これらの解釈法則が存在します。追加するには、これらは普遍的に適用され、その使用は常に例外や失敗なしで容易であるという点です。

13番目の推論として、$V$ が独立して解釈可能な論理関数である場合、それはシンボル的な法則 $V(1-V) = 0$ を満たすことを示します。

独立して解釈可能な論理関数とは、その符号を含むシンボルによって表されるものの間の関係を前提としないで解釈できるものを指します。したがって $x(1-y)$ は独立して解釈可能です、しかし $x-y$ はそうではありません。後者の関数は、その解釈のために $y$ によって表されるクラスが $x$ によって表されるクラスの完全な部分であるという条件を前提としています。一方、前者の関数はそのような要件を含みません。

現在、$V$ が独立して解釈可能であり、$w$ が $V$ 内に含まれる個体の集合を表す場合、方程式
\[
w = V,
\]
は $V$ の展開中の任意の構成要素の消滅を伴う結果を引き起こさないことを示します。なぜなら、そのような構成要素の消滅は、シンボル $V$ によって表されるもののクラス間の関係を意味するからです。したがって、$V$ の展開は以下の形式になります：
\[
a_1t_1 + a_2t_2 + \textrm{等}
\]
係数 $a_1, a_2, 等$ はすべて次の条件を満たします：
\[
a_1(1-a_1) = 0, \quad a_2 (1-a_2) = 0, \quad \textrm{等}
\]

\#\#\# 翻訳

したがって、第5章の命題4の論理に従って、関数$V$は以下の法則に従う：
\[
V(1-V) = 0.
\]

この結果は、$V$が何かのクラスや集合を表すものであるという事実から$\textit{a priori}$明らかであるが、その構成要素の性質からも導かれることが確認される。条件$V(1-V) = 0$は「論理関数の解釈可能性の条件」と呼ぶことができる。

14. 最後の命題で得られた一般形式の解または論理的結論を、「項間の関係」として表すことができる。以前と同じく、私は「項」を述語の部分、単純なものであれ複雑なものであれ、接続詞「は」や「である」で結ばれた部分として定義する。個々の符号によって表される事物のクラスを、「述語の要素」と呼ぶことができる。

15. 例題1．「清められた動物」（VI.6）の定義が与えられ、それから「不清められた動物」の説明を求められた。

ここで、$x$が「清められた動物」を表し、$y$が「蹄を分ける動物」を表し、$z$が「反芻する動物」を表すとすると、
\[
x = yz;
\]
従って，
\[
1-x = 1-yz ;
\]
そして、右辺を展開すると、

\[
     1-x = y(1-z)+z(1-y)+(1-y)(1-z);
\]
これは以下のように解釈できる：「不清められた動物は蹄を分けるが反芻しないすべてのもの、反芻するが蹄を分けるしないすべてのもの、そして蹄を分けても反芻してもいないすべてのものである。」

例題2．同じ定義が与えられ、蹄を分けることのない動物の説明を求められた。

方程式 $x = yz$ から、
\[
y = \frac{x}{z};
\]
従って，
\[
1 - y = \frac{z - x}{z};
\]
そして，第二項を開発すると，
\[
1 - y = 0 \cdot xz + \frac{-1}{0}x(1 - z) + (1 - x)z + \frac{0}{0}(1 - x)(1 - z).
\]

ここで，ルールに従って，係数が $\frac{-1}{0}$ の項は別途 $0$ と等しいと仮定する。これにより，
\begin{eqnarray*}
1 - y &=& (1 - x)z + v(1 - x)(1 - z),\\
x(1 - z) &=& 0;
\end{eqnarray*}

これらの式の最初の方程式は，解釈によって次のように命題を与えます：
「蹄を分かつ動物はすべて反刍する不潔な動物であり，また反刍しない不潔な動物の不定な余剰（あるもの，ないもの，またはすべて）である。」

二つ目の方程式は次の命題を与えます：「反刍しない清浄な動物は存在しない。」これは先に述べた独立関係の一つです。我々は蹄を分かつ動物に対する直接的な関係を探求しましたが，清浄な動物と反刍する動物との間には，蹄を分かつ動物とは独立して，前提によって与えられる別々の関係があります。この関係も過程により必然的に与えられます。

例題 3．--次のような定義を考えましょう：「責任者はすべて自由に行動できる理性的存在または行動の自由を自発的に犠牲を払った理性的存在である」。そして，その分析を適用します。

\begin{tabular}{c l c l}
$x$ &は&責任者を, \\
$y$ &は&理性的存在を, \\
$z$ &は&自由に行動できる存在を, \\
&w &は&行動の自由を自発的に犠牲を払った存在を表す.
\end{tabular}

この定義の表現において，二つの互いに排他的な選択肢「自由に行動できる理性的存在」および「行動の自由を自発的に犠牲を払った理性的存在」が与えられていますから，両方の分類には同時に個体は存在しないと仮定します。これにより，命題を文字で解釈することができます：

\[
x = yz + yw.
\]

まず、「理性的な存在」が「責任を持つ存在」、自由を有する存在、そして自由の放棄を自発的に選択した存在に対してどのような関係にあるかを定めるべきです。おそらくこの目的は、前提の要素間の関係を以下のように表現することで達成されるでしょう：理性的性が責任から推論できる範囲、行動の自由、自由の放棄、それらの反対面からどれほど推測できるか。

(6)より、

\[
y = \frac{x}{z + w}
\]

を導きます。第二項を開発し、係数が0である項は省略しますと、

\begin{eqnarray*}
y = \frac{1}{2}xzw + xz(1-w) + x(1-z)w + \frac{1}{0}x(1-z)(1-w)\\
                                       + \frac{0}{0}(1-x)(1-z)(1-w),
\end{eqnarray*}

となります。係数が$\frac{1}{2}$と$\frac{1}{0}$の項を0に等しいとするとき、

\begin{equation}
y = xz(1-w) + xw(1-z) + v(1-x)(1-z)(1-w);
\end{equation}

\begin{equation}
xzw = 0;
\end{equation}

\begin{equation}
x(1-z)(1-w)= 0
\end{equation}

となり、これらを解釈すると、

\textsc{直接結論}：「理性的な存在はすべて自由を放棄せずに行動の自由を持つ責任者であるか、または自由を放棄したにもかかわらず行動の自由を持たない責任者であり、またそれ以外の責任を持たず、自由を持たず、自由の放棄も自発的に行わない無数の存在（一部、全部、あるいは全く）を含む。」

\textsc{第一独立関係}：「責任を持つ者は同時に行動の自由を持つ者でもあり、かつ自由の放棄を自発的に選択した状態にある者でもない。」

\textsc{第二}：「責任を持つ者は同時に行動の自由を持たない者でもあり、かつ自由の放棄を行わなかった状態にある者でもない。」

上記の独立関係は別の形で表現され、より便利な形式に置くことができます。したがって(8)から、

\[
xw = \frac{0}{z} = 0z + \frac{0}{0}(1-z)
\]

となり、開発すると

\begin{equation}
xw = v(1-z);
\end{equation}

また(9)から、

\[
x(1-w) = \frac{0}{1-z} = \frac{0}{0}z + 0(1-z)
\]

となり、開発すると

\begin{equation}
x (1 - w) = vz;
\end{equation}

となります。そして(10)と(11)を解釈すると以下の命題が得られます。

1. 責任ある者が自己の自由を自発的に犠牲にした場合、彼は自由ではない。

2. 責任ある者が自己の自由を自発的に犠牲にしなかった場合、彼は自由である。

これらの関係は、前で定めたものと異なる表現形態に過ぎない。

16. これらの結果を検討する際には、読者に以下の点を念頭に置いていただきたい。つまり、推論法や解析の唯一の目的は、原始命題における語彙の関連性から導かれる必然的な関係を決定することである。したがって、この目標の達成度合いを評価する際には、語彙の意味から暗示される他の関係とは一切関わりなく、その表現された関連性のみに関心を持つべきである。したがって、「自己の自由を自発的に犠牲にした者は自由ではない」という関係は、語彙の意味から直接導かれるものであり明らかに明白である。これにより、方法によって割り当てられた二つの独立した関係のうちの一つは一方では無駄な制限を受けることなく、他方では不要な要素として見なされる可能性があるように思われる。しかし、原始前提における語彙の関連性のみに焦点を当てる場合、その関係がこれらの批判から免れていることが明らかになる。この解釈は直接の結論と独立した関係の両方を通じて完全であり、何ら余分な要素も含んでいない。

上記で述べた暗黙的な関係、「自己の自由を自発的に犠牲にした者は自由ではない」というものを取り入れる場合、これは異なる命題として扱うことができる。その適切な表現は以下の式である：
\[
w = v (1 - z).
\]

この方程式は原始前提を表す方程式と併用して使用する必要がある。このような検討を行う方法は、将来的の章で理論体系の議論に入ることによって明らかになるだろう。分析が導く結果の唯一の違いは、上記で導かれた独立した関係の一つが代替されることである。

17. 例4．例2と同様の定義を仮定し、無理人についての説明を求めることとする。

$$1 - y = 1 - \frac{x}{z+w} = \frac{z + w - x}{z + w}$$

ここで、

\begin{eqnarray*}
1 - y &=& \frac{1}{2} xzw + 0 \cdot xz(1 - w) + 0 \cdot x(1 - z)w - \frac{1}{0} x (1 - z)(1 - w)\\
&+& (1-x)zw + (1-x)z(1-w) + (1-x)(1-z)w + \frac{0}{0}(1-x)(1-z)(1-w)\\
&=& (1-x)zw + (1-x)z(1-w) + (1-x)(1-z)w + v(1-x)(1-z)(1-w)\\
&=& (1-x)z + (1-x)(1-z)w + v(1-x)(1-z)(1-w),
\end{eqnarray*}

ここで $xzw = 0, \quad x(1-z)(1-w) = 0$。

これらの独立な関係は、以前に得たものと同一であり、これは明らかにその通りでなければならない。なぜなら、与えられたクラス $y$ の存在に関わらず独立的に存する何らかの関係が、その対立クラス $1 - y$ の存在に関わらず独立的に存することであるから。

最初の方程式による直接的な解は以下のようになる：「無理人とは、自由に行動できる者や、自己放棄により自発的に行動の自由を放棄した者で、行動の自由がない者全体。さらに、行動の自由がないが自己放棄をしなかった未定の残りの無責任な者も含まれる」。

18. この章で分析された命題は、定義と呼ばれる種類のものである。私は、第二項または述語項が個別のものではなく、一般型が $Y = vX$ であるようなものを一つも議論していない。ここで $Y$ と $X$ は論理符号 $x, y, z, 等$ の関数であり、$v$ は未定のクラス符号である。このような命題の分析は、符号 $v$ の除去により（これは重要なステップではないが）著しく容易になり、この過程は次の章で説明する方法に依存している。また、単独の命題理論を完成させるために必要な別の重要問題についても考慮しないが、その分析は未来のページで開発される命題系の簡略化の方法に依存している。

\newpage

\chapter{消去理論（Of Elimination）}
\label{ch:vii}

\begin{center}
\large\textit{連立論理方程式からの媒介変数の完全消去法・論理的帰結の純化}
\end{center}
\vspace{4mm}

\#\#\# 消去法について

1. 前章で述べた例では、元の前提のすべての要素が結論に再出現し、ただ順序や接続関係が異なるだけでした。しかし、一般的な推論において特に前提が複数ある場合によく起こるように、一部の要素は結論には現れません。このような要素、つまり通常「中間項」と呼ばれる要素は、他の要素との間に助けることで形成される接続を確立するために元の命題に導入されるものとみなすことができます。これらの他の要素だけが結論の表現に入ることを目指します。

2. 中間要素や中間項について、誤った概念が広まっています。一般的には、推論はこれらの要素の排除に特有のものであるとされていますが、最後の章で示された例により、その見解は否定されます。推論は二つの前提から一つの中間項を排除して、その項が含まれない単一の結論を導くという過程の基本型態であるとされるのです。従って、一般的には、この観点に基づいて「三段論法」がすべての推論の基礎または共通型態であると主張されています。しかし、推論の形式や構造がいかに複雑であっても、それは一連の三段論法に分解されるものとされています。

そのような見解の妥当性については、後述する章で検討します。現在は、記号によって表現された論理システムと一般的な代数システムにおける排除に関する重要な点について注意を喚起したいと思います。代数的なシステムでは、二つの方程式から一つの符号、三つの方程式から二つの符号、一般に$n$個の方程式から$n-1$個の符号が排除されることが可能です。したがって、与えられた独立方程式の数と、それらから排除できる量の符号の数との間には明確な関係性があります。

しかし、論理システムではそのような固定された関係性は存在しません。与えられた前提を表す方程式の数と、それを排除することができる典型的な符号の数との間に固有の関係性はありません。一つの方程式から無限に多くのような符号が排除されることがあります。一方で、無限に多くの方程式から単一のクラス符号のみが排除されます。

論理システムにおいては、この特殊な状況では、どのような場合でも排除問題を解決する方法が存在します。これは論理の記号が従う驚くべき対称法則の結果です。与えられた前提方程式から導かれる方程式に加えて、思考の基本法則から抽出された別の方程式または方程式系が追加されます。これにより問題の解決に必要な手段が提供されます。

論理の記号が従う対称法則から派生する多くの結果の中でも、おそらく最も注目すべき結果はこの点です。

3. アルゴリズムにおいてしばしば起こる通り、与えられた方程式のシステムから符号を排除することで、単なる恒等式 $0 = 0$ のような結果が得られ、残された符号に独立した関係が存在しない場合があります。同様な結果は論理のシステムでも現れ、その解釈も類似しています。このような状況は以前述べた原理の一般性を損なうものではありません。我々がこれから着手する方法の目的は、任意の数の符号から任意の数の論理方程式を排除し、結果で残された実際の関係を明らかにすることです。しかし、そのような残存する関係が存在しない場合もあります。そのような場合において、方法の真理性は恒等的な命題へと導かれる事実によって示されます。

4. 以下の命題では採用される表記法は前章と同じです。$f(x)$ とは、論理符号 $x$ を含む任意の式を意味し、他の論理符号も含まれる場合もあります。$f(1)$ とは、$x$ がその中で 1 に置き換わったときの $f(x)$ の値を表します。同様に、$f(0)$ は $x$ が 0 に置き換わったときに同じ関数がどのようになるかを示します。

\textsc{命題 I}

5. クラス符号 $x$ を含む（または含まない）任意の論理方程式 $f(x) = 0$ の場合、方程式

\[
f(1)f(0)=0
\]

は $x$ の解釈によらず真であり、これは上記の方程式から $x$ を排除した完全な結果となります。

つまり、任意の与えられた方程式 $f(x) = 0$ から $x$ を排除するためには、その方程式において順次 $x$ を 1 に置き換え、その後 $x$ を 0 に置き換えて得られる2つの結果の方程式を乗算することで実現されます。

同様に、任意のクラス符号 $x$, $y$ 等から任意の形式 $V = 0$ の方程式を排除する完全な結果は、最初の項を与えられた符号の成分に完全に展開し、すべての係数を乗算してその積を 0 と等しいものとする方法で得られます。

方程式 $f(x) = 0$ の最初の項を展開すると、(V.10)より

\begin{eqnarray}
f(1)x + f(0)(1 - x) = 0; \nonumber \\
\textrm{または, }
[f(1) - f(0)]x + f(0) = 0.
\therefore x = \frac{f(0)}{f(0) - f(1)}; \nonumber\\
\textrm{および }
1 - x = -\frac{f(1)}{f(0) - f(1)}. \nonumber
\end{eqnarray}

これらの式を $x$ と $1 - x$ に代入し、基本方程式
\[ x(1 - x) = 0,
\]
に入れる。

すると次のようになる：
\begin{eqnarray}
-\frac{f(0)f(1)}{[f(0) - f(1)]^2} = 0; \nonumber \\
\textrm{または, }
f(1)f(0) = 0,
\end{eqnarray}

これが求められる形である。

6. この過程から、実際には与えられた方程式 $f(x) = 0$ と普遍的に真の方程式 $x(1 - x) = 0$（論理記号の基本法則）との間で消去が行われていることがわかります。したがって、項を消去するためには一つの前提または方程式だけが必要であり、思考の必要法則が別の前提または方程式を供給します。この結論の証明は他の形でも示すことができますが、その要素は心によって提供されます。次のように進めることが可能です：

(1) を $x$ に掛けると
\[
f(1)x = 0,
\]
となります。そして、通常の代数学の形式で $x$ を消去してこの方程式と (1) から $x$ を消去する方法を探します。

もし二つの代数方程式が次の形であれば
\[
ax + b = 0,
\]
\[
a'x + b' = 0;
\]
そのとき、$x$ の消去結果は
\[
ab' - a'b = 0
\]
と知られています。しかし、上記の方程式を (1) と (3) にそれぞれ比較すると
\[
a = f(1) - f(0), \quad b = f(0);
\]
\[
a' = f(1), \quad b' = 0;
\]
となります。これらを (4) に代入すると
\[
f(1)f(0) = 0,
\]
となり、先程と同じ結果になります。この証明の形では、基本方程式 $x(1 - x) = 0$ が (3) の導出過程で現れます。

7. この重要な定理の証明をさらに別の形式で示します。これは半論理的な性格を持ち、この定理の証明をより明確にするのに役立ちます。

先に
\[
f(1)x + f(0)(1 - x) = 0.
\]
としました。まずこの方程式を $x$ で掛けて、次いで $1 - x$ で掛けると次のようになります：
\[
f(1)x = 0, \quad f(0)(1 - x) = 0.
\]
これらの式から解き出して展開すると

\begin{eqnarray*}
     f(1) &=& \frac{0}{x} \\
          &=& \frac{0}{0}(1-x),\:\textrm{展開上},\\
     f(0) &=& \frac{0}{1-x} \\
          &=& \frac{0}{0}x.
\end{eqnarray*}
これらの式の直接的な解釈は以下の通りです：

1. すべての $f(1)$ を表すクラスに含まれる個体は、$x$ の個体ではありません。

2. すべての $f(0)$ を表すクラスに含まれる個体は、$x$ の個体です。

従って、一般的な論理によれば、$f(1)$ と $f(0)$ のクラスに同時に属する個体はありません。つまり、$f(1)f(0)$ のクラスには個体が存在しません。したがって，
\[
     f(1)f(0) = 0.
\]
または、展開された式を乗算すれば、結果はすぐに得られます。

8. 定理（5）から次のルールを得ます：

提案された方程式から任意のシンボルを消去するためのルール

ルール：方程式の項が必要に応じて移動させた後、シンボルに順次 1 と 0 を与え、結果となる方程式を乗算します。

この命題の最初の部分は証明されました。

9. 次に一般的な方程式を考えます：
\[
     f(x,y) = 0;
\]
その第一項は $x$，$y$ および他のシンボルの任意の関数を表します。

先述した通り、この方程式から $y$ を消去すると
\[
     f(x,1)f(x,0) = 0;
\]
となります。これは与えられた方程式で順次 $y$ を 1 と 0 に置き換え、結果を乗算することで導かれる形です。

さらに、得られた結果から $x$ を順次 1 と 0 に置き換え、結果を乗算すると
\[
     f(1,1)f(1,0)f(0,1)f(0,0) = 0;
\]
が得られます。これは最終的な消去結果です。

この方程式の第一項の四つの因子は原始方程式の第一項 $f(x,y)$ の完全展開の四つの係数に等しいので、命題の第二部分は明らかになります。

\begin{center}
\textsc{例題}.
\end{center}

10. 例題 1．「すべての人間は死ぬ」を示す命題とそのシンボル表現
\[
     y = vx,
\]

$y$ が「男」を、$x$ が「人間」を表すとき、不定なクラス符号 $v$ を除去し、その結果を解釈することが求められます。

まず、項を最初の側に移動すると、以下の式を得ます：
\[
y - vx = 0.
\]
$v = 1$ の場合、これは
\[
y - x = 0;
\]
$v = 0$ の場合、これは
\[
y = 0;
\]
となります。これらの2つの方程式を乗算すると、以下の式を得ます：
\[
y - yx = 0,
\]
または $y(1 - x) = 0$ となり、ここで $y^2 = y$ と注意します。

上記の方程式は要求される除去結果であり、その解釈は「存在しない男は人間ではない」--これは明らかな結論です。

最後に得られた方程式から非人間を描写するものを探すと、以下のようになります：
\[
x = \frac{y}{y},
\]
\[
\therefore 1 - x = \frac{0}{y}.
\]

したがって、展開すると $1 - x = \frac{0}{0}(1 - y)$ となり、これを解釈すると「存在しないものは男ではない」ことが分かります。これは一般的な論理で「対比的変換」または「否定的な変換」と呼ばれる例です（WhatelyのLogic, 第2巻. 第2章. 第4節）。

例題 2.--「男は完璧なものではない」という命題を表す方程式
\[
y = v(1 - x),
\]
において $y$ が「男」、$x$ が「完璧な存在」を表すとき、$v$ を除去し、その結果から「完璧な存在」および「不完全な存在」の描写を求めることです。以下のようにします：
\[
y - v(1 - x) = 0.
\]
これを除去法により解くと、以下の式を得ます：
\[
\{y - (1 - x)\} \times y = 0,
\]
または
\[
y - y(1 - x) = 0,
\]
または
\[
yx = 0;
\]
これは「完璧な男は存在しない」という命題により解釈されます。上記の方程式から以下の式を得ます：
\[
x = \frac{0}{y} = \frac{0}{0}(1 - y)\:（展開により）；
\]
これを解釈すると「完璧なものは男ではない」ことが分かります。同様に、
\[
1 - x = 1 - \frac{0}{y} = \frac{y}{y} = y + \frac{0}{0}(1 - y),
\]
を解釈すると、「不完全な存在はすべての男と、男でないものの中から不定な余剰部分を持つ存在」であることが分かります。

11. 一般に、命題の論理を扱う際には、対応する方程式において$v$という不定なクラス符号が現れる場所から最初にこれを消去することが最も便利である。これにより方程式の形式は変化するが、その意味は損なわれない。

次に、第四章の例を適用してみましょう。命題「男たちは高貴なる地位に置かれても忌望の目から免れる者はいない」に対応する式は
\[
y = v(1 - xz)
\]
であり、その等価な命題「高貴なる地位にある男たちは忌望の目から免れない」に対応する式は
\[
yx = v(1 - z)
\]
です。そして、$v$が不定なクラス符号である場合、これらの方程式自体も等価であることが観察されました。この事実を証明するために、各方程式から$v$という符号を消去するだけが必要です。

最初の方程式は
\[
y - v(1 - xz) = 0
\]
であり、これを満たすためにまず $v = 1$ とし、次に $v = 0$ とし、その結果を乗じることで得られます。
\begin{eqnarray*}
(y - 1 + xz)y &=& 0,\\
\textrm{または } yxz &=& 0.
\end{eqnarray*}

次に与えられた方程式の第二式は移項により
\begin{eqnarray*}
yx - v(1 - z) &=& 0; \\
\textrm{よって } (x - 1 + z)yx &=& 0,\\
\textrm{または } yxz &=& 0,
\end{eqnarray*}
となり、先に得られた結果と同じです。読者はこの結果を容易に解釈できるでしょう。

12. 例3．本章の一般方法の主題として、Mr. Senior の財富の定義を再検討します。「財 Wealthは移転可能な物質であり、供給が制限され、楽なことの生産者または苦痛の防止者である」です。この定義に従って、以前の議論に合わせて、定義の最後の部分に共通する性質を持つすべてのものを考慮に入れると、代表方程式は
\begin{eqnarray*}
w &=& st \{pr + p(1 - r) + r(1 - p)\}, \\
\textrm{または } w &=& st \{p + r(1 - p)\},
\end{eqnarray*}
となります。

ここで、
\begin{tabular}{c l l}
$w$ & は & 財 Wealthを表す．\\
$s$ & は & 供給が制限された物質を表す．\\
$t$ & は & 移転可能な物質を表す．\\
$p$ & は & 楽なことの生産者を表す．\\
$r$ & は & 苦痛の防止者を表す．\\
\end{tabular}

上記の式から、不要なシンボルを排除し、主と述語の任意の提案された配置に従って解と展開で結果を表現することができます。

まず、要素 $r$（痛みの防止に関連するもの）を取り除いた場合の $w$（財富）の表式を考えましょう。方程式の項を最初の側に移すと、

\[
w - st(p + r - rp) = 0
\]

となります。

$r = 1$ とすると、最初の項は $w - st$ となり、$r = 0$ とすると $w - stp$ になります。従って、ルールにより、

\[
(w - st)(w - stp) = 0
\]

または

\[
w - wstp - wst + stp = 0;
\]

より、

\[
w = \frac{stp}{st + stp - 1}
\]

となります。この方程式の第二項を展開すると、

\[
w = stp + \frac{0}{0}st(1 - p)
\]

となります。

したがって、結論を得ます：—「財富は供給が制限され、移転可能で楽を生み出すものすべてに加えて、楽を生み出さないものも無限の余剰分である。」これは十分明らかです。

注意すべき点として、式（7）の乗算を実行し、その方程式を式（8）の形に簡略化する必要はなくとも、$w$ に関する他のシンボルの表式を求めることができます。展開法は常に乗算法を代用することができます。例えば、$w$ を中心として（7）を展開すると、

\[
(1 - sf)(1 - stp)w + stp(1 - w) = 0
\]

となり、

\[
w = \frac{stp}{stp - (1 - st)(1 - stp)}
\]

となります。この方程式を展開しても、先ほどの結果と同じになります：

\[
w = stp + \frac{0}{0}st(1 - p)
\]

13. 次に、財富、移転可能性、楽を生み出す傾向との関係から、供給が制限されるものの説明を求めることを考えましょう。痛みの防止に関する一切の参照を省きます。

方程式（8）は、元の方程式から$r$を消去した結果であり、

\[
w - s \left(wt + wtp - tp\right) = 0;
\]

より、

\begin{eqnarray*}
s = \frac{w}{wt + wtp - tp}\\
= wtp + wt(1-p) + \frac{1}{0} w(1-t)p + \frac{1}{0}w (1-t)(1-p)\\
+ 0 (1-w) tp + \frac{0}{0} (1-w) t (1-p) + \frac{0}{0} (1-w)(1-t) p\\
+ \frac{0}{0} (1-w) (1-t) (1-p).
\end{eqnarray*}

この解を逐語的に解釈し、その後にその解から導かれる一般的な観察を述べ、最後には結論の表現を多少簡略化する方法を示すことにします。

まず最初に、逐語的な解釈は以下の通りです。供給が限られた物事は、

\textit{すべて移転可能な財産と楽を生み出すもの--すべて移転可能な財産でないものであり、楽を生み出すものでもない--無限の量の財産ではないが、移転可能でないものであり、楽を生み出すものでもないか、移転可能でないものであり、楽を生み出すものであるか、移転可能でもなく、楽を生み出すものでもない。}

ここで、係数が$\frac{1}{0}$の項に加えることができる独立な関係は以下の通りです：

1. \textit{移転不可能で楽を生み出す財産は存在しない。}
2. \textit{移転不可能で楽を生み出すものではない財産は存在しない。}

14. この解について、以下のような観察がなされることが予想されます。

まず最初に、得られた「供給が限られた物事」の表現において、「すべて移転可能な財産と楽を生み出すもの」という項は部分的に冗長であると考えられるかもしれません。なぜなら、すべての財産は（元の命題から暗示され、独立な関係で直接主張されるように）必然的に移転可能であるからです。

私は答えます。普通の言葉で言えば、「財産」という用語に「移転可能な」を付加する必要はないと考えられますが、他の部分での議論が結論を表現し、すべての財産は必ず移転可能であるということが示されたとしても、この方法の完全性には、各結論の項が代表する対象を、選択した質や要素との関係を明確に述べることで、全ての場合において完全に定義することがあります。これは異なる部分が実際的に区別され独立していることを保証し、冗長さを防ぐため必要です。

例えば、我々が考慮してきた対象のペアが「移転可能」の言葉を含まなかった場合、「すべての財産」という単項語は「すべての財産が移転可能である」と「すべての財産が移転不可能である」という二つの項に分解されることが論理的に可能です。しかし、後者の項は「独立した関係」によって消失することが示されます。従って、それは必要性を欠き、冗長となります。残された項は実際に得られた結論と一致します。

論理的な分割によって生じる余分な項や冗長な項が生成されない解法は、「純粋な解法」と呼ばれるべきです。これは先述した展開法と消去法により得られる全ての解法に該当します。また、現在のシステムにおいてその方法が可能である場合でも、一般的な代数的な消去法を採用しても、得られた解法の純粋性を保証することはできません。その方法の一般性不足は唯一の欠点ではなく、他の欠点も存在するのです。

15. 二つ目の点として、結論は二つの項からなり、これら二つの項の総合的な意味を一つの項で表現するのがより便利であることが指摘されるべきである。つまり、「すべての生産的で移転可能な財」および「すべての非生産的で移転可能な財」という代わりに、「すべての移転可能な財」だけで十分である。この指摘は完全に正当であるが、このような簡略化が可能となると、それは必ずしも方程式の形式によって提示される形で即座に提案されるものであり、その方程式を最も単純な形に簡約した場合、解釈される結果もまたその最適な形となる。元の解決策において、係数が1である項 $wtp$ と $wt(1 - p)$ の合計は $wt$ となり、係数が $\frac{1}{0}$ である項 $w(1 - t)p$ と $w(1 - t)(1 - p)$ の合計は $w(1 - t)$ となる。そして、係数が $\frac{0}{0}$ である項 $(1 - w)(1 - t)p$ と $(1 - w)(1 - t)(1 - p)$ の合計は $(1 - w)(1 - t)$ となる。したがって、完全な解は

\[
s = wt + \frac{0}{0}(1-w)(1-t) + \frac{0}{0}(1-w)t(1-p)
\]

であり、独立関係としては、

\[
w(1-t) = 0, \text{または} w = \frac{0}{0}t
\]

である。

この解釈は以下のように立訳されるべきである：

1. 物資の供給が限定されているものはすべて移転可能な財であり、その余剰部分は財でないかつ移転不可能なものと、移転可能だが財でもないかつ生産的でないものである。
2. すべての財は移転可能である。

これは一般の結論と付随する条件を最も単純な形で表現できる方法である。

16. 二つ以上の符号を除去する必要があるとき、提案された方程式からそれらを除去するために（6）命題Iを使用したり、または順次に除去することができます。過程の順序は関係ありません。方程式
\[
w = st \left(p + r - pr\right)
\]
から $r$ を除去し、結果を得ました：
\[
w - wst - wstp + stp = 0.
\]
もし $r$ と $t$ の両方を除去する必要があったら、上記が過程の最初のステップとして残り、最後の方程式から $t$ を除去することになります。ここで $t = 1$ のとき、その方程式の第一項は
\[
w - ws - wsp + sp
\]
となり、$t = 0$ のとき同じ項は $w$ となります。したがって
\[
w \left(w - ws - wsp + sp\right) = 0,
\]
または
\[
w - ws = 0
\]
が求められる結果となります。

最後の結果から $w$ を決定すると、

\[
w=\frac{0}{1-s}=\frac{0}{0}s
\]

となり、「すべての財は供給が限られている」という結論を得ます。$p$ が方程式に含まれていないことに注意すると、上記は結論の要素間のどのような関係に関わらず、任意の条件のもとでも真であることが明らかです。

元の方程式に戻ります。もし $s$ と $t$ を除去する必要があるなら、

\[
w=st(p+r-pr)
\]

となります。しかし、それぞれの符号を別々に除去する代わりに、$st$ を一つの符号として扱うだけで十分です。なぜなら、基本的な符号の法則により

\[
st (1 - st) = 0
\]
が成り立つからです。

したがって与えられた方程式を以下のように形変形します：
\[
w - st (p + r - pr) = 0;
\]

そして $st$ を順次に 1 と 0 に置き換え、その結果の積を求めると、

\begin{eqnarray*}
(w - p - r + pr) w = 0,\\
\textrm{または }  w - wp - wr + wpr = 0,
\end{eqnarray*}

が求められる結果となります。

特に、「快楽を生むもの」（$p$）を「財」（$w$）と「苦痛の防止をもたらすもの」（$r$）の関係で表すために、方程式を解くと、

\[
p = \frac{w(1-r)}{w(1-r)}
\]
\[
= \frac{0}{0}wr + w(1 - r) + \frac{0}{0}(1 - w)r + \frac{0}{0} (1 - w) (1 - r)
\]

\[
= w(1 - r) + \frac{0}{0}wr + \frac{0}{0}(1 - w).
\]

以下の結論を得ます：—「快楽を生むものとは、すべての財産でなく、痛みを予防するもの、そして一定量の非財産である。」

同じ方程式から得られるものは、

\[
1-p = 1 - \frac{w(1-r)}{w(1-r)} = \frac{0}{w(1-r)},
\]

が展開すると、

\begin{eqnarray*}
w(1-p)
= \frac{0}{0}wr + \frac{0}{0}(1-w) \cdot r + \frac{0}{0}(1-w) \cdot (1-r)\\
= \frac{0}{0}wr + \frac{0}{0}(1-w).
\end{eqnarray*}

したがって、—「快楽を生まないものは、財産であり、痛みを予防するもの、または一定量の非財産である。」

類似の例についての議論も同様に容易です。

17. 最後の消去の例では、与えられた方程式から複合符号 $st$ を単一の符号として扱い、それを消去しました。同じ方法は、個々の符号が基本的な符号の法則を満たす任意の符号の組み合わせにも適用可能です。したがって、表現 $p + r - pr$ は自身と等しくなるため、これを単一の符号 $y$ として表すと、

\[
y = y^2, \text{または} y(1-y) = 0,
\]

という基本的な法則が満たされます。消去法のルールは、各個々の符号がその法則に従うことを前提としているため、この法則を満たす任意の関数やその組み合わせから方程式を消去するには、一度の操作でそれが可能となります。

本章と前章で採用した解釈の形式は、おそらく最も直接的な意味を示すものであり、一般的な言語表現に適応したより快適な表現法も、等しく真実かつ適切に使用されるべきである。例えば、方程式（9）は次のようないくつかの方法で解釈できる：「財産は供給が制限され、移転可能であり、喜びを生み出すものか、または供給が制限され、移転可能であり、喜びを生み出さないものか」。逆に、「何かが供給が制限され、移転可能であり、喜びを生み出すものは財産である」という解釈も提供される。最終的な展開で項が単位係数を持つとき、このような反対の解釈は常に与えられる。

18. 注記：基本方程式 $f(1)f(0) = 0$ は、任意の方程式 $f(x) = 0$ の符号 $x$ を消去した結果を表すものであり、驚くべき解釈を持つ。

$f(x) = 0$ という方程式は、クラス $x$（例えば「人間」）を代表する個体についての主張を含むことを意味し、他の個体も含まれる可能性がある。我々の目的は、その主張に他の個体間にどのような関係が含まれているかを確認することである。方程式 $f(1) = 0$ は、「人間」が宇宙全体を構成する場合の元の主張になることを表し、方程式 $f(0) = 0$ は「人間」が存在しなくなる場合の元の主張になることを表す。したがって、方程式 $f(1)f(0) = 0$ は、元の主張に基づいて、どちらの場合でも真実であることを表す。つまり、「人間」が「すべてのもの」または「何もない」という場合でも等しく真実であることを表す。したがって、この定理は次のことを示している：「与えられた対象のクラスが存在するか消失しても等しく真実なものは、そのクラスに全く依存せず、逆もまた然りである」。ここでは、数学的思考の法則から直接導き出される形式的な結果を哲学的な一般命題として解釈するもう一つの例を見ることができる。

\newpage

\chapter{命題体系の還元法（Of Reduction）}
\label{ch:viii}

\begin{center}
\large\textit{連立方程式の単一方程式への還元 $V = 0$・二乗和と同値性}
\end{center}
\vspace{4mm}

\chapter{還元論}
{\large 系統の命題に対する還元について}

1. 前章では、単独の基本命題、すなわち単一の方程式で表される基本命題の理論を十分に確立し、その上で適切な方法を構築しました。任意の与えられた方程式系中の要素が除去され、残存する要素間の関係が提案された形（否定、肯定、または通常の主辞と述語の関係など）で導き出せる方法を示しました。しかし、命題の系統に対する考察に進むべき余地があります。つまり、与えられた命題系の方程式から任意の数の符号が除去できることを調査し、結果の解釈によって残存する符号間の全関係を導き出し、特に他の要素を用いて特定の要素またはその組み合わせの表現を決定し、適切な形で結論を提示できる方法を見つけることです。これらの質問は、任意の方程式系やその中の方程式を等価な単一の方程式に還元できることを示すことで回答されます。そして、この還元過程において、以前の章で扱った基本命題の理論が重要な拡張を受けていることも理解できます。これは性質上特異なものではありません。科学の多くの専門分野は内部から完全に把握することができず、外部からの視点でも研究し、他の類似の主題と関連して考慮することが必要で、その全体像を理解するために不可欠です。

本章では、方程式の系を等価な単一の方程式に簡約する二つの異なる方法を示します。これらのうち最初の方法は任意定数倍数の使用に基づいています。理論上は十分単純ですが、その後の消去や展開などの過程が煩雑になるという欠点があります。しかし、これは最初に発見された方法であり、またその単純さから、この方法を維持することが適切であるとされています。第二の方法では任意定数の導入は必要ありませんし、前記の方法よりもほぼすべての面で優れています。したがって、本作業後の調査において一般的に採用されるでしょう。

2. まず最初の方法について考慮します。

\centering
\textsc{命題I}.

任意の論理方程式の系は、最初の方程式以外をそれぞれ異なる任意定数倍数に乘じて加えることにより、等価な単一の方程式へ簡約されることが可能です。

第VI章の命題2によれば、任意の一つの方程式$f(x, y..) = 0$の解釈は、最初の項の係数が消えない成分を0に等しいものとして得られます。したがって、与えられた二つの方程式$f(x, y..) = 0$と$F(x, y..) = 0$の両者の合意は、そのような成分（両方または一方の与えられた方程式で現れる）を0に等しいシステムによって得られます。第VI章の規則に基づいて展開した結果です。例えば、

\begin{equation}
xy - 2x = 0,
\end{equation}

\begin{equation}
x - y = 0;
\end{equation}

を考えると、最初の方程式の展開は

\[
- xy - 2x (1 - y) = 0
\]

となり、

\begin{equation}
\textrm{従って, }xy = 0, x (1 - y) = 0.
\end{equation}

二番目の方程式の展開は以下のようになります：
\[
x(1-y) - y(1-x) = 0;
\]
\begin{equation}
\textrm{従って，} x(1-y) = 0, \quad y(1-x) = 0.
\end{equation}

両方の展開において係数が消えない成分は $xy$, $x(1-y)$, $(1-x)y$ であり、これらは以下のシステムを導きます：
\begin{equation}
xy = 0, \quad x(1-y) = 0, \quad (1-x)y = 0;
\end{equation}

これは、両方の展開で与えられた方程式が生じる成分と係数が消えないものと同じであるため、二つのシステムに等しいです。二項方程式系の場合を限定して考えると、次のように一つの方程式を求めることになります：その展開によって与えられた方程式と同じ成分と、かつ係数が消えないものを生じます。

与えられた方程式 $V_1 = 0, V_2 = 0$ を表すために
\[
V_1 = 0, \quad V_2 = 0,
\]
を用います。ここで $V_1$ と $V_2$ は論理符号 $x, y, z, 等$ の関数です。このとき、任意の定数量 $c$ を用いて
\begin{equation}
V_1 + cV_2 = 0,
\end{equation}

という一つの方程式が要求される目的を達成します。

ここで，完全な展開 $V$ の中で任意成分 $t$ とその数的係数 $A$ を用いて
\[
At \quad (t \textrm{は成分，} A \textrm{は数的係数})
\]

を表し，それに相当する成分 $Bt$ を用いて
\[
Bt,
\]

を展開 (6) の中で表すと，
\[
(A + cB)t.
\]

成分 $t$ の係数が消えるのは $A$ と $B$ がともに消えるときのみであり，それ以外の場合は消えない。もし $A$ と $B$ がともに消えないことを仮定し，
\begin{equation}
A + cB = 0,
\end{equation}

を満たすとする場合、以下の状況だけが現れます。

1. $A$ が消え，$B$ が存在する場合。この場合，上記の方程式は
\[
cB = 0,
\]

となり，$c = 0$ を要求します。しかし，これは $c$ が任意の定数であるという仮定と矛盾します。

2. $B$ が消え，$A$ が存在する場合。この仮定は (7) を
\[
A = 0,
\]

に帰着させ，自身の仮定を破ります。

3. $A$ や $B$ のどちらも消えない場合。方程式 (7) から，
\[
c = \frac{-A}{B}
\]

という具体的な値が得られ，これは $c$ が任意であるという仮定と矛盾します。

したがって、係数 $A + cB$ は $A$ と $B$ がともに消滅するときのみ消滅し、それ以外の場合は消滅しない。従って、同じ構成要素は式 (6) の展開において、$V_1 = 0, V_2 = 0$ の方程式群における係数が消えないかのように現れる。そして、方程式 $V_1 + cV_2 = 0$ は方程式群 $V_1 = 0, V_2 = 0$ と同等である。

類似の論理により、一般の方程式系
\[
V_1 = 0, V_2 = 0, V_3 = 0, \ldots;
\]
は単純な方程式
\[
V_1 + cV_2 + c'V_3 + \ldots = 0,
\]
$c, c', \ldots$ が任意の定数であることで置き換えることができる。この方程式は、前章で述べた通常の論理方程式と同様に扱うことができる。任意の定数 $c_1, c_2, \ldots$ は論理的な記号ではない。これらの定数は法則
\[
c_1 (1 - c_1) = 0, c_2 (1 - c_2) = 0,
\]
を満たさない。しかし、それらの導入は第II章15節と第V章6節で述べられた一般原則によって正当化され、ほぼすべての後の調査において例示されているものであり、論理に関連する思考体系における解釈可能な形になるまで、論理の記号としてのこれらの符号が量の記号として扱われる法則 $x (1 - x) = 0$ に従って扱われることを前提とする。

3. 以下は上述した方法を例示するものである。

例題1。ある特定の物質類の性質に関する分析から、以下の一般的な結論が得られたと仮定しよう：

1. 性質 $A$ と $B$ の組合せでは、性質 $C$ や性質 $D$ のいずれかは必ず存在するが、両方同時に存在しない。

2. 性質 $B$ と $C$ の組合せでは、性質 $A$ と $D$ はどちらもそれらと共に存在するか、またはどちらも存在しない。

3. 性質 $A$ と $B$ が両方とも欠如している場合、性質 $C$ と $D$ も両方とも欠如しており、逆に、性質 $C$ と $D$ が両方とも欠如している場合、$A$ と $B$ も両方とも欠如している。

以上の原著英文を上記から、ある特定の事例において性質 \(A\) の存在が性質 \(B\) と \(C\) の存在または非存在について何らかの結論を導くことができるかどうかを決定する必要がある場合を考えます。性質 \(D\) については考慮しません。

\[
\begin{array}{cccccc}
性質 & A を x で表す; \\
     & B を y で表す; \\
     & C を z で表す; \\
     & D を w で表す.
\end{array}
\]

すると、前提のシンボル表現は以下のようになります。

\[
xy - v (w(1 - z) + z (1 - w));
yz = v (xw + (1 - x ) (1 - w));
(1-x)(1-y)=(1-z)(1-w).
\]

これらの方程式の最初の2つから、不定なクラス符号 \(v\) を個別に消去すると、

\[
xy (1 - w (1 - z) - z (1 - w)) = 0;
\]
\[
yz (1 - xw - (1 - x)(1 - w)) = 0.
\]

ここで、展開により

\[
1 - w(1 - z) - z(1 - w) = wz + (1 - w)(1 - z),
\]

および

\[
1 - xw - (1 - x) ( 1 - w) = x(1 - w) + w(1 - x),
\]

と表すことができ、これらの式で、

1 - x を $\bar{x}$、1 - y を $\bar{y}$ などと簡略化して表すと、最後の3つの方程式から
\setcounter{equation}{0}
\begin{eqnarray}
xy (wz + \bar{w}\bar{z}) = 0; \\
yz (x\bar{w} + \bar{x}w) = 0; \\
\bar{x}\bar{y} = \bar{w}\bar{z};
\end{eqnarray}

これらの式から \(w\) を消去する必要があります。

上記の2つ目の方程式に \(c\) を乗じ、3つ目の方程式に \(c'\) を乗じて加算すると、

\[
xy (wz + \bar{w}\bar{z}) + cyz (x\bar{w} + \bar{x}w) + c'(\bar{x}\bar{y} - \bar{w}\bar{z}) = 0.
\]

ここで \(w\) を1とし、従って \(\bar{w}\) を0とすると、上記方程式の最初の項は

\[
xyz + c\bar{x}yz + c'\bar{x}\bar{y}.
\]

また、同じ項で \(w\) を0とし、\(\bar{w}\) を1とすると、

\[
xy\bar{z} + cxyz + c'\bar{x}\bar{y} - c'\bar{z}.
\]

したがって、\(w\) の消去結果は以下の形で表すことができます。

\begin{equation}
(xyz + c\bar{x}yz + c'\bar{x}\bar{y}) (xy\bar{z} + cxyz + c'\bar{x}\bar{y} - c'\bar{z}) = 0;
\end{equation}

この方程式から \(x\) を決定します。

もし今度も以前の例と同様に進めると、上記方程式の最初の項の因子を全て乗算することになりますが、このような手順は必ずしも必要ではないことを示すことも良いでしょう。$x$ という符号の表現を求めることを目的として $x$ について最初の項 (4) を展開すると、

\[
yz(y\bar{z} + cyz - c'\bar{z}) x + (cyz + c'\bar{y})(c'\bar{y} - c'{z})(1 - x) = 0;
\]
\[
\textrm{または, } cyzx + (cyz + c'\bar{y})(c'\bar{y} - c'\bar{z})(1 - x) = 0;
\]

から、

\[
x = \frac{(cyz + c'\bar{y})(c'\bar{y} - c'\bar{z})}{(cyz + c'\bar{y})(c'\bar{y} - c'\bar{z}) - cyz};
\]

さらに、$y$ と $z$ について第二項を展開すると、

\[
x = 0yz + \frac{0}{0} y\bar{z} + \frac{c'^2}{c'^2}\bar{y}z + \frac{0}{0} \bar{y}\bar{z};
\]
\[
\textrm{または, } x = (1 - y)z + \frac{0}{0} y(1 - z) + \frac{0}{0}(1 - y)(1 - z);
\]
\[
\textrm{または, } x = (1 - y)z + \frac{0}{0}(1 - z);
\]

これらの結果は、$A$ の性質が存在する場所では、$C$ が存在し $B$ が存在しないか、あるいは $C$ が存在しないという解釈になります。逆に、$C$ の性質と $B$ の性質が両方とも存在しない場所では、$A$ の性質が存在します。

これらの結果は、次の方法で得られることが遥かに容易です。しかし、異なる方法を持つことは互いの検証のためにも満足すべきことです。

4. 次の方法へ移ります。

\begin{center}
\textsc{命題 II.}
\end{center}

任意の方程式 $V_1 = 0$, $V_2 = 0$, $等$ が、その最初の項の展開がすべて正の係数を持つ成分のみで構成されている場合、これらの方程式は加算によって単一の等価な方程式に組み合わせることができます。

以前と同様、$At$ を関数 $V_1$ の展開における任意の項、$Bt$ を関数 $V_2$ の展開における対応する項とします。それらは、それぞれの与えられた方程式を加えた結果、方程式
\setcounter{equation}{0}
\begin{equation}
V_1 + V_2 + \textrm{等} = 0,
\end{equation}
の展開における対応する項は
\[
\left(A+B+\textrm{等}\right)t.
\]
となります。しかし、仮定によれば、係数 $A$, $B$, 等 のいずれもが負でないため、導出された方程式の合計係数 $A + B$ などは、それぞれの係数 $A$, $B$, 等 が同時に零となるときのみ零になります。したがって、元の方程式系 $V_1=0$, $V_2=0$, 等 の各方程式を総合すると、方程式 (1) の展開における同種の成分は、元の方程式系から導かれた各方程式の解釈と等価となります。

\begin{center}
\textsc{命題 III.}
\end{center}

5. 任意の方程式系 $V_1=0, V_2=0, 等$ の項が移項により同一側に整理された場合、その総合的な解釈は以下の単一行の方程式によって表されます：

\[
V_{1}^2 + V_{2}^2 + 等 = 0,
\]

これは与えられた方程式の平方を加えた結果です。

任意の方程式系 $V_1=0$ の展開により得られる方程式
\[
a_1 t_1 + a_2 t_2 + 等 = 0
\]
において、$t_1, t_2, 等$ は成分であり、$a_1, a_2, 等$ はそれらの対応する係数です。すると、方程式 $V_{1}^2=0$ の展開により得られる方程式
\[
a_{1}^2 t_1 + a_{2}^2 t_2 + 等 = 0,
\]
が得られます。これは展開の法則からも、または条件
\[
t_{1}^2=t_1, \quad t_{2}^2=t_2, \quad t_1 t_2=0
\]
のもとで関数 $a_1 t_1 + a_2 t_2, 等$ を平方した結果からも証明できます。

第3章の定理3において指定されたものである。したがって、方程式 $V_{1}^2=0$ の展開における構成要素は、方程式 $V_1=0$ の展開における構成要素と同一であり、これらには正の係数が付いています。同様に、方程式 $V_2=0, 等$ についても同じことが言えます。したがって、最後の定理により、方程式

\[
V_{1}^2 + V_{2}^2 + 等 = 0
\]

は、

\[
V_1=0, \quad V_2=0, \quad 等
\]

という方程式系と解釈上等価です。

補遺。--最初の項が既に条件

\[
V^2=V \textrm{, または } V(1-V)=0
\]

を満たす任意の方程式 $V=0$ は、平方化の過程によっても影響を受けません。このような方程式は、第5章の定理4により、直ちに構成要素の級数展開が可能であり、各係数は1である。

以下、第四定理。

6. 方程式系において、平方化または他のどの過程でも方程式を形として簡略化し、その発展において示されるすべての構成要素が正の係数を持つ場合、導出された方程式（消去によって得られるもの）もまた同じ特性を持ち、他の方程式と加算して組み合わせることができます。

任意の方程式 $V = 0$ の符号を含む項の全展開において、すべての構成要素が負の係数を持つことがない場合を考えます。その展開は次のように書くことができます：

\[
V_{1} x + V_{0}(1-x) = 0
\]

ここで $V_{1}$ と $V_{0}$ はそれぞれ次の形を持ちます：

\[
a_{1}t_{1} + a_{2}t_{2} + \cdots + a_{n}t_{n},
\]

$t_{1}, t_{2}, \ldots, t_{n}$ は他の符号の構成要素であり、$a_{1}, a_{2}, \ldots, a_{n}$ は各々正または0の量です。消去の結果は

\[
V_{1}V_{2} = 0
\]

となります。そして $V_{1}$ と $V_{2}$ の係数がすべて負でないため、積 $V_{1}V_{2}$ も負の係数を持つことはありません。したがって、構成要素の係数が正である他の方程式にこの方程式を加えることができ、結果の方程式は元の方程式から得られた全意味を持ちます。

以下、第五定理。

7. 既存の命題から論理を表す方程式系の実用的な則や方法を導き出す。

前節の調査により、以下の点が確立された：

1. 方程式 $V = 0$ の形で、$V$ が双対法則 $V(1 - V) = 0$ を満たすとき、これらの方程式は単純な加算によって組み合わせることができる。

2. 他の形式の方程式 $V = 0$ も平方化という過程により、単純な加算による組み合わせが適用可能な形に簡略化できる。

したがって、実際の論理表現におけるどの方程式が前者の分類に属し、どの方程式が後者の分類に属するかを決定することが残されている。

第四章の結論で述べた一般型の命題は次のようである。それらを表す命題の分類は以下の通り：

1. 主辞が普遍的であり、述語が特定的な命題。
その象徴的な型（IV. 15）は

\[
X = vY,
\]

$X$ と $Y$ が双対法則を満たす。$v$ を消去すると、

\setcounter{equation}{0}
\begin{equation}
X(1 - Y) = 0,
\end{equation}

となり、これは同じ法則を満たすことも確認できる。平方化によるさらなる簡略化は必要ない。

2. 両項が普遍的であり、その象徴的な型は

\[
X = Y,
\]

$X$ と $Y$ がそれぞれ双対法則を満たす。方程式を $X - Y = 0$ の形に書き換えて平方化すると、

\begin{eqnarray}
X - 2XY + Y = 0,\nonumber\\
\textrm{または } X(1 - Y) + Y(1 - X) = 0.
\end{eqnarray}

この方程式の最初の項は、その形から明らかに双対法則を満たす。

同じ方程式を得る別の方法も存在する。方程式

\[
X = Y
\]

は、

\[
X = vY, \quad Y = vX,
\]

と等価である（$X$ が $Y$ と同一視されることは、すべての $X$ が $Y$ であり、すべての $Y$ が $X$ であることを意味する）。これらの方程式を $v$ を消去することで、

\[
X(1 - Y) = 0, \quad Y(1 - X) = 0,
\]

となり、これらを加えると (2) の結果を得る。

3番目の命題は、両方の項が特別な場合のものである。そのような命題の形式は以下のようになります：
\[
vX = vY,
\]
ただし$v$は完全に任意ではなく、従って消去することはできません。$v$は「ある」という代表を表しますが、これは意味において「すべて」を含むかもしれませんが、「何もなし」を含みません。したがって、第二項を第一項の側に移動し、結果の等式を規則に従って平方する必要があります。
その結果は明らかに以下のようになります：
\[
vX(1 - Y) + vY(l - X) = 0.
\]
これらの結論は、将来的の方向のためにルールとして要約することが便利かもしれません。

8. \textbf{ルール}。---等式が次の典型的な形式で項$X$と$Y$に従う法則を遵守するように表現されている場合、以下の等式を変換します：
\[
X = vY \quad \text{を} \quad X(1 - Y) = 0,
\]
\[
X = Y \quad \text{を} \quad X(1 - Y) + Y(1 - X) = 0.
\]
\[
vX = vY \quad \text{を} \quad vX(1 - Y) + vY(1 - X) = 0.
\]

任意の等式が形$X = 0$で与えられている場合、変換は必要ありません。また、形$X = 1$で現れる任意の等式は、上の変換の2つ目から明らかに$1 - X = 0$に置き換えられます。

これらのシステムの等式がこのような形に簡略化された後、それらまたはその過程によって導かれた任意の等式を加算することにより結合することができます。

9. \textbf{注}。---第四章で見たように、命題の項を実際の意味に関心なく記号の言語に直して翻訳すると、項$X$と$Y$が法則を遵守しない等式を作り出すことがあります。参照された章の第3節の(3)番目の命題で与えられた等式$w = st(p + r)$はそのようなものであり、このような等式は確かに意味を持っています。それらが好奇心のためにもしくは他の何らかの理由のために使用されるべきである場合、ルール（VI. 5）により簡略化することが最善です。

10. 例題2。---次のような初等幾何学の命題を考えてみましょう：

1番目の命題。相似な図形は、対応する角が等しく、対応する辺が比例しているすべての図形から成るものです。

2番目の前提は、対応する角が等しい三角形は、その対応する辺の比も等しく、逆もまた然りである。

これらの前提を表現するために、次のように定義しましょう：
\begin{tabular}{c l l}
$s$ &= &相似性. \\
$t$ &= &三角形. \\
$q$ &= &対応する角が等しい. \\
$r$ &= &対応する辺の比も等しい. \\
\end{tabular}

すると、これらの前提は以下の式で表現されます：
\setcounter{equation}{0}
\begin{eqnarray}
s & = & qr, \\
tq & = & tr.
\end{eqnarray}

これらの方程式を規則に従って簡約化するか、つまり同じ結果を得るためには、これらの方程式の項を最初の側に移し、各方程式を平方し、それらを加えると、以下の式が得られます：
\begin{equation}
s + qr - 2qrs + tq + tr - 2tqr = 0.
\end{equation}

非相似な図形についての説明を求めます。これは三角形という要素からなる図形で、対応する角が等しく、対応する辺の比も等しいものから形成されるものです。

(3) から得られる式は以下の通りです：
\begin{eqnarray}
s = \frac{tq + qr + rt - 2tqr}{2qr - 1}, \nonumber \\
\therefore 1-s = \frac{qr - tq - rt + 2tqr - 1}{2qr - 1}.
\end{eqnarray}

二つ目の項を完全に展開すると、以下の式が得られます：
\begin{eqnarray}
1 - s = 0tqr + 2tq(1 - r) + 2tr(1 - q) + t(1 - q)(1 - r) \nonumber \\
+ 0(1 - t)qr + (1 - t)q(1 -r) + (1 - t)r(1 - q)\nonumber \\
+ (1 - t)(1 - q)(1 - r).
\end{eqnarray}

上記の展開において、係数が2である2つの項は規則により0と等しいとします。その後、係数が0である項を省略すると、以下の式が得られます：
\begin{eqnarray}
1-s=t(1-q)(1-r)+(1-t)q(1-r)+(1-t)r(1-q)\nonumber\\
+(1-t)(1-q)(1-r); \\
tq(1-r)=0; \\
tr(1-q)=0;
\end{eqnarray}

これらの式の直接的な解釈は以下の通りです：

1. 非相似な図形は、対応する角が等しくないかつ辺の比も等しくない三角形と、三角形でない図形（その角が等しく且つ辺の比が等しくない、または対応する辺の比が等しく且つ角が等しくない、またはどちらの条件も満たさない）から成る。

2. 対応する角が等しくかつ辺の比が等しくない三角形は存在しない。

3番目の定理。対応する辺が比例しているが、角度は等しくない三角形は存在しない。

11. この最終方程式の直截的な解釈は次のとおりである。一般理論に従って見ると、特定のクラスの物体（異なる図形）を、元の前提の一部要素を使って説明する際に得られる独立関係も同時に得られることがわかる。そして、この情報が直接的調査の対象に関連している点でも余分な情報ではないことは、容易に示すことができる。例えば、独立した関係は常に必要に応じて、求められた関係をより簡潔な形で表現するために使用できる。たとえば、(7)を
\[
0 = t q (1 - r)
\]
という形に書き換えて(6)に加えると、
\begin{eqnarray*}
t (1 - q) (1 - r) + t q (1 - r) &=& t (1 - r),\\
1 - s &=& t (1 - r) + (1 - t) q (1 - r) + (1 - t) r (1 - q) \\
&+& (1 - t) (1 - q) (1 - r),
\end{eqnarray*}
を得る。この式を解釈すると、異なる図形の説明の最初の項として「対応する辺が比例していない三角形」ではなく、「対応する辺が比例していない三角形」というより簡潔な表現が得られることがわかる。便利性に配慮してこのような簡略化を行うべきかどうかは、常に便宜性に基づいて判断されるべきである。

12. 常に有利な簡略化（VII. 15）は、求められた直截的な説明の項を可能な限り少ないグループに集約することである。たとえば、(6)の第二項の第三および第四項を加えると単項 $(1 - t)(1 - q)$ となる。この簡略化を最後のものと組み合わせると
\[
1 - s = t (1 - r) + (1 - t) q (1 - r) + (1 - t) (1 - q)
\]
を得る。これを解釈すると、
\textit{異なる図形は対応する辺が比例していないすべての三角形と、三角形でないすべての図形（その図形が三角形でない場合）において、対応する角度が等しくないか、または対応する角度が等しいが対応する辺が比例していない図形から成る。}

一般解の完全さは、余分なものではなく、その全貌を提供することで有用である。
我々が求めるすべての情報を与えつつも、最適な形でその情報を表現する手段も提供してくれるからである。

13. 既に述べた原理を説明するもう一つの観察について述べる必要がある。式(5)における $1 - s$ の完全展開において、二つの項は係数が2であり、$\frac{1}{0}$ ではない。後で示すように、この事実は二つの前提が独立していなかったことを示している。これを確認するために、これらの前提の方程式を簡約形に戻して再検討する。

\begin{eqnarray*}
s(1-qr) + qr(1-s) &= 0,\\
tq(1-r) + tr(1-q) &= 0.
\end{eqnarray*}

これらの方程式の最初の項が共通の要素を持っている場合、それらを乗算することでその要素が現れる。これを実行すると、

\[
stq(1 - r) + str(1 - q) = 0
\]

を得る。

従って、

\[
stq (1 - r) = 0, \quad str (1 - q) = 0
\]

という方程式が得られ、これらは原始的な方程式から導かれる方程式である。これらの解釈は次の通りである：

相似な三角形は対応する角が等しいとき、その対応する辺の比も等しい。

相似な三角形は対応する辺の比が等しいとき、その対応する角も等しい。

これらの結論はそれぞれの前提のみを用いて導くことができる。この点で、定義によれば、前提は独立していない。

14. 最後に、この章最初の方法で説明した問題について再検討し、現在の方法を使って性質 $C$ の存在が性質 $A$ と $B$ に対するどのような結論を導くかを確認する。

シンボル $v$ を消去することで次の方程式を得た：

\begin{eqnarray}
xy(wz + \bar{w}\bar{z}) = 0, \\
yz (x\bar{w} + \bar{x}w) = 0,\\
\bar{x}\bar{y} = \bar{w}\bar{z}.
\end{eqnarray}

これらの式から、$w$ を消去し、$z$ を決定する必要があります。既に (1) と (2) は条件 $V(1 - V) = 0$ を満たしています。第三の方程式を項を第一辺に移項し、平方すると、

\[
\bar{x}\bar{y}(1 - \bar{w}\bar{z}) + \bar{w}\bar{z}(1 - \bar{x}\bar{y}) = 0
\]

を得ます。次に (1)、(2)、(4) を加えると、

\[
xy(wz + \bar{w}\bar{z}) + yz(x\bar{w} + \bar{x}w) + \bar{x}\bar{y}(1 - \bar{w}\bar{z}) + \bar{w}\bar{z}(1 - \bar{x}\bar{y}) = 0
\]

となります。$w$ を消去すると、

\[
(xyz + yz\bar{x} + \bar{x}\bar{y}) \{xy\bar{z} + yzx + \bar{x}\bar{y}z + \bar{z}(1 - \bar{x}\bar{y})\} = 0
\]

を得ます。次に、第二因子の項を第一因子の項順に乗じて計算し、

\[
x\bar{x} = 0, \quad y\bar{y} = 0, \quad z\bar{z} = 0
\]

からほとんど全てが消え、残るのは

\[
xyz + \bar{x}\bar{y}z = 0;
\]

より、

\[
z = \frac{0}{xy + \bar{x}\bar{y}}
= 0 \cdot xy + \frac{0}{0} \cdot x\bar{y} + \frac{0}{0} \cdot \bar{x}y + 0 \cdot \bar{x}\bar{y}
= \frac{0}{0} \cdot x\bar{y} + \frac{0}{0} \cdot \bar{x}y
\]

となり、

\[
\text{性質 $C$ が見つかった場合、性質 $A$ や性質 $B$ のいずれかは共に見つかるが、両方同時に存在しない。}
\]

方程式 (5) から、任意定数乗法の方法で以前の調査で得られた結果を容易に導くことができますし、$x, y, z$ 間の関係の他の提案形も導くことができます。例えば、

\[
\text{性質 $B$ が欠けている場合、性質 $A$ と $C$ は共に存在するか、または $C$ は欠ける。}
\]

逆に、

\[
\text{性質 $A$ と $C$ が共に存在する場合、$B$ は欠ける。}
\]

この結論の逆は、$\bar{y}$ の展開値において $xz$ に 1 が係数として含まれていることから成り立っています。

\newpage

\chapter{記号計算の省略法・簡便法（Methods of Abbreviation）}
\label{ch:ix}

\begin{center}
\large\textit{対称性・同値変形・論理式の計算量削減アルゴリズム}
\end{center}
\vspace{4mm}

\#\#\# 第1章 簡略化の方法

\#\#\#\# 1. 前章で確立され、説明された3つの基本的な展開法、消去法、簡約法は論理上のすべての実用目的に十分ですが、それらには特定の場合において重要な程度で簡略化が可能であり、特に消去法に関してはそのような場合があります。これらについては本章で注意を喚起したいと思います。まず、前章の方法による簡略化の原理を含むいくつかの命題を証明し、その後、具体的な例に適用します。

任意の項が基本法則 \( V(1-V) = 0 \) を満たす場合、その項は個々に構成要素となります。しかし、これらが同時に現れるときには、展開項のように必ずしも同じシンボルを含むわけではありません。例えば、\( ax + bxy + cyz \) は3つのクラス項 \( x \), \( xy \), および \( yz \) をそれぞれ係数 \( a \), \( b \), \( c \) で乗じた式と説明できます。次の2つの命題に適用される原理は、不要なクラス項の排除であり、それが最終結果の構成要素を増やすことのないクラス項が無駄であると見做されます。

\#\#\#\# 命題I

任意の方程式 \( V = 0 \) において、\( V \) が正の係数を持つクラス項の系列で構成されている場合、その中から他の項を因子とする項を排除し、すべての正の係数を1に変更することが許されます。

この級数の正項の意味は、その最終展開における成分の数と性質にのみ依存し、それらが含むすべての符号に対する展開に依存するだけであり、実際の係数の値には一切依存しない（VI. 5）。いま、$x$ を級数の任意の項とし、$xy$ を $x$ の因子を持つ別の任意の項とする。$x$ を符号 $x$ および $y$ について展開すると、

\[
xy + x(1 - y)
\]

となり、項 $x$ および $xy$ の和を展開すると、

\[
2xy + x(1 - y)
\]

となる。

しかし、先に述べた通り、これらの式が方程式の左辺において 0 である方程式の係数すべてが正であるとき、その両辺における表現は等価であり、最終的な展開にはただ二つの成分 $xy$ および $x(1 - y)$ のみが存在するから、結果的には次のような方程式が生じる

\[
xy = 0, \quad x(1 - y) = 0.
\]

したがって、項の合計 $x + xy$ は単なる項 $x$ に置き換えることができる。

この論理は命題で考慮されるすべての場合にも適用される。たとえば、項 $x$ が重複する場合、合計 $2x$ は $x$ に置き換えることができ、その状況では最終的な簡約において方程式 $x = 0$ が現れるからである。

\begin{center}
\textsc{命題 II.}
\end{center}

3. 除去法の過程で二つの因子をそれぞれ正項のみで構成され、論理符号の基本法則を満たすもの同士で乗算する必要があるとき、両方の因子から共通の項を排除したり、一方の因子から他方の因子に含まれる項で割り切れる項を排除することは許される。ただし、その除外された項は結果の因子の積に加えられることが条件である。

この命題の文において、「割り切れる」という言葉は便宜上の用語であり、代数的な意味において $xy$ および $x(1 - y)$ が $x$ で割り切れるという意味である。

この命題の意味をより明確にするため、乗算する因子が \( x + y + z \) と \( x + yw + t \) であるとする。すると、これらの2つの因子から \( x \) の項を排除し、第2の因子から \( yw \) の項を排除できるという主張がある。ただし、これらの項は最終的な積に移されるとする。

したがって、結果の因子は \( y + z \) と \( t \) となり、これら2つの因子の積 \( yt + zt \) に \( x \) の項と \( yw \) の項を加えると、

\[
x + yw + yt + zt
\]

が与えられる。これは与えられた因子 \( x + y + z \) と \( x + yw + t \) の積に等しい表現であり、具体的には消去過程において等しい。

まず、2つの因子が共通項 \( x \) を持つ場合を考える。このとき、因子をそれぞれ \( x + P \), \( x + Q \) と表し、一方では \( P \)，他方では \( Q \) を \( x \) 以外の正の項の和とする。

すると、

\[
(x + P)(x + Q) = x + xP + xQ + PQ
\]

となる。しかし、消去過程は特定の因子を乗算し、その結果を0に等しいとすることである。したがって、上記方程式の右辺が0に等しくなるか、またはそのような表現の因子が存在する。

前者の場合を選んだ場合、最後の命題によれば、\( xP \) および \( xQ \) の項は別の項 \( x \) の因子を持つ正の項であるため、排除できる。結果の式は

\[
x + PQ
\]

となり、両方の因子から \( x \) を排除し、残った因子の積に加えることで得られる式と同じになる。

後者の場合を選んだ場合、(1) 式の右辺が最終的な消去結果にどのように影響するかは、その成分の数と性質により決まる。これらの成分は \( xP \) および \( xQ \) の項を排除した際に変化しない。なぜなら、\( x \) の展開には \( x \) の因子であるすべての可能性のある成分が含まれるからである。

最後に、1つの因子が他方の因子で割り切れる項 \( xy \) を含む場合を考える。

$x + P$ と $xy + Q$ を因子とする。すると、

\[
(x+P)(xy+Q) = xy+xQ+xyP+PQ.
\]

しかし、前命題の論理によれば、項 $xyP$ は別の正の項 $xy$ の因子を含むため、これを省略できる。従って、

\begin{eqnarray*}
xy + xQ + PQ\\
= xy + (x + P) Q.
\end{eqnarray*}

この式は第二因子から項 $xy$ を排除し、その値を最終的な積に移すことを表している。したがって、命題は明らかである。

\begin{center}
\textsc{命題 III.}
\end{center}

4. 任意の符号 $t$ が何らかの方程式系から他の符号を消去した結果において保持される場合、その消去結果は次のように表すことができる：

\[
Et + E(1-t) = 0,
\]

ここで $E$ は、提案された方程式系において $t = 1$ とし、同じ他の符号を消去したものを用いる。また $E'$ は、提案された方程式系において $t = 0$ とし、同じ他の符号を消去したものを用いる。

証明：$\phi(t) = 0$ を最終的な結果とする。この式を展開すると、

\[
\phi(1)t + \phi(0)(1-t) = 0.
\]

任意の過程で $\phi(t)$ を提案された方程式系から導き出す場合、同じ過程で $t = 1$ に変更したとき $\phi(1)$ を導き出し、また同じ過程で $t = 0$ に変更したとき $\phi(0)$ を導き出す。従って、命題の真偽は明らかである。

5. 最後に証明された3つの命題について、厳密な理論の開発や応用には全く必要でないものの、実用的な目的を達成する重要な役割があることに注意しなければならない。命題 1 により加算結果の簡略化が可能となり、命題 2 により乗算結果の簡略化が可能となる。また命題 3 により複雑な消去過程は二つの独立した過程に分割され、通常はより単純な形式となる。この方法は最終的な目的が $t$ の値を他の符号で表す場合によく用いられる。

6. 例1。アリストテレスの『ニコマケー Ethics』2巻、第3章では、行為が徳を有するとは、その自己に特定の性格を持つことではなく、実行者がそれを知的かつ慎重な選好で、目的自体として、そして固い行動原理に基づいて実行することによってある特定の心状態を持っていることを意味すると定義した後、次の2章で、徳は情動、能力、習慣のいずれかに属するのか、それとも他の関連する点とともに、その問題を検討します。

\begin{center}
前提。
\end{center}

1. 徳は情動（〔ギリシャ語原典引用〕）、能力（〔ギリシャ語原典引用〕）または習慣（〔ギリシャ語原典引用〕）のいずれかである。

2. 情動は、称賛や非難の対象ではなく、慎重な選好を伴うものでもない。

3. 能力は、称賛や非難の対象ではなく、慎重な選好を伴うものでもない。

4. 徳は、称賛や非難の対象であり、慎重な選好を伴うものである。

5. 何かの技術または科学がその作品を良好な状態に保つためには、極端を避けて、人間性（〔ギリシャ語原典引用〕$\tilde{\alpha}\varsigma$）に対する中庸を見据えることである。

6. 徳はあらゆる科学や芸術よりも正確で優れています。
これは論理的により強い議論です。科学と真の芸術が欠点と過度性を避けるなら、徳は均衡の線に沿って調和を追求するべきであることはもっともです。科学や芸術が作業を良好な状態にするなら、徳がその特有の作業を「良好な状態」にする理由があるはずです。最終的な前提を次のように解釈しましょう。称賛や非難について言及しないことにします。これらの言及は意図的な優先性のみに伴うものであり、これを保持することを目的としています。従って、次の代表符号を仮定できます：

\[
\begin{array}{cccccc}
& v = 徳. \\
& p = 感情. \\
& f = 能力. \\
& h = 習慣. \\
& d = 意図的な優先性が伴うもの. \\
& g = 作業を良好な状態にするもの. \\
& m = 人間の性質に対する均衡を見据えるもの.
\end{array}
\]

ここで、$q$ を不定なクラス符号として用いることにします。前提は次の等式で表現されます：

\begin{eqnarray*}
v &=& q \left\{ p (1 - f) (1 - h) + f (1 - p) (1 - h) + h (1 - p) (1 - f) \right\}. \\
p &=& q (1 - d). \\
f &=& q (1 - d). \\
v &=& qd. \\
g &=& qm. \\
v &=& qg.
\end{eqnarray*}

これらの式から符号 $q$ を消去すると、次のような等式が得られます：

\begin{eqnarray}
v \left\{ 1 - p (1 - f) (1 - h) - f (1 - p) (1 - h) - h (1 - p) (1 - f) \right\} = 0. \\
pd = 0. \\
fd = 0. \\
v (1 - d) = 0. \\
g (1 - m) = 0. \\
v (1 - g) = 0.
\end{eqnarray}

次に、(2), (3), および (4) の符号 $d$ を消去し、$p$, $f$, および $h$ との関係で $v$ を決定します。まず、(2), (3), および (4) を加算すると、

\[
(p + f)d + v(1 - d) = 0.
\]

この式から符号 $d$ を通常の方法で消去すると、

\[
(p + f)v = 0.
\]

これを (1) に加えると、$v$ を決定すると、

\[
v = \frac{0}{p + f + 1 - p(1 - f)(1 - h) - f(1 - p)(1 - h) - h(1 - f)(1 - p)}.
\]

これを展開すると、

\[
v = \frac{0}{0}h (1 - f) (1 - p).
\]

この等式的解釈は：「徳は習慣であり、能力や感情ではない」です。

次に、元の方程式系から $f$, $p$, および $g$ を消去し、その後で $v$ を $h$, $d$, および $m$ との関係で決定することを考えます。この場合、$p$ と $f$ を一緒に消去します。式 (1), (2), および (3) の和を取ると、

\begin{eqnarray*}
v \left\{ 1 - p(1-f)(1-h) - f(1-p)(1-h) - h(1-p)(1-f) \right\} \\
+ pd + fd = 0.
\end{eqnarray*}

この式を $p$ と $f$ について展開すると、

\begin{eqnarray*}
(v + 2d)pf + (vh + d)p(1-f) + (vh + d)(1-p)f \\
+ v(1-h)(1-p)(1-f) = 0.
\end{eqnarray*}

したがって、消去の結果は

\[
(v + 2d)(vh + d)(vh + d)v(1 - h) = 0
\]

となります。ここで $v + 2d = v + d + d$ であり、これは命題 I. により $v + d$ に簡約されます。この結果と二番目の因子の積は

\[
(v + d)(vh + d)
\]

となり、これは命題 II. により $d + v(vh)$ または $vh + d$ に簡約されます。

同様にして、この結果を三番目の因子で掛けると単純に $vh + d$ となります。最後に、この結果を四番目の因子 $v(1-h)$ で掛けると、最終的な方程式

\begin{equation}
vd(l - h) = 0
\end{equation}

が得られます。

次に、式 (5) および (6) から $g$ を消去します。その結果は

\begin{equation}
v(1-m) = 0
\end{equation}

となります。

最後に、方程式 (4), (8), および (9) の和を取ると、

\[
v(1-d) + vd(1-h) + v(1-m) = 0
\]

となり、

\[
v = \frac{0}{1 - d + d(1-h) + 1 - m}
\]

が得られます。この結果の展開は

\[
v = \frac{0}{0}hdm
\]

となります。これは、徳は人間性に対する平均的な意図に基づく意志決定と連携する習慣であるという意味です。

適切に言えば、これは定義ではなく、徳の説明である。しかし、前提から正確に導き出すことができるものは、これだけである。アリストテレスは、この点において慎重さ（節度）が安全で中庸な行動の線を決定する必要性と結びつけており、疑いの余地なく、従来の徳の理論はより知的な性格を持っていたことが多かったと言えるだろう。徳は単なる良心の優越性や道徳的な容易さではなく、全体の心の正しい状態と習慣にこそ含まれていると考えられた。そして、これらの理論にはある程度の正しさがあった。実際、無条件の良心の命令への服従は徳ある行動の重要な要素であるが、それらの命令が不変の正義（$\alpha\iota\acute{\omega}\nu\iota\alpha\;\delta\acute{\iota}\kappa\alpha\iota\alpha$）と一致していること、すなわち物事の構成そのものが基礎となっているもしくはそれ自体が物事の構成の基礎となっている不変の正義原則に従っていることがまた一つの要素である。そして、この一致は、特に高い程度においては無知や精神的な退廃と矛盾していることが多い。

アリストテレスの特定の理論に戻ると、おそらく多くの人々はその性格が過度に否定的であると考えるだろうし、極端を避けることはより優れた能力の発揮の余地を十分には提供しないと言える。アリストテレス自身もこの欠点について不完全な意識があったと思われる。彼の第七書の冒頭で「英雄的な徳」（$\tau\acute{\epsilon}\nu\;\upsilon\acute{e}\mu\;\alpha\mu\iota\;\eta\acute{a}\rho\omega\acute{i}\kappa\acute{o}\nu$）という言葉を使い、人間の性質を超えるものとして述べていたからである（ニコマコス倫理学書7）。

7. 既に述べたことがある（VIII.1）が、単独の方程式や命題に関する理論は、その完全な解答を得るためには、方程式系の理論と結びつけて考察する必要があります。この指摘は、与えられた単独の方程式から、ある単純な基本クラスではなく、より複雑なクラス（その表現に2つの要素以上のものが含まれている）との関係を決定することを求められる場合に現れます。以下に特定の例とその後の一般化された問題が挙げられます。

例題2.---私たちは第7章で使用した財富の定義の象徴的な表現を再考しましょう、すなわち
\[
w = st \left\{p + r\left(1 - p\right)\right\},
\]
ここで、以前と同じように \\
\begin{tabular}{c l}
&w = 財富,\\
&s = 供給が制限されたもの,\\
&t = 移転可能なもの,\\
&p = 楽しみを生むもの,\\
&r = 痛みを防ぐもの;
\end{tabular}\\
そして、移転可能なものと楽しみを生むものの関係を、他の定義要素（すなわち財富、供給が制限されたもの、痛みを防ぐもの）にどのように結びつけるかを求めることとする。

移転可能なものと楽しみを生むものの表現は $tp$ である。これを新しい符号 $y$ で表すことにしよう。すると、方程式
\begin{eqnarray*}
w &= st \left\{p + r\left(1 - p\right)\right\},\\
y &= tp,
\end{eqnarray*}
を得る。ここで $t$ と $p$ を消去することで、$y$ を $w$, $s$, および $r$ の関数として決定できる。その結果を解釈すると、求めた関係が得られる。

これらの式の項を第一項に移動させると、以下のようになります：

\begin{eqnarray}
w - stp - str (1 - p) = 0, \nonumber \\
y - tp = 0.
\end{eqnarray}

これらの方程式の平方和を計算すると、

\begin{equation}
w + stp + str (1-p) - 2wstp - 2wstr(1-p) + y + tp - 2ytp = 0
\end{equation}

となります。

$t$と$p$に関して最初の項を展開し、これらの符号を消去するために、

\begin{eqnarray}
(w+s-2ws+1-y)tp + (w+sr-2wsr+y)t(1-p) \nonumber \\
+(w+y)(1-t)p + (w+y)(1-t)(1-p)
\end{eqnarray}

となります。そして、$t$と$p$を消去する結果は、

これらの係数の積を0に等しいものとして得られます。

または、Prop. 3により、上記の方程式から$t$と$p$を消去した結果は以下の形になります：

\[
Ey + E'(1-y)
\]

ここで $E$ は与えられた方程式において $y$ を $1$ に置き換えて $t$ と $p$ を消去した結果、そして $E'$ は同じ方程式において $y$ を $0$ に置き換えて $t$ と $p$ を消去した結果です。それぞれの場合における $t$ と $p$ の消去方法は、$tp$, $t(1-p)$, その他の4つの成分の係数を乗算することです。

$y = 1$ とすると、係数は以下のようになります：

1. $w(1-s) + s(1-w)$
2. $1 + w(1-sr) + s(1-w)r$, これはProp. Iにより1に等しい。
3. および4. $1 + w$, これはProp. Iにより1に等しい。

したがって、$E$ の値は

\[
w(1-s) + s(1-w)
\]

となります。また、$y = 0$ とすると、

1. $1 + w(1-s) + s(1-w)$, これは1に等しい。
2. $w(1-sr) + sr(1-w)$
3. および4. $w$

これらの係数の積は

\[
E' = w(1-sr)
\]

となります。

方程式から $y$ を決定するための式は次のようになります：
\begin{eqnarray*}
\left\{w(1-s)+s(1-w)\right\}y + w(1-sr)(1-y)=0,\\
\therefore y=\frac{w(1-sr)}{w(1-sr)-w(1-s)-s(1-w)};
\end{eqnarray*}
そして、第二項を展開すると：
\begin{eqnarray*}
y&=\frac{0}{0}wsr+ws(1-r)+\frac{1}{0}w(1-s)r+\frac{1}{0}w(1-s)(1-r)\\
&\quad+0(1-w)sr+0(1-w)s(1-r)+\frac{0}{0}(1-w)(1-s)r\\
&\quad+\frac{0}{0}(1-w)(1-s)(1-r);
\end{eqnarray*}
これより、式を簡略化すると：
\begin{eqnarray}
y=ws(1-r) + \frac{0}{0}wsr + \frac{0}{0}(1-w)(1-s),\\
\textrm{with }w (1-s) = 0.
\end{eqnarray}

この解釈は次の通りです：

1. 移転可能で楽を生む財産とは、供給が制限されているが痛みを防ぐものすべて、供給が制限されていないが痛みを防ぐ無限量のもの、そして財産ではないかつ供給が制限されないものすべてである。

2. すべての財産は供給が制限されている。

上記の解法では、括弧内に（付随する独立関係7によって暗黙的に含まれている）全説明の一部を書き込みました。

8. 次の問題はより一般的であり、最後の問題のような問題に対する簡単な実用的なルールを提供します。

\begin{center}
\textsc{一般問題}.
\end{center}

与えられた符号 $x$, $y..w$, $z..$ を結びつける任意の方程式について、

残りの符号 $w$, $z$, 等 で表される任意のクラスを、符号 $x$, $y..$ で論理的に表現する方法を求めること。

まず、符号 $x$, $y$ と符号 $w$, $z$ の2つずつの場合に限定します。これは一般的なルールを決定するために十分です。

与えられた方程式 $V=0$ について、クラスの表記が $\phi (x,y)$ で表されるとする。

$\phi (x,y)=t$ と仮定すると、上記の2つの方程式から $x$ と $y$ を消去する必要がある。

したがって、方程式 $V=0$ は次の形で展開できます：
\setcounter{equation}{0}
\begin{equation}
Axy + Bx(1-y) + C(1-x)y + D(1-x)(1-y) = 0,
\end{equation}

ここで、$A$, $B$, $C$, $D$ は符号 $w$ と $z$ の関数です。

また、$\phi (x,y)$ がクラスまたは事物の集合を表すのであれば、それは1の係数を持つ構成要素やその系列でなければなりません。

従って、$\phi(x, y)$ の完全な展開が以下の形で表されるとき、

\[
axy + bx(1 - y) + c(1 - x)y + d(1 - x)(1 - y),
\]

係数 $a$, $b$, $c$, $d$ はそれぞれ 1 または 0 である必要があります。

次に、方程式 $t = \phi(x, y)$ を移項し、平方して、

\[
t{1-\phi(x,y)} + \phi(x,y)(1-t) = 0;
\]

とします。この式を $x$ および $y$ について展開すると、

\begin{eqnarray*}
{t(1 -a ) + a (1 - t)} xy + {t(1 - b) + b(1 - t)} x(1 - y) \\
+ {t(1-c) + c(1-t)} (1-x)y\\
+ {t(1-d) + d(1-t)} (1-x) (1-y) = 0;
\end{eqnarray*}

これを式（1）に加えると、

\begin{eqnarray*}
{A + t(1-a) + a(1-t)}xy\\
+ {B + t(l-b) + b(l-t)} x(1-y) + 等 = 0.
\end{eqnarray*}

$x$ および $y$ を消去した結果を、

\[
Et+E' (1-t)=0,
\]

という形とするとき、$E$ は既述により、上記展開の係数が $t = 1$ のときに得られる約分された積であり、$E'$ は同様に $t = 0$ の条件で約分された同じ因子の積である。

したがって、

\[
E = (A+1-a)(B+1-b)(C+1-c)(D+1-d).
\]

この式の任意の因子、例えば $A + 1 - a$ を考えると、$a = 1$ のとき $A$ になり、$a = 0$ のとき $1 + A$ となり、これは命題 I.により 1 に約分される。したがって、$E$ は $\phi(x, y)$ の展開の係数が 1 である構成要素の係数の積であると推測できる。

また、

\[
E' = (A+a)(B+b)(C+c)(D+d).
\]

この式の任意の因子、例えば $A + a$ を考えると、$a = 0$ のとき $A$ になり、$a = 1$ のとき 1 に約分される。他の因子についても同様である。したがって、$E'$ は $\phi(x, y)$ の展開の係数が 0 である構成要素の係数の積である。

これらの状況を総合すると、以下のルールを確立できる：

9. \textit{論理方程式から与えられた符号 $x$, $y$ 等の任意の組み合わせによって表されるクラスと、その他の符号によって表されるクラスとの関係を導く}。

\#\#\# 原則

与えられた方程式を $x$, $y$ に関する項について展開し、次のように方程式を作り出す：
$$Et + E'(1-t) = 0,$$
ここで $E$ は、上記の展開において与えられたクラスの表現における係数が 1 の各成分の係数の積であり、$E'$ は、同様に、その係数が 0 の各成分の係数の積である。方程式から $t$ を解き、解釈することで求められる表現を得る。

\#\#\# 注

この原則の証明において $V$ は正の項のみで構成されていると仮定されますが、容易に示すことができるように、この条件は不要であり、原則は一般化され、与えられた方程式の準備が実際には必要ないことが明らかになります。

10. 例3

同じ財産の定義が与えられている場合、$t(1 - p)$（移転可能だが楽しまないもの）を他の要素 $w$, $s$, $r$ の形で表す式を求めよ。

与えられた方程式
$$w - stp - str (1 - p) = 0,$$

を平方すると、
$$w + stp + str (1 - p) - 2wstp - 2wstr (1 - p) = 0;$$

これを $t$ と $p$ について展開すると、

\begin{eqnarray*}
(w + s - 2ws) tp + (w + sr - 2wsr) t (1 - p) + w (1 -t)p\\
+ w(1-t)(1-p) = 0.
\end{eqnarray*}

これらの係数を以下のように表示するのが最も適切である：
\begin{eqnarray*}
(w(1-s)+s(1-w))tp+(w(1-sr)+sr(1-w))t(1-p)+w(1-t)p\\
+w(1-t)(1-p)=0.
\end{eqnarray*}

関数 $t(1-p)$ を定めるものとして $z$ と表すことにします；
すると、$t$ および $p$ に関する $z$ の完全な展開は次のようになります。
\[
z = 0 \cdot tp + t (1 - p) + 0 \cdot (1 - t) p + 0 \cdot (1 - t) (1 - p).
\]
したがって、最後の問題により
\begin{eqnarray*}
E z + E' (1 - z) &=& 0;\\
\textrm{ここで } E &=& w (1 - sr) + sr (1 - w);\\
E' &=& \left( w (1 - s) + s (1 - w) \right) \times w \times w = w (1 - s); \\
\therefore \quad w (1 - sr) + sr (1 - w) z + w (1 - s) (1 - z) &=& 0.
\end{eqnarray*}
したがって，
\begin{eqnarray*}
z &=& \frac{w (1 - s)}{2wsr - ws - sr}\\
&=& \frac{0}{0} wsr + 0 w(1 - r) + \frac{1}{0} w(1 - s)r + \frac{1}{0} w(1 - s)(1 - r)\\
&+& 0 (1 - w) sr + \frac{0}{0} (1 - w)s(1 - r) + \frac{0}{0} (1 - w)(1 - s)r \\
&+& \frac{0}{0} (1 - w)(1 - s)(1 - r).\\
\textrm{または, } z &=& \frac{0}{0} wsr + \frac{0}{0} (1 - w)s(1 - r) + \frac{0}{0} (1 - w)(1 - s), \\
\textrm{ここで } w (1 - s) &=& 0.
\end{eqnarray*}

したがって、\textit{移転可能で生産的でないものは、または財（供給が限られているもので、苦痛を防ぐもの）；または財ではないが、供給が限られているもので、苦痛を防がないもの；または財ではないが、供給が無制限なものである。}

類似の方法で導き出された以下の結果は容易に検証できます：

\textit{供給が限られているもので生産的でないものは移転不可能である。}

\textit{生産的でない財は、移転可能で、供給が限られているもので、苦痛を防ぐものである。}

\textit{供給が限られているもので、財または生産的であるが両方ではないものは、移転可能で苦痛を防ぐものまたは移転不可能なものである。}

11. 自然史の領域から多くの興味深い例が選択されることが可能である。しかし、私はそのような応用は独立した価値を持つとは考えない。これらは例えば動物学における分類の真の原理に光を当てるものではない。これらの発見のために必要なのは、肯定的な知識の基礎であり、器官構造や目的論的適応に関する理解である。これは外部観察と分析の手段を使用することからしか得られない種類の知識である。

ただし、自然史におけるどのような命題の集合を取っても、分類体系に関係なく多くの論理的な問題が浮かび上がる。おそらくこのような例を作る際には、そのクラスや種の名前によって即座に想起される性質について言及することは不要であり、例えば動物界の環状構造を持つアーネリダ類の場合、土壌 worm とヒゲ worm の一種である。

4. 例題。
1. アーネリダは柔らかい体を持ち、裸で覆われていないか、または管に包まれている。

2. アーネリダは赤血球を持つ二重循環血管を持つすべての無脊椎動物から成る。

\begin{tabular}{c l}
$a$ と仮定する &アーネリダ； \\
$s$ と仮定する &柔らかい体を持つ動物； \\
$n$ と仮定する &裸である； \\
$t$ と仮定する &管に包まれている； \\
$i$ と仮定する &無脊椎動物であり、赤血球を有する 等；
\end{tabular}

すると与えられた命題は以下の式で表現される。

\begin{eqnarray}
a &=& vs (n(1-t) + t(1-n))； \\
a &=& ir；
\end{eqnarray}

さらに、以下が暗黙の条件として加わる：

\begin{equation}
nt = 0。
\end{equation}

$v$ を消去し、システムを単一の方程式に簡略化すると、

\begin{equation}
a [1 - sn(1-t) - st(1-n)] + a (1-ir) + ir (1-a) + nt = 0。
\end{equation}

もし赤血球を持つ、外部被覆を持つ、および脊椎を持つという要素に対して、柔らかい体を持ち管に包まれている動物が（前提に基づいて）どの関係にあるかを得たいと仮定するなら、

まず $a$ を消去しなければならない。その結果は

\[
ir(1 - sn(1-t) - st(1-n)) + nt = 0.
\]

次に、(IX.9)で $s$ と $t$ について展開し、最初の係数を命題1により簡略化すると、

\begin{equation}
nst + ir(1-n)s(1-t) + (ir+n)(1-s)t + ir(1-s)(1-t) = 0.
\end{equation}

したがって、$st=w$ のとき、

\[
nw + ir(1-n) \times (ir+n) \times ir(1-w) = 0;
\]

または、

\begin{eqnarray*}
nw + ir(1-n)(1-w) &=& 0; \\
\therefore w &=& \frac{ir(1-n)}{ir(1-n)-n} \\
&=& 0 \cdot irn + ir(1-n) + 0 \cdot i(1-r)n + \frac{0}{0}i(1-r)(1-n) \\
&+& 0 \cdot (1-i)rn + \frac{0}{0}(1-i)r(1-n) + 0 \cdot (1-i)(1-r)n \\
&+& \frac{0}{0}(1-i)(1-r)(1-n); \\
\therefore w &=& ir(1-n) + \frac{0}{0}i(1-r)(1-n) + \frac{0}{0}(1-i)(1-n).
\end{eqnarray*}

したがって、軟体動物は管状に包まれており、赤血を有し裸でないすべての無脊椎動物であり、さらに赤血を有しない裸でない無脊椎動物と脊椎動物（裸でない）の不確定な部分から成る。

同様の方法により、次のような簡略化された方程式が導き出され、その解釈は読者に任せます。これらは(5)の展開から得られました。

\begin{eqnarray*}
s(1-t) &=& irn + \frac{0}{0}i(1-n) + \frac{0}{0}(1-i); \\
(1-s)t &=& \frac{0}{0}(1-i)r(1-n) + \frac{0}{0}(1-r)(1-n); \\
(1-s)(1-t) &=& \frac{0}{0}i(1-r) + \frac{0}{0}(1-i).
\end{eqnarray*}

上記の例では、私は特にこの方法の力をどのように示すかという特定の方法を提示することを目指していません。その力は、確率論の数学理論と結びついたときにのみ完全に示されるでしょう。しかし、読者がこの点について正確な見解を持つことを望むなら、通常の論理の規則により次の問題を検討し、その後その解決策を見ることをお勧めします。同時に覚えておくべきは、どんな複雑さであっても、その方法による解決がより煩雑になるだけであり、それほど確実には達成可能であるということです。

例題5．ある自然現象の観察によって得られた以下の一般結果を想定します。

まず、これらの生成物のいずれにおいても、性質 $A$ と $C$ が欠如している場合、性質 $E$ は存在し、一方で性質 $B$ または $D$ のいずれかが存在しますが、両方とも存在しない。

次に、これらの生成物のいずれにおいても、性質 $A$ と $D$ が存在しつつ性質 $E$ が欠如している場合、性質 $B$ および $C$ はどちらも存在するか、またはどちらも欠如します。

さらに、これらの生成物のいずれにおいても、性質 $A$ が性質 $B$ または $E$ または両方に共存する場合、性質 $C$ または性質 $D$ のいずれかが存在しますが、両方とも存在しない。逆に、これらの生成物のいずれにおいても、性質 $C$ または $D$ が単独で存在する場合、性質 $A$ は性質 $B$ または $E$ または両方に共存します。

次に、性質 $A$ の存在から、性質 $B$, $C$, および $D$ に対する結論を導き出すことが求められます。また、独立して性質 $B$, $C$, および $D$ の間に関係が存在するかどうかも確認します。さらに、同様の方法で性質 $B$ と性質 $A$, $C$, および $D$ に対する結論を導き出すことも求められます。

ここで注意すべきは、各データにおいて、性質 $A$, $B$, $C$, および $D$ の情報が別の要素 $E$ と共に複雑に絡み合っている点です。この要素 $E$ についての結論を述べる必要はありません。したがって、与えられた命題を表す方程式系から符号 $v$ を消去する必要があります。

性質 $A$ を $x$, 性質 $B$ を $y$, 性質 $C$ を $z$, 性質 $D$ を $w$, そして性質 $E$ を $v$ と表記します。データは以下のようになります。
\begin{align}
(1 - x)(1 - z) &= q v (y (1 - w) + w (1 - y));\\
(1 - v) x w &= q (yz + (1 - y)(1 - z));\\
xy + xv (1 - y) &= w (1 - z) + z (1 - w);
\end{align}

ここで $\bar{x} = 1 - x$, 等, および $q$ は不定のクラス符号です。まず、最初と第二の方程式から $q$ をそれぞれ消去し、その結果を第三の方程式（5）章で示した方法によって簡略化したものを加算すると、

\begin{align}
(1 - x)(1 - z) (1 - v y (1 - w) - v w (1 - y)) + (1 - v) x w (y (1 - z) + z (1 - y)) \\
+ (xy + xv (1 - y)) (wz + (1 - w)(1 - z)) + (w (1 - z) + z (1 - w)) (1 - xy - xv (1 - y)) = 0.
\end{align}

この方程式から$v$を消去し、$x$の値を求めることになります。これを達成するためには、現在の章の命題3の方法を使用するのが便利です。

まず、消去後の結果を次の方程式で表します：

\[
Ex + E' (1-x) = 0.
\]

$x = 1$ を(4)の最初の項に代入して $E$ を求めます。すると、

\[
\bar{v}w(y\bar{z} + z\bar{y}) + (y + v\bar{y})(wz + \bar{w}\bar{z}) + (w\bar{z} + z\bar{w})\bar{v}\bar{y}
\]

となります。$v$ を消去すると、

\[
(wz + \bar{w}\bar{z}) \left\{ w(y\bar{z} + z\bar{y}) + y(wz + \bar{w}\bar{z}) + \bar{y}(w\bar{z} + z\bar{w}) \right\}
\]

となります。実際の乗算を行い、条件 $w\bar{w} = 0$, $z\bar{z} = 0$ 等に従って計算すると、

\[
E = wz + y\bar{w}\bar{z}
\]

となります。

次に、$x = 0$ を(4)に代入して $E'$ を求めます。すると、

\[
z(1 - vy\bar{w} - v\bar{y}w) + w\bar{z} + z\bar{w}
\]

となります。$v$ を消去し、命題1と2を使用して結果を簡略化すると、

\[
E' = w\bar{z} + z\bar{w} + \bar{y}\bar{w}\bar{z}
\]

となり、したがって最終的に

\begin{equation}
(wz + y\bar{w}\bar{z})x + (w\bar{z} + z\bar{w} + \bar{y}\bar{w}\bar{z})(1-x) = 0;
\end{equation}

から、

\[
x =
\frac{w\bar{z} + z\bar{w} + \bar{y}\bar{w}\bar{z}}
{w\bar{z} + z\bar{w} + \bar{y}\bar{w}\bar{z} - wz - y\bar{w}\bar{z}}
\]

となります。これを展開すると、

\begin{eqnarray*}
x = 0yzw + yz\bar{w} + y\bar{z}w + 0y\bar{z}\bar{w}\\
  + 0\bar{y}zw + \bar{y}z\bar{w} + \bar{y}\bar{z}w + \bar{y}\bar{z}\bar{x};
\end{eqnarray*}

または、項を縦に並べると、

\begin{equation}
x = z\bar{w} + \bar{z}w + \bar{y}\bar{z}\bar{w};
\end{equation}

この解釈は次の通りです：

「$A$ の性質が見つかる任意の物質において、$C$ の性質または $D$ の性質が見つかるか、それらの両方ではないか、または $B$, $C$, および $D$ のすべての性質が欠けているかのいずれかが見つかります。逆に、$C$ の性質または $D$ の性質が単独で見つかるか、または $B$, $C$, および $D$ のすべての性質が一緒に欠けている場所では、$A$ の性質が見つかるでしょう。」

また、$B$, $C$, および $D$ の性質には独立な関係は存在しないことが明らかです。

次に、$y$ を求めます。これには(5)をその符号について展開します。

\[
(xwz + x\bar{w}\bar{z} + \bar{x}w\bar{z} + \bar{x}z\bar{w})y + (xwz + \bar{x}w\bar{z} + \bar{x}z\bar{w} + \bar{x}\bar{z}\bar{w})\bar{y} = 0;
\]

したがって、以前の方法で進めて、

\[
y = \bar{x}\bar{w}\bar{z} + \frac{0}{0}(\bar{x}wz + xw\bar{z} + xz\bar{w});
\]

次に、

\begin{eqnarray}
xzw &= 0;\\
\bar{x}z\bar{z}w &= 0;\\
\bar{x}z\bar{w} &= 0;
\end{eqnarray}

(10) を解いて簡略化すると、

\[
\bar{x}z = \frac{0}{0}w;
\]

より、独立な関係は以下のようになります：—「もし性質 $A$ が欠如し、かつ性質 $C$ が存在するならば、性質 $D$ は存在する。」また、加算と解消により (8) および (9) を追加すると、

\[
xz + \bar{x}\bar{z} = \frac{0}{0} \bar{w}.
\]

したがって、一般の解と残る独立な関係は以下のようになります：

1. 「もし生産物中のいずれかに性質 $B$ が存在するならば、性質 $A$, $C$, および $D$ のうち少なくとも一つは欠如しており、あるいはそれらのうちの一つだけが欠如している。」逆に、「これら全てが欠如する場合、性質 $A$ が存在することが結論づけられる」と (7)。

2. 「もし性質 $A$ と $C$ の両方が存在するか欠如するならば、性質 $D$ は欠如し、これは性質 $B$ の存在または欠如とは独立である。」(8) および (9)。

これらの結論を検証することは試みませんでした。

\newpage

\chapter{完全なる論理体系の条件（Conditions of a Perfect Method）}
\label{ch:x}

\begin{center}
\large\textit{記号的手法の普遍性・形式的一貫性と論理的必然性}
\end{center}
\vspace{4mm}

\#\# 完璧な方法の条件

\#\#\# 1. 主題：基本原理についての議論が長く行われ、我々は次に重要な原理についての考慮に入ります。これらの二つの科学論理学の大分野間の移行期は、過去の進歩のいくつかのステップを振り返りながら、そのような主題において完璧な方法とは何かを問い直す機会となるかもしれません。ここで私は、効果性だけではなく、適切さと美しさに基づく概念も含めて、我々がこれまで占めてきた主題における方法の完璧性について述べます。精密な分析を行うことで、以下のような結論に至る可能性が高いでしょう：完璧な方法は、その設計目的を達成するための効率的なものであるだけでなく、全体として各部分と過程において一定の統一と調和を示すべきです。この概念が最も完全に現れるのは、方法の形式自体が基礎となる原理を想起し、可能であれば一つの基本原理を想起することです。これらの考慮事項を論理学の科学に適用する際には、単なる分析過程だけでなく、推論の形態や方法だけでなく、導出を行うためのデータまたは前提の体系についても最善とは何かを問い直すことが適切かもしれません。

2. 簡単な力においては、第8章で開発された最初の方法は、その適切な範囲内では完璧なものであることが疑いようがない。任意定数の導入により、前提形式への依存や、それらを表現する方程式間の条件から独立することが可能となる。しかし、これは外的な要素を導入することになり、同じ章で示された第2の削減法と比較して、より労働集約的かつ洗練度が低い解法形態であるように思われる。しかしながら、ある条件下では後者の方法はそれ以外とは異なるより完璧な形式を呈する。方程式

\[
x(1-x) = 0,
\]

によって表される基本的な条件を普遍的な形式として採用すれば、両方の過程と結果に一貫性が与えられ、そうでなければ得られないものとなる。方法が簡潔さや便利さのみを価値ある質素とするなら、そのような原理の導入は利点を生むことはないだろう。なぜなら、上記の特徴を各解法ステップに強制することは、場合によっては予め準備された簡略化作業に大きな労力を伴うことがあるからである。しかし、このことが可能であることを知ることは興味深いし、このような解法形態が自然と現れる条件についても知ることは重要な点でもある。これらの点のいくつかについては本章で考慮する。

\begin{center}
\textsc{命題I.}
\end{center}

3. \textit{論理符号間の任意の方程式を形式 $V=0$ に簡約し、その中で $V$ が対偶法則を満足するようにする}

\[
V(1 - V) = 0.
\]

第5章の命題4においては、上記条件は $V$ が構成要素の和であるときには満たされることを示している。また、第6章の命題2から明らかにすることは、方程式が移項によって $V=0$ の形に簡約され、最初の項の展開により同じ系列の構成要素と非零係数を生み出すとき、これらはすべて等価であるという点である。これらの特定の数値的な係数の具体的な値は重要ではない。

次に述べる命題の目的は、常に方程式のすべての項を第一辺へ移動させ、そのメンバーを完全に展開し、結果において消え去らない係数全てを1に変更することで達成される。ただし、既にその値を持つ係数は例外とする。

しかし、多くの符号を含む関数の展開は、表現が非常に長いため不便であるため、実際には遭遇する唯一のケースでのこの複雑さを避ける方法を示すことが望ましい。

論理系を適切に表した方程式の主要な基本形は既に第8章で議論されている。それらは以下のようである：

\begin{eqnarray*}
X &=& vY, \\
X &=& Y, \\
vX &=& vY.
\end{eqnarray*}

条件 $X(1-X) = 0, Y(1-Y) = 0$ を満たすとき、上記の最初の2つの方程式は以下の形へと変形される：

\begin{eqnarray}
X(1-Y) &=& 0, \\
X(1-Y) + Y(1-X) &=& 0;
\end{eqnarray}

また、同じ条件下では最後の方程式は以下のように変形されることが示される：

\begin{equation}
v \left( X(1-Y) + Y(1-X) \right) = 0;
\end{equation}

これらの結果は明らかに最初の項において条件 $V(1-V) = 0$ を満たす。

以上は、論理系の方程式が適切に表されるとき実際に現れる形と条件であり、上記の方法により常にその法則に従うように簡略化することが可能である。ただし、個々の方程式はこの法則を満たしてもよいが、それらの等価な合計（第8章4節）は必ずしもそうではないため、その合計にも必要な条件を課す方法を示さなければならない。

したがって、システムの各簡略化された方程式を加えた方程式を以下のように表す：

\begin{equation}
v + v' + v'' + \cdots = 0,
\end{equation}

この方程式は元のシステムから得られた单一等価な方程式である。ここで $v, v', v'', \ldots$ をクラス項（第9章1節）とし、以下の条件を満たすとする：

\[
v(1-v) = 0, \quad v'(1-v') = 0, \quad \cdots
\]

（4）の完全な解釈は、その最初の項が含むすべての基本的な符号$x$, $y$，etc.について展開することで得られます。その後、係数が消えない全ての成分を0と等しくします；すなわち、$v$, $v'$, $v''$，etc.に見られる全ての成分です。これらの成分は以下の3つから成ります：1) $v$に見られるもの；2) $v$には見られないが$v'$に見られるもの；3) $v$と$v'$には見られないが$v''$に見られるもの、そしてそのようなものなどが続きます。したがって、これらの成分は以下の式
\[
v + (1 - v)v' + (1 - v)(1 - v)v''+ \cdots,
\]
に見られるようなものとなります。この式には重複のない成分しか含まれず、明らかに$V(1-V)=0$という法則を満たします。

例えば、以下の表現を考えます：
\[
(1 - t) + v + (1 - z) + tzw,
\]
ここで項$(1-t)$と$(1-z)$は括弧で囲ってあり、これらが単一のクラスの項として取り扱われるべきであることを示しています。この場合、（5）に従って、条件$V(1-V)=0$を満たす式へと簡略化するために、最初の項後の全ての項を$t$で乗じ、次に二番目の項後の全ての項を$(1-v)$で乗じ、最後に残った項を$z$で乗じます。その結果は以下のようになります：
\[
1 - t + tv + t(1 - v)(1 - z) + t(1 - v)zw.
\]

4. すべての論理方程式は、次のように簡約されるべきである：
$$V=0$$,
ここで $$V$$ は対称法則を満たしている。しかし、このような形式で常に方程式が現れることがなければ、より高い段階の完璧さがあるだろう。それはあらかじめ何らの準備もせずに方程式がそのような形を呈示し、かつ、システム全体の方程式を単一の相当形に簡約するための追加手順を行った後でも、その形式を保持し続けることがなければ、明らかに完璧さは高まらない。これらの特徴が自然と現れる理由は、方法論上の欠陥ではなく、前提が常に完全で正確かつ独立的であるとは限らないという事実の結果である。前提が完全ではない場合、それは具体的な関係（形式的な関係とは対照的に）を表していない材料的な関係を含んでいる場合である。前提が正確でない場合、それは意図しない関係を含んでいる場合である。しかし、これらの点については、現在の状況ではあまり考慮する必要がないので、それらが独立的でないという意味での欠如について考えてみよう。

5. ある命題の体系が独立であるとは、その体系のどの部分からも他の部分から導き出される結論を導き出すことが不可能なときを指すことができる。これらの命題を表す方程式をすべて形として
\[
V=0
\]
に簡約したと仮定すると、上記の条件は特定の関数 $V$ の展開において現れる任意の構成要素が、同じ体系の他の関数 $V'$ の展開において現れることのないことを意味する。この条件が満たされない場合、体系の方程式は独立ではない。これは様々なケースで起こる可能性がある。

すべての方程式が最初の項において対称法則を満足すると仮定し、ある一方の方程式の展開に正の項 $x$ が現れ、他方の方程式の展開に項 $xy$ が現れている場合、これらの方程式は独立ではない。なぜなら、項 $x$ はさらに $xy + x(1 - y)$ と展開でき、体系の両方の方程式から
\[
xy=0
\]
という方程式が導かれるためである。

また、一方の方程式に項 $xy$ が現れ、他方の方程式に項 $xz$ が現れている場合、これらは共通の構成要素 $xyz$ をもたらすように展開できる。このような他のケースも容易に想像できる。最初見れば独立に思える前提であっても実際にはそうではない場合がある。

このような形の方程式 $V = 0$ が真に独立でない場合、個々にそれぞれ対称法則
\[
V(1 - V) = 0
\]
を満足していても、それらを加えた等価な方程式はその条件を満たさない可能性がある。一方、体系の方程式がすべて上記の法則を満足し、互いに独立である場合、それらの和も同じ法則を満足する。

これらの点については、必要以上に詳しく述べたが、私は将来の努力の目標と指南针となる理想的な完璧さの模範を作り出すこと、そしてすべての研究においてそれを常に前に持つことが重要だと考えている。現在の研究分野では、方程式の変形をいかに簡潔性を保ちつつも容易にするかという手法改善の主な目標は、上記の条件が普遍的に適用されることを目指すべきである。

これらの点について詳しく述べたのは、私は将来の努力の模範となる理想的な完璧さを作り出すこと、そしてすべての研究においてそれを常に前に持つことが重要だと考えているからである。現在の研究分野では、方程式の変形をいかに簡潔性を保ちつつも容易にするかという手法改善の主な目標は、上記の条件が普遍的に適用されることを目指すべきである。

この主題に関連して、以下の命題は注意に値します。

\begin{center}
\textsc{命題 II.}
\end{center}

任意の方程式 $V = 0$ の最初の項が条件 $V(1-V) = 0$ を満たし、その方程式中の任意の符号 $t$ の表現が他の符号の展開関数として決定されている場合、展開係数はしかして $1$, $0$, $\displaystyle\frac{0}{0}$, $\displaystyle\frac{1}{0}$ の形をとるだけである。

なぜなら、方程式を $t$ に関して展開すると、

\begin{equation}
Et + E'(1 - t),
\end{equation}

を得ます。ここで $E$ および $E'$ は、$t$ を順に $1$ と $0$ に置き換えたときに $V$ がどのようになるかによって定義されます。したがって $E$ および $E'$ は以下の条件を満たします。

\begin{equation}
E(1 - E) = 0,\quad E'(1 - E') = 0.
\end{equation}

今、式 (1) より、

\[
t = \frac{E'}{E' - E},
\]

右辺は他の符号の関数として展開されます。係数を計算する際に $E$ および $E'$ が受け取る唯一の数值値は $1$ と $0$ のみであることが明白です。従って以下の状況のみが生じ得ます：

1. $E' = 1$, $E = 1$ のとき、$\displaystyle\frac{E'}{E'-E} = \frac{1}{0}$。

2. $E' = 1$, $E = 0$ のとき、$\displaystyle\frac{E'}{E'-E} = 1$。

3. $E' = 0$, $E = 1$ のとき、$\displaystyle\frac{E'}{E'-E} = 0$。

4. $E' = 0$, $E = 0$ のとき、$\displaystyle\frac{E'}{E'-E} = \frac{0}{0}$。

以上より命題の真偽は明らかです。

6. 方程式の解法において、形式 \(1\)、\(0\)、および \(\frac{0}{0}\) は条件 \(V(1 - V) = 0\) に直接関連付けられることなく現れることがあります。しかし、係数 \(\frac{1}{0}\) についてはそうではありません。上記の条件が満たされるときに関連付けられる項は、その条件が満たされない場合でも、\(1\)、\(0\)、および \(\frac{0}{0}\) のいずれか以外の値を受けることができます。この係数が付く任意の展開の係数が上記の四形式のいずれにも現れない場合は、\(\frac{1}{0}\) に変更することが許されます。これは、その係数が \(0\) と等しいことを示す適切な符号としてこれを用いることで一貫性を高めることができます。この方法は私はしばしば採用します。

\centerline{\textsc{命題 III}}

7. \textit{任意の符号 \(x, y, 等\) を含む方程式 \(V = 0\) の中で、第一項が二重性の法則を恒等的に満足する場合、すなわち
\[
V(1 - V) = 0,
\]
から任意の符号 \(x, y, 等\) を消去した結果は、他の符号に関して展開した方程式の各成分の係数が \(1\) であるものを合計し、それを \(0\) と等しいことで得られます。}

与えられた方程式 \(V = 0\) が三つの符号 \(x, y, t\) を含むと仮定します。そして \(x\) および \(y\) を消去することを考えます。\(t\) に関して方程式を展開すると、
\setcounter{equation}{0}
\begin{equation}
At + B(1 - t) = 0,
\end{equation}

ここで \(A\) と \(B\) は符号 \(t\) から独立しています。

第九章の命題3によれば、\(x\) および \(y\) を消去した結果は以下の形になります：
\begin{equation}
Et + E'(1 - t) = 0,
\end{equation}

ここで \(E\) は方程式 \(A = 0\) の符号 \(x, y\) を消去した結果、\(E'\) は方程式 \(B = 0\) の符号 \(x, y\) を消去した結果です。

そして \(A\) および \(B\) には以下の条件が満たされる必要があります：
\[
A(1 - A) = 0, \quad B(1 - B) = 0
\]

したがって、$A$（現在はこの係数についてのみ論じる）は0であるか、1であるか、または構成要素の一つであり、あるいは$x$と$y$を含む一部の構成要素の和である。もし$A = 0$であるなら、明らかに$E = 0$となる；もし$A$が単一の構成要素であるか、または$x$と$y$を含む一部の構成要素の和であれば、$E$も0となる。なぜなら、$A$の完全な展開（$x$と$y$に関する）には係数が消滅する項が含まれており、$E$はすべての係数の積であるからである。したがって、$A = 1$の場合、$E$は$A$に等しくなるが、他の場合では$E$は0となる。同様に、$B = 1$の場合、$E$も$B$に等しくなり、他の場合では$E$は消滅する。したがって、式（2）は、もし存在すれば、式（1）のうちで$A$と$B$の係数が1である部分を含む。この論理は一般化できる。例えば、$V$が$x, y, z, t$の符号を含み、$x$と$y$を消去する必要がある場合を考える。もし$V$を$z$と$t$に関する展開すると、

\[
zt + xz(1 - t) + y (1 - z) t + (1 - z) (1 - t)
\]

となるなら、求めたい結果は

\[
zt + (1 - z) (1 - t) = 0
\]

となり、これは係数が1である部分の展開である。

したがって、任意の方程式系から単一の等価な方程式$V = 0$を導き出し、この方程式が条件

\[
V(1 - V) = 0
\]

を満たすなら、通常の消去法は完全に不要となり、一方の展開過程がその役割を果たすことができる。

8. この方法には実用的な利点があるかどうかは不明だが、理論的には一貫性と完備性があり、それゆえ注意を払う価値がある。したがって、将来的な章（第XIV章）ではこの方法の説明に専念する。応用力学の進歩は他の例も提供しており、問題系や方程式のシステムを支配する中心的な法則のもとで統一される。

9. これまでの議論から明らかであるように、ある特定の命題類には上記の準備方法の特別な道具が不要である。それは以下の条件によって特徴付けられるものである：

まず、これらの命題は共通の種類であり、コプラ「is」または「are」の使用によって暗示されるもので、述語が個別のものである。

次に、命題の項は、その項を表す要素間の理解された関係の仮定なしでも理解可能である。

さらに、これらの命題は独立している。

もしこのような推測が完全には無駄ではないとすれば、私たちは言語を表現や思考の道具として使用する際に行われるべき条件を予想することができます。不動の存在が単に意図することを述べる場合でも、一方で省略なく、他方で繰り返しも行わないように。これらの条件は、完璧な言語の概念と正確な方法の過程との両面から見ても、学習者の注意に値するものである。

\newpage

\chapter{二次命題とその記号表現（Of Secondary Propositions）}
\label{ch:xi}

\begin{center}
\large\textit{時間・条件・仮言命題のブール代数的定式化}
\end{center}
\vspace{4mm}

次章は「二次命題」について以及その記号的表現の原理について述べます。

{\large 二次命題およびその記号的表達の原理に関するもの}

1. 第四章で既に確立された教義によれば、論理的な命題は二大類のいずれかに属することができる。すなわち、原始的命題と派生的命題である。これらの最初の一類については、本著作の前半の章ですでに議論が行われており、現在私たちはその他の命題について、すなわち真偽として認識される他の命題に関連する命題を考慮することに進むことになる。我々が入る調査は、その全体的な順序と進行において、既に行われたものと類似しているだろう。これらの二つの調査は、認識される思考の対象が異なる点で異なり、形式的かつ科学的な法則を明らかにするか、またはそれらの法則に基づく方法や過程が異なる点で異なるわけではない。確率はそのような結果を期待する一端を持ち得るだろう。我々が知る自然の均一性と、我々が信じる自然の創造者の不変の恒久性とは一致し、このような高度な能力を持つ心において、周囲の情景との対話だけでなく、自己に対する知識と自己の構成法則についての反省においても、より実在的な調和と均一性が存在するべきであると考えることができる。このような予想は我々の研究の最初の基準にはならないし、それらは我々が特定の調査領域に提出されたものの中での実際的な構造を確定するために行う耐え忍びの研究から逸れるようなことはない。しかし、類似性の基礎が適切かつ独立して決定されると、その類似性を Meditation の対象とし、その範囲を追跡し、まだ明らかでない真実の暗示を受け入れることは、純粋に科学的な目的においても不矛盾である。事実への最終的な訴えはこのように排除されることなく、類推の使用はその適切な領域を超えて拡大されるわけではなく、独立した調査が検証しなければならない関係を提示するか、または拒否させることになる。

2. 二次的な命題とは、命題が真偽として考えられるときにその関係を扱う命題である。私たちは事物の関係を表現するために一次的な命題を使用することができる。しかし、命題自体についても思考の対象とし、それらに対する判断を表すこともできる。これらの判断の表現が二次的な命題となる。どのような知識の程度でも、その命題が真か偽かのいずれか一方の主張ができるような命題は存在しない。それぞれのこれらの主張も二次的な命題である。「太陽が照っているのは真実である」、「惑星が自らの光で輝くとは真実ではない」など、これらはその種類の例である。前者では「太陽が照っている」という命題が真であると主張される。後者では「惑星が自らの光で輝いている」という命題が偽であると主張される。二次的な命題には、命題間の関係や依存性を表現するすべての判断も含まれる。このクラスまたは分類に属すものは条件付きの命題、「もし太陽が照っているなら、その日は晴れる」などがある。また、ほとんどの排他的命題、「太陽が照るか、事業が延期されるか」などもある。前者では「その日は晴れる」という命題の真偽が「太陽が照っている」という命題の真偽に依存していることを表現する。後者では二つの命題、「太陽が照る」「事業が延期される」間の関係を表現し、一方が真である場合、他方が真ではないことを意味する。同じクラスの二次的な命題には、命題の同時の真偽を主張するすべての命題も含まれる。「「太陽が照る」と「旅程が延期される」両方とも真実ではない」などがある。我々が注意した区別要素は、しばしば同一の二次的な命題内で混在している可能性がある。それは「または」、「もし」によって表現される排他的要素と、「もし」によって表現される条件付き要素を含むことができる。さらに、関連する命題自体も複合的な構造を持つことができる。例えば、「太陽が照る」かつ「余裕があれば」、「事業を開始するか」「前段階の手順を取るか」いずれか一方である。この例では、多くの命題が適切で意味のある方法で結び付けられ、それらの間の明確な関係を表現している。したがって、我々の定義によれば、この組み合わせは二次的な命題となる。

これらの区別要素は、同一の二次的な命題内で混在することもある。それは「または」、「もし」によって表される排他的要素と、「もし」によって表される条件付き要素を含むことができる。さらに、関連する命題自体も複合的な構造を持つことができる。例えば、「太陽が照る」かつ「余裕があれば」、「事業を開始するか」「前段階の手順を取るか」いずれか一方である。この例では、多くの命題が適切で意味のある方法で結び付けられ、それらの間の明確な関係を表現している。したがって、我々の定義によれば、この組み合わせは二次的な命題となる。

次級命題の理論は、その多様な応用と、既に述べたように一次命題の理論との密接で調和した類似性から、注意深く研究する価値がある。これらの各点について、いくつか追加の観察をしたいと思う。

3. まず最初に述べたいのは、日常的な論理は、少なくとも一次命題よりも次級命題の形式において多く示されるということである。また、倫理学者や形而上学者の言説も、おそらく事物とその性質についてより頻繁に議論するものではなく、原則と仮定について、真実と真実間の相互関係について議論することが多いかもしれない。私たちの経験がまだ解決されていない倫理や社会に関する大問題に対する結論は、人間的な起源の制限を多面的に示している。そして、普遍的原理の存在は疑う余地がないとしても、その適用に関する当時の知識を含む部分公式は条件、例外、そして失敗にさらされている。

したがって、その主題自体から最も興味深い研究分野であるべき領域においても、私たちの実際的な知識の大部分は仮説的である。思慮深く哲学を論じる著者の思考形式が同じ形態を採用する傾向があったことは、後で明らかになるだろう。したがって、仮定的命題その他の次級命題の議論に対する一般的手法の導入は、これまで遭遇したよりももっと興味深い応用分野を開くことになるだろう。

4. 二次命題の理論についての議論は、一次命題の理論と密接な類似性を持つことから興味深いです。両者の場合で、心の作用に課せられる形式的な法則は表現において同一であることが明らかになります。これらの法則に基づく数学的過程もまた同一であり、したがって、前者の部分で調査された方法論は新しい適用分野でも引き続き有効です。しかしながら、方法論の法則と過程が変わらぬままではあるものの、解釈の規則は新たな条件に適応する必要があります。クラスという概念ではなく命題を用いることになり、クラスや個体との関係性ではなく、命題または事象間の関連性を考えることになります。

それでも、目的と解釈が異なる2つのシステムにおいても、貫くべき調和的な関係や類似性を見ることができます。これはまだ残る難題を克服する手助けとなりつつ、自身で興味深い研究対象となり、人間の能力構造の統一性を証明する結論的な証拠でもあります。

\begin{center}
\textsc{命題I.}
\end{center}

5. \textit{二次命題と時間との関連性について調査する。}

この問いに取り組む前に、二次命題と一次命題の類似性を結ぶ性質について明確に述べる必要があります。

一次命題は、ある宇宙内での事物間の関係を表現します。その宇宙は実際の宇宙の限界を超えるか、またはそれと一致するかもしれません。私たちの言及対象となる物質は、この宇宙の限界内で制限されます。これらの関係は本質的に \textit{主観的} です。与えられたクラスの一部または全部が別のクラスのメンバーである場合があります。一次命題が参照する主題とその間の表現される関係はすべて、上記の特性を持っています。

ただし、二次的な命題を扱うと、我々は新たなクラスの主題と関係に直面することになります。ここで扱う主題自体が命題であるため、質問が生じます：これらの主題も「物」だと見做し、前例に比喩的に類推して自己の宇宙に参照できるでしょうか？

また、これらの主題となる命題間の関係は共存する真偽の関係であり、本質的な等価性の関係ではありません。異なる二つの命題を結びつけるとき、「一つが他である」という主張ではなく、「Xが真である場合、Yも真である」や「XとYは共に真である」といった表現を使用します。

これらの関係について考えると、その全般的な意味の範囲を調査する必要はありません（これには原因論的な形而上学的問題が含まれる可能性があり、科学の適切な範囲を超えるかもしれません）；しかし、論理的推論のために十分な意味を見つけることが十分です。例えば、「Xが真である場合、Yも真である」という条件命題について検討しましょう。この命題の確実な意味は、Xが真となる時間はYが真となる時間と一致することです。これは共存の関係に過ぎませんが、必ずしもその命題の全般的な意味を網羅するわけではありません。しかし、この命題の表明には確かに含まれており、論理的推論のために十分な関係であることが確認できます。

日常生活の言語は、二次的な命題が時間という概念と密接に結びついているという見解を正当化します。したがって、一次的な命題の適用範囲を「ある」（some）という単語で制限する一方で、二次的な命題の適用範囲は「時々」（sometimes）という単語で制限されます。「時々不正が勝利を収める」という言葉は、「不正が現在勝利を収めている」という命題が真である時間があることを主張しているに等しいです。確かに、ある命題の真偽は特定の時期や状況に限定されるものではなく、常に真である命題があり、これらは「永遠の真理」と呼ばれます。

この区別は、プラトンとアリストテレスを読んだ人なら誰でも知っていることでしょう。特にアリストテレスによって用いられるこの区別は、科学における抽象的な真理（幾何学の命題など）が常に真であるのに対し、事象や現象の関係性が時々真で時々偽であるという条件的または観測上の真理を表すために使われます。しかし、両種類の命題が表現される言語形式は、時間という概念に共通の依存性を持っていることを示しています。一方では有限な期間に限定され、他方では永遠に延びている形で現れています。

6. 実際には、普通の論理では我々はしばしば、そのような時間という概念が言葉そのものに含まれていることに気づいていないかもしれません。しかし、その指摘こそが、我々はよく構築された言語の形式を用いて論理を行うが、その形式が建立されている理由に注意を向けずにいることを示すだけであることを示しています。現在の調査の進行過程で、この原則の例を見ることができます。私は時間という概念を利用して、二次的な命題の表現法則と、それらを表すシンボルの組合せ法則が同一の法則に従うことを決定します。しかし、これらの法則と形式が一度決定されれば、そのような時間という概念（私はそれが上記目的にとって不可欠であると考えています）は実用的に不要となります。その後、一般的な言語の形式から、ここであげられる思考のシンボル的な道具の類似した形式へ移行し、その過程を利用し、その結果を解釈することができるようになります。時間という概念について何ら意識することなく。

\begin{center}
\textsc{命題 II.}
\end{center}

7. 二次的な命題の表現法則を確立し、そのシンボルが一次的な命題の表現において使用される対応するシンボルと同じ組合せ法則に従うことを示す。

我々は、真偽について主張したいか、またはその間の関係を二階の命題として表現したい Elementary Propositions $X$, $Y$, $Z$ を大文字で表し、それらに関する心理的操作を小文字 $x$, $y$, $z$ で表すことにしましょう。ここで、$x$ は次のように解釈されます：$x$ は命題 $X$ が真となる時間の部分を我々の注意に固定するという心的行為を表します。そして、この意味は「$x$ は命題 $X$ が真となる時間を表す」と主張されるときには理解されます。

さらに、接続記号 $+$, $-$, $=$ 等を以下のように解釈します：$x+y$ は命題 $X$ と $Y$ がそれぞれ真となる時間の部分の合計を表し、これらの時間は完全に分離されています。同様に、$x-y$ は命題 $X$ が真となる時間の部分から命題 $Y$ が真となる時間の部分を取り除いた残りの時間を表します。また、$x=y$ は命題 $X$ と $Y$ が真となる時間の部分が同一であることを表します。$x$ を「命題 $X$ の代表符号」と呼びます。

以上の定義から、常に

\[
x + y = y + x
\]

が成り立つことがわかります。両辺は同じ時間を表すためです。

さらに、$xy$ を次のように解釈します：$y$ と $x$ それぞれを順に実行するという心的行為の全体を表します。つまり、以下の要素からなる心的過程全体を指します：
1. 命題 $Y$ が真となる時間の部分を選択する心理的操作。
2. 上記の時間の部分から命題 $X$ が真となる時間の部分を選択する心理的操作。

これらの連続的なプロセスの結果は、命題 $X$ と $Y$ がともに真となる時間の全体を我々の注意に固定することです。

この定義から、常に

\[
xy = yx
\]

が成り立つことがわかります。

我々が最初に命題$Y$が真である部分の時間を心の中で選び、その結果から命題$X$が真である部分の時間を選び、または最初に命題$X$が真である部分の時間を選び、その結果から命題$Y$が真である部分の時間を選びます。この両方の場合で最終的な結果は同じであり、それは命題$X$と$Y$がともに真である部分の時間です。

この論理を続けることで、符号$x, y, z, 等$の組合せ法則がここでの解釈によって与えられるものと、本文の最初の部分で与えられた解釈によって与えられるものの表し方において同一であることを証明することができます。この最終的な同一性の理由は明らかです。両者の場合でも、研究しているのは同じ能力、またはその組合せであり、これらの操作の本質的な特性は、これらが存在全体やすべての事象が現実化される時間全体に取り組んでいると仮定するか、あるいは少なくとも一部を含む主張、真理、命題に関連している一部の部分に取り組んでいるとしても変わりません。

したがって、上記の法則に加えて、第2章の(4)で述べた法則は、

\begin{equation}
x(y + z) = xy + xz;
\end{equation}

という表現を持ちます。また特に、二重性の基本法則（第2章の(2)）は、

\begin{equation}
x^2 = x, または, x (1 - x) = 0;
\end{equation}

という表現を持ちます。この法則は論理における思考体系と量の科学における思考体系を区別する役割を果たし、前者にはそれ自身が持つ完璧さと普遍性を与え、それらがそれほど持たない場合とは異なります。

8. さらに、この法則（3）（他の法則も同様に）は符号$0$と$1$によって満足されるため、これらの符号が現在の思考体系において解釈できるかどうかを以前のように問い立てることができます。以前に行なった論理と同じような推論により、これらは解釈可能であり、以下の2つの立場を支持することができます：

1. 二次的な命題において、$0$は時間の要素に対する「何もなし」を表す。

2. 同じ体系において、$1$は議論がいかなる形でも関連する全体の時間、すなわち宇宙を表す。

初等的な命題において、議論の宇宙はしばしば実際の事物の一部に限定される一方で、しばしばその宇宙と一致する場合もある。同様に、二次的な命題においても、議論の宇宙は単一日または瞬間まで限定されることがあるし、全体の時間まで含むこともある。最も字義通りに解釈すれば「永遠」である。実際、議論の性質によって明示的か暗黙的に制限が加えられない限り、二次的な命題におけるシンボル$1$の適切な解釈は「永遠」であり、初等的な体系においては実際に存在する宇宙を表す。

9. 代数符号$x$, $y$, $z$を命題の真理に適用するのではなく、それらを事象の発生に適用することも同様に適切である。実際、事象の発生はその発生を主張する命題の真偽を含み、逆にその命題の真偽は事象の発生を含む。符号$x$のこれらの意味は必然的に互いに関連している。問題の状況によって最も望ましい解釈として提示されるこれら2つの等価な解釈を使えることは便利である。必要に応じてこの自由を利用することにする。純粋論理の問題では、符号$x$, $y$などは前提において関係を表す基本的な命題を表すと考える。確率の数学理論については、先に示したように（I. 12）、論理に基づいており、この後半部分で扱う予定であるので、これらの単純な事象を表すために同じ符号を使用する。

\begin{center}
\textsc{命題 III.}
\end{center}

\#\#\# 10. 次等比誼の表現に対する一般の法則を導く

この命題から生じる様々な調査においては、既に完成された初等比誼に関する類似の調査と比較的正確な類推関係があるため、証明の完全さは必要ないでしょう。まず項の表現について考えます；次に、それらを結びつける命題の表現について考えます。

1 を時間全体の長さとし、 \( x \) を命題 \( X \) が真である部分の時間とします。すると \( 1 - x \) は命題 \( X \) が偽である部分の時間を表します。

また、 \( xy \) を命題 \( X \) と \( Y \) がともに真である部分の時間を表すとするとき、この観察と組み合わせることで以下の解釈を得ることができます：

項 \( x(1-y) \) は命題 \( X \) が真であり、命題 \( Y \) が偽である時間期間を表します。項 \( (1-x)(1-y) \) は命題 \( X \) と \( Y \) が同時に偽である時間期間を表します。

項 \( x(1-y) + y(1-x) \) は、命題 \( X \) が真であり、または命題 \( Y \) が真であるが両方とも真ではない時間期間を表します。これは、それぞれ単独で且つ排他的に真である時間の合計です。項 \( xy + (1-x)(1-y) \) は、命題 \( X \) と \( Y \) がどちらも真であるか、またはどちらも偽である時間期間を表します。

別の符号 \( z \) が現れた場合でも、同じ原理が適用されます。すなわち、 \( xyz \) は命題 \( X \), \( Y \), および \( Z \) が同時に真である時間の長さを表し、 \( (1-x)(1-y)(1-z) \) はこれらが同時に偽である時間の長さを表します。これらの項の合計は、それらが共に真か偽である時間の長さを表します。

上記例で示された解釈の一般的な原理には、さらなる説明やより明確な声明は必要ありません。

11. 現在、命題の表現法に関する法則を、その命題が現れる明確な場合に分け、展示し、研究することができます。しかし、最初に注意を向けるべき重要な原則があります。これまで定めた表現法の原理は完全に一般性を持ち、主張する真偽についての断言を、私たちの論述が参照する時間全体（それが無限の永遠であるか、開始と終了が明確に決定された一定の期間であるか、あるいは瞬間的な段階であるかに関係なく）のどの部分でも行うことができるようにします。しかし、人間の推理における実際の手続きでは、このような制限は通常適用されません。命題が真であるという主張をするとき、一般的にはその命題が論述が参照する時間全体を通じて真であることを意味することが多いです。そして、無条件の真偽について異なる断言が論理的証明の前提として一緒に述べられた場合、これらの断言は同じ時間の範囲に参照され、特定の制限された部分には参照されません。正確な科学の対象や領域である必要事項においては、それぞれの真偽に関する断言は「永遠の真理」を主張することができるのです。瞬間的な現象（例えば社会的状況）について推理を行う場合、各断言は現在の時間に即した限定的な参照を持つことができます、「今」という言葉で表現します。しかし、どちらの場合でも、逆に明示的に異なる時期を指す表現がない限り、それぞれの命題は同じ期間に関連しています。

それによって生じる状況は以下の通りです：

1. 「命題 $X$ が真である」という命題を表す場合

ここでは、論述が参照する時間範囲に限定され、命題 $X$ が真であることを表現することを求められます。命題 $X$ が真である期間は $x$ で表され、論述が参照する時間の長さは1で表されます。したがって、

\[
x = 1
\]

という式が求められるのです。

2. 「命題 $X$ が偽である」という命題を表す場合

時間の範囲について、私たちが議論する限界内において、命題$X$は偽であることを表明します；またはその限界内で、それが真となる部分の時間はありません。したがって、命題$X$が真となる部分の時間$x$は求められます。従って必要な方程式は、

\begin{equation}
x = 0.
\end{equation}

この結果はまた、命題$X$が偽である時間を表す式 $1-x$ を全体の時間に等しいと仮定することで得られます。これにより

\begin{eqnarray*}
1 - x &=& 1, \\
従って \quad x &=& 0.
\end{eqnarray*}

3つ目の、排他的命題「命題$X$が真であるまたは命題$Y$が真である」を表す方法は、命題$X$と命題$Y$が互いに排他的であることを含意します。つまり、そのうちの一つだけが真です。

命題$X$が真であり且つ命題$Y$が偽であるまたは命題$Y$が真であり且つ命題$X$が偽である時間は、式 $x(1-y) + y(1-x)$ で表されます。従って必要な方程式は、

\begin{equation}
x(1-y) + y(1-x) = 1
\end{equation}

です。

上記の命題において、「または」が絶対的な排他的性を持たないと仮定し、命題$X$と$Y$が同時に真である可能性を排除しない場合、上記方程式の第一項に $xy$ を加える必要があります。従って、

\begin{eqnarray}
xy + x(1-y) + (1-x)y &=& 1, \nonumber \\
従って \quad x + (1 - x)y &=& 1.
\end{eqnarray}

4つ目の、条件付き命題「もし命題$Y$が真であるなら、命題$X$も真である」を表す方法は、命題$Y$が真であるとき、命題$X$も真であるという事実から、命題$Y$が真となる時間は命題$X$が真となる時間と一致することを表現すれば十分です。つまり、全体の時間$x$の中で命題$X$が真となる部分の時間の一部であるということです。

命題$Y$が真となる時間は $y$ であり、命題$X$が真となる全体の時間は $x$ です。符号$v$を不定な時間を表すものとすると、$vx$は全体の時間$x$の中から不定な部分を表します。従って、

\[
y = vx
\]

命題の表現として。

12. ここで $v$ を無限の時間の符号とみなすとき、$vx$ は全体、または全体の一部、あるいは全くない時間を$x$の代表するものと理解されることが可能である。これらの意味のいずれかが実現するために、任意の符号$v$の特定な決定が必要となる。例えば、$v$が全体時間$x$を含む一定の時間で表すように決定された場合、$vx$は全体時間$x$を代表するものとなる。また、$v$が時間$x$の一部を含む一定の時間で表すように決定され、その時間が時間$x$の全量を満たさない場合、$vx$は時間$x$の一部を代表するものとなる。最後に、$v$が全く共通部分を持たない時間で表される場合、$vx$は0の値を持つことになり、「時間なし」または「決して」と等価となる。

注意すべき点として、命題「もし $Y$ が真であれば、$X$ も真である」には、$X$ と $Y$ のいずれか一方も真であるという主張は含まれていない。この命題は、$Y$ の真実性が $X$ の真実性の必要条件である場合（そのような場合は $v=1$）、または $Y$ が実現されたときに $X$ の真実性を保証する条件でありながら、$X$ が真であってもその条件が満たされない可能性がある場合（そのような場合は $v$は時間$x$の一部を含む符号である）、または全く $Y$ が偽である場合（そのような場合は $v$は時間$x$と全く共通部分を持たない符号である）に一致する可能性がある。これらのすべてのケースは、$v$が無限の時間の符号であるという一般的な仮定の中に含まれている。

5. 条件付きと排他的論理和の特性を含む命題を表現する。

一般形の条件命題は「$Y$ が真なら $X$ が真」であり、その表現は前節で述べたように $y = vx$ です。このように、基本的な命題における慣習と同様に、$Y$ および $X$ を条件命題の項として指定することが適切です。また、通常の論理の言葉遣いを採用し、 particle「if」が付いた項 $Y$ を「前提」と呼び、項 $X$ を「結論」と呼ぶことにします。

次に、先ほどの例のように項が単純な命題ではなく、それぞれまたはどちらかが異なる項を particle「either」や「or」で結びつけて形成した排他的析取命題である場合を考えます。以下はその例です。ここで $X$、$Y$、$Z$ などは単純な命題を表します。

1. 「$X$ が真または $Y$ が真なら $Z$ が真」
2. 「$X$ が真なら $Y$ が真または $Z$ が真」
3. 「$X$ が真または $Y$ が真なら $Z$ と $W$ の両方が真または両方偽」

明らかに、これらの項のうちの一つまたは二つが排他的析取命題であるという事実は前提と結論との関係を変えるものではありません。したがって、既に確立された原理に従って前提と結論の適切な表現を得て、後者に不定記号 $v$ を付けて等式を形成するだけです。したがって、先述した命題に対する方程式は以下のようになります。

\begin{eqnarray*}
&1. \quad x(1 - y) + (1 - x)y = vz.\\
&2. \quad x = v \{y(1-z)+z(1-y)\}.\\
&3. \quad x (1 - y) + y (1 - x) = v \{ zw + (1 - z) (1 - w) \}
\end{eqnarray*}

この例で示されたルールは一般に適用されます。

異なる方法で排他的命題と条件命題の要素が複合命題の表現に参入する場合があることは想像できるものの、私は自然な人間の理性の要請によってそれらが現れることが稀であると考えており、したがってこれらについて議論することは避けます。この省略は大きな問題を引き起こすことはありません。なぜなら、上記の応用に基礎となった一般的な原理は完全に一般化されたものであり、少し考えればどのような想像できる場合にも適用できます。

13. 上述の表現法則には解釈に関する法則が暗黙に含まれています。方程式
\[
x = 1
\]
は命題 $X$ が真であることを表すと理解されるべきであり、方程式
\[
x = 0,
\]
は命題 $X$ が偽であることを表します。方程式
\[
xy = 1
\]
は命題 $X$ と $Y$ が両方とも真であることを表し、方程式
\[
xy = 0
\]
は命題 $X$ と $Y$ が両方とも真でないことを表します。

同様に、方程式

\begin{eqnarray*}
x(1-y) + y(1-x) &=& 1,\\
x (1 - y) + y (1 - x) &=& 0,
\end{eqnarray*}

はそれぞれ排他的命題「$X$ が真であるか $Y$ が真である」の真偽を主張します。方程式

\begin{eqnarray*}
y &=& vx,\\
y &=& v(1-x)
\end{eqnarray*}

はそれぞれ命題「もし命題 $Y$ が真であるなら、命題 $X$ も真である」、「もし命題 $Y$ が真であるなら、命題 $X$ は偽である」とを表します。

この作品の次の章でしばしば例として現れる場合があります。特定の方程式の一部の項が不定符号 $v$ の影響を受けている一方、他の項はそのような影響を受けない場合です。以下の例は説明のために適しています。もし

\[
y = xz + vx (1 - z)
\]

であると仮定するとします。

ここで、命題 \( Y \) が真である時間は、命題 \( X \) と \( Z \) が同時に真であるすべての時間に加えて、\( X \) が真で \( Z \) が偽である時間の無限大の部分を含むものとされています。これから、1つ目には、もし \( Y \) が真であれば、\( X \) と \( Z \) が同時に真であるか、または \( X \) が真で \( Z \) が偽であるいずれかであることがわかります。2つ目には、もし \( X \) と \( Z \) が同時に真であれば、\( Y \) は真です。これらの中の後者は逆解釈と呼ばれ、これは方程式の第二項から前件を引き出し、第一項から後件を取り出すことによって得られます。第二項に係数が1である項が存在することにより、この後者の解釈方法が可能になります。その原理は次のようになっています：

14. \textbf{原則}．\textit{方程式の特定の項において、係数が1である任意の構成項またはそれらを前提とし、他の項すべてが結論となる命題として取り出すことができます。}

例えば、

\[
y = xz + vx (1 - z) + (1 - x) (1 - z)
\]

という方程式に対して、以下の解釈があります：

\textbf{直接解釈}．\textit{もし命題 \( Y \) が真であれば、\( X \) と \( Z \) がともに真であるか、または \( X \) が真で \( Z \) が偽であるか、または \( X \) と \( Z \) がともに偽であるいずれかです}。

\textbf{逆解釈}．\textit{もし \( X \) と \( Z \) がともに真であるか、または \( X \) と \( Z \) がともに偽であるであれば、\( Y \) は真です}。

これらの部分的な解釈の合計は与えられた方程式の全体的な意味を表現します。

15. ここでは、再び注意を向けたいのは、二次命題の解釈理論において時間という概念が重要な要素であるように思われるが、表現法と解釈法の法律が明確に確立されると、実際にはその要素として無視できる可能性があるということです。これらの法則によって生じる形は、完璧な言語の形式に似ているように見えます。どのような既知または存在する言語でも、慣用句から解放し、余分な部分を取り除き、その言語で任意の命題を最も単純かつ直截的な方法で表現するとします——すべての言語が基礎とする純粋かつ普遍的な思考の原則に従って。それらはすべてその現れを示していますが、どの言語もそれから多少なりとも逸脱しています。このような言語から解析記号への移行は、一つの符号セットを別のものに置き換えること以外には、形式や性格において根本的な変更はありません。関係を表す要素、つまり事物または命題について、文字を使用します。析取接続については、+で書き表します。連接コピュラまたは関係記号については、=で書き表します。この類似性については深く追求する必要はありません。その実在と完全性は、後の章での理論の適用によって示される表現形式から、まだ不完全ながらも優れた思考の道具である英語言語との直接比較を通じてより明らかになります。

16. 一次命題論理と二次命題論理の一般的な類似性について、本章の主題を離れる前にいくつかの追加的な注釈を述べたいと思います。

私たちは、 Primary Propositions の理論を単なる空間の概念に基づいて確立することができたと断言してもよいでしょう。これは、Secondary Propositions の理論が時間という概念に基づいて確立されたのと同じ方法で行うことができたかもしれません。おそらく、これが行われていたら、我々が考慮している類似性は、空間と時間を「人間の理解の形式」であると考える人々の見解に、より密接な一致を示していたでしょう。空間と時間は、知覚によって与えられるものに対する知識の条件として、心の構造によって課せられたものだと考えられているからです。

しかし、この見解は一方では証明不能であり、他方では「場所」、〔ギリシャ語原典引用〕 を認識する必要性に私たちを縛るものです。実際、それがそうであるかどうかの問題は、私には能力を超えたものだと考えられます。しかし、Primary Propositions の論理過程が空間での現象の表出を必須条件とするわけではなく、そのような存在形態がある場合でも、同等の厳密さで適用可能であることを、私は確立されたと理解しています。

幾何学や動力学における多くの既知の例では、感覚によって与えられる空間概念とは異なる、ある知性的な空間概念に基づいて問題を分析することができることもまた事実です。これらの形式的な解析は、想像力によって実現されるものでもあります。

空間は、我々の感知において三つの次元、すなわち長さ、幅、深さを持つと提示されます。しかし、曲面の性質に関する多くの問題、固体体の軸周りの回転、弾性媒質の振動などに関連する大規模な問題群では、この制限は分析的な調査において任意の性格を示し、もし解法過程にのみ注意が向けられれば、空間が四次元やそれ以上の次元で存在する理由を見出すことができないことが明らかになります。このような仮想世界での知性的手続きは、類推の最も明確な光線によって理解することができます。

アリストテレスは、三次元空間の存在と古代の宗教的・哲学的な見解について如実に述べています：

「三次元空間の存在と、そのような観点から古代の宗教的・哲学的精神を描くためには、以下のように記述される：
〔ギリシャ語原典引用〕`--\textit{De Caelo}, 1.

したがって、私は空間の概念が基本的な命題の理論の発展に必須であるとは考えていません。しかし、極めて困難な問題について慎重に話すことを望むにもかかわらず、時間の概念が二次的な命題の理論を確立するのに必須であると考える傾向があります。

空間に対する人間の認識がどのような変化があったとしても、論理に関与する能力には影響を与えない可能性があると私は考えます。しかし、時間に関する認識は私たちがそれを感知する方法と同じでなければならず、そのような根拠は存在しないように思われます。

これらの推測を可能であるとは限らない前提として捨て去り、以下のように主張します：一次命題においてシンボル $1$ が宇宙全体を表す理由は、対応する等式の各項間の同一性符号 $=$ が示すように、それらが同一であることを意味し、単に空間上の同一部分に存在することを意味していないからです。同様に、二次命題においてシンボル $1$ が事件全体ではなく、その順序的な瞬間と期間で進化する永遠性を表す理由は、対応する等式の論理的各項間の同一性符号 $=$ が示すように、それらが同一であることを意味せず、その発生時間が同じであることを意味しているからです。

これらの理由は私にとって直ちに解釈に関する問題を決定づけるものと見なされます。以前の論文（Mathematical Analysis of Logic, p. 49）では、ウォリスの仮定命題の簡略化理論に従ってシンボル $1$ を「状況」または「状況の組合せ」の全体と解釈しましたが、この見解は「状況」という用語の意味を定義する必要性を伴います。そしてそれは明らかに、その用語を超えた時間以外の概念は論理の対象から外れ、形式的論理の過程を制約することを示しています。

\newpage

\chapter{二次命題の処理法（Methods in Secondary Propositions）}
\label{ch:xii}

\begin{center}
\large\textit{仮言的三段論法・選言的三段論法の代数方程式解法}
\end{center}
\vspace{4mm}

\#\# 二次命題に対する方法と過程

1. 前の研究（XI.7）から明らかになったように、論理符号の組合せ法則は、これらの符号が一次命題や二次命題を表す際にも同じである。両者の唯一の違いは解釈の違いにあるのみである。また、第五章第六節で確認されたように、同一形式法則（つまり符号の組合せと使用に関する法則）に関連する異なる思考体系と解釈体系が存在すると、これらの関連過程は問題の一次条件の表記からその象征的解決の解釈までの間において両者で同じである。したがって、一次命題と二次命題の二つの形式に現れる思考体系の間に、この形式法則の一貫性が存在し、私たちは以前の議論で示した一次命題クラスに対する手続きは、二次命題に対して修正せずに適用可能である。

2. したがって、二つの基本的な過程である消去と展開の法則は、二次命題の系においても一次命題の系においても同じです。また、一次命題において（第六章第2節で）どのように任意の分数形式を含まない方程式の解釈が、その方程式を成分の系列に発展させて、係数が消えないすべての成分に対して0と等しいとする方法によって実現されるかを見たことがありました。二次命題の方程式に対しても同じ手法が適用され、最終的に導く解釈された結果は、前記の場合（第六章第6節）と同じく共存する否定を含む系となります。しかし、前記の場合ではこれらの否定の力は特定の類の存在に対して発揮されるのに対し、後者の場合では、与えられた前提の語句に含まれる基本命題の組み合わせの真偽に関連します。そして、一次命題においても（第六章第7節など）のように、否定系をさまざまな他の命題形式に変換することが可能であることが確認されたように、ここでも同様の変換が可能です。唯一の違いは方程式の形ではなく、その解釈の性質にあるのみです。

3. さらに、一次命題において任意の要素が方程式系の残りの要素やそれらの一部を選び出した数だけの要素で表される方法（第六章第10節）を見つけ、その表現を論理的な推論に解釈することができたように、二次命題の系においても同じ目的を達成するために同じ手段が使用され、解釈の違いを除いては変わりありません。最終的に解決から排除したい要素の消去、方程式系の単一化、代数的解法とその解釈された形への展開は、一次命題の論議における対応するステップとは何ら異nessはありません。

ただし、二次命題の扱いにおいて現れる異なるケースを集約し、一般的なルールとして整理することにより、どのような困難も排除できるかもしれません。

---

**則**

与えられた命題（XI.11）を符号的に表す。

不定符号$v$（VII.5）が含まれるそれぞれの方程式から、個別にこの符号を除去する。

残りの符号は最終解から排除したいものであるので、それらを逐次的に除去する前に、符号またはその他の符号が含まれている方程式を一つの方程式に簡約化する（VIII.7）。結果を得た方程式を集めて、単一の方程式$V = 0$とする。

その後、最終的な関係を特定の形式で表す方法に応じて次の手順に従う：

1. **否定形または否定系の場合**：関数$V$を展開し、係数が消滅しないすべての成分を0と等しくする。
2. **析取命題の場合**：係数が消滅したすべての成分の和を1と等しくする。
3. **単純要素を持つ条件命題の場合（例えば$x$または$1-x$）**：その要素の代数的表現を求め、それを展開する。
4. **複合式を持つ条件命題の場合（例えば$xy$、$xy+(1-x)(1-y)$など）**：その式を新しい符号$t$と等しくし、後述する符号が出現する結果に$t$の展開された関数としてそれを決定する。これは通常の方法または特別な方法（IX.9）により行う。
5. **結果を解釈する**（XI.13, 14）。

特定の基本命題$x$が真か偽かを確認するだけであれば、すべての符号を除き$x$のみを残す。その場合、方程式$x = 1$は命題が真であることを示し、$x = 0$は偽であることを示し、$0 = 0$は前提が真または偽を決定するのに十分でないことを示す。

---

4. 例1.—シケルロの不完全な論文『運命について』において、次のような予測が議論の対象となっています：「もしファビウスが犬星の昇る方に生まれたなら、彼は海で死なない」。この予測については、本章の方法を適用します。$y$ を「ファビウスは犬星の昇る時に生まれた」という命題を表すものとし、$x$ を「ファビウスは海で死ぬ」という命題を表すものとします。$x$ が「ファビウス、…」という命題を表しているとは、$x$ がその命題に特定され（XI.7）、方程式 $x = 1$ はその命題の真偽を宣言し、方程式 $x = 0$ はそれを否定することだけを意味します。議論すべき方程式は以下のようになります：
\[
y = v(1-x).
\]
まず、与えられた命題を否定または否定の系に帰着する（XII.3）ことを要求しましょう。移項すると次のようになります：
\[
y - v(1-x) = 0.
\]
$v$ を消去すると以下の式が得られます：
\begin{eqnarray}
y \left\{ y - (1-x) \right\} &=& 0, \nonumber \\
\textrm{または } y - y(1-x) &=& 0, \nonumber \\
\textrm{または } yx &=& 0.
\end{eqnarray}
この結果の解釈は次のようになります：「ファビウスが犬星の昇る時に生まれたという命題と、彼が海で死ぬという命題は両方とも偽である」。シケルロはこのような形の命題を「相反するものによる結びつき」と呼んでいます；そして彼は、将来に関する偶然な主張に対する困難を回避するためにチリッパスがこの方法を使用したことを述べています：「ここではチリッパスは、占星学者や神託者を欺くために、このような結びつきを使ってその受け取り方を否定するだろうと期待している。もし誰かが犬星の昇る方に生まれたなら、彼は海で死なない；しかし、それよりもむしろ次のように言うべきである：「ファビウスが犬星の昇る方に生まれたことと、彼が海で死ぬことは両方とも偽である。… これは冗談の許容性を示すものだ！ … 言法には多くの種類があり、そのどれもこれほど歪曲されたものはない。チリッパスは占星学者たちをこのようにしてストイク派のために満足させると期待しているのだ」とシケルロは述べています。（Cic. De Fato, 7, 8）

5. 与えられた命題を排他的形式に簡約する。
(2) 式の最初の項に入らない因子は次のようになる：
\[
x(1-y), (1-x)y, (1-x)(1-y).
\]
したがって、以下の式を得る：
\begin{equation}
y(1-x) + x(1-y) + (1-x)(1-y) = 1.
\end{equation}
この式の解釈は次のようになる：「ファビウスは犬星の昇りとともに生まれ、海で死ぬことなく生き残る；または、犬星の昇りとともに生まれなかったが、海で死ぬ；または、犬星の昇りとともに生まれなかったが、海で死ぬことなく生き残る。」

しかしながら、このような場合においては、異なる因子によって区別される因子が存在する場合、それらを一つの項にまとめること（VII. 15）は最も便利であることが以前にも述べた通りである。したがって、(3) 式の最初と第三の項を結びつけると、以下のようになる：
\[
(1 - y)x + 1 - x = 1;
\]
また、第二と第三の項も同様に結びつけば、以下のようになる：
\[
y(1-x)+1-y=1.
\]
これらの式はそれぞれ次の解釈を与える：

「ファビウスは日星の昇りとともに生まれなかったが、海で死ぬ；または、彼は海で死なない。」

「ファビウスは日星の昇りとともに生まれたが、海で死ぬことなく生き残る；または、彼は犬星の昇りとともに生まれなかった。」

これらの解釈は以前のものと厳密に等価であることが明らかである。

次に、「ファビウスは海で死ぬ」という仮定から導かれる結果を条件命題として確認する。

元の方程式から$v$を消去した結果を表す(2) 式において、$x$を$y$の関数として決定することが必要である。

以下のように展開すると：
\[
x = \frac{0}{y} = 0y + \frac{0}{0}(1-y),
\]
または
\[
x = \frac{0}{0}(1-y).
\]
この式の解釈は次のようになる：「ファビウスが海で死ぬなら、彼は犬星の昇りとともに生まれていなかった。」

これらの例は、前節で述べたように、基本的な命題と派生的な命題との間の関連性を部分的に示している。その関連性の特徴は、過程の同一性と解釈の類似性である。

6. 例2.--プラトンの『共和論』第二書において、神の性質の不変性を証明するための興味深い議論があります。これは、プラトンが一般的な原理に到達するために用いる具体的事例から慎重に導き出す精密な帰納法と、それらから彼の主張すべき特定の推論を導く明瞭かつ連続的な論理の非常に優れた例です。この議論は以下の対話の中に含まれています：

「あるものが本来の形から逸脱するとき、それは自己または他の何かによって変化する必要がありますか？もちろんです。最高の状態にあるものは最も変化しないで済むはずです。体は食べ物や飲み物、労働などによって、植物は熱や風などの影響によって最も変化しないでしょう。健康で強く、最も変化しないものとは何でしょうか？もちろん、最も強いと最も賢明な魂は外部からのどんな困難も最小限にしか揺るがず変化させないでしょう。また、すべての作られた器物、家具、衣服についても同じ原理が適用されます。これらが良好な状態にあるものは時間や他の事故によって最も変化しないでしょう。もちろんです。そして、自然または技術または両方により適切に整備されたものは、他の何かから最小限の変化を受けるべきです。そう思われますね。しかし、神と神聖なものはすべてにおいて最高の状態にあるのです。もちろん。このようにすれば、神は最も多くの形を被らないはずですか？もちろんです。また、彼自身を変えるのでしょうか？彼が変化するなら明らかにその通りです。そして、彼はより良いものや美しいものへと変化するのでしょうか、それともより悪いものや卑しいものへと変化するのでしょうか？彼が変化すれば必ずしも悪い方に変化します。なぜなら、神には美しさや徳の不足などと言えることはないからです。私はあなたが最も正しく述べたと思いますね、こう言いました。そして、そのようにであれば、アディマンタスよ、神または人間が何かにおいて自己をより劣ったものにすることを好むとは思われますか？不可能だと彼は言いました。そうですね、と私は言いました。それでは、神が自らを変えることを望むことはあり得ませんね；常に最も美しいと最善の形で存在し続けるべきです。」

上記議論の前提は以下の通りです：

1. 神が変化する場合、神は自己自身によって変化するか、他の何かによって変化する。

2. 神が最良の状態にある場合、神は他の何かによって変化しない。

3. 神は最良の状態にある。

4. 神が自己自身によって変化する場合、神はより悪い状態になる。

5. 神が意志を持って行動する場合、神はより悪い状態にはならない。

6. 神は意志を持って行動する。

これらの前提を要素として表すと次のようになります：

$x$ を「神が変化する」という命題を表すとする。

$y$ は「神は自己自身によって変化する」ことを表す。

$z$ は「神は他の何かによって変化する」ことを表す。

$s$ は「神は最良の状態にある」ことを表す。

$t$ は「神はより悪い状態になる」ことを表す。

$w$ は「神が意志を持って行動する」ことを表す。

すると、これらの前提をシンボル言語で表現したものは以下の等式となる：

\begin{align}
   xyz + x(1-y)(1-z) &= 0,\\
   sz &= 0,             \\
   s &= 1,              \\
   y(1-t) &= 0,         \\
   wt &= 0,             \\
   w &= 1.
\end{align}

$x$ を保持し、$z$, $s$, $y$, $t$, そして $w$（これらのシンボルが上記のシステムで出現する順序に従って）を次々と消去すると、それぞれの結果を解釈する。

(1) と (2) の $z$ を消去すると、

\begin{equation}
     xs(1-y) = 0.
\end{equation}

(3) と (7) の $s$ を消去すると、

\begin{equation}
     x(1-y) = 0.
\end{equation}

(4) と (8) の $y$ を消去すると、

\begin{equation}
     x(1-t) = 0.
\end{equation}

(5) と (9) の $t$ を消去すると、

\begin{equation}
     xw = 0.
\end{equation}

(6) と (10) の $w$ を消去すると、

\begin{equation}
     x = 0.
\end{equation}

これらの等式は、(8) から始めて以下の結果を導く：

(8) より $x = \frac{0}{0}y$ であるから、「神が変化する場合、神は自己自身によって変化する」。

(9) より $x = \frac{0}{0}t$ であるから、「神が変化する場合、神はより悪い状態になる」。

(10) より $x = \frac{0}{0}(1-w)$ であるから、「神が変化する場合、神は意志を持って行動しない」。

(11) より「神は変化しない」。これはプラトンの結果である。

現在、私は以前に述べた通り、消去の順序は関係ありません。この場合、$w$ から始めて逆順で同じシンボルを消去することによってこの事実を確認しましょう。結果得られる方程式は以下のようになります。

\[
t=0, y = 0, x(1-x) = 0, z = 0, x = 0;
\]

これらの方程式から以下のような解釈を得ます：

神は劣化しない。
自己によって変化しない。
変化するなら、他のものによって変化する。
他のものによって変化しない。
変化しない。

異なる道で同じ結論に至ります。

これらの例は方法の力の展示としては微々たるものですが、より複雑な例と同様に、その性質や特徴を説明するために十分です。

7. 次のように、一次命題系と二次命題系との間の類似性の一例として述べることができます。後者の系においても基本方程式
\[
x(1 - x) = 0,
\]
が解釈可能であり、これは命題が同時に真で偽であることはないとする公理を表します。これと III. 15 の対応する解釈を比較すると、この方程式は展開して
\[
x=\frac{0}{1-x}=\frac{0}{0}x,
\]
という形で解くと、それぞれの公理を得ることができます：「何かはそのものである」、「真なる命題は真である」という形式。これらは「同一性の原理」と呼ばれるものと同型です。これらの公理について最も対立する見解が存在します。一部の人々はこれらを哲学の本質的な部分と見なし、ロックはこれらについて一章を費やし、「軽微な命題に関する」章としてタイトルをつけました。
\footnote{人間理解論, 第四書, 第八章}

両者の視点には真実と誤りが混在しているように思われます。経験を排すか、または学校の虚飾的な議論の材料とするかという観点から見れば、これらの命題は軽微以上に無意味です。他方で、思考の法則や条件と密接に関連していると見なされるなら、これらは推測的な重要性を獲得します。

\newpage

\chapter{サミュエル・クラークおよびスピノザの哲学的論証の分析（Clarke and Spinoza）}
\label{ch:xiii}

\begin{center}
\large\textit{存在論的証明・神の必然的実在性・スピノザ『エチカ』の記号論理的厳密批判}
\end{center}
\vspace{4mm}

\chapter[クラークとスピンノザ]
{\large 以下章における調査の一般的な順序は次の通りです。まず、上記論文のいくつかの一般命題の証明において実際に含む前提条件を検討します。これらの前提条件が明示的に述べられているか否かに関わらず、証明過程で仮定され、最終的な結論の基礎となる部分として使用される任意の命題を指します。これらの決定された前提条件はシンボル語で表現し、それらから、この作品の前章で開発された方法を使用して、実際に著者が導き出した特定の推論に加えて、さらに重要な推論も導き出します。一部の場合には、前提条件からある事実や原理を省略したり、新たな命題を追加または代入することで前提条件を修正し、そのような変更によって最終的な結論がどのように影響されるかを決定します。これらの目的の追求においては、これら有名な作品で定められた形而上学的原理がどれほど信頼に値するかについて尋ねる必要はありません。しかし、与えられた前提条件から正当に導けるどのような結論が抽出できるかを確認することだけが必要です。そしてこれを行うことで、最終的なまたは結論的な命題の要素に関する選択と順序に関して我々が持つ完全な自由度を例示します。すなわち、基本的な命題が真偽であるかを決定し、与えられた制約や組み合わせ下で真偽であるかを決定することです。}

2. この調査における主要な実用的な困難は、方法が既に決定された前提に対して適用される際に生じるのではなく、どのような前提であるかを確定する上で存在します。特に形而上学的な問題に対する最も厳格な例として考えられる領域においては、異なる思考の流れが混在していることが時々見られ、証明の特定の但し重要な部分が括弧内に記載されたり、主な議論の外側で述べられていることがあります。また、前提の意味が多少不明瞭である場合や、厳格な形式的論理の法則から見て不正確または結論的でない議論があることもしばしば見られます。証明の真の前提を決定し、それらを明確に示す困難は、特定の結論が証明されているかどうかを確認しようとする人々にとって非常に重大なものとなるかもしれません。しかし、この方法によって特定の結論が証明されているかどうかを確定するためには、誰もがそのような困難を克服しなければなりません。そのため、この方法に関連してよりはっきりと現れますが、他のどの方法よりも特有のものではありません。実際、この関連性は非常に密接であるため、この論文の方法を使用してさえ、前提に基づいて結論の妥当性についてすら推測を立てることが不可能となるのです。つまり、証明のすべての前提に対する明確な理解と正確な記述がなければ、結論の妥当性については何らの推測も立てられません。

通常の手続きにおいては、議論の一部のステップを検討し、それから全体の範囲について漠然とした一般的な印象を形成することが非常に一般的であり、前提の前段階的な完全な分析を行うことなくそのような予備的な分析が含まれていることもあります。

厳格な前提の決定が求められる必要性は、必ずしも悪 thing とは考えられず、特にその任務が達成されると、疑念や曖昧さのすべての源泉が取り除かれるという点で、その価値がある。

3. 「神の存在と属性の証明」は一連の命題または定理からなり、それぞれが主として二つの異なるクラスに分類される前提によって証明されている。これらの前提は観察事実、例えば物質世界の存在、運動の現象など、そして先験的に認識される権威と普遍性を持つ仮説的原理である。当然、その証明の有効性は後者の真実性、論理が正しければという前提を除いて、推論の正確性に依存する。しかし、その主張に対する評価はいかなるものであれ、少なくとも知的な業績としては非常に高い merits があることは間違いない。

彼等が含む議論の列は一般的には明確に整理されておらず、それでもほとんど常に正規な論理の例示であり、理解の鋭さと推理力の強さを示し、それらは稀有なもので、おそらく超えることがない。その中で、形而上学的探求の領域で目覚めた知的な努力の集大成を見ることができる。

これらの命題が存在した時代には、仮説的原理の支配が現在ほど疑問視されなかった時期であり、数学物理学の新興学派による厳格な証明がその方向性を示唆するモデルとなった時期でもあった。このため、形而上学的な探求の領域で扱われる主題の尊さを考慮に入れるなら、これらの命題は高い評価に値すると私は考えている。そして、その一部について精密な分析を行うことが無駄であるとは思わない。

4. ベネディクト・スパノザの『倫理学』は、神と宇宙の同一性を証明し、その教義に基づいて倫理学および哲学の体系を確立することを目的とする論説である。主な議論の分析は非常に困難であり、これは個々の命題の複雑さではなく、曖昧な定義と不明瞭な公理を使用することで生じる混乱が原因である。スパノザの論理は部分的に文脈的であり、クラーク博士の論理よりも著しくその程度が高い；おそらくこれが、一部の人々にとって彼の論理が形式的な妥当性を有するように見えた理由かもしれない。しかし実際には、そのような正当な主張ではない。これらの点は適切な場所で考慮される。

\begin{center}
\textsc{クラークの証明}

\textsc{命題I}
\end{center}

5. 「何かが永遠に存在していた。」

この証明は以下の通りである：

「何かが現在存在していることは明らかであり、したがって何かは常に存在していた。もし何かが存在していないと仮定すると、現在存在するものは何もなければ生じたことになる。つまり、全く原因なく何かが生じることになる。これは言葉の矛盾である。何かが生成されたにもかかわらず、その生成の原因が全くないという状況は、何かが何もないときに生成されることを意味し、それは同時に生成されないときと同じ時刻に生成されることを意味する。存在するものは、自己の性質による必然性から存在し続けた場合、または他の存在体の意志によって存在した場合、いずれの場合もその存在は自体で永遠である。」

3番目の主張は、現在存在するものが何も無かった状態から生じた場合、何かが実際に作用し、同時にその作用は全く何にも依存せずに達成されたことになる。

4番目の主張では、何かが全く何にも依存せずに達成された場合、それは全く達成されなかったことになる。

この議論の第二部分は、証明すべき点を仮定する単なる推測か、あるいは異なる言葉でその点を明確にする試みに過ぎないように思われる。

第三項と最後の項では、元々証明すべき主張、「何か常に存在していた」という主張、または第二部分での証明された点、「現在存在するものが何も無かった状態から生じた」という仮定が成り立たないという点の真実性についての明確な証明が含まれている。これは次のように分解できる：

5番目の主張は、何かが存在する場合、それは自己の必然性によって存在するか、あるいは他の存在の意志によって存在することである。

6番目の主張では、何かが自己の必然性によって存在する場合、何か常に存在していたことになる。

7番目の主張では、何かが他の存在の意志によって存在する場合、「現在存在するものが何も無かった状態から生じた」という命題は偽である。

この最終的な主張は、クラーク博士のテキストでは明示的に述べられていないが、彼の前存在する別の存在という結論は明らかに「現在存在するものが何も無かった状態から生じた」という命題を否定することと等価であることが意図されている。

したがって、この証明は二つの独立した議論の流れで構成されていることが明らかになる。一つの流れは私が指定した証明の第一部分と第二部分を含むものであり、もう一つの流れは第一部分と第三部分を含むものである。後者の流れについて考えることにしましょう。

前提は以下の通りである：

1番目の前提：何かが存在する。

2番目の前提：何かが存在する場合、何か常に存在していたか、あるいは現在存在するものが何も無かった状態から生じたかのいずれかである。

3番目の前提：何かが存在する場合、それは自己の必然性によって存在するか、あるいは他の存在の意志によって存在することである。

4番目の前提：何かが自己の必然性によって存在する場合、何か常に存在していたことになる。

5番目の前提：何かが他の存在の意志によって存在する場合、「現在存在するものが何も無かった状態から生じた」という仮定は偽である。

現在、上記の命題を象徴的に表す必要があります。

次のように定義します：
\[
\begin{array}{cccccc}
& x = & \text{何かは存在する。}\\
& y = & \text{何かは常に存在した。}\\
& z = & \text{現在存在するものは何もから生まれた。}\\
& p = & \text{それが自己の性質による必然性で存在する} \\
& & (\text{i.e. 上記で述べた何か)。}\\
& q = & \text{それが他の存在の意志によって存在する}。
\end{array}
\]

「$x=$ 何かは存在する」という表現は、$x$ がその命題（XI.7）を表す代表的な符号であることを意味し、方程式 $x=1$, $x=0$ 分別でその真偽を宣言します。

前提の方程式は以下の通りです：
\[
\begin{array}{cccccc}
1つ目： & x = 1;\\
2つ目： & x=v[y(1-x)+z(1-y)];\\
3つ目： & x=v[p(1-q)+q(1-p)];\\
4つ目： & p=vy;\\
5つ目： & q=v(1-z);
\end{array}
\]

不定符号 $v$ を消去すると、以下の方程式を得ます：
\begin{eqnarray}
1-x=0;\\
x[yz+(1-y)(1-z)]=0;\\
x[pq+(1-p)(1-q)]=0;\\
p(q-y)=0;\\
qz=0.
\end{eqnarray}

6つ目、まず、符号 $y$, $z$, $p$, $q$ それぞれを表す単独の命題について、上記から真偽に関するどのような結論が導かれるかを検討します。すなわち、「何かは常に存在した」、「現在存在するものは何もから生まれた」、「何かは自己の性質による必然性で存在する」、「何かは他の存在の意志によって存在する」という命題についてです。

この目的のために、$y$ 以外のすべての符号を消去し、$z$ 以外のすべての符号を消去し、そのような単独の関係が存在するかどうかを決定する結果を得ます。

$x$ を (1), (2), (3) から消去するために必要なのは、(2) と (3) の $x$ の値を (1) から導き出し、それを代入することです。これにより以下の結果が得られます：
\begin{eqnarray}
yz+(1-y)(1-z)=0;\\
pq+(1-p)(1-q)=0.
\end{eqnarray}

$p$ を消去するために (4) と (7) の和を取ると、以下の方程式を得ます：
\[
p(1-y)+pq+(1-p)(1-q)=0;
\]
これより
\[
(1-y)(1-q)=0.
\]

$q$ を (5) と (9) から消去すると、以下の方程式が得られます：
\[
qz + (1-y)(1-q) = 0;
\]
これより
\[
x(1-y) = 0.
\]

残る方程式（6）と（10）について、これらを加えると、

\[
yz + 1 - y = 0
\]

となります。この式から $z$ を消去すると、

\begin{equation}
1 - y = 0, \quad \text{または} \quad y = 1
\end{equation}

を得ます。同じ式から $y$ を消去すると、

\begin{equation}
z = 0
\end{equation}

となります。

（11）の解釈は

「何か常に存在した。」

であり、（12）の解釈は

「何かが存在するときには、それは何もなければなり得ない。」

です。

次に、方程式（6）、（7）と方程式（4）、（5）を用いて $p$ および $q$ を含む二つの方程式を求めます。

まず $y$ を消去するために、（4）と（6）から、

\[
p(1 - y) + yz + (1 - y)(1 - z) = 0
\]

を得ます。これより、

\begin{equation}
(p + 1 - z)z = 0, \quad \text{または} \quad pz = 0
\end{equation}

となります。

次に、（5）と（13）から $z$ を消去するために、

\[
qz + pz = 0
\]

を得ます。これより、

\[
0 = 0
\]

となります。

残る方程式は（7）のみで、この方程式から $q$ を消去すると、

\begin{equation}
0 = 0
\end{equation}

が得られます。

したがって、前提として「何かが存在するときには、それは自己の性質によって必然的に存在することである」という命題の単純な真偽を肯定または否定することはできないことが分かります。また、$p$ を消去すると同様に、

\begin{equation}
0 = 0
\end{equation}

が得られるため、「何かが存在するときには、それは他の存在の意志によって存在することである」という命題の単純な真偽についても前提から結論を導くことはできないことが分かります。

一つ以上の命題を表す基本的な符号で表現される関係について、考慮すべきは方程式（7）のみです。残りの符号で表現される命題は、ここで述べられているような関連性に関わらず絶対的に真または偽であるからです。方程式（7）の解釈を、

\[
p(1 - q) + q(1 - p) = 1
\]

と表すと、

「何かが存在するときには、それは自己の性質によって必然的に存在することか、他の存在の意志によって存在することである。」

となります。

以上の分析の詳細を示したのは、上記のように過剰に冗長であるかもしれないが、その後の例では結果を得る手順を示すことに重点を置きたいと考えており、そのような繰り返しの頻度が退屈になることを避けるためだ。前述の方法による適用から得られた結論は、通常の論理によって容易に検証できる。

読者は、ドクター・クラークの最初の命題に関連する他の前提を用いて、その方法を適用し、それらから上記分析が導いた二つの結論のうちの最初の二つを導き出すことが難しくないだろう。

**命題II**

7. ある不変で独立した存在は永遠に存在している。

上記の命題を証明するための前提は以下の通りである：

1. 何かしら常に存在している。
2. 何かしら常に存在している場合、その存在は一つの不変で独立した存在か、または全ての存在事物が可変で依存的な系列の中に含まれているかのいずれかである。
3. すべての存在事物が可変で依存的な系列の中に含まれている場合、その系列は外部から原因を受け取るか、あるいは内部から原因を受け取るかのいずれかである。
4. 外部からの原因は受け取っていない（なぜなら、仮定により全ての存在事物を含んでいる）。
5. 内部からの原因も受け取っていない（なぜなら、どの部分も必要でないし、どの部分も必要でない場合、全体も必要でないからである）。

第四と第五の前提の括弧内の補助証明は省略するが、以下の前提を示すことができる：

\begin{eqnarray*}
x &=& \text{何かしら常に存在している。} \\
y &=& \text{一つの不変で独立した存在が存在している。} \\
z &=& \text{可変で依存的な系列が存在している。} \\
p &=& \text{その系列は外部から原因を受け取っている。} \\
q &=& \text{その系列は内部から原因を受け取っている。}
\end{eqnarray*}

次に、以下の連立方程式があります：
\begin{eqnarray*}
\textrm{1番目の方程式 } x &=& 1 ; \\
\textrm{2番目の方程式 } x &=& v \{y(1 - z) + z(1 - y)\} ; \\
\textrm{3番目の方程式 } z &=& v \{p(1 - q) + (1 - p)q\} ; \\
\textrm{4番目の方程式 } p &=& 0; \\
\textrm{5番目の方程式 } q &=& 0 :
\end{eqnarray*}
これらの式から不定符号 $v$ を逐次消去すると、以下の連立方程式が得られます：
\setcounter{equation}{0}
\begin{eqnarray}
l - x = 0; \\ % (1)
x \{yz + (1 - y)(1 - z)\} = 0; \\ %(2)
z \{pq + (1 - p)(1 - q)\} = 0; \\ %(3)
p = 0; \\ %(4)
q = 0.  % (5)
\end{eqnarray}

上記の連立方程式から $x$, $p$, $q$, および $y$ を逐次消去すると、方程式
\[z = 0.\]
が得られます。また、$x$, $p$, $q$, および $z$ を逐次消去すると、同様の方法で方程式
\[y = 1.\]
が得られます。

これらの方程式の解釈は以下の通りです：

1番目の解釈 \textit{存在するすべてのものは変化と依存性によって構成される連続した系列の中に完全には収まらない。}

2番目の解釈 \textit{ある一つの不変で独立な存在が存在する。}

これらの中からDr.クラークが証明する命題は後者です。

上記の分析により、文字符号 $x$, $y$, $z$, $p$, $q$ によって表されるすべての命題が絶対的に真または偽であると決定されたため、これらの命題をさらに結びつける関係が存在するかどうかについて調査する必要はありません。

別の証明がProp. II.について与えられますが、その簡略化のために省きます。クラークの議論の一環として部分的に構成される「無限順列の不可能性」については、一般的に形而上学の基本原理とみなされ、因果関係以外の問題にも拡張されてきました。アリストテレスはこの原理を用いて、証明の第一原理の必要性を立証しました。\footnote{『形而上学』III. 4; 『前件論理』I, 19, \textit{et seq}.}また、人間の行為における目的（善）の必要性なども含めています。\footnote{『ニコマコス倫理学』第一書 第二章}彼の著作の中で最も頻繁に引用される原理の一つかもしれません。中世の哲学者たちは、無限の種類と属概念の階層構造の不可能性を立証するために同様に適用し、これにより普遍的存在の必然的な存在が導かれました。おそらく無限の系列について明確かつ完全な概念を形成することができないという事実、すなわち全体として理解できないことから、論理的矛盾や概念自体における論理的不整合性と混同されることがありました。

8. 以下の議論の分析は基本的な命題の理論に依存しています。

\textsc{命題III}：

「変化しないかつ独立した存在は自己生成であるべきである」という主張があります。

前提は次の通りです：

1. すべての存在は、何も無から生じたか、外部原因によって生じたか、または自己生成である必要があります。

2. 何も無から生じた存在はありません。

3. 変化しないかつ独立した存在は外部原因によって生じたものではありません。

以上の前提をシンボルで表すと、

\begin{eqnarray*}
x &=& \text{何も無から生じた存在。} \\
y &=& \text{外部原因によって生じた存在。} \\
z &=& \text{自己生成である存在。} \\
w &=& \text{変化しないかつ独立した存在。}
\end{eqnarray*}

すると、

\setcounter{equation}{0}
\begin{eqnarray}
x(1-y)(1-z) + y(1-x)(1-z) + z(1-x)(1-y) = 1, \\ % (1)
x = 0, \\ % (2)
w = v(1-y),  % (3)
\end{eqnarray}

最後の式から$v$を消去すると、

\begin{equation}
wy = 0. % (4)
\end{equation}

ある符号の値が0または1として与えられた場合、単純な代入によってその符号を省略するのが最も適切です。したがって、$x$ を省略すると、

\begin{eqnarray}
y(1-z) + z(1-y) &=& 1; \\ % (5)
\textrm{または, } yz + (1-y)(1-z) &=& 0. % (6)
\end{eqnarray}

次に、(4)と(6)を加え、$y$ を省略すると、

\begin{eqnarray*}
w(1-z) = 0, \\
\therefore w = vz;
\end{eqnarray*}

この式の解釈は、「不変で独立な存在は必然的に自己存続するものである。」

(5) の実際的な形での解釈は、「すべての存在は外部の原因によって生じたか、または自己存続するものである。」

9. 上記命題についてドクター・サミュエル・クラーク博士の注では、驚くべき議論が展開されており、それが物質世界が上記で述べた自己存続する存在ではないことを証明することを意図しています。私は引用する文脈は以下の通りです：

「もし物質が必然的に存在すると仮定すると、その必然的な存在の中に重力の能力が含まれているか否かということが考えられます。重力を含んでいない場合、知性を持つ存在が支配していない単なる物質的世界においては、運動は存在しないはずです。なぜなら、運動は既に示され、問題文でも認められているように、自己そのものには必要ないからです。しかし、重力の能力が物質の仮定された必然的な存在の中に含まれている場合、重力を含むことから真空が必要不可欠である（これは不朽のSir Isaac Newtonが十分に証明しています）、またこれにより物質は必然的な存在ではないことが導かれるのです。なぜなら、真空が実際に存在する場合、物質が存在しないことは明らかに不可能ではないからです。」——(pp. 25, 26)

上記の論理を慎重に検討すると、その実際的な前提は以下の通りであることがわかります：

1. 物質が必然的な存在である場合、重力の性質が必然的に存在するか否か。

2. 重力の性質が必然的に存在しない場合で、世界が支配する知性の下に置かれていなければ運動は存在しない。

3. 重力の性質が必然的に存在する場合、真空の存在が必然的である。

4番目、真空の存在が必要であるならば、物質は不要な存在ではない。

5番目、物質が不要な存在であるならば、世界には監督的な知性が存在しない。

6番目、運動が存在する。

以上の前提の中で、最初の四つは証明において表現され、第五のものは最初と第二の文の関連から暗に示されています。第六のものは事実を述べており、著者がそれを必要とするとは思っていなかったかもしれませんが、明らかに彼の論理の基礎の一環です。これらの基本的な命題を以下のように表すことにしましょう：

\begin{eqnarray*}
x &=& \text{物質は不要な存在である。} \\
y &=& \text{重力は必然的に存在する。} \\
t &=& \text{重力は必然的に存在しない。} \\
z &=& \text{世界は単なる物質であり、監督的な知性の下にない。} \\
w &=& \text{運動が存在する。} \\
v &=& \text{真空は必要である。}
\end{eqnarray*}

すると、前提のシステムは以下の式で表現されます。ここで $q$ は不定な時間のシンボルとして使用されています：

\begin{eqnarray*}
x &=& q\{y(1-t) + (1-y)t\}. \\
tz &=& q(1-w). \\
y &=& qv. \\
v &=& q(1-x). \\
x &=& qz. \\
w &=& 1.
\end{eqnarray*}

これらの式から $q$ を消去すると、以下のシステムが得られます：

\begin{eqnarray}
x\{yt+(1-y)(1-t)\} &=& 0.\\
tzw &=& 0. \\
y(1-v) &=& 0. \\
vx &=& 0. \\
x(1-z) &=& 0. \\
1-w &=& 0.
\end{eqnarray}

これらの式から $w$, $v$, $z$, $y$, および $t$ を消去すると、以下の等式が得られます：
\[ x = 0, \]
これは「物質は不要な存在ではない」という命題を表しています。これがドクター・クラークの結論です。他の任意の五つのシンボル（$v$, $z$, $y$, $t$, および $x$ の組み合わせを除く）を消去しようとすると、結果は $0=0$ の形式になります。
したがって、以下のことが明らかになります：
\textit{単一のシンボルで指定される基本的な命題の絶対的な真偽を表す他の結論は存在しない。}

結論を含む方程式が2つのシンボルを含む場合、存在するものは以下のもののみである：「世界が単に物質であり、監督的な知性の下でないなら、重力は必ずしも欠けていない。」この結論は次のような方程式によって表現される：

\[tz = 0, \textrm{ すなわち } z = q (1 - t).\]

上記の分析において、運動の存在を示す最終前提を省略し、残りの前提に含まれる仮定原理を変更せずに考察する場合、いくつかの興味深い結論が導き出される。これらについては次の章で注意深く説明する。

10. ドクター・クラーク氏の議論の残り部分について、その主な内容と関連性を簡潔に述べる。特に他の部分よりも複雑な構造を持つ部分に焦点を当てる。

命題 \textsc{iv} においては、「自存する存在の本質や実体は理解不能である」と示される。用いられた論理の流れは、他のすべてのものについてその根本的な性質を知る能力が欠けていることを我々は知らないから、さらに自存する存在の本質も理解できないという点に焦点を当てる。

命題 \textsc{v} においては、「自存する存在の本質や実体自身は我々にとって絶対的に理解不能であるとしても、その性質のいくつかは厳密な証明によって示されるべきであり、また彼の存在も同様である」と主張される。

命題 \textsc{vi} においては、「自存する存在は必然的に無限で omnipresent である」と主張され、さらに「その無限性は満たし尽くされた無限と広大な無限の両方である」と主張される。証明が進行する根拠は、絶対的な存在は時間、場所、状況から独立しており、制約を受けることなく、不完全性を排除することである。これより推論され、「自存する存在は最も単純で変化しない、腐敗しない存在であり、部分を持たず、形も動きもなく、物質に見られるような他の特性を持っていない」と主張される。

実際に用いられた前提は以下のようになる：

1. 有限の存在が自己存続的であるなら、その存在しないことを仮定することは矛盾である。

2. 有限の存在は、一つの場所から欠如することを仮定しても矛盾ではない。

3. 一つの場所から欠如することが矛盾でないものについては、全ての場所から欠如することも矛盾ではない。

4. 全ての場所から欠如することが矛盾でないものについては、その存在しないことを仮定することも矛盾ではない。

以下は、以上の文を表すために導入した変数を用いた式です：

\begin{eqnarray*}
x &=& \text{有限の存在。}\\
y &=& \text{自己存続的なもの。}\\
z &=& \text{その存在しないことを仮定することが矛盾であるもの。}\\
w &=& \text{一つの場所から欠如することが矛盾でないもの。}\\
t &=& \text{全ての場所から欠如することが矛盾でないもの。}
\end{eqnarray*}

これらを表すと、不定記号を省略した式は以下のようになります：

\begin{eqnarray}\setcounter{equation}{1}
&xy(1-z) &= 0.\\
&x(1-w)  &= 0.\\
&w(1-t)  &= 0.\\
&tz = 0.
\end{eqnarray}

$t$, $w$, そして $z$ を順に消去すると、以下の式を得ます：

\[
xy = 0,
\]

\[
\therefore y = \frac{0}{0}(1-x).
\]

この解釈は、「何らかの存在が自己存続的であるならそれは無限である」という結論を導きます。

命題 VII では、自己存続的な存在は必然的に一つであると主張します。証明の順序は以下の通りです：まず、自己存続的な存在は「必然的に存在する」ことである、そして「絶対的な必要性は単純で一貫しており、どのような可能性のある差異や変化も含まない」ということ、さらに、「存在の差異や変化」は依存を意味することから、つまり、「何らかの存在が必然的に存在するならそれは自己存続的な存在の一贯した本質である」という結論に至ります。

この結論は以下の前提からも導かれます：

1. 二つ以上の独立な必然的存在がある場合、そのいずれかを単独で存在させることが可能である。

2. そのいずれかが単独で存在させられるなら、他方の存在しないことを仮定することは矛盾ではない。

3. そのような仮定が矛盾ではないなら、二つの独立な必然的存在は存在しない。

次のように基本命題を表すことにしましょう：
\begin{eqnarray*}
x &=& \textrm{二つの必要で独立な存在が存在すること。}\\
y &=& \textrm{いずれか一方だけが存在することが矛盾しないこと。}\\
z &=& \textrm{他方が存在しないことが矛盾しないこと。}
\end{eqnarray*}

それでは、以前と同様に進めると、次のような式を得ます：
\begin{eqnarray}\setcounter{equation}{1}
x(1-y) &=& 0.\\
y(1-z) &=& 0.\\
zx = 0.
\end{eqnarray}

$y$ と $z$ を消去すると、
\[
x=0.
\]
従って、\textit{二つの必要で独立な存在は存在しない}。

11. 二つの前段落の前提に依拠する点について、公刊当時非正統派学院で学徒を務めていたブリストル主教が、「証明」出版時に、有名な手紙の中で反論を呈したことはよく知られている。これらの手紙とクラーク博士の返答は通常、著作の版元に付属している。争点となる実質的な問題は、「すべてにおいて絶対的に必要であるものは空間のどの部分においても、時間軸上のどの点においても絶対的に必要である」という原理の妥当性にある。この原理はクラーク博士の論理で仮定され、彼の初回手紙に対する返答の中で明示的に述べられている。ブリストル主教の二番目の通信では、「普遍性という概念が自己存在の概念に含まれているわけではなく、あるいは直接的にその結果であるとは限らない」と述べている。「つまり、絶対的な存在は空間のどの部分においても存在することを意味するが、必ずしも全ての部分で存在することを意味しない。」クラーク博士はこの反論を効果的に処理することはできなかったように見える。全般的な通信内容は非常に興味深いものである。ブリストル主教の反論は、第一原理の洗練に依拠するよりも、論理的な機敏さよりも真実の解明に主要に依存することを確信した鋭い思考が生み出す疑問と一致している。クラーク博士の返答は十分に納得させるものとは認められずとも、議論が発端となった著作において明らかになった特異な知性の力を見せる。

12. 第8段落では、万物の自己存在し且つ原始的な原因は知性的存在であると主張している。

この命題を支持する主な論拠は、原因が効果よりも優れていることから、自己存続する存在はすべてのもののもととなる原因であり原点であるため、自身の中にすべてのものの完璧性を含んでいるべきであるという点です。そして、知性は創造の中の一部で示される完璧性の一つであると主張されます。

さらに、この完璧性が図形の変化、可分性、または物質の既知の特性のいずれかの改変ではないという理由も述べられます。これらのものは完璧性ではなく、限界であるからです。そして、宇宙の枠組みにおける設計の現れから帰納的に存在する知性的自己存続する存在の証拠が追加されています。

ただし、宇宙全体において運動の実在性に基づく独立した論拠も付けてあります。この証明を簡潔に示し、本稿の方法で適用します。

論拠は重要な説明を省略して次のように述べられます：

「世界には運動があることが明らかであり、その運動は時間と共に始まったか、あるいは永遠に存在していたかのいずれかです。もし運動が時間と共に始まったとするなら、最初の原因が知性を持つ存在であるという点については議論を認める必要があります……逆に、もし運動が永遠に存在していたとするなら、それは永遠に知性を持つ存在によって引き起こされたのか、あるいはそれが自己存続的で必要ないのか、それとも無数の連続的な通信によって存在したのかです。もし運動が永遠に知性を持つ存在によって引き起こされたとするなら、これも現在の論争における最初の原因が知性を持つ存在であるという点について議論を認めるものとなります。もしそれが自己存続的で必要ないとするなら、それは自明な矛盾となるでしょう——すなわち、運動が自己存続的であると仮定すると、物質は動いているわけではないはずです。しかし同時に、この自己存続的な運動の決定が全てにおいて一貫しているのであれば、その結果は無限の静止状態となります……もし運動が自己存続的で必要ない存在によって無数の連続的な通信によって永遠に存在していたとするなら（スパノザはこれについて矛盾を犯すように主張している）、これは以前述べた通り（この議論の第二の一般命題の証明において）明らかなる矛盾となります。したがって、運動は必然的に知性を持つ何かによって最初から引き起こされるべきです。」

以上の議論の前提は以下のようになります：

1. 運動が時間と共に始まったとするなら、最初の原因は知性を持つ存在である。

2. 運動が永遠に存在していたとするなら、それは永遠に知性を持つ存在によって引き起こされたか、あるいは自己存続的で必要ないか、それとも無数の連続的な通信によって存在したかのいずれかです。

3. 運動が永遠に知性を持つ存在によって引き起こされたとするなら、最初の原因は知性を持つ存在である。

4. 自己存続的で必要ないとするなら、物質は動いているわけではないし、同時に静止しているわけでもありません。

5. その運動は無限の連続的な伝達によって存在し、同時にそれが自己存在于内に存在せず、かつ永遠の知性のある存在により永遠に引き起こされたこともないことは偽りである。

これらの命題を表すために、次のように仮定します：
\[
\begin{aligned}
x &= \text{運動は時間から始まった（従って）}\\
1-x &= \text{運動は永遠に存在していた。}\\
y &= \text{最初の原因は知性を持つ存在である。}\\
p &= \text{運動は永遠の知性を持つ存在によって永遠に引き起こされていた。}\\
q &= \text{運動は自己存続的である。}\\
r &= \text{運動は無限の連続的な伝達によって存在していた。}\\
s &= \text{物質は静止している。}
\end{aligned}
\]
前提の等式は以下の通りです：
\[
\begin{aligned}
x &= vy.\\
1-x &= v \left\{p(1-q)(1-r) + q(1-p)(1-r) + r(1-p)(1-q)\right\}.\\
p &= vy.\\
q &= vs(1-s) = 0.\\
r(1-q)(1-p) &= 0.
\end{aligned}
\]
第4の等式より $q = 0$ であるため、$q$ の値を残りの等式に代入すると以下のシステムを得ます：
\[
\begin{aligned}
x &= vy,\\
p &= vy,
\end{aligned}
\quad
\begin{gathered}
1-x = v \left\{p(1-r) + r(1-p)\right\}, \\
r(1-p) = 0.
\end{gathered}
\]
不定符号 $v$ を消去すると、最終的な簡約システムが得られます：
\begin{eqnarray}
x(1-y)=0,\\
(1-x)\left\{pr + (1-p)(1-r)\right\}=0,\\
p(1-y)=0.\\
r(1-p)=0.
\end{eqnarray}
まず、$y$ の値を求めます。$y$ が含まれるクラーク博士の結論を求めるために、まず (1) と (2) から $x$ を消去すると以下の等式を得ます：
\[
(1-y)\left\{pr + (1-p)(1-r)\right\} = 0.
\]
次に、(4) と (5) から $r$ を消去すると以下の等式を得ます：
\[
\begin{gathered}
r(1-p) + (1-y)\left\{pr + (1-p)(1-r)\right\} = 0, \\
\therefore \left\{1 - p + (1-y)p\right\} \times (1-y)(1-p) = 0;
\end{gathered}
\]
これより
\[
(1-y)(1-p) = 0.
\]

最後に、(3)式と(6)式から$p$を消去すると、

\begin{eqnarray*}
1 - y &=& 0, \\
\therefore y &=& 1,
\end{eqnarray*}

これは必要な結論を表し、「第一の原因は知性を持つ存在である」と述べています。

次に、前提から導かれる他の結論を検討しましょう。

$y$の値を(1)式、(2)式、(3)式、(4)式に代入すると、これらは以下の連立方程式になります：

\begin{equation*}
(1-x) \lbrace pr + (l-p)(l-r) \rbrace = 0, \quad r(l-p) = 0.   \tag{7}
\end{equation*}

これらの式から$x$を消去すると、

\begin{equation*}
r(1-p) = 0, \text{ したがって } r = vp,
\end{equation*}

これは「無限に継続された運動は、永遠の知性を持つ存在によって永遠に原因付けられていた」という結論を表しています。

また、与えられた連立方程式から$r$を消去すると、

\begin{gather*}
(1-x)(1-p) = 0, \\
\text{または } 1 - x = vp,
\end{gather*}

これは「永遠に存在した運動は、ある永遠の知性を持つ存在によって永遠に原因付けられていた」という結論を表しています。

最後に、同じ原始連立方程式から$p$を消去すると、

\begin{equation*}
(1-x)r = 0,
\end{equation*}

これを解くと、

\begin{equation*}
1 - x = r(1 - r),
\end{equation*}

となり、「永遠に存在した運動は、無限に継続された通信によって存在していなかった」という結論を導きます。

この式を以下のように解くと、

\begin{equation*}
r = vx,
\end{equation*}

上記の等式は「無限に継続された運動が存在した場合、それは時間において始まった」という同値な結論を導きます。

13. 上記の四つの結論の最初と最後は、読者には互いに矛盾しているように思われるでしょう。運動が無限に連続して存在するという仮定から導かれた二つの結果、すなわち 1) 運動は永遠の知性によって永遠に原因付けられてきたこと、2) 運動は時間において始まったこと、は明らかに矛盾しています。しかし、これらは元の前提から厳格に導き出された結論です。これらの対立は論理的なものではなく、技術的には「材料的」なものと呼ばれるものです。しかしながら、この対立は前提を形式的に述べることによって公式に表現されるべきでした。つまり、「何かが永遠の知性によって外部から原因付けられた場合、それは時間において始まらない」という命題を追加すれば、現在のような結論間の対立は発生しなかったでしょう。形式的論理学は、形式的に表された関係のみに注意を払うことができます（VI.16）。したがって、特定の例では、前提間の関連性を明示的に述べる必要があるかもしれません。

ここで述べたことを説明するために、方程式 (2) と (4) に以下の方程式を追加します：
\[
px = 0
\]
これは先に述べた条件を表しています。これにより次の式が得られます：
\[
(1-x) \{ pr + (1-p)(1-r) \} + r(1-p) + px = 0. \tag{8}
\]

この式から $p$ を消去すると、単純に
\[
r = 0,
\]
が得られ、これは「運動は無限に連続して存在しない」という命題を表します。次に (8) 式における $r$ の値を代入すると、次の式を得ます：
\[
(1-x)(1-p) + px = 0, \quad \text{または} \quad 1 - x = p.
\]

これにより、解釈は「運動が永遠に存在してきた場合、それは永遠の知性によって永遠に原因付けられてきた」という命題とその逆命題が得られます。

第9命題においては、「万物の自存する原始的な原因は必然的な作用主ではなく、自由と選択を持つ存在である」と主張されています。その証明は主に彼の知性の所有と、目的的存在（設計と選択を含む）の存在に基づいています。最高の原因が最善な結果を生み出すために必然的に作用するという反論に対しては、「これは適合性と智恵による必然性であり、自然的な必然性ではない」と返答されています。

14. 第10命題においては、「自存する存在、万物の最高の原因は必然的に無限の力を有している」と主張されています。この証明の基礎は、「彼からすべてのものに与えられた力が全てを彼の意志の実行に対して困難や抵抗を生じさせることはない」ことにあると定義されています。自存する存在の無限の力とは、矛盾を含むものを「作る」という力や、「そのような力を付与された存在にとって自然的または道徳的な不完全性を示すものを行う」という力を含まないことを意味し、しかし物質の創造と、非物質的で思考能力を持つ主体の創造（運動を開始する力と意志の自由を選択する力を持ち得る）には及ぶと定義されています。この自由に関する教義に基づいて、「古い大問題、ギリシャ語で言うところの『パジョン・タ・カコン』、悪の原因や起源は何であるか？」について満足な解答を与えることができると主張しています。この論点に関しては、私は簡潔にその主張を示します。

すべての悪は、他の生物が持っている某些の能力や優れた性質の欠如である不完全な悪、または痛みや死のような自然の悪、あるいは倫理的な悪（全てのvice）から成ります。これらのうち最初のものは本来の悪ではありません。何らかの創造物が所有する能力、機能、または完璧さは神によって自由に与えられたものであるため、特定の能力や完璧さの欠如はその種類の生物にとって本来の性質から外れたものであれば、それほど大きな悪とは言えません。また、彼らが存在することもなければ生まれることもなくなること自体が本来の悪と呼ばれるべきものではありません。

私たちが自然の悪として呼ぶ第二の種類の悪は、前項の結果であるか、または全体としてはより大きな善と相殺されているか、または倫理的な悪の結果であるかによって分けられます。前者の場合、それは前項と同じく本来の悪とは言えません。後者の場合、それはより大きな善と相殺されており、したがって本来の悪とは言えません。最後に、これは罰であり、これは第三かつ最終的な種類の悪、すなわち倫理的な悪の結果です。

そして、この悪は神が他の目的のために与えた自由の濫用から生じます。それは全体の創造物の完璧さと秩序を達成するために合理的で適切なものでした。しかし彼らは神の意図と命令に反し、全体の完璧さに必要なものを自己の腐敗と破壊に使用しました。そしてこのようにして世界にはすべて種類の悪が入り込み、創造主および管理者の無限の善に何ら影響を与えずにいました。

上記の議論の主要な前提は以下のようになります：

1. すべての悪は不完全な悪、自然の悪、または倫理的な悪から成ります。
2. 不完全な悪は絶対的な悪ではありません。
3. 自然の悪は不完全な悪の結果であるか、より大きな善と相殺されているか、または倫理的な悪の結果であるかによって分けられます。
4. これは不完全な悪の結果であるか、またはより大きな善と相殺されているものは絶対的な悪ではありません。

5番目の前提は、すべての絶対的な悪が推定される悪に含まれているということです。

これらの前提を表現するために、以下のように仮定します：
\begin{align*}
  w &= \text{ 推定された悪} \\
  x &= \text{ 不完全な悪} \\
  y &= \text{ 自然的な悪} \\
  z &= \text{ 道徳的な悪} \\
  p &= \text{ 不完全な悪の結果} \\
  q &= \text{ 大きい善に補償されるもの} \\
  r &= \text{ 道徳的な悪の結果} \\
  t &= \text{ 絶対的な悪}
\end{align*}

これらの前提を原始命題とし、すべての述語が個別であり、接続詞「または」が絶対的に排他的であるとして、以下の等式があります：
\begin{eqnarray*}
 w &=& v \left\{ x(1-y)(1-q) + y(1-x)(1-z) + z(1-x)(1-y) \right\} \\
 x &=& v (1 - t) \\
 y &=& v \left\{ p (1-q)(1-r) + q (1-p)(1-r) + r(1-p)(1-q) \right\} \\
 p(l-q) + q(l-p) &=& v(1-t) \\
 t &=& vw
\end{eqnarray*}

これらの式から、符号 $v$ を個別に消去すると以下の式が得られます：
\begin{eqnarray}
 w \left\{
   1 - x(1-y)(1-z) - y(1-x)(1-z) - z(1-x)(1-y)
 \right\} &=& 0 \\
 xt &=& 0 \\
 y \left\{
   1 - p(1-q)(1-r) - q(1-p)(1-r) - r(1-p)(1-q)
 \right\} &=& 0 \\
 \left\{ p(1-q) + q(1-p) \right\}t &=& 0 \\
 t(1-w) &=& 0
\end{eqnarray}

まず、絶対的な悪が道徳的な悪にどのように依存するか、および道徳的な悪の結果について何らかの結論を形成できるかどうかを探求します。

この目的のために、$t$ を $z$ と $r$ の関数として決定しなければなりません。

したがって、符号 $w$, $x$, $y$, $p$, $q$ を消去する必要があります。この過程は簡単で、任意の式の集合を加算することで単一の式に簡略化できます。

(1)と(5)から $w$ を消去すると，
\[
t \left\{
   1 - x(1-y)(1-z) - y(1-x)(1-z) - z(1-x)(1-y)
 \right\} = 0. \tag{6}
\]
(3)と(4)から $p$ を消去すると，
\[
yqr + yqt + yt(1-r)(1-q) = 0. \tag{7}
\]
この式から $q$ を消去すると，
\[
yt(1-r) = 0. \tag{8}
\]
(2)と(6)から $x$ を消去すると，
\[
t \left\{
   yz + (1-y)(1-z)
 \right\} = 0. \tag{9}
\]
(8)と(9)から $y$ を消去すると，
\[
t(1-z)(1-r) = 0.
\]
これが要素 $t$, $z$, および $r$ の間にある唯一の関係式である。

次に、
\begin{align*}
  t &= \frac{0}{(1-z)(1-r)}   \\
    &= \frac{0}{0}zr + \frac{0}{0}z(1-r)
     + \frac{0}{0}(1-z)r + 0(1-z)(1-r)   \\
    &= \frac{0}{0}z + \frac{0}{0}(1-z)r;
\end{align*}
これは「絶対の悪は、道徳的な悪であるか、またはそれが道徳的悪でない場合も、道徳的悪の結果である」という解釈を与える。

上記の解法から得られた任意の結果には解釈が与えられる。例えば、(8)から
\begin{multline*}
  t = \frac{0}{y(1-r)} = \frac{0}{0}yr
    + \frac{0}{0}(1-y)r + \frac{0}{0}(1-y)(1-r)   \\
  = \frac{0}{0}yr + \frac{0}{0}(1-y);
\end{multline*}
これは「絶対の悪は、不完全な悪の結果であるか、または全くそのような悪でもない」という解釈を与える。

与えられた方程式から任意に提案された質問に対する様々な結論が導き出される。以下にいくつかの例を挙げる：
1. 質問 1：絶対的な悪、不完全な悪の結果、そしてそれらに対する補償のある悪との要素間に関係性は何か？
答え：「関係性は存在しない」。すべての符号を $z$, $p$, $q$ に限定すると、結果は $0 = 0$ となる。

2. 質問 2：絶対的な悪、不完全な悪、そしてそれらの結果との間にはどのような関係性が存在するか？
答え：「関係性は存在しない」。

答えは次のようになります：

最終的な関係式は、$x$, $t$, および $p$ の間の関係が以下のようになります：
\[
xt + pt = 0;
\]
従って，
\[
t = \frac{0}{p + x} = \frac{0}{0}(1 - p)(1 - x).
\]
したがって、絶対的な悪は不完全な悪の悪や不完全な悪による結果ではない。

---
質問 3. --- 自然的悪と不完全な悪およびより大きな善で補償された悪との関係を求めよ。

次のように得られる：
\begin{gather*}
pqy = 0, \\
\therefore y = \frac{0}{pq} = \frac{0}{0}p(1 - q) + \frac{0}{0}(1 - p).
\end{gather*}
したがって、自然的悪は不完全な悪による結果で補償されないものや、不完全な悪による結果ではないものである。

---
質問 4. --- 自然的悪のうち、道徳的悪でないものは絶対的な悪と道徳的悪による結果との関係は？

$y(1 - z) = s$ の場合、消去を経て次のように得られる：
\begin{gather*}
ts(1 - r) = 0; \\
\therefore s = \frac{0}{t(1 - r)} = \frac{0}{0}tr + \frac{0}{0}(1 - t).
\end{gather*}
したがって、道徳的悪でない自然的悪は絶対的な悪のうち、道徳的悪による結果であるものや、絶対的な悪ではないものである。

---
次に類似の方法で導き出された結論を示す。各結論の主題は、その説明が求められた特定の物事を示し、述語は含むべき要素を示す：

絶対的な悪が道徳的悪による結果ではないものについては、それは道徳的ではなく自然的悪である。

絶対的な悪が道徳的悪でないものは自然的悪であり、それが道徳的悪による結果である。

自然的悪が道徳的悪による結果ではないものは絶対的な悪ではない。

最後に、絶対的でない悪についての説明を求め、自然的悪と道徳的悪との関係を用いて表そう。

最終的な方程式を得ることができます：
\[
1 - t = yz + \frac{0}{0}y(1-z) + \frac{0}{0}(1-y)z + (1-y)(1-z).
\]
この方程式の直接的解釈は必要真理ですが、逆の解釈は興味深いものです。つまり、「自然かつ道徳的な悪」と「自然でも道徳的でもない悪」は絶対的な悪ではない。

この結論は必ずしも真実を表すものではありませんが、与えられた前提の中で形式的に含まれていることは確かです。

15. 同じ議論からより詳細な前提系を取り上げます。これらの前提において、「または」という粒子が絶対的な排他的性質を持たないことを仮定し、「悪の不完全性」や「自然の悪」、あるいは「道徳的悪」のいずれかであるという表現の中でこれらを一つもしくはそれ以上の特性を持つものを含むことにします。

前提として以下を採用します：

1. すべての悪（$w$）は悪の不完全性（$x$）、自然の悪（$y$）、または道徳的悪（$z$）のいずれかである。
2. 悪の不完全性（$x$）は絶対的な悪（$t$）ではない。
3. 自然の悪（$y$）は悪の不完全性（$p$）の結果であるか、より大きな善（$q$）で補償されるか、または道徳的悪（$r$）の結果である。
4. 悪の不完全性（$p$）の結果であるものは絶対的な悪（$t$）ではない。
5. より大きな善（$q$）で補償されるものは絶対的な悪（$t$）ではない。
6. 道徳的悪（$z$）は自由の濫用（$u$）の結果である。
7. 道徳的悪（$r$）の結果であるものは自由の濫用（$u$）の結果である。
8. 絶対的な悪は被疑された悪に含まれる。

通常の表現でこれらの前提を表すと、不定記号$v$を消去した後、以下の方程式を得ます：
\begin{gather*}
w(1-x)(1-y)(1-z) = 0, \tag{1} \\
xt = 0,                \tag{2} \\
y(1-p)(1-q)(1-r) = 0, \tag{3} \\
pt = 0,                \tag{4} \\
qt = 0,                \tag{5} \\
z(1-u) = 0,            \tag{6} \\
r(1-u) = 0,            \tag{7} \\
t(1-w) = 0.            \tag{8}
\end{gather*}
これらの各方程式は条件 $V(1-V) = 0$ を満たします。

自然の悪は、絶対的な悪であり、それが道徳的悪から生じたものであるか、あるいは絶対的な悪でないものである。

すべての悪は、自由の濫用から生じた絶対的な悪であり、それらが絶対的な悪でないものである。

自然の悪は、不完全な悪であり、それが絶対的な悪でないものであるか、あるいは絶対的な悪でもないものである。

絶対的な悪は、自由の濫用から生じた自然の悪であり、それらが絶対的な悪でないかつ不完全な悪でもないものであるか、あるいは絶対的な悪でもないかつ不完全な悪でもないものである。

自由の濫用の結果としての悪は、絶対的な悪であり、それが不完全な悪でないものであり、それらが自然の悪ではない余剰分野を含むすべての自然の悪である。

16. これらの例は説明のために十分です。必要であれば、読者は他の例も容易に思いつくでしょう。現在、スパニオラの証明の最も重要な部分を検討します。

定義：

1. 自己原因（causa sui）とは、その本質が存在を含むもの、すなわちその性質が存在しないことを考えることができないものを指します。

2. ある種において有限または制限された（in suo genere finita）とは、他の同じ種のものによって制限されるものを指し、例えば体は常に別の体より大きい体を思い描くことができることから有限とされます。また、思考も他者の思考によって制限されます。しかし、体は思考によって制限されず、思考も体によって制限されません。

3. 物質とは、自己内在（in se）であり、自己概念化される（per se concipitur）もの、すなわちその概念が他の何かの概念から形成されることを必要としないものを指します。

4. 属性とは、知性が物質において感知し、その本質を構成するものです。

5. 模様とは、物質の中にある感情、または他者のことによって存在し、そのことによっても認識されるものです。

6. なるほど、無限の存在を理解している。それは無限の属性を持つものであり、それぞれが永遠で無限の本質を表す。

**解説:** 私は「絶対的に無限」と言っているが、「種類においてのみ無限」であるとは言わない。なぜなら、種類においてのみ無限なものはいくつかの無限属性を持たないことを否定することもできるからだ。しかし、絶対的に無限なものは、否定を含まない何かがその本質に属するものと解釈される。

7. あるものが単独の自己の性質による必然性によって存在し、自身のみによって行動を決定されるとき、それを「自由」と呼ぶ。一方で、他者の決定によって存在と行動が一定の方法で決定されるとき、それを「必要」または「制約」と呼ぶ。

8. 永恒とは、永遠なものの単独の定義から必然的に導かれる存在そのものを指す。

**解説:** このような存在は、永遠なる真実として考えられ、それは事物の本質であるため、単なる経過時間や無始末の時間を用いて説明することはできない。

\centerline{公理}

1. すべての存在するものは、自己内在的に \emph{in se} に存在するか、または他のものの中に存在する。

2. 何かが他者によって考えられないとき、それは自身で考えられるべきである。

3. 特定の原因から与えられたとき、その効果は必然的に生じるし、逆に特定の原因を与えられなければ、効果は不可能である。

4. 效果に対する知識は原因に対する知識に依存し、またそれを含む。

5. 何かが共通点を持たないものは互いによく理解できない。または一つの概念にはもう一つの概念を導き出すことはない。

6. 真実な概念はその対象と一致するべきである。（\emph{idea vera debet cum suo ideato convenire.}）

7. 存在しないことができるものは、その本質において存在を含まない。
%-----------------------File: 016.png----------------------------
他の定義は一部の証明において暗黙的に含み、また他の公理は推論の中で仮定されている。例えば、命題I「物質はその性状よりも自然に優先する」の証明はDef.3とDef.5への単なる参照から成るが、これには以下の公理が暗黙的に仮定されているように見える：「何かを概念化するものは、そのものよりも自然に優先する」という公理。また、命題Vの証明においてはこの公理の逆も仮定されている。このような例は他にも多数存在する。したがって、倫理学の第一書の前提とされる公理や定義のみで論理的な過程だけでその全結論を導くことは不可能である。以下に続く簡潔な分析において、私は私自身が考えている実際の前提を適切な順序で提示し、これらの前提がスパノザが導き出した結論にどのように関連しているかを示すことを試みる。

17. そして私は、スパノザの証明過程においていくつか並行して行われた可能性存在の宇宙に対する分類を次のように考えています：

1. 自身の中に存在するもの$x$と他の何かの中に存在するもの$x'$に分けます。これら2つのクラスが合わせて宇宙全体を形成することから、
\[
  x + x' = 1; \quad \text{(公理I)}
\]
または
\[
  x = 1 - x'.
\]

2. 自身で概念化されるもの$y$と何か他のものによって概念化されるもの$y'$に分けます。これにより、
\[
  y = 1 - y'. \quad \text{(公理II)}
\]

3. 物質$z$とそのモード$z'$に分けます。これにより、
\[
  z = 1 - z'. \quad \text{(定義III-V)}
\]

4. 自由なもの$f$と必然的なもの$f'$に分けます。これにより、
\[
  f = 1 - f'. \quad \text{(定義VII)}
\]

5番目の定義では、原因であり自存するもの $e$ と他の何かによって引き起こされるもの $e'$ の区分が行われます。これにより、

\[
e = 1 - e'. \quad (\text{定義 I. 公理 VII.})
\]

%-----------------------File: 017.png----------------------------
そして、彼の論理は明示されたまたは仮定された原理に基づいて進められます：これらの区分は並行であり、かつ等価である。したがって、定義 III.において、本質は自身によって考えられるものと等価になります。これにより、

\[
z = y.
\]

また、公理 IV. はスパニオラが実際に適用するところでは、原因はそのことによって考えられるものと同一であることを確立します。これにより、

\[
y = e.
\]

さらに、定義 VII.において、自由なものは自存するものと同一視されます。これにより、

\[
f = e.
\]

最後に、定義 V. ではモードは他の何かによって存在するものと同一視されます。これにより、

\[
z' = x',
\]
従って、

\[
z = x.
\]

これらのすべての結果をまとめると、以下の等式列になります：

\[
x = y = z = f = e = 1 - x' = 1 - y' = 1 - f' = 1 - z' = 1 - e'.
\]

この系列内の任意の2つの項が等号で結ばれると、スパニオラによって導かれたかどうかに関わらず、彼の体系から正当な推論を表す結論が表現されます。例えば、

\[
z = 1 - e',
\]

これは彼の体系の第六の命題を表します：一つの本質は別のものによって生じることはできません。同様に、

\[
\]

$$z = e,$$
は彼の第七の命題を表しており、「物質の性質には存在することが属する」という主張を示しています。この推論過程については詳細に追及する必要はありません。スピノザは、神の性質の特性に基づく推論に特に焦点を当て、まず唯一の物質は神であることを証明しようとしました。この過程の手順において、私はいくつかの誤りを見出すことができ、これらは主に言葉の曖昧な使用によるものであり、ここでは注意が必要です。

18. 命題\textsc{v}においては、「同一性質または属性を持つ二つ以上の物質が存在することは不可能である」と示そうと試みています。その証明は以下のようになっています：もし異なる物質が一つ以上存在する場合、それらは特性またはモードによって区別されるか、します。特性によって区別される場合、同じ特性を持つ物質はただ一つだけ存在することになります。モードによって区別される場合、これらを非本質的なものとして取り除くと、実際の区別の根拠が残らないことになります。したがって、同一性質を持つ物質はただ一つだけ存在します。ここで仮定されている前提は、他の部分の論文で見られる前提とは矛盾しています。すなわち、物質、定義\textsc{iv}によれば、知性を通じて属性を介して理解されます。定義\textsc{vi}によれば、それは多くの特性を持つことができます。したがって、一つの物質は、他のいくつかの特性が同じままである一方で、その一部の特性によって別の物質と区別されることが可能であると推測できます。

命題\textsc{viii}においては、「すべての物質は必然的に無限である」と示そうと試みています。その証明は以下のようになっています：存在する唯一の物質は一つの特性を持つものであり、命題\textsc{v}によれば；そしてそれは性質上存在することを属有しており、命題\textsc{vii}によれば。したがって、その性質上は有限または無限のいずれかで存在することになります。しかし有限であるとは言えません。定義\textsc{ii}によれば、それは同じ性質を持つ別の物質によって境界を必要とするからです。これもまた必然的に存在しなければならない、命題\textsc{vii}によれば。したがって、二つの同じ特性を持つ物質が存在することになります。これは矛盾である、命題\textsc{v}によれば。したがって、物質は無限であると結論付けることができます。

この証明において、「有限」という言葉は「自体有限」（Def. ii）という表現と混同されている。つまり、何かが自体有限であるためには、他の同じ種類のものによって境界付けられている必要があるとする仮定がなされている。これは通常の用語の意味とは一致していない。

スパinoザの用語「有限」の使用は空間のみを物質の形式とし、すべての存在するものは空間の影響に過ぎないとする傾向があり、これこそが彼の体系の最終的な基礎の一端であると考えられる。

上記命題に対する最初の補遺は興味深い。その原文を引用する：「有限とは実際には否定の側面から見ればなり、絶対的な存在の主張とは無限であるという肯定の側面から見ればなり、したがって単に命題 VII からもすべての物質は無限であるべきであると結論付けることができる。」

この主張はクラークによって主張され、バトラー（XIII.11）によって反駁された原理である。おそらくこれは、先験的に無限なる存在の存在を証明しようとするすべての試みの根底にある原理であることが発見されるであろう。

上記の物質に関する一般的な性質と、神を無限の属性を持つ物質として定義することから、スパニオラの神の本性に関する特有な教義は必然的に導かれる。物質は自己存続的であり、自由であり、その性質自体が因果律を伴い、他のものたちが存在し、それらを概念化するものであるため、これらの性質は神にも主張される。彼は自己存続的で、命題\textsc{xi}に述べられている；不分割であり、命題\textsc{xiii}に述べられている；唯一の物質であり、命題\textsc{xiv}に述べられている；すべてのものが存在し、それらを概念化するものである、命題\textsc{xv}に述べられている；自由であり、命題\textsc{xvii}に述べられている；すべてのものの内在的な原因であり、命題\textsc{xviii}に述べられている。神が唯一の物質であることを証明するのは、定義\textsc{vi}から引き出されるもので、これは「神は絶対無限の存在であり、物質の性質を表す属性を否定することはできない」という宣言に解釈される。定義によって彼に与えられたすべての考えられる属性がこのように指定され、命題\textsc{v}において二つの同じ属性を持つ物質は存在しないと決定されているため、神が唯一の物質であることが導かれる。

スパニオラの「エチカ」は、彼の他の多くの著作と同じく幾何学的形式で提示されるが、この作品のシンボリカルな方法には良い実践を提供していない。もちろん、論理の列が最終的な前提が正しく決定されたときには、その方法で扱えるが、現在の場合においては、そのような扱いは、文字を使用する代わりにアルファベットを使用する以外では、元の証明で使用されるプロセスとほとんど異なる。定義として等価と定義された用語の間での遊戯のような論理はあまり見られず、主題に関連する興味のために、それらの証明がどれほど高く評価されようと、ここに数ページを割いてその解説に専念した。

19. クラーケーとスパインザの議論を読むことは、無限なる存在やその属性、そして宇宙との関係を完全に先験的に確立しようとするすべての試みが徒労であるという深い信服を得ることなしには不可能だと考えています。

このような推測の基礎となる根本原則は、我々が明確に理解できるものは存在するべきである、すなわち「\textit{\`{a} priori}」に存在を確立しようとする試みが失敗することです。その真実性が認められたとしても、有限なものがあらゆる無限の理解をすることができるでしょうか？しかし、無限なるものに対する可能性の理解は、現象的存在の制限を単なる取り除く以上の意味で認められるべきであり、それなしでは無限と永遠のことを先験的に知ることの基礎となる何らかの実質的な根拠が確立されません。スパインザはこのような知識の現実性について明示的かつ明白に主張しています：「人間の心は神の永遠で無限なる存在に関する認識を適切に持っている」（第2部、命題\textsc{xlvii}）。これは第2部、命題\textsc{xxxiv}と比較してみてください：「我々の中で絶対的な概念であっても完全な概念であっても真実であるすべてのアイデアは真実である」；また第1部、公理\textsc{vi}によれば、「真実なアイデアはその対象と一致するべきである」。さらに、このような知識は他のすべての知識の要因として位置付けられています：「理性の性質は永遠なる種類において何かを感知することである」（第2部、命題\textsc{xliv}、推論\textsc{ii})。もし人間の心には特定から普遍へと、有限なものから無限なものへと、一時的なものから永遠なるものへと見上げる傾向があるとし、この傾向は我々が感知する以上の存在や理解を超えるものについて高い確率で提示することを示唆するとすれば、その主張は少なくとも多くの心に受け入れられるかもしれません。しかし、このような推測の性格は他の原因によって部分的に説明されるべきであり、可能性的または限定的な知識に対する不満、それが我々が本当に達成できるものであると頻繁に見られるように；絶対的な確信を求め、道徳的な義務を示す十分な暗示だけが与えられたところでのみ提供されたもので、推測的な知性の要求を満たすことのないもの；あるいは現状の存在に対する不満なども含めると、これらの動機が過度に優位にあるとすれば、類推や可能性的帰納法というより慎重な手続きが軽視されるかもしれません。しかし、後者は明らかに我々の現在の状況において最も適した道です。周囲の設計の豊かな証拠から知性のある原因を推論し、自己の性質と道徳的な規定を研究することから世界の道德的指導者への概念を提起する--これらは知性がその能力や知識の材料に限られているにもかかわらず、達成可能な絶対的な確信を超える野心的な試みよりも有用です。これらの推論は古代から現在まで、啓示を除いて世界の運命が偶然と不変の運命によって放棄されていないという信念の最も堅固な基礎となっています。

第十四章。

\setcounter{chapter}{13}

\newpage

\chapter{還元法による方程式体系の分析実例（Example of Analysis）}
\label{ch:xiv}

\begin{center}
\large\textit{複合的論理条件の完全還元と具体的演繹の展開}
\end{center}
\vspace{4mm}

\#\#\# 第1章 [例の解析]

\#\#\#\# 大見出し：等式系の分析の例：単一の相当な等式 \( V = 0 \) の方法による約束化，ここで \( V \) は条件 \( V(1-V) = 0 \) を満たす。

1. 前章で用いた「物質は必然的な存在ではない」と証明するために、特筆すべき前提群を取り上げます。6番目の前提（運動が存在する）を省略し、残りの前提から導かれるいくつかの結果を検討します。これは実際にはドクター・クラークの仮定的な原理を真に受け入れるということであり、運動の存在について我々が知識を持っていないという状況を想定することです。このような前提群の一部を選択することで興味深い結果が得られる場合もありますが、現在の例は他の観点から再考されました。実際に用いられた前提は以下の通りです：

1. 物質が必然的な存在であると仮定すると、重力の性質が必ずしも存在するか否かで分岐します。

2. 重力が必ずしも存在しない場合、そして世界に支配的な知性が存在しない場合、運動は存在しません。

3. 重力が必ずしも存在する場合、真空が必要となります。

4. 真空が必要である場合、物質は必然的な存在ではありません。

5. 物質が必然的な存在であると仮定すると、世界には支配的な知性が存在しないという状況が必然的に生じます。

これらの前提を表すために以下の記号を使用します：

\begin{eqnarray*}
  x &=& \text{物質は必然的な存在です。} \\
  y &=& \text{重力の性質は必ずしも存在します。} \\
  w &=& \text{運動が存在します。} \\
  t &=& \text{重力の性質は必ずしも存在しません。} \\
  z &=& \text{世界は単なる物質であり、支配的な知性に支配されていません。} \\
  v &=& \text{真空が必要です。}
\end{eqnarray*}

（続く）

前提の表現と不定なクラス符号（$q$）の除去により、以下の方程式系を得ます：
\begin{eqnarray*}
xyt + x\bar{y}\bar{t} &= 0,   \\
                  tzw &= 0,   \\
             y\bar{v} &= 0,   \\
                   vx &= 0,   \\
             x\bar{z} &= 0;
\end{eqnarray*}
ここで、簡略化のため$\bar{y}$は$1-y$、$\bar{t}$は$1-t$と表記し、同様に他の符号も用います；従って、$1-\bar{t} = t$、$1-\bar{y} = y$などとなります。

これらの方程式の最初の項はすべて正の項のみを含むため、それらを加算することで一つの方程式を作ることができます（VIII. 定理2）：
\[
xyt + x\bar{y}\bar{t} + y\bar{v} + vx + x\bar{z} + tzw = 0,
\]
そして最初の項を条件$V(1-V) = 0$を満たすように簡略化する必要があります。

この目的のために、まず$x$と$y$に関する展開を求めます。その結果は以下のようになります：
\[
\begin{split}
  (t + \bar{v} + v + \bar{z} + tzw) xy
   + (\bar{t} + v + \bar{z} + tzw) x\bar{y}   \\
   + (\bar{v} + tzw) \bar{x}y + tzw\bar{x}\bar{y} = 0.
\end{split}
\]
そして、この展開の四つの係数を要求される条件を満たす形に簡略化することで目的が達成されます。

まず最初の係数は、$v + \bar{v} = 1$より
\[
1 + t + \bar{z} + tzw,
\]
これは単位（IX. 定理1）に簡略化されます。

二番目の係数は
\[
  \bar{t} + v + \bar{z} + tzw;
\]
そしてその簡略形（X. 3）は
\[
  \bar{t} + tv + t\bar{v}\bar{z} + t\bar{v}zw.
\]

三番目の係数$\bar{v} + tzw$は同じ方法で$\bar{v} + tzwv$に簡略化されます；そして最後の係数$tzw$には簡略化の必要はありません。したがって、展開は以下のようになります：
\[
  xy
+ \left(\bar{t} + tv  + t\bar{v}\bar{z} + t\bar{v}zw\right)x\bar{y}
+ \left(\bar{v}  + tzwv\right) \bar{x}y
+ tzw\bar{x}\bar{y} = 0; \tag{1}
\]
これが求められる簡略化の形です。

2. 現在、第\textsc{iii}章の主張した原理に従って、上記方程式において特定のシンボルの組全体が持つ関係は、その方程式を問題となる特定のシンボルに関して展開し、結果において係数が1である項のみを残すことによって導き出すことができます。ここで、$x$ と $y$ を選んだ場合、すぐに得られる方程式は
\[
xy = 0
\]
であり、これから
\[
y = \frac{0}{0}(1 - x)
\]
という関係が得られ、これは「万有引力が必ず存在するなら、物質は必要なる存在ではない」という解釈を与えます。

次に $x$ と $w$ の間の関係を求めてみましょう。方程式 (1) をこれらのシンボルに関して展開すると
\begin{multline*}
  \left(y + \bar{t}\bar{y} + tv\bar{y} + t\bar{v}\bar{z}\bar{y}
      + t\bar{v}z\bar{y}\right)xw
+ \left(y + \bar{t}\bar{y} + tv\bar{y}
      + t\bar{v}\bar{z}\bar{y}\right)x\bar{w}
\\
+ \left(\bar{v}y + tzvy + tz\bar{y}\right)\bar{x}w
+ \bar{v}y\bar{xw} = 0
\end{multline*}
を得ます。この式において $xw$ の係数のみが1に簡約されます。具体的には、
\[
t\bar{v}\bar{z}\bar{y} + t\bar{v}z\bar{y} = t\bar{v}\bar{y}
\]
\[
tv\bar{y} + t\bar{v}\bar{y} = t\bar{y}
\]
\[
\bar{t}\bar{y} + t\bar{y} = \bar{y}
\]
\[
y + \bar{y} = 1
\]
と簡約されます。このような簡約は常にこの形で行われます。したがって、得られる方程式は
\begin{gather*}
xw = 0, \\
\therefore w = \frac{0}{0}(1 - x)
\end{gather*}
であり、これは「運動が存在するなら、物質は必要なる存在ではない」という解釈を与えます。

同様に $x$ と $z$ に関して方程式 (1) を展開すると
\begin{align*}
x\bar{z} &= 0, \\
\therefore x &= \frac{0}{0}z,
\end{align*}
という関係が得られ、これは「物質が必要なる存在であるなら、世界は単に物質的であり、支配的な知性がない」という解釈を与えます。

このことは実際には第5前提の再現に過ぎませんが、これにより二つの要素を結ぶ他の関係が存在しないことが示されます。

もし要素 $x$, $w$, および $y$ の関係全体を求めると、式（1）をこれらの符号について展開し、以前の方法で進めるうちに、次のような関係が得られます。
\begin{equation*}
xy + xw\bar{y} = 0.
\end{equation*}

「運動は存在しない」という仮説のもとでの質問に対する結果を決定するために、物質の必要性と重力の必然的な存在について考察します。その結果、次のような関係が得られます。
\begin{gather*}
w = \frac{-xy}{x\bar{y}},
\\
\therefore 1 - w = \frac{x}{x\bar{y}} = \frac{1}{0}xy + x\bar{y} + \frac{0}{0}\bar{x};
\\
\text{または, } \hfill
1 - w = x\bar{y} + \frac{0}{0}\bar{x}, \text{ ただし } xy = 0.
\hfill
\end{gather*}

最初の式の直接的な解釈は、「運動が存在しない場合、物質は必然的な存在であり、重力は必ずしも存在する必要がないか、あるいは物質は必然的な存在ではない」となります。

逆に解釈すると、「物質が必然的な存在であり、重力は必ずしも存在する必要がない場合、運動は存在しない」となります。

同様の方法で、要素 $x$, $z$, および $w$ の関係全体を求めると、次のような関係が得られます。
\begin{equation*}
xzw + x\bar{z} = 0.
\end{equation*}

この式から、次のように導けます。
\begin{equation*}
z = x\bar{w} + \frac{0}{0}\bar{x}, \text{ ただし } xw = 0.
\end{equation*}

したがって、「世界は単なる物質であり、支配的な知性の影響を受けない場合、物質は必然的な存在であり、運動は存在しないか、あるいは物質は必然的な存在ではない」となります。

逆に解釈すると、「物質が必然的な存在であり、運動は存在しない場合、世界は単なる物質である」となります。

3. 当然，我々は同じ方法を、任意の複雑な仮説 \( u \) に対して適用することができます。例えば、「世界は単なる物質であり、支配的な知性 \( z \) がないが、運動が存在する」というような仮説 (\( w \)) について、元の前提に含まれる他の疑問や推測要素（例：「物質は必然的な存在である」\( x \)，「重力は物質の必然的性質である」\( y \)）に対してです。この目的のために，一般の方程式 (1) に新しい方程式を追加します：
\[
u = wz.
\]
このシステムを単一の方程式 \( V=0 \) に簡約し、\( V \) が条件 \( V(1-V) = 0 \) を満たすようにして、上記の方法で \( u \)，\( x \)，および \( y \) の関係を決定し，最後には \( u \) を \( x \) と \( y \) の展開した関数として定義します。しかし、より良い方法は第 VIII 章と第 IX 章の手法を採用することです。ここでは、いくつかの結果を簡略に示し，その導出の主要な手順を説明し，検証は読者に任せるつもりです。

先述した問題で \( x \)，\( y \)，\( w \)，および \( z \) の関係を見つけると、
\[
xw + x\bar{w}y + x\bar{w}\bar{y}\bar{z} = 0.
\]
となります。そして，\( u = xy \) と書き，一般の問題（第 IX 章）で \( x \) および \( y \) の符号を消去すると，
\[
xu + xy\bar{u} = 0,
\]
から
\[
u=\frac{1}{0}xy+0x\bar{y}+\frac{0}{0}\bar{x};
\]
よって
\[
wz=\frac{0}{0}\bar{x}\text{ かつ }xy=0.
\]
したがって，「世界は単なる物質であり、支配的な知性がないが、同時に運動が存在するならば、物質は必然的な存在ではない。」

これまでに示されたように、運動が存在する場合、物質は必要なる存在でないことが示され、したがって、上記の結論は以前より少々少ない情報しか提供しない。しかし、その結論は提案された問題に対する適切かつ完全な答えである。問題は特定の複合仮説の結果を単純に決定することであった。

従って、限られた前提から容易に、多くの追加的な推論を導くことができる。与えられた条件の中で、基本命題の任意の組み合わせが含まれる推論を導くことができる。もし条件が前提と矛盾するものである場合、その事実は解の形式によって示される。方法により、提案された条件を表現するシンボルの組み合わせに対する値は0となる。一方で、条件の満足が関係を探している命題群に制約を加えない場合、つまりこれらの命題の任意の組み合わせが同様に可能である場合、その事実も解の形式によって示される。それぞれの場合の例は添えられている。

通常の方法で、物質が必要なる存在であり、かつ運動が存在するという条件から、世界は単なる物質であり、支配的な知性がない、そして重力は必然的に存在するという命題に対する結果を求めると、方程式
\begin{equation*}
xw = 0,
\end{equation*}
を得る。これは提案された条件が前提と矛盾していることを示し、従って満たすことができないことを意味する。

もし物質が必然的な存在でないことを条件とし、同時に運動が存在することを考慮に入れ、先に述べた要素に対してすなわち、監督する知性の欠如と重力の必要性に対し探求すると、以下の結果を得ます：

\begin{equation*}
(1-x)w = \frac{0}{0}yz + \frac{0}{0}y(1-z) + \frac{0}{0}(1-y)z + \frac{0}{0}(1-y)(1-z),
\end{equation*}

これは直訳すると以下のようになります：

もし物質が必然的な存在でないかつ運動が存在するならば、世界はただ物質的であり監督する知性の欠如と重力の必要性を伴うか、あるいはこれらの結果のいずれかが成立し他方は否認され、またはこれら全てが真ではない。従って、四つの可能な組み合わせのうち少なくとも一つは必然的に真であり、その中でただ一つのみが真であるという状況において、どの組み合わせも真であることが主張されます。このような結果は単なる真理であり、必要な真理に過ぎません。しかし、提出された問題に対する唯一の答えを含んでいます。

これらの適用について過怠を非難する必要はありません。自己の詳細を十分に掘り下げることで、一般原則と方法への信頼性を確立することが必要な場合もあります。現在の調査の主題においてこの目的が達成されると、それによって得られたが後に無用となった全ては忘れられましょう。

%-----------------------File: 028.png----------------------------
%-----------------------File: 029.png----------------------------
第15章

\newpage

\chapter{アリストテレス論理学とその現代的拡張の批判的検討（Aristotelian Logic）}
\label{ch:xv}

\begin{center}
\large\textit{伝統的三段論法全式のブール代数による完全包摂と限界の超越}
\end{center}
\vspace{4mm}

\#\# アリストテレス論理とその現代の拡張

1. アリストテレスの論理体系は、詳細は修正されても、基本的な特徴は変わらず、学術教育において極めて重要な位置を占めています。そのため、この体系の性質についての説明や、主な問題点についての簡潔な議論が本著作の中で求められるべきであると考えられます。私は、狭き門や厳格な批判的態度でアリストテレスの理論に接するつもりはありません。むしろ、真実を全ての他の考慮事項よりも重視し、古代の学説の実質的な性質を明らかにし、その範囲と充足性に関する一般的な誤解を取り除くことを目指すのです。

アリストテレスやその類似の論理体系における精神と手続きにおいて、 essentials として重要視すべきは、推論可能な形態の分類を試みたことと、これらの形態が個別または総合的に、ある公理的な性質を持つ一般原理に参照されることである。アリストテレスの「命題」：属について肯定または否定されたものは、その属下に含まれる任意の種にも同じ意味で肯定または否定される。このような一般原理については、私は次のように述べることができる。これらの原理は直接的に、あるいは抽象的な形で、その原理を説明するための推論を述べており、その推論を述べることにより、その有効性を主張しているか、あるいは表現の中に技術的な用語が含まれており、定義を行うことで最終的には同じ点に到達する。すなわち、推論可能な形態の抽象的な声明である。したがって、分類の概念はこれらの体系において支配的な要素となっている。

さらに、これらは論理を二つの大きな分野に分解することを示している。一つは命題的推論の扱いに関し、もう一つは仮定的または条件付き推論に関しである。この区別は本作における一次と二次の推論との区別のほぼ同じものである。命題的推論の理論についての議論は、通常の論理学のすべての一般的な著作において、仮定的推論の理論よりもはるかに充実しており、主に古代の学説的な区分と、帰納的推論の canon に関することである。ここでは後者の適用のみに注意を向けなければならない。

2. 分類的な命題は、以下の四種類に分類される：
\begin{center}
\begin{tabular}[h!]{ll@{\hspace{0in}}cl}
 & & &\multicolumn{1}{c}{{\textsc{型}}}\\
第1. & 全称命題（肯定）  & 命題： & 全ての $Y$ のものは $X$ のものである。\\
第2. & 全称命題（否定）  & 同様に & 何も $Y$ のものは $X$ のものではない。\\
第3. & 特称命題（肯定）  & 同様に & 何か $Y$ のものは $X$ のものである。\\
第4. & 特称命題（否定）  & 同様に & 何か $Y$ のものは $X$ のものではない。\\
\end{tabular}
\end{center}

これらの形式に対して、最近では四つの新たな形式が加えられ、全体として八つの形式を形成するが、これらは六つに簡約され、論理学者であるデ・モラン教授によって詳細かつ能力的に議論された関係の下で存在する（次節を参照）。上記の体系に対し、Sir W. ハミルトンによって世界に提示された異なるスキームが提供されており、彼の方法は論理学の専門家たちから非常に高い評価を得ている推論法の基礎となっている。

形式論理の過程は、上記の命題系に関連して二種類と説明される：「変換」および「三段論法」である。変換とは、上記のどの命題も等価な形で表現し、項の順序を逆転させるという意味である。一方、三段論法とは、共通項が主題または述語として存在する二つの命題から、推論的に含意される第三の命題を導き出し、それが「結論」であるという意味である。多くの論理学者は、これらの過程、そして一部では単一の三段論法プロセスさえもが、推理の普遍的な型を提供し、論理学者の著作に分類された特定のモデルに従って、意識的にまたは無意識的に思考が進行することであると主張している。

3. 私が採りたい道は、現在の論文の原理に基づいて推論と換位の過程を最も一般的な方法で展開し、そのような推論と換位を論理の体系に関連付けて見ることにより、最終的な思考法則に基礎を置く論理システムの構築を目指して、それらの真の位置と本質的な性格を決定することである。

8つの基本命題の表現は以下の通りである：
\begin{center}
\begin{tabular}[h!]{ll@{\hspace{0in}}ll}
1. &すべての $Y$ が $X$ である,              &$y = vx$.\\
2. &すべての $Y$ が $X$ でない,               &$y = v(1-x)$.\\
3. &一部の $Y$ が $X$ である,             &$vy = vx$.\\
4. &一部の $Y$ が $X$ でない,         &$vy = v(1-x)$.\\
5. &すべての $Y$ でないものが $X$ である,          &$1-y = vx.$&(1)\\
6. &すべての $Y$ でないものが $X$ でない,           &$1-y = v(1-x)$.\\
7. &一部の $Y$ が $X$ である,         &$v(1-y) = vx$.\\
8. &一部の $Y$ が $X$ でない,     &$v(1-y) = v(1-x).$\\
\end{tabular}
\end{center}

これらの形式を参照する際には、論理量の概念、「普遍的」および「個別的」という形容詞と、「肯定的」および「否定的」という質の概念を、命題の語句に適用することにより、直ちに説明される方法で便利である。これらはそれぞれの命題自体ではなく、論理量の語句に対して使用する。

したがって、「すべての $Y$ の」は普遍的肯定的な語句とみなす。「$Y$ の」または「ある $Y$ の」は個別的肯定的な語句とみなす。「すべての非-$Y$ の」は普遍的否定的な語句とみなす。「ある非-$Y$ の」は個別的否定的な語句とみなす。表現「ない $Y$ の」は命題の語句として適切なものではない。実際、命題「ない $Y$ のは $X$ のである」という真の意味は「すべての $Y$ のは非-$X$ のである」と解釈されるからである。したがって、その命題の主詞は普遍的肯定的な語句であり、述語は個別的否定的な語句とみなす。

「人」および「非人」の概念間の実在する区別は明らかである。この区別の適用については、上記のように肯定的および否定的な指定を使用するが、命題の語句に対する用語への語源学的な適切性には強調しない。肯定的および否定的という指定はより適切であったかもしれないが、前者は本稿の他の部分で既に固定された意味を持つために使用されていない。

これらの符号化された形式から直ちに導かれる変換法則が得られる。例えば、「すべての $Y$ のは $X$ のである」という命題を表す式
\[ y = vx \]
から、$v$ を消去することで
\[ y(1-x) = 0 \]
を得る。これを $1-x$ について解くと
\[ 1-x = \frac{0}{0}(1-y) \]
となる。この解釈は「すべての非-$X$ のは非-$Y$ のである」という命題を意味する。

これは「負の変換」と呼ばれる例である。

同様に、方程式
\[ y = v(1-x), \]
が「No $Y$'s are $X$'s」（つまり、「何も$Y$は$X$ではない」）を表すとき、
\[ x = \frac{0}{0}(1-y), \]
と解釈され、これは「No $X$'s are $Y$'s」（つまり、「何も$X$は$Y$ではない」）を意味する。これは単純な変換の例であり、実際には前記の例で示された変換と同じ種類である。論理学者が指摘したすべての変換の例は、上記のいずれかの種類または命題の項の単なる移動（質を変更しない）という過程のみを含むものであり、例えば
\begin{align*}
  vy &= vx \quad (\text{つまり、「Some $Y$'s are $X$'s」（一部の$Y$は$X$である）}) \\
\intertext{を}
  vx &= vy \quad (\text{つまり、「Some $X$'s are $Y$'s」（一部の$X$は$Y$である）})
\end{align*}
と変換する場合などである。あるいは、これらの過程と補助的な制限過程の組み合わせが含まれる場合もある。例えば、方程式
\begin{align*}
  y &= vx \quad (\text{つまり、「All $Y$'s are $X$'s」（すべての$Y$は$X$である）}) \\
\intertext{を$v$で乗じることにより}
  vy &= vx \quad (\text{つまり、「Some $Y$'s are $X$'s」（一部の$Y$は$X$である）}) \\
\intertext{と導き、これより}
  vx &= vy \quad (\text{つまり、「Some $X$'s are $Y$'s」（一部の$X$は$Y$である）})
\end{align*}
を得る場合などである。

この例では、制限過程が移動過程よりも先に実行される。

これらの例から、変換は論理においてもっと広範な過程の特定の応用であることが見て取れます。この作業の中で多くの例が挙げられており、その過程は任意の命題中のどの要素も、いかなる複雑さを有しても、残りの要素に対する論理的関数として決定することを目指しています。

従来の注意を主に主題と述語に限定するのではなく、これらいずれかに入力される任意の要素やその組み合わせを取り上げることができます。新しい命題の「主題」としてそれらを選択し、与えられたデータに基づいてその述語を決定することができます。また、従来の変換の法則が認識しているよりも、上記で列挙された単純な命題の形式にもこのような方法の適用が一部認められるということも述べることができます。

例えば、方程式
\[
y = vx, \text{ これはすべての } Y \text{'s が } X \text{'sであることを表します},
\]
はさらに以下の3つの命題を導きます：
1. $y(1-x) = 0.$ つまり、存在しない $Y$'s はすべて $X$'s でない。
2. $\displaystyle 1 - y = \frac{0}{0}x + (1 - x).$ これは、存在しない $Y$'s はすべて存在しない $X$'s であるとともに、ある不定な量の $X$'s の残りを含む。
3. $\displaystyle x = y + \frac{0}{0}(1 - y).$ これは、すべての $X$'s はすべての $Y$'s とある不定な量の存在しない $Y$'s の残りを含む。

これらの結論は確かに与えられた命題を他の等価な形に置き換えているに過ぎませんが、従来の「否定的変換」の方法の役割はこれだけです。

さらに、これらの変換過程は基本的なものではなく、より単純な過程の組合せである。これらは、最終的な法則や公理に直接依存する符号的道具を使用する方法よりも一層単純である。この観察は、次に考察する三段論法の場合にも同様に適用される。

5. 三段論法の本質は具体的な例から最もよく理解できる。
例えば、以下の命題を考える：
\begin{align*}
  &\text{すべての }X\text{'s は }Y\text{'s である},\\
  &\text{すべての }Y\text{'s は }Z\text{'s である}.
\end{align*}

これらの命題から、以下の結論を導くことができる：
\begin{equation*}
\text{すべての }X\text{'s は }Z\text{'s である}.
\end{equation*}
これは三段論法による推論である。項 $X$ と $Z$ を極端と呼び、$Y$ を中間項と呼ぶ。一般に三段論法の機能を定義することができる。与えられた二つの命題が（1）表で示される種類であり、一つの中間または共通項 $Y$ が含まれており、その一つの命題ではこの中間項 $Y$ は極端 $X$ と結び付けられ、もう一つの命題では極端 $Z$ と結び付けられている場合、極端 $X$ と $Z$ の関係を求めよ。中間項 $Y$ は肯定的な形で現れる可能性があり、すべての $Y$'s, 一部の $Y's$; または否定的な形で現れる可能性があり、すべての非-$Y$'s, 部分的に非-$Y's$ であり、これらのいずれかの命題においては、他の命題における具体的な形式に関係なく、そのいずれかの形を取ることができる。

特に具体的な例において、この論文の方法で三段論法を解決することは非常に簡単です。その解決は、命題の系統を簡略化する過程の特定の応用形に過ぎません。先ほど与えられた例では、

\begin{gather*}
x = vy,\\
y = v'z;
\end{gather*}

この式から代入すると、

\begin{equation*}
x = vv'z,
\end{equation*}

これが、

\begin{equation*}
すべてのXはZである。
\end{equation*}

と解釈されます。または、(VIII.7)で開発された方法に従って厳密に進めると、

\begin{equation*}
x(1-y)=0,\quad y(1-z)=0.
\end{equation*}

これらの式を加え、$y$を消去すると、

\begin{equation*}
x(1-z)=0;
\end{equation*}

となり、

\begin{equation*}
x = \frac{0}{0}z,
\end{equation*}

すなわち「すべてのXはZである」と解釈されます。同様にして、他の任意の場合も扱うことができます。

6. 特定の例を考慮から去り、

三段論法が全て簡略化される一般形について検討しましょう。

\begin{center}
\textsc{命題I}

\emph{三段論法の一般的な規則を導く。}
\end{center}

ここでは、三段論法の一般的な規則とは、主語と述語が量と質においてどのような組み合わせでも可能である前提に対して適用可能な規則を指します（二つの普遍的項が同一の命題内で組み合わされる場合を除く）。従って、許容される命題の形式は、(1)で表される表から得られるものです。

$X$ と $Y$ を最初の前提に入力する要素または対象とするもの、$Z$ と $Y$ を第二前提に入力する要素または対象とするものとして考えます。二つの基本的に異なる場合が現れます。
$Y$ を含む語句は、似た質的性質を持つか、異なる質的性質を持つかのいずれかです。両方とも「$Y$」について述べているときや、両方とも「非-$Y$」について述べているときは似た質とみなします；一方が「$Y$」について述べており、他方が「非-$Y$」について述べているときは異なる質とみなされます。これらの前提の組み合わせにおいて最初の条件を満たす任意のペアは、以下の方程式で表現できます：
\begin{align*}
 \tag{1} vx &= v'y,\\
 \tag{2} wz &= w'y;
\end{align*}
なぜなら、シンボル $y$ を「すべての $Y$」または「すべての非-$Y$」を表すために使用できるからです。このシンボルの解釈は完全に慣習的であるためです。もし $y$ を「すべての非-$Y$」として使用する場合、$1-y$ が「すべての $Y$」を表します。そして、他の変更はありません。また、シンボル $x$ と $z$ も同様に自由に扱うことができますので、方程式 (1) と (2) を適切な解釈を与えることで、それぞれの語句の質によって依存する前提のすべての組み合わせを表現できます。さらに、シンボル $v$, $v'$, $w$, $w'$ の適切な解釈を仮定することで、量的な観点からもすべての種類を表現することができます。例えば、$v=1$ とし、$v'$ を不定なクラスとして表す場合、方程式 (1) は通常の意味での普遍的命題を表します。つまり、普遍的主語と部分的述語を持つ命題を表します。実際には、前提のいずれかに主語や述語を与えたい量（学術的な意味でのこの用語を使用）は何でも与えることができます。ただし、仮定によれば、両方とも普遍的であることはないとします。したがって、システム (1), (2) は、中間項が似た質を持つ前提のすべての組み合わせを完璧な一般性で表現します。

7. 分析を適用される方程式と同様に一般性を持たすため、本書の方法により（1）と（2）から$y$を消去し、$x$、$1-x$、および$vx$を$z$と$v$、$v'$、$w$、$w'$の符号で表す表現を求めることにしましょう。これにより、結論の主題に関するすべての可能な形が含まれます。$v(1-x)$の形式は除外されます。なぜなら、前提から与えられた「$vx=$一部の$X$」という解釈からは、「$v(1-x)=$一部の非-$X$」を解釈することはできないからです。符号$v$が「一部」という意味で使用される場合、その符号は前提と結びついた項のみに適用されます。

分析の結果は以下の通りです：
\begin{multline*}
  \tag{I.}
    x = \bigl[vv'ww' + \frac{0}{0}\{
     vv'\bigl(1-w\bigr) \bigl(1-w'\bigr)
   + ww'\bigl(1-v\bigr) \bigl(1-v'\bigr)
   +    \bigl(1-v\bigr) \bigl(1-w\bigr)\}\bigr]z\\
  + \frac{0}{0}\{vv'\bigl(1-w'\bigr) + 1-v\} \bigl(1-z\bigr),
\end{multline*}
\begin{multline*}
  \tag{II.}
    1-x = \bigl[
     v\bigl(1-v'\bigr) \{ww' + \bigl(1-w\bigr) \bigl(1-w'\bigr)\}
   + v\bigl(1-w\bigr)w'\\
+ \frac{0}{0}\{
     vv'\bigl(1-w\bigr) \bigl(1-w'\bigr)
   + ww'\bigl(1-v\bigr) \bigl(1-v'\bigr)
   + \bigl(1-v\bigr) \bigl(1-w\bigr) \} \bigr]z\\
  +\bigl[v\bigl(1-w\bigr)w'
 + \frac{0}{0}\{vv'\bigl(1-w'\bigr) + 1-v\} \bigr] \bigl(1-z\bigr),
\end{multline*}
\begin{multline*}
  \tag{III.}
    vx = \{vv'ww' + \frac{0}{0}vv'\bigl(1-w\bigr) \bigl(1-w'\bigr)\}z
   + \frac{0}{0}\bigl(1-w'\bigr) \bigl(1-z\bigr).
\end{multline*}

これらの各表現の第二項には、$z$が因子である一つと、$1-z$が因子であるもう一つの項が含まれています。しかし、論理的推論において形式上は結論で対立するクラス（$Z$と非-$Z$）を一緒に含めることは許されません。

したがって、これらの方程式の第二項が単項式に簡約される条件を特定し、その種類の推論の規則を定めるために、それらの第二項が単項式に簡約される条件を明らかにする必要があります。

最も単純な形は（III）であり、$w'=1$ の場合に一項に簡約される。その式は次のようになります：
\[ vx = vv'wz, \]
左辺は最初の前提の外延と同一であり、右辺は第二前提の外延と同じ量と質を持つ。

なぜなら $w'=1$ であるから、（2）の右辺は中項 $y$ を含むが普遍的である。従って仮定により左辺は特殊であり、したがって（3）の右辺も同じ符号 $w$ の係数を含むため、その右辺もまた特殊である。これより以下の法則を得る。

推論の条件。
少なくとも一つの中項が普遍的であること。

推論の規則。
外延を等しくする。

方程式（I）と（II）の分析からさらに明らかになるように、中項が同一の量である場合に上記は推論の唯一の条件であることが示されます。すなわち、（I）の右辺は $w'=1$ かつ $v=1$ の場合に一項に簡約され、（II）の右辺は $w'=1$, $v=1$, $w=1$ の場合に一項に簡約されます。これらの各々の場合において、$w'=1$ であることが必要であり、それが先に指定された唯一の十分条件です。

次に、中項が異なる量である場合について考慮します。その前提は以下の形で表すことができます：
\begin{align*}
 vx &= v'y, \tag{4} \\
 wz &= w'(1-y). \tag{5}
\end{align*}
また、$y$ を消去し、$x$, $1-x$, および $vx$ の表現を決定すると、次のようになります：
\begin{multline*}
 x = \bigl[ vv'(1-w)w' + \frac{0}{0} \bigl\{ ww'(1-v) + (1-v)(1-v')(1-w) \\
  + v'(1-w)(1-w') \bigr\} \bigr] \\
+ \bigl[ vv'w' + \frac{0}{0} \left\{ (1-v)(1-v') + v'(1-w') \right\} \bigr](1-z). \tag{IV}
\end{multline*}

%-----------------------File: 038.png----------------------------
\begin{align*}
1 - x = & \bigl[ww'v + v(1-v')(1-w) + \frac{0}{0} \{ww'(1-v) \\
& + (1-v)(1-v')(1-w) + v'(1-w)(1-w')\}\bigr]z \\
& + \bigl[v(1-v') + \frac{0}{0} \{v'(1-w') + (1-v)(1-v')\}\bigr](1-z).
\tag{V.}
\end{align*}
\begin{multline*}
vx = \{vv'(1-w)w' + \frac{0}{0} vv'(1-w)(1-w')\}z \\
+ \{vv'w' + \frac{0}{0} vv'(1-w')\}(1-z). \hfill \tag{VI.}
\end{multline*}
今、（VI）の第二項は $z$ に関して単項式になります。$w = 1$ の場合、
\[
vx = \{vv'w' + \frac{0}{0} vv'(1-w')\}(1-z);
\]
この第二項は第一前提における極端 $wz$ に対して量と質において逆です。
なぜなら $w = 1$ より、$wz$ は普遍的であることがわかります。しかし因子 $vv'$ が示すように、その付属する項は特別であり、仮定により $v$ および $v'$ が両方とも 1 に等しいわけではありません。したがって、以下の推論の条件を導きます：

\textsc{第一推論条件}---\emph{少なくとも一つの普遍的極端。}

\textsc{推論法則}---\emph{その極端の量と質を変更し、結果を他の極端に等しくする。}

さらに、（V）の第二項は $v' = 1$, $w' = 1$ の場合に単項式になります。このとき、
\[
1 - x = \{vw + \frac{0}{0}(1-v)w\}z.
\]
$v' = 1$, $w' = 1$ より、前提の中項は両方とも普遍的であり、したがって極端 $vx$, $wz$ は特別です。しかし結論において、$x$ を含む極端は第一前提における極端 $vx$ に対して量と質において逆であり、$z$ を含む極端は第二前提における対応する極端 $wz$ と一致します。したがって、以下の一般の法則を導きます：

%-----------------------File: 039.png----------------------------
\textsc{第二推論条件}---\emph{二つの普遍的中項。}

\textsc{推論法則}---\emph{いずれか一方の極端の量と質を変更し、結果を他方の極端に等しくする。}

異なる中項の場合には、論理推論の他の条件やルールは上記のもの以外存在しない。したがって、方程式（IV.）は$s=1$および$v = 1$のとき、$syllogistic$ $conclusion$ の形に帰着されても、これによって新たな推論の条件を確立するわけではない。なぜなら、これまで述べた通り、単なる条件$s = 1$または$v = 1$が十分であるからである。

8. 以下の例は、上記で与えられた三段論法の一般規則を十分に例証するものである：
\begin{enumerate}[label=\textbf{\arabic*.}]
\item 前提：すべての $Y$ は $X$ である。すべての $Z$ は $Y$ である。\\
これは第1の場合に属する。すべての $Y$ が普遍的中項である。両極端を等置することにより、結論として以下を得る：
\begin{center}
すべての $Z$ は $X$ である。
\end{center}
普遍項「すべての $Z$」が主語となり、特殊項「（一部の）$X$」が述語となる。

\item 前提：すべての $X$ は $Y$ である。いかなる $Z$ も $Y$ ではない。\\
これらの前提の適切な表現は次の通りである：
\begin{center}
すべての $X$ は $Y$ である。\\
すべての $Z$ は非 $Y$ である。
\end{center}
これらは第2の場合に属し、推論の第1条件を満たす。第1前提の中項 $Y$ は肯定・特殊であり、第2前提の中項「非 $Y$」は否定・特殊である。「すべての $Z$」を普遍的極端として選び、規則に従ってその量と質を変更すると「一部の非 $Z$」が得られ、これを他方の極端「すべての $X$」と等置すれば、
\begin{center}
すべての $X$ は非 $Z$ である（すなわち、いかなる $X$ も $Z$ ではない）。
\end{center}
他方の普遍的極端から始めて同様に進めれば、同値な結果として以下を得る：
\begin{center}
いかなる $Z$ も $X$ ではない。
\end{center}
\end{enumerate}

3. すべての $Y$ は $X$ である。
   すべての非-$Y$ は $Z$ である。

ここでも中項は質が異なります。従って、前提は第2例に属し、二つの普遍的中項があるため、推論の第二条件が満たされます。規則により最初の外延を量と質を変更すると、(一部の) $X$ すべてが非-$X$ となります。これを他の外延と等しいと仮定すると、

\begin{center}
  すべての非-$X$ は $Z$ である。
\end{center}

逆順でも同様の結果を得られます：

\begin{center}
  すべての非-$Z$ は $X$ である。
\end{center}

最後の二つの例の結論は、学術的な論理体系において有効と認められません。この体系では、命題の主項が肯定的であることが求められます。しかし、これらの結論自体は完全に正当であり、それらを決定する規則は推量法の最も一般的な原理を形成します。これらの結論を導く調査過程は、不要な複雑さのように思われるかもしれませんが、事実上はより容易な方法で得られ、何らかの記号的道具を使用せずにでも得られる可能性があります。私の目的は、最も一般的な方法で調査を行い、完全に Exhaustive な分析を行うことでした。この目的を達成するために、手法の簡潔さや冗長さは問題になりません。実際、分析は推量法そのものではなく、より一般的な命題の組み合わせに対する分析です。符号 $v$, $v'$, $w$, $w'$ に任意のクラス解釈を与えることが許されます。

この言及を例示するために、以下のような仮想的なケースを適用します：

ある種類の布地が白と緑で縞模様である場合と黒と黄色で縞模様である場合について、調査を行い、次のように観察しました：

1. すべての白と緑で縞模様である布地はまた黒と黄色で縞模様であり、逆もまた然り。

2番目には、赤とオレンジで縞模様が施されたすべての部分も青と黄色で縞模様が施されており、その逆もまた真実である。

次に、緑色でマークされた部分が白、黒、赤、オレンジ、青に関連してどのような状態にあるかを決定する必要があるとする。ここで、$v=$ 白、$x =$ 緑、$v' =$ 黒、$y =$ 黄、$w =$ 赤、$z =$ オレンジ、$w' =$ 青と仮定すると、前提の表現は以下のようになる：

\begin{align*}
  vx &= v'y,   \\
  wz &= w'y,
\end{align*}

これはシステム (1) (2) に一致する。方程式 (I.) から次の結論が導かれる：

緑色で縞模様が施された部分は、オレンジ、白、黒、赤、青のいずれかまたはすべてを含む。これらの特性を持つ全ての部分は緑色で縞模様が施された部分に含まれる；あるいは、オレンジ、白、黒のみを含み、赤や青を含まない；あるいは、オレンジ、赤、青のみを含み、白や黒を含まない；あるいは、オレンジのみを含み、白や赤を含まない；あるいは、白と黒のみを含み、青やオレンジを含まない；あるいは、白もオレンジも含まない。

この結論の性質を考えると、その表現 (I.) は必要以上に複雑なものではなく、またその表現が導かれる分析も無駄に複雑なものではないと評価できる。

9. この章で示した三段論法の形式は重要な観察を可能にする土台を提供している。私たちは、その二つの主要分野のそれぞれにおいて、規則や方法の決定に関連する議論全体が、前提が同じ中項を持つ場合の系（1）、（2）の方程式対と、中項が異なる場合の系（4）、（5）の方程式対を分析することによって簡素化できることが確認された。さらに、その分析は言語の構造から直接導き出される一般法則に基づく方法によって行われたものであり、第II章で確認され、第III章で人間の心の操作の研究により支持され、そして第VIII章で任意の方程式系を分析するための適用可能性が証明された。したがって、三段論法は実質的な推理の基本型であるかどうかという問いに対する解決策を得ることができる。すなわち、その法則の研究が推演的論理学の全体的研究と一致しているかどうかを確認できる手段がある。

ここで、我々が取り組んできた方程式系の分析において、その事実が示唆される兆候がないことは考えられない。三段論法が唯一の重要な推理過程であるにもかかわらず、その現象がこの高貴な優越性を示す兆候を提供していないと考えることは不可能だ。しかし、この章で検討した特定の方程式系の議論において、すべての命題を表現する方程式系の分析が含まれているという兆候も見られない。しかしながら、論理学に関する著者は、三段論法の優越性だけでなく、推演的推理における普遍的な充足性についてほぼ一致して主張している。

アーサー・ウィテイエル大主教の言葉は常に明瞭で確定的であり、論理学に関しては特に注意を要する。彼は次のように述べている。「論理学は、推理の文法として存在し、規則的な三段論法を別の議論方法に置き換えるための特別な形態として提示されるのではなく、すべての正しい推理が最終的にその形式に還元できる形である」と。

また、ミル卿は『論理学体系』の一章「推論または三段論法」において、通常の三段論法の形式を列挙した後で述べている。「すべての有効な推論、つまり一般命題が以前に認められたものから他の同等またはそれ以下の一般性の命題を導くための推理は、これらの上のいずれかの形で示すことができる」と。さらに、「従って、我々は一般的な論理学者の見解に従って、第一格の二つの基本形式をすべての有効な推論の普遍的な型と考えることができる」と述べている。

これらの見解に基づき、論理学の科学は他のすべての科学が抱える不完全性や進歩の条件から免れるものであると主張されている。そして、その起源は昔の一つの大いなる心の産物であり、大胆な比喩としてパラスの神話的誕生に例えられている。

三段論法は排除の一種であるため、私共が直面している問題は明らかに以下の二つに分岐します：1つ目は、すべての排除は必ずしも三段論法に帰着できるか否か；2つ目は、推論は適切には常に排除のみで構成されているものと見なせるか否か。慎重に検討した結果、私は最初の問題に対する真実なる答えを以下のように考えます：理論的には、任意の命題をどのように解釈し組み合わせるかによりますが、三段論法の原理によって排除が可能となるように分解と統合は常に可能です。しかし、その過程はしばしば制約があり自然でなく、非三段論法的な操作を伴うことがあります。二つ目の問題に対する私の答えは、推論は任意の意味での制限なしに排除のみで限定されるべきではないということです。定義が存在するとしても、それは思考の法則に基づく方法の総和よりも少ないものにはならないべきです。そしてその方法の中では、非常に重要な排除の過程も三段論法的な操作の一部に過ぎません。

多くの誤解は、論理の性質や推論の範囲についての理解において、私はそれを避けてはならないと思いますが、これら全ての真理を単純さと直感的な確信に特徴づけられるものとして扱う傾向にあるように思われます。他の科学における真理との関係や、論理的推論の芸術における一般的方法との関係について調査することなくです。

アリストテレスの「すべてと何も」に関する定説は自明な原理ですが、それは推理能力の最終的な法則の中に見出されるものではありません。他のいかなる自明な法則も追跡できる普遍的な法則から派生し、厳格な科学的進化の秩序で導き出すことができるものでもありません。確かに、どの学問においても基本的な真理は通常最も単純なものですが、その単純さがそれらを基本的なものとして扱うべき理由となるわけではないことが明らかです。その理由は、それらが支える構造の性質と範囲に求められるべきです。

この視点から見れば、レーニッツは「矛盾の原理」が論理において重要な位置を持つことを適切に判断したように思います（Nouveaux Essais sur l'entendement humain. Liv.\textasciitilde{}IV, cap.\textasciitilde{}2. Theodicée Pt.\textasciitilde{}I, sec.\textasciitilde{}44）。実際、その思考法の法則の転換表現である原理の結果はすでに示されています（III. 15）。しかし、推論の性質やその構成要素について十分な議論がなされたのであれば、導入法に関する主題については別の部分で少し注意を向けられるかもしれません。

本章において、仮定的命題の通常の扱いは、推量的命題よりも著しく欠陥があることが指摘されています。一般的に「仮定的三段論法」と呼ばれるものは、実際には真の三段論法ではないようです。

以下のような議論を考えてみましょう：

\begin{tabbing}
前提--- \= \\
         \>もし $A$ が $B$ であるなら、$C$ は $D$ である,   \\
         \>しかし $A$ が $B$ である,   \\
         \>従って $C$ は $D$ である;
\end{tabbing}
ここで、$X$ は「$A$ が $B$ である」という命題を表し、$Y$ は「$C$ が $D$ である」という命題を表します。すると、前提には2つの要素しか含まれていませんが、三段論法は必ず3つ以上の要素を含むべきです。

真の仮定的三段論法としては以下のようなものがあります：

\begin{tabular}{rl}
            & もし $X$ が真なら、$Y$ が真である;   \\
            & もし $Y$ が真なら、$Z$ が真である;   \\
$\therefore$ & もし $X$ が真なら、$Z$ が真である。
\end{tabular}

この作品の前半部で二次的な命題について議論した後では明らかですが、仮定的三段論法の形態は、推量的三段論法と完全に一致するべきです。個別の命題、例えば「時々もし $X$ が真なら、$Y$ が真である」というようなものも導入できますし、この章で推量的三段論法について述べた条件や推理の規則は、仮定的命題に対応する場合にも完全に適用可能であり、その解釈は変更することなく通用します。

11. それでは、学術的論理の性質と範囲に関する最終的な結論は何に達することができるでしょうか。私は以下のようだと考えます：それは科学ではなく、科学的な真実の集まりであり、自己で体系を形成するには不完全であり、完璧な体系の基礎として十分に重要な役割を果たすほど根幹的でもありません。しかし、学術論理が正当な主張を持たない属性を与えられたからといって、それゆえに注目価値がないとは必ずしも言えないでしょう。言語の成長と共に発展したシステムであり、哲学における最も重要な問題や有名な証明にその影響を与えた体系は、完全には無視されるべきではないでしょう。また、記憶と使用頻度も知性のプロセスにおいて大きな役割を果たしており、古代論理のいくつかの原理が文化ある心に深く組み込まれていることも認めなければなりません。しかし、特定の記憶形式が推移や三段論法の特別な規則を示す方法で持つ実質的な有用性は問われ、そのような技能が与えられるよりも、自己の資源に頼る未練習の心がより多くの精神力を養うために有利であるかどうかは、まだ検討する価値があるかもしれません。この章で導き出された特定の結果については、それらは普通または学術的論理の性質と、より完璧な推論理論との関係に関する調査を助けるためだけに設計されていることに注意すべきです。

\newpage

\chapter{確率論の基礎理論（On the Theory of Probabilities）}
\label{ch:xvi}

\begin{center}
\large\textit{論理学から確率論への自然な架橋・事象の代数と確率測度}
\end{center}
\vspace{4mm}

\#\# 機率論について

1. また一年が経つと、 Pascal が最初の確率に関する問題を解決し、その解答によって確率論という科学の基礎を築いたからには、もう二百年になるだろう。この科学は、より抽象的な数学的考察よりも共通して魅力を持つものではないが、その論理や語調、数式に

1. また一年が経つと、 Pascal が最初の確率に関する問題を解決し、その解答によって確率論という科学の基礎を築いたからには、もう二百年になるだろう。この科学は、より抽象的な数学的考察よりも共通して魅力を持つものではないが、その解明に於いても、それに関連する新しい解析方法の開発にも、そして実用的な生活における重要な役割を果たすためにも、その最初の問題である Chevalier de Méré が Port Royal の隠者に提案した問題は、長い系列の一環として存在し続けている。この主題に対する興味は、単なる数学的関連性や有用性の関連だけでなく、独立した対象としての確率論への考察、その基本的な概念方法、当時の力学と同様に徐々に付随して加わった重要な補助原理、そして実際に達成された完成度の評価にも集中する。ここでは、その完成度が一貫性と過程の調和にある場合について述べる予定である。

2. 有名な著者%
\footnote{Poisson, Recherches sur la Probabilités des Jugemens.}は、科学の基本的な定義を次のように述べている：

``事件の確率とは、その事象が起こったか、または起こる可能性があるという信頼度である。''

「ある事象の確率の尺度は、その事象に有利な状況の数と、有利または不利で全て等可能性のある総合的な状況の数との比である」（全て同程度の可能性を持つ同等に起こりうる状況における有利または不利の状況の数との比）。

これらの定義から、数学的確率という用語は、事象が発生するか否かに関連する状況に対する私たちの知識の状態を参照していることが分かります。私たちはどの程度の情報を持っているかによって、その事象が起こる可能性があると推測する理由が変化し、単純に言えば、その事象に対する「期待」も変化します。確率は部分的な知識に基づく期待です。ある事象が発生する影響を全て理解した完全な知識があれば、期待は確定性へと変わり、確率の理論が必要なくなります。

3. ある事象に対する期待は、その発生に有利な既知のケースの数が全体の等可能性のあるケースの数に対して増加するにつれて強まりますが、それが心の感情としての期待の強度を何らかの数値基準に参照できるものであると断言することは哲学的に適切ではありません。気性の強い人は希望を高く持ち、胆怯な人たちは絶望し、意志が弱い人は疑問に陥ります。科学的な調査対象としては、意見と感覚にはある類似性があります。温度計や精密に準備された写真板は熱や光の感覚の強度を示さない代わりに、それらの感覚が生じる際に伴う特定の物理的状況を示します。同様に、確率論は期待に基づく状況の数値的な尺度を対象とし、この目的は確率論の正当な応用全体を含みます。生命保険や他の確率論の統計的应用において使用する規則は、期待の心的現象とは全く独立しています。これらの規則は過去が未来に似ているという前提に基づいています。同じ状況下では、特定の数値的な頻度で同じ事象が再発生することが傾向にあると仮定します。これは人間の希望や恐怖の強さを計算する試みではなく、未来が過去に類似することを前提としています。

現在の経験は、特定の種類の出来事が、特定の状況下において、一定の頻度で再発生する傾向があることを証明しており、その真の原因が私たちに知られているか否かに関わらず、これは実際には確認されている。もちろん、この傾向は一般的には、観察範囲が十分大きいときのみに現れる。

大規模な国の司法記録、出生と死亡の登録簿、年齢や性別との関係において、年々驚くべき一貫性を示す。特定の言語や言語家族においても、同じ音、音列、そして書かれた言語の場合には、同じ文字と文字列が一定の頻度で再発生する。

アイルランド各地に見つかった原始的なオーガムの彫刻において、最初は単語の区別がつかなかったものも、この原理を厳密に適用することで鍵を見出し、その解読が可能となった。この発見はダブリン大学の名誉博士チャールズ・グレーヴス教授によるもので、彼は欧州すべての言語に対してシーケンス表を作成してこの原理を検証した。

同じ方法は学術的な東洋学者エドワード・ヒンクス博士によって、 Layard氏の企画により最近ネバフの遺跡から発見された楔形文字記録の解読に適用されていると理解される。

以上の声明と定義から、確率論の正当な対象を概念化することを試みる。

確率とは、特定の知識に基づく期待のことである。この知識は、事件が発生する頻度に関する相対的な頻度を知ることにある。したがって、事象の確率、または組み合わせの確率は、その発生と失敗の過去の系列観察から得られるものであっても、特定の事物の構成に関する知識から導き出されるものであっても、すべての場合において、当該データとなる。いくつかの関連事象や、それらとの組み合わせの確率は、その他の事象と関連している別の事象の確率から求める対象（\textit{qu\ae{}situm}）となる。

5. 事象は単純なものか複合的なものとに区別される。後者の用語は、単純な事象の組み合わせで構成される事象に適用される（I.13）。このように、理論の目的は、他の事象と関連している与えられた事象の確率から、その他の事象が単純なものか複合的なものであっても、ある事象の確率を決定することである。

例えば、サイコロの構造から、特定の投げで1（アース）が出る確率は$\displaystyle\frac{1}{6}$であることが分かっているとする。そして、2回続けて投げるとき、その中で1がちょうど1つ出る確率を求めたいとする。この問題をデータと対象（\textit{qu\ae{}situm}）の順序で述べると次のようになる：

初めのデータ：最初の投げでアースが出る確率 $\displaystyle = \frac{1}{6}$

第二の前提は、二度目の投げがジャックとなる確率が $\frac{1}{6}$ である。

問題は、次の2つの事象の確率を求めよ：
1. 第一次投げがジャックで、第二次投げがジャックでない。
2. 第一次投げがジャックでない、第二次投げがジャックである。

ここでは、二つのデータは最初の投げでジャック（エース）が出る確率と、二番目の投げでジャックが出る確率として定義される単純な事象の確率である。問題の核心（qu\ae{}situm）は複合事象の確率であり、それはデータに含まれるまたは暗示された単純な事象の特定の排他的組み合わせ又は結合である。

おそらく一般的には、問題の数値的な条件が存在する事物の構造から導き出される場合、データは単純な事象の確率であり、核心（qu\ae{}situm）はそのような単純な事象に依存した複合事象の確率となるでしょう。このような問題のクラスは、確率理論を学ぶ人々が過度に注目した可能性があります：賭博や技術的なゲームなどです。前者では、ダイスの構造などの物理的状況が各可能なゲームステップの確率を決定し、その結果はこれらのステップの特定の組み合わせである；後者では、プレイヤーの相対的な技巧度が既知であり、それが同じ要素を決定します。しかし、統計的な問題の場合、我々の知識の要素は事物の構造からではなく、自然や人間社会の記録された観察から抽出される場合、このような観察によって提供されるデータが単純な事象の確率である理由はありません。むしろ、特定の組み合わせでの出来事または状況の発生（因果関係のある秘密的なつながりを示唆）は、その頻度を数値的に記録して将来的な指導に有用であることを我々に示唆します。

このような観察によって提供されるデータは複合事象の確率です。このようなデータが含まれる問題に対する一般的方法による解法は、確率論の完全な発展にとって重要であり、またその理論を実際的に重要な多くの問いへの適用にとって必要条件であるかもしれません。

6. 以前の問題のような問題に対する既知の方法がどの程度適用可能かを評価する前に、確率論におけるすべての問題は別の形で見ることができるという事実に注意を向けた方が良いでしょう。この形態は、事件ではなく、そのイベントが発生したことを主張する命題や発生する可能性があることを主張する命題を用いることで構成されます。そして、数値的な確率の要素はこれらの命題の真偽に関連付けられ、イベントの発生とは関係ありません。

すなわち、イベント $E$ の発生確率を表す数値分数 $p$ を考慮する代わりに、そのイベントが発生すると主張する命題 $X$ の真偽の確率として見ることにします。同様に、任意の確率 $q$ は、イベント $E$ と $F$ の共存を主張する命題の真偽の確率として解釈されます。また、イベントの排他的かつ仮定的な組合せから排他的かつ条件付きの命題への変換も同様に行います。

新しい確率に対するこの適用は古い適用の必然的伴走であるが、その採用には既に論理的な命題について議論し、主な種類を定義し、互いの依存関係の量と性質を決定するための方法を確立したという実用上の利点がある。この基礎や等価の基礎に基づいてどのような一般的な確率理論も構築されるべきである。

私はここで言及しないが、適用理論における特定の場合には他の考慮事項が必要となるかもしれません。データが明確な解決に十分でない場合、これは仮定によって補完される必要があるかもしれません。また、大量の数値を含む声明が含まれる場合、解決後も特別な処理方法が必要となる可能性があり、そのような問題から生じる困難に対処するための特殊な方法が必要となるかもしれません。

しかし、どの場合でも、上記（4章）で述べた一般問題の何らかの形式が含まなければならないし、その問題についての議論において理論的な適切かつ特有の作業が行われるべきである。私はこの目的のために次の各章での調査が主に devoted されていることを強調したい。

補助仮定を含む問題には、例外的に言及するつもりはありませんが、このような問題（I.17）に対する一般的な方法論を確立することが最初の重要性があるからです。

7. 以下は、ラプラスから主に引用した確率問題の解決に適用された原理の要約である。これらは既に述べた基本的な定義の結果であり、これらの定義がいかに利用可能となったかを示す指標として見なされる。

1. 原則第1条。ある事象が起こる確率を \( p \) とするとき、その事象が起こらない確率は \( 1-p \) となる。

2. 第2条。二つの独立な事象の共現確率は、それぞれの事象の確率の積に等しい。

3. 第3条。二つの依存する事象の共現確率は、そのうちの一つの事象が起こる確率と、その事象が起こった場合に他の事象も起こる確率の積に等しい。

4. 第4条。ある事象 \( E \) が発生したとき、別の事象 \( F \) も同時に発生する確率は、事象 \( E \) と \( F \) の共現確率を、事象 \( E \) の発生確率で割ったものに等しい。

5. 第5条。二つの共存不能な事象のいずれかが起こる確率は、それぞれの事象の個別の確率の和に等しい。

6. 第6条。観察された事象が \( n \) 個の異なる原因のうちの一つからしか生じないことが先ず知られている場合、その中のいずれかの原因が起こる確率は、その原因の存在下での事象の確率を分子とし、すべての原因に対する類似した確率の和を分母とする分数である。

7. 第7条。将来の事象の確率は、各原因の確率をそれぞれの原因が存在する場合にその将来の事象が発生する確率で掛けて得られる積の総和に等しい。

8. これらの原理の範囲と相対的な有用性について、以下にいくつかの観察を述べる。

1. これらの原理を適切に組み合わせることにより、独立した単純な事象のそれぞれが与えられた明確な確率を持つ独立した事象のどのような方法であっても依存する複合事象の確率を表現することが常に可能である。実際に解決された問題の大半はこのような種類であり、それらの問題に対する解法に伴う困難性は、確率論による理論によって提供される示唆が不十分なものではないことを指摘すべきである。その代わりに、大量の項や複雑な項を含む関数や級数が導入されることにより、これらの示唆を利用するために必要な分析が必要であったからである。したがって、独立した単純な事象の確率がデータとなるすべての問題は、受け入れられた方法の範囲内にあると全面的に認めることができる。

2. 以上の原理のうちいくつかは、特に第六と第七の原理は、すべてのデータが単純な事象の確率である必要性を前提としていない。これらの原理が因果関係に関する問題に特有の適用においては、それらが考慮する原因が互いに排他的であることを仮定しており、その結果、効果の生起に影響を与えるためのこれらの原因の組み合わせは不可能であるとされている。以前説明したように、これらの数値的な確率を関連する事象から表現する命題に移すと、これらの原理が適用可能な最も一般的な問題は以下のデータと求めるもの（qu{\ae}sita）の順序で述べることができる。

\begin{center}\textsc{データ}.\end{center}

1. $n$ つの条件付き命題の確率：

\[
\begin{array}{cccccc}
\text{「原因」 $A_1$ が存在するとき、事象 $E$ が生じる；}\\
      \text{「原因」 $A_2$ が存在するとき、事象 $E$ が生じる；}\\
\multicolumn{5}{c}{\dotfill}\\
      \text{「原因」 $A_n$ が存在するとき、事象 $E$ が生じる；}
\end{array}
\]

2. これらの命題の前件が互いに排他的であるという条件。

\begin{center}\textsc{要求}.\end{center}

事象 $E$ の発生という命題が真である確率；また、その命題が既知の真であるとき、それぞれの原因 $A_1, A_2 \dotsc A_n$ の発生という命題の真偽確率。

ここで、与えられたデータは一連の複合事象の確率であり、これらは条件付き命題によって表されます。しかし、このシステムは明らかに非常に限定的で特定的なものです。なぜなら、命題の前提は互いに排他的であるという条件を満たしており、全体として唯一の後果が事象 $E$ であるからです。このようなシステムについて議論する能力があるからといって、任意の一連の条件付き命題の確率に関する問題を解決できるとは限らないことは明らかであり、ましてや任意の命題の体系の確率に関する問題を解決できるとは限りません。

主題を実質的な観点から見ると、排他的因果関係という仮定はしばしば現実の世界で達成されるものではないことが明白です。現象は通常、複雑な原因の協働によって生じるものであり、その量と特性は不明であることが多いようです。このような状況は、自然や社会に関する多くの分野において確率論が有用な応用を提供する可能性がある場合でも、ほぼすべての場合に当てはまります。

9. これらの原理には、オックスフォード大学の天文学教授によって以下のように述べられたもう一つの原理を加えることができます。
\footnote{確率理論に関する特定の問題について；W.F. ドンキン、M.A., F.R.S., 等による Philosophical Magazine, 1851年5月。}

8の原理：互いに排他的な仮説 $h_1, h_2, h_3, \dotsc$ が存在し、特定の情報状態に対するこれらの仮説の確率は $p_1, p_2, p_3, \dotsc$ であり、新しい情報がいくつかの仮説の確率を変更するものとします（特に $h_{m+1}$ およびそれ以降の仮説について）で、それらに対する他の仮説には一切関係ないとして；このとき、これらの後者の仮説の確率は新しい情報を受け入れた後に互いに同じ比を保ち、すなわち
\[
  p'_1 : p'_2 : p'_3 \ldots : p'_m = p_1 : p_2 : p_3 \ldots : p_m,
\]
ここで、重ね線で示した文字は新しい情報を受け入れた後の値を表します。

この原理は、これまで述べた多くのものよりも根本的な性格を有するものであり、おそらく教授ドンキンが提案したように、公理的とみなされるべきものであるかもしれません。実際、確率の測度の定義そのものに基礎を置いているかのように思われます。「有利な結果の数に対する全体的に有利または不利とみなされる等可能性の状況の総数との比」がその定義です。この定義を採用すると、等可能性の状況の数が減少する割合において、有利な状況の数が変化しない限り、これらの状況に関連するどのイベントの確率も同じ割合で増加します。新しい仮説として、可能状況の数が減少し、その中で関心のあるイベントに有利な状況の数が影響を受けない場合、これらのイベントの確率は一定の比率で増加します。したがって、これらの相対的な確率の尺度は変化しません。

もし我々が考慮している原理が、確率そのものの定義そのものに不可分に含まれているとすれば、この原理自体は、定義を注意深く研究することによって得られる解決策よりもさらに進むことは稀であると思われます。これらの観察から、他の種類の助けなしには、確率論の理論や実践的な問題の解法における重要な進歩が期待できないという疑問が生じるかもしれません。そして、上記のような原理に基づき、与えられた命題の数やその性質に関係なく、問題を解決するための一貫した方法を確立することは不可能であると思われます。

このような目標の達成には他の要素が必要であり、それらの考慮事項は次の章で扱われる予定です。

\newpage

\chapter{確率論における一般解法の証明（General Method in Probabilities）}
\label{ch:xvii}

\begin{center}
\large\textit{任意条件付き確率問題の代数方程式体系への帰着定理}
\end{center}
\vspace{4mm}

\#\# 総則に関する確率論の一般方法の証明

1. 事件の確率は、その事件に有利な場合の数と、有利または不利で全て等可能性のある総合的な場合の数との比である（XVI.2）。以下の調査では、「確率」という用語は上記の「確率の尺度」の意味で使用される。

以上の定義から、以下のような結論を導くことができる。

I. 事件が確定的に起こるとき、その事件の数学的な確率は1である。有利な場合と可能な場合がこの場合同じであるためである。

したがって、ある事件 \( x \) の発生確率を \( p \) とするとき、\( 1 - p \) はその事件が発生しない確率である。この結果を直接定義から導くために、有利な場合の数を \( m \)、可能な場合の数を \( n \) とすると、不利な場合の数は \( n - m \) となる。したがって、定義により、

\[
\frac{m}{n} = x の発生確率
\]
\[
\frac{n - m}{n} = x の発生しない確率
\]

であり、

\[
\frac{n - m}{n} = 1 - \frac{m}{n} = 1 - p
\]

である。

II. 任意の2つの事件が共に起こる確率は、そのいずれかの事件の確率と、そのイベントが発生した場合に他のイベントも発生する確率との積である。

$m$ を最初の事象が起こる有利な状況の数とし、$n$ を同様の可能性を持つ最初の事象が起こらない不利な状況の数とする。このとき、最初の事象が起こる確率は定義により $\dfrac{m}{m + n}$ である。$m$ 個の有利な状況の中で $l$ 個が最初と第二の事象の両方が起こる状況であるとする。このとき、最初の事象が起こった場合に第二の事象も起こる確率は定義により $\dfrac{l}{m}$ であり、これらの分数を乗じると、

\[
\dfrac{m}{m + n} \times \dfrac{l}{m} = \dfrac{l}{m + n}
\]

を得る。この結果の分数 $\dfrac{l}{m + n}$ の分子は両方の事象が同時に起こる有利な状況の数であり、分母は可能な全体の状況の数 $m + n$ であるから、これは両方の事象が同時に起こる確率を表している。

したがって、任意の事象 $x$ の確率を $p$、そして $x$ が起こった場合に $y$ も起こる確率を $q$ とするとき、両方の事象の連続性 $xy$ の確率は $pq$ である。

III. 最初の事象 $x$ が起こった場合に第二の事象 $y$ が起こる確率は、その共通な発生確率の分子と最初の事象 $x$ が起こる確率の分母からなる分数を表す。

これは第2原理の直接的な結果である。

IV. 排他的な一連の事象の少なくとも一つが起こる確率は、その個々の確率の合計と等しい。

$n$ を可能な状況の数とする。また、$m_1$ を最初の事象が起こるための状況の数、$m_2$ を次の事象が起こるための状況の数とし、その他の同様な場合も同様に定義する。すると、各事象の個別の確率は $\frac{m_1}{n}$, $\frac{m_2}{n}$ などとなる。

また、これらの事象が排他的であるとする。すなわち、一つの事象が起こるための状況は他の事象が起こるための状況にはならない。したがって、ある一連の事象のうち少なくとも一つが起こるための状況の数は $m_1 + m_2 \dotsc$ であり、そのような一連の事象のいずれかが起こる確率は $\frac{m_1 + m_2 \dotsc}{n}$ となる。これは先に述べた分数の和である $\frac{m_1}{n}$, $\frac{m_2}{n}$ などであるから、この原理は明らかである。

---
2. 定義：二つの事象が独立であるとは、そのいずれかの事象が起こる確率が他方の事象の発生または不発によって影響を受けないことをいう。

この定義と原理 II を組み合わせることで、以下の結論を得る：

5. 二つの独立な事象の共存確率は、それぞれの事象の個別の確率の積に等しい。

すなわち、ある事象 $x$ の確率を $p$, その他の事象 $y$（$x$ とは独立）の確率を $q$ とする。このとき、$q$ もまた $x$ が起こった場合に $y$ が起こる確率である。したがって、原理 II により、$pq$ は事象 $x$ と $y$ の共存確率となる。

同じ状況下において、$x$ が起こり $y$ が起こらない確率は $p(1-q)$ である。これは $x$ が起こる確率 $p$ と $y$ が起こらない確率 $1-q$ の積である。同様に、両方の事象が起こらない確率は $(1-p)(1-q)$ となる。

3. 上記とは異なる種類の原理が存在し、この章の後の調査のために必要である。それについて明示的な説明に入る前に、いくつかの前駆的な観察を述べたい。

まず最初に、単純な事象と複合的な事象との区別は、事象そのものの性質ではなく、それらが心にどのように提示されるかによります。例えば、「健康である」、「成功する」といった用語にはいくつの個別の要素が含まれているでしょうか？それぞれを「単純な事象」と見做すこともできるでしょう。

言語の伝統的な使用法は、特定の組み合わせの事象に单一且定的名称を与え、他の無数の組み合わせは異なる用語や短い表現で表すことを許可してきました。これは根本的に任意のものです。したがって、単一動詞または文法学者が「単純な句」と呼ぶものによって表される事象を「単純な事象」と呼ぶとき、それらが実際にはどの程度単純であるかを意味するわけではありません。その用語は主に文法的な表現に対して使用されます。

4. したがって、事象を単純なものと複合的なものとに分類する区別が、事象の実質的な性質に基づいていない場合、それはそれらの相互依存性や独立性に関する問題には影響を与えないはずです。もし私が二つの単純な事象についての知識が特定の事実に限定されているとします。すなわち、一つの事象が発生する確率が \( p \) であり、もう一つの事象が発生する確率が \( q \) であるという事実だけである場合、私はそれらを独立と見做し、その両方が同時に発生する確率は \( pq \) であると主張します。この主張の根拠は、これらの事象が単純なものであるからではなく、与えられたデータがそれらの間に関連性や依存性に関する何らの情報も提供していないからです。

事象の確率が与えられても、その依存性についての情報が一切提供されない場合、心はそれらを独立と見做します。そして、この思考方法は、実際的な表現によれば単純な事象か複合的な事象かに関わらず、同じです。すなわち、それぞれが単一動詞で表されるか複数の動詞で表されるかに関係なく、同様に正確です。

5. 然而，假定我们不仅拥有某些事件的確率信息，而且还拥有一些关于它们可能组合的特定信息。例如，假設我们知道某些组合是不可能发生的，因此仅剩下的组合才是可能的。在这种情况下，同样的基本原則仍然适用。我们在計算任何可設想的结果確率时利用这种信息的方式，并不取决于所给確率及其可能关联的具体性質。这些事件可以是简单的也可以是复杂的；我们获取它们的確率和关联关系的知識来源或方法无关紧要。我们必须将这些事件视为除了我们已知的联系之外，彼此独立的事件。这并不重要的是这样的信息是否明确地包含在数据中，还是通过逻辑推理推导出来的。这种観点引导我们陈述这一基本原則：

VI. 既定確率的事件应被视为仅与数据中明确表达或必然蕴含的关系相互独立；而我们在利用这些关系知識的方式則与其来源无关。

上の原理の実用的重要性は、その性質いかんにかかわらず、与えられた確率の事象に対してどのようなものであれ、求められる確率の事象についてどのようなものであれ、常に論理の計算を適用することで、後者の事象が前者の事象の明確な組み合わせとして表現されること、そしてその関連性全体を明示的に特定することができるという点にあります。言い換えれば、必要な組み合わせを見つけることができ、それが可能である組み合わせのみが存在することも特定できます。したがって、上記の原理から、これらの事象について論理的な推論を行うことができるようになり、そのような事象は直接経験によって与えられた可能性のある組み合わせ条件を持つ単純な事象として扱えることになりますし、その中から特定の組み合わせの確率を求めることもできます。確率における一般方法の可能性は、この簡素化に依存しています。

6. これから入る調査は論理の計算を用いることで基礎づけられていますので、この適用から派生したいくつかの用語と表現方法について説明する必要があります。

$x$という事件とは、その発生を主張する命題が符号的に式で表現される
等式
\[
x = 1
\]
によって表される事件を指します。$\phi(x, y, z, \dotsc)$という事件とは、その発生が式
\[
\phi(x, y, z, \dotsc) = 1
\]
によって表現される事件を指します。このような事件は、その概念に含まれる単純な事件 $x$, $y$, $z$ に関連して複合的な事件と呼ばれることがあります。例えば、「雨が降る」というイベントを表す $x = 1$，「雷が鳴る」というイベントを表す $y = 1$ のそれぞれの発生が論理式で表現される場合、
\[
x(1 - y) + y (1 - x)
\]
は「雨が降ったり雷が鳴ったりするが両方ともに起こらない」という命題を表すイベントまたは状態を表します。その状態の表現は
\[
x(1 - y) + y(1 - x) = 1
\]
です。

上記の意味で関数 $\phi(x, y, z, \dotsc)$ を $V$ と表す場合、（VI.13）から明らかに二重性の法則
\[
V(1 - V) = 0,
\]
が恒等的に満たされることがわかります。

単純な事件 $x$, $y$, $z$ は、「自由度」がないと定義され、これはこれらのイベントが任意の組み合わせで発生する自由が制限されることを意味します。言い換えれば、それらに依存する複合的なイベントが発生することを排除される場合です。例えば、「雨が降る」という命題と「雷が鳴る」という命題は、「雷が鳴るが雨は降らない」という命題によって排除されると、「雨が降る」や「雷が鳴る」のイベントは制限されます。これは可能な範囲が想像される組み合わせの範囲よりも狭くなることを示します。単純な無条件のイベントは定義により独立です。

任意の複合的な事件も同様に、その発生が特定の条件下でのみ可能であると仮定された場合に「制限」されると言います。つまり、他のイベントとの組み合わせによってその条件を満たす状況下でしか発生しない場合です。

7. 我々は自然な思考順序から単純かつ無条件の事件へと複合的かつ条件付きの事件へと進めていきます。
\newline
\newline
\centering \textsc{命題 I.}

1. もし$x$, $y$, $z$という任意の無条件単純な事象に対応するそれぞれの確率が$p$, $q$, $r$であるとき、複合事象$V$の確率は$[V]$とし、この関数$[V]$は関数$V$において$x$, $y$, $z$を$p$, $q$, $r$などに置き換えたものから形成されます。

2. 同じ状況下で、事象$V$が発生した場合に他の事象$V'$も発生する確率は$\frac{[VV']}{[V]}$とし、ここで$[VV']$は論理関数$V$と$V'$の積を計算し、その結果において$x$, $y$, $z$などを$p$, $q$, $r$などに置き換えたもので表されます。

まず最初に、二つの単純な事象$x$と$y$の可能性のある組み合わせについて注目します。それぞれの確率が$p$と$q$であるこれらの事象の基本的な組み合わせ（V.11）およびそれに対応する確率は以下の通りです：
\[
\begin{array}{cccccc}
 \textsc{ 事象.}& &\textsc{ 確率.}\\
 xy,         & \text{ $x$と$y$の共存,}     & pq.\\
 x(1-y),     & \text{ $x$のみの発生,}      & p(1-q). \\
 (1-x)y,     & \text{ $y$のみの発生,}      & (1-p)q. \\
 (1-x)(1-y), & \text{ $x$と$y$の共存または両方の失敗,}  &(1-p)(1-q).\\
\end{array}
\]
これらの状況では、複合事象を表す構成要素の確率はその論理表現が$x$と$y$に対するものと同じ関数であることがわかります。そしてこれは、与えられた確率を持つ任意の単純な事象の数に関わらず、それらの事象の共存または両方の失敗が複合事象の確率に含まれる場合にも明らかに適用されます。

次に、以上の成分のいずれかを組み合わせた任意の\emph{排他的}組合せについて考慮する。普通の言葉で表現される「イベント$x$がイベント$y$なしに起こるか、またはイベント$y$がイベント$x$なしに起こる」この複合事件は、符号的に $x(1-y) + y(1-x)$ の形で表され、原理 \textsc{iv} と \textsc{v} によって決定される確率は $p(1-q) + q(1-p)$ となる。これらの後者の表現は前者が $x$ と $y$ に対して持つ関数と同じである。

この表現もまた、より一般的な法則の一部であることが明らかだ。任意の二つ以上の成分で表されるイベントは互いに排他的であり、それらの唯一の組み合わせ方は、普通の言葉では「または」で表され、シンボル論理では '+' の記号によって表される排他的な組合せである。互いに排他的な一連のイベントが起こる確率は、その個々の確率の和であり、それぞれの個々の確率を '+' という記号で結びつけることで表現される。したがって、上記の法則は一般的なものであることが確認できる。

任意の無条件のイベント$V$の確率は、$V$の中でシンボル$x, y, z, \dotsc$ を $p, q, r, \dotsc$ に置き換えることで求められる。

8. また、原理 \textsc{iii} によれば、イベント$V$が起こるときイベント$V'$も一緒に起こる確率は、分母がイベント$V$の起こり確率で、分子がイベント$V$と$V'$の共に起こり確率を表す分数で表現される。

現在、事象 $V$ と $V'$ の共存によって構成される状況やイベントの表現は、論理計算の法則に従って $V$ と $V'$ の表現を乗算することで形成されます。なぜなら、両方の $V$ と $V'$ に含まれる要素のみが積に現れ、それ以外のものは現れないからです。また、先ほど示した通り、その積で表されるイベントの確率は $x, y, z$ を $p, q, r,\dotsc$ に置き換えることで決定されます。したがって、求められる分子は定義により $[VV']$ によって表されます。そして分母は $[V]$ となりますから、

\[
\text{もし $V$ が起こると $V'$ も同時に起こる確率} = \frac{[VV']}{[V]}
\]

となります。

9. 例えば、単純な事象 $x, y, z$ の確率がそれぞれ $p, q, r$ であるとき、$x$ または $y$ が起こると $y$ または $z$ も同時に起こる確率を求めることを考えます。この場合、前件と後件の論理的な表現は以下のようになります。

1. 前件：$x$ または $y$ が起こる
\[
x(1-y) + y(1-x)
\]

2. 後件：$y$ または $z$ が起こる
\[
y(1-z) + z(1-y)
\]

これらの二つの表現を論理計算の法則に従って乗算すると、前件と後件の共存を表す式は以下のようになります。
\[
xz(1-y) + y(1-x)(1-z)
\]

この結果において $x, y, z$ を $p, q, r$ に置き換え、また前件の表現も同様に変換すると、求められる確率は次のようになります。
\[
\frac{pr(1-q) + q(1-p)(1-r)}{p(1-q) + q(1-p)}
\]

この計算の特殊な機能は、上記のような状況において、後件と前件に同時に含まれる結びつきを一括して決定することです。しかし、この応用の利点はほとんどが将来性を持ち、その明確化は後の命題で示されます。

\begin{center}\textsc{命題 II.}\end{center}

10. ある特定の条件 \( V \) が満たされるとき、確率 \( x, y, z, \dotsc \) のそれぞれが \( p, q, r, \dotsc \) であることが知られている。ここで \( V \) は \( x, y, z, \dotsc \) の関数として表現されている。これらの事象 \( x, y, z, \dotsc \) の絶対的な確率を求めよ、すなわち、条件 \( V \) が無視された場合のそれぞれの事象の発生確率を求めよ。

求められる確率を \( p', q', r', \dotsc \) とする。つまり、単純な事象だけでなく独立した事象としての \( x, y, z, \dotsc \) の確率である。このとき、命題 \textsc{i} より、条件 \( V \) が満たされる（論理式で表すと \( V = 1 \)）ときにこれらの事象が発生する確率は
\[
  \frac{\left[ xV \right]}{\left[ V \right]},\quad
  \frac{\left[ yV \right]}{\left[ V \right]},\quad
  \frac{\left[ zV \right]}{\left[ V \right]},\text{ 等,}
\]
である。ここで \( \left[ xV \right] \) は論理計算の規則に従って \( V \) を \( x \) で掛けて得られる結果であり、その結果において \( x, y, z \) をそれぞれ \( p', q', r' \) などに置き換えたものである。しかし、これらの条件付き確率は仮定により \( p, q, r, \dotsc \) と等しい。したがって、
\[
  \frac{\left[ xV \right]}{\left[ V \right]} = p,\quad
  \frac{\left[ yV \right]}{\left[ V \right]} = q,\quad
  \frac{\left[ zV \right]}{\left[ V \right]} = r, \text{ 等,}
\]
という方程式の系が得られ、これらの方程式の数と等しい \( p', q', r', \dotsc \) の値を求めることができる。

今、\( xV \) は \( V \) 内の \( x \) の因子となる構成要素の合計である。これを \( V_x \) と表すことにし、同様に \( yV \) を \( Vy \) などと表す。すると、我々の方程式は
\[
  \frac{\left[ V_x \right]}{\left[ V \right]} = p,\quad
  \frac{\left[ V_y \right]}{\left[ V \right]} = q, \text{ 等,}
\]
という形になる。ここで \( \left[ V_x \right] \) は \( V_x \) 内の符号 \( x, y, z, \dotsc \) をそれぞれ \( p', q', r' \) などに置き換えた結果を表すものである。

一般問題の意味とその解決原理をより明確にするために、以下の例を用いてみましょう。

urnから球を取り出すという試験において、取り出した球が特定の色である「白」か、または特定の成分である「大理石」か、あるいは両方の特徴を持つ場合のみに注目し、それ以外の球が取り出されたかどうかを記録していないと仮定しましょう。そして、その仮定が満たされる度合いにおいて、白い球が取り出される確率が $p$ であり、大理石の球が取り出される確率が $q$ であることが見つかったとします。これらのデータのみから、上記の条件に関係なく次の試験で白い球が取り出される確率を求め、また大理石の球が取り出される確率を求めることを要求されます。

ここで、$x$ を白い球の取り出し、$y$ を大理石の球の取り出しとします。条件 $V$ は論理関数
\[
xy + x(1-y) + (1-x)y
\]
で表されます。したがって，
\[
V_x = xy + x(1-y) = x, \quad V_y = xy + (1-x)y = y;
\]
より
\[
[V_x] = p; \quad [V_y] = q;
\]
となり、問題の最終方程式は
\[
\frac{p'}{p'q' + p'(1-q') + q'(1-p')} = p, \quad
\frac{q'}{p'q' + p'(1-q') + q'(1-p')} = q;
\]
となります。これらから
\[
p' = \frac{p + q - 1}{q}, \quad
q' = \frac{p + q - 1}{p}
\]
が得られます。

$p'$ と $q'$ はそれぞれ $p$ と $q$ に比例することが確認され、これはドンキン教授の原理によれば当然のことです。この種の問題の解法は、その原理を直接適用することで得られるかもしれません。

可能な反論に対処するために、ここでは述べます。上記の議論は、白球と大理石球の図面が物理的な構成によって独立である必要があるとは限らないことを指摘します。その独立性の仮定は確かに解法に含まれていますが、それは球体の性質や関係、あるいは形状、色、構造などからの事前の仮定には依存していません。その仮定は、これらのすべてについて我々が完全な知識を持っていないという我々の全般的な無知に基づいています。確率は常に私たちの実際の知識の状態に参照し、その数値的な値は情報の変化とともに変動します。

[**注：以下には命題の主張と証明が続きます
命題のヘッダーは小文字で太字にして中央揃えとなり、主張は別段落で斜体で記述されます。
このようにして適切にカスタマイズした \section を使用して模倣することができます（例：\newtheoremstyle{boole}{}{}{\itshape}{}{\upshape}{.}{0pt} %{\thmname{\#1}\thmnote{\#3}\thmnumber{ \#2}} \theoremstyle{boole} \newtheorem*{Propn}{} \newcommand\Proposition[2]{\section*{Proposition \#1.}\Propn[\#2]} ）

以下はその例です：

\textbf{命題 III. 11.} 任意の確率問題において、質問（qu\ae{}situm）と与えられたデータとの論理的な関連性を決定すること；すなわち、求められる確率のある事象が、与えられた確率のある事象の論理的関数として割り当てる。

任意の複合事象 $S$, $T$ などがあり、その確率は与えられており、$S$ と $T$ は符号 $x$, $y$, $z$ などの単純な事象を表す符号で明示的に知られた関数として表現されます。同様に、求められる任意の事象 $W$ があり、$W$ も $x$, $y$, $z$ などの符号で表現された明示的な関数として与えられています。$S, T,\dotsc W$ は基本的な二重性法則を満たさなければならないため、これらは単一の論理的符号 $s, t,\dotsc w$ に置き換えることができます。したがって、

\[
s = S, \quad t = T, \quad w = W.
\]

これらの式は、定義により \( S, T, \dotsc W \) から \( s, t, \dotsc w \) の符号と \( x, y, z, \dotsc \) の符号を結びつける論理方程式の系列となる。

論理計算法の方法を使用して、上記のシステムから任意の符号 \( x, y, z, \dotsc \) を消去し、確率が与えられていない事象を表す符号を除去し、残りの符号 \( s, t, \dotsc w \) およびその他の符号 \( x, y, z, \dotsc \)（もし存在する場合）から \( w \) を開発した関数として決定できる。結果は以下の形となる：
\[
w = A + 0B + \frac{0}{0} C + \frac{1}{0} D,
\]
ここで、\( A, B, C, \) および \( D \) はすべての可能な「成分」を含み、すなわち符号 \( s, t, \dotsc \) から形成されるすべての符号を表す事象の確率が与えられたものに対する成分を含む。

上記は明らかに求めた関係の完全な表現となる。なぜなら、これは単一の符号 \( w \) を表す事象 \( W \) の関数または組み合わせとして、符号 \( s, t, \dotsc \) で同様に表される事象を完全に決定し、論理計算法の法則により、これらの事象 \( s, t, \dotsc \) の間に関連付けられる条件
\[
D - 0,
\]
を定めるからである。したがって、原理 \textsc{vi.} に基づき、\( s, t, \dotsc \) を「単純な」事象と見做し、その組み合わせ \( w \) および関連する条件 \( D \) は明確に決定される。

論理的簡約過程における一貫性を望むのであれば、ここでは一般的に行われる順序を述べる。

12. （viii. 8）によれば、原始方程式は以下の形に簡約される：
\begin{align*}
&s(1-S) + S(1-s) = 0;\\
&t(1-T)+ T(1-t) = 0;\tag{1}\\
\multispan2\dotfill\\
&w(1-W)-W(1-w) = 0;
\end{align*}
これらの方程式は、その意味を損なうことなく加算することが可能である。その後、シンボル$x$、$y$、$z$を個別にまたは同時に消去することができる。後者の場合、系の方程式を加算した後に、すべての消去すべきシンボルについて展開し、その結果の各成分の係数の積を0とする必要がある（vii. 9）。

求められる表現が$w$であるため、Prop. iii. Chap. ix.により、消去後の結果は以下の形で表すことができる：
\[
Ew + E'(1-w) = 0.
\]

$E$と$E'$は、一般系（1）において$w=1$と$w=0$を代入し、シンボル$x$、$y$、$z$などを消去することで順次決定される。したがって、$E$と$E'$をそれぞれ決定するための単純な方程式は以下のようになる：
\begin{align*}
&s(1-S)+S(1-s)+t(1-T) + T(1-t)\ldots+1-W=0;\\
&s(1-S)+S(1-s)+t(1-T) + T(1-t)+W=0.
\end{align*}
これらの方程式から、シンボル$x$、$y$、$z$などを消去し、$w$を決定するための後続の展開を行うことで、$E$と$E'$がそれぞれ決定される。

必要に応じて、Prop. i. および ii. Chap. ix. の簡略化を利用することが可能である。

13. 資料は、事象の確率を教えてくれるだけでなく、それら間に明示的な関係性を指定する場合がある。そのような関係性は論理的に表現され、その方程式や方程式群は一般系に導入されるべきである。

Proposition IV.

14. 任意の事象の確率が与えられたとき、一般的な方法により他の任意の事象の結果的または派生的な確率を決定する。

前節と同様にして、与えられた確率を持つ事象 \( S \), \( T \), などがあり、求められる確率を持つ事象 \( W \) が存在します。これらの確率は \( x \), \( y \), \( z \), などに関する既知の関数であり、以下のようにデータを表記します：
\[
\begin{aligned}
S の確率 &= p;\\
T の確率 &= q;
\end{aligned}\tag{1}
\]
そして同様に、\( p \), \( q \), などは既知の数値です。次に、事象 \( S \) を \( s \), \( T \) を \( t \)，そして \( W \) を \( w \) と表記すると、前節の結果により以下の式が得られます：
\[
w = A + 0B + \frac{0}{0} C + \frac{1}{0} D; \tag{2}
\]
ここで、\( A \), \( B \), \( C \)，および \( D \) は \( s \), \( t \)，などに関する関数です。さらに、データ (1) は以下の形に変換されます：
\[
s の確率 = p, \quad t の確率 = q, \text{ など} \tag{3}
\]

式 (2) は以下のように解釈できます：
\[
\left.
\begin{gathered}
w = A + qC\\
D=0,
\end{gathered}\quad\right\}\tag{4}
\]
ここで \( q \) は不定なクラス符号（VI.12）です。また、構成要素の性質により（V.定理 iii.）、以下の関係が成立します：
\[
A + B + C + D = 1,
\]
したがって、上記のシステムの第2式は以下のように表されます：
\[
A + B + C = 1.
\]

%-----------------------File: 069.png----------------------------
\( A + B + C \) の関数を \( V \) と表記すると、システム (4) は以下のようになります：
\begin{gather}
w = A + qC;   \tag{5}   \\
V = 1.        \tag{6}
\end{gather}

この結果について一時的に考えてみましょう。\( V \) は \( s \), \( t \)，などに関する構成要素の系列の和であり、それらが含む単純な事象は \( s \), \( t \)，などによって表されるものです。したがって、式 (6) は \( s \), \( t \)，などとその確率 \( p \), \( q \)，などが条件付きであることを示しています。式 (5) は同様に条件付きの組み合わせとして \( w \) を表現します。

現在、原理\textasciitilde{}\textsc{vi}によれば、この知識がどのように得られたかという方法は、その知識が得られたらどのような方法で使用されるべきかの方法に影響を及ぼしません。我々は経験によって事象$s$、$t$などが単純な事象として提供され、任意の組み合わせに入ることの自由があるとみなして、実際に条件$V$に曝露された場合にはそれぞれ確率$p$、$q$などを持つというように論理を進めるべきです。

次に、制約$V$が取り除かれたときのこれらの事象の対応する確率を$p', q', \dotsc,$とする。現章の定理\textasciitilde{}\textsc{ii}により、これらの量は以下の方程式のシステムによって決定されます：
\[
  \frac{[V_s]}{[V]} = p, \qquad
  \frac{[V_t]}{[V]} = q, \text{ 等} \tag{7}
\]
また、定理\textasciitilde{}\textsc{i}により、同じ条件$V$のもとでの事象$w$の確率は以下のようになります：
\[
  \text{Prob. } w = \frac{[A + cC]}{[V]}. \tag{8}
\]
ここで、$V_s$は$V$中の$s$の因子となる構成要素の和を表し、$[V_s]$は$s, t, \dotsc,$がそれぞれ$p', q', \dotsc,$に置き換わったときのその和となります。定数$c$は無限事象$q$の確率を表し、従って任意であり、0から1までの任意の値を持つことができます。

現在、次のように観察されるべきである。$p'$, $q'$, などは（7）からしか値が決定されておらず、それらを（8）に代入することで、問題 Prob. $w$ を既知の量 $p$, $q$ などの関数にするためである。したがって、明らかに $p'$, $q'$, などという文字ではなく、$s$, $t$, などと同一の定量的な意味を持つ他の文字を使用してもよいことが明白である。この特定の手順は、シンボル $s$, $t$, などが論理的から定量的に変化することを伴うに過ぎない。したがって、システム（7）と（8）は次のようになる：
\[
  \frac{V_s}{V} = p, \quad
  \frac{V_t}{V} = q, \text{ など;} \tag{9}
\]
\[
  \text{Prob. } w = \frac{A + cC}{V}. \tag{10}
\]

これらの式を使用するには、（9）から $s$, $t$, などが定量的なシンボルとして $p$, $q$ などの項で表される値を求め、それを（10）に代入すればよい。したがって、$s$, $t$, などは確率を表すため、正の真分数となることが明白である。

システム（9）はより対称的に表現できる：
\[
  \frac{V_s}{p} = \frac{V_t}{q} \dotso = V. \tag{11}
\]

または、（9）と（10）を一緒に表すために、次の対称的なシステムを使用することができる：
\[
  \frac{V_s}{p} = \frac{V_t}{q} \dotso = \frac{A + cC}{u} = V; \tag{12}
\]
ここで $u$ は Prob. $w$ を表す。

15. 次に、無限大の事象 $q$ の確率を表す定数 $c$ の解釈を考える。逆向きに解釈した論理的な方程式
\[
  w = A + qC,
\]
は、事象 $A$ が起こることか、または事象 $C$ と事象 $q$ の関連性がある場合にのみ、事象 $w$ が起こると示唆している。したがって、$q$ は事象 $C$ が起こったときに事象 $w$ が起こる条件を表す。そして、$q$ の確率は $c$ であるから、$c$ は事象 $C$ が起こったときに事象 $w$ が起こる確率を表している。

従って、原理\textasciitilde{}\textsc{ii}により，
\[
  c = \frac{\text{Cとwの一致確率}}{\text{Cの確率}}
\]
となります。

%-----------------------File: 071.png----------------------------

したがって、実際にcの値を決定するための新たな経験の性質は定まる。もし$C$における$s$, $t$, 等 を単純な事象$x$, $y$, $z$, 等 の原始的な関数として置き換えると、分母となる確率を構成する事象の表現が形成される。そして、その表現にwの原始的な表現を乗じると、分子となる確率を構成する事象の表現が形成され、\emph{この事象の頻度比を新たな観察によって決定した以前のものと比較}することでcの値を得ることができる。ここで注意すべきは、定数$c$は問題の解法において必ず現れるわけではなく、求められる確率がデータから確定的に得られない場合にのみこの任意要素が現れることである。そのような場合、最終的な論理方程式（2）または（5）はそれをどのように決定するかを教えてくれる。

もしcを決定するための新たな経験が得られなければ、cの極限値0と1を指定することで、wの確率の上限と下限を求めることができる。これらは
\[
\begin{array}{cccccc}
  \text{wの確率の最小限} &= \frac{A}{V}, \\
  \text{wの確率の最大限} &= \frac{A+C}{V}.
\end{array}
\]
これらの間には、過去の経験が絶対的に矛盾しない限り、求められる確率は必ず存在する。

もし事象$C$の表現が複数の成分から成る場合で、論理的なwの値が以下のような形の場合
\[
  w = A + \frac{0}{0}C_1 + \frac{0}{0}C_2 + \text{等}
\]
であれば、それらの総和を上記のように使用する代わりに、最終的な解を以下のような形で表すことができる：
\[
  wの確率 = \frac{A + c_1C_1 + c_2C_2 + \text{等}}{V}.
\]

ここで、$c_1 =$ 事象 $C_1$ が発生した場合に事象 $w$ が発生する確率であり、他のものも同様です。便宜上、どの形式を用いるべきか決定します。

16. 上記は提案された問題の完全な理論的解法です。また、その原始的な文书中で言及されている任意の事象が条件付けられている場合でも適用可能です。例えば、データの一つが事象 $x$ が発生した場合に事象 $y$ が発生する確率 $p$ であるという情報がある場合、事象 $x$ の発生確率が与えられていないのであれば、Prob. $x = c$（任意の定数）と仮定します。これにより Prob. $xy = cp$ となり、これらの条件をデータに導入し、以前の方法に従って使用します。また、事象 $x$ が発生した場合に事象 $y$ が発生する確率を求めたい場合、解は以下の形になります：
\[
  \text{求める確率} = \frac{\text{Prob. }xy}{\text{Prob. }x}
\]
分子と分母は、以前の一般的手法により個別に決定します。

17. この研究結果により、確率に関する問題を解決するための一般的なルールを確立することができます。

\centerline{\textsc{一般的なルール。}}

\textsc{ケースⅠ}---\emph{すべての事象が条件付けされていない場合}

与えられたまたは求められる確率を持つ事象のシンボリカル表現を作ります。
複合事象に関連するこれらの表現を新しいシンボル列 $s$, $t$ などに等しくし、これらは事象を単純なものとして表すものとみなします。

論理符号をすべて除去し、それぞれの確率 $p$, $q$, 等 が与えられた事象 $s$, $t$, 等 の論理符号だけを残し、または求められる確率 $w$ の事象の論理符号のみを除いて、$w$ を $s$, $t$, 等 の開発された関数として定めます。その形は以下のようになります：
\[
  w = A + 0B + \frac{0}{0}C + \frac{1}{0}D.
\]
%-----------------------File: 073.png----------------------------
$A + B + C = V$ とし、$V_s$ を $V$ の中で $s$ の因子を含む成分の合計、$V_t$ を $t$ の因子を含む成分の合計、そして与えられた確率を持つすべての符号について同様に定義します。

次に論理から代数へと移り、以下の方程式を作ります：
\begin{align*}
  &\frac{V_s}{p} = \frac{V_t}{q} = V,  \tag{1}  \\
  &w の確率 = \frac{A + cC}{V}. \tag{2}
\end{align*}
(1) を用いて $s$, $t$, 等 を $p$, $q$, 等 の関数として決定し、その値を (2) に代入します。結果は求められる解を表します。

または対称的な方程式系を作ります：
\[
  \frac{V_s}{p} = \frac{V_t}{q} \dotsc =  \frac{A+cC}{u}
 = \frac{V}{1},  \tag{3}
\]
ここで $u$ は求められる確率を表します。

$c$ が解に現れた場合、その解釈は以下のようになります：
\[
  c = \frac{w の確率 C w}{c の確率},
\]
この解釈はその発見のために必要な経験の性質を示しています。

\textsc{Case\textasciitilde{}II.}---\emph{事象が一部条件付けられている場合}

事象 $X$ が起こると事象 $Y$ が起こる確率 $p$ が与えられたとき、先ず以下の2つの命題に分解します：
\[
\begin{array}{cccccc}
  X の確率 &= c,   \\
  XY の確率 &= cp.
\end{array}
\]
質問が事象 $W$ が起こると事象 $Z$ が起こる確率であるとき、先述のケースでそれぞれの項を別々に決定し、その分数全体は求められる確率を表します。%-----------------------File: 074.png----------------------------

この場合、$X$, $Y$, $W$, $Z$ はどのような複合的事象でもとします。記号表現の積を表す $XY$ および $WZ$ は論理計算の規則に従って $X$ と $Y$ の符号表現の積、そして $W$ と $Z$ の符号表現の積を示します。

定数 $c$ の決定は、特定の場合には、または便宜上、任意の定数列 $c_1, c_2, \dotsc$ の決定に分解されたり置き換えられたりすることがあります。

18. 以前に述べた通り（I.12）、確率理論を構成するための二つの異なる定義、または概念方法が存在します。一つは数と直接関連している一方で、もう一つは論理と直接関連しています。現在までには、数値的定義から導かれる結果について考慮し、それが論理的な結びつきを必要とする点に到達したことを示しました。その性質の一部について見た上で、論理過程と量的代数の一般的な方法が互いに補完的に作用する方法で両方の目的を果たすことがわかりました。

逆向きの調査順序を設け、論理的な関係がより直接的に含まれる確率の定義から出発し、その結果として数値的定義がどのように導かれるか、そしてどちらの要素にも依存する同じ一般的な方法が異なる手続きの順序で最終的に確立されることを示す必要があります。

論理計算の象徴的表現と代数学の象徴的表現との間に存在する類似性は、本稿を通じてしばしば指摘された事実である。甚至可以说，これらには形式的な共通点があり、相当程度ではあるが法則面でも共通点がある。後者の観点において唯一の例外を除き、その違いは解釈によるもののみである。つまり、同じ表現は論理的解釈や量的解釈のいずれかによって受け取ることができる。前者の場合、$xy$ は $x$ と $y$ を表す事象が共存することを意味し；後者の場合、それは $x$ と $y$ を表す数または量の積を意味する。したがって、単純な事象や複合的な事象を表すすべての表現は、別の解釈システムによって純粋に量的意味を持つことができる。

ここで問題となるのは、論理的解釈と数値的解釈の同一符号式に対する移行原理が存在するかどうかである。そして、この質問に対する答えは以下の考慮事項から得られる。

19. 意外性という感情があると仮定し、その対象は事件の発生であり、この感情の強度は異なる程度を持つことができる。また、この意外性の感情は、事件が生成される状況に対する我々の知識を伴い、その知識の程度や種類に応じて変化すると仮定する。そして、意外性の感情の強度自体が、純粋な心的反応として正確な数値測定が可能であると仮定せずに、これらの条件を満たすような数値評価法の可能性について調査することは完全に正当である。これらの条件は、事件の既知の状況がより強い期待を説明するものと感じられるときにはその数値的価値が増加し、弱い期待を要求されるときには減少し、等しい程度の期待を必要とする明らかな状況では不変であることを満たす。

これらの条件は、意外性という感情が数値要素として表現される法則と、期待される事件が論理要素として表現される法則が同一であるという基本原則を仮定することで満たされます。したがって、任意の無条件のイベント $\phi(x, y, z)$ が事件 $x$, $y$, $z$ のどのような方法で組み合わさっているかに関係なく、その上記の原理に従って同じ表現 $\phi(x, y, z)$ を期待値と定義することができます。ここで $x$, $y$, $z$ は単純なイベントではなく、これらのイベントの期待であることを表します。

まず、この仮定のもとでは明らかに、互いに排他的な一連の事象が起こる確率は、それぞれの事象の個別の確率の合計と等しいことが証明できる。したがって、互いに排他的な$n$つの事象 $\phi_1(x, y, z), \phi_2(x, y, z), \dotsc, \phi_n(x, y, z)$ の交替が存在する場合、その交替の表現は以下のようになる：

\[
  \phi_1(x, y, z) + \phi_2(x, y, z) + \dotsb + \phi_n(x, y, z) = 1
\]

ここで、文字式 $x$, $y$, $z$ は論理的なものであり、三つの選択肢が組み合わさった単純な事象に関連している。そして、仮定によれば、これらの選択肢のいずれかが起こる確率の表現は

\[
  \phi_1(x, y, z) + \phi_2(x, y, z) + \dotsb + \phi_n(x, y, z)
\]

であり、ここで $x$, $y$, $z$ は上記の単純な事象の確率を表している。

この表現は、他の条件が一定である限り（$\emph{c\ae teris paribus}$）、交替が含まれる選択肢の数の増加とともに増大し、その数の減少とともに減少する。これは述べた条件に合致している。

さらに、上記の仮定による確率の測度の定義から出発して、必然的な推論または次々と提示されるステップを通じて、受容されている数値的定義へと導かれる。まず、確実性の適切な数値的測度として1（$1$）を認識する。なぜなら、どのような事象 $x$ またはその反対事象 $1-x$ も必ず起こるからである。この命題の表現は以下のようになる：

\[
  x + (1 - x) = 1
\]

そして、仮定によれば、$x + (1 - x)$ の確率測度が確実性の測度となる。しかし、その式の値はどのような $x$ の特定の値に関わらず常に $1$ である。したがって、仮定のもとでは $1$ または $1$ が確実性の測度となる。

次の場所で、互いに排他的でありながら等可能性の $n$ つの事象が存在します。これらを $t_1, t_2, \dotsc, t_n$ と表記します。

%-----------------------File: 077.png----------------------------

「これらのうち少なくとも一つは起こる」という命題は、以下の等式で表現されます：
\[
t_1 + t_2 + \dotsb + t_n = 1
\]
そして、前節の論理に従って数値的な確率へと移行する際に、この等式は形式的に依然として真であり、また確率 $t_1, t_2, \dotsc, t_n$ が等しいことから、以下の等式も成り立ちます：
\[
nt_1 = 1
\]
これより $t_1 = \frac{1}{n}$ であり、同様に $t_2 = \frac{1}{n}, t_n = \frac{1}{n}$ となります。次に、部分的な系列 $t_1, t_2, \dotsc, t_m$ のうち少なくとも一つが起こる確率を求めることを仮定します。求められる表現は以下のようになります：
\begin{align*}
t_1 + t_2 + \dotsb + t_m &= \frac{1}{n} + \frac{1}{m} + \dotsb &&(m つの項)\\
                        &= \frac{m}{n}.
\end{align*}

したがって、$n$ 個の可能な等可能性の事象の中から $m$ 個の有利な場合があるとき、特定の事象の発生確率は分数 $\frac{m}{n}$ で表されます。

どのような事象も異なる等可能性の方法で起こる場合、それは実質的に互いに排他的な選択肢の一つが起こることと同値です。したがって、任意の事象の発生確率は、その有利な場合の数を分子、総計可能な場合の数を分母とする分数で表すことができます。これは確率の厳密な数値的定義であり、この定義はより特異的な論理的定義に含まれており、我々が追跡しようとした結果を含んでいます。

20. 上記の調査から明らかに、1つには、通常の数値的な確率の定義から出発するか、または確率の数値的測度に与えられる価値の法則から出発するかに関わらず、同じ実践的な結果へと導かれることが明らかになります。2つには、これらのいずれの定義もその結論まで追求し、不可分的に含む関係を考慮に入れた場合、推論や暗示によって他方の定義へと到達します。確率論の科学的な観点から言えば、両原則が互いに影響を与え合い、依存する関係において一緒に見ることこそが重要です。

\newpage

\chapter{一般解法の初等的実例（Elementary Illustrations）}
\label{ch:xviii}

\begin{center}
\large\textit{独立試行・排反事象・複合事象の確率計算の検証}
\end{center}
\vspace{4mm}

\#\#\# 第1章：基本的な例による一般方法の説明

\#\#\#\# 小見出し：
\large 一般的な確率法を示すための基本的な例

$1$. この章では、前章で証明した一般方法を、基本的な例を通じて説明する。
選択された例は主に単純性から、得られた解の容易な検証が可能なものを中心にする。
しかし、後ほど詳しく考慮されるより高度な問題についても、一部の暗示が現れる。
これらの問題を分析するためには、データ間の必要条件を明確な方法で決定することが欠かせない。
それらが可能な経験を表し、最終的な解に対応する制限を指定するためにである。
その方法とその応用についてのより詳しい考慮は、次の章に残されている。

2. 計例1。---与えられた日において雷が鳴る確率は \( p \)、雷と雪がともに降る確率は \( q \) であるが、雷と雪の現象との関連についてはさらに何も知らないとする。要求されるのは、提案された日に雪が降る確率である。

\begin{align*}
&\text{令 $x$ を---雷が鳴るという事象を表す。}\\
&\text{令 $y$ を---雪が降るという事象を表す。}
\end{align*}
すると \( xy \) は---雷と雪がともに降るという事象を表し、問題のデータは次の通りである：
\[
\mathrm{確率}, x = p, \quad \mathrm{確率}, xy = q.
\]
ここでは複合事象 \( xy \) のみが含まれているので、規則に従って
\[
xy = u.  \tag{1}
\]
このデータは次のようになる：
\[
\mathrm{確率}, x = p, \quad \mathrm{確率}, u = q; \tag{2}
\]
そして求めるべきは \( \mathrm{確率}, y \) である。ここで (1) より
\[
y = \frac{u}{x} = ux + \frac{1}{0}u(1-x) + 0(1-u)x + \frac{0}{0}(1-u)(1-x).
\]
したがって (XVII. 17) より
\begin{align*}
V &= ux + (1-u)x + (1-u)(1-x), \\
V_x &= ux + (1-u)x = x, \quad V_u = ux;
\end{align*}
そして一般の規則の方程式は、すなわち
\[
\frac{V_x}{p} = \frac{V_u}{q} = V,
\]
が次のようになる：
\[
\mathrm{確率}, y = \frac{A + cC}{V}.
\]
これらを代入すると \( A = ux \), \( C = (1-u)(1-x) \), そして \( V \) は \( x + (1-u)(1-x) \) に簡略化されるので、
\begin{align*}
&\frac{x}{p} = \frac{ux}{q} = x + (1-u)(1-x), \tag{3}\\
&\mathrm{確率}, y = \frac{ux + c(1-u)(1-x)}{x + (1-u)(1-x)}, \tag{4}
\end{align*}
これらから \( x \) と \( u \) を消去することで、
\[
\mathrm{確率}, y = q + c(l-p). \tag{5}
\]
この結果において \( c \) は、事象 \( (1-u)(1-x) \) が起こった場合に事象 \( y \) が起こる未知の確率を表す。ここで
\[
(1-u)(1-x) = (1-xy)(1-x) = 1-x,
\]
実際の乗算により得られる。したがって \( c \) は、雷が鳴らない場合に雪が降る未知の確率を表す。

一般的な解（5）は、次のように解釈することができます： hailが降る確率は、thunderとhailがともに起こる確率$q$と、thunderが起こらない確率$1-p$を乗じたもの、すなわちthunderが起こらないときにもhailが降る確率$c$の積です。そして一般的な論理はこの結果を確認します。

もし$c$が数値的に決定できない場合、$0$と$1$という極限値に$c$を設定することで、$\mathrm{Prob.},y$の下限と上限を求めることができます：
\begin{align*}
  &\text{下限 $= q$.}   \\
  &\text{上限 $= q + 1-p$.}
\end{align*}

3. 問題2：二つの事象のうち少なくとも一つが起こる確率は$p$であり、両方の事象が失敗する確率は$q$です。これらのうちちょうど一つだけが起こる確率は何ですか？

$x$と$y$をそれぞれの事象として表すと、データは以下のようになります：
\[
\begin{array}{cccccc}
\text{Prob. }xy + x(1-y) + (1-x)y = p, \\
\text{Prob. }x(1-y) + (1-x)y + (1-x)(1-y) = q;
\end{array}
\]
そして、求める確率は
\[
\text{Prob. }x(1-y) + y(1-x).
\]
ここで、すべての事象が複合事象であると仮定し、以下の関係を導入します：
\[
\begin{array}{cccccc}
xy + x(1-y) + (1-x)y = s, \\
x(1-y) + (1-x)y + (1-x)(1-y) = t, \\
x(1-y) + (1-x)y = w.
\end{array}
\]
$x$と$y$を消去し、$w$を$s$と$t$の展開された関数として決定すると、
\[w = st + 0 s(1-t) + 0 (1-s)t + \frac{1}{0} (1-s)(1-t).\]
したがって、$A = st, C=0, V=st + s(1-t) + (1-s)t = s + (1-s)t,
V_s=s, V_t=t$；そして、一般の法則（XVII. 17）の方程式は
\begin{equation}
\frac{s}{p} = \frac{t}{q} = s + (1-s)t, \tag{1}
\end{equation}
\[
\text{Prob. }w = \frac{st}{s + (1-s)t};
\]
となり、$s$と$t$を消去すると
\[
\text{Prob. }w = p + q - 1.
\]
したがって、求めた確率は$p + q - 1$である。この結果は以下の方法で検証できる：$p$は与えられた事象のいずれかまたは両方が起こる確率であり、$1-p$はそれらが両方とも失敗する確率である。また、$q$はどちらも失敗する確率であり、$1-q$は両方が起こる確率である。したがって、$1-p + 1-q$, すなわち $2-p-q$は両方とも起こるかまたは両方とも失敗する確率である。しかしながら、残りの可能な選択肢は一方だけが起こることであり、その確率は$1-(2-p-q)$, つまり$p + q - 1$である。

4. 問題3：証人Aが真実を話す確率が$p$で、証人Bが真実を話す確率が$q$であり、彼らの主張が異なる確率が$r$である。もし彼らが一致した場合、その主張が真実である確率は何ですか？

$x$ を $A$ が真実を話すという仮定で表し、$y$ を $B$ が真実を話すという仮定で表すとします。すると、$A$ と $B$ の声明が異なるという仮定は $x(1-y) + y(1-x)$ で表現されます。一方、彼らの声明が一致するという仮定は $xy + (1-x)(1-y)$ で表現され、真実についての一致を表す仮定は $xy$ です。したがって、以下のデータがあります：
\[
\text{Prob. } x = p, \quad \text{Prob. } y = q, \quad
\text{Prob. } x(1-y) + y(1-x) = r,
\]
これらのデータから、次を求めます：
\[
\frac{\text{Prob. } xy}{\text{Prob. } xy + (1-x)(1-y)}.
\]
しかし、$\text{Prob. } x(1-y) + y(1-x) = r$ であることが明白なので、$\text{Prob. } xy + (1-x)(1-y)$ は $1-r$ となります。したがって、求めなければならないのは
\[
\frac{\text{Prob. } xy}{1-r}
\]
です。

現在の関心事象が複合的なイベントであることに注意すると、$x(1-y) + y(1-x)$ と $xy$ を仮定します：
\[
\left.
\begin{array}{cccccc}
  x(1-y) + y(1-x) &= s   \\
  xy              &= w
\end{array}
\right\} \tag{1}
\]
これらのデータから、$\text{Prob. } x = p$, $\text{Prob. } y = q$, $\text{Prob. } s = r$ であり、求めなければならないのは $\text{Prob. } w$ です。

式（1）は簡略化すると、次のようになります：
\begin{multline*}
  \{x(1-y) + y(1-x)\}(1-s) + s\{xy + (1-x)(1-y)\}   \\
  + xy(1-w) + w(1-xy) = 0;
\end{multline*}
これより
\begin{gather*}
w  = \frac{x(1-y)(1-s) + y(1-x)(1-s) + sxy + s(1-x)(1-y) + xy}{2xy-1}\\
   = \frac{1}{0}xys + xy(1-s) + 0x(1-y)s + \frac{1}{0}x(1-y)(1-s) \\
   + 0(1-x)ys + \frac{1}{0}(1-x)(1-y)s + \frac{1}{0}(1-x)y(1-s)\\
   + 0(1-x)(1-y)(1-s).
\end{gather*}
この展開式における係数 $\dfrac{1}{0}$ は、（X.6）の等価な形で置き換えられています。ここで
\[
V = xy(1-s) + x(1-y)s + (1-x)ys + (1-x)(1-y)(1-s);
\]
これより論理から代数へ移行すると，
\begin{gather*}
    \frac{xy(1-s) + x(1-y)s}{p} = \frac{xy(1-s) + (1-x)ys}{q}   \\
  = \frac{x(1-y)s + (1-x)ys}{r}\\
= xy(1-s) + x(1-y)s + (1-x)ys + (1-x)(1-y)(1-s).   \\
  \textrm{Prob. }w
= \frac{xy(1-s)}{xy(1-s) + x(1-y)s + (1-x)ys + (1-x)(1-y)(1-s)},
\end{gather*}
これから容易に得られるのは
\[
  \textrm{Prob. }w = \frac{p+q-r}{2};
\]
したがって
\[
  \frac{\text{Prob. }xy }{1-r} = \frac{p+q-r}{2(1-r)}  \tag{3}
\]
となります。

同様に、$A$ と $B$ が一致した場合のその主張が偽である確率を求めることを考えます。このとき式（1）の第2方程式を次のように置き換える必要があります：
\[
  (1-x)(1-y) = w;
\]
最終的な論理方程式は
\begin{gather*}
  w = \frac{1}{0}xys + 0xy(1-s) + 0x(1-y)s + \frac{1}{0}x(1-y)(1-s)\\
\hfill + 0(1-x)ys + \frac{1}{0}(1-x)y(1-s) + \frac{1}{0}(1-x)(1-y)s \\
\hfill + (1-x)(1-y)(1-s);   \tag{4}
\end{gather*}
これより先に述べた手順を踏むと、最終的に
\[
  \text{Prob. }w = \frac{2 - p - q - r}{2}.   \tag{5}
\]
したがって
\[
\frac{\text{Prob. } (1-x)(1-y)}{1-r} = \frac{2-p-q-r}{2(1-r)} \tag{6}
\]
となります。

これらの結果は相互に整合性があります。なぜなら、$A$と$B$の共有点を述べた命題が真か偽のいずれかであることは確実であり、したがって（3）と（5）の第二項は加算によって1となるべきです。
現在、恒等的に
\[
\frac{p+q-r}{2(1-r)} + \frac{2-p-q-r}{2(1-r)} = 1
\]
が成り立ちます。

結果（5）と（6）の単純さから推測すると、それらは既知の原理を適用することで容易に導き出せるかもしれません。しかし、これらは既知の方法の範囲には直接含まれていないことに注意しなければなりません。データの数が、その中に含まれる単純な事象の数を超える場合があります。
M. Cournot の優れた著作「確率論の解説」では、このような状況において原始前提から異なるデータセットを選択し、それぞれのデータセットから独立性を仮定して個々の単純な事象の確率を求め、それらの値を比較することで、独立性の仮定がどの程度正当化されているかを判断することを提案しています。この方法は、上記の基準により所述する単純な事象がほぼまたは完全に独立であることが証明される場合にのみ正確さに近づくことができるでしょう。
証拠と判断に関する問題においてこのような仮定が採用されると、実際の経験によってその正当性が確認されているかどうかは疑問視されます。

5. 例題4。---特定地域における一連の観測から、任意に選んだ家が発熱を引き受ける確率が$p$であり、コレラを引き受ける確率が$q$であり、両方の病気から逃れ、清潔さと換気に関する衛生状態が欠如している確率が$r$でした。任意に選んだ家が衛生状態が欠如している確率は何ですか？

ある家の観点から、次のようないっそうの符号を用います：
%-----------------------File: 085.png----------------------------

\begin{tabular}{rl}
  符号 $x$ & 熱病の発生。\\
             $y$ & \hfill「」\hfill 霍乱。 \\
             $z$ & 卫生状態の欠如。
\end{tabular}

確率が与えられた事象は、それぞれ符号 $x$, $y$, および $(1-x)(1-y)(1-z)$ で表され、求められる確率は $z$ である。したがって，
\[
(1-x)(1-y)(1-z) = w;
\]
と仮定する。すると、我々のデータは次のようになる：
\[
\text{Prob. } x = p,\quad \text{Prob. } y = q,\quad \text{Prob. } w = r,
\]
そして $z$ の確率を求めることになる。今，
\begin{multline}
z = \frac{(1-x)(1-y)-w}{(1-x)(1-y)} \\
= \frac{1}{0}xyw + \frac{0}{0}xy(1-w) + \frac{1}{0}x(1-y) + \frac{0}{0}x(1-y)(1-w) \\
+ \frac{1}{0}(1-x)yw + \frac{0}{0}(1-x)y(1-w) + 0(1-x)(1-y)w \\
\hfill + (1-x)(1-y)(1-w). \tag{1}
\end{multline}
上記から導かれた $V$ の値は
\[
\begin{array}{cccccc}
V = xy(1-w) + x(1-y)(1-w) + (1-x)y(1-w) \\
+ (1-x)(1-y)w + (1-x)(1-y)(1-w) = 1-w+w(1-x)(1-y);
\end{array}
\]
また、同様に $V_x, V_y, V_w$ を簡略化すると，
\[
V_x = x(1-w),\quad V_y = y(1-w), \quad V_w = w(1-x)(1-y);
\]
これらは代数方程式を形成し
\[
\frac{x(1-w)}{p} = \frac{y(1-w)}{q} = \frac{w(1-x)(1-y)}{r} = 1 - w + w(1-x)(1-y). \tag{2}
\]
上記の項 $\frac{0}{0}$ の係数を持つ項について、私は単一の項にまとめることではなく、多様性のために（第17章参照），
\[
\frac{0}{0}x(1-w) + \frac{0}{0}(1-x)y(1-w); \tag{3}
\]
を用いて表すことにする。すると $z$ の確率は

%-----------------------File: 086.png----------------------------

\[
\textrm{Prob. }z = \frac{(1-x)(1-y)(1-w) + cx(1-w) + c'(1-x)y(1w)}
                        {1 - w + w(1 - x)(1 - y)}. \tag{4}
\]

（2）と（4）から、
\[
\textrm{Prob. }z = \frac{(1-p-r)(1-q-r)}{1-r}
+ cp + c'\frac{q(1-p-r)}{1-r},
\]
が求められる確率の表現となる。この結果において $c = 0$ および $c' = 0$ とすると、その下限値は
\[
\frac{(1-p-r)(1-q-r)}{1-r}
\]
となり、$c=1$, $c'=1$ とすると上限値は
\[
1-r
\]
となる。

6. この解法を検討すると、選択された前提が非常に欠陥多きものであることが明らかです。定数 $c$ と $c'$ はこれを見せるし、最終的な論理方程式の対応する項（3）は欠陥を補う方法を示しています。つまり，
\[
\begin{array}{cccccc}
x(1-w) = x \{1 - (1-x)(1-y)(1-z)\} = x, \\
(1-x)y(1-w) = (1-x)y \{1 - (1-x)(1-y)(1-z)\} = (1-x)y,
\end{array}
\]
から、$c$ は発熱を引き起こした家が衛生状態が悪かった確率であり、$c'$ は発熱なしのコレラを引き起こした家が衛生状態が悪かった確率であることが分かります。

論理展開に影響を与えた係数 $\frac{0}{0}$ の項を集めた場合、最終的な解は以下の形になります：
\[
\text{Prob. } z = \frac{(1-p-r)(1-q-r)}{1-r} + c\left(p+q-\frac{pq}{1-r}\right)
\]
ここで $c$ は、いずれかまたは両方の病気を引き起こした家が衛生状態が悪かった確率を表します。この結果は以前のものと完全に一致しており、実際にはこのような場合に提示された異なる形式の解の必要なる等価性は形式的に証明できます。

上記の解法は特定の場合でも検証可能です。例えば、第二形で $c = 1$ の場合、Prob. $z = 1-r$ となり、これは正しい結果です。もし発熱またはコレラのいずれかの存在が衛生状態が悪かったことを絶対的に示すとすると、任意の一軒の家が衛生状態が悪い確率は単にその家が $z$ というカテゴリに入らない確率に等しくなり、これはデータによれば $1-r$ となります。

上記の解法を一般化して検証することは困難かもしれません。

上記の解法における定数 $p$, $q$, および $r$ は以下の条件を満たすべきです：
\[
p + r \stackrel{=}{<} 1, \quad q + r \stackrel{=}{<} 1.
\]

7. 問題5。仮説三段論法の前提の確率が与えられたとき、結論の確率を求めよ。

裸の形で表すと以下のようになります：

\begin{tabular}{ll}
  主大前提: & 命題 $Y$ が真であるならば、命題 $X$ も真である。\\
  小前提: & 命題 $Z$ が真であるならば、命題 $Y$ も真である。\\
  結論:    & 命題 $Z$ が真であるならば、命題 $X$ も真である。
\end{tabular}

大前提の確率を $p$, 小前提の確率を $q$ とします。

すると、以下のデータがあります。命題 $X$ を $x$ などと表し、任意定数 $c$ および $c'$ を仮定します：
\begin{align*}
  \text{確率 } y &= c, & \text{ 確率 } xy &= cp;\\
  \text{確率 } z &= c', & \text{ 確率 } yz &= c'q;
\end{align*}
これらのデータから、以下の確率を求める必要があります：
\[
  \frac{\text{確率 } xz}{\text{確率 } z} \quad \text{または} \quad \frac{\text{確率 } xz}{c'}.
\]

まず、
\[
  xy = u, \quad yz = v, \quad xz = w,
\]
と仮定します。通常の方法で $w$ を $y$, $z$, $w$, $v$ の展開関数として求めると、与えられた確率に対応する命題の符号から次の式を得ます：
\[
\begin{split}
  w &= uzvy + 0u(1-z)(1-v)y + 0(1-u)zvy   \\
    &+ \frac{0}{0}(1-u)z(1-v)(1-y) + 0(1-u)(1-z)(1-v)y   \\
    &+ 0(1-u)(1-z)(1-v)(1-y)
     + \text{係数が }\frac{1}{0} の項;
\end{split}
\]
ここで、論理から代数へ移行すると、
\begin{multline*}
  \frac{uzvy + u(1-z)(1-v)y}{cp}
= \frac{uzvy + (1-u)zvy + (1-u)z(1-v)(1-y)}{c'}   \\
= \frac{uzvy + (1-u)zvy}{c'q}   \\
= \frac{uzvy + u(1-z)(1-v)y + (1-u)zvy + (1-u)(1-z)(1-v)y}{c} = V. \\
  \text{確率 } w = \frac{uzvy + a(1-u)z(1-v)(1-y)}{V},
\end{multline*}
ここで
\[
\begin{split}
  V &= uzvy + u(1-z)(1-v)y + (1-u)zvy + (1-u)z(1-v)(1-y)   \\
    &+ (1-u)(1-z)(1-v)y + (1-u)(1-z)(1-v)(1-y),
\end{split}
\]
この方程式系の解は
\[
  \text{確率 } w = c'pq + ac'(1-q)
\]
を導き、従って
\[
  \frac{\text{確率 } xy}{c'} = pq + a(1-q),
\]
が得られます。この式における任意定数 $a$ は、命題 $Z$ が真でありながら $Y$ が偽である場合に $X$ も真となる確率を表します。つまり、小前提が偽であるときの結論が真となる確率です。

この調査は、命題 \( Z \) を真と仮定し、その後で \( X \) の確率を求めることによって大幅に簡素化される可能性がありました。その場合のデータは単純に以下のようになります：
\[
  \text{Prob. }y = q, \quad \text{Prob. }xy = pq;
\]
これより \( \text{Prob. }x = pq + a (1-q) \) を得ることができます。明らかに、そのような状況下ではこの方法を用いることが許されることがわかりますが、私は直接的かつ条件なしの方法を適用して解を導き出すことにしました。その結果は、通常の論理で検証できるものであり、計算が必要ないほど普通の方法で得られるものです。一般的な手法は、その助けが最も必要とされないケースにおいて特に煩雑に見える傾向があります。

上記の方法は、 kittygora 三段論法だけでなく、すべての推論形式にも適用可能です。任意の論理列の前提に個別に付与される確率が与えられた場合、その方法によって結論が正当に導き出された後の結果の確率を決定することが常に可能となります。結論の真偽とその正確性は異なることについて読者には再確認する必要はありません。

8. これらの適用例において特に注目に値するのは、ある重要な事実です。それは、真のとき等価である命題が、確率としてのみ考慮される場合に必ずしも等価でないということです。この原理は以下の例によって説明されます。

例6。与えられた確率 \( p \) の排他的論理和命題「命題 \( Y \) が真か、または命題 \( X \) と \( Y \) が両方偽である」の下で、条件付き命題「命題 \( X \) が真ならば \( Y \) は真である」という確率を求めよ。

$x$ と $y$ をそれぞれ命題 $X$ および $Y$ に適応するものとする。すると、次のように表される：
\[
\text{Prob. } y + (1-x) (1-y) = p,
\]
から $\frac{\text{Prob. } xy}{\text{Prob. } x}$ の値を求めることを要求される。

まず、
\[
y + (1-x) (1-y) = t \quad \tag{1}
\]
と仮定する。$y$ を消去すると、
\[
(1-x) (1-t) = 0,
\]
となり、これより
\[
x = \frac{0}{0} t + 1 - t;
\]
そして通常の方法で進めていくと、
\[
\text{Prob. } x = 1-p + cp \quad \tag{2}
\]
となる。ここで $c$ は、$Y$ が真であるか、または $X$ および $Y$ が偽である場合に $X$ が真である確率を表す。

次に $\text{Prob. } xy$ を求めることを考える。まず
\[
xy = w \quad \tag{3}
\]
と仮定する。式 (1) および (3) の $y$ を消去すると、
\[
z (1-t) = 0,
\]
となり、これより
\[
\text{Prob. } z = cp,
\]
$c$ は前と同じ意味を持つ。したがって
\[
\frac{\text{Prob. } xy}{\text{Prob. } x} = \frac{cp}{1-p+cp},
\]
となる。

上記の論理的推論から、純粋な論理学において、排他的論理和と条件付き命題は等価であることがわかる。これらはともに真であり、ともに偽である。しかしながら、上記の調査から、排他的論理和が確率 $p$ を持つとき、条件付き命題は異なる部分的には不確定な確率 $\frac{cp}{1-p+cp}$ を持つことが確認される。しかし、これらの表現は、どちらか一方が 1 または 0 の値を取るとき、他方も同じ値を取ることを示す。したがって、結果は完全に整合的であり、論理的な変換は確率理論から導き出された公式の正確性を確認する。

読者は類似の分析により容易に証明できるだろう：条件付き命題の確率が $p$ であるとき、排他的論理和の確率は $1-c+cp$ となる。ここで $c$ は命題 $X$ の真偽確率を表す任意の値である。

9. 計算問題 7：与えられた以下のデータのいずれか（またはそれらの組み合わせ）から、事象 $x$ の確率を求めよ。具体的には、次のデータの最初のもの、最初と第二のもの、または最初、第二および第三のものである：

\[
\text{Prob. } y + (1-x) (1-y) = p, \quad
xy = w, \quad
y + (1-x) (1-y) = t.
\]

1. 事件 \(x\) の発生確率、または三つの事件 \(x, y, z\) のうちのいずれも発生しない確率は \(p\) である。

2. 事件 \(y\) の発生確率、または三つの事件 \(x, y, z\) のうちのいずれも発生しない確率は \(q\) である。

3. 事件 \(z\) の発生確率、または三つの事件 \(x, y, z\) のうちのいずれも発生しない確率は \(r\) である。

\#\#\# 第一の場合の解法

ここでは、上記のデータの中で最初のみが与えられていると仮定する。
\[
\textrm{Prob. }\{x + (1 - x)yz\} = p
\]
を満たす確率 \(x\) を求める。

まず，
\[
  \textrm{Let}\quad x + (1 - x)yz = s,
\]
とおく。次に \(yz\) を単独の符号として消去すると、以下の式を得る：
\begin{equation*}
x(1 - s) = 0.
\end{equation*}

したがって
\begin{equation*}
x = \frac{0}{1-s} = \frac{0}{0}s + 0(1-s),
\end{equation*}

従って、規則に従って計算すると，
\[
\textrm{Prob. } x = cp,
\]
が得られる。ここで \(c\) は、事件 \(x\) の発生または単独の失敗により最初の二つの選択肢のうちどちらが実現するかという確率である。解法の限界は明らかに \(0\) と \(p\) である。

この解法は、未補強の良識によっても伝えられる情報以上のものは提供しないように見える。しかし、ここでの単一のデータが実際に推論できる限りのすべてを提供しているだけである。次の場合の解法では、追加されたデータにより最終的な解法がより狭い範囲に制限されることを見ることができる。

\#\#\# 第二の場合の解法

ここでは、以下の式をデータとします：
\begin{align*}
  \textrm{Prob. }\{ x + (1 - x)yz \} &= p, \\
  \textrm{Prob. }\{ y + (1 - y)xz \} &= q.
\end{align*}

これらの式を
\begin{align*}
  x + (1 - x)yz &= s, \\
  y + (1 - y)xz &= t,
\end{align*}
と書き換えてみましょう。最初の式から、(VIII.7)により次の式を得ます：
\begin{gather*}
  \{x + (1 - x)yz\} (1 - s)
+ s \{1 - x - (1 - x)yz\} = 0, \\
  \textrm{または}\quad
  \bigl( x + \bar{x}yz \bigr)\bar{s}
+ s\bar{x}\bigl( 1 - yz \bigr) = 0;
\end{gather*}
ここで、$\bar{x}$ を $1 - x$、$\bar{y}$ を $1 - y$ と表記します。次に、$1 - yz$ の値を成分で書き換えると次の式を得ます：
\begin{equation*}
\bigl( x + \bar{x}yz \bigr)\bar{s}
+ s\bar{x}\bigl( y\bar{z} + \bar{y}z + \bar{y}\bar{z} \bigr) = 0,
\end{equation*}
これはすべての項が正である方程式です。

同様に、二番目の式から次の式を得ます：
\[(y + \bar{y}xz) \bar{t} + t\bar{y} (x\bar{z} + \bar{x}z + \bar{x}\bar{z}) = 0;\]

これらの二つの式の和から $y$ と $z$ を消去します。

この和において、$y = 1, z = 1$ と置くと、結果は $\bar{s} + \bar{t}$ となります。

同様に、$y = 1, z = 0$ と置くと、結果は
\[ x\bar{s} + s\bar{x} + \bar{t}. \]

さらに、$y = 0, z = 1$ と置くと、結果は
\[ x\bar{s} + s \bar{x} + x\bar{t} + t \bar{x}. \]

最後に、$y = 0, z = 0$ と置くと、結果は
\[ x\bar{s} + s \bar{x} + tx + t\bar{x}, \text {または} x\bar{s} + s\bar{x} + t. \]

これらの4つの式を乗算する必要があります。まず、最初の2つは次の方法で乗算することができます：
\begin{gather*}
(\bar{s} + \bar{t})(x\bar{s} + s\bar{x} + x\bar{t} + t\bar{x}) \\
= x\bar{s} + x\bar{t} + (\bar{s} + \bar{t})(s\bar{x} + t\bar{x}) \quad \text{(IX.定理Ⅱ)}\\
= x\bar{s} + x\bar{t} + \bar{s}\bar{x}t + s\bar{x}\bar{t}.       \tag{2}
\end{gather*}

次に、第二と第四の式は次のようになります：
\begin{gather*}
(x\bar{s} + s\bar{x} + \bar{t})(x\bar{s} + s\bar{x} + t) \\
= x\bar{s} + s\bar{x}.           \tag{3}
\end{gather*}

最後に、(2)と(3)を乗算すると次のようになります：
\begin{gather*}
(x\bar{s} + s\bar{x})(x\bar{s} + s\bar{x}\bar{t} + x\bar{t} + t\bar{x}\bar{s}) \\
= x\bar{s} + s\bar{x}(s\bar{x}\bar{t} + x\bar{t} + t\bar{x}\bar{s})\\
= x\bar{s} + s\bar{x}\bar{t}.
\end{gather*}

したがって、最終的な方程式は次のようになります：
\[(1-s)x + s(1-t)(1-x) = 0,    \]
これを$x$について解くと次のようになります：
\begin{gather*}
x = \frac{s(1-t)}{s(1-t)-(1-s)} \\
= \frac{0}{0}st + s(1-t) + 0(1-s)t + 0(1-s)(1-t)  ,
\end{gather*}

この結果をルールに従って進めると、最終的に次の式が得られます：
\begin{equation*}
  \text{Prob. }x = p (1-q) + cpq.      \tag{4}
\end{equation*}
ここで$c$は、事象$st$が起こった場合に$x$が起こる確率を表します。そして、$st$の展開式を$s$と$t$の式で乗算することで得られるものは次のようになります：

$c =$ 確率\quad 事象$x$と$y$が同時に起こり、または事象$x$と$z$が同時に起こり、$y$が失敗する場合、または事象$y$と$z$が同時に起こり、$x$が失敗する場合、事象$x$が起こる確率。

事象$x$の確率の限界は明らかに$p(1-q)$と$p$です。

この解法は以前のものよりも明確であると言えます。なぜなら、それは任意定数によらない項を含むからです。

\begin{center}\textsc{第三の場合の解法}\end{center}

ここで与えられた情報は次の通りです：
\[
\begin{array}{cccccc}
  \text{Prob.} \{x + (1-x)yz\} = p, \\
  \text{Prob.} \{y + (1-y)xz\} = q, \\
  \text{Prob.} \{z + (1-z)xy\} = r.
\end{array}
\]

以前と同様に、$\bar{x}$を$1-x$、$\bar{y}$を$1-y$、$\bar{z}$を$1-z$と表記し、次のように仮定します：
\begin{align*}
  x + \bar{x}yz = s,  \\
  y + \bar{y}xz = t,  \\
  z + \bar{z}xy = u.
\end{align*}

（VIII.8）により簡約すると、方程式は次のようになります。
\setcounter{equation}{4}
\begin{equation}
\begin{split}
  &(x + \bar{x}yz)\bar{s}
+ s\bar{x}(y\bar{z} + \bar{y}z + \bar{y}\bar{z})
\\
  &+ (y + \bar{y}xz)\bar{t}
+ t\bar{y}(z\bar{x} + x\bar{z} + \bar{x}\bar{z})
\\
  &+ (z + \bar{z}xy)\bar{u}
+ u\bar{z}(x\bar{y} + \bar{x}y + \bar{x}\bar{y}) = 0.   %\tag{5}
\end{split}
\end{equation}

次に、上記方程式から $y$ と $z$ を直接消去するのではなく、（IX.公理、iii）に従って、その消去結果を
\begin{equation*}
Ex + E'(1-x) = 0,
\end{equation*}
とするものとします。すると、与えられた方程式において $x = 1$ と置き換えて $y$ と $z$ を消去した結果から $E$ を求めることができますし、$x = 0$ と置き換えて $y$ と $z$ を消去した結果から $E'$ を求めることができます。まず最初に、$x = 1$ と置き換えると
\[
\bar{s} + (y + \bar{y}z)\bar{t} + t\bar{y}\bar{z} + (z + y\bar{z})\bar{u} + u\bar{y}\bar{z} = 0,
\]
となります。この方程式の最初の項において、順次 $y = 1, z = 1, y = 1, z = 0$ などと置き換えて結果を乗算すると
\[
(\bar{s} + \bar{t} + \bar{u})(\bar{s} + \bar{t} + \bar{u})(\bar{s} + \bar{t} + \bar{u})(\bar{s} + t + u),
\]
という表現を得ます。これは
\[
(\bar{s} + \bar{t} + \bar{u})(\bar{s} + t + u)
\]
に等しいです。これが $E$ の表現です。この形を保持します。例（VIII.3）では既に示したように、このような表現の実際的な乗算による簡約は便利であっても必要ではありません。

再び、式(5)において \(x = 0\) と置くと、

\[
yz\bar{s} + s(y\bar{z} + \bar{y}z + \bar{y}\bar{z}) + y\bar{t} + t\bar{y} + z\bar{u} + u\bar{z} = 0;
\]

この式から、同様の消去法により \(E'\) の表現は、

\[
(\bar{s} + \bar{t} + \bar{u})(s + \bar{t} + u)(s + t + \bar{u})(s + t + u)
\]

となる。

したがって、\(y\) と \(z\) を消去して得られる最終的な結果は、

\[
(\bar{s} + \bar{t} + \bar{u})(\bar{s} + t + u)x + (\bar{s} + \bar{t} + \bar{u})(s + \bar{t} + u)(s + t + \bar{u})(s + t + u)(1 - x) = 0
\]

である。

従って、

\[
x = \frac{(\bar{s} + \bar{t} + \bar{u})(s + \bar{t} + u)(s + t + \bar{u})(s + t + u)}
         {(\bar{s} + \bar{t} + \bar{u})(s + \bar{t} + u)(s + t + \bar{u})(s + t + u) - (\bar{s} + \bar{t} + \bar{u})(\bar{s} + t + u)};
\]

または、第二項を展開すると、

\begin{equation}
\begin{split}
x = \frac{0}{0}stu + \frac{1}{0}s\bar{t}u + \frac{1}{0}st\bar{u}
+ s\bar{t}\bar{u}
\\
+ \frac{1}{0}\bar{s}tu + 0\bar{s}\bar{t}u + 0\bar{s}t\bar{u}
+ 0\bar{s}\bar{t}\bar{u}. %\tag{6}
\end{split}
\end{equation}

したがって、論理から代数へ移行すると、

\begin{equation}
\begin{split}
  \frac{stu + s\bar{t}\bar{u}}{p}
= \frac{stu + \bar{s}t\bar{u}}{q}
= \frac{stu + \bar{s}\bar{t}u}{r}
\\
= stu + s\bar{t}\bar{u} + \bar{s}\bar{t}u + \bar{s}t\bar{u}
+ \bar{s}\bar{t}\bar{u}. %\tag{7}
\end{split}
\end{equation}

%-----------------------File: 095.png----------------------------
\[
  \text{Prob. }x = \frac{s\bar{t}\bar{u} + cstu}{stu + s\bar{t}\bar{u} + \bar{s}\bar{t}u + \bar{s}t\bar{u} + \bar{s}\bar{t}\bar{u}}, \tag{8}
\]

この方程式系を簡略化するために、${\displaystyle \frac{s}{\bar{s}}}$ を ${\displaystyle s}$ に、${\displaystyle \frac{t}{\bar{t}}}$ を $t$ とし、その他の類似な変換を行い、この変換後、$\lambda$ を $stu + s + t + 1$ としておく。すると、

\[
  \text{Prob. } x = \frac{s + cstu}{\lambda} \tag{9}
\]

とあり、関係式は

\[
  \frac{stu + s}{p} = \frac{stu + t}{q} = \frac{stu + u}{r} = stu + s + t + u + 1 = \lambda \tag{10}
\]

となる。これらの方程式から、

\begin{align*}
                      stu + s &= \lambda p,           \tag{11} \\
                      stu + s &= \lambda - t - u - 1, \\
         \therefore \lambda p &= \lambda - u - t - 1. \\
                        u + t &= \lambda(1 - p) - 1.  \\
\text{同様に},\quad u + s &= \lambda(1 - q) - 1,  \\
\text{および}\quad        s + t &= \lambda(1 - r) - 1.  \\
\end{align*}

これらの方程式から、

\begin{gather*}
  s = \frac{\lambda(1 + p - q - r) - 1}{2},  \quad
  t = \frac{\lambda(1 + q - r - p) - 1}{2},  \\
  u = \frac{\lambda(1 + r - p - q) - 1}{2}.  \tag{12}
\end{gather*}

次に、(10) より、

\[
  stu = \lambda p - s.
\]

上記の $s$, $t$, $u$ の値をこの方程式に代入すると、

\begin{multline*}
  \{(1 + p - q - r)\lambda - 1\}
  \{(1 + q - p - r)\lambda - 1\}
  \{(1 + r - p - q)\lambda - 1\}\\
=4\{(p + q + r - 1)\lambda + 1\}, \tag{13}
\end{multline*}

という方程式が得られ、$\lambda$ の値を決定する。$s$, $t$, $u$ の値は (12) より与えられるので、(9) にこれらを代入することで問題の解が完成する。

%-----------------------File: 096.png----------------------------
10. 現在、読者に明確に提示するためにこの研究の一目的として挙げられた難題が生じる。上記の方程式のどの根が $\lambda$ の値として選ばれるべきかをどのように決定するかである。この困難については、本章の冒頭で一部言及し、その完全な考慮は次の章に残されており、その結果から以下のように得られる。

以上が原著英文の

問題のデータが可能な経験から導き出されるためには、量 $p$, $q$, および $r$ は以下の条件を満たさなければなりません：
\[
\begin{aligned}
1 + p - q - r & \geqq 0, \\
1 + q - p - r & \geqq 0, \\
1 + r - p - q & \geqq 0.
\end{aligned}                        \tag{14}
\]
さらに、一般解において用いる $\lambda$ の値は以下の条件を満たさなければなりません：
\[
\lambda > \frac{1}{1 + p - q - r}, \quad
\lambda > \frac{1}{1 + q - p - r}, \quad
\lambda > \frac{1}{1 + r - p - q}.      \tag{15}
\]

これらの2つの条件セットは、一般解の制限に十分です。中心方程式 (13) から $\lambda$ の値がただ一つしか得られないことが示され、その値は必要なものと一致します。

$1 + p - q - r$ を上記の $\lambda$ の係数の中で最小とするとき、$\dfrac{1}{1 + p - q - r}$ はそれらの値の中で最大となります。この値より大きい $\lambda$ の存在を示すために、$\dfrac{1}{1 + p - q - r}$ を用いて $\lambda$ の値がただ一つであることを確認します。

(13) 式を以下のように書き換えてみましょう：
\begin{gather*}
\{(1 + p - q - r)\lambda - 1\} \{(1 + q - p - r)\lambda - 1\} \{(1 + r - p - q)\lambda - 1\} \\
-4\{(p + q + r - 1)\lambda + 1\} = 0;   \tag{16}
\end{gather*}
この最初の項を $V$ と表記します。

$\lambda = \dfrac{1}{1 + p - q - r}$ を仮定すると、$V$ は以下のようになります：
\[
-4\left(\frac{p + q + r - 1}{1 + p - q - r} + 1\right)
= -4\left(\frac{2p}{1 + p - q - r}\right),
\]
これは \emph{負の値} です。

また、$\lambda = \infty$ のとき、$V$ は \emph{正の無限大} となります。

さらに，
\begin{multline*}
\frac{d^2 V}{d\lambda^2}
= (1 + p - q - r)(1 + q - p - r)\{(1 + r - p - q)\lambda - 1\} \\
+ \text{類似の正の項},
\end{multline*}
この式は $\lambda = \dfrac{1}{1 + p - q - r}$ から $\lambda = \infty$ の間で正となります。

もし、$\lambda$ を横座標として測定し、その縦座標を $V$ とする曲線を作ると、指定された限界内では、この曲線は $\lambda$ の横座標の下側から上側へと移動し、常に凹正在します。したがって、その曲線は指定された限界内で一度だけ $\lambda$ の横座標と交わります。従って、方程式 $(16)$ は一つの根に対応します。

この解は、条件 $(15)$ を満たす $(13)$ の根である $\lambda$ によって表されますし、$s, t,$ および $u$ は $(12)$ で与えられます。$c$ の意味は通常の方法で導くことができます。

以上の説明から、問題はすべての場合において多少なりとも不定的であることがわかります。

\newpage

\chapter{統計的条件と極限不等式（Of Statistical Conditions）}
\label{ch:xix}

\begin{center}
\large\textit{経験的データ間の論理的不等式・確率の上下限決定法}
\end{center}
\vspace{4mm}

\#\# 様式的条件について

\#\#\# 1. 様式的条件とは

「様式的条件」という用語は、問題の数値データが互いに一致するためには欠かせない条件を指します。つまり、これらのデータが統計的な観察によって実際に提供される可能性があるような条件です。このような条件の決定は重要な問題であり、この作品のために十分な程度まで解決することを目的として、本章で取り組もうとします。私はこれを単に独立した考察の対象とするだけでなく、既に部分的に例示された確率論の理論に対する必要不可欠な補完としても捉えています。

両者の関連性の性質は以下の通りです：

\#\#\# 2. 確率論と代数方程式

無数の例がありますし、そのうちの一例として前章で述べた練習問題7（Ex. 7）が挙げられます。確率論における問題の解が最終的には高次方程式の解に依存することもあります。このような場合、適切な根を選択するためには特定の条件が必要であり、これはデータで割り当てられた数値の関係や必要とする要素に対する適切な制限に関連しています。これらの条件を発見することは時として独自の論理によって達成されることがあります。例えば、事件 $A$ の発生確率が $p$ であり、同じ事件 $A$ と別の事件 $B$ の共存確率が $q$ の場合、明らかに
\[
p \ge q
\]
という関係が成り立ちます。

しかし、このような関係の一般的な決定には特定の方法が必要です。そして、そのような方法を次のように確立します。

実験から得られた経験によれば、どのような事象の確率も近似プロセスの結果であり、それは観察された事象が起こる回数に対するその事象が等可能性のあるすべての事象全体の中でどの程度見られるかという比の極限です。この比は観測の回数が増えるにつれてより正確に近づきます。

現在、任意のクラスに含まれる個体の数を表すためにシンボル $n$ を使用します。したがって、$x$ が男性を、$y$ が白い生物を表すと仮定すると、

\[
\begin{array}{cccccc}
nx &=& \text{男性の数。}\\
nxy &=& \text{白い男性の数。}\\
nx(1-y) &=& \text{白くない男性の数；その他の場合も同様。}
\end{array}
\]

この表記法によると、$n(1)$ は議論の対象となる宇宙全体に含まれる個体の数を表し、$\frac{n(x)}{n(1)}$ は $n(1)$ で示される存在のうち任意の個体が男性である確率を表します。観察によって $n(x)$ と $n(1)$ の総計値について知らされていない場合、その確率は個々の観測回数が増えるにつれて $\frac{n(x)}{n(1)}$ に近づく極限です。

同様にして、この章で一般的に仮定されるように $x$ が特定の種類の事象を表すとします。そのような事象の発生回数を $n(x)$、すべての種類の等可能性のある観察事象の数を $n(1)$、そして $\frac{n(x)}{n(1)}$ またはその極限が事象 $x$ の発生確率となります。

従って、量 $n(x)$, $n(y)$, $n(1)$ などに関する比について得られる結論は、イベント $x$, $y$ などの表す事象の確率についての結論に変換できることが明らかです。例えば、次のような関係式が見つかった場合、

\[
n(x) + n(y) < n(1),
\]

すなわち、イベント $x$ の発生回数とイベント $y$ の発生回数の合計が可能な発生回数 $n(1)$ よりも少ないという関係式が得られた場合、

\[
\frac{n(x)}{n(1)} + \frac{n(y)}{n(1)} < 1,
\]

または

\[
\text{Prob. } x + \text{Prob. } y < 1.
\]

となる結論を導くことができます。

一般に、このような統計的関係は、$n(1)$ を $1$ に変更し、他のシンボル $n(x)$ を Prob. $x$ に変更することで、関連するイベントの確率との間の関係式に変換されます。

3. まず、クラスやイベントの数値的な関係を決定する方法を探求します。特に、その数値的範囲の主界と次界についてです。次に、この方法を確率論における問題の解の制限に適用します。

明確に、シンボル$n$はその操作において分配的であることが明らかです。したがって、

\begin{gather*}
n\{xy + (1-x)(1-y)\} = nxy + n(1-x)(1-y)\\
nx(1-y) = nx - nxy,
\end{gather*}

と続きます。どのクラスも、異なるグループや部分に分解できるとき、その中のものの数はそれぞれの部分のもの数の合計に等しいです。さらに、論理シンボル$x$、$y$などから形成された任意の表現は、シンボルの法則に従ってどのような方法でも展開または拡大することができますし、結果の各項に対してシンボル$n$を適用しても構いません。ただし、その定数倍因子が現れた場合、それはシンボル$n$の外側に配置されるべきです。これにより結果の値は影響を受けません。表現$n(1)$が現れた場合、当然のことながらそれは宇宙に含まれる個体の数を表します。したがって、

\begin{gather*}
n(1-x)(1-y) = n(1-x-y+xy)\\
= n(1) - n(x) - n(y) + n(xy).
\end{gather*}

また、

\begin{gather*}
n\{xy + (1-x)(1-y)\} = n(l-x-y+2xy) \\
= n(1) - nx - ny + 2nxy.
\end{gather*}

最後の項において、項$2nxy$はクラス$xy$に含まれる個体の数の二倍を示します。

4. 次に、与えられた論理的表現を持つクラスの数的な限界について調査します。この調査では、以下の原理が基本的な重要性を持ちます：

1つ目。与えられたクラスのすべてのメンバーが特定の属性$x$を持っている場合、その属性$x$をもつ個体の総数は与えられたクラスに含まれる個体の数の上界となります。

2つ目。任意のクラス$y$に含まれる個体の最小限の数は、対立クラス$1-y$の主要な数的な限界から宇宙に含まれる個体の数を引いたものとなります。

これらの原理を例示するために、以下の問題に対して適用します：

\textsc{問題}。与えられた$n(1)$、$n(x)$、および$n(y)$において、$nxy$の上界と下界を求めよ。

当たるデータは、論理の宇宙に含まれる個体数、クラス \( x \) に属する個体数、およびクラス \( y \) に属する個体数です。そして、一度にクラス \( x \) とクラス \( y \) に属する個体数を含むクラスが作られるとき、その個体数の限界を決定することが求められます。

原理\textasciitilde{}\textsc{i.}により、この個体数はクラス \( x \) に含まれる個体数を超えることはなく、またクラス \( y \) に含まれる個体数を超えることもありません。それらの主要な限界は、二つの値 \( n(x) \) と \( n(y) \) のうちの最小値となります。

原理\textasciitilde{}\textsc{ii.}により、クラス \( xy \) の小限界は以下の式で与えられます：
\[
n(l) - \text{主要な限界}\{x(1-y) + y(1-x) + (1-x)(1-y)\}, \tag{1}
\]
ここで \( x(1-y) + y(1-x) - (1-x)(1-y) \) はクラス \( xy \) の補集合であり、つまりそれが論理の宇宙を形成するために必要な部分です。

また \( x(1-y) + (1-x)(1-y) = 1-y \) であるため、式（1）は以下のようになります：
\[
\begin{split}
n(l) - \text{主要な限界}\{1-y + y(1-x)\} \\
= n(l) - n(1-y) - \text{主要な限界}\{y(1-x)\}.  %\tag{2}
\end{split}
\]
\( y(1-x) \) の主要な限界は、二つの値 \( n(y) \) と \( n(1-x) \) のうちの最小値となります。もし \( n(y) \) が最小である場合、式（2）は以下のようになります：
\[
\begin{split}
n(l) - n(1-y) - n(y) \\
= n(l) - n(1) + n(y) - n(y) = 0.
\end{split}
\]

次に、\( n(1-x) \) が \( n(y) \) より小さい場合、式（2）は以下のようになります：
\[
\text{主要な限界}\{n(y)(1-x)\} = n(1-x);
\]
したがって、式（2）は以下のように変形されます：
\begin{gather*}
n(l) - n(1-y) - n(1-x) \\
= n(l) - n(1) + n(y) - n(1) + n(x) \\
= nx + ny - n(1).
\end{gather*}

したがって、クラス \( nxy \) の小限界は、\( n(y) \) が \( n(1-x) \) より小さいか大きいかによって、0 または \( n(x) + n(y) - n(1) \) のいずれかとなります。これは等価な条件として、\( n(x) \) が \( n(1-y) \) より大きいか小さいかによっても表されます。

現在、任意のクラスの数値の限界として$0$は必ず最小限界であるため、$n(xy)$の最小限界的には上の式の2番目のものを考慮すれば十分です。従って、

\begin{align*}
  n(xy) の大限界 &= n(x), n(y) の最小値 \\
  n(xy) の小限界 &= n(x) + n(y) - n(1).
\end{align*}

この$nxy$の小限界的な式は、デ・モラン教授によって初めて与えられたもので、論理学的形式として次のように適用されます：

\begin{verse}
  特定の会社内の大多数の人々がコートを持っている。\\
  同じ会社内の大多数の人々がジャケットを持っている。\\
  よって、会社内の一部の人々はコートとジャケットを所有している。
\end{verse}

次に、命題Iについて述べます。

5. 任意の符号$x, y, z, 等$ の成分によって表されるクラスの大限界と小限界を表現する方法。$n(x), n(y), n(z), 等, n(1)$ の値が与えられたとき、

まず、成分$xyz$について考えます。

明らかに、数値の大限界は$n(x), n(y), n(z)$の最小値となります。

小限界的な数値は前問と同様に導くことができますが、その問題の解からも導くことができます。したがって、

\[
  n(xyz) の小限界 = n(xy) + n(z) - n(1). \tag{1}
\]

これは$n(xyz)$は少なくとも$n(xy) + n(z) - n(1)$以上であることを意味します。また、$n(xy)$は少なくとも$n(x) + n(y) - n(1)$以上です。従って、

\[
  n(xyz) \geq n(x) + n(y) - n(1) + n(z) - n(1),
\]

または

\[
  n(xyz) \geq n(x) + n(y) + n(z) - 2n(1).
\]

したがって、

\[
  n(xyz) の小限界 = n(x) + n(y) + n(z) - 2n(1).
\]

この論理を拡張することで、以下の結論を得ます：

1. 任意の成分によって表されるクラスの大限界的な数値は、各因子に$n$を個別に付加し、結果の値の最小値を取ることで求められます。

2番目の限界は、上記のすべての値を加算し、その結果から$n(1)$の値をいくつ（少なからず）引くことで求められます。

したがって、次のように表されます：
\begin{multline*}
 \begin{aligned}
  nxy(1-z) の-major limit &= nx, ny, n(1-z) の最小値 \\
  nxy(1-z) の-minor limit &= n(x) + n(y) + n(1-z) - 2n(1)
 \end{aligned}
\\= nx + n(y) - n(z) - n(1).
\end{multline*}

一般的なシンボルの使用においては、すべての値 $n(x)$, $n(y)$, $n(1-z)$ を $n\{xy(1-z)\}$ の-major limit とみなす方が良いかもしれません。実際、これらどれよりも超えることはありません。次に述べる調査では、この表現方法を採用します。

\begin{center}\textsc{命題 II.}\end{center}

6. \emph{与えられた $n(x)$, $n(y)$, $n(z)$ などの値と $n(1)$ の条件下で、シンボル $x$, $y$, $z$ などからなる一連の成分によって表されるクラスの-major 数値限界を決定する方法。}

このような限界的数値を決定する一つの方法は、各成分の-major limit の最小和を形成することです。したがって、式
\[
 n\{xy + (1-x)(1-y)\}
\]
の-major limit は、最初の成分から得られる二つの最小値 $nx$, $ny$ を加え、第二の成分から得られる二つの最小値 $n(1-x)$, $n(1-y)$ を加えることで求められます。どの値が最小であるかを知らなければ、四つの可能な和を形成し、それらの中で $n(1)$ 以上または等しいものを除外します。したがって、上記の例では以下のようになります：
\begin{align*}
&nx + n(l-x) = n(l).\\
&n(x) + n(1-y) = n(1) + n(x)-n(y).\\
&n(y) + n(l-y) = n(l) + n(y)-n(x).\\
&n(y) + n(1-y) = n(1).
\end{align*}
最初と最後の値を除外すると、必要な式は以下のようになります：
\[n(1) + n(x)-n(y), \text{ および } n(1) + n(y)-n(x).\]
これらのうちの一つ（$nx = ny$ の場合を除く）が $n(l)$ よりも小さいものであり、もう一つは大きいものです。最小値は当然、それらの中から選ばれます。

二つ以上の成分が共通の因子を持つ場合、例えば $x$ という因子は、原則 {\sc i} より明らかに、最終的な主要制限式において単純な項 $n(x)$ のみを提供することができます。したがって、$n(x)$ が二つ以上の成分の主要制限として現れる場合、これらの制限を加算する際に $nx + nx$ を $nx$ に置き換えるなど、そのような操作が必要となります。

例えば、式 $n\{xy + x(1-y)z\}$ の主要制限はすぐに以下のようになります：
\begin{alignat*}{2}
\&1.\; nx. \&4.\;\& ny + nx. \\
\&2.\; nx + n(1-y). \qquad \&5.\;\& ny + n(1-y).  \\
\&3.\; nx + nz. \&6.\;\& ny + nz.
\end{alignat*}
これらのうち、最初と第六項のみを保持する必要があります；第二、第三、第四の項は最初の項よりも大きく、第五項は $n(1)$ に等しいです。したがって、主要制限は
\[ n(x) \text{ と } n(y) + nz\]
となります。これらの二つの値の中で最後のもの、それが $n(1)$ より小さい場合、それを採用します。

これらの考慮から以下のルールを導きます：

\textsc{ルール}---各成分から一つの因子を取り出し、その前に $n$ の符号を付けて、それらの項や結果を加算し、同じ項の重複は排除します。この合計が得られる主要制限であり、その中で最も小さいものが使用されるべき主要制限となります。

例えば、式
\[ xyz + x(1-y)(1-z) + (1-x)(1-y)(1-z)\]
の主要制限は以下のようになります：
\begin{align*}
     n(x) + n(1 - y), &\text{ と } n(x) + n(1 - z),\\
\text{または}\hfill
  n(x) + n(1) - n(y), &\text{ と } n(x) + n(1) - n(z).\hfill
\end{align*}

もし $n(y)$ を最初の項から選んだ場合、第二項から $n(x)$ を選ぶと、第三項からは $n(1-y)$ を選ばなければなりません。これにより以下の結果が得られます：
\[
n(y) + n(x) + n(1 - y), \text{ または } n(1) + n(x).
\]
しかし、この結果は $n(1)$ より大きいので、これはすべてのクラスに対する明らかなる主要制限より大きいため、考慮に入れる必要はありません。

\begin{center}
\textsc{命題 III.}
\end{center}

7. 任意の級を表すシンボル$x, y, z$の成分からなる最小数的限界を求めること。この級は与えられた$n(x), n(y), n(z) \dotsc n(1)$によって表現されているものとする。

この目的は、前節の命題と第\textsc{ii}原理を組み合わせることにより達成されるが、より適切には以下の方法で達成される：

任意の2つの成分を考え、これらが互いに異なる唯一の因子だけ異なっている場合、それらを加算して一つの級項として形成する。例えば$x(1 - y) + xy$は$x$という級項を形成する。このような集約を可能な限り行い、つまり与えられた成分列からできるだけ多くのこの種類の和を抽出し、各和が単一の項に収斂するまで行うと、残された成分は加算できないものとする。これらの最終的な集約と、それらを加算できない成分をそれぞれ異なる項として書き出す。

これらの項に対する命題Iによる最小数的限界は、与えられた表式の最小数的限界であり、その中からただ一つが正となる。

次のように、式 $xy + (1 - x)y + (1 - x)(1 - y)$ から集合 $y$ と $1 - x$ を形成することができます。これは最初の項と第二項を加え、第二項と第三項を加えることにより得られます。したがって、与えられた式の下限は $n(y)$ および $n(1 - x)$ となります。

また、与えられた式が以下のようである場合、

\begin{multline*}
xyz + x (1-y) z + (1-x) yz + (l-x) (1-y) z\\
+ xy(1-z) +(1-x)(1-y)(1-z),
\end{multline*}

最初の四項を加えると単項 $z$、最初の項と第五項を加えると単項 $xy$、第四項と第六項を加えると単項 $(1-x) (1-y)$ を得ることができます。これ以外の集合を単項に集める方法はなく、また未処理の集合も存在しません。これらの三つの結果を得た値が与えられた式の下限となります。

\begin{gather*}
n(z), \; n(x) + n(y) - n(1),\\
\text{および}\quad
n(1-x) + n(1-y) - n(1), \text{または} \quad n(1) - n(x) - n(y).
\end{gather*}

8. 上記の法則の証明は、以下の原理の適切な適用にかかっている：1つ目。任意の構成要素の集合（これらの構成要素が一項に加えられることが可能である）の最小限界は、明らかにその一項の最小限界となる。この点で法則の最初の部分が説明される。2つ目。異なる構成要素または異なる構成要素からなる任意の二項の和の最小限界は、これらを加えることが可能でない場合、それぞれの最小限界的数の合計になる（これらの最小限界的数が両方とも正である場合）。しかし、一方が正であり、他方が負の場合には、その正な最小限界的数だけとなる。なぜなら、負の項が追加された場合、限界的数の値は減少し、すなわち二項の合計よりも一項の限界的数の方が小さくなるからである。現在、二つの構成要素が一つ以上の因子で異なるため、加算することはできない場合、二つの最小限界的数は両方とも正とはならない。したがって、項を $xyz$ と $(1-x)(1-y)z$ とする。これらは二つの因子で異なる。第一項の最小限界は $n(x + y + z - 2)$ であり、第二項の最小限界は $n(1 - x + 1 - y + z - 2)$ または
\[
\text{1つ目。} n\{x + y - 1 - (1 - z)\}. \quad \text{2つ目。} n\{1 - x - y - (1 - z)\}.
\]
もし $n(x + y - 1)$ が正であれば、$n(1 - x - y)$ は負であり、したがって第二項は負である。もし $n(x + y - 1)$ が負であれば、第一項も負であり、多くの因子を含む場合でも同様である。このようにして、法則の適用方法に応じて、異なる項の最小限界的数は加算されないことが示される。なぜなら、これらの項が共通の構成要素を含んでいる場合、その最小限界的数を加算することは不可能であり、あるいは二つの最小限界的数が両方とも正でない場合、加算は無意味であるからである。

\textsc{第四定理}

9. 任意のクラス $s$, $t$ 等に属する個体の数が与えられたとき、論理的に定義された他のクラス $w$ の数的限界系を導き出す。

本章で議論する最も一般的な問題は、前節の調査がその導入に過ぎず、その後の調査がその応用に専念しているため、この問題の性質と目的を明確に述べることが望ましい。

クラス $s$, $t \dotsc w$ が論理的に定義されるとは、それらのクラスが符号的な表現を書き下すことを可能にするような定義であることを意味し、そのクラスが単純であれ複合であれ、この方法論によって符号 $w$ は直接的に符号 $s$, $t$, 等 の発展した関数として決定される。具体的には
\begin{equation*}
w = A + 0B + \frac{0}{0}C + \frac{1}{0}D,      \tag{1}
\end{equation*}
と表される。ここで $A$, $B$, $C$, および $D$ は符号 $s$, $t$, 等 の構成要素から形成される。

このような表現から、$w$ の数値的な制限が最も一般的な方法でどのように決定できるかについては、後ほど考慮する。現在のところ、我々の目的は、前節の命題によって開発された原理に基づいてこの目標をどの程度達成できるかを示すことに留まる。このような調査は本作業の目的に十分である。

簡潔性のために、私の議論は特に次の展開に基づくものとする：
\begin{equation*}
w = st + 0s(1 - t) + \frac{1}{0}(1 - s)t + \frac{0}{0}(1 - s)(1 - t),
\tag{2}
\end{equation*}
ここでは、すべての係数の種類が現れている。

$(1-s)(1-t)$という構成要素が、その係数が$\frac{0}{0}$であるにもかかわらず、$w$の中にどの程度含まれているのかは不明であり、それが該当するクラスの一部がすべて含まないか、部分的に含まないか、または完全に含んでいるかを意味します。明らかに、$(1-s)(1-t)$を表すクラスのすべてのメンバーが$w$の中に存在すると、$n(w)$はその数値的な最大値を持つことになります。さらに、$s(1-t)$や$(1-s)t$を表すクラスに含まれる個体が$w$には存在しないため、$w$の上位の数値的制限はクラス$st + (1-s)(1-t)$と同一となります。したがって、これらの制限は以下のようになります：
\[
ns + n(1-t) \quad \text{および} \quad nt + n(1-s).
\]

同様に、展開式$A + 0B + \frac{0}{0}C + \frac{1}{0}D$の上位の数値的制限は、定理2から$A + C$の上位の数値的制限を用いて求めることができます。

また、$w$の任意の下位の数値的制限は原理\textsc{ii.}により以下の式で与えられます：
\[
n(1) - \text{主要制限}(n(1-w)).
\]
ここで、$w$の展開が（1）によって与えられているため、$1-w$の展開は明らかに以下のようになります：
\[
1 - w = 0A + B + \frac{0}{0}C + \frac{1}{0}D.
\]
これは、第\textsc{x.}章の定理2の方法によって直接証明することができます。したがって
\begin{align*}
\text{下位制限}(n(w)) &= n(1) - \text{主要制限}(B + C) \\
&= \text{下位制限}(A + D),
\end{align*}
原理\textsc{ii.}により、クラス$A + D$と$B + C$が補遺であるためです。したがって、（2）の第二項の下位制限は$n(t)$となります。この論理を一般化すると以下の結果を得ます：

\emph{$A + 0B + \frac{0}{0}C + \frac{1}{0}D$の展開の下位制限の系は、$A + D$の下位制限から与えられます。}

この結果は直接的に導くこともできます。なぜなら、下位の数値的制限を求めることには最大のものを求めることであるため、一般的に、クラス$D$を$w$の表記式において$A$と組み合わせることにより得られる制限が、それ以外の場合よりも高い下位制限を得ることができます。これは我々が論理学の様々な例で許可されている組み合わせです。

最後に、$w$ の展開の終了項が方程式 $D = 0$ を示すことに注意すると、明らかに $n(D) = 0$ である。したがって，
\[
  n(D) \stackrel{=}{<} 0
\]
という不等式が成り立つことが分かります。この方程式は定理3で扱われることにより、必要な条件を数的要素 $n(s)$, $n(t)$ 等について得ることができます。これらの条件は実際的な問題を満たし、そのデータに可能な経験の結果を含めることを保証します。

したがって、(2) の第二項における $\frac{1}{0}(1-s)t$ から
\begin{gather*}
  n(1-s) + n(t) - n(1) \stackrel{=}{<} 0,   \\
  \therefore n(t) \stackrel{=}{<} n(s).
\end{gather*}
これらの結論は次のルールとして表現することができます：

10. ルール．---$w$ のクラスを符号 $s$, $t$, 等の論理関数として展開した式
\[
  w = A + 0B + \frac{0}{0}C + \frac{1}{0}D.
\]
とすると，
\begin{align*}
  w の主極限 &= A の主極限 + C の主極限.   \\
  w の最小極限 &= A の最小極限 + D の最小極限.
\end{align*}
データ間の必要な数的条件は
\[
  D の最小極限 \stackrel{=}{<} n(1)
\]
という不等式によって与えられます。

上記の方法を確率問題における解の制限に適用するには、各 $n(x)$ を Prob. $x$，$n(y)$ を Prob. $y$ 等と置き換えて、最後に $n(1)$ を 1 に置き換えるだけです。しかし、この応用は非常に重要であるため、変形の主な結果をルールとして示すことが望ましいかもしれません。

11. 任意の事象 $s$, $t$ 等の確率が与えられ、その他の事象 $w$ が符号 $s$, $t$ 等の論理関数として展開されているとき
\[
  w = A + 0B + \frac{0}{0}C + \frac{1}{0}D,
\]
とすると、確率 Prob. $w$ の上極限系と下極限系を求め、データ間の条件を求めることを要求します。

\#\#\# 解答

Prob. $(A + C)$ の上界と Prob. $(A + D)$ の下界は、求められるような2つのシステムを形成します。データ中の定数間の条件は以下の不等式によって与えられます：
\[
  \text{Inf. lim. Prob. } D \stackrel{=}{<} 0.
\]

これらの原理に基づく応用において、常に次の関係が成り立ちます：
\[
  \text{Inf. lim. Prob. } x_1 x_2 \dotsc x_n = \text{ Prob. } x_1 + \text{ Prob. }x_2 \dotsc + \text{ Prob. }x_n - (n-1).
\]

また、下界は単項語、つまり与えられたものまたは集合によって形成されたもののみから決定できます。上界は $\sum \text{ Prob. } x$ の形で与えられ、Prob. $x$ は異なるシンボルに適用され、求められる上界の式の異なる項から取り出されます。したがって、Prob. $xyz + x(1-y)(1-z)$ の上界は次のようになります：
\[
  \text{Prob. }x, \quad \text{Prob. }y + \text{ Prob. }(1-z), \quad \text{and Prob. }z + \text{ Prob. }(1-y).
\]

注意すべき点として、最後のケースで Prob. $z$ を最初の項から Prob. $(1-z)$ を二番目の項から取り出すことは許されますが、その合計が 1 となるため、結果として無駄な \textit{\`{a} priori} 必然的なものとなります。明らかに、0と1の間に入らない限界は排除できます。

最後の章のCase \textsc{iii.}に対するこの方法を適用します。

最終的な論理的解は次のようになります：
\[
\begin{split}
  x &= \frac{0}{0}stu + \frac{1}{0}s\bar{t}u
     + \frac{1}{0}st\bar{u} + s\bar{t}\bar{u}   \\
    &+ \frac{1}{0}\bar{s}tu + 0\bar{s}\bar{t}u
     + 0\bar{s}t\bar{u} + 0\bar{s}\bar{t}\bar{u},
\end{split}
\]
データは次の通りです：
\[
 \text{Prob. }s = p, \quad \text{Prob. }t = q, \quad \text{Prob. }u = r.
\]

$x$ の数値的限界と $p$, $q$, および $r$ の関係を両方求めます。

$x$ の上界は、ルールに従って $stu + s\bar{t}\bar{u}$ の上界で与えられます。したがって、それは次のようになります：
\[
  p, q + 1 - r, \quad r + 1 - q.
\]

$x$ の下界は次の式の下界から与えられます：
\[
  s\bar{t}u + st\bar{u} + s\bar{t}\bar{u} + \bar{s}tu.
\]

これらの成分のうち、最初と第三を一つの項として $s\bar{t}$ にまとめ、第二と第三を一つの項として $s\bar{u}$ にまとめることができます。$x$ の下限は、$s(1-t)$, $s(1-u)$, $(1-s)tu$ の各項から別々に導き出されます。これらにより、

\begin{multline*}
  p + 1 - q - 1,\quad p + 1 - r - 1,\quad 1 - p + q + r - 2,   \\
\text{または }\hfill p - q,\quad p - r,\quad \text{および }q + r - p- 1.
\end{multline*}

が得られます。

最後に、定数 $p$, $q$, および $r$ の間の条件は、

\[
  s\bar{t}u, \quad st\bar{u}, \quad \bar{s}tu
\]

から与えられ、これらの項からルールにより、

\begin{align*}
& p + 1 - q + r - 2 \stackrel{=}{<} 0, \quad
  p + q + 1 - r - 2 \stackrel{=}{<} 0, \quad
  1 - p + q + r - 2 \stackrel{=}{<} 0. \\
&\text{または } \hfill
  1 + q - p - r \stackrel{=}{>} 0, \quad
  1 + r - p - q \stackrel{=}{>} 0, \quad
  1 + p - q - r \stackrel{=}{>} 0. \hfill
\end{align*}

が導き出されます。これらは最終的な解の分析で用いられる制限条件です。$\lambda$ の制限条件は、$s$, $t$, および $u$ が正であるという条件から単純に決定されました。それらよりさらに狭い制限を推定する可能性は十分にあります。

12. 次の応用例は重要な問題の一例であり、その解法は次の章で述べられます。与えられた条件は、

\[
  \text{Prob. }x = c_1,   \quad
  \text{Prob. }y = c_2,   \quad
  \text{Prob. }s = c_1p_1,\quad
  \text{Prob. }t = c_2p_2,
\]

および論理方程式

\[
\begin{split}
  z = stxy + s\bar{t}x\bar{y} + \bar{s}t\bar{x}y + 0\bar{s}\bar{t} \\
  + \frac{1}{0}\left\{
  \begin{split}
    stx\bar{y} + st\bar{x}y + st\bar{x}\bar{y} + s\bar{t}xy
      + s\bar{t}\bar{x}y   \\
    + s\bar{t}\bar{x}\bar{y} + \bar{s}txy + \bar{s}tx\bar{y}
      + \bar{s}t\bar{x}\bar{y};
\end{split}
\right.
\end{split}
\]

から、定数 $c_1$, $c_2$, $p_1$, $p_2$ および $z$ の主要制限と最小制限を決定することです。

まず、定数間の条件を求めてみましょう。係数が $\frac{1}{0}$ の項に限定して項を集積することで、以下の項を得ることができます：
\[
s(1-x), \qquad t(1-y), \qquad sy(1-t), \qquad tx(1-s);
\]
これらの項以外の項は作ることができません。したがって、定数間の条件は次のようになります。
\begin{align*}
&n(s) + n(1-x) - n(1) \stackrel{=}{<} 0,\\
&n(t) + n(1-y) - n(1) \stackrel{=}{<} 0,\\
&n(s) + n(y) + n(1-t) - 2n(1) \stackrel{=}{<} 0,\\
&n(t) + n(x) + n(1-s) - 2n(1) \stackrel{=}{<} 0.
\end{align*}

次に、$n(x)$ を $c_1$, $n(y)$ を $c_2$, $n(s)$ を $c_1 p_1$, $n(t)$ を $c_2 p_2$, $n(1)$ を $1$ と置き換えると、わずかな簡略化の後で以下の式を得ます：
\[
\begin{array}{cccccc}
  c_1 p_1 \stackrel{=}{<} c_1, \qquad
& c_2 p_2 \stackrel{=}{<} c_2,\\
  c_1 p_1 \stackrel{=}{<} 1 - c_2(1-p_2), \qquad
& c_2 p_2 \stackrel{=}{<} 1 - c_1(1-p_1).
\end{array}
\]
これらが定数間の必要条件です。

また、$z$ の主要限界は以下の式
\[
stxy + s(1-t) x(1-y) + (1-s)t(1-x)y;
\]
の限界と一致します。上記で導かれた条件
\[
n(s) \stackrel{=}{<} n(x), \qquad n(t) \stackrel{=}{<} n(y)
\]
を考慮すると、これらの限界は以下のようになります：
\begin{align*}
n(s) + n(t),   &\text{ または, } c_1p_1 + c_2p_2,\\
n(s) + n(1-x), &\text{ または, } 1 - c_1 (1-p_1),\\
n(t) + n(1-y), &\text{ または, } 1 - c_2 (1-p_2).
\end{align*}

最後に、$z$ の小限界を確認するために、係数が $1$ や $\frac{1}{0}$ の項から単純な項 $s$ と $t$ を形成します。これらの項以外の項は選択や集積によって作ることができません。したがって、$z$ の小限界は $c_1p_1$ と $c_2p_2$ です。

13. 上記で開発された方法は必ずしも可能な最も狭い限界を常に指定するわけではありませんが、確率論における問題の解の制限としては十分であると考えられます。

数値拡張のクラスの最狭限を決定する問題は、常に純粋な代数学的形式に帰着される。\footnote{著者は四年前に書いた手紙を失ったことを悔やみます。彼はこの方法が相当長く発展されていると信じています。彼の記憶には、内容自体よりもその方法の原理が上記のように説明されたものと同じで、その充足性が証明されたという印象だけが残っています。この章の前段階の方法は、ほとんど言うまでもなく容易ですが、もちろんより一般的ではありません。}

したがって、方程式
\[
  w = A + 0B + \frac{0}{0}C + \frac{1}{0}D,
\]
を再考し、$w$ の最高下限の数値限界を $an(s) + bn(t) \dotsc + dn(1)$ で表すとする。ここで $a, b, c,\dotsc d$ は定数であり、$s$, $t$, 等 は $A$, $B$, $C$, $D$ の構成要素である論理的シンボルである。

したがって
\[
\begin{split}
an(s) + bn(t)\dotsc + dn(1)=\text{最小限の $A$ を条件 $D = 0$ に従って}\\
\text{定める}。
\end{split}
\]

したがって、関数
\[
  as + bt \dotsc + d,
\]
を展開し、結果から $D$ 内の構成要素を見つけたものをすべて取り除き、残った構成要素の係数について、それらが $A$ 内に存在する場合、個々の値は 1 を超えないべきであり、それらが $A$ 内に存在しない場合、個々の値は 0 よりも大きくならないべきである。したがって、不等式の系列
$f \stackrel{=}{<} 1$
と
$g \stackrel{=}{<} 0$
が得られ、$f$ および $g$ は $a, b, c,\dotsc d$ の線形関数である。これらの定数 $a, b \dotsc d$ の値を決定し、上記の条件を満たしながら
\[
  an(s) + bn(t) \dotsc + dn(1),
\]
の最大値を与えるものを見つけ、その最大値が $w$ の最高最小限となる。さらに、与えられた定数 $n(s), n(t) \dotsc n(1)$ の関係を同様に決定するための関係も加えることができる。

14. 次のやや複雑な例は、問題が任意の要素を含む場合、その要素がデータが一致するが求める値（qu{\ae}sita）は無制限である一連の問題の代表となるとき、解の制限がどのように行われるかを示します。

**問題。** $n$ つの事象 $x_1, x_2, \dotsc, x_n$ について、次の特定情報が知られています：

1. 事象 $x_1$ の発生またはすべての事象の失敗の確率は $p_1$ です。
2. 事象 $x_2$ の発生またはすべての事象の失敗の確率は $p_2$ です。他のものも同様です。

任意の単一の事象または事象の組合せ、一般的な関数形 $\phi(x_1, \dotsc, x_n)$ または $\phi$ の確率を求めることを要求します。

以前の記号を使用して、問題のデータは以下のようになります：
\[
  \text{確率} (x_1 + \bar{x}_1 \dotsc \bar{x}_n) = p_1, \quad
  \text{確率} (x_n + \bar{x}_1 \dotsc \bar{x}_n) = p_n.
\]
そして $\phi(x_1, \dotsc, x_n)$ の確率を求めます。

一般に、次のように仮定します：
\begin{gather*}
  x_r + \bar{x}_1 \dotsc \bar{x}_n = s_r,    \tag{1}   \\
  \phi = w.                            \tag{2}
\end{gather*}

これにより、問題の集団論理方程式が得られます：
\[
  \sum \bigl\{\bigl( x_r + \bar{x}_1 \dotsc \bar{x}_n \bigr)\bar{s}_r
    + s_r \bigl( \bar{x}_r - \bar{x}_1 \dotsc \bar{x}_n \bigr)\bigr\}
  + \phi\bar{w} + w\bar{\phi} = 0.  \tag{3}
\]

この方程式から、シンボル $x_1, \dotsc, x_n$ を消去し、$s_1, \dotsc, s_n$ の展開された論理関数として $w$ を決定する必要があります。

以前の消去結果を以下のように表すとします：
\[
  Ew + E'(1-w) = 0;
\]
すると $E$ は以下の方程式から同じシンボルを消去した結果となります：
\[
  \sum \bigl\{\bigl( x_r + \bar{x}_1 \dotsc \bar{x}_n \bigr)\bar{s}_r
    + s_r \bigl( \bar{x}_r - \bar{x}_1 \dotsc \bar{x}_n \bigr)\bigr\}
  + 1 - \phi = 0.  \tag{4}
\]

現在、$E$ は上記方程式の最初の項の展開において見つかったすべての成分（$x_1, x_2, \ldots, x_n$ の符号に関して）の係数の積となる。また、$\phi$ および従って $1 - \phi$ も、それぞれの係数が 1 のような成分からなる級数で構成される。方程式 (4) の最初の項の展開における係数の形を決定する際には、次のように並べる方が便利である：

1. $1 - \phi$ 内に見つかった成分の係数。
2. $\phi$ 内に見つかった $\bar{x}_1, \bar{x}_2, \ldots, \bar{x}_n$ の係数（もし存在する場合）。
3. $\phi$ 内に見つかった他の成分の係数、$\bar{x}_1, \bar{x}_2, \ldots, \bar{x}_n$ を除く。

これらは明らかに完全な分類である。

まず第一に、$1 - \phi$ 内に見つかった任意の成分の係数は、方程式 (4) の最初の項の展開において次の形となる：
\[
  1 + \text{正の項}\left(\sum\right)
\]
従って、このようなすべての係数は 1 で置き換えることができる。命題 \textsc{i.} 第 IX 章。

次に第二に、$\phi$ 内に見つかった $\bar{x}_1, \bar{x}_2, \ldots, \bar{x}_n$ の係数（もし存在する場合）は方程式 (4) の最初の項の展開において次の形となる：
\[
  \sum \bar{s}_r = \bar{s}_1 + \bar{s}_2 + \cdots + \bar{s}_n
\]

最後に第三に、$\phi$ 内に見つかった他の任意の成分 $x_1, x_2, \ldots, x_i, \bar{x}_{i+1}, \ldots, \bar{x}_n$ の係数は方程式 (4) の最初の項の展開において次の形となる：
\[
  \bar{s}_1 + \bar{s}_2 + \cdots + \bar{s}_i + s_{i+1} + \cdots + s_n
\]

現在、$E$ はすべての係数の積であることが確認される。しかし、$\phi$ に存在しない因子の係数は1となるため、$E$ は $\phi$ 内に存在する因子の係数の積として考えることができる。これらの係数が形成される方法から、$E$ の決定法を以下のルールで導き出すことができる：$\phi$ 内に見られる各因子（$\bar{x}_1 \; \bar{x}_2 \dotsc \bar{x}_n$ 以外）において、$x_1$ に対して $\bar{s}_1$ を、$\bar{x}_l$ に対して $s_1$ を書き換えて結果を加える。また、$\phi$ 内に $\bar{x}_1, \bar{x}_2 \dotsc \bar{x}_n$ が存在する場合、その因子に対して $\bar{s}_1 + \bar{s}_2 \dotsc + \bar{s}_n$ を書き換える。これらの和の積が $E$ となる。

$E'$ の求め方については、(3) 式において $w = 0$ とし、$x_1, x_2 \dotsc x_n$ を除去した式から導き出す必要がある。その結果得られる式は以下のようになる：
\[
  \sum \bigl\{
    \bigl( x_r + \bar{x}_1 \dotsc + \bar{x}_n \bigr)\bar{s}_r
         + s_r\bigl(\bar{x}_r - \bar{x}_1 \dotsc \bar{x}_n \bigr)
  \bigr\} + \phi = 0.  \tag{5}
\]
%-----------------------File: 116.png----------------------------
したがって、$E'$ は $1 - \phi$ の因子から形成される。つまり、$\phi$ 内に存在しない因子の同様な方法で $E$ が形成されるように $E'$ を形成する。

次に、方程式を考慮する：
\[
  Ew + E'(1-w) = 0.
\]
これにより
\[
  w = \frac{E'}{E'-E}.   \tag{6}
\]

$E$ と $E'$ は符号 $s_1, s_2 \dotsc s_n$ の関数であるため、$w$ の展開値はこれらの符号から形成されるすべての因子を含み、その各因子に適切な係数が付加されることで決定される。

さらに、$E$ と $E'$ はそれぞれ因子の乗積で形成され、これらがいずれも消滅するためには、その構成要素のうち少なくとも一つが消滅しなければならない。また、任意の因子 $\bar{s}_1 + \cdots + \bar{s}_n$ は、その形成に用いられたすべての項が同時に値を0と1とするときのみ消滅する。すなわち、これらの項がそれぞれ0または1となるためである。さらに明らかに、異なる因子は同時に消滅しない。したがって、因子 $\bar{s}_1 + \bar{s}_2 + \cdots + \bar{s}_n$ と $s_1 + \bar{s}_2 + \cdots + \bar{s}_n$ は同時に消滅しない。なぜなら、前者は $\bar{s}_1 = 0$ または $s_1 = 1$ のときのみ消滅し得るが、後者は $s_1 = 0$ のときのみ消滅するからである。

まず、$w$ の展開値における因子 $\bar{s}_1 \bar{s}_2 \cdots \bar{s}_n$ の係数を決定する。$\bar{s}_1 = 1$, $\bar{s}_2 = 1, \dotsc, \bar{s}_n = 1$ という同時仮定が成り立つと、因子 $s_1 + s_2 + \cdots + s_n$ は $E$ や $E'$ の中に含まれる場合に消滅する。このとき、他のどの因子も消滅しない。しかしながら、$s_1 + s_2 + \cdots + s_n$ はどちらの $E$ や $E'$ の中にも因子として現れず、したがってこれらの量はいずれも消滅しない。よって、比 $\frac{E'}{E' - E}$ は1、0、または $\frac{0}{0}$ ではない。

ゆえに、$w$ の展開値における因子 $\bar{s}_1 \; \bar{s}_2 \cdots \bar{s}_n$ の係数は $\frac{1}{0}$ で表される。

次に、成分 $s_1 s_2 \dotsc s_n$ の係数を決定しましょう。

まず、仮定 $s_1 = 1, s_2 = 1, \dotsc, s_n = 1$ をすると因子 $\bar{s}_1 + \bar{s}_2 + \dotsb + \bar{s}_n$ は消えてしまいます。この因子が $E$ に存在し、$E'$ にはない場合、$\phi$ が成分 $x_1 x_2 \dotsc x_n$ と $\bar{x}_1 \bar{x}_2 \dotsc \bar{x}_n$ の両方を含むとき、$\frac{E'}{E'-E}$ は $\frac{E'}{E'} = 1$ となります。この因子が $E'$ に存在し、$E$ にはない場合、$\phi$ がこれらの成分のどちらも含まないとき、$\frac{E'}{E'-E}$ は $\frac{0}{-E} = 0$ となります。最後に、一方の成分 $x_1 x_2 \dotsc x_n$ と $\bar{x}_1 \bar{x}_2 \dotsc \bar{x}_n$ のうちのどちらかのみが $\phi$ 内に存在し、もう片方は存在しない場合、この因子は $E$ と $E'$ 両方に含まれます。このとき、$\frac{E'}{E'-E}$ は $\frac{0}{0}$ となります。

したがって、成分 $s_1 s_2 \dotsc s_n$ の係数は以下のようになります：
- $\phi$ が成分 $x_1 x_2 \dotsc x_n$ と $\bar{x}_1 \bar{x}_2 \dotsc \bar{x}_n$ の両方を含むとき、その係数は $1$。
- $\phi$ がこれらの成分のどちらも含まないとき、その係数は $0$。
- 一方の成分のみが存在し、もう片方は存在しないとき、その係数は $\frac{0}{0}$。

次に、他の任意の成分 $s_1 \dotsc s_i \bar{s}_{i+1} \dotsc \bar{s}_n$ の係数を決定します。

仮定 $s_1 = 1, \dotsc, s_i = 1, s_{i+1} = 0, \dotsc, s_n = 0$ をすると因子 $\bar{s}_1 + \bar{s}_2 + \dotsb + \bar{s}_i + s_{i+1} + \dotsb + s_n$ は消えてしまいます。この因子が $E$ 内に存在し、$E'$ 内にはない場合、成分 $x_1 x_2 \dotsc x_i \bar{x}_{i+1} \dotsc \bar{x}_n$ が $\phi$ 内に存在するとき、$\frac{E'}{E'-E}$ は $\frac{E'}{E'} = 1$ となります。この因子が $E'$ 内に存在し、$E$ 内にはない場合、成分 $x_1 x_2 \dotsc x_i \bar{x}_{i+1} \dotsc \bar{x}_n$ が $\phi$ 内に存在しないとき、$\frac{E'}{E'-E}$ は $\frac{0}{-E} = 0$ となります。

\emph{したがって、任意の他の成分 $s_1 \dots s_i, \; \bar{s}_{i+1} \dots \bar{s}_n$ の係数は、類似の成分 $x_1 \dots x_i, \; \bar{x}_{i+1} \dots \bar{x}_n$ が $\phi$ 内に存在するか否かにより 1 または 0 となる。}

従って、$\phi$ の与えられた表現から成分 $x_1, \; x_2 \dots x_c$ および $\bar{x}_1, \; \bar{x}_2 \dots \bar{x}_n$ を除外し、これらがいずれかまたは両方とも含まれている場合、符号 $x_1, \; x_2, \dotsc x_n$ をそれぞれ $s_1, \; s_2, \dotsc s_n$ に変更します。成分 $s_1 \; s_2 \dots s_n$ および $\bar{s}_1 \; \bar{s}_2 \dots \bar{s}_n$ の係数を上記で与えられた特別なルールに基づいて決定し、他のすべての成分は係数 0 とします。その結果が $w$ の値となります。

特に、$\phi = x_1$ の場合を考えます。与えられたデータからイベント $x_1$ の確率を決定することが求められます。

符号 $x_1$ を符号列 $x_1, \; x_2, \dotsc x_n$ で展開すると、$x_1$ の因子を持つすべての成分が生成されます。その中には成分 $x_1 \; x_2 \dots x_n$ は含まれますが、成分 $\bar{x}_1 \; \bar{x}_2 \dots \bar{x}_n$ は含まれません。

したがって、符号 $s_1, \; s_2, \dotsc s_n$ を用いた $x_1$ の展開値において、成分 $s_1 \; s_2 \dots s_n$ の係数は $\frac{0}{0}$ であり、成分 $\bar{s}_1 \; \bar{s}_2 \dots \bar{s}_n$ の係数は $\frac{1}{0}$ となります。

$x_1$ から成分 $x_1 \; x_2 \dots x_n$ を除外すると、結果は $x_1 - x_1x_2 \dots x_n$ となり、その中で $x_1$ を $s_1$ に置き換えると、$s_1 - s_1s_2 \dots s_n$ が $x_1$ の符号 $s_1, \; s_2, \dotsc s_n$ に対する表現の対応する部分となります。

$x_1$の最終的な表現は以下のようになります：
\begin{equation*}
\begin{split}
  x_1 = s_1 - s_1s_2 \dotsc s_n
      + \frac{0}{0}s_1s_2 \dotsc s_n
      + \frac{1}{0}\bar{s}_1\bar{s}_2 \dotsc \bar{s}_n   \\
  + \text{係数が }0 \text{ の成分}.
\end{split}
  \tag{7}
\end{equation*}

上記展開中のすべての成分の合計、つまり係数が1, 0, $\frac{0}{0}$である成分の和は
$1 - \bar{s}_1\bar{s}_2 \dotsc \bar{s}_n$となります。

したがって、確率$x_1$を決定するための以下の代数系を得ます：
\[
  \text{Prob. }x_1 = \frac{s_1 - s_1s_2 \dotsc s_n + cs_1s_2 \dotsc s_n}{1 - \bar{s}_1\bar{s}_2 \dotsc \bar{s}_n},  \tag{8}
\]
ただし、以下の関係が成り立ちます：
\begin{equation*}
\begin{split}
  \frac{s_1}{p_1} = \frac{s_2}{p_2} \dotsc = \frac{s_n}{p_n}   \\
  = 1 - \bar{s}_1 \bar{s}_2 \dotsc \bar{s}_n = \lambda.
\end{split}
   \tag{9}
\end{equation*}

これらの関係は、関数$\phi$の形に依存せず、一度定まったこれらの量は、データが同じであればどのような問題であっても同様に使用できます。求めたい事象またはその組合せの性質は、決定された$s_1 \; s_2 \dotsc s_n$の値を置き換える関数の形だけ影響を与えます。

(9)から得られる式は以下のようになります：
\[
  s_1 = p_1 \lambda, \quad
  s_2 = p_2 \lambda, \dotsc
  s_n = p_n \lambda.
\]
従って，
\[
  1 - (1 - p_1 \lambda)(1 - p_2 \lambda) \dotsc (1 - p_n \lambda)
 = \lambda.
\]
または，
\[
  1 - \lambda
  = (1 - p_1 \lambda)(1 - p_2 \lambda) \dotsc (1 - p_n \lambda);  \tag{10}
\]
この方程式から$\lambda$の値を求めます。

この$\lambda$の値が決定すれば、確率$x_1$は以下のようになります：
\[
  \frac{p_1\lambda - (1-c)p_1p_2 \dotsc p_n\lambda^n}
       {1 - (1-p_1\lambda)(1-p_2\lambda)\dotsc (1-p_n\lambda)},
\]
または、(10)を用いて簡略化すると
\[
  \text{Prob. }x_1 = p_1 - (1-c)p_1p_2 \dotsc p_n \lambda^{n-1}. \tag{11}
\]

次に、確率$x_1$の値が与えられるべき条件を求めてみましょう。その条件は定数$p_1, p_2, \dotsc, p_n$と確率$x_1$の限界値との間で満たされるべきものであることが求められます。

係数が$\frac{1}{0}$となる項は一つだけなので、定数に対する条件も一つだけあります。すなわち、

\begin{gather*}
  \text{小限界, } (1-s_1)(1-s_2) \dotsc (1-s_n)
    \stackrel{=}{<} 0. \\
  \begin{array}{cccccc}
   \text{または,} &
     n(1-s_1) + n(1-s_2) \dotsc + n(1-s_n) - (n-1)n(1)
     \stackrel{=}{<} 0. \\
   \text{または,} & n(1) - n(s_1) - n(s_2) \dotsc - n(s_n)
     \stackrel{=}{<} 0. \\
   \text{従って} & p_1 + p_2 \dotsc + p_n
     \stackrel{=}{>} 1,
  \end{array}
\end{gather*}

これが求められる条件です。

確率$x_1$の大限界は、係数が1または$\frac{0}{0}$の各成分の和の大限界と一致します。その和は$s_1$であるため、

\[
  \text{大限界, 確率 }x_1 = \text{ 大限界 }s_1 = p_1.
\]

次に、確率$x_1$の小限界は以下の式の小限界的と一致します：

\[
  s_1 - s_1 s_2 \dotsc s_n + (1-s_1)(1-s_2) \dotsc (1-s_n).
\]

上記から形成できる異なる集合体、すなわち最も多くの成分を含む各項は以下の通りです：

\[
  s_1(1-s_2),\quad s_1(1 - s_3), \dotsc, s_1(1 - s_n),\quad
  (1 - s_2)(1 - s_3) \dotsc (1 - s_n).
\]

これらの式から、小限界は以下のようになります：

\[
  p_1 - p_2,\quad p_1 - p_3, \dotsc, p_1 - p_n,\quad
  1 - p_2 - p_3 \dotsc - p_n.
\]

したがって、確率$x_1$の値はこれらのいずれかを下回ることなく、かつ$p_1$を超えることはありません。

しかし、これらの条件を用いることなく、確率の一般方法の証明において述べられた原理を利用することができます。$s_1, s_2, \dotsc s_n$ がそれぞれ 1 より小さいという条件は、$\lambda$ がそれぞれの量 $\frac{1}{p_1}, \frac{1}{p_2}, \dotsc \frac{1}{p_n}$ よりも小さいことを意味します。また、$s_1, s_2, \dotsc s_n$ がそれぞれ 0 より大きいという条件は、$\lambda$ が 0 より大きいことを意味します。

ここで $p_1 p_2 \dotsc p_n$ が適切な分数であり、条件 $p_1 + p_2 \dotsc + p_n > 1$ を満たすとします。このとき、中心方程式（10）から正の $\lambda$ の値を導き出し、それぞれの量 $\frac{1}{p_1}, \frac{1}{p_2}, \dotsc \frac{1}{p_n}$ よりも小さいものがただ一つ存在することが示されます。その $\lambda$ の値が求められるものです。

これを証明するために、方程式
\[
  (1 - p_1\lambda)(1 - p_2\lambda) \dotsc (1 - p_n\lambda)
  - 1 + \lambda = 0
\]
を考慮します。

$\lambda = 0$ のとき、最初の項は消え、方程式が満たされます。$\lambda = 0$ から $\lambda = \frac{1}{p_1}$ の範囲での最初の項の変化を調べるために、$p_1$ が $p_1, p_2, \dotsc p_n$ の最大値であると仮定します。

方程式の最初の項を $V$ と表すと、
\[
  \frac{dV}{d\lambda}
  = -p_1(1-p_2\lambda) \dotsc (1-p_n\lambda)
    - p_n(1-p_1\lambda) \dotsc (1-p_{n-1}\lambda) + 1,
\]
$\lambda = 0$ のとき、この式は $-p_1 - p_2 \dotsc - p_n + 1$ となり、負の値を取ることになります。

また、
\[
  \frac{d^2V}{d\lambda^2}
  = p_1p_2(1-p_3\lambda)(1-p_n\lambda) + \text{等,}
\]
と表され、与えられた制限の下では各項が正であることが示されます。

最後に、$\lambda = \frac{1}{p_1}$ のとき、
\[
  V = -1 + \frac{1}{p_1},
\]
となり、これは正です。

すべてのこれらのことが示唆するところでは、曲線を構築し、その縦座標が$\lambda$に対応する$V$の値を表すとき、その曲線は原点を通ります。また、$\lambda$の小さな値に対しては、その曲線は横座標の下にあります。この曲線の凸性は、$\lambda = 0$と$\lambda = \frac{1}{p_1}$の間で下向きであり、極限値$\lambda = \frac{1}{p_1}$においては、その曲線は横座標の上にあり、縦座標が正となります。この説明から、指定された範囲内では、つまり$\lambda = 0$と$\lambda = \frac{1}{p_1}$の間で、一度かつしか交差しないことがわかります。

問題の解は、したがって(11)式によって表され、$\lambda$の値は方程式(10)の根であり、その根が指定された範囲内、すなわち$0, \frac{1}{p_1}, \frac{1}{p_2}, \dotsc, \frac{1}{p_n}$の間にあるものである。

定数$c$は明らかに、すべての事象$x_1, x_2, \dotsc, x_n$が同時に起こるか、または同時に失敗する確率を表します。

この$\lambda$の値の決定は、考慮された問題においてデータが同じであれば全ての問題に対して十分です。これは、以前の議論から予想される通り、関数$\phi$の形に依存しない独立な決定となります。

別の例として、以下の条件を考えましょう：
\[
  \phi = \text{ or }
  w = x_1(1-x_2) \cdots (1-x_n) + x_n(1-x_1) \cdots (1-x_{n-1}).
\]
これは、イベント$x_1, x_2, \ldots, x_n$のうち一つだけが発生する確率を要求することと等価です。$w$の値は明らかに以下のようになります：
\[
  w = s_1(1-s_2) \cdots (1-s_n) + s_n(1-s_1) \cdots (1-s_{n-1}) + \frac{1}{0}(1-s_1) \cdots (1-s_n),
\]
これより、確率は以下のようになります：
\begin{gather*}
  \text{Prob. }
    \left\{ x_1(1-x_2) \cdots (1-x_n) + x_n(1-x_1) \cdots (1-x_{n-1}) \right\} \\
  = \frac{s_1(1-s_2) \cdots (1-s_n) + s_n(1-s_1) \cdots (1-s_{n-1})}
         {1 - (1-s_1) \cdots (1-s_n)} \\
  = \frac{p_1\lambda(1-p_2\lambda) \cdots (1-p_n\lambda) + p_n\lambda(1-p_1\lambda) \cdots (1-p_{n-1}\lambda)}
         {\lambda} \\
  = \frac{p_1(1-\lambda)}{1-p_1\lambda}
  + \frac{p_2(1-\lambda)}{1-p_2\lambda} \cdots
  + \frac{p_n(1-\lambda)}{1-p_n\lambda}.
\end{gather*}

この解法は、序論（I.16）で述べた注釈を説明するのに適しています。問題の主要な困難性はデータの性質に由来し、その解答（qu{\ae}sita）の性質には依存していません。$\lambda$が決定される中心方程式とそれに伴う特殊な議論は、元の文书中の任意要素の解釈を変えることで問題がどのような形を取り得るかに関わらず、同等に適用されます。

\newpage

\chapter{原因と結果の連鎖に関する問題（Problems on Causes）}
\label{ch:xx}

\begin{center}
\large\textit{原因の確率・観測結果からの事後推論・ベイズ逆確率の代数的再構築}
\end{center}
\vspace{4mm}

\#\# 因果関係に関する問題

科学の最終的な目的は、すべての個別の事例において原因と結果の関係を理解し、自然の中でどのようにつながっているかの順序に従って、我々には未知である作用機序（\textit{modus operandi}）によって思考の中で両端を結ぶことです。この論文では、知的機能の構成においてこのような目的に対する特別な関心があることを示しました。

考えられる領域には、事物が宇宙の一瞬間的な部分としてのみ理解される領域がありますが（第十一章参照）、またある領域（第十二章）では、これらが途切れない連鎖の一環として認識され、人間の目から見れば無限に続くように見える連鎖の一環として認識されます。そのような系列を観察するときには、原因と結果の関係が他のすべての関係よりも優越していることを避けては通れません。

ここでは、上記の関係が含まれる抽象的な問題について考慮することを提案します。このような問題は、その関係に関連する性質から予想されるように、広範な多様性があります。

事前確率または経験によって与えられた原因の確率と、その効果に対するそれぞれの関連確率から、要求されるのはその効果の確率を決定することです。そしてこれは、1つには絶対的に、2つには特定の条件のもとで行われます。このような目的に、以下の問題のいくつかは関連しています。

一方で、観測された効果と所述する原因の単独または結合によるそのような効果の生産に対する既知の傾向から、特定の原因の確率や原因システム内の特定の関係の確率を決定することも要求されることがあります。この種類の問題は、章の後半で考慮されますし、一般的な調査の他の形態についても触れられます。

私はここで述べたいと思いますが、これらの例は主に方法の説明のために設計されたものであり、その方法の適用における容易さや便利性については考慮されていません。むしろ、それらは成功いかんに関わらず、より複雑な現実と視覚的な表現から原因と効果の関係を研究する際に生じる問題のクラスのタイプとして設計されました。

2. 以下に示す問題は、二つの原因 \(A_1\) と \(A_2\) の確率がそれぞれ \(c_1\) および \(c_2\) であるとする。原因 \(A_1\) が現れた場合に、事件 \(E\) が伴う確率（この原因 \(A_1\) の結果としてか否かを問わず）は \(p_1\) であり、また原因 \(A_2\) が現れた場合に、そのイベント \(E\) が伴う確率（この原因 \(A_2\) の結果としてか否かを問わず）は \(p_2\) である。さらに、両方の原因 \(A_1\) および \(A_2\) が存在しない限り、イベント \(E\) は現れない。

\footnote{このようなデータが経験によってどのように提供されるかという方法は、容易に想像できる。私は書いている部屋の窓の外には、リーヴァー川の上流からの洪水と海洋の潮汐による二つの原因で溢れやすい田圃がある。観測が \(N\) 回行われた結果として、次のようなデータを得たとする：洪水によって河が増水した回数は \(A\) 回であり、そのうち河が氾濫した回数は \(P\) 回である。潮汐によって河が増水した回数は \(B\) 回であり、そのうち河が氾濫した回数は \(Q\) 回である。そして、両方の原因が存在しない限り、この田圃は氾濫しないとする。これらの条件のもとで、この田圃が全体として何の確率で氾濫するかを決定することを求められる。

ここでは、一般問題における要素 \(a\), \(b\), \(p\), \(q\) は比
\[
  \frac{A}{N},\; \frac{P}{A},\; \frac{B}{N},\; \frac{Q}{B}
\]
を表し、あるいはこれらの比が無限に増加するときの値に近づく値を表す。}

事件 \(E\) の確率を求めよ。

原因 \(A_1\) と \(A_2\) が互いに排他的である場合におけるこの問題の解は、広く知られており
\[
  \text{Prob. } E = c_1p_1 + c_2p_2;
\]
であり、これは受容された確率論の基本的な重要な原理の特定の一例を表している。ここでは、上記の制限から解放された問題を解くことが提案される。

まず、以下のように表すことにしましょう：
\[
\begin{array}{cccccc}
  A_1 の原因を x とする。\\
  A_2 の原因を y とする。\\
  E の結果を z とする。
\end{array}
\]
すると、以下の数値データがあります：
\begin{align*}
  &x = c_1 の確率,     &y = c_2 の確率,   \\
  &xz = c_1p_1 の確率, &yz = c_2p_2 の確率。 \tag{1}
\end{align*}

また、A\_1, A\_2 の両方が存在しない場合、結果 E は発生しないという条件が与えられていますから、以下の論理方程式が成り立ちます：
\[
(1-x)(1-y) = v(1-z).
\]
これを v を消去すると、
\[
z(1-x)(1-y) = 0. \tag{2}
\]

次に、以下を仮定します：
\[
xz = s, \quad yz = t. \tag{3}
\]

これらの方程式を簡略化し、結果を (2) と結びつけると、
\[
xz(1-s) + s(1-xz) + yz(1-t) + t(1-yz) + z(1-x)(1-y) = 0. \tag{4}
\]

この方程式から z を x, y, s, t の発展した論理関数として決定し、その確率を (5) のデータを使用して導き出します：
\[
x = c_1 の確率, \quad y = c_2 の確率, \quad
s = c_1p_1 の確率, \quad t = c_2p_2 の確率。 \tag{5}
\]

現在、式（4）を \( z \) について展開し、\(\bar{x}\) を \(1 - x\)、\(\bar{y}\) を \(1 - y\)、そして同様に置き換えると、

\begin{gather*}
(x\bar{s} + s\bar{x} + y\bar{t} + t\bar{y} + \bar{x}\bar{y})z
+ (s + t)\bar{z} = 0, \\
\therefore z + \frac{s+t}{s + t - x\bar{s} - s\bar{x} - y\bar{t} - t\bar{y} - \bar{x}\bar{y}} = \\
stx + \frac{1}{0}stx\bar{y}
+ \frac{1}{0}st\bar{x}y
+ \frac{1}{0}st\bar{x}\bar{y} \\
+ \frac{1}{0}s\bar{t}x + s\bar{t}x\bar{y}
+ \frac{1}{0}s\bar{t}\bar{x}y
+ \frac{1}{0}s\bar{t}\bar{x}\bar{y}\\
+ \frac{1}{0}\bar{s}tx + \frac{1}{0}\bar{s}tx\bar{y}
+ \bar{s}t\bar{x}y
+ \frac{1}{0}\bar{s}t\bar{x}\bar{y} \\
+ 0\bar{s}\bar{t}x + 0\bar{s}\bar{t}x\bar{y}
+ 0\bar{s}\bar{t}\bar{x}y
+ 0\bar{s}\bar{t}\bar{x}\bar{y}. \tag{6}
\end{gather*}

この結果から、論理から代数へと移行し、

(XVII.17) の式を得ます：
\begin{gather*}
V = stxy + s\bar{t}x\bar{y} + \bar{s}\bar{t}\bar{x}y + \bar{s}\bar{t}\bar{x}\bar{y}\\
= stxy + s\bar{t}x\bar{y} + \bar{s}\bar{t}.
\end{gather*}

したがって、求める確率 \( u \) を表す式は以下の連立方程式から得られます：
\begin{equation*}
\tag{7}
\begin{split}
\frac{stxy + s\bar{t}x\bar{y} + \bar{s}\bar{t}x}{c_1}
= \frac{stxy + \bar{s}t\bar{x}y + \bar{s}\bar{t}y}{c_2} \\
= \frac{stxy + s\bar{t}x\bar{y}}{c_1 p_1}
= \frac{stxy + \bar{s}t\bar{x}y}{c_2 p_2} \\
= \frac{stxy + s\bar{t}x\bar{y} + \bar{s}\bar{t}x\bar{y}}{u}
= \frac{stxy + s\bar{t}x\bar{y} + \bar{s}\bar{t}x\bar{y} + \bar{s}\bar{t}}{1} = V,
\end{split}
\end{equation*}

ここで \( s, t, x, y, V \) を消去する必要があります。

もし等しい分数の数列があるとすれば、例えば

\[
\frac{a}{a'} = \frac{b}{b'} = \frac{c}{c'} \dotso = \lambda,
\]

であれば、

\[
\frac{la + mb + nc}{la' + mb' + nc'} = \lambda
\]

が成り立ちます。したがって、上記の方程式系から

\[
\frac{\bar{s}t\bar{x}y}{u - c_1 p_1} = \frac{s\bar{t}x\bar{y}}{u - c_2 p_2} = \frac{\bar{s}\bar{t}}{1 - u} = V
\]

が導けます。そして、最初の3つの項の積を最後の項の立方と等しくすると、

\[
\frac{s \bar{s}^2 t \bar{t}^2 x \bar{x} y \bar{y}}{(u - c_1 p_1)(u - c_2 p_2)(1 - u)} = V^3
\]

が得られます。

また、方程式系（7）から、

\[
\frac{\bar{s}\bar{t}\bar{x}}{1 - u - c_1 + c_1 p_1} = \frac{\bar{s}\bar{t}\bar{y}}{1 - u - c_2 + c_2 p_2} = \frac{stxy}{c_1 p_1 + c_2 p_2 - u} = V
\]

が導けます。同様に、

\[
\frac{s \bar{s}^2 t \bar{t}^2 x \bar{x} y \bar{y}}{(1 - c_1 + c_1 p_1 - u)(1 - c_2 + c_2 p_2 - u)(c_1 p_1 + c_2 p_2 - u)} = V^3
\]

が得られます。

（8）と（9）における $V^3$ の値を等しくすると、

\[
(u - c_1 p_1)(u - c_2 p_2)(1 - u) = \{1 - c_1(1 - p_1) - u\}\{1 - c_2(1 - p_2) - u\}(c_1 p_1 + c_2 p_2 - u)
\]

が得られ、これはより便利な形で、

\[
\frac{(u - c_1 p_1)(u - c_2 p_2)}{c_1 p_1 + c_2 p_2 - u} = \frac{\{1 - c_1(1 - p_1) - u\}\{1 - c_2(1 - p_2) - u\}}{1 - u}
\]

と書き換えることができます。

この方程式から $u$ の値を求めることができます。残るは、その根のうちどのものを用いるべきかを決定することです。

3. 以前に示された（XIX.12）ように、求める確率を表すために量 $u$ は各量 $c_1 p_1$, $c_2 p_2$ より大きくなければならず、また各量 $1 - c_1(1 - p_1)$, $1 - c_2(1 - p_2)$, $c_1 p_1 + c_2 p_2$ より小さくなければなりません。さらに、定数間の条件として、最後の3つの量はそれぞれ最初の2つの量より大きくなければならないことが要求されます。この条件が満たされた場合、最終的な方程式（10）は指定された範囲内にただ一つの根を持つことになります。その根は従って求める $u$ の値となります。

まず、下限 $c_1 p_1$, $c_2 p_2$ をそれぞれ $a$, $b$ と表し、上限 $1 - c_1(1 - p_1)$, $1 - c_2(1 - p_2)$, $c_1 p_1 + c_2 p_2$ をそれぞれ $a'$, $b'$, $c'$ と表します。すると一般の方程式は次の形で表現できます。
\begin{gather*}
(u-a)(u-b)(1-u) - (a'-u)(b'-u)(c'-u) = 0 \tag{11}\\
\text{または}\quad (1-a'-b')u^2 - \{ab - a'b' + (1-a'-b')c'\} u + ab - a'b'c' = 0. \tag{11}
\end{gather*}

上記方程式の最初の項を $V$ と表すと、
\begin{equation*}\tag{12}
\frac{d^2 V}{du^2} = 2(1 - a' - b').
\end{equation*}

ここで、$a$ を $u$ の下限の最大値、$a'$ を $u$ の上限の最小値と仮定し、$u = a$ から $u = a'$ までの範囲で $V$ の値を追跡します。

$u = a$ のとき、（11）式の最初の項から $V$ は負であり、$u = a'$ のとき $V$ は正です。これらの限界間では $V$ は連続的に変化し、無限大にはならないことがわかります。さらに $\frac{d^2 V}{du^2}$ は常に同じ符号を持ちます。

したがって、$u$ を $x$ 軸の座標とし、$V$ を $y$ 軸の座標とすると、曲線は $u = a$ に対応する $u$ 軸以下の点から $u = a'$ に対応する $u$ 軸以上の点へと移動します。この曲線は連続であり、常に同じ方向に凸または凹の向きを持ちます。このような状況では、曲線は $u$ 軸を一度だけ交差することになります。

ゆえに、$u = a$ と $u = a'$ の間には、方程式（11）を満たす $u$ の値が一つだけ存在します。さらに、曲線を想像して描くと、他の $u$ の値は以下の条件によって決まります：もし量 $1 - a' - b'$ が正のとき、その他の $u$ の値は $a$ よりも小さいことがわかります。また、$1 - a' - b'$ が負のとき、その他の $u$ の値はいずれかの量 $a', b', c'$ よりも大きいことがわかります。したがって、方程式（11）を解く際には平方根の正の符号を使用する必要があります。これにより、以下の式を得ることができます：
\begin{equation*}\tag{13}
u = \frac{ab - a'b' + (1 - a' - b')c' + \sqrt{Q}}{2(1 - a' - b')}
\end{equation*}
ここで $Q = \{ab - a'b' + (1 - a' - b')c'\}^2 - 4(1 - a' - b')(ab - a'b'c')$。

4. この調査の結果は、ある程度検証することができます。まず、事象 $E$ の確率が一般に、事象 $E$ と原因 $A_1$ や $A_2$ の共起確率よりも大きいことが明らかです。したがって、解を考慮すると以下の不等式が成り立ちます：
\begin{equation*}
u > c_1 p_1, \quad u > c_2 p_2.
\end{equation*}

また、事象 $E$ の確率は一般に、原因 $A_1, A_2$ が互いに排他的である場合の確率よりも小さいことが明らかです。したがって以下の不等式が成り立ちます：
\begin{equation*}
u \leq c_1 p_1 + c_2 p_2.
\end{equation*}

さらに、事象 $E$ の欠如と原因 $A_1$ の存在との共起確率は、事象 $E$ の絶対的欠如確率よりも小さいことが一般に成り立つため、以下の不等式が成り立ちます：
\begin{gather*}
c_1(1 - p_1) \leq 1 - u, \\
\therefore u \leq 1 - c_1(1 - p_1).
\end{gather*}

同様に、
\begin{equation*}
u \leq 1 - c_2(1 - p_2).
\end{equation*}

これらの条件は、一般解の制限条件が確認されました。

また、$p_1 = 1$, $p_2 = 1$ とします。これは、原因 $A_1$ や $A_2$ のいずれかが存在するとき、事象 $E$ が発生することを仮定することを意味します。すると、$a = c_1$, $b = c_2$, $a' = 1$, $b' = 1$, $c' = c_1 + c_2$ であり、式 (13) に代入すると、

\begin{gather*}
u = \frac{c_1 c_2 - c_1 - c_2 - 1
         + \sqrt{(c_1 c_2 - c_1 - c_2 - 1)^2
                 + 4(c_1 c_2 - c_1 - c_2)} }{-2} \\
= c_1 + c_2 - c_1 c_2 \quad (\text{簡略化})\\
= 1 - (1 - c_1)(1 - c_2).
\end{gather*}

この式は、少なくとも一つの原因が存在する確率を表しており、その状況下では明らかに事象 $E$ の確率と一致します。

最後に、$c_1$ および $c_2$ が非常に小さいと仮定し、それらの積を無視できる場合を考えます。すると、$u$ の表式は $c_1 p_1 + c_2 p_2$ に簡略化されます。各原因の確率が小さくなるほど、原因の組合せの確率は著しく小さくなります。したがって、このような簡略化を続けることで、事象 $E$ の確率は互いに排他的な原因として考えたときと同じになります。

この解法について詳しく説明しましたが、これはその後の解法のモデルとなり、それらはより複雑であるため、このような準備なしでは難解に思えるかもしれません。

5. 問題 II．前問で採用した仮定、すなわち両方の原因 $A_1$ と $A_2$ が存在しないとき事象 $E$ は発生しないという仮定を変更し、原因 $A_1$ と $A_2$ は両方とも存在しないことはないとします。その他の状況は前問のとおりとします。このとき、事象 $E$ の確率を求めよ。

ここでは、前回の解法における式(2)に代えて以下の式が用いられます：
\begin{equation*}
(1-x)(1-y) = 0.
\end{equation*}

$z$ の展開された論理表現は以下のように求められます：
%-----------------------File: 130.png----------------------------
\begin{equation*}\begin{split}
& z = stxy + \frac{1}{0} stx\bar{y} + \frac{1}{0}st\bar{x}y + \frac{1}{0} st\bar{x}\bar{y} \\
& + \frac{1}{0}s\bar{t}xy + s\bar{t}x\bar{y} + \frac{1}{0}s\bar{t}\bar{x}y + \frac{1}{0} s\bar{t}\bar{x}\bar{y} \\
& + \frac{1}{0}\bar{s}txy + \frac{1}{0} \bar{s}tx\bar{y} + \bar{s}t\bar{x}y + \frac{1}{0}\bar{s}t\bar{x}\bar{y} \\
& + 0 \bar{s}\bar{t}xy + 0 \bar{s}\bar{t}x\bar{y} + 0\bar{s}\bar{t}\bar{x}y + \frac{1}{0}\bar{s}\bar{t}\bar{x}\bar{y};
\end{split}
\end{equation*}

最終的な解は以下の式によって与えられます：
\begin{equation*}
\operatorname{Prob. } E = u;
\end{equation*}
ここで、$u$ は方程式
\begin{equation*}\tag{1}
\frac{(u-a)(u-b)}{a+b-u}=\frac{(a'-u)(b'-u)}{u-a'-b'+1},
\end{equation*}
の解で、ここで $a=c_1 p_1$, $b=c_2 p_2$, $a'=1 - c_1(1-p_1)$, $b'=1-c_2(1-p_2)$ と定義されます。

制限条件は以下の通りです：$u$ の値はそれぞれの量 $a$, $b$, および $a' + b'-1$ より大きい一方で、それぞれの量 $a'$, $b'$, および $a + b$ より小さい必要があります。

前回の問題の解法と同様に示すことができます：方程式(1)はこれらの条件を満たす唯一の根を持つ二次方程式です。これらの条件自体は、以前と同じルールで導き出されましたが、小限界 $a' + b'-1$ は $1-z$ の大限界を求めることによって得られました。

さらに付け加えると、データ中の定数は、上記の条件を満たすだけでなく、個々に $a'$, $b'$, および $a + b$ より大きいこと、そして $c_1 + c_2 \stackrel{=}{>} 1$ の条件も満たさなければなりません。これはルールの適用からも明らかです。

6. 問題 III．二つの事象 \( A \) と \( B \) の確率はそれぞれ \( a \) および \( b \) であり、事象 \( A \) が発生した場合に伴う事象 \( E \) の確率は \( p \) で、また事象 \( B \) が発生した場合に伴う事象 \( E \) の確率は \( q \) である。\( A \) が発生した場合に \( B \) も発生する確率または逆に \( B \) が発生した場合に \( A \) も発生する確率を求めよ。

事件 \( A \) を \( x \), 事件 \( B \) を \( y \), 事件 \( E \) を \( z \) で表すことにしましょう。すると、与えられたデータは以下のようになります：
\begin{align*}
\operatorname{Prob. } x & = a, & \operatorname{Prob. } y & = b. \\
\operatorname{Prob. } xz & = ap, & \operatorname{Prob. } yz & = bq.
\end{align*}
求めるべき値は以下のようになります：
\[
\frac{\operatorname{Prob. } xy}{\operatorname{Prob. } x} \text{ または } \frac{\operatorname{Prob. } xy}{\operatorname{Prob. } y}.
\]

ここで、
\[
xy = s, \quad yz = t, \quad xy = w.
\]
\( z \) を消去すると、以下のように簡約化できます：
\[
s \bar{x} + t \bar{y} + sy \bar{t} + xt \bar{s} + xy \bar{w} + (1 - xy)w = 0,
\]
従って
\[
w = \frac{s \bar{x} + t \bar{y} + sy \bar{t} + xt \bar{s} + xy}{2xy - 1}.
\]

以上を論理から代数へ移すと、以下のように表されます：
\[
\operatorname{Prob. } xy = \frac{xyst + xy \bar{s} \bar{t}}{V},
\]
ここで \( x, y, s, t \) は以下の連立方程式によって決定されます：
\begin{gather*}
\frac{xyst + x \bar{y} s \bar{t} + xy \bar{s} \bar{t} + x \bar{y} \bar{s} \bar{t}}{a} = \frac{xyst + \bar{x} y \bar{s} t + xy \bar{s} \bar{t} + \bar{x} y \bar{s} \bar{t}}{b} \\
= \frac{xyst + x \bar{y} s \bar{t}}{ap} = \frac{xyst + \bar{x} y \bar{s} t}{bq} \\
= xyst + x \bar{y} s \bar{t} + \bar{x} y \bar{s} t + xy \bar{s} \bar{t} + x \bar{y} \bar{s} \bar{t} + \bar{x} y \bar{s} \bar{t} + \bar{x} \bar{y} \bar{s} \bar{t} = V.
\end{gather*}

上記の連立方程式をより便利な形に簡約化するために、各項を \( \bar{x} \bar{y} \bar{s} \bar{t} \) で割り、結果として以下のように定義します：
\[
\frac{xs}{\bar{x} \bar{s}} = m, \quad \frac{yt}{\bar{y} \bar{t}} = m', \quad \frac{x}{\bar{x}} = n, \quad \frac{y}{\bar{y}} = n'.
\]

これにより、以下のように表されます：
\begin{gather*}
\frac{mm' + m + nn' + n}{a} = \frac{mm' + m' + nn' + n'}{b} \\
= \frac{mm' + m}{ap} = \frac{mm' + m'}{bq} \\
= mm' + m + m' + nn' + n + n' + 1.
\end{gather*}

また、
\[
\operatorname{Prob. } xy = \frac{mm' + nn'}{mm' + m + m' + nn' + n + n' + 1}.
\]

これらの式は以下のように簡約化できます：
\begin{gather*}
\frac{mm' + m}{ap} = \frac{mm' + m'}{bq} \\
= \frac{nn' + n}{a(1 - p)} = \frac{nn' + n'}{b(1 - q)} \\
= (m + 1)(m' + 1) + (n + 1)(n' + 1) - 1. \\
\operatorname{Prob. } xy = \frac{mm' + nn'}{(m + 1)(m' + 1) + (n + 1)(n' + 1) - 1}.
\end{gather*}

ここで、
\[
(m + 1)(m' + 1) = \frac{\mu}{\nu + \mu - 1}, \quad (n + 1)(n' + 1) = \frac{\nu}{\nu + \mu - 1}.
\]

すると、以下のように表されます：
\[
mm' + m = \frac{m(m' + 1)(m + 1)}{m + 1} = \frac{m \mu}{(m + 1)(\nu + \mu - 1)},
\]
他の項についても同様に計算します。これらを代入し、各項を \( \nu + \mu - 1 \) で掛けると以下のようになります：
\begin{gather*}
\frac{m \mu}{(m + 1)ap} = \frac{m' \mu}{(m' + 1)bq} \\
= \frac{n \nu}{(n + 1)a(1 - p)} = \frac{n' \nu}{(n' + 1)b(1 - q)} = 1.
\\
\operatorname{Prob. } xy = (mm' + nn')(\nu + \mu - 1).
\end{gather*}

上記の連立方程式から以下のように表されます：
\begin{gather*}
\frac{m}{m + 1} = \frac{ap}{\mu}, \text{ 従って } m = \frac{ap}{\mu - ap}.
\end{gather*}

同様に、
\[
m' = \frac{bq}{\mu - bq}, \quad n = \frac{a(1 - p)}{\nu - a(1 - p)}, \quad n' = \frac{b(1 - q)}{\nu - b(1 - q)}.
\]

従って、
\[
m + 1 = \frac{\mu}{\mu - ap}, \quad n + 1 = \frac{\nu}{\nu - a(1 - p)}, \text{ 他も同様}
\]

これらの値を (2) の簡約形に代入すると、
\[
\nu + \mu - 1 = \frac{\mu}{(m + 1)(m' + 1)} = \frac{\nu}{(n + 1)(n' + 1)},
\]
以下のようになります：
\[
\nu + \mu - 1 = \frac{(\mu - ap)(\mu - bq)}{\mu} = \frac{(\nu - a(1 - p))(\nu - b(1 - q))}{\nu}.
\]

また、 (3) の値を代入すると以下のようになります：
\[
\operatorname{Prob. } xy = \left[ \frac{abpq}{(\mu - ap)(\mu - bq)} + \frac{ab(1 - p)(1 - q)}{\{\nu - a(1 - p)\} \{\nu - b(1 - q)\}} \right] (\nu + \mu - 1).
\]

ここで、
\[
\nu + \mu - 1 = \mu - ap - bq + \frac{apbq}{\mu}.
\]
従って、
\[
\frac{apbq}{\mu} = \nu - 1 + ap + bq.
\]

同様に、
\[
\frac{ab(1 - p)(1 - q)}{\nu} = \mu - 1 + a(1 - p) + b(1 - q).
\]

これらを加えると、最初の項が Prob. \( xy \) の表現と一致するので、以下のようになります：
\[
\operatorname{Prob. } xy = \nu + \mu + a + b - 2.
\]

Prob. \( xy \) を \( u \) で表し、\( a + b - 2 = m \) とすると、
\[
\mu + \nu = u - m.
\]

また、 (5) より以下のようになります：
\[
\mu \nu = abpq - (ap + bq - 1)\mu.
\]

同様に、 (4) の最初と第三の項を比較すると以下のようになります：
\[
\mu \nu = ab(1 - p)(1 - q) - \{a(1 - p) + b(1 - q) - 1\}\nu.
\]

ここで \( ap + bq - 1 \) を \( h \), \( a(1 - p) + b(1 - q) - 1 \) を \( h' \) と表すと、以下のようになります：
\[
h \mu - h' \nu = ab(pq + (1 - p)(1 - q)) = ab(p + q - 1).
\]

ここで \( ab(p + q - 1) = l \) とすると、
\[
h \mu - h' \nu = l.
\]

(6) と (8) より以下のようになります：
\[
\mu = \frac{h'(u - m) + l}{m}, \quad \nu = \frac{h(u - m) - l}{m}.
\]

これらを (7) の簡約形に代入すると、
\[
\mu(\nu + h) = abpq,
\]
以下のようになります：
\[
(hu - l)\{h'(u - m) + l\} = abpqm^2.
\]

これは二次方程式で、これを解くことで Prob. \( xy \) の値を求めることができます。

解は次のように表すことができます：
\[
h = \frac{lh' + h(h'm - l) \pm \sqrt{\{lh' - h(h'm - l)\}^2 + 4hh'abpqm^2}}{2hh'}.
\]
しかし、さらに
\[
lh' - h(h'm - l) = l(h + h') - hh'm = (1 - hh')m,
\]
が成り立つことに注意すると、
\[
h = ap + bq - 1, \quad h' = a(1 - p) + b(1 - q) - 1,
\]
より
\[
h + h' = a + b - 2 = m,
\]
となるため、
\[
\operatorname{Prob.} xy = \frac{lh' + h(h'm - l) \pm m \sqrt{\{(l - hh')^2 + 4hh'abpq\}}}{2hh'}.
\]
ここで、平方根の符号はどのようになるべきかを決定する必要があります。一般的な方法でこれを確認することは可能ですが、既知のケースとの比較によって符号が決定されることが容易であり、現在の場合には十分です。例えば、事象 $A$ が常に起こり、事象 $B$ が決して起こらないという確実な状況を考えると、事象 $A$ と $B$ は決して共に起こることはありません。したがって、$p = 1$, $q = 0$ のとき、$u = 0$ であるべきです。

$p = 1$, $q = 0$ の仮定により
\[
h = a - 1, \quad h' = b - 1, \quad l = 0, \quad m = a + b - 2,
\]
が得られます。式 (10)に代入すると、
\[
\operatorname{Prob.} xy = \frac{(a-1)(b-1)(a+b-2) \pm (a+b-2)(a-1)(b-1)}{2(a-1)(b-1)},
\]
となり、この式は符号が下のものであるとき零になります。したがって、一般問題の最終的な解は
\[
\frac{\operatorname{Prob.} xy}{\operatorname{Prob.} x}
= \frac{lh' + h(h'm - l) - m \sqrt{\{(l - hh')^2 + 4hh'abpq\}}}{2ahh'},
\]
と表されます。ここで、
\[
h = ap + bq - 1, \quad h' = a(1-p) + b(1-q) - 1,
\]
\[
l = ab(p + q - 1), \quad m = a + b - 2.
\]

最終的な論理解法に影響を与える項は最初の問題と同様であり、定数間の条件も同じです：
\[
ap < 1 - b(1-q), \quad bq < 1 - a(1-p).
\]

7. 上記の解が正しくなることを確認するには、得られた式が $p$, $q$ と $1 - p$, $1 - q$ の二つの集合に関して対称であることが必要です。つまり、$p$ を $1 - p$ に、$q$ を $1 - q$ に置き換えると式は変化しないことから明らかです。これは最終結果を導く際に用いた方程式
\[
\operatorname{Prob.} xy = ab \left\{ \frac{pq}{\mu} + \frac{(1-p)(1-q)}{\nu} \right\}
\]
からも明らかです。もし事象 $E$ の確率 $p$, $q$ が、事象 $A$ と $B$ の発生状況に基づいて存在するなら、その否定事象の確率 $1 - p$, $1 - q$ も同じ条件のもとで存在します。したがって、データは $E$ の発生または非発生に関係なく変化しないため、解は $p$, $q$ および $1 - p$, $1 - q$ の対称的な関数であるべきです。

$p = 1$, $q = 1$ の特別な場合を検討しましょう。このとき、以下の関係が得られます：
\begin{equation*}
h = a + b - 1, \quad h' = -1, \quad l = ab, \quad m = a + b - 2,
\end{equation*}

そして、次のように置き換えて計算します：
\begin{gather*}
\frac{\operatorname{Prob. } xy}{\operatorname{Prob. } x} 
= \frac{-ab + (a+b-1)(2-a-b-ab)-(a+b-2)(ab-a-b+1)}{-2a(a+b-1)} \\
= \frac{-2ab(a+b-1)}{-2a(a+b-1)} = b.
\end{gather*}

したがって、この場合、事象 $A$ と $B$ は互いに実質的に独立であることがわかります。他の事象 $E$ との恒常的な関連性を仮定し、その発生頻度について何も知るものがなければ、事象 $A$ と $B$ 自体の相互依存性は確立されません。私はこの結論が理にかなっていると考えますが、特定の例を見ると最初の印象では異なる結果を示すように思えるかもしれません。

たとえば、ある種類の雑草が岸辺で見つかる確率と、特定の種類の動物植物が海辺で見つかる確率が別々に決定され、これらの事象はともに嵐によって引き起こされる波動中にしか発生しないことが確認された場合、私はこれら的事象が独立であるとは結論付けるべきだと思います。つまり、特定種類の雑草を拾うと、その動物植物を見つける確率がより高くなると考えられるでしょう。しかし、この事実は、実際の問題の条件に含まれていない別の状況に対する我々の知識によるものだと考えます。すなわち、嵐は偶然的な現象であるということです。

観察範囲を常に嵐で荒れる海に限定すると、私は雑草と動物植物が岸辺で見つかる事象が独立であるべきであると考えるでしょう。

より一般的には、すべての自然現象には共通する条件があります。これらの条件は互いの独立性に影響を与えないことが感じられます。したがって、我々の記録された経験を通じて常に存在した特定の条件は、その経験が及ぶ範囲内の現象に対する上記の特性を持つと考えることができます。

8. 問題 IV。---上記の観察を多少示すために、最後の問題に追加されたデータとともに、事象 $E$ の絶対的な確率が与えられた場合を考えます。完成したデータの系は以下のようになります：
\begin{gather*}
\operatorname{Prob. } x = a, \quad \operatorname{Prob. } y = b, \quad \operatorname{Prob. } z = c, \\
\operatorname{Prob. } xz = ap, \quad \operatorname{Prob. } yz = bq,
\end{gather*}
そして、$\operatorname{Prob. } xy$ を求めることを求められます。

$xz = s$, $yz = t$, $xy = w$ と仮定すると、最終的な論理方程式は以下のようになります：
\begin{equation*}
\begin{split}
w &= xystz + xy\bar{s}\bar{t}\bar{z} \\
& \quad + 0(x\bar{y}s\bar{t}z + \bar{x}y\bar{s}tz
+ x\bar{y}\bar{z}\bar{s}\bar{t} + xy\bar{z}\bar{s}\bar{t}
\\ & \quad
+ \bar{x}\bar{y}z\bar{s}\bar{t} + \bar{x}\bar{y}\bar{z}\bar{s}\bar{t})
\\
&+ \text{項の係数が } \frac{1}{0} \text{ の項}  %\tag{1}
\end{split}
\end{equation*}

代数系を形成した後、前の問題で採用された変換により、その後の消去は簡略化されます。最終的な結果は以下のようになります：
\begin{equation*}
\operatorname{Prob. } xy = ab \left\{ \frac{pq}{c} + \frac{(1-p)(1-q)}{1-c} \right\}.  \tag{2}
\end{equation*}

定数間の条件は以下のようになります：
\begin{equation*}
c \stackrel{=}{>} ap, \quad
c \stackrel{=}{>} bq, \quad
c \stackrel{=}{<} 1 - a(1 - p), \quad
c \stackrel{=}{<} 1-b(1-q).
\end{equation*}

もし $p = 1$, $q = 1$ の場合、次のように得られます：
\begin{equation*}
\operatorname{Prob.} xy = \frac{ab}{c},
\end{equation*}
ここで $c$ は $a$ や $b$ よりも小さい値を取ることはできません。したがって、事象 $E$ が偶発的なものであることが分かっている場合、その事象 $x$ と $y$ の恒常的な共存は、その頻度の逆数比例で $xy$ の確率を増加させます。

公式 (2) は多くの例で検証できます。特に、$q = c$ とする場合、次のように得られます：
\begin{equation*}
\tag{3} \operatorname{Prob.} xy = ab.
\end{equation*}

この仮定 $q = c$ は、定義（第16章）により事象 $B$ と $E$ の独立性を含みます。したがって、$B$ と $E$ が独立である場合、$A$ と $E$ の間に存在するいかなる関係も、$A$ と $B$ の間の関係を確立することはできません。従って、上記の式 (3) より $A$ と $B$ も独立であることが示されます。

(2) 式から容易に確認できますが、$\operatorname{Prob.} z$ の値は $\operatorname{Prob.} xy$ を最小にするとき、次のように与えられます：
\begin{equation*}
\operatorname{Prob.} z = \frac{\sqrt{pq}}{\sqrt{pq} + \sqrt{(1-p)(1-q)}}.
\end{equation*}

もし $p = q$ の場合、これは
\begin{equation*}
\operatorname{Prob.} z = p
\end{equation*}
を与えます。この結果は、(3) 式に対する同じ論理から正しさを確認できます。

問題 V．任意の三つの事象の確率とそれらの連続確率が与えられたとき、任意の二つの事象の連続確率を求めよ。

与えられたデータは以下の通りです：
\begin{equation*}
\operatorname{Prob.} x = p, \quad
\operatorname{Prob.} y = q, \quad
\operatorname{Prob.} z = r, \quad
\operatorname{Prob.} xyz = m,
\end{equation*}

求めるものは $\operatorname{Prob.} xy$ です。

$xyz = s$, $xy = t$ を仮定すると、最終的な論理方程式は次のようになります。
\[
t = xyzs + xy\bar{z}\bar{s} + 0(x\bar{y}\bar{s} + \bar{x}\bar{s}) + \frac{1}{0}(\text{他のすべての成分の総和});
\]
これより、最終的に
\[
\operatorname{Prob. } xy = \frac{H - \sqrt{(H^2 - 4pq\bar{r}^2 - 4\bar{p}\bar{q}\bar{r}m)}}{2\bar{r}},
\]
ここで
\[
\bar{p} = 1 - p, \text{ 他も同様}. \quad H = \bar{p}\bar{q} + (p + q)\bar{r}.
\]
この式は、$p = 1$, $q = 1$, $r = 0$ のとき、またはしたがって $m = 0$ のときなどに検証可能である。

条件 $\operatorname{Prob. } z = r$ を省略しても、解は確定的です。その場合、
\[
\text{Prob. } xy = \frac{pq(1 - m) + (1 - p)(1 - q)m}{1 - m};
\]
となります。また、上記の結果が $m = pqr$ のとき一致することを最後に確認できます。

9. 以下は問題Iの一般化であり、その解法はより複雑になりますが、類似の分析により得られます。

\textsc{問題VI}---ある事象が特定の原因 $A_1, A_2, \dotsc A_n$ のいずれかまたはそれらの組み合わせによってのみ起こる場合、そして一般に $c_i$ は原因 $A_i$ の確率を表し、$p_i$ は原因 $A_i$ が存在するとき事象 $E$ が発生する確率を表すとします。このとき、与えられた $c_i$ および $p_i$ の系列の値から、事象 $E$ の確率を求めることを求められます。\footnote{
上記の問題は、著者が1851年11月のケンブリッジとダブリン数学雑誌に提出し、以下の観察を添えて提示しました。}

---

多くの考慮の結果、私はこの問題に関して、現在の慣習とは異なる道を選択し、その理由は次の通りです：まず第一に、私はこの問いを提出することで、既存の方法が十分であるかどうかをテストすることを提案します。第二に、この問いの議論は、純粋解析の重要な分野に関する知識を多少なりとも増加させることを予想しています。しかし、私があえてこの理由のみに基づいて弁明したいのは、前者だけです。

私は、以前の研究から逸脱することなく、この問いが彼らの注意を引くものであると信じる人々を見つけることを願っています。しかしながら、個人的な技能や知識を試すためのものではなく、公的かつ科学的な目的のために提出されたものとしてのみ見なされるべきだと断言します。

著者は、上記の問題の発表がいくつかの興味深い私的な通信につながったことを追加的に述べるべきであると信じています。しかしながら、その結果、解が得られなかったことは否めません。

原因 $A_1, A_2, \dotsc, A_n$ を $x_1, x_2, \dotsc, x_n$ で表し、事件 $E$ を $z$ で表します。すると、一般に、

\[
\text{Prob. }x_i = c_i,\qquad \text{Prob. }x_i z = c_i p_i.
\]

さらに、$E$ は原因 $A_1, A_2, \dotsc, A_n$ のいずれかまたはそれらの組み合わせと関連してのみ発生するという条件は、

\[
z(1-x_1)(1-x_2) \dotsm (1-x_n) = 0.   \tag{1}
\]

という論理的な条件を確立します。

---

現在、一般に次のように仮定します：
\[
x_i z = t_i,
\]
これは以下のような形になります：
\[
x_i z(1 - t_i) + t_i (1 - x_i z) = 0.
\]
この式は、問題の論理条件を表す $n$ 個の方程式のタイプであり、これらは（1）と合わせて使用されます。これらの方程式全てを足し合わせると、以前の簡略化により許されるように、次のようになります：
\[
\sum \{ x_i z(1 - t_i) + t_i (1 - x_i z) \} + z(1 - x_1)(1 - x_2) \cdots (1 - x_n) = 0,   \tag{2}
\]
ここで、$\sum$ の和は $i = 1$ から $i = n$ までとされています。この単純で十分な論理方程式と、一般的な方程式
\[
\text{Prob. } x_i = c_i, \quad \text{Prob. } t_i = c_i p_i   \tag{3}
\]
によって表される $2n$ のデータと共に、Prob. $z$ を決定するための要素が与えられます。

次に（2）を $z$ について展開します。すると次のようになります：
\begin{multline*}
\left[ \sum \{ x_i (1 - t_i) + t_i (1 - x_i) \} + (1 - x_1)(1 - x_2) \cdots (1 - x_n) \right] z \\
+ \sum t_i(1 - z) = 0,
\end{multline*}
従って
\[
z = \frac{\sum t_i}{\sum t_i - \sum \{ x_i (1 - t_i) + t_i (1 - x_i) \} - (1 - x_1)(1 - x_2) \cdots (1 - x_n)}.   \tag{4}
\]
上記方程式の第二項の展開における任意の成分は、$2n$ 個の因子から成り立っています。そのうち $n$ つは $x_1, x_2, \dotsc, x_n, 1 - x_1, 1 - x_2, \dotsc, 1 - x_n$ の集合から、そして残りの $n$ つは $t_1, t_2, \dotsc, t_n, 1 - t_1, 1 - t_2, \dotsc, 1 - t_n$ の集合から選ばれます。ただし、$x_1(1 - x_1), t_1(1 - t_1)$ シンボルの組み合わせは許可されません。次に、$(1 - t_1), (1 - t_2) \dotsc (1 - t_n)$ が $t$-因子である、すなわち $t_1, \dotsc, 1 - t_1$ の集合から得られる因子について考えます。

どのような成分の係数も、（4）式の第二項においてそれぞれ$t_1, t_2, \dots, t_n$を$0$に置き換えて求めることができます。その後、$x_1, x_2, \dots, x_n$の値は成分の$x$因子の性質によって依存しているので、それらの値を与えます。

今、$t_1, t_2, \dots, t_n$をそれぞれ$0$に置き換えると、第二項は次のようになります：
\[
  \frac{0}{ -\sum x_i - (1-x_1)(1-x_2) \dotsm (1-x_n) }.
\]
%-----------------------File: 141.png----------------------------
この式は、その後に$x_1, x_2, \dots, x_n$に対して与えるどの値（0または1）でも消えます。なぜなら、これらの値がすべて0でない場合、項$\sum x_i$は消えないし、すべて0の場合、項$-(1-x_1) \dotsm (1-x_n)$は$-1$となり、したがって分母は常に0にならず、分数も0になります。

従って、$(1-t_1) \dotsm (1-t_n)$の因子であるすべての成分の係数は0であり、全ての可能な$x$成分の合計が1なので、$z$の展開には総和項$0(1-t_1)\dotsm (1-t_n)$が含まれます。

次に、$t$-因子が $t_1\; t_2 \dotsc t_r(1-t_{r+1}) \dotsc (1-t_n)$ である任意の成分を考える。ここで $r$ は 1 以上であり、等しくまたはそれ以上である。式（4）の第二項において、$t_1 = 1,\; t_2 = 1, \dotsc, t_r = 1, t_{r+1} = 0, \dotsc, t_n = 0$ と置き換えると、以下の表現を得る：
\[
  \frac{r}{x_1 + x_2 + \dotsb + x_r - x_{r+1} - \dotsb - x_n
           - (1-x_1)(1-x_2) \dotsc (1-x_n)}
\]
ここで、シンボルの許容される唯一の値は 0 と 1 のみであるから、上記の表現は $x_1 = 1,\; x_2 = 1, \dotsc, x_r = 1,\; x_{r+1} = 0,\; \dotsc,\; x_n = 0$ のとき 1 となることが明らかであり、他のすべての値の組合せではその表現は 1 より大きい値をとることになる。したがって、最終展開における $x$-因子が $t$-因子に類似する形
\[
  x_1 \dotsc x_r(1-x_{r+1}) \dotsc (1-x_n)\, t_1 \dotsc t_r(1-t_{r+1}) \dotsc (1-t_n)
\]
であるすべての成分に対して、係数 1 が適用されることになる。一方で、現在の場合に含まれる他の成分は仮想的な係数 $\frac{1}{0}$ を持つことになる。

また、この論理は個々のシンボルの特定の配置や順序に依存しないことも明白である。

したがって、$z$ の完全な展開は以下の形になります：
\begin{multline*}
  z = \sum(XT) + 0(1-t_1)(1-t_2)\dotsc (1-t_n)   \\
      + \text{成分のうち係数が }\frac{1}{0} のもの，
  \tag{5}
\end{multline*}
ここで $T$ は $(1-t_1)\dotsc (1-t_n)$ を除く任意の $t$-成分を表し、$X$ は $x_1\dotsc x_n$ の対応するまたは類似の成分を表します。
%-----------------------File: 142.png----------------------------
例えば、$n = 2$ の場合、以下のようになります：
\begin{equation*}
\sum (XT) = x_1 x_2 t_1 t_2 + x_1 \bar{x}_2 t_1 \bar{t}_2 + \bar{x}_1 x_2 \bar{t}_1 t_2,
\end{equation*}
ここで $\bar{x}_1$, $\bar{x}_2$ は $1 - x_1$, $1 - x_2$ を表します。従って
\begin{gather*} \tag{6}
\begin{split}
z = x_1 x_2 t_1 t_2
+ x_1 \bar{x}_2 t_1 \bar{t}_2
+ \bar{x}_1 x_2 \bar{t}_1 t_2 \\
+ 0(x_1 x_2 \bar{t}_1 \bar{t}_2
+ x_1 \bar{x}_2 \bar{t}_1 \bar{t}_2
+ \bar{x}_1 x_2 \bar{t}_1 \bar{t}_2
+ \bar{x}_1 \bar{x}_2 \bar{t}_1 \bar{t}_2) \\
+ \text{成分のうち係数が }\frac{1}{0} のもの.
\end{split}
\end{gather*}
この結果は、記号の違いを考慮すると、本章の問題Iにおいて $z$ の開発された形と一致し、これは明らかにすべき条件です。

10. 簡潔性を保つため、上記の方程式(6)から、$n = 2$ の特定の場合における $\operatorname{Prob. } z$ の決定のために必要な条件を求めることを計画します。そして、同じ調査の線に従って考えを進めることにより、任意の整数値を持つ $n$ に対するより一般的な場合における対応する条件を指定することを目指します。

次に、$n = 2$ の場合を考えると、

\begin{gather*}
\begin{split}
V = x_1 x_2 t_1 t_2
+ x_1 \bar{x}_2 t_1 \bar{t}_2
+ \bar{x}_1 x_2 \bar{t}_1 t_2
+ x_1 x_2 \bar{t}_1 \bar{t}_2
+ x_1 \bar{x}_2 \bar{t}_1 \bar{t}_2
\\
+ \bar{x}_1 x_2 \bar{t}_1 \bar{t}_2
+ \bar{x}_1 \bar{x}_2 \bar{t}_1 \bar{t}_2.
\end{split}
\\
\operatorname{Prob. } z = \frac{x_1 x_2 t_1 t_2
+ x_1 \bar{x}_2 t_1 \bar{t}_2
+ \bar{x}_1 x_2 \bar{t}_1 t_2}{V},
\end{gather*}

$x_1$, $t_1$ などの値の決定条件は、

\[
\begin{split}
\frac{x_1 x_2 t_1 t_2
+ x_1 \bar{x}_2 t_1 \bar{t}_2
+ x_1 x_2 \bar{t}_1 \bar{t}_2
+ x_1 \bar{x}_2 \bar{t}_1 \bar{t}_2}{c_1}
\\
=\frac{x_1 x_2 t_1 t_2
+ \bar{x}_1 x_2 \bar{t}_1 t_2
+ x_1 x_2 \bar{t}_1 \bar{t}_2
+ \bar{x}_1 x_2 \bar{t}_1 \bar{t}_2}{c_2}
\\
=\frac{x_1 x_2 t_1 t_2 + x_1 \bar{x}_2 t_1 \bar{t}_2}{c_1 p_1}
=\frac{x_1 x_2 t_1 t_2 + \bar{x}_1 x_2 \bar{t}_1 t_2}{c_2 p_2} = V
\end{split}
\]

を満たす。

これらの方程式の各項を $\bar{x}_1 \bar{x}_2 \bar{t}_1 \bar{t}_2$ で割り、$\operatorname{Prob. } z$ の分子と分母も同じ量で割って、

\begin{equation*}\tag{7}
\frac{x_1 t_1}{\bar{x}_1 \bar{t}_1} = m_1, \quad
\frac{x_2 t_2}{\bar{x}_2 \bar{t}_2} = m_2, \quad
\frac{x_1}{\bar{x}_1}=n_1, \quad \frac{x_2}{\bar{x}_2}=n_2;
\end{equation*}

$$
\begin{gathered}
zの確率は \frac{m_1 m_2 + m_1 + m_2}{m_1 m_2 + m_1 + m_2 + n_1 n_2 + n_1 + n_2 + 1} である。また、
\frac{m_1 m_2 + m_1 + n_1 n_2 + n_1}{c_1} = \frac{m_1 m_2 + m_2 + n_1 n_2 + n_2}{c_2}
= \frac{m_1 m_2 + m_1}{c_1 p_1} + \frac{m_1 m_2 + m_2}{c_2 p_2}
= m_1 m_2 + m_1 + m_2 + n_1 n_2 + n_1 + n_2 + 1,
\end{gathered}
$$
したがって、$m_1+1$と$m_2+1$を$M$、$n_1+1$と$n_2+1$を$N$と仮定すると、

\begin{equation*}
(m_1 + 1) (m_2+1)=M, \quad (n_1 + 1) (n_2 + 1) = N,
\end{equation*}

$$
\begin{gathered}
zの確率は \frac{M-1}{M+N-1} である。また、
\frac{n_1 (n_2 + 1)}{c_1 (1-p_1)} = \frac{n_2 (n_1 + 1)}{c_2 (1-p_2)} = \frac{m_1 (m_2 + 1)}{c_1 p_1} = \frac{m_2 (m_1 + 1)}{c_2 p_2}
= M+N-1;
\end{gathered}
$$
または、

\[
\begin{split}
\frac{m_1 M}{(m_1 + 1)c_1 p_1}=\frac{m_2 M}{(m_2 + 1)c_2 p_2} =\frac{n_1 N}{(n_1 + 1)c_1(1-p_1)}\\
=\frac{n_2 N}{(n_2 + 1)c_2(1-p_2)}=M+N-1.
\end{split}
\]

今、最終的な論理方程式（5）に対して類似の変換と簡略化を精神的に行います。これにより$z$の確率$\operatorname{Prob. } z$を求めます。

$$
\operatorname{Prob. } z = \frac{M-1}{M+N-1}
$$

ここで、

\[
\begin{gathered}
M=(m_1 + 1) (m_2 + 1)\dotsc (m_n + 1), \\
N = (n_1 + 1) (n_2 + 1)\dotsc (n_n + 1),
\end{gathered}
\]

$m_1, \dotsc, m_n, n_1,\dotsc,n_n$は方程式系（11）によって与えられます。

\[
\begin{split}
\frac{m_1 M}{(m_1 + 1) c_1 p_1} = \frac{m_2 M}{(m_2 + 1) c_2 p_2} \dotso = \frac{m_n M}{(m_n + 1) c_n p_n}
= \frac{n_1 N}{(n_1 + 1)c_1 (1-p_1)}\dotso
= \frac{n_n N}{(n_n+1)c_n(1-p_n)}
=M+N-1.
\end{split}
\]

さらに結果を簡略化するため、次のように仮定します。
\[
\frac{M+N-1}{M} = \frac{1}{\mu}, \quad \frac{M+N-1}{N} = \frac{1}{\nu};
\]
これより
\[
M = \frac{\mu}{\mu + \nu - 1}, \quad N = \frac{\nu}{\mu + \nu - 1}.
\]

次に、以下の式を得ます。
\[
\frac{m_1}{(m_1+1)c_1 p_1} = \frac{m_2}{(m_2+1)c_2 p_2} = \cdots = \frac{m_n}{(m_n+1)c_n p_n} = \frac{1}{\mu},
\]
\[
\frac{n_1}{(n_1+1)c_1(1-p_1)} = \frac{n_2}{(n_2+1)c_2(1-p_2)} = \cdots = \frac{n_n}{(n_n+1)c_n(1-p_n)} = \frac{1}{\nu};
\]
これより
\[
m_1 = \frac{c_1 p_1}{\mu - c_1 p_1}, \quad \dotsc, \quad m_n = \frac{c_n p_n}{\mu - c_n p_n};
\]
さらに
\[
m_1 + 1 = \frac{\mu}{\mu - c_1 p_1}, \quad \dotsc, \quad m_n + 1 = \frac{\mu}{\mu - c_n p_n},
\]
\[
n_1 + 1 = \frac{\nu}{\nu - c_1(1-p_1)}, \quad \dotsc, \quad n_n + 1 = \frac{\nu}{\nu - c_n(1-p_n)}.
\]

これらの値を式 (9) の $M$ と $N$ に代入すると、以下の式を得ます。
\[
\frac{\mu^n}{(\mu - c_1 p_1)(\mu - c_2 p_2) \cdots (\mu - c_n p_n)} = \frac{\mu}{\mu + \nu - 1},
\]
\[
\frac{\nu^n}{(\nu - c_1(1-p_1))(\nu - c_2(1-p_2)) \cdots (\nu - c_n(1-p_n))} = \frac{\nu}{\mu + \nu - 1}.
\]

これらは対称的な形に簡略化できます。
\[
\mu + \nu - 1 = \frac{(\mu - c_1 p_1) \cdots (\mu - c_n p_n)}{\mu^{n-1}} = \frac{(\nu - c_1(1-p_1)) \cdots (\nu - c_n(1-p_n))}{\nu^{n-1}},
\]
最終的に
\[
\operatorname{Prob.} z = \frac{M-1}{M+N-1} = 1 - \nu.
\]

ここで $1 - \nu = u$ と仮定すると、以下の式を得ます。
\[
\mu - u = \frac{(\mu - c_1 p_1) \cdots (\mu - c_n p_n)}{\mu^{n-1}} = \frac{(1 - c_1(1-p_1)) \cdots (1 - c_n(1-p_n) - u)}{(1-u)^{n-1}}.
\]

さらに、以下のように定義します。
\[
c_1 p_1 = a_1, \quad c_n p_n = a_n, \quad 1 - c_1(1-p_1) = b_1, \text{ etc.}
\]
これらの式は以下の形で表すことができます。
\[
\mu - u = \frac{(\mu - a_1) \cdots (\mu - a_n)}{\mu^{n-1}},
\]
ここで
\[
\mu = u + \frac{(b_1 - u) \cdots (b_n - u)}{(1-u)^{n-1}}.
\]

この$\mu$の値を(14)に代入すると、$u$のみを含む方程式が得られ、その解は$\operatorname{Prob. } z$を決定します。なぜなら、(13)より$\operatorname{Prob. } z = u$だからです。残るは$u$の範囲を定めるだけです。

11. 今、$n = 2$の特別な場合で限界が決定された同じ分析法により、現在の場合には以下の結果に到達します。量$u$は$\operatorname{Prob. } z$を表すために、その下限は$a_1, a_2, \dotsc a_n$であり、上限は$b_1, b_2,\dotsc b_n$, $a_1 + a_2 \dotsc + a_n$です。したがって、方程式(14)を満たすこれらの条件を満たす根が常に存在し、かつただ一つであることが推測できます。実用的な検証のために、方程式(14)の根が限界$a_1, a_2, \dotsc a_n$と$b_1, b_2,\dotsc b_n$の間でただ一つ存在することを示すだけで十分です。

まず、（15）において$\mu$が変化する性質について考えます。ここでは、$u$が最小の上限である$a_1$から最大の下限である$b_1$まで変化するとき、$\mu$にどのような影響を与えるかを考察します。$u = a_1$のとき、明らかに$\mu$は正で$a_1$より大きいことがわかります。また、$u = b_1$のとき、$\mu = b_1$であり、これは正です。$a_1 \leq u \leq b_1$の範囲では、$\mu$が$u$と共に増加することが示されます。したがって、

\[
\frac{d\mu}{du} = 1 - \frac{(b_2 - u)(b_3 - u) \dotsm (b_n - u)}{(1 - u)^{n-1}} - \frac{(b_1 - u)(b_3 - u) \dotsm (b_n - u)}{(1 - u)^{n-1}} \dotsm + (n-1)\frac{(b_1 - u)(b_2 - u) \dotsm (b_n - u)}{(1 - u)^n}
\]

となります。

ここで、

\[
\frac{b_1 - u}{1 - u} = x_1, \quad \frac{b_2 - u}{1 - u} = x_2, \quad \dotsm, \quad \frac{b_n - u}{1 - u} = x_n
\]

と置きます。明らかに$x_1, x_2, \dotsc, x_n$は真分数であり、

\[
\begin{split}
\frac{d\mu}{du} &= 1 - x_2 x_3 \dotsm x_n - x_1 x_3 \dotsm x_n \dotsm - x_1 x_2 \dotsm x_{n-1} + (n-1)x_1 x_2 \dotsm x_n \\
&= 1 - (1 - x_1) x_2 x_3 \dotsm x_n - x_1 (1 - x_2) x_3 \dotsm x_n \dotsm \\
&\quad - x_1 x_2 \dotsm x_{n-1} (1 - x_n) - x_1 x_2 \dotsm x_n
\end{split}
\]

となります。

二項の負の項は（論理的開発の言葉を借りれば）$x_1, x_2, \dotsc, x_n$からなる\emph{成分}であり、これらの和は従って1を超えることはありません。したがって$\displaystyle \frac{d\mu}{du}$は正であり、指定された範囲の$u$において$\mu$は増加します。

現在、式(14)を次のように書き換える：

\begin{equation}\tag{17}
\frac{(\mu-a_1)\cdots(\mu-a_n)}{\mu^{n-1}} - (\mu-u) = 0,
\end{equation}

ここで $u = a_1$ を仮定する。すると最初の項は次のようになる：

\begin{equation}\tag{18}
(\mu-a_1)\left\{\frac{(\mu-a_2)\cdots(\mu-a_n)}{\mu^{n-1}} - 1\right\},
\end{equation}

この表現は負の値を持つ。なぜなら、(15)式において同じ仮定をすると、

\begin{equation*}
\mu - a_1 = \frac{(b_1 - u)\cdots(b_n - u)}{(1-u)^{n-1}}
= \text{正の量},
\end{equation*}

となるからである。一方で、

\begin{equation*}
\frac{(\mu-a_2)\cdots (\mu - a_n)}{\mu^{n-1}}
=\frac{\mu-a_2}{\mu} \cdots \frac{\mu-a_n}{\mu},
\end{equation*}

であり、第二項の因子はすべて正の分数であるため、その項は 1 よりも小さい。従って (18) 式は負の値を持つ。

したがって、$u = a_1$ の仮定により(17)式の最初の項は負となる。

次に $u = b_1$ とする。すると(15)式より $\mu=u=b_1$ であり、(17)式の最初の項は正となる。

最後に、$u=a_1$ と $u=b_1$ の間では、(17)式の最初の項が連続して増加する。その表現を次のように書き換える：

\begin{equation*}
(\mu-a_1)\frac{\mu - a_1}{\mu}\cdots \frac{\mu-a_n}{\mu},
\end{equation*}

ここで $\mu$ が増加すると最初の項も増加し、それに含まれる各因子も増加する。また、負の項 $(\mu-u)$ は $u$ の増加に伴って減少し、

\begin{equation*}
\frac{(b_1-u)\cdots(b_n-u)}{(1-u)^{n-1}}
\end{equation*}

から得られる値により明らかである。したがって、$u=a_1$ と $u=b_1$ の間では、(17)式の最初の項は連続して増加し、負の値から正の値へと変化する。従って、与えられた制限内で $u$ の一つだけの値が存在し、その値により与えられた方程式が満たされる。

12. これらの結果を集めて一般問題の解を次のように得る：

\begin{equation*}
\end{equation*}

事件 \(E\) の確率は、方程式
\[
\mu - u = \frac{(\mu - c_1 p_1) \cdots (\mu - c_n p_n)}{\mu^{n-1}}
\]
から導かれる値 \(u\) に等しいもので、ここで
\[
\mu = u + \frac{(1 - c_1 (1 - p_1) - u) \cdots (1 - c_n (1 - p_n) - u)}{(1 - u)^{n-1}}
\]
であり、この値は以下の二つの量の範囲内にあります：
\[
c_1 p_1, c_2 p_2, \ldots, c_n p_n \quad \text{と} \quad 1 - c_1 (1 - p_1), 1 - c_2 (1 - p_2) \ldots, 1 - c_n (1 - p_n),
\]
前者が下限、後者が上限です。

一般の場合においても、特に \(n = 2\) の場合で証明されたように、上記の方法で決定される値 \(u\) は、その量を超えないことが推測できます：
\[
c_1 p_1 + c_2 p_2 \ldots + c_n p_n.
\]

13. 特定の検証結果は以下の通りです。

1つ目。\(p_1 = 1, p_2 = 1, \ldots, p_n = 1\) の場合、これは各事件 \(A_1, A_2, \ldots, A_n\) が発生したとき、事件 \(E\) も必ず発生すると仮定します。この場合、事件 \(E\) の発生確率は、イベントまたは原因 \(A_1, A_2, \ldots, A_n\) がすべて発生しない確率に等しくなりますから、その表現は \(1 - (1 - c_1)(1 - c_2) \cdots (1 - c_n)\) となります。

一般解（19）は
\[
\mu - u = \frac{(\mu - c_1) \cdots (\mu - c_n)}{\mu^{n-1}}
\]
を導き、ここで
\[
\mu = u + \frac{(1 - u)^n}{(1 - u)^{n-1}} = 1.
\]
したがって，
\begin{gather*}
1 - u = (1 - c_1) \cdots (1 - c_n), \\
\therefore u = 1 - (1 - c_1) \cdots (1 - c_n),
\end{gather*}
これは先に述べた事前決定と等価です。

2つ目。\(p_1 = 0, p_2 = 0, \ldots, p_n = 0\) の場合、（19）は
\begin{gather*}
\mu - u = \mu, \\
\therefore u = 0,
\end{gather*}
と推測される通りの結果を導きます。

3番目の点。$c_1, c_2, \ldots, c_n$ を小さな量とし、それらの平方や積を無視できるようにします。そうすると、方程式（19）の第二項を展開して、

\begin{gather*}
\begin{split}
  \mu - u
  &= \frac{\mu^n - (c_1p_1 + c_2p_2 + \cdots + c_np_n)\mu^{n-1}}{\mu^{n-1}} \\
  &= \mu - (c_1p_1 + c_2p_2 + \cdots + c_np_n),
\end{split}   \\
  \therefore u = c_1p_1 + c_2p_2 + \cdots + c_np_n.
\end{gather*}

この結果は、原因 $A_1, A_2, \ldots, A_n$ が互いに排他的である場合の解と一致します。しかし、これらの原因の確率が小さいほど、それらが互いに排他的になる条件に近づきます。なぜなら、どの原因も同時に現れる確率は小さくなるからです。したがって、得られた結果は、$E$ の発生確率に対する表現の極限形であることが明らかです。

4番目の点。特に $n = 2$ の場合において、$\mu$ を一般解から除去することができます。その結果は、

\[
\frac{(u - c_1p_1)(u - c_2p_2)}{c_1p_1 + c_2p_2 - u} = \frac{\{1 - c_1(1 - p_1) - u\}\{1 - c_2(1 - p_2) - u\}}{1 - u}
\]

となり、これは前の問題で得られた特定の解と一致します。

方程式（19）によって、一般的な解は高次の方程式を解くことに依存することが一般的ですが、実際的な困難は大きくありません。なぜなら、極限条件により $u$ の近似値を選択できるためであり、方程式（19）の形態は逐次近似法に適しています。

14. 問題7。---最後の問題で与えられたデータを用いて、原因 $A_1, A_2, \ldots, A_n$ の特定の組み合わせが現れた場合に、事象 $E$ が実現する確率を求めよ。

これらの原因 $A_1, A_2, \ldots, A_n$ を以前と同様に $x_1, x_2, \ldots, x_n$ で表す。問題の文脈において言及された特定の組み合わせを $\phi(x_1, x_2, \ldots, x_n)$ で表し、その実際的な発生は論理方程式

\[
\phi(x_1, x_2, \ldots, x_n) = 1
\]

によって表現される。

データは次のようになっています：
\begin{align}\tag{1}
\operatorname{Prob.} x_1 &= c_1, & \operatorname{Prob.} x_n &= c_n, \\
\operatorname{Prob.} x_1 z &= c_1 p_1, & \operatorname{Prob.} x_n z &= c_n p_n;
\end{align}して、調査の対象は次の比です：
\begin{equation}\tag{2}
\frac{\operatorname{Prob.} [\phi(x_1, x_2 \dotsc x_n)z]}{\operatorname{Prob.} [\phi(x_1, x_2 \dotsc x_n)]}.
\end{equation}
まず、分子の値を求めます。

次のように仮定します：
\begin{gather*}\tag{3}
x_1 z = t_1, \dotsc, x_n z = t_n,
\\\tag{4}
\phi(x_1, x_2 \dotsc x_n) z = w.
\end{gather*}
または、単純化のために $\phi(x_1, x_2 \dotsc x_n)$ を $\phi$ と表す場合、最後の式は
\begin{equation}\tag{5}
\phi z = w,
\end{equation}
となり、これに加えて次の式が必要です：
\begin{equation}\tag{6}
\bar{x}_1 \bar{x}_2 \dotsc \bar{x}_n z = 0.
\end{equation}

現在のシステム (3) の任意の方程式 $x_r z = t_r$ を次のように簡略化できます：
\begin{equation*}
x_r z \bar{t}_r + t_r(1 - x_r z) = 0.
\end{equation*}
同様に、方程式 (5) を簡略化し、結果を全て加算すると、論理方程式が得られます：
%-----------------------File: 150.png----------------------------
\[
\Sigma \{x_r z \bar{t}_r + t_r(1 - x_r z)\} + \bar{x}_1 \dotsc \bar{x}_n z + \phi z \bar{w} + w(1 - \phi z) = 0, \tag{7}
\]
ここで $z$ を消去すると、$w$ は $x_1, \dotsc, x_n, t_1, \dotsc, t_n$ の展開された論理関数として決定されます。

次に、上記の方程式において順次 $z = 1, z = 0$ を代入し、結果を乗算すると次のようになります：
\[
\{\Sigma (x_r \bar{t}_r + \bar{x}_r t_r) + \bar{x}_1 \dotsc \bar{x}_n + \phi \bar{w} + w \bar{\phi}\} \times (\Sigma t_r + w) = 0.
\]
この方程式を $w$ に関して展開し、Prop. i. Chap. ix. の定理に従って $\sum t_r + 1$ を $1$ に置き換えると次のようになります：
\[
Ew + E'(1 - w) = 0;
\]
ここで
\begin{gather*}
E = \Sigma (x_r \bar{t}_r + t_r \bar{x}_r) + \bar{x}_1 \dotsc \bar{x}_n + \bar{\phi}, \\
E' = \Sigma t_r \{\Sigma (x_r \bar{t}_r + t_r \bar{x}_r) + \bar{x}_1 \dotsc \bar{x}_n + \phi\}.
\end{gather*}
したがって
\[
w = \frac{E'}{E' - E}. \tag{8}
\]

この方程式の第二項は、現在、符号の重複系列 \( x_1, x_2, \dotsc, x_n, t_1, t_2, \dotsc, t_n \) に対して展開する必要があります。
この目的を達成するために、結果の展開の要素を三つの異なるクラスに分類し、それぞれのクラスに対応する係数を個別に決定することが最も便利です。

まず、\(\bar{t}_1 \dotsc \bar{t}_n\) の因子である要素を考えましょう。\( t_1 = 0, \dotsc, t_n = 0 \)と置くと，
\[
E' = 0,\quad E = \Sigma x_r + \bar{x}_1 \dotsc \bar{x}_n + \bar{\phi}.
\]
明らかに、\(x\) の符号に対してどのような値（0 または 1）を与えても \(E\) は消えないことが分かります。したがって、\(\bar{t}_1 \dotsc \bar{t}_n\) を含むすべての要素の係数は 0 です。

次に、\(\bar{t}_1 \dotsc \bar{t}_n\) の因子を含まないかつ \(x_1, \dotsc, x_n\) と \(t_1, \dotsc, t_n\) の二つの符号系列に対して対称的な要素を考えます。ここで、対称的要素とは、\(x_1\) を \(t_1\) に、\(x_2\) を \(t_2\) に変換したときに不変となるものを指します。要素 \( x_1 \dotsc x_n\ t_1 \dotsc t_n, \bar{x}_1 \dotsc \bar{x}_n\ \bar{t}_1 \dotsc \bar{t}_n, \dotsc \) はこの意味での対称的です。

これらのすべての対称的要素について、明らかに
\[
\sum (x_r \bar{t}_r + t_r \bar{x}_r)
\]
が消えます。また、\(\bar{t}_1 \dotsc \bar{t}_n\) を含まない要素については、さらに \(\bar{x}_1 \dotsc \bar{x}_n\) も消えるので，
\begin{gather*}
E = \bar{\phi} \qquad E' = \Sigma t_r(\phi),   \\
w = \frac{\sum t_r(\phi)}{ \sum t_r(\phi) - \bar{\phi} }.
\end{gather*}
\(\phi\) 内に \(x\) の因子が含まれる要素については、上記の式の第二項は 1 になります。一方、\(\bar{\phi}\) 内に \(x\) の因子が含まれる要素については、0 になります。したがって、対称的で \(\bar{t}_1 \dotsc \bar{t}_n\) を含まない要素について、\(x\) の因子が \(\phi\) 内に存在するものは係数が 1 であり、\(\phi\) 内に存在しないものは係数が 0 です。

最後に、二つの符号の組に対して非対称でありながら$\bar{t}_1 \dotsm \bar{t}_n$を含まない成分について考慮する。

明らかに、$E$ および $E'$ は消えないため、(8) 式における $w$ の分数値の分子が分母を超えることがわかる。その値は従って 1, 0, もしくは $\frac{0}{0}$ で表すことができない。したがって、論理的な展開においてはこの値は $\frac{1}{0}$ で表されるべきである。このような成分の係数となる。

15. これより $w$ を $x_1, \dotsc x_n$, $t_1, \dotsc t_n$ の論理的な関数として展開した最終的な論理方程式は以下の形をとる：
\[
  w = \sum\nolimits_1 (XT)
  + 0\{ \sum\nolimits_2(XT) + \bar{t}_1\dotsm \bar{t}_n \}
  + \frac{1}{0}
  \left(\text{他の成分の和}\right) \tag{9}
\]
ここで、$\sum_1 (XT)$ は因子 $X$ が $\phi$ 内に含まれるすべての対称的な成分の和を表し、$\sum_2(XT)$ は因子 $X$ が $\phi$ 内に含まれないすべての対称的な成分の和を表す。また、成分 $\bar{x}_1 \dotsm \bar{x}_n\; \bar{t}_1 \dotsm \bar{t}_n$ が現れた場合でも、どちらの場合もこれを拒否する。

注：上記のコードは元の形式を再現するために醜いものですが、以下はより簡潔なコードで、基本的には同じ内容です。
\[
  w = \sum\nolimits_1 (XT)
  + 0\{ \sum\nolimits_2(XT) + \bar{t}_1\dotsm \bar{t}_n \}
  + \frac{1}{0} \left(\text{他の成分の和}\right) \tag{9}
\]

論理から代数へ移行すると、以下のように観察できる。%-----------------------File: 152.png----------------------------ここで以及すべて類似の例において、方程式系$x_1, \dotsc x_n, t_1 \dotsc t_n$の値を決定するための代数的なシステムが形成される際に用いられる関数$V$は、問題に含まれる任意の関数$\phi$の性質によらず独立である。ただし、これはデータの表現におけるものではなく、問題の求められる結果（\textit{qu{\ae}situm}）の表現においてである。したがって、現在の例では、

\begin{multline*}
  \hfill \text{問題 }w = \frac{\sum\nolimits_1 (XT)}{V}, \hfill \\
  \text{ここで }\hfill
  V = \sum\nolimits_1 (X T) + \sum\nolimits_2 (X T)
      + t_1 \dotsc \bar{t}_n \hfill \\
  \hfill = \sum (XT) + \bar{t}_1 \dotsc t_n. \hfill \tag{10}
\end{multline*}

ここで$\sum (X T)$は$x$および$t$の符号のすべての対称的な成分を表し、ただし$\bar{x}_1 \dotsc \bar{x}_n,\; \bar{t}_1 \dotsc \bar{t}_n$の成分は含まれていない。この$V$の値は前回の問題の解法において実質的に用いられたものと同一であり、したがって、その結果を活用することが可能である。

もし、先に述べた解法において
\[
\frac{x_1 t_1}{\bar{x}_1 \bar{t}_1} = m_1,\quad
\frac{x_n t_n}{\bar{x}_n \bar{t}_n} = m_n,\quad
\frac{x_1}{\bar{x}_1} = n_1, \text{ など}
\]
と仮定すると、得られる結果は次のように表すことができます：
\[
\text{Prob. }w = \frac{M_1}{M + N - 1},   \tag{11}
\]
ここで $M_1$ は関数 $\phi$ から成分 $x_1 \dotsc \bar{x}_n$ を除き、残りを同じ成分 $\bar{x}_1 \dotsc \bar{x}_n$ で割った後、$\frac{x_1}{\bar{x}_1}$ を $m_1$、$\frac{x_2}{\bar{x}_2}$ を $m_2$ と置き換えたものとなります。$M$ および $N$ の値は前問題と同じです。これらの値に戻り、対応する $m_1$, $m_2$, 等 の値を用いて
\[
\text{Prob. }w = M_1(\mu + \nu - 1),
\]
を得ます。一般の $m_r$ および $n_r$ の値は
\[
m_r = \frac{c_r p_r}{\mu - c_r p_r},\qquad
n_r = \frac{c_r(1-p_r)}{\mu - c_r(1-p_r)},
\]
であり、$\mu$ および $\nu$ は方程式の系を解くことで与えられます。
\[
\mu + \nu -1 = \frac{(\mu-c_1p_1)\dotsc (\mu-c_np_n)}{\mu^{n-1}}
= \frac{ \{\nu-c_1(1-p_1)\} \dotsc \{\nu-c_n(1-p_n)\} }{\nu^{n-1}} .
\]
上記の Prob. $w$ の値は式 (2) の分数の分子となります。その分母を決定するために、
\[
\phi(x_1,x_2 \dotsc x_n)=v,
\]
または
\[
\phi = v ;
\]
と仮定します。これにより
\[
\phi \bar{v} + v \bar{\phi} = 0.
\]

この方程式の最初の項を（7）式において $\phi z\bar{w}+w(1-\phi z)$ の対応する形として代入することで，主な論理方程式を得ます：
\[
\sum \{x_r z \bar{t}_r + t_r(1-x_r z)\} + \bar{x}_1 \dotsc \bar{x}_n z + \phi \bar{v} + v \bar{\phi} = 0,
\]
ここで $z$ を消去し，第II章第九章の命題により簡略化すると，
\[
\phi \bar{v} + v \bar{\phi} + \sum t_r \{\sum (x_r\bar{t}_r + t_r\bar{x}_r) + \bar{x}_1 \dotsc \bar{x}_n\}=0.
\]
したがって，
\[
v = \frac{\phi + \sum t_r \{\sum(x_r\bar{t}_r + t_r\bar{x}_r) + \bar{x}_1 \dotsc \bar{x}_n\} }{2\phi-1}
\]
となり，同様に展開すると，
\[
\begin{split}
v = \sum\nolimits_1(XT)+\bar{t}_1\dotsc\bar{t}_n \sum\nolimits_1(X) + 0\{\sum\nolimits_2(XT)+\bar{t}_1\dotsc\bar{t}_n\sum\nolimits_2(X)\}
\\
+ \frac{1}{0}(\text{他の成分の和}).
\end{split}  \tag{12}
\]

ここで $\sum_1(X)$ は $\phi$ 内に含まれるすべての成分の合計を表し，$\sum_2(X)$ は $\phi$ 内に含まれないすべての成分の合計を表します。これらの表現は $\phi$ と $1-\phi$ の代わりに使用され，対称性を保つためです。

したがって，
\[
\sum_1(X)+\sum_2(X)=1
\]
であり，かつ以前と同様に，
\[
\sum_1(XT)+\sum_2(XT) = \sum(X T).
\]
これにより $V$ は以前と同じ値を取ることになり，
\[
\operatorname{Prob. } v
= \frac{\sum_1(XT)+\bar{t}_1\dotsc\bar{t}_n \sum_1(X)}{V},
\]

または前のケースと同様に変形すると，
\[
\operatorname{Prob. } v = \frac{M_1+N_1}{M+N-1},  \tag{13}
\]

ここで $N_1$ は $\phi$ を $\bar{x}_1\dotsc\bar{x}_n$ で割った結果を $n_1, n_2, 等$ に置き換えたものから形成されます。

問題の最終的な解は，分数
\[
\frac{\operatorname{Prob. } \phi(x_1,\dotsc x_n)z}
     {\operatorname{Prob. } \phi(x_1,\dotsc x_n)}, \mathrm{ or }
\frac{\operatorname{Prob. } w}{\operatorname{Prob. } v}.
\]
の各項に決定された値を代入することにより得られます。したがって，（11）と（13）より，
\[
\textrm{求める確率} = \frac{M_1}{M_1 + N_1}
\]

関数 $M_1$ および $N_1$ の形成過程についての非常に僅かな注意から、その手続きは著しく簡略化できることがわかります。実際、一般問題の解を次のルールとして表すことができます：

\begin{itemize}
    \item $\phi(x_1, x_2, \dotsc, x_n)$ の中で含まれている成分 $\bar{x}_1 \dotsm \bar{x}_n$ を拒否し、残りのすべての成分から因子 $\bar{x}_1$, $\bar{x}_2$, 等 を削除し、結果において $x_r$ を $\frac{c_r p_r}{\mu - c_r p_r}$ に変更します。これを $M_1$ と呼びます。
    \item 再び、$\phi(x_1, x_2, \dotsc, x_n)$ の中で見つかった成分 $\bar{x}_1 \dotsm \bar{x}_n$ を 1 に置き換え、残りのすべての成分から因子 $\bar{x}_1$, $\bar{x}_2$, 等 を削除し、結果において $x_r$ を $\frac{c_r (1-p_r)}{\nu - c_r (1-p_r)}$ に変更します。
\end{itemize}

これらの手順により求められる解は以下の式で表されます：
\[
\frac{M_1}{M_1 + N_1} \tag{14}
\]
ここで、$\mu$ および $\nu$ は以下の方程式の解によって定められます：
\begin{multline*}
    \mu + \nu - 1
    = \frac{(\mu - c_1 p_1) \dotsm (\mu - c_n p_n)}{\mu^{n-1}} \\
    = \frac{\{\nu - c_1 (1 - p_1)\} \dotsm \{\nu - c_n (1 - p_n)\}}{\nu^{n-1}}.
\tag{15}
\end{multline*}

$\mu$ および $\nu$ の範囲は、前回の問題と同一です。これは連続性の一般原則から推測できますが、他の考慮事項により十分な制限条件も確立することができます。

したがって、確率論における第17章の定理ivより、代数方程式の基本系において $x_1, \dotsc x_n, t_1, \dotsc t_n$ は正の真分数であることが示されます。これにより
\[
    \frac{x_r}{1-x_r} = n_r = \frac{c_r (1-p_r)}{\nu - c_r (1-p_r)}
\]
が成り立ちます。従って、一般に $n_r$ は正の量であり、したがって
\[
    \nu > c_r (1-p_r)
\]
が成り立ちます。

同様に、
\[
    \frac{x_r t_r}{(1-x_r)(1-t_r)} = m_r = \frac{c_r p_r}{\mu - c_r p_r}
\]
が成り立つため、一般に
\[
    \mu > c_r p_r
\]
が成り立ちます。

16. このように表現される二つの条件のクラスは、$\mu$と$\nu$を決定する方程式の根の中からどの根を選ぶべきか一般に決定するのに十分である可能性が高い。特に$n = 2$の場合を考える。このとき、

\begin{gather*}
  \mu + \nu - 1 = \frac{(\mu - c_1 p_1)(\mu - c_2 p_2)}{\mu} 
= \mu - (c_1 p_1 + c_2 p_2) + \frac{c_1 p_1 c_2 p_2}{\mu},
\\
  \therefore\nu = 1 - c_1 p_1 - c_2 p_2 + \frac{c_1 p_1 c_2 p_2}{\mu}
= 1 - c_1 p_1 - \frac{(\mu - c_1 p_1) c_2 p_2}{\mu}.
\end{gather*}

ここで$\mu > c_1 p_1$であるから、一般に

\[
  \nu < 1 - c_1 p_1
\]

が成り立つ。同様にして、

\[
  \nu < 1 - c_2 p_2, \quad
  \mu < 1 - c_1 (1 - p_1), \quad
  \mu < 1 - c_2 (1 - p_2)
\]

が得られる。

既に示されたように、これらの条件に対する$\mu$の値はただ一つであり、

\[
  \mu > c_r p_r, \quad
  \mu < 1 - c_r (1 - p_r),
\]

である。したがって、この場合の解は定まっている。そして私は、同じ方法が一般に適用可能で十分であると考える。しかし、これはさらに明確な光を投げかけるべき問題である。

上記の結果を検証するために、$\phi(x_1, \dotsc, x_n) = 1$と仮定する。これは前問で考慮された状況に相当する。このとき、

\[
  M_1 = \frac{\mu^n}{(\mu - c_1 p_1)(\mu - c_2 p_2)} - 1
= \frac{\mu}{\mu + \nu - 1} - 1 \quad (15)，
\]

\[
  N_1 
= \frac{\nu^n}{(\nu - c_1(1 - p_1))(\nu - c_2(1 - p_2))} - 1
= \frac{\nu}{\mu + \nu - 1} - 1.
\]

これらを(14)に代入すると、

\[
  \operatorname{Prob. } z = 1 - \nu,
\]

となり、これは前回の解と一致する。

再び、$\phi(x_1, \dotsc x_n) = x_1$ とします。この式を展開し、因子 $\bar{x}_2, \dotsc \bar{x}_n$ を省略すると、

$$x_1 (x_2 + 1)\dotsc (x_n+1)$$

となりますから、

\begin{align*}
M_1 &= \frac{c_1 p_1 \mu^{n-1}}{(\mu-c_1 p_1)\dotsc (\mu-c_n p_n)} = \frac{c_1 p_1}{\mu + \nu - 1} \quad \text{式(15)より}, \\
N_1 &= \frac{c_1 (1-p_1)\nu^{n-1}}{\{\nu - c_1(1-p_1)\}\dotsc \{\nu - c_n(1-p_n)\}} = \frac{c_1(1-p_1)}{\mu + \nu - 1}.
\end{align*}

これらを代入すると、

\begin{equation*}
Eが起こる確率（A_1が起こったとき） = p_1
\end{equation*}

を得ます。この結果はデータによって検証されます。類似の検証も容易に追加できます。

次に、

\begin{equation*}
\phi(x_1, \dotsc x_n) = x_1 \bar{x}_2 \dotsc \bar{x}_n + x_2 \bar{x}_1 \bar{x}_3 \dotsc \bar{x}_n + \dotsb + x_n \bar{x}_1 \dotsc \bar{x}_{n-1}
\end{equation*}

の場合を検討します。このとき、

\begin{gather*}
M_1 = \frac{c_1 p_1}{\mu - c_1 p_1} + \dotsb + \frac{c_n p_n}{\mu - c_n p_n}, \\
N_1 = \frac{c_1 (1-p_1)}{\nu - c_1(1-p_1)} + \dotsb + \frac{c_n (1-p_n)}{\nu - c_n(1-p_n)};
\end{gather*}

から、

\[
\left. \text{\parbox{0.45\linewidth}{$E$ が生じる確率（$A_1, A_2, \dotsc, A_n$ のいずれかが起きたとき）}} \right\}
= \frac{\sum \displaystyle \frac{c_r p_r}{\mu - c_r p_r}}{\sum \displaystyle \frac{c_r p_r}{\mu - c_r p_r} + \sum \displaystyle \frac{c_r(1-p_r)}{\nu - c_r(1-p_r)}}
\]
を得ます。

ここで注意すべきは、この場合が互いに排他的な性質を持つ原因 $A_1, A_2, \dotsc, A_n$ の要素をデータの一部とする有名なケースとは異なる点です。後者のケースでは、原因 $A_1, A_2, \dotsc$ の確率が決定された観測条件が互いに排他的であるという条件を含んでいます。

最後に、$\phi(x_1,\dotsc x_n)=x_1 x_2 \dotsc x_n$ の場合について考慮します。
このとき、

\begin{gather*}
M_1 = \frac{c_1 p_1 \dotsc c_n p_n}
           {(\mu - c_1 p_1)\dotsc (\mu-c_n p_n)}
=\frac{c_1 p_1 \dotsc c_n p_n}{\mu^{n-1}(\mu + \nu - 1)}, \\
N_1 = \frac{c_1(1-p_1)\dotsc c_n(1-p_n)}
           {\{\nu - c_1(1-p_1)\} \dotsc \{\nu-c_n(1-p_n)\}}
=\frac{c_1(1-p_1)\dotsc c_n(1-p_n)}{\nu^{n-1}(\mu + \nu -1)}.
\end{gather*}

したがって、以下の結果を得ます：\[
\left. \text{\parbox{0.45\linewidth}{すべての原因 $A_1, \dotsc, A_n$ が共に作用した場合に、事象 $E$ が生じる確率}} \right\}
= \frac{p_1 \dotsc p_n \nu^{n-1}}{p_1 \dotsc p_n \nu^{n-1} + (1-p_1)\dotsc (1-p_n)\mu^{n-1}}.
\]
この表現は、どの $p_1, \dotsc p_n$ のいずれかが 1 に等しいとき、その値を 1 とするべきである。

17. 問題 VIII.---特定の原因 $A_1, A_2,\dotsc A_n$ が制限されており、これらはすべて失敗することができず、しかし一定の組み合わせでしか現れることができ、その組み合わせは方程式
\begin{equation*}
\phi(A_1, A_2 \dotsc A_n)=1,
\end{equation*}
によって表されます。また、これらの原因の個々の確率 $c_1,\dotsc c_n$ と、それぞれの原因が現れた場合に事象 $E$ が生じる確率 $p_1, \dotsc p_n$ が与えられています。このとき、事象 $E$ の確率を求めよ。

この問題は、最後に考慮した問題とはいくつか異なる点がありますが、特に方程式 $\phi(A_1, \dotsc A_n)=1$ によって表される制限がデータの一つとして与えられること、そしてその制限は問題の数値要素に関する知識を導き出す経験と一致していると仮定されることにあります。

事件 \(A_1, \dotsc A_n\) をそれぞれ \(x_1, \dotsc x_n\) で表し、事件 \(E\) を \(z\) で表すと、以下の関係が成り立ちます：
\[
\operatorname{Prob. } x_r = c_r, \quad
\operatorname{Prob. } x_r z = c_r p_r.
\]
一般に、次のように仮定します：
\[
x_r z = t_r,
\]
そして、この方程式の系と以下の式群を組み合わせます：
\[
\bar{x}_1 \dotsc \bar{x}_n = 0, \quad \phi(x_1, \dotsc x_n) = 1, \quad \text{または } \phi = 1,
\]
これにより、\(z\) の展開表式を得ることができます：
\[
z = \sum (XT) + 0 \bar{t}_1 \bar{t}_2 \dotsc \bar{t}_n \sum (X),
\]
ここで \(X\) は \(\phi\) 内の各成分を順に表し、\(T\) は符号 \(t_1, \dotsc t_n\) の同じ数列の成分を表します；\(\sum(XT)\) には、両方の符号群に対する対称的な成分のみが含まれます。

現在の場合における簡約法は、以前の問題で例示した方法と非常に類似しているため、私はその結果だけを示します。以下のように得られます：
\[
\operatorname{Prob. } z = \frac{M}{M+N}
\]
そして関係式は以下のようになります：
\[
\frac{M_1}{c_1 p_1} \dotsc = \frac{M_n}{c_n p_n} = \frac{N_1}{c_1 (1-p_1)} = \frac{N_n}{c_n(1-p_n)} = M+N.
\]
ここで、\(M\) は \(\phi(x_2, \dotsc x_n)\) 内のすべての因子 \(\bar{x}_1, \dotsc \bar{x}_n\) を削除し、結果で \(x_1\) を \(m_1\)、\(x_n\) を \(m_n\) に置き換えたものから形成されます。また、\(N\) は \(M\) 内の \(n_1\) に \(m_1\)、\(等\) を置き換えたものから形成されます；さらに、\(M_1\) は \(M\) のうち \(m_1\) の因子となる部分で構成され、\(N_1\) は \(N\) のうち \(n_1\) の因子となる部分で構成されます。

例として、原因 $A_1, \dotsc, A_n$ が互いに排他的である場合を考えましょう。このとき、関数 $\phi(x_1, \dotsc, x_n)$ は以下のようになります：
\[
\phi(x_1, \dotsc, x_n) = x_1 \bar{x}_2 \dotsm \bar{x}_n + x_n \bar{x}_1 \dotsm \bar{x}_{n-1}.
\]
%-----------------------File: 159.png----------------------------
したがって、
\begin{align*}
M &= m_1 + m_2 + \dotsb + m_n,\\
N &= n_1 + n_2 + \dotsb + n_n,\\
M_1 &= m_1, N_1 = n_1, & \text{etc.}
\end{align*}
これらを代入すると、以下の式が得られます：
\[
\frac{m_1}{c_1 p_1} \dotsb = \frac{m_n}{c_n p_n}
= \frac{n_1}{c_1 (1-p_1)} \dotsb
= \frac{n_n}{c_n (1-p_n)} = M + N.
\]
したがって、以下の式を得ます：
\[
\frac{m_1 + m_2 + \dotsb + m_n}
     {c_1 p_1 + c_2 p_2 + \dotsb + c_n p_n}
= M + N,
\]
または
\[
\frac{M}{c_1 p_1 + c_2 p_2 + \dotsb + c_n p_n} = M + N.
\]
したがって、式 (3) から以下の確率を得ます：
\[
\operatorname{Prob. } z = c_1 p_1 + c_2 p_2 + \dotsb + c_n p_n,
\]
これは既知の結果です。

他の特定の場合でも、システム (4) の解は容易に得られます。しかし、ほとんどの場合において、この式が非常に複雑な結果を導くことは明らかです。また、一般的問題におけるデータで仮定された原因間の関数的な関係が実際の経験で頻繁に現れるとは思われません；ただし、先述した特定の場合を除きます。

一般の問題が制約によって変更され、事件 $E$ が起こらないという制約が導入され、すべての原因 $A_1, \dotsc, A_n$ が欠如している場合（既に述べた制約とは異なる）で、方程式 $\phi(A_1, \dotsc, A_n) = 1$ の残りの条件を保持する場合、最終的な論理式は以下のようになります：
\[
z = \sum_1(XT) + 0 \sum(X),
\]
ここで $\sum(X)$ は以前と同様に $\phi(x_1, \dotsc, x_n)$ に等しく、$\sum_1(XT)$ は $\phi$ の特定の成分 $\bar{x}_1 \dotsm \bar{x}_n$（その中に含まれている場合）を排除し、その後の各 $x$-成分にそれぞれの $t$-成分を乗じて形成されます。原因が互いに排他的である特定の場合においては、この方法で導かれる確率 $\operatorname{Prob. } z$ の値は以前と同一です。

18. 問題 IX.---前問のデータを仮定し、ある事象 \( E \) が現れたとき、その事象と特定の原因 \( A_r \) の関連性を決定する必要があるとする。言い換えれば、観察された事象 \( E \) が発生した場合の原因 \( A_r \) の \textit{\`{a} posteriori} 機会を定める必要がある。

この場合、以下の分数の値を求めなければならない：
\[
\frac{\operatorname{Prob. } x_r z}{\operatorname{Prob. } z}, \quad \text{または} \quad \frac{c_r p_r}{\operatorname{Prob. } z}
\]
前問と同様に、原因の関連性や無関連性に関する仮定に基づいて \( \operatorname{Prob. } z \) の値が割り当てられているため、式 (1) においてその要素の値を代入することで、すべての場合においてこの問題は確定的な解を持つことが明白である。

原因の \textit{\`{a} priori} 機会が等しい場合、つまり \( c_1 = c_2 \dotso = c_r \) のとき、各原因に対する式 (1) の値は量 \( p_r \) に比例する。したがって、原因の関連性のどのような形であれ、各原因の \textit{\`{a} posteriori} 機会は観察された事象 \( E \) の機会と比例する。この定理の特別な場合として、互いに排他的な原因の場合がよく知られている。そのとき、以下の式が成り立つ：
\[
\frac{\operatorname{Prob. } x_r z}{\operatorname{Prob. } z} = \frac{c_r p_r}{{\textstyle \sum} c_r p_r} = \frac{p_r}{p_1 + p_2 \dotso + p_n},
\]
ここで \( c_1, \dotsc, c_n \) の値は等しいものとする。

これらの定理や類似の定理について、原因が互いに排他的な場合における証明では、関数記号 $\phi$ の導入は必要ありません。実際、それは私たちはすべて可能かつ考えられる仮説を選択する権利を主張しているのです。しかし、どの形でも、この作品の方法による問題解決法は、与えられたデータの数が無限大であるような問題に対する解法には、ある程度複雑な性格を持つことになります。そのシステム的な展開が、まず論理的関係、次に数値的关系を最初から順序立てて提示するかどうかに関わらず、その過程は長く、時として退屈であることがありますが、それらの長さと退屈さが問題解決に対する補償となるかどうかについては、私は問うことはありません。この方法の最大の価値は、人間の理性に無力な質問でも私たち自身を支配する力を与える点にあることでしょう。そのため、その適用例として、困難から見れば拒否されるかもしれないが直ちなる利益を見出せないような問題に対する適用も、必要とされていないとは言えません。

19. 以下の問題は以前のものよりもずっと簡単な説明です。

**問題 X**：ある自然現象が与えられた状況下で発生する確率は $p$ である。また、その現象を恒久的な原因によって存在する確率は $a$ であると観測結果も記録されている。すなわち、そのような状況下では常にその事件が発生させる原因がある。ある自然現象が与えられた状況下で連続して $n$ 回発生した場合、次の $(n+1)$ 回目にも同様に発生する確率は何か？また、このような観測後、その恒久的な原因の存在確率は何か？

**最初の場合**：$t$ を恒久的因 素の存在を表し、$x_1, x_2, \dots, x_{n+1}$ を連続して発生する自然現象の各回を表します。

もし永久的な原因が存在すれば、事象$x_1, x_2, \dots, x_{n+1}$は必然的な結果となる。従って，
\begin{equation*}
t = v x_1, \quad t = v x_2, \text{ およびその他のもの},
\end{equation*}
不定記号を消去すると，
\begin{equation*}
t(1-x_1) = 0, \quad t(1-x_2) = 0, \quad t(1-x_{n+1}) = 0.
\end{equation*}

今、$x_1 x_2 \dotsm x_n$が起こった場合に、事象$x_{n+1}$が起こる確率を求めることを考える。すなわち、以下の分数の値を求めることになる：
\begin{equation*}
\frac{\operatorname{Prob. } x_1 x_2 \dotsm x_{n+1}}
     {\operatorname{Prob. } x_1 x_2 \dotsm x_n}.
\end{equation*}

まず、$\operatorname{Prob. } x_1 x_2 \dotsm x_n$の値を求めることから始める。
%-----------------------File: 162.png----------------------------

組合せ$x_1 x_2 \dotsm x_n$を$w$と表すと、以下の論理方程式が得られる：
\begin{gather*}
t(1-x_1) = 0, \quad t(1-x_2) = 0, \dotsc t(1-x_n) = 0, \\
x_1 x_2 \dotsm x_n = w.
\end{gather*}

最後の式を以下のように変形する：
\[
(x_1 x_2 \dotsm x_n)(1-w) + w(1 - x_1 x_2 \dotsm x_n) = 0,
\]
これを先ほどの方程式に加えると、
\[
\sum t(1-x_i) + x_1 x_2 \dotsm x_n (1-w) + w(1- x_1 x_2 \dotsm x_n) = 0, \tag{1}
\]
ここで$\sum$はすべての$i$（$i$は1から$n$まで）についてを表す。この論理方程式を数値条件と結びつけるために、
\[
\operatorname{Prob. }x_1 = \operatorname{ Prob. }x_2 \dotsc = \operatorname{ Prob. }x_n = p, \quad \operatorname{Prob. }t = a;
\]
を用いる。我々の目的は$\operatorname{Prob. }w$を求めることである。

(1)から、
\begin{gather*}
w = \frac{\sum t(1-x_i) + x_1 x_2 \dotsm x_n}{2x_1 x_2 \dotsm x_n - 1} \\
= \frac{\sum (1-x_i) + x_1 x_2 \dotsm x_n}{2x_1 x_2 \dotsm x_n - 1} t
+ \frac{x_1 x_2 \dotsm x_n}{2x_1 x_2 \dotsm x_n - 1} (1-t), \tag{2}
\end{gather*}
$t$について展開すると、さらに$x_1, x_2, \dotsc, x_n$について展開する必要がある。

もし $x_1 = 1, x_2 = 1, \dots, x_n = 1$ とすると、$t$ の係数と $1-t$ の係数はともに 1 になります。同じ記号に 0 と 1 の交換によって形成された他の値を与えると、明らかに $t$ の係数が負になる一方で、$1-t$ の係数は 0 になります。したがって、完全な展開 (2) は以下のようになります。

\begin{multline*}
  w = x_1 x_2 \dotsm x_n t + x_1 x_2 \dotsm x_n(1-t)
    + 0(1 - x_1 x_2 \dotsm x_n)(1-t)   \\
  + \text{構成項のうち係数が $\frac{1}{0}$ であるか、または $\frac{1}{0}$ に等しいもの}
\end{multline*}

ここで、

\begin{multline*}
  V = x_1 x_2 \dotsm x_n t + x_1 x_2 \dotsm x_n (1-t)
    + (1 - x_1 x_2 \dotsm x_n) (1-t)   \\
  = x_1 x_2 \dotsm x_n t + 1 - t;
\end{multline*}

したがって、論理から代数へ移行すると、

\begin{multline*}
    \frac{x_1 x_2 \dotsm x_n t + x_1(1-t)}{p}
  = \frac{x_1 x_2 \dotsm x_n t + x_2(1-t)}{p} \dotso
\\ \hfill
  = \frac{x_1 x_2 \dotsm x_n t + x_n(1-t)}{p}
  = \frac{x_1 x_2 \dotsm x_n t}{a} = x_1 x_2 \dotsm x_n t + 1 - t.
\\ \hfill
  \text{確率 }w = \frac{x_1 x_2 \dotsm x_n}{x_1 x_2 \dotsm x_n t + 1 - t}.
\end{multline*}

以上の方程式の形から明らかに、$x_1 = x_2 = \cdots = x_n$ であることが分かります。これらの量をすべて $x$ で置き換えると、システムは次のようになります：
\begin{align*}
  \frac{x^n t + (1-t)x}{p} &= \frac{x^n t}{a} = x^n t + 1 - t, \\
  \text{Prob. }w &= \frac{x^n}{x^n t + 1 - t};
\end{align*}
これから容易に導くことができ、
\[
  \text{Prob. }w = \text{ Prob. }x_1 x_2 \cdots x_n
  = a + (p-a)\left( \frac{p-a}{1-a} \right)^{n-1}
\]
が得られます。この結果において $n$ を $n+1$ に置き換えると、
\[
  \text{Prob. }x_1 x_2 \cdots x_{n+1}
  = a + (p-a)\left( \frac{p-a}{1-a} \right)^n
\]
が得られます。したがって、
\[
  \frac{\text{Prob. }x_1 x_2 \cdots x_{n+1}}
       {\text{Prob. }x_1 x_2 \cdots x_n}
  =
  \frac{a + (p-a)\left( \frac{p-a}{1-a} \right)^n}
       {a + (p-a)\left( \frac{p-a}{1-a} \right)^{n-1}}
\tag{3}
\]
が得られ、これは現象が $n$ 回繰り返された場合に、その $(n+1)$ 番目の出現の確率を表します。第XIX章の方法によれば、$a$ の値は $p$ を超えることはできません。

次の検証は明らかです：

1. まず、$a = 0$ の場合、式は $p$ に簡約され、これは当然のことです。確実性が存在しない永久的な原因のない状況では、現象の連続的な出現は独立となります。

2. 次に、$p = 1$ の場合、式は 1 に簡約され、これも当然のことです。

3番目、もし \( p = a \) である場合、式は 1 になりますが、\( a = 0 \) の場合はこの例外を除きます。ある事象の確率が与えられた状況下でその事象を常に生み出す原因がある確率と等しいとき、その事象が過去にその状況下で観察されたという事実は、同じ状況下での再出現を確定的にします。\footnote{我々は再び幼少期に戻ることも思い出することもできず、上記のような結果が若者の事件の関連性から信頼性を持つ理由をどのように説明するかについて言及することはできません。しかし、一般的には、その期待の強さは、因果関係を推論（理性の一部として）する必要性と、その能力を鍛えるための狭いながら印象的な経験に由来すると推測できます。したがって、すべての連続性は原因と効果の連続性への参照です。著者の友人の幼少期の出来事で、彼が寝る前に聞こえた彼の兄との会話は、「なぜ夜に眠ることによって朝は明るくなるのか？」という関連性のある質問でした。兄は年齢が1つ上だったため、答えることができました。「夜に子供たちが寝ないとしても朝は明るくなる」と言いました。}%終脚注

4番目、\( n \) が増加するにつれて、式の値は 1 に近づきます。これは、事象の頻度が増えると再出現の確率も増加することを示唆します。これは自然な期待の法則に合致しています。

\textsc{第二の場合}。現在、我々は恒久的な原因の存在の後天的確率を求めます。これには既に決定された分母を用いて、以下の分数の値を求めることが必要です。
\[
  \frac{\operatorname{Prob. } t x_1 x_2 \dotsc x_n}
       {\operatorname{Prob. }   x_1 x_2 \dotsc x_n}
\]

numerator を決定するために、次のように仮定します：
\[
t x_1 x_2 \dotsc x_n = w.
\]
その後、以前と同様の手順で進めると、論理的な開発から次の式を得ます：
\[
w = t x_1 x_2 \dotsc x_n + 0 (1-t).
\]
次に、論理から代数へ移行すると、すぐに次の結果を得ることができます：
\[
\operatorname{Prob. } w = a.
\]
この結果は予想できるものでした。その後、上記分数の分子と分母の値をそれぞれ代入すると、恒久的原因に対する事後確率の表式は以下のようになります：

%-----------------------File: 165.png----------------------------
\begin{equation*}
\frac{a}{a+(p-a)\left(\frac{p-a}{1-a}\right)^{n-1}}.
\end{equation*}

この表式の値は、$n$ の値が増加するにつれて明らかに増大します。

上記の結果（3）を検証するために、学識ある相手から感謝を表します。彼は確率論に関する原始的な貢献によりすでに言及されており、

\[
\frac{a}{a+(p-a)\left(\frac{p-a}{1-a}\right)^{n-1}}
\]

という表式の最初の結果を検証していただきました。

``事件（状況下において）が起こる確率 \textit{\`{a} priori} が \( p \)、ある原因 \( C \) の存在によってその事件が必然的に生じることの確率が \( a \) とする。また、\( C \) のような原因がない場合にその事件が起こる確率を \( x \) とすると、以下の等式が成り立ちます：
\[
p = a + (1-a)x
\]
従って，
\[
x = \frac{p - a}{1 - a}
\]

観察された現象は、その事件が \( n \) 回起こることです。\( C \) の存在を仮定した場合の \textit{\`{a} priori} 確率は、
\[
1,
\]
\( C \) の存在を仮定しない場合の \textit{\`{a} priori} 確率は、
\[
x^n
\]
となります。従って、\( C \) が存在する \textit{\`{a} posteriori} 確率は
\[
\frac{a}{a + (1-a)x^n}
\]
であり、\( C \) が存在しない \textit{\`{a} posteriori} 確率は
\[
\frac{(1-a)x^n}{a + (1-a)x^n}
\]
となります。したがって、次の再発生の確率は
\[
\frac{a}{a + (1-a)x^n} \times 1 + \frac{(1-a)x^n}{a + (1-a)x^n} \times a
\]
または
\[
\frac{a + (1-a)x^{n+1}}{a + (1-a)x^n}
\]
となります。ここで \( n \) を \( \displaystyle \frac{p - a}{1 - a} \) に置き換えると、これは式（3）と一致します。

同様の検証は、一般方法を直接適用した結果についてもおそらく見つけることができるでしょう。

同じ状況下において、\( n \) 回の機会中 \( r \) 回起こり、残り \( n-r \) 回失敗した場合に、次の \( (n+1)^{\mathit{th}} \) 回にその事件が起こる確率は
\[
\frac{a + m(p - a)\left(\frac{p - la}{1 - a}\right)^r}
     {a + m(p - la)\left(\frac{p - la}{1 - a}\right)^{r-1}}
\]
となります。ここで、\( m = \displaystyle \frac{n(n-1) \cdots (n-r+1)}{1 \cdot 2 \cdots r} \) で、\( l = \frac{r}{n} \)。

永久的な原因（ \( r \) が \( n \) より小さい場合）の確率は 0 です。これは容易に検証できます。''

もし事件が発生する確率を$p$、その事件が発生した場合に恒久的な原因によるものである確率を$c$とするとき、$n$回連続して観察された発生後、${(n+1)}^{\mathit{th}}$の類似な機会で再発する確率は以下の式により与えられます：
\[
\frac{c + (1-c)x^n}{c + (1-c)x^{n-1}},
\]
ここで $\displaystyle x = \frac{p(1-c)}{1-cp}$.

20. 前述の問題の解は任意の要素を含んでいません。データの外見から、このような状況を予想することはほとんどありませんでした。しかし、これらの問題においては、すべての原因の確率が事前に知られていると仮定されています。この仮定された知識の要素が欠けている場合、最終的な解に任意の常数が必ず現れる可能性があると思われます。この観察は、重要な注目を集めている一連の問題によって一部確認されています。結論として、これらを簡潔に検討します。

天界には非常に密接な双星が多数存在しています。これらのつながりの表れは物理的な根拠があるのかないのか、いずれか一方であるべきです。もし物理的な根拠がない場合、我々は双星現象を恒星が天頂にランダムに配置された結果と考えることができます。つまり、その分布は恒星のどちらのメンバーがどの場所に現れるかが等確率であることを意味します。この仮定が成り立つとし、必要な明るさを持つ恒星の数が分かっている場合、互いの距離内で観測された現象を構成する2つの恒星が存在する確率を求めることができます。ミッチェルは\footnote{Phil. Transactions, An. 1767.}、明るさが$\beta$ カンパルコニイと同等の天界の恒星が230個あると推定し、これらの恒星がランダムに配置された場合、そのような組み合わせが現れる確率は8対1であると算出しました。このような組み合わせが観測されたとき、そのメンバー間には物理的なつながりがある確率を求めることになります。

再び、1801年の年次において、十個の知られた惑星の軌道傾斜角の合計は、フランスの測定によれば $91^{\circ}\cdot$4187 でした。すべての傾斜角が等しく可能な場合、ラプラスは（『解析的確率論』、第276ページ参照）以下のように決定しました：平均的な傾斜角がこの限界内に収まる確率は極めて小さな .00000011235 しかないということです。そして彼は、惑星の軌道傾斜角がこのような狭い範囲内で制限されることに対する配置する原因がある非常に高い確率があると結論付けました。

デ・モラン教授（『メトロポリターニカ大百科』、確率論の記事参照）は、傾斜角の合計を $92^{\circ}$ として、上記の確率を .00000012 とし、「太陽系の形成において惑星の傾斜角がそのようなものである必要がある原因は $1\colon .00000012$」という結論を導き出しました。

観測された一致性から、岩晶石における円偏光の方向とその結晶構造中の特定の斜面の方向との間の一致も同様な決定的な結論が導き出されました（エディンバラ評論、第185号、第32ページ参照。この記事は完全には誤りがないわけではないものの、注意に値するものです）。

これらの問題はみな類似した性質を持っています。ある仮説が立てられ、その仮説から生じ得る様々な可能性のある結果の確率を完全に厳密に定めることが可能です。実際に観測された結果がそのうちの一つと一致し、その結果の仮説上の確率が既知となった場合、求められるのはその仮説自体の確率又はその否定の確率です。

ミッチェルの問題では、仮説は「星のランダムな分布」であるとされます——観測結果としての可能性のある結果は、「近接双星」の出現です。このような結果が非常に小さな確率であることは、仮説の確率も同様に小さなものであることを意味すると考えられ、あるいは少なくとも同じ程度の小さな値であると考えられます。そして、このことから、恒星配置における整理作用の高い確率、そして一部の人々はそれを決定的なものと見做すこともあります。

類似した議論は他の例にも適用されます。

21. どのような形で提示されようと、一般の問題には限りなく解をもたらすもののみが存在する。$x$ を提案された仮説とし、$y$ をその仮説のもとで起こり得る可能性のある現象の一つとして、そしてその確率は仮説が真であるという前提のもとで計算される$p$とする。そして、現象$y$が観察された場合に仮説$x$が真である確率を求めることを要求する。この問題のデータ自体には任意要素の導入が必要であり、しかしながら
\begin{equation}\tag{1}
\operatorname{Prob. } x = a, \quad \operatorname{Prob. } xy = ap;
\end{equation}
と書くことができる。$a$は完全に任意で、0と1の範囲内であることが唯一の制限条件である。そして、求める条件付き確率を$P$とするとき
\begin{equation}\tag{2}
P=\frac{\operatorname{Prob. } xy}{\operatorname{Prob. } y} =\frac{ap}{\operatorname{Prob. } y}.
\end{equation}
が成り立つ。次に$\operatorname{Prob. } y$を決定する必要がある。

%-----------------------File: 169.png----------------------------

$t = xy$と置くと
\begin{equation}\tag{3}
y = \frac{t}{x}=tx + \frac{1}{0}t(1-x)+0(1-t)x+\frac{0}{0}(1-t)(1-x).
\end{equation}

ここで$\operatorname{Prob. } x = a$、$\operatorname{Prob. } t = ap$であることから、論理から代数へと移行すると
\begin{equation*}
\operatorname{Prob. } y = \frac{tx + c(1-t)x}{tx+1-t},
\end{equation*}

そして関係式として
\begin{equation*}
\frac{tx + (1-t)x}{a}=\frac{tx}{ap}=tx+1-t.
\end{equation*}

これらから容易に$\operatorname{Prob. } y$を求めることが可能であり、それは
\begin{equation}\tag{4}
\operatorname{Prob. } y = ap + c (1-a).
\end{equation}

次に(3)に戻ると$c$は、$(1-t)(1-x)$が発生した場合に現象$y$が発生する確率であることがわかる。しかし
\begin{equation*}
(1-t)(1-x)=(1-xy)(1-x)=1-x.
\end{equation*}

よって$c$は、仮説$x$が発生しない場合に現象$y$が発生する確率である。

$\operatorname{Prob. } y$の値を(2)に代入すると、以下の定理を得る。

任意の現象 \( y \) が、仮定された物理的仮説 \( x \) 上で確率 \( p \) である場合、その現象が観察されたときに物理的仮説の \textit{\`{a} posteriori} 確率 \( P \) は以下の式で表されます：
\[
P = \frac{ap}{ap + c(1 - a)}
\]
ここで、\( a \) と \( c \) は任意の定数であり、前者は仮説の \textit{\`{a} priori} 確率を表し、後者は仮説が偽である場合に現象 \( y \) が発生する確率を表します。

上記の定理から導かれる主な結論は、他の条件が同じであるとき、\( P \) の値は \( p \) の値と共に増加し、減少することです。したがって、仮説が仮定されたときに現象の確率が大きいほど、または小さいほど、その現象が観察されたときに仮説の確率も大きいほど、または小さいほどになります。\( p \) が非常に小さい場合、通常 \( P \) も小さくなります。ただし、\( a \) が大きいか、または \( c \) が小さい場合はこの限りではありません。

したがって、第二に、仮説が仮定されたときに現象の確率が非常に小さい場合、その現象が観察されたときに仮説の確率も非常に小さいです。ただし、仮説の \textit{\`{a} priori} 確率 \( a \) が大きいか、または他の仮説上での現象の確率が小さい場合はこの限りではありません。

公式（5）は、様々な場合において正確な検証が可能であり、例えば$c = 0$のときや$a = 1$のとき、または$a = 0$のときなどにそのような検証が行われる。しかし、$a$と$c$の値を決定する方法がない限り、公式（5）は特定の値$P$を与えることにはならないことが明白である。この目的を達成しようとする任意の解法は、原則的に誤りであるか、もしくは上記の任意要素に関する暗黙の仮定が含まれている。

デ・モラン氏によるラプラス問題の解決策、すなわち惑星軌道の傾きが黄道面に限定された狭い範囲内で保たれる理由がある場合、そのような傾きの存在確率$p = 0.00000012$であるという解法は、私には後者のタイプに属すると考えられる。彼は、各軌道における傾きが等しく確率的であると仮定した場合、全体の傾きの和が$92^\circ$未満となる確率$p = 0.00000012$を見つけ出し、次のように述べている：「もし傾きがそのようにある理由があるなら、その事象の確率は$1$である。したがって、太陽系の形成において傾きが現在のものである必要のある原因があったという確率は$1 \colon 0.00000012$（すなわち$1 \colon p$）である。」この結果は、$p$に上記の値を割り当て、$c = 1$, $a = \frac{1}{2}$と仮定した場合に公式（5）が実際に与えるものと同じである。

ここで、$P$は各傾きが等しく確率的であるという事後的な確率を表すとする。すると$1 - P$はそのような状況ではない確率であり、デ・モラン氏の選択肢によれば、決定的な原因がある確率である。したがって、方程式（6）はデ・モラン氏の結果と一致する。

\begin{gather*}
  P = \frac{\frac{1}{2}p}{\frac{1}{2}p + \frac{1}{2}} = \frac{p}{1+p}   \\
  \therefore 1 - P \colon P \colon\colon 1 \colon p.    \tag{6}
\end{gather*}

22. しかし、$a$ および $c$ に特定の値を割り当てる正当性はありますか？
私はその正当性が認められないという見解を持っています。
%-----------------------File: 171.png----------------------------
この問題自体における重要性よりも、その将来的な影響の方がより重要な点です。確率論の既存の応用では、任意定数は明示的に現れていません。しかし、上記の例や最高の権威によって承認された多くの他の例においては、これらの定数に対する仮定的な決定が試みられています。そしてこの状況は、理論の結果に、特に因果関係に関する問題に関して、明確で精密な特徴を与えました。一方では、これは数値が世界を統治し、その権限と領域を拡大するという観念を強調しました；他方では、推測が唯一の基礎となる領域への数値の侵入に対する強い反対意見を引き起こしました。このような問題をこの手法で解く際に現れる任意定数の存在自体は、明確な解決が不可能であることを示唆し、調査が停止すべき点を指摘しています。私たちはこれらの定数の意味を解釈する手段を持っていますが、そのような経験は完全に無視される必要があるような経験と同様に、私たちには及ばないものであることが明白です。

これらの問題に関連するもう一つの困難は、おそらく私たちの知覚能力そのものの構造に内在しているものであり、それが何かを正確に定義することです。秩序という原則の現象については、非常に複雑な例以外では、それらを検出することは難しくありませんし、法則に基づく作用要因がその基礎であると推測することにも躊躇しません。しかし、その秩序に数値的な基準を設定しようとすると、無駄な努力とも言えるでしょうし、甚至は傲慢な試みかもしれません。しかしながら、確かな確率を weighs するためには、この試みが行われなければなりません。

異なる宇宙の構成形態の可能性を正確に評価し、問題に対する明確な解を確立するための要素を決定するためにです。

23. 確率論において不十分なデータに基づく問題の解法における必要に迫られた任意性を回避しようとする最も一般的な方法は、実際の確率が未知である要素に対してすべての可能性のある確率度を与え、これらの確率度が等しく確率的であると仮定することです。そして、観察された現象の多くの原因としてこれらを考慮に入れ、いくつかの互いに排他的で等しく確率的な原因のうちの一つによる効果を表す定理（問題9）を適用します。

例えば、一定期間の暗闇後に太陽が昇るという観察が $m$ 回連続して行われた場合、そのような状況下で再び昇る確率は受け入れられた原理に基づいて次のように決定されます。未知の確率 $p$ が0と1の間にあるとき、暗闇後の太陽が昇る確率が $p$ から $p + dp$ の範囲に存在する確率を $c$（無限小で定数）とする。すると、太陽が連続して $m$ 回昇る確率は
\[
  c \int_0^1 p^m \, dp;
\]
となります。そして、太陽がこれを行う確率と再び昇る確率（つまり、連続して $m + 1$ 回昇る確率）は
\[
  c \int_0^1 p^{m+1} \, dp,
\]
となります。

したがって、太陽が連続して $m$ 回昇った場合に、次回も昇る確率は
\[
  \frac{c \int_0^1 p^{m+1} \, dp}{c \int_0^1 p^m \, dp} = \frac{m+1}{m+2},
\]
と知られており、一般的な解となります。

上の解法は、黒球と白球の混合物を含むurnからボールを取り出す問題との仮定された類似性に基づいて通常構築されます。そして、取り出されるボールが置き換えられない有限個の場合や、無限個の場合で取り出されるボールが置き換えられるか否かに関わらず、すべての数的比率が等可能性を想定されることに注意してください。

特に、$m$回連続して白球が取り出された場合、$(m+1)$回目の取り出しで白球が出る確率は
\[
  \frac{m + 1}{m + 2}
\]
となります（参考文献：エディンバラ哲学学会誌第20巻第4部、テロット主教の論文）。

上記と似た応用における原理は、我々の知識や、正確には無知の等しい分布というものです。つまり、何も知らない状態について等しい確率を割り当てるということです。しかし、私はこの方法が任意的な手続きであると考えます。異なる仮定が同じ最終結論に至る例も存在し、そのような例は既に提示されています。しかしながら、そのような例は例外的であり、特定の問題に関しては引用された論文で示されています。

特に問題に関連しては、ボールが有限個であり置き換えられないという仮定では異なる結論を導く別の仮定が存在することが証明されます。また、「知識や無知の等しい分布」を厳密に含む他の仮定も存在し、これらは互いに矛盾した結果を導くことが容易に理解できます。

24. たとえば、太陽が昇るという状況を、袋の中に無限個のボールがあり、そのすべてが黒か白のいずれかであるとするとき、白いボールを引き出すことを表すものとして考えましょう。そして、仮定された原理は「システム内のボールのすべての可能な構成は等確率である」というものです。ここで、「構成」とは、システム内の各ボールに特定の色（黒または白）が割り当てられた配置を指します。それから、$m$回の引き出しで$m$個の白いボールが引かれた場合、$(m+1)$回目の引き出しで白いボールが引かれる確率を求めましょう。

まず、ボールの数を$\mu$と仮定し、次の方法で符号$x_1, x_2, \dotsc x_\mu$にそれぞれ割り当てます。$x_i$は$i$番目のシステム内のボールが白いという事象を表し、その状態を宣言する命題は$x_i = 1$です。同様に、符号$1-x_i$は$i$番目のボールが黒であることを表します。明らかに、すべての符号$x_1, x_2,\dotsc x_\mu$からなる各構成要素は、システム内のボールの黒さと白さに関するすべての可能な構成を同様に表し、そのような構成の数が$2^\mu$であるため、仮定に従って各構成の確率は$\dfrac{1}{2^\mu}$となります。これは、任意の構成要素に対する符号$x_1, x_2,\dotsc x_\mu$をそれぞれ$\dfrac{1}{2}$で置き換えることで求められる値です。したがって、上記の説明に従った系列の各事象の確率は、その表現に符号$x_1, x_2,\dotsc x_\mu$をそれぞれ$\dfrac{1}{2}$で置き換えることで求められます。

現在、$\mu$が大きいほど、どのボールも一度引いて返し入れた後で再度引かれる確率は低くなります。$\mu$が無限大に近づくにつれて、この確率は0に近づきます。そしてその場合、実際に引かれたボールの状態は$m$個のシンボル$x_1, x_2, \dotsc x_\mu$の論理関数として表現できます。したがって、展開することで、これらの$m$個のシンボルの成分の系列となります。したがって、その確率は未展開または展開された形での各シンボルに$\dfrac{1}{2}$を代入することにより求められます。しかし、これは、各回引くたびに白いボールが出る確率が事前に$\dfrac{1}{2}$であることが分かっている場合に必要なものであり、十分なものです。

したがって、ボールの数が無限大で、システムのすべての構成が等しく確実である場合、連続して$m$個の白いボールを引く確率は$\dfrac{1}{2^m}$となります。そして、連続して$m + 1$個の白いボールを引く確率は$\dfrac{1}{2^{m+1}}$となります。

したがって、$m$個の白いボールを引き続けた後で次の引くものがまた白いボールである確率は$\dfrac{1}{2}$となります。つまり、過去の経験はこの場合未来の期待に影響を与えないということです。

25. この結果を通常の方法で確認することは満足感があります。これを達成するためには、以下のように探求します：

まず：袋の中に$\mu$個のボールがあり、各回引いた後でボールを返し入れる条件下で$p$試行中に$r$個の白いボールと$p-r$個の黒いボールを引く確率を求めます。システムのすべての構成が事前に等しく確実である場合。

次に：$\mu$が無限大になるとき、この確率はどのようになりますか？

さらに：連続して$m$個の白いボールを引き続けた後で、次の引くものがまた白いボールである確率を求めます。

urn 内に含まれる \(\mu\) 個の球がそれぞれ白球か黒球であるものとして、各球を試験後に入れ戻し、すべての構成が等確度であるとき、\(p\) 回の試験で \(r\) 個の白球と \(p-r\) 個の黒球を引く確率は、urn 内に白球が 0 個、1 個、---最後に \(\mu\) 個あるという個別の仮説における同じ結果の確率の和に等しい。したがって、urn 内に \(n\) 個の白球があるという仮説におけるこの結果の確率の和であり、\(n\) は 0 から \(\mu\) まで変化する。

現在、$n$個の白球があると仮定すると、一度の引く作業で白球を引く確率は $\frac{n}{\mu}$ であり、$p$回の引く作業中 $r$ 回白球、$p-r$ 回黒球を特定の順序で引く確率は
\[
  \left( \frac{n}{\mu} \right)^r
  \left( 1 - \frac{n}{\mu} \right)^{p-r}
\]
となります。しかし、このような順序が $p$ 個のものから $r$ 個を選ぶ組合せ数と等しいため、$\mu$ 個の球中 $n$ 個が白球であるシステムから $r$ 回の引く作業で白球を $r$ 個引く総確率は
\[
  \frac{ p(p-1) \dotsc (p-r+1) }{ 1 \cdot 2 \dotsc r }
  \left( \frac{n}{\mu} \right)^r
  \left( 1 - \frac{n}{\mu} \right)^{p-r} \tag{1}
\]
となります。

また、$\mu$ 個の球から $n$ 個が白球であるような構成数は
\[
  \frac{\mu(\mu-1) \dotsc (\mu-n+1)}{1 \cdot 2 \dotsc n},
\]
であり、可能なすべての構成数は $2^\mu$ です。したがって、ちょうど $n$ 個が白球である確率は
\[
  \frac{\mu(\mu-1) \dotsc (\mu-n+1)}{1 \cdot 2 \dotsc n 2^\mu}
\]
となります。

式 (1) をこの表現で掛けて、$n = 0$ から $n = \mu$ まで総和を取ると、
\[
  \frac{p(p-1) \dotsc p-r+1}{1 \cdot 2 \dotsc r}
  \sum_{n=0}^{n=\mu}
  \frac{\mu(\mu-1) \dotsc (\mu-n+1)}{1 \cdot 2 \dotsc n 2^\mu}
  \left( \frac{n}{\mu} \right)^r
  \left( 1-\frac{n}{\mu} \right)^{p-r} ,   \tag{2}
\]
が得られます。これは、$\mu$ 個の球から全ての構成が等確率であるシステムにおいて $p$ 回の引く作業で $r$ 個の白球が出る総確率を表します。

次に
\begin{gather*}
  \sum_{n=0}^{n=\mu}
  \frac{\mu(\mu-1) \dotsc (\mu-n+1)}
       {1 \cdot 2 \cdot n \cdot 2^\mu}
  \left( \frac{n}{\mu} \right)^r
  \left( 1-\frac{n}{\mu} \right)^{p-r}
\\
= \sum_{n=0}^{n=\mu}
  \frac{\mu(\mu-1) \dotsc (\mu-n+1)}
       {1 \cdot 2 \cdot n \cdot 2^\mu}
  \left( \frac{n}{\mu} \right)^r
  \left( 1-\frac{n}{\mu} \right)^{p-r}
  \varepsilon^{n\theta} \dotso (\theta = 0)
\\
= \frac{1}{2^\mu}
  \left( \frac{D}{\mu} \right)^r
  \left( 1-\frac{D}{\mu} \right)^{p-r}
  \sum_{n=0}^{n=\mu}
  \frac{\mu(\mu-1) \dotsc (\mu-n+1)}
       {1 \cdot 2 \cdot n}
  \varepsilon^{n\theta}
\\
= \frac{1}{2^\mu}
  \left( \frac{D}{\mu} \right)^r
  \left( 1-\frac{D}{\mu} \right)^{p-r}
  (1 + \varepsilon^\theta)^\mu,        \tag{3}
\end{gather*}
となります。ここで $D$ は $\frac{d}{d\theta}$ のシンボルを表し、$\phi(D) \varepsilon_{n\theta} = \phi(n) \varepsilon^{n\theta}$ です。

既知の定理により
\begin{gather*}
  t^m = 1 + \Delta 0^m t + \frac{\Delta^2 0^m}{1 \cdot 2} t(t-1)
      + \frac{\Delta^3 0^m}{1 \cdot 2 \cdot 3} t(t-1)(t-2).%*
\\
  \therefore D^m (1+\varepsilon^\theta)^\mu
= \{1 + \Delta 0^mD + \frac{\Delta^2 0^m}{1 \cdot 2} D(D-1)
    + \text{ 等}\} (1+\varepsilon^\theta)^\mu.                   %*
\end{gather*}
第二項において $\varepsilon^\theta = x$ と置くと
\[
  D^m (1+\varepsilon^\theta)^\mu
= (1 + \Delta 0^m x\frac{d}{dx}
    + \frac{\Delta^2 0^m}{1 \cdot 2} x^2 \frac{d^2}{dx^2}
    + \text{ 等}) (1+x)^\mu,                                     %*
\]
となります。ここで
\[
  D(D-1) \dotsc (D-i+1) = x^i \left( \frac{d}{dx} \right)^i.
\]

これらの訂正は既に適用されているようです。
3、5、および6番目の誤りについては驚いており、これらが何を指しているのか分かりません。

---

**補遺。--- 3、5、および6番目の式の下部から、1を0m に読み替えてください。**
-----------------------ファイル：177.png----------------------------

上記の等式的第二項において微分を行い、x = 1（θ = 0）とすると、次のようになります。
\[
D^m(1 + \varepsilon^\theta)^\mu
= \mu(\Delta 0^m)2^{\mu-1}
+ \frac{\mu(\mu-1)}{1 \cdot 2}(\Delta^2 0^m)2^{\mu-2}
+ \text{ 等}
\]

この等式的第二項の最後の項は以下のようになります。
\[
\frac{\mu(\mu-1)\cdots (\mu-m+1)\Delta^m 0^m}{1 \cdot 2 \cdots m}2^{\mu-m}
= \mu(\mu-1)\cdots (\mu-m+1)2^{\mu-m};
\]
ここで、\(\Delta^m 0^m = 1 \cdot 2 \cdots m\) です。μが大きな値であるとき、この項は他のすべての項よりも大きい値をとります。また、μが無限大に近づくにつれて、他の項と比べて無限大に近づきます。さらに、この項は μm2{\mu-m} の形を取りますので、極限を取ると
\[
D^m(1 + \varepsilon^\theta)^\mu = \mu^m 2^{\mu-m}
= \left(\frac{\mu}{2}\right)^m 2^\mu.
\]
したがって、シンボル D の任意の関数 φ(D) を表す場合、D の昇べき級数展開が可能であるとき、次のようになります。
\[
\phi(D)(1 + \varepsilon^\theta)^\mu = \phi\left(\frac{\mu}{2}\right)2^\mu,  \tag{4}
\]
θ = 0 および μ = ∞ の場合です。厳密には、これは上記の等式的両辺の比が μ が無限大に向かって増加するにつれて等式の状態に近づくことを意味します。

この定理により、(3) の最後の項は以下の形になります。
\[
\frac{1}{2^\mu} \left(\frac{1}{2}\right)^r
\left(1 - \frac{1}{2}\right)^{p-r} 2^\mu
= \left(\frac{1}{2}\right)^p.
\]
したがって、(2) は以下の式を導きます。
\[
\frac{p(p-1)\cdots (p-r+1)}{1 \cdot 2 \cdots r} \left(\frac{1}{2}\right)^p,
\]
これは、無限個の黒球と白球が含まれる urn から、すべての構成が等確率である場合に、p 回の引き出しで r 個の白球が出る確率を表します。

したがって、$p = m$, $r = m$ とすると、$m$ 回の引きで全て白い球が出る確率は $\left( \frac{1}{2} \right)^m$ であり、さらに $m+1$ 番目の引きでも白い球が出る確率は $\left(\dfrac{1}{2}\right)^{m+1}$ となる。したがって、最初の $m$ 回の引きで全て白い球が出た場合に、$m+1$ 番目の引きも白い球が出る確率は
\[
  \left( \frac{1}{2} \right)^{m+1} \div
  \left( \frac{1}{2} \right)^m = \frac{1}{2}
\]
となる。一般に、任意の提案された結果は、各特定の引きが白い球か黒い球かという確率が等しい偶数の確立と同じ確率を持つことになる。これは先に得られた結論と一致する。

26. これらの結果は、欠如したデータを仮定によって補うとき、解は一般に仮定の性質により異なることを示している。従って、理論確率の原理がどの程度まで我々をこのような仮定の選択へ導くべきかという問題は、より明確な形で残されている。私は既に、これらの原理がそうではないと確信していることを述べたことがあるし、それ以外にもそのような確信を強める理由がある。つまり、この著作の方法によって特定の問題に対する解を見つけることができた場合でも、あるデータ（例えば Prob. $x = p_1$）が省略されたときには同様の方法で別の明確な解を得ることもしばしば可能である。しかし、第十九章に定められた範囲内で $p$ を変数と仮定した場合の積分によって最初の解から第二の解を導く方法を見つけることはできなかった。この推論は、第23条で言及された原理が有効である場合に可能になると考えられる。しかし、それでも私は一般的な数学的な見解とは異なる点について懐疑的であり、特にラプラスの精神と手法を最も深く理解した英語の著者の一人から異論を唱えることにためらいを感じる。しかしながら、確率理論において二番目に重要な問題として評価されるべき質問が慎重に検討されることを希望する。

第二十一章

---

以上の内容は「%CHAPTER XXI.」というタイトルのみでしたので、章の名前や内容については推測せざるを得ません。具体的な章題や内容を含む原著英文が提供された場合、より正確な翻訳を行うことができます。ご提供いただければ幸いです。

\newpage

\chapter{判断・陪審・証言の確率への特殊応用（Probability of Judgements）}
\label{ch:xxi}

\begin{center}
\large\textit{裁判官・陪審員の合議・複数証言の信憑性と冤罪確率}
\end{center}
\vspace{4mm}

\#\# 判決の確率に関する特定の適用：前章で十分に説明された本論文の一般方法への実用的な応用

1. 本論文で理論確率問題に対する解法として提示した一般的な方法が、前章で充分に説明されていると仮定すれば、ここではその方法を社会統計学の広大な分野から選りすぐった一つの実用的な応用例に適用することを考える。すなわち、判決の確率の推定である。この応用が直近の結果によって評価すれば、最良のものとは言えないかもしれない。しかし、経験データだけでは最も重要な調査目的を達成するためには必要な仮説の必要性を明確に示し、任意の仮説を仮定した場合の推論結果を厳格に導き出すことができることで、この方法はその範囲内において全ての役割を果たす。また、人間の行動に関連する問題では、仮説がより正当化される傾向があると考えられる。前章末尾の節で述べられたような他の問題とは異なり、我々の人間性に関する経験的な知識は統計記録の不足と不完全さを補完する役割を果たす。

2. 犯罪裁判に関連する問題における要素は以下の通りである：

1. 特定の陪審員が事案について正しく判断する確率。
2. 被告者が有罪である確率。
3. 有罪判決が下されるか、無罪放免となる確率。
4. その有罪判決または無罪放免が公正である確率。
5. 陪審員の構成。
6. 経験から得られるデータ、例えば一致した決定を達成した相対的な事例数や特定の多数決を得た事例数、上級裁判所によって覆された決定の数など。

また、現在考慮している問題のクラスは直接的または逆方向的なものと見なすことができます。確率に関する直接的な問題は、裁判員の一人ひとりが正しい判断を行う確率や被告が有罪である確率が事前に知られている場合、そして特定の種類の決定あるいは明確な多数決の確率を求められる場合を指します。逆問題は、経験から得られたデータから原因と結果の関係にある要素を決定する必要があるもので、直接観察することはできないものであり、例えば裁判所の判決記録から、裁判員が正しい判断を行う確率を求めるような問題です。このような種類の問題は全体の中で最も困難かつ重要なものであるため、ここでは特にそのことに注目します。

3. 以上の列挙に含まれる直接的な問題を解くことは難しくありません。仮に一人だけの裁判員がいる場合を考えましょう。$k$ を被告が有罪である確率、$x$ を裁判員が正しい意見を持つ確率、$X$ を被告が判決で有罪となる確率とします。すると、

\begin{align*}
  kx &= \text{被告が有罪であり、裁判員が彼を有罪とする確率}\\
  (1-k)(1-x) &= \text{被告が無罪であり、裁判員が彼を有罪とする確率}
\end{align*}

これらの状況が判決で有罪となる唯一の場合であり、さらに互いに排他的であるため、

\[
  X = kx + (1-k)(1-x) \tag{1}
\]

類似の方法で、$n$人の陪審員がそれぞれ正しく判断する確率が$x_1, x_2, \dotsc, x_n$である場合、一致した有罪判決を下す確率$X$は以下のようになります：
\[
  X = k x_1 x_2 \dotsm x_n + (1-k)(1-x_1)(1-x_2) \dotsm (1-x_n).
\]
%-----------------------File: 181.png----------------------------
したがって、それぞれの確率$x_1, x_2, \dotsc, x_n$が等しく、各々を$x$で表す場合、以下の式になります：
\[
  X = kx^n + (1-k)(1-x)^n.   \tag{2}
\]
後者の場合における被告が有罪である確率は次のようになります：
\[
  \frac{kx^n}{kx^n + (1-k)(1-x)^n}
\]
これらの結果は、$k, x_1, x_2, \dotsc$を表す確率のイベントが独立であるという仮定に基づいています。この仮定は、当該事件が唯一与えられた確率を持つイベントであり、それら以外のイベントは考慮されないという事実から導かれています。

特定の票数による有罪判決の確率も同様の原理で求めることができます。陪審員が3人で構成されている場合、それぞれの正しく判断する確率が$x, x', x''$であるとします。被告が2人のうちの1人によって有罪と宣告される確率$X_2$は以下のようになります：
\[
  \begin{split}
    X_2 = k \{xx'(1-x'') + xx''(1-x') + x'x''(1-x)\}   \\
        + (1-k)\{(1-x)(1-x')x'' + (1-x)(1-x'')x' + (1-x')(1-x'')x\},
  \end{split}
\]
$x = x' = x''$の場合は、以下の式になります：
\[
  3kx^2(1-x) + 3(1-k)x(1-x)^2.
\]

同じ論理の方法により、以下のように示されます。ある陪審員数が \(n\) 人の陪審団から \(i\) 個の声で有罪と宣告される被告の確率
\begin{equation}\tag{3}
X_i = \frac{n(n-1) \cdots (n-i+1)}{1 \cdot 2 \cdots i} \left\{ kx^i(1-x)^{n-i} + (1-k)x^{n-i}(1-x)^i \right\}.
\end{equation}

一定の多数 \(a\) による有罪判決の確率を求めたい場合、次のようにします：
$$
i - a = n - i,
$$
これより
$$
i = \frac{n+a}{2},
$$
この値は上記の式に代入する必要があります。もちろん、\(a\) の値は \(i\) が整数となるように制限されます。もし \(n\) が偶数であれば、その値は 0, 2, 4, 等 であり、奇数であれば 1, 3, 5, 等 となります。

多数の判決（少なくとも \(m\) 個の声）による有罪判決の確率は、上記のように求めた以下の確率を加算することで得られます：

\[
\begin{minipage}{.8\linewidth}
1つ目：正確な多数 \(m\) による有罪判決の確率；
\end{minipage}
\]

\[
\begin{minipage}{.8\linewidth}
2つ目：次に大きい多数 \(m + 2\) による有罪判決の確率；
\end{minipage}
\]
そして、その系列の最後要素は全員一致の有罪判決の確率です。したがって、12人の陪審員から少なくとも4人の声で有罪と宣告される確率は、

$$
X_8 + X_9 \cdots + X_{12},
$$

これらの項の値は (3) 式において \(n = 12\) に代入することで得られます。

4. 裁判所ではなく、単なる議論集団（\(n\) 人の参加者）を考慮する場合、それぞれの個々の正誤判断確率が \(x\) で表されるとき、上記の式は \(k\) の項を取り除いてやや簡略化された形になります。

一致の確率は以下の式で表されます：
\[
X = x^n + (1-x)^n.
\]

全体の中から$i$人の合意の確率は次の式で表されます：
\[
X_i = \frac{n(n-1)\cdots (n-i+1)}{1 \cdot 2 \cdots i} \{x^i(1-x)^{n-i} + x^{n-i}(1-x)^i\}.
\]

このような調査については、さらなる詳細な説明は不要です。 Condorcet, Laplace, Poissonなどの他の著者は、$n$と$i$が大きな数である場合に必要な計算方法や簡略化の手順を詳しく調査してきました。全般的には、この研究は非常に推測的なものであり、$x$と$k$の値は直接観察によって決定されるわけではありません。

これらの値は一般的には$\frac{1}{2}$を超えることが予想されます。また、公的知識が進展するにつれて$x$は増加し、一方で$k$は犯罪者を国際裁判所に提出するための法的手続きにより大きく依存すると考えられます。Poissonは、革命期（フランスの恐怖時代など）において、政治的な罪を含むと仮定すると$k$の値が$\frac{1}{2}$未満になる可能性があると指摘しています。

近現代ヨーロッパの歴史は、この観察を支持し、無差別な民主主義だけではなく、他の要因から無罪の主張を懸念する必要性も示唆しているかもしれません。

Laplaceは、$x$の値が以下の範囲内である場合すべて等確率であると仮定しました：
\[
x = \frac{1}{2} \text{ から } x = 1.
\]
彼は、一人ひとりの判断の正確性の確率が制限されている範囲内で、$x$のどの値を優先する理由がないと仮定しました。この仮説は完全に任意であり、その結果についてはここで検討することは無意味です。

ポワソンは、経験から推論的に$x$と$k$の値を導き出すことを最初に試みたようです。1825年から1830年までの6年間（含む）において、フランスの裁判所前に犯罪に対する罪状で告訴された個人の数は11,016人であり、そのうち判決で有罪となった者は5,286人でした。陪審員は各12人で構成されており、判決は単純多数で下されるものとしました。これらの数字が確率を推定するのに十分に大きいと仮定すると、告訴された個人が単純多数の賛成により有罪となる確率は分数$\dfrac{5286}{11016}$または約$0.4782$となります。したがって、以下の等式が得られます：
\[
X_7 + X_8 \dotsc + X_{12} = 0.4782 \tag{5}
\]
%-----------------------File: 184.png----------------------------
ここで、$X_i$の一般表現は(3)を$n=12$に置き換えたものとします。

1831年の法律が改正され、有罪判決には少なくとも4人の賛成が必要となったため、その年における犯罪に対する罪状で告訴された個人の数は2,046人であり、そのうち有罪となった者は743人でした。したがって、単純多数により個々人が有罪となる確率は$\dfrac{743}{2046}$または約$0.3631$となります。これより以下の等式が得られます：
\[
X_8 + X_9 \dotsc + X_{12} = 0.3631 \tag{6}
\]

前六年間と同様の値であると仮定すると、これらの二つの等式（5）と（6）は$k$と$x$の値を近似して決定するのに役立ちます。ポワソンは次のように求めました：
\[
k = 0.5354, \quad x = 0.6786
\]

同じ期間における財産犯罪についても、類似の分析により彼は以下のような値を見つけました：
\[
k = 0.6744, \quad x = 0.7771
\]

等式（5）と（6）の解はそれぞれ2つの$k$の値と2つの$x$の値を導き出し、各ペアの一つが$\dfrac{1}{2}$よりも大きい値で、もう一つがそれよりも小さい値となります。これらの値について、各ケースにおいて大きい方の値を選択することが仮定されました。これは告訴された被告が有罪である可能性が高いと考えられ、また陪審員が証拠に基づく正解を形成する可能性が高いと考えられるからです。

5. ポイソンが用いたデータ、特に1831年以前のものは明らかに不完全で、上記のxとkの決定値に対して大きな信頼性をもたらすことは難しい。これらのデータは方法論的な観点から主に紹介されるに過ぎない。1825-30年の6年間やそのような期間において、各可能な投票多数による処刑判決の数を記録することは可能だった。政策的理由がない限り、$X_8, X_9, \dotsc X_{12}$ の各要素の値は正確に測定されることが可能だった。このようにして、ポイソンの問題が特定の一例である一般的な問題の概念が浮かび上がる。このような表面上余分なデータから、xとkの値をどのように決定することができるだろうか？

ポイソンや他の著者が採用した確率$x$と$k$に対応する事象間の絶対的独立性という仮説が維持される場合、5つの方程式のタイプ（3）からなるシステムを形成し、これらの方程式の中から最も有利と思われる特定の組を選択するか、または最小二乗法によりこのシステムの方程式を組み合わせることができる。xとkの値が与えられた方程式の異なる組から得られるものと非常に異なる場合、後者の方法が厳密に適用可能かどうかについて疑問が生じる可能性がある。

M. コルヌトは、4人の判事で構成される裁判所における個々の投票記録から、各判事が正しい判決を出す確率を調査する問題を考えた。要素$x, x', x'', x'''$ の決定のために彼は8つの方程式を得ており、これらを2つの4つの方程式のセットに分けており、x, x', x'', x''' の値がこれらのセットから得られるものと非常に異なる場合、判事たちの意見の独立性の仮説が特定の場合において支持されないことを示す可能性があるとしている。この調査方法の原理は（XVIII.4）で触れられている。

6. 以下に述べる問題のクラスに対して、本論文の方法を適用し、まず、裁判所や議会の記録のみが個々のメンバーの判断の正確性に関する確率についてどのような情報も提供するかどうかを調査します。その結果、それらの記録から個々の判断の正確性に関する明確な結論を導くことは不可能であることが明らかになった場合、第二に必要な仮定の種類と量は何であるかを検討します。

\centerline{\textsc{命題I}}

裁判所や議会の決定の記録のみから、個々のメンバーの判断の正確性に関する明確な結論を導くことは不可能です。

この命題は、未加味の良識の信念を表現しているように見えるかもしれませんが、その正当性を厳格に証明することも興味深いことであることが明らかになります。
%-----------------------File: 186.png----------------------------

議会が $n$ 人のメンバーから成る場合を考えます。また、彼らの判断が相互に依存するか否かについては仮定しません。論理的記号 $x_1, x_2, \dotsc, x_n$ を以下のように定義します：一般的な符号 $x_i$ は、第 $i$ のメンバー、すなわち議員 $A_i$ が正しい意見を述べたという事象を表します。Prob. $x_1$, Prob. $x_2, \dotsc, $Prob. $x_n$ の値を考える問題の解答（\textit{qu{\ae}sita}）として、その可能性のあるデータの表現について次に調査する必要があります。

現在のデータは、符号$x_1, x_2, \dotsb, x_n$を特定の論理関数で表現できる事象の確率です。これらの論理関数を$X_1, X_2, \dotsb, X_m$と表し、実際のデータ系は以下のようになります：
\[
  \text{Prob. }X_1 = a_1,\quad
  \text{Prob. }X_2 = a_2,\quad
  \text{Prob. }X_m = a_m.
\]
この状況から、$x_1, x_2, \dotsb, x_n$がそれぞれ$1-x_1, 1-x_2, \dotsb, 1-x_n$に置き換わった場合でも、$X_1, X_2, \dotsb, X_m$は不変であることが示されます。例えば、投票が全員一致だったという事実が記録されている場合、その確率が十分に高いと仮定して決定される事象は論理関数
\[
  x_1 x_2 \dotsm x_n + (1-x_1)(1-x_2)\dotsm (1-x_n)
\]
で表現されます。この関数も上記の条件を満たします。

また、ある個人（例：$A_1$）の投票が他のメンバー全員と異なるという事実が記録されている場合、その確率が与えられる事象は論理関数
\[
  x_1(1-x_2)\dotsm (1-x_n) + (1-x_1)x_2\dotsm x_n
\]
で表現されます。この関数も上記の条件を満たします。

したがって、意見の一致や相違は真実と誤りのいずれかに一致する可能性があり、意見が分かれている場合でもどちらの側も正しい可能性があるため、議会での特定の状態（一貫性または不一致）を表現するのは符号$x_1, x_2, \dotsb, x_n$の論理関数であり、これには否定的な符号
\[
  1-x_1, 1-x_2, \dotsb, 1-x_n
\]
も含まれます。議会の記録では、どの意見群が正しいかは明示されません。したがって、$X_1, X_2, \dotsb, X_m$は上記の種類のみであると推定されます。

7. 次に、一般的な方法に従ってProb. \(x_1\) の値を決定するには、まず関数 \(X_1, \dotsc X_m\) を新しい符号 \(t_1, \dotsc t_m\) に等しくします。これらの方程式
\[
  X_1 = t_1, X_2 = t_2, \dotsc X_m = t_m,
\]
から、シンボル \(x_2, x_3,\dotsc x_n\) を消去し、その後、\(t_1, t_2, \dotsc t_m\) の符号を用いて与えられた確率を持つ事象を表す開発された論理関数として \(x_1\) を決定します。上記の消去結果は
\[
  Ex_1 + E'(1-x_1) = 0;    \tag{1}
\]
と表され、ここで \(E\) および \(E'\) は \(t_1, t_2, \dotsc t_m\) の関数です。すると
\[
  x_1 = \frac{E'}{E'-E}.   \tag{2}
\]

現在、関数 $X_1, X_2, \dotsc, X_m$ は符号 $x_1, \dotsc, x_n$ および $1-x_1, \dotsc, 1-x_n$ に関して対称である。したがって、方程式において $E'$ は $E$ と同一であることが明らかである。
従って (2) 式から
\[
x = \frac{E}{0}
\]
を得る。また、上記方程式の右辺の展開における唯一の係数は $\frac{0}{0}$ および $\frac{1}{0}$ であることが明白である。前者は、ある成分の係数を決定する際に与えられた $t_1, \dotsc, t_m$ の値が $E = 0$ を満たすときに出现在する。後者はそれ以外のすべての場合において出現在し、等価な結果をもたらす。したがって、次の形で展開を表すことができる：
\[
x_1 = \frac{0}{0}C + \frac{1}{0}D    \tag{3}
\]
ここで $C$ および $D$ は符号 $t_1, t_2, \dotsc, t_m$ の成分または成分の集合である。

次に論理から代数へと移行し、
\[
\text{Prob. }x_1 = \frac{cC}{C} = c,
\]
となる。一般ルール (XVII. 17) の関数 $V$ は現在の場合において $C$ に簡約される。したがって Prob. $x_1$ の値は、0 や 1 を超えないという条件を除いて任意であることが明白である。Prob. $x_2, \dotsc, x_n$ の個々の値も同様に任意である。従ってこれら任意の値が互いに関連付けられている必要条項がないわけではないことが推察される。しかし、このような条件を調査することは、第十九章の方法に適切に属するものである。

最後の論理方程式に戻り、$c$ の解釈を求めると、元々の問題を再述するだけになります。なぜなら、$C$ と $D$ が $t_1, t_2, \dotsc t_m$ のすべての可能性のある成分を含むのであれば
\[
  C + D = 1;
\]
であり、また $D$ は係数 $\frac{1}{0}$ に影響を受けているため、$t_1, t_2, \dotsc t_m$ にこれらの $x_1, x_2, \dotsc x_n$ の表現を代入すると $D = 0$ となります。したがって、同じ代入により $C = 1$ となります。ルールによれば、$c$ は $C$ が発生した場合に $x_1$ が発生する確率です。よって $C = 1$ であり、すべての可能性を含むことから、$c$ は $x_1$ の絶対的な発生確率として解釈されるべきです。

特定の場合における最終論理方程式の実際の形を決定することも興味深いでしょう。例えば、データが独立かつ互いに排他的な事象から得られるものとする場合を考えましょう。例えば、$a_1, a_2, \dotsc a_m$ がそれぞれ異なる明確な多数派の確率であるとします。このとき、論理関数 $X_1, X_2, \dotsc X_m$ は互いに排他的であり、条件を満たす必要があります
\[
  X_1 X_2 = 0,\quad X_1 X_m = 0,\quad X_2 X_m = 0, \text{ etc.}
\]
これより
\[
  t_1 t_2 = 0,\quad t_1 t_m = 0, \text{ etc.}
\]

このような状況では、$x_1$ の開発された論理値は容易に以下のように示すことができます。
\begin{equation*}\begin{split}
x_1 &= \frac{0}{0}(\bar{t}_1 \bar{t}_2 \dotsc \bar{t}_m + t_1 \bar{t}_2 \dotsc \bar{t}_m \dotsc + t_m \bar{t}_1\dotsc \bar{t}_{m-1}) \\
&+ 0/0 \text{ の係数を持つ成分}.
\end{split}\end{equation*}
上記の式において $\bar{t}_1$ は $1 - t_1$ を表します。

これらの調査は、陪審員の判決確率が意見の一致と不一致に関連してデータとなる場合にも同様に適用されます。ここで、論理的シンボル \( w \) は被告の罪状が成立する事態を表します。すると、現在の問題における関数 \( X_1, X_2, \dots X_m \) は、\( w, x_1, x_2, \dots x_n \) を同時に \( \bar{w}, \bar{x}_1, \bar{x}_2, \dots \bar{x}_n \) に変換しても問題内の何らかの変化が生じないような関数です。したがって、\( w \) の最終的な論理的値および \( x_1, x_2, \dots x_n \) の各々の最終的な論理的値は同じ形式（3）で示され、それから類似の一般的な結論が導き出されます。

したがって、単なる統計資料からは、判事または弁護士個人の意見の正確性、被告の罪状、争点の是非に関するいかなる推論も導くことができません。このような要素（上記のもの）の決定が科学の範囲内に収まることを可能にするためには、真実の尺度としての仮定的な基準により新たなデータを得るか、個々の判断の相互関係または独立性に関する仮説により新しい調査形式を許可する必要があります。異なる仮説の結果を検討する際には、以下の一般命題が重要となります。

\begin{center}
\textsc{命題} II.
\end{center}

8. \( n \) の単純な事象 \( x_1, x_2, \dots x_n \) の確率は次のようである：
\[
\operatorname{Prob. } x_1 = c_1, \quad \operatorname{Prob. } x_2 = c_2, \dotsc \operatorname{Prob. } x_n = c_n;
\]
また、\( m-1 \) の複合事象 \( X_1, X_2, \dots X_{m-1} \) の確率は次のようである：
\[
\operatorname{Prob. } X_1 = a_1, \quad \operatorname{Prob. } X_2 = a_2, \dotsc \operatorname{Prob. } X_{m-1} = a_{m-1};
\]
これらの複合事象 \( X_1, X_2, \dots X_{m-1} \) は互いに異なるかつ排他的であるとします。このとき、任意の他の複合事象 \( X \) の確率を求めることを要求します。

この命題においては、$X_1, X_2, \dotsc X_{m-1}$ および $X$ が符号 $x_1, x_2, \dotsc x_n$ のみを用いて定義される関数であると仮定する。さらに、互いに排他的な事象 $X_1, X_2, \dotsc X_{m-1}$ に対しては
\begin{equation}\tag{3}
X_1 X_2 = 0, \quad X_1 X_{m-1} = 0, \quad X_2 X_3 = 0, \text{ 等;}
\end{equation}
が成り立つ。すなわち、システムの任意の二つの要素の積はゼロとなる。このとき
\begin{equation}\tag{4}
X_1 = t_1, \quad X_{m-1} = t_{m-1}, \quad X = t
\end{equation}
と仮定する。すると $t$ は、$x_1, x_2, \dotsc x_n, t_1, t_2, \dotsc t_{m-1}$ の論理関数として決定される。

次に (3) より
\begin{equation}\tag{5}
t_1 t_2 = 0, \quad t_1 t_{m-1} = 0, \quad t_2 t_3 = 0, \text{ 等;}
\end{equation}
が成り立つ。すなわち、$t_1, t_2, \dotsc t_{m-1}$ の二項積は全てゼロとなる。したがって $t$ の展開式においては、係数が 1、0、または $\displaystyle \frac{0}{0}$ を持つような要素のみが含まれる。これらの要素は以下の因子のいずれかを含む：
\begin{equation}\tag{6}
\bar{t}_1 \bar{t}_2 \dotsc \bar{t}_{m-1}, \quad t_1 \bar{t}_2 \dotsc \bar{t}_{m-1}, \dotsc, \bar{t}_1 \dotsc \bar{t}_{m-2} t_{m-1};
\end{equation}
ここで $\bar{t}_1$ は $1 - t_1$ の意味である。残る部分は、各要素において符号 $x_1, x_2, \dotsc x_n$ を含む部分と、その係数を指定する。

任意の整数 $i$（$1 \leq i \leq m-1$）に対して $X_i = t_i$ であるから、
\begin{equation*}
X_i \bar{t}_1 \dotsc \bar{t}_{m-1} = 0
\end{equation*}
が成り立つ。したがって、上記方程式の第一項に含まれる任意の要素は $\displaystyle \frac{1}{0}$ の係数を持つことになる。

今、次のように定義します：

\[
X_m = 1 - X_1 \dotsm - X_{m-1};
\]

このような項が $\bar{t}_1 \dotsc \bar{t}_{m-1}$ を因子として含みながら、係数が 1, 0, もしくは $\displaystyle \frac{0}{0}$ の形であるものは、

\[
X_m \bar{t}_1 \dotsc \bar{t}_{m-1}
\]

という表現に含まれる必要があります。

ここで $X_m$ を二つの部分に分解することができます：すなわち、$X X_m$ と $(1-X)X_m$。前者は $X$ 内に見られる $X_m$ の項の合計であり、後者は $X$ 内に見られない項の合計です。$t$ の展開式が $X$ に等しいとき、前者 $X X_m$ の各項の係数は 1 になりますが、後者 $(1-X)X_m$ の各項の係数は 0 になります。したがって、現在考慮している要素は $t$ の展開に次の項をもたらします：

\[
X X_m \bar{t}_1 \dotsc \bar{t}_{m-1}
+ 0(1-X)X_m \bar{t}_1 \dotsc \bar{t}_{m-1}.
\]

また、$X_1 = t_1$, $X_2 t_1 = t_2 t_1 = 0$, と次のように与えられるとき、

$t_1 \bar{t}_2 \dotsc \bar{t}_{m-1}$ を因子として含みながら、係数が 1, 0, もしくは $\displaystyle \frac{0}{0}$ の形である項は、

\[
X_1 t_1 \bar{t}_2 \dotsc \bar{t}_{m-1}
\]

に含まれることが明らかです。同様の議論により、これらの項は $t$ の展開に次の項をもたらします：

\[
X X_1 t_1 \bar{t}_2 \dotsc \bar{t}_{m-1}
+ 0(1-X)X_1 t_1 \bar{t}_2 \dotsc \bar{t}_{m-1}.
\]

以上を残りの項に適用することで、$t$ の最終的な表現は、

\[
t = X X_m \bar{t}_1 \dotsc \bar{t}_{m-1}
+ X X_1 t_1 \bar{t}_2 \dotsc \bar{t}_{m-1} \dotsm
+ X X_{m-1} \bar{t}_1 \dotsc \bar{t}_{m-2} t_{m-1}
\]
\[
+ 0(1-X)X_m \bar{t}_1 \dotsc \bar{t}_{m-1}
+ 0(1-X)X_1 t_1 \bar{t}_2 \dotsc \bar{t}_{m-1} + \text{等.}
\]
\[
+ \text{係数が }\frac{1}{0} の項.
\]

となります。

この表現において、$X\,X_m$は$X$と$X_m$に共通する成分の和を表し、その和は論理計算の則により$X$と$X_m$を掛け算することで得られる。

論理から代数へ移行すると、($X\,X_m$)は上記の積が、乗算を実施した後または共通成分を選択した後に、シンボル$x_1, \dotsc x_n$に量的意味を与えることで表されます。

この理解に基づき、一般則（XVII. 17）により次の式を得ます：
\begin{gather*}
  \text{Prob. }t
\\
= \frac{(X X_m) \bar{t}_1 \dotsc \bar{t}_{m-1} + (X X_1)t_1 \bar{t}_2 \dotsc \bar{t}_{m-1} + (X X_{m-1})\bar{t}_1 \dotsc \bar{t}_{m-2} t_{m-1}}{V}, \tag{9}
\\
  V = X_m \bar{t}_1 \dotsc \bar{t}_{m-1} + X_1 t_1 \bar{t}_2 \dotsc \bar{t}_{m-1} + \dotsb + X_{m-1} \bar{t}_1 \dotsc \bar{t}_{m-2} t_{m-1}, \tag{10}
\end{gather*}

これにより、$x_1, \dotsc x_n, t_1, \dotsc t_{m-1}$の関係は次のタイプになります（$i$は1から$n$まで変化）：
\begin{gather*}
  \frac{(x_i X_m) \bar{t}_1 \dotsc \bar{t}_{m-1} + (x_i X_1)t_1 \bar{t}_2 \dotsc \bar{t}_{m-1} + (x_i X_{m-1})\bar{t}_1 \dotsc \bar{t}_{m-2} t_{m-1}}{c_i}
\\
= \frac{X_1 t_1 \bar{t}_2 \dotsc \bar{t}_{m-1}}{a_1} \dotsb = \frac{X_{m-1} \bar{t}_1 \dotsc \bar{t}_{m-2} t_{m-1}}{a_{m-1}} = V. \tag{11}
\end{gather*}

上記のシステムから、シンボル$t_1, \dotsc t_{m-1}$を次に消去します。

私たちは
\[
  t_1 \bar{t}_2 \dotsc \bar{t}_{m-1} = \frac{a_1 V}{X_1}, \quad \bar{t}_1 \dotsc \bar{t}_{m-2} t_{m-1} = \frac{a_{m-1} V}{X_{m-1}}. \tag{12}
\]

これらの値を（10）に代入すると、
\[
  V = X_m \bar{t}_1 \dotsc \bar{t}_{m-1} + a_1 V + \dotsb + a_{m-1} V.
\]
従って，
\[
  \bar{t}_1 \dotsc \bar{t}_{m-1} = \frac{(1 - a_1 - \dotsb - a_{m-1})V}{X_m}.
\]

現在、次のように定義します：
\[
a_m = 1 - a_1 \cdots - a_{m-1},    \tag{13}
\]
すると、次の関係が成り立ちます：
\[
\bar{t}_1 \cdots \bar{t}_{m-1} = \frac{a_m V}{X_m}.   \tag{14}
\]

次に、(12)と(14)を用いて方程式(9)と、(11)の最初の行をシンボル$V$で等しい方程式を簡約します。またProb. $X$をProb. $t$で表すことにします。すると以下の関係が得られます：
\[
\text{Prob. } X
= \frac{a_1(X X_1)}{X_1}
+ \frac{a_2(X X_2)}{X_2} \cdots
+ \frac{a_m(X X_m)}{X_m},             \tag{15}
\]
\[
\frac{a_1(x_i X_1)}{X_1}
+ \frac{a_2(x_i X_2)}{X_2} \cdots
+ \frac{a_m(x_i X_m)}{X_m} = c_i;    \tag{16}
\]
ここで、$X_m$と$a_m$は(7)と(13)で与えられています。

これらの方程式は、現在考慮している問題の直接的な解法を含んでいます。方程式(16)では、$i$に順次1, 2, \dots, $n$を与えることにより$n$個の方程式が形成され、これより$x_1, x_2, \cdots, x_n$の値を求めることができます。これらの値を方程式(15)に代入することで、$\text{Prob. } X$を定数$a_1$, $c_1$, などから関数として求めることができます。

以上の解法から導き出される注目すべき結論は、複合事象$X_1, \cdots, X_{m-1}$の確率が、事象$x_1, \cdots, x_n$が完全に独立であり、それぞれの確率が与えられた$c_1, \cdots, c_n$である場合と同じであれば、事象$X$の確率も、同じ仮定のもとで計算された絶対的な独立性を持つ事象$x_1, \cdots, x_n$を前提とした場合と同じであるということです。この仮定のもとでは、量的システムにおける推定$x_1 = c_1$, $x_n = c_n$により$X_1 = a_1$, $X_m = a_m$が得られ、これより方程式(15)と(16)は以下のようになります：
\[
\operatorname{Prob. } X = (X X_1) + (X X_2) \cdots + (X X_m),    \tag{17}
\]
\[
(x_i X_1) + (x_i X_2) \cdots + (x_i X_m) = c_i.     \tag{18}
\]

しかし、$X_1 + X_2 + \dotsb + X_m = 1$ であるから、(17) 式の第二項は $X$ に含まれるすべての成分を取って、それらに代数的な意味を与えることによって形成されることが明白である。同様の論じ方は (18) 式にも適用できる。従ってこれらの式はそれぞれ
\begin{gather*}
\operatorname{Prob. } X \text{（論理的）} = X \text{（代数的）}, \\
x_i = c_i.
\end{gather*}

したがって、$X = \phi(x_1, x_2, \dotsc, x_n)$ であるときには
\begin{equation*}
\operatorname{Prob. } X = \phi(c_1, c_2, \dotsc, c_n),
\end{equation*}

という結果を得る。これが求められた結果である。

したがって、もし量 $c_1,\dotsc c_n$ が不定であり、平均的な共通値を持つ仮定がなされていなければ、システム (15) (16) はこれらの量に任意の値 $x_1, x_2, \dotsc x_n$ を与えることで満足される。ここで、これらの値は方程式
\begin{equation*}
X_1 = a_1, \quad X_{m-1} = a_{m-1}, \quad X=a,
\end{equation*}
を満たすものであり、要素 $a$ の値は経験によって与えられるものとします。

9. 一般解 (15) (16) を判断の確率について適用する前に、以下の変換を行うことが便利です。データが
\begin{align*}
x_1 = c_1 & \dots x_n=c_n, \\
\operatorname{Prob. } X_1 = a_1 & \dots
\operatorname{Prob. } X_{m-2} = a_{m-2};
\end{align*}
であり、未知数 $\operatorname{Prob. } X_{m-1}$ の値を $a_{m-1}$ と表すとします。このとき、(15) および (16) では \\
\begin{center}
\begin{tabular}{ll}
$X$ を $X_{m-1}$ に, & $\operatorname{Prob. } X$ を $a_{m-1}$ に, \\
$X_{m-1}$ を $X_{m-2}$ に, & $\quad \quad a_{m-1}$ を $a_{m-2}$ に, \\
$X_m$ を $X_{m-1} + X_m$ に, & $\quad \quad a_m$ を $a_{m-1} + a_m$ に,
\end{tabular}\end{center}
と変更します。これらの変換を行い、$(X_{m-1} X_r)=0$ が $r = m-1$ のときのみ成り立つことを考慮すると、方程式 (15) (16) は
\begin{equation}\tag{19}
a_{m-1}=\frac{(a_{m-1} + a_m)X_{m-1}}{X_{m-1}+X_m},
\end{equation}
\begin{equation}\tag{20}
  \frac{a_1(x_i X_1)}{X_1} \dotsc
+ \frac{a_{m-2}(x_i X_{m-2})}{X_{m-2}}
+ \frac{(a_{m-1} + a_m)(x_i X_{m-1} + x_i X_m)}{X_{m-1} + X_m}.
\end{equation}
を導きます。

(19) から
\begin{equation*}
\frac{X_{m-1}}{a_{m-1}}=\frac{X_m}{a_m} = \frac{X_{m-1}+X_m}{a_{m-1}+a_m},
\end{equation*}
が得られることから、(20) の最後の項は
\begin{equation*}
\frac{a_{m-1}(x_i X_{m-1})}{X_{m-1}}+\frac{a_m(x_i X_m)}{X_m}.
\end{equation*}
と簡略化されます。これらの簡略化により、システム (17) および (18) は以下の対称的な形に置き換えることができます：
\begin{equation}\tag{21}
\frac{X_{m-1}}{a_{m-1}}=\frac{X_m}{a_m},
\end{equation}
\begin{equation}\tag{22}
\frac{a_1(x_i X_1)}{X_1} +
\frac{a_2(x_i X_2)}{X_2}\dotsc +
\frac{a_m(x_i X_m)}{X_m} = c_i.
\end{equation}

これらの方程式は (7) および (13) と組み合わせることで、経験によって提供される数値データ $c_1 \dotsc c_n, a_1 \dotsc a_{m-2}$ の関数として $a_{m-1}$ を決定することができます。これにより、(17) および (18) のシステムは次の対称的な形で置き換えることができます。

これらの方程式と (7) 及び (13) を組み合わせることにより、$c_1 \dotsc c_n, a_1 \dotsc a_{m-2}$ の関数として $a_{m-1}$ を決定することができます。経験によって提供される数値データに基づいて、これらの方程式を適用します。

\#\#\# 命題 III

10. 与えられた一致の事例や議決で特定された数的多数決から抽出された任意の確率体系について、ある確定的な仮定のもとで議会の一員が正しい判断を行う平均確率を求めよ。

以前の章の一部で、一致の決定や特定された数的多数決の確率を経験からどのように決定するかについては触れられている。Prop. \textsc{i} の表記法を採用し、与えられた確率を表す関数 $X_1, X_2, \dotsc X_{m-1}$ で表すイベントを示す。この状況の性質から明らかになったように、これらのイベントは互いに排他的であり、それらを表す関数は符号 $x_1, x_2, \dots x_n$ に関して対称であることが示された。これらの符号は Prop. \textsc{i} で用いたとおりの意味を持ち続ける。すなわち、$x_i$ は$i$番目の議員が正しい意見を形成するというイベントを表す。

現在の経験からの直接的なデータは以下の通りである：
\[
\operatorname{Prob. } X_1 = a_1, \quad
\operatorname{Prob. } X_2 = a_2, \dotsc
\operatorname{Prob. } X_{m-2} = a_{m-2},
\]
\[
\operatorname{Prob. } X_{m-1} = a_{m-1}.
\]

$X_1\dots X_{m-1}$ は論理符号 $x_1, \dotsc x_n$ の関数であり、これらのイベントを表す確率に対して、未定値 $c$ を割り当てる。したがって、以下の式が成り立つ：
\[
\operatorname{Prob. } x_1 = \operatorname{Prob. } x_2 \dots = \operatorname{Prob. } x_n = c.
\]

現在、第\textsc{i}定理において確認したように、直接与えられた情報（1）（2）は仮説なしに独自の推論を進めるものであり、問題の再述に過ぎません。一方で、ラプラスとポワソンの方法を使用し、方法論におけるデータとしてシステム（3）のみを利用し、最終的に得られた結果を実験的システム（1）（2）と比較するという手法は、個々の判断の独立性という疑問のある仮説に完全に依存していることが明白です。しかし、この仮説に完全に依存することは適切ではありませんが、それを完全に排除することもできませんし、その代わりとなる何かを用いることもできません。

そこで、個々の判断の独立性が制限されている場合の結果を検討しましょう。その制限条件は余計な情報によって提供されます。システム（1）（3）から、本方法論の手法を使用してProb. $X_{m-1}$ を決定し、それを（2）と比較して定数$c$を求めます。ここで任意の選択権が主張されていますが、第\textsc{i}定理から明らかに、何かこのようなものが存在することは、何らかの明確な解を得るためには避けられないことがわかります。

選択原則については、最終的に比較のために残された方程式（2）は、その数値要素$a_{m-1}$が最も重要な経験から得られる一次系列の中から最大であると評価されるべきものであると考えます。

現在、関数 $X_1, X_2, \dots, X_{m-1}$ を表す事象が互いに排他的であるという性質から、前節の命題の結論的な方程式は適用可能となります。同じ関数の対称性と $c_i$ によって表される値の系が単一の値 $c$ に簡約されることにより、式 (22) 表現する方程式は同一となり、$x_1, x_2, \dots, x_n$ の値は等しくなり、単一の値 $x$ で置き換えることができます。そして、

\begin{gather*}
  \frac{X_{m-1}}{a_{m-1}} = \frac{X_m}{a_m},   \tag{4} \\
  \frac{a_1(x X_1)}{X_1} + \frac{a_2(x X_2)}{X_2} + \dotsb + \frac{a_m(x X_m)}{X_m} = c.                  \tag{5}
\end{gather*}

このような解法の性質は以下の通りです：

与えられたデータにおいて、関数 $X_1, \dots, X_{m-1}$ と値 $a_1, \dots, a_{m-1}$ が与えられているため、まず、

\begin{align*}
  X_m &= 1 - X_1 - \dotsb - X_{m-1}, \\
  a_m &= 1 - a_1 - \dotsb - a_{m-1}.
\end{align*}

与えられたまたは決定された関数 $X_1, X_2, \dotsc, X_m$ のそれぞれから、特定の符号（例えば因子としての $x_1$）を含む成分を選択する必要があります。これにより、関数 ($x X_1$), ($x X_2$), などが決定され、その後すべての関数において、$x_1, x_2, \dotsc, x_n$ を個別に $x$ に置き換える必要があります。

あるいは、システム（4）（5）における任意の代数的関数 $X_i$ を、論理的な符号 $x_1, x_2, \dotsc, x_n$ が独立した事象を表し、これらの共通の確率が $x$ であるという仮定のもとで、その事象列 $x_1, x_2, \dotsc, x_n$ のいずれかの特定の事象と $X_i$ の同時発生の確率を表すものと考えることができます。同様の仮定のもとでは、($x X_i$) はその系列 $x_1, x_2, \dotsc, x_n$ のいずれかの特定の事象と $X_i$ の同時発生の確率を表します。

$X_i$, ($x X_i$), などの方程式の形が決定されると、方程式（4）は $x$ の値を与えることになります。そしてこの値を方程式（5）に代入することで、必要な要素 $c$ の値を求めることができます。その解が提供する2つの値のうち、1つは $\frac{1}{2}$ より大きく、もう1つは $\frac{1}{2}$ より小さいものがある場合、一般的な考慮に基づいて議会の一員が正しく判断するよりも誤った判断をする可能性が高いと推測される場合には、大きい方の値を選択します。

したがって、仮定された原理によれば、方程式（2）における最終的な比較のために残すべき最大の値 $a_{m-1}$ が与えられることにより、提案された問題に対する明確な解を得ることができます。そして、システム内の他の方程式がより優れた精度を持つなど、別の理由で同様に方程式（2）の位置に残される場合でも、同じ形式の解は依然として適用されます。

以上の原著英文を11. 上記の制限が最終的な解にどの程度影響を与えたかを検討しましょう。方程式（5）は、選択肢の選択とは全く独立しています。同様に、式（4）の第二項も独立しています。もし私たちは関数 $X_1$ を制限したとしても、$X_{m-1}$ ではなく、方程式を決定するための $x$ の値は
\[
  \frac{X_1}{a_1} = \frac{X_m}{a_m}   \tag{6}
\]
となるでしょうが、その結果得られた $x$ の値は同じ最終方程式（5）に代入しなければなりません。もし $x_1, x_2, \dotsc, x_n$ が本当に独立している場合、式（4）、（6）およびそれらの他の類似方程式は等価であり、これらいずれかから得られる $x$ の値は真の $c$ の値となることが知られています。これにより、最終方程式（5）を検証する手段が提供されます。もしその方程式が正しければ、上記の状況において $c = x$ という仮定によって満たされるべきです。つまり、方程式
\[
  \frac{a_1(x X_1)}{X_1}
+ \frac{a_2(x X_2)}{X_2} \dotsc
+ \frac{a_m(x X_m)}{X_m} = x     \tag{7}
\]
が解くと、式（4）または（6）から得られる $x$ の値と同じになるはずです。実際、これが成り立つでしょう。なぜなら、仮定により
\[
  \frac{X_1}{a_1} = \frac{X_2}{a_2} \dotso = \frac{X_m}{a_m}
\]
が成り立つので、既知の定理により
\[
  \frac{X_1}{a_1} = \frac{X_2}{a_2} \dotso = \frac{X_m}{a_m}
= \frac{X_1 + X_2 \dotsc + X_m}{a_1 + a_2 \dotsc + a_m}
= 1
\]
となります。したがって、$xX_1$, $xX_2$, $\dotsc$, $xX_m$ を代入すると（7）は
\[
  (xX_1) + (xX_2)\dotsc (xX_m) = x
\]
という恒等式となります。

従って、事件$x_1, x_2, \dotsc, x_n$が本当に独立であるならば、式(4)と(5)のシステムは正しいものであり、そのシステムは、それが得られた過程における最初の一歩の任意性から独立します。これらの事件が互いに独立でない場合、最終的な方程式のシステムには、上記の選択原則を保留した上で任意の要素があります。しかし、方程式(5)の持続的な形式から推測すると、その解は個々の判断が絶対的に独立であるという仮定に基づく解よりも、より少ない程度に任意性を持つことが推測されます。$c$の値の限界が$x$の値の限界に依存する議論によって、このような点が決定されるでしょう。

これらの考慮は、方程式（7）が、イベント$x_1, x_2, \dotsc x_n$が本当に独立であるとき、厳密に正確であり、関数$X_1, X_2, \dotsc X_m$に関する対称性を持つ、任意の要素を含まない、という特徴を持つことから、その方程式を問題の平均的な一般解として受け入れるべきではないかと我々に問い立てます。この点が適切であるかどうかを確認するための適切な方法は、方程式（7）によって与えられる$x$の値が、システム（4）、（5）で示されるタイプを持つ方程式系によって決定された$c$の範囲内に一般的に含まれるかどうかを確認することだと考えます。通常の場合には、これが事実である可能性が高いと思われます。

このような考慮とは独立して、（7）は特定の解法原理の表現と見做すこともできます。すなわち、$X_1, X_2, \dotsc X_m$をその確率が$x$であるイベントの排他的な原因として扱い、経験からこれらの原因の確率$a_1, a_2, \dotsc a_m$を得て、個々の判断の独立性の仮定のもとでイベントの条件付き確率をそのような原因に依存して形成し、

\[
  \frac{(x X_1)}{x_1},\quad \frac{(x X_2)}{X_2},
  \text{ 等 (XVII. Prop~\textsc{i.})}
\]

を導き、方程式（7）を得ることです。しかし、私はこの方法がより緩和的であると見做す一方で、実際的な操作ではより便利な解法の手段であると考えます。これは一般解法に採用された方法よりも厳密性は欠けますが、その代わりに実践的な利便性を提供する可能性があります。

12. 以上の方程式において、論理関数 \(X_i\) が議会の \(r\) 名のメンバーが一方の投票をし、その他の \(n-r\) 名のメンバーが他方の投票を行うという事象を表す場合に、代数的関数 \(X_i\), (\(x X_i\)) 等がどのような形を取るかを特定するだけが残されている。明らかにこの場合、代数的関数 \(X_i\) は、事象 \(x_1, x_2, \dotsc, x_n\) が独立であると仮定し、その共通の確率を \(x\) で測るときの予想される事象の確率を表している。したがって、第3節の式により
\[
X_i = \frac{n(n-1)\dotsc (n-r+1)}{1\cdot 2\dotsc r} \{ x^r + (1-x)^{n-r} \}
\]
となる。

同じ状況下において (\(x X_i\)) は、議会の特定のメンバーが上記で述べた投票状態と一致する正確な判断を行う確率を表す複合事象を表す。したがって、特定のメンバーが正しく投票し、残りの \(n-1\) 名のメンバー中 \(r-1\) 人が正しく投票しているか、または同じメンバーが正しく投票し、残りの \(n-1\) 名のメンバー中 \(r\) 人が誤って投票している確率を表す。したがって
\[
(x X_i) = \frac{(n-1)(n-2)\dotsc (n-r+1)}{1\cdot 2\dotsc r-1} x^r + \frac{(n-1)(n-2)\dotsc (n-r)}{1\cdot 2\dotsc r} x^{n-r}.
\]

\begin{center}\textsc{第四定理}\end{center}

13. 任意の確率系が記録された一致事例や刑事裁判所での決定における指定された多数決から与えられた場合、最後の定理と類似した仮定のもとで必要となる議員の正確な判断の平均確率 \(c\) および被疑者の一般の罪の確率 \(k\) を求めることを要求する。

この問題の解法は、前回のものと僅かな違いしかありませんし、同じ一般的な公式（4）および（5）、または（7）に従うことができます。二つの要素 $c$ と $k$ の決定が必要であることに注意すべきです。したがって、$X_1, X_2,\dotsc X_{m-1}$ のうちの二つを最終的な比較のために保存します。例えば、残りの $m-3$ 関数をデータの表記に使用するか、または二つのそれぞれのセット $X_2, X_3,\dotsc X_{m-1}$ と $X_1, X_2,\dotsc X_{m-2}$ を使用します。いずれの場合も、元の独立したデータが少なくとも二つ必要であることが仮定されます。方程式（7）のみを使用すると、現在の場合には二つの方程式を提供し、次のように書き表すことができます：

\begin{align*}
  \frac{a_1(xX_1)}{X_1}
+ \frac{a_2(xX_2)}{X_2} \dotsc
+ \frac{a_m(xX_m)}{X_m} &= x,     \tag{1}
\\
  \frac{a_1(kX_1)}{X_1}
+ \frac{a_2(kX_2)}{X_2} \dotsc
+ \frac{a_m(kX_m)}{X_m} &= k.     \tag{2}
\end{align*}

これらの方程式は次の方法で使用されます。$x_1, x_2, \dotsc x_n$ をそれぞれ裁判員が正しい意見を形成する各事象とします。また、$w$ を被告者が有罪であるという事象とします。これらのシンボルを利用して、与えられた確率の関数 $X_1, X_2,\dotsc X_{m-1}$ および関数 $X_m$ を論理的に表現できます。次に、関数 $X_1$ の中から、任意の特定のシンボル $x_1, x_2, \dotsc x_n$ の因子を含む構成要素を選択し、また $w$ の因子を含む構成要素も選択します。これらの結果において、それぞれの $x_1, x_2, \dotsc x_n$ を $x$ に、そして $w$ を $k$ にそれぞれ置き換えます。この操作により、$(xX_1)$ と $(kX_1)$ が得られます。同様の変換を全般的に行うと、方程式（1）、（2）は特定の仮定に基づいて $x$ および $k$ を決定します。

14. 上記の調査から以下の事実と結論をまとめることができます：

1. ある集団の意見に含まれる一致と不一致の記録から、個々のメンバーが正誤を判断する確率や、提出された問題に対するその集団の価値についての明確な数値的な結論は導き出せない。

2. しかし、様々な仮定に基づいてこのような結論を導くことができる。1つ目として、個々の判断が絶対的に独立であるという通常の仮定；2つ目として、実際のデータによって正当化される特定の仮定の修正；3つ目として、上記の修正により得られた解に共通する形式を示唆する独自の解決原理。

最後に。最終的な結果に対する疑念は、方法の不完全性ではなく、仮定自体の不確実性に依存している。

ただし、仮定が最大限の範囲で妥当し、経験上のすべてのデータを考慮に入れる限り、得られた結果の偏差は小さく、その平均値は第3章の方法によって非常に信頼性のある形で決定できる可能性が高い。これらの方法の中から私は最初の方法に優先順位をつけたい。

この平均解の原理が合意された場合、他の考慮事項によれば、特定の構造を持つ裁判所や議会、そして固定ルールに基づいて議論を行う集団に対する c と k の値はほぼ一定であり、ただしこれらは知識と正義の進歩に依存する小さな長期的な変動を伴う可能性がある。現在、その定量化に適したデータはほとんど存在しないと思われる。

22章

\newpage

\chapter{科学の本質と人間の知性の構造について（Constitution of the Intellect）}
\label{ch:xxii}

\begin{center}
\large\textit{数学的論理学が照らし出す人間精神の認識論的本質と哲学的展望}
\end{center}
\vspace{4mm}

\#\# 知の構造について

\#\#\# 1. システムの構造とは

私がいう「システムの構造」は、その観察上の特性を生み出す原因や傾向の総合であり、これらの原因や傾向が干渉なく作用し、システムが考えられる適応状態にあるときの条件において機能するときのものです。このような適応性についての判断は、システムが最も自由に動作し、最も調和した結果を生み出す、またはその構築の明らかなる目的を果たす他の方法で動作する状況における研究に基づくべきです。

明確な原因と条件、そして目的を知っている場合があります。これは、当該主題に関する最完璧な知識と言えます。また、原因が部分的にしか理解できていない場合もありますが、その作用の法則をある程度決定し、さらに一般的な傾向を見出し、推論することができる場合もあります。

このような観点から、私たちは私たちの性質に道徳的な能力があると結論付けるべきであると考えています。これは、視覚の能力に関連する器官と同様に特殊な機器を理解できるわけではないし、また人々が普遍的な行動規範を採用することに一致しているという理由でもありません。それは、どの形であれ、道徳的承認または非難の感情が全員で現れる一方で、人間社会の進歩が見られる場所では、一定の行動クラスにその傾向が一貫して且つ法的な性質を示すように付随するからです。

常にそしてどこでも秩序の表れは、目的の達成または意図の存在を実際的（XX.22）な証拠として提供します。それは測定可能ではないものの、その存在は確実です。常にそしてどこでも秩序の現れは、目的の達成または意図の存在を実際的（XX.22）な証拠として提供します。

2. 意識の構成に関する特定の問題は、ほとんど言うまでもなく、毎時代において思慮深い想像力を引きつけた。これは、古代思想の最も偉大な指導者たちが我々の種族に内在すると信じていた知識への欲求を対象とし、通常のこの動機に人間的で個人的な関心という誘因を加えるからである。実際、真の真理への献身はしばしば偏見なくその目的を追求するが、それは外なる観察の美しい領域を広く探求することを促す一方で、自分たちの能力に対する研究を怠らないようにと警告する。

果たして、物質的利益に最も熱心な時代でも、思考の流れは部分的に内向きになり、すべてを理解するために理解されるものを理解しようとする欲求は挫折しても再燃するばかりである。

この努力の執拗さは、真実を追求する誠実な探求者たちにおいて、このような推測が希望的に結実しないという確信が一般的になければ維持されなかった可能性が高い。私たちは、このような調査結果に常に付随する誤りと不確定性があることを認識しながらも、それでも一定量の肯定的な知識が残る可能性があると考えることができる。遭遇する矛盾はしばしば言葉によるものよりも実質的なものであるということも感じられるだろう。特に、その主題から得られる推定的な結論には、物理学のより明確で華やかな結果と比較しても不十分ではない価値と興味が存在する。

このような考慮は完全に正当化されているように思われる。多くの問題について、知性の本質と構成に関する調査において不可解決であると思われるものもあるとしても、限定的ながしかし疑いの余地のない知識を得られる他の問題も存在している。また、高確度な類推に基づく結論を達成できる他の問題も存在する。

昼と夜の領域は厳密には重複していないが、一方の光が他方の闇に徐々に消えていく境界線を通じて区別されるように、肯定的な知識の各領域はその影響や明かりを一部拡張する形で議論の余地のある推測の地域を囲んでいると言える。したがって、知性の構成に関する質問がある種の絶対的な決定を受け入れることができないとしても、現在の知識状態において反射的に得られる情報によりその推定的な解決は容易になる可能性がある。また、科学の本質に関連する質問や、すなわち特定の真実と科学の教義についても、この著作の一般原則を受容した者は現代で通常受け入れられている見解とは異なる肯定的な意見を持つことになる。

以下に例を挙げるとしよう：

現代の文法の最高权威者の一部は、接続詞は主張を結びつけるだけであると主張しています。現在、この見解が正確かどうか直接検討することなく、科学的な原理を認めている本稿の作者にとっては、このような見解を一貫して持つことは不可能であることが明らかです。なぜなら、そのような見解は、両方の関係から観察した精神的操作の法則について調査し適用された現在の著作において、\emph{実行}または\emph{表現}の可能性を否定することに他なりません。

（ラザム『英語論』；サー・ジョン・ストドラーの普遍文法など）

次に述べる内容では、本稿の前半部で確立されたより明確な結論を要約し、それから先述した一般的な目的に関連する一つか二つの思考の流れを示すことに取りかかります。

3. 以下の結論は、知性の構造に関連し、肯定的知識の領域に属するものとして考えられる。これらには、この論文の一般的方法を基礎とし、その適用が決定された科学的な思考法則と推理法則がある。また、第5章で述べられた原則も含まれる。これらの法則は思考の領域を二つの異なるながら共存する領域に分解すること（IV.XI.）によって解決される。それらの領域内の知性の操作が共通の法体系のもとで支配されること（XI.）についての法則がある。これらの法則の一般的な数学的な性格と、その実際的な表現（II.III.）がある。これらは数の領域における思考法則との類似性と、それらとの相違点についても言及される。また、論理学が関心を持つすべての思考対象において支配的な性格を有する二つの極限的概念である宇宙と永遠に関する概念があり、これらの概念は数の科学における基本的な単位概念と密接に関連している（II.IX.）。

これらのような多くの類似した結果は、単なる推測的または類推的な結論ではなく、科学の真実として扱われる資格がある。これらを形而上学と呼ぶかどうかは問題ではない。これらの真実が基づく証拠の性質は種類こそ異なるものの、物理科学の一般の真理を支持するためのどの証拠よりも価値が高いものではない。

また、すべてのより正確な知識の進行に見られる一定の順序があることは認められるべきである。
物理学の各分野の研究は観察から始まり、事実を比較検討することで仮定的な法則との関連性についての認識を進める。その仮定の妥当性は、仮説が正当であれば確率を著しく増加させるための新たな実験によってテストされる。そして、現象の法則が十分に決定された後、仮説と推論の適切な混合による原因の調査が研究を締めくくる。
この知識の進行順序においては、本書で考慮した性質や法則もその役割を持つ。従って、科学の本質を公正に調査するか、またはそのような科学との関係における知的な構成を調査するためには、上記の二つの系列のいずれもが孤立して扱われるべきではないことが明白である。
特に、互いに関連性のある真理は、相互の影響と関連性において考慮されるべきである。

4. すべての肯定的知識に対する実験的基础の必要性を、そのような精神的な法則や原理、そして操作の体系という観点から考察すると、まだ大きな程度で議論が分かれる推測的思想の世界における重要な問題に光を当てる傾向がある。経験によって提示される特定の事実から一般の科学的命題へどのように到着するのか？これらの命題の性質とは何か？それは単なる経験の集積であるか、それとも精神が自ら提供する結びつける原理があるか？要するに、科学的な真実の本質は何か、そしてその信頼性の根拠は何なのか？

以上の原著英文をそのような問いに対する単一で一般的な答えは存在しないことは明らかです。場合によっては、それらの問いすら議論する必要がないこともあります。一般の命題が実際に観察された事実をただ数えること（\textit{per enumerationem simplicem}）によって確認された事実を表現しているだけである例はよく知られています。天文学者はこの観察に基づいて、すべての既知の惑星は太陽を中心に西から東へと回っていると主張します。しかし、一般の命題がその真実性を特定の例から推測して得られたものから、未観察の事例にもその真実性を拡張することもあります。単なる帰納的論理原則だけではこのような手続きは正当化されません。

火星に関する多くの観測からケプラーが楕円軌道上での回転を推論した場合、その結論は彼の前提よりも広範なものでしたし、実際には単なる観察によって得られるどんな前提よりも広範なものでした。それでは、このような一般化にさえ将来的な正当性を与えるためにはどのような要素が必要なのでしょうか？それは我々の本質的な能力である秩序を認識する力と、自然現象が秩序に基づいて結びついているという予測的仮定（その根拠はいかなるものであれ）です。これらの要素なしでは、物理学の一般化された真実性は決して確認されることはなかったでしょう。

この手続きによって得られる結論が確実なものであると認めるとすれば、物理科学における多くの一般的な命題は確実なものではなく、概ねの真実性を有するものに過ぎないという結論しか出てこないかもしれません。知識の保証はここにあります：このような結論が自然の構造を正確に反映している場合、その真実性に対する信頼は無限に確認され、最終的には確信と区別せずに定着します。

上記のように推論の基礎として一般命題が必要であるという必要性は、帰納的推論の原理が示すものであり、論理学者は特定の前提からどのような結論も導くことができないと主張しています。高名な現代の著者は、すべての推論は特定の事実から他の特定の事実へと進むものであると主張します。彼らは、あるクラスの一部のメンバーが持っている性質を所有していることを結論づける際に、その性質が全員に共有されているという中間的な一般命題を確立する必要がないことを例示しています。

しかし、推論が行われる秩序や類似性の原理が明示されたり、理解されたりしなければ、最終的な結論は正当化されません。この形で少なくとも、帰納的推論の基礎としての一般命題の必要性は確認されます。しかしながら、私はこれらの推論法の存在そのもの、さらにそれらの特性（本作で調査した法則が示す特有の）に深く関わっていると見なします。

もし論理過程を精密に分析すると、結論が参照する個体のすべての特殊性が抽象化され、そのクラスへの属する性質のみに関心が集中することになります。

5. しかし、経験の比較された事実から帰納的に導き出される一般的な命題だけでなく、必要真理と呼ばれる領域に属する他の命題も存在します。そのようなものは算術の一般命題や、本論文の一般方法を基礎とする思考法の原理などです。これらの命題は特定の例で厳格に検証可能であり、個々の事例からその全体性が明らかになります。また、感覚の不完全性から正確に検証できない必要真理を表現する一般的な命題も存在します。幾何学の一部の命題はこのようなものでしょうが、このような命題は単に幾何学の領域に限られていません。その性質と起源に関する問題は非常に古いもので、知覚の構造についての調査よりも他のどの問題よりも密接に関連しているため、ここで考慮することも無意味ではありません。

この問題に対する主な見解には以下のようなものがあります。まず、これらの命題は経験に依存せず、心の中に独立して存在すると主張されています。それらの命題が対象とする概念は永遠の原型の模倣であるとされます。プラトンはそのような原型を理想の世界として提示しましたが、これらは感覚的対象全体が薄い影や曖昧な暗示に過ぎない真実性を持つものだと考えられています。

一方で、これらの命題の対象は個々の経験的対象のコピーであると主張する者もいます。あるいは、単なる名詞であり、個々の経験的対象そのものであるとも主張します。それらに関連する命題は、その対象が不完全なため部分的に真実であるとされます。最後に、これらは感覚的観察から抽象化された知的な製品であり、個々の事物とは異なり科学の対象となるべきもの、つまり厳密で一般的な命題を主張できるようなものであるとも主張します。

おそらく、他の見解も存在し、その中には感覚的要素が優位にあるものと、知的な或いは理想的要素が優位にあるものがあります。

現在、上記の視点を採用する（純粋に理想的な見解や純粋に名目主義的な見解については考慮しない）場合、科学において真偽混在なしで成り立つ命題が正しい対象とは何かという問いに対して、一つの答えしか考えられないでしょう。それは、経験上の個々の対象もしくはそれらを提示する心理的像も、この称号に厳密な権利を持つものではないと考えられます。

自然や芸術において、直線、三角形、円などの幾何学的な定義と完全に一致するものは存在しないことが確実であり、感覚で認識できないほどの誤差が存在する場合もあります。また、完璧な心理的像を形成し、その一致がより正確なものであるかどうかは疑問視されるべきでしょう。

しかし、不完全であっても、このような概念は何か超越的なものに指向するという点では確実です。不完全性が消失する方向への漸近過程において、これらの概念は存在します。完璧な三角形や正方形、円は自然の中に存在せず、私たちの代表的想像力を逃れ、考えの中だけに現れます。しかし、理解力の驚くべき能力により、これらは絶対的に真なる命題の主題となることができます。

理性的領域は、想像力の領域よりも広大であることが明らかになります。確かに、科学的な要素としての形態、方向、大きさなど、これらの概念を厳密に思い描けると信じる者はいるかもしれませんが、そのような見解からは、すべての試みがこれらの要素を正確に現実化することができないというより正当な意見が浮かび上がります。つまり、外部の対象が実際にそれらに近づくほど、または想像力による抽象化によって特定の「極限状態」に近づくほど、科学に関する一般命題は絶対的な真実に近づき、その実際的な誤差は消失します。

物理科学についても同様な観察が適用されます。力、極性、結晶化などと呼ばれるこれらの「基本概念」は、私が考えると、経験から独立した知的製品でもなく、外部のものからの単なるコピーでもありません。一方では、経験に必然的な前駆体がある一方で、形成には抽象化の力が要求され、それが我々の性質の一部である一般的な能力や傾向によって導かれます。

秩序という概念について深い洞察によれば、それは自己正当化または自己制御を内包しており、私たちの能力の信頼性はその結果が理由に満足する秩序と一致することにより評価されます。「秩序の概念は驚くべきものであり、それが自身の中に自己正当化や自己制御を持っていることである。他の能力が私たちを欺くかどうかを見つけるためには、その結果が満足する秩序に従って連続性があるかどうかを見る必要があります。」（コルヌ，《知識の基礎に関する試論》）

これらの原理に基づき、我々が所有している抽象化の力を考えるとき、科学の基本概念が経験から導けないものであると仮定することは非哲学者的です。確かに、現実世界を理解する能力は、秩序の支配が維持される形で他の構成要素に適用されれば有用となります。このような新しい推測の舞台において、知的な手続きの法則が変わらなければ、科学の基本概念は現在我々が認識しているものとは異なるかもしれません。

天体の運動に対する重力の効果を近似的に研究するためには、惑星が完璧な球体または楕円体であるという「極限仮定」に参照します。光線が大気中を通る経路を近似的に決定するためには、温度による干渉要因を抽象化した過程があります。そして、人間の知識のすべての高次の領域においても同様の手順が適用されます。

これらの知的なプロセスに関連して興味深いのは、感覚的表現と個々の経験の多様性から一般化された真理や不変の真理への昇華傾向があることです。この傾向と能力が一致する場所では、自然の各系列の関連事実はより正確な秩序を示唆します。それは、その法則を理解力によって捉えるための基礎と機会となり、変化する現象の世界で完璧に忠実ではないが、実際にはその真理が例示される法則があります。

6. 思想の最終的な法則がその形式において数学的であり、誤りの可能性という事実と結びつけると、驚くべき結論の基礎となる。科学的事実にのみ注目した場合、知性の操作と外部自然の運動との間にはほぼ正確な類似性があると考える傾向があるかもしれない。物理科学について詳しいが、自己の能力の本質を深く考えることに慣れていない人が、その能力の法則が数学的であることが証明されたという報告を受けたと仮定する。不信任感が薄れると、印象として生じるのは、知性の順序は物理宇宙の順序と同じくらい必要不可欠であるということだろう。

自然科学の領域において、問題の初期要素と最終要素との絶対的な関連性を数学的形式で示すことは、効果と原因が結びついている物理的必然性を象徴する。無機世界における状態と条件の必要な順列と、そのような証明プロセスにおいて前提と結論との必要な関連性は、平行であるように見える。この「必要」の用語が最初の系列に適用されたかどうか、あるいはすべての正当な理由付けによる彼女の作品に対する推論における命題間の不溶的な関連性に適用されたかどうかについては、疑問があるかもしれない。歴史的には前者を優先するべきであり、哲学的には後者を優先すべきであると考えられるが、その結びつきという事実は争いがないし、指摘される類似性も明白である。

したがって、正当な論理の法則が一貫して遵守された場合、知性的操作と外部自然の操作との間には非常に密接な類似性が存在するだろう。形式と表現において数学的な法則に従うことで、絶対的な服従さえも含めると、二つの系列は共通の特性を持つことになる。思想界と物理界における必然性の支配は完全かつ普遍的である。

しかし、外部の自然観察は、その作用の一貫性と約束された法則への不変な従順が彼女の全領域にわたって支配しているという事実を、ますます確固たる証拠として示し続けています。しかし、知的な世界における過程に対する僅かな注意は、別の状況を明らかにします。正確には、論理の数学的法則は、正しい論理のみの法則であり、その実際的な違反は永続的に再発する現象です。誤りは、物質的なシステムにおいて場所を持たないものですが、ここでは大きな場所を占めています。これを科学が決定する能力の範囲を超えた最終的な事実の一端として受け入れなければなりません。数学的形式の厳格さからも免れない違反を伴う法則が存在することを認めなければなりません。これらの法則には、力によるものではないその本質的な権威があり、自然界の不変な秩序の類推によっても理解することはできません至上性があります。

如所指出的区别是真实的，那么它不会因我们対智力构成其他部分的看法而受到影响。如果我们认为理智是自由的，这是这些推测普遍精神中最符合的観点，那么这种自由必须被视为与必然性的支配相対立，而不是与真理的确切优越性相対立。正确的推理法則可能会被违反，但这并不因此就不真实地存在。即使在极端形式的必然性假設下，它们的性質和权威也不会受到影响。

假定如本著作中所确定的有效推理定律，或更广泛地说，在彻底分析的结论中最终会显现出来的那些定律，只是实际推理过程（无论是正确还是錯誤）所受法則系统的一部分。假定如果我们完全了解这个系统，我们会发现整个智力程序是必然的，就像无机世界的运动一样。最后，基于这一假設的结果，假定不正确的推理或錯誤现象无论出现在何处，都是其他法則干扰了産生正确推理的那些定律所致。

即便如此，仍然存在一些智力法則，它们与其他法則有此特殊特征：即在它们单独指导下完成的所有智力系统的运动都是正确的；任何対它们的干涉不仅是干涉，而且是违反。这必然会使这些法則具有独特和优越的特性。它们会更加明显地表明一个未常实现的最终目的，并拥有内在公正但未必总是能获得服从的权威。

以上便是原文的忠实翻译。

少し考慮すると、自然界の自然法則によって世界が支配されるという状況において類似性を見出すことは不可能です。無機物の自然界には偏見や差別は存在しません。彼女の法則のどの部分も他の部分から切り離して特別に尊ばれるべきものとして宣言することはできませんし、それが最高の目的を果たす唯一の法則であると断定することもできません。むしろ、すべての法則が等しく並立しているように見え、それらの法則の協調的な作用が体系全体の調和そして、私たちが理解できる限りでは、その一般的な\emph{設計}に必要不可欠であることが、我々の知識が広がるにつれてますます明らかになります。無機物世界における秩序の表面上の乖離、例えば惑星系の擾乱、結晶化過程への外力の侵入による中断など、このような現象はしばしばより高いレベルの秩序の概念に収斂するか、あるいはより注意深く見れば通常とは異なる観点を失うことがあります。これらの事実に対する唯一の解釈は、真と偽、正しくと誤りの区別が知性の過程において存在することであるが、物理的必然性の領域には存在しないということです。私たちは有機的な生物の低レベルから意識的な知性の高次の段階に進むにつれて、この対比が徐々に明らかになります。生命現象が表れる場所では、一定の一般的な法則は神秘的な活動原理によって部分的に覆われることがあります。つまり、動物界の構造は特定の一般型に従っており、しかし、その型は美と比例の最高基準に関連して時として不完全に現実化されることがあります。現代の芸術においては、個々の形態を正確に模倣する一方で、すべてのそれらから抽象することで理想的な優雅さや表現性の概念を追求しようとする二つの選択肢が揺れ動き続けます。それは地上の形態では完璧には現われないかもしれません。また、有機型を犠牲にすることなく種が新たな条件や住処に適応するための目的論的な適応はゆっくりと達成されます。しかし、これは外的状況の呪い的な力ではなく、内側からエネルギーを引き出すことによって達成されるようです。生命はそのすべての形態において無機物の静止した固定性とは対照的に存在しますが、それが体現する型の完璧さは多少理想的な性格を持っているかもしれません。最高の\emph{提示された}優美さと適応性を正確に定義することは不可能であるため、この対比は知性的過程と純粋な物理的世界のプロセスとの間の差異よりもはるかに明確ではありません。数学的法則によって両者が統治される確定的な技術的な性格が強調されることにより、それぞれがその役割として行使する権威の基本的な違いがより鮮明になります。

7. この章の一般的な目的に関連する別の例があります。異なる源から得られた真実や事実を比較することによって、教益的な思考の流れが提示されることもあります。その例は、思想の法則が科学的に表現された形と、古代において物理学的考察が、またすべての時代において metaphysical（形而上学的）考察が実際に採用した形との比較にあります。この観察を二つの例で説明し、特にここではその注意を向けたいと思います。

1. 数学の古典において、数における単位性と論理における宇宙との間には科学的な関連があることが示されている（III. 13）。これらはそれぞれの体系において同じ相対的位置を占め、同一形式の法則に従っている。現在、ギリシャ人の心はその活動の初期段階にある。この段階は Bacon の時代や Newton の時代と比べても、人間の知性の発展における重要な時期であり、自然の偉大な問題が開花する一方で観測手段が不足し、その必要性がまだ理解されていなかったときである。このため、「宇宙」および「一」という用語はほとんど同一視されたように見えた。すべての存在がその現れとして考えられる単位性を定義することこそが哲学の最初の目標であった。

（脚注）アリストテレスの『メタ物理学』第1巻にいくつかの文がある。（参照）

2. タレスはこの基本的な単位性を水と見なした。アナキメネスやディオゲネスはそれを空気と見た。メタポント出身のヒッパルコスとエフェソス出身のヘラクレスは火と見なした。パルメニデスは彼の見解が明確でないか、自信がないかのために、存在全体が一つであるという主張をした。メリシスは宇宙が無限であり、変化や運動を受け付けず、自己同一であり、運動自体が存在しないと述べた。（脚注）「エドオケイデートゥス ポアン エペロン アイネア，
カイ アナロイオトオン，カイ アキナイトン，
カイ ヘン，ホモイオン ヘウプス カイ プラエル。
カイン ケイナエスイン アイネア。」（ディオゲネス・ラルティオス、第9巻、4章）

3. アナクサゴラスはクラゾメネス出身であり、おそらく同郷のヒモティムスに影響を受けたように思われるが、知性を世界とその秩序の原因として探求した。（脚注）「ノウヌヌ ドゥテ テウト エプネウ，
デュテ ユソス，デュテ Θυσει，
Το Κόσμον και την Τάξιν παντοδαπής
Αἰτίαν εὑρίσκων ἐκ τῆς σύνεσεως.」（アリストテレス『メタ物理学』第1巻、3）

4. これらの推測の多くに流れるパンテイシズムの傾向は Xenophanes の言葉から明らかである。エレート派学派の創始者である彼は、「天を眺め、一神が存在する」と宣言した。（脚注）「ξενοφάνης δὲ ... εἰς τὸ ὅλον οὐρανὸν ἀποβλέψας, τὸ ἕν εἶναι φαίνεται τὸν θεόν.」（同上）

5. 数多くの単位性の概念を抽象的にも具象的にも、静的なものとしても活発な原理としても考えることができる形の中で、これらの古代の教えにはどれも応用が見られるかもしれない。アリストテレスの著作に私は主に言及しているが、そのような推測に一度は没頭した多くの人々の中でも、おそらく名前だけが残っている可能性が高い。奇妙で示唆的な真実であるが、自然は天体以外の観点からはほとんど説明できない混沌の場面のように見えたし、一般的な信仰はその多神教の中で混乱していた。原始的な単位性の概念さえも、粗雑な形であっても深く根付いた。失望と無限の探求を経て多くの賢者の心に生き延びた。（脚注）「テウト Καί Εγώ Ουκ Περιορίζω Την Απόψιν Μου
Ουκ Σκέπτομαι Την Σκέψιν，
ΕΝΤΑΥΘΑ ΚΑΙ ΕΠΙΣΤΡΕΦΩΝ ΠΕΡΙ ΤΗΝ ΑΛΗΘΕΙΑΝ
ΒΙΟΥ ΦΩΝΕΙΝΤΑ ΜΕΝ ΣΚΕΠΤΟΜΑΙ，
ΕΝΤΑΥΘΑ ΚΑΙ ΕΠΙΣΤΡΕΦΩΝ ΠΕΡΙ ΤΗΝ ΑΛΗΘΕΙΑΝ
ΒΙΟΥ ΦΩΝΕΙΝΤΑ ΜΕΝ ΣΚΕΠΤΟΜΑΙ。」（シクストゥス・エミリウスによって保存され、ティモンの著作とされる）以下はその翻訳である：「私に与えられたものは偏見から解放されたもので、疑問から解放されたものであり、すべてを見渡すことができるものである。長い間欺かれていたが、真実への探求をやめることはなかった。人生の残りの時間は無秩序な旅の中で過ごされ、一つと全てを見つけることはできなかった。」

2. その思考法の等しく密接な関連性は、二次方程式を表す法則と一致し、本稿において「対称法則」と呼ばれるものである。古代の思想が、デュアルイズムという名前の哲学的考察に傾倒する傾向も同様である。エピメネidesの理論\footnote{アリストテレス『物理学』Ⅰ. 4. 6.}は、自然の表面上の矛盾を「争い」と「友情」の2つの対立原理によって説明したものであり、レウキッパの理論\footnote{アリストテレス『物理学』Ⅰ. 4, 9.}は、すべての存在を充満と真空の2つの要素に分解した。これらの理論はその性質を持つ。

宇宙をリューザや弓に例える有名な比較\footnote{$ \pi\alpha\lambda\acute{\iota}\nu\tau\rho{o}\pi{o}\varsigma\ \grave{\alpha}\rho\mu{o}\nu\acute{\iota}\eta\  \ddot{o}\kappa\omega\varsigma\ \pi\varepsilon\rho\ \tau\acute{o}\xi{o}\upsilon\ \kappa\alpha\grave{\iota}\ \lambda\acute{\upsilon}\rho\eta\varsigma.
$---ヘラクレイト、『オリジネスの哲学メンナ』Ⅸ. 9。またプラUTERS, 『イシスとオシリスについて』}は、その「再現する調和」が対立した緊張状態の産物であることを示唆している。

ピタゴラス派の体系において、デュアルイズムを他の要素と組み合わせたものとして認識される10つの基本的な反対概念が認められている：有限と無限、偶数と奇数、単一と多数、右と左、男性と女性、静止と運動、直線と曲線、光と闇、善と悪、正方形と長方形。アルケモンの体系においても同じ基本的なデュアルイズムが受け入れられているが、ピタゴラス派の体系における明確な数値制限は存在しない。

この思想の大きな発展は、古代マニ教の教えに見られる。この教えはペルシャの宗教システムの基礎となり、他の東方地域にも広まり、キリスト教会の歴史においても重要な位置を占めた。デュアルイズムが哲学的意見としての形で発展し始めた段階では、まだ悪の原理の人格化は存在していなかったが、自然の原始的な単一性を前提とする諸理論に続くものであった。アリストテレスは次のように述べている：「明らかに存在する自然には善とは対照的なものが存在しており、秩序と美しさだけでなく、無秩序と醜さも存在した。そして悪なものの方が善なものよりも多く、醜いものは美しいものよりも多かったため、誰かが争いと友情をそれぞれこれらの異なる現象の原因とした。」

ギリシャにおいては、デュアルイズム理論は主に哲学的見解やその補足として認識された。アリストテレスによる要素の導入は、「暖かさ」と「乾燥」そしてそれらの反対概念から始まる有名な例である。プラトンも数学的な原理と関連付けてこれらの生成を説明している。

プラトン派の大きな哲学的対比、「存在と非存在」において、前者が善と真実に関連付けられ、永遠の理念や原型世界と結びついている一方で、後者が悪、誤り、現在の現象の虚幻的な形態に関連付けられている。これらの2つの哲学的考察はここであいまいに融合している。

この思想傾向はギリシャの哲学の若さと成熟期だけでなく、模倣の時代でも見られる。特にプラトンの天才が支配した場所では、その影響は強く感じられた。すべての実在の統一性、その真実と善悪の本質的な同一性、現象的存在の幻覚的虚無性といった視点は、思考の傾向を大きく形成した。

この哲学的見解が促進された時代には、有名な学問や社会活動の日が過ぎ去った後、ギリシャの哲学学校で神秘主義の強力な傾向が現れた。アレオパジテスのディオンニシウスの仮説的な著作は、アリストテレスの教義との接触によって修正され、中世の学術論争に影響を及ぼした。

この哲学的見解は真の信仰と健全な哲学とは一致しないという主張には根拠がない。人間の愛情の優しい影響や日常生活の普通の事象との親密な関わりが、我々の性質の本当の目的である真実の教育の一部であることを考えると、最高の思考対象を過度に抽象化して理解しようとすることは不適切である。

これらの哲学的考察は特にヘリベートの「アルファとオメガ、偉大なる神」（トレントの聖なるラテン詩）で見られる。これらすべての考察が基盤とする原理は次のように述べられている：「何もかもを完全に包摂し、全体を見渡す眼では、何一つ欠けているものがない。」

これらの哲学的考察はさらに、イヨハネ1:5のより高い純朴さと対照されるべきである。アンセルムの著作もまた、宇宙論の理論を思考と存在の類推に基づいて確立するものである。

その原始的な統一性は永遠の真理の中に住んでいるという点で表現されている。自然法則に従って行動し、道徳的主体が正義の指針に従うことは、その真実が行動として現れる形であり、世界自体はその創造主の自己反省する思考の表現である。

16世紀と17世紀には、より古い哲学的考察が再び復活した。ロレンツォ大公の友人や関係者、イギリスで知られるケンブリッジプラトン派、ヨーロッパ各地に散在する思索的な精神は、古代の問題「一、真実、善」を解決するために再び努力したか、それともそのような調査が無駄であることを証明するために努力した。

これらの考察の論理的要素は、より現代的な体系の輪郭から追跡することが容易であり、特にヘーゲルの哲学において絶対の統一性と自己と非自己の区別を型とする自然の関連性との結びつきが基礎となっている。

16世紀と17世紀には、より明確な復活が見られた。ロレンツォ大公の友人や関係者、イギリスで知られるケンブリッジプラトン派、ヨーロッパ各地に散在する思索的な精神は、古代の問題「一、真実、善」を解決するために再び努力したか、それともそのような調査が無駄であることを証明するために努力した。

これらの考察の論理的要素は、より現代的な体系の輪郭から追跡することが容易であり、特にヘーゲルの哲学において絶対の統一性と自己と非自己の区別を型とする自然の関連性との結びつきが基礎となっている。

8. 人間の意見の極めて恣意的な性質を考えると、それらが存在問題の推測的解釈を形成する唯一の形ではないことを期待することも、ここでは主張することもありません。特定の影響下で他の教義の形式が生じますが、これらはしばしば上記のものを隠蔽することがあります。

\footnote{この声明の証拠は、最近出版された「\textit{Origenis Philosophumena}」というタイトルの著書に見つけることができます（その正確性については疑問が残ります）。その中で記述されている初期のキリスト教の Corruption は、多くの場合、オフィテの特性から明らかに、古代の最も底辺的な部分から派生したものであり、特定の哲学的推測の形式を示しています。大抵の場合、これらは二元論的なスキームに属するか、または二つの間で二元論的な関係が追跡できる三原則を認めるものである---\textit{Orig.\textasciitilde{}Phil.}、pp.135, 139, 150, 235, 253, 264。}%

しかし、我々が考慮した特定の理論の広範な普及と、それらが思想の明示された法則と明らかに類似していることから、両者の系統間に関連性があることを正当に想像することができる。他のすべての精神的な行為や手続きは、それぞれ固有的の誤謬に満ちているように、この著作で「二重性の法則」と呼ばれる思想の法則の作用も、単なる一致の欠如を純粋な反対とすることの自己特有の傾向を持つ可能性があり、これにより互補的な部分から成る世界が形成され、対立する力によって基礎的に分割された世界の概念が形成される。このような形で、一部の大規模な但し素早い推論に基づき、二元論は古代の思索形式としての起源を持っていたと考えられるが、二元論がいかなる形でも真実の理論であるかどうかという問題とは無関係に。しかし、ここでは、これらの古代の思索形式が現代において再び現れた場合、実際的な知識の進歩に対する影響について詳細に考慮することがより重要である。前節で述べたことを踏まえて、この点に関して次のように補足的意見を加えたいと思う。

初めに、すべての健全な哲学は、このような推測が実際的なものであることを決定するための手段としての見地から、その判断を下します。自然知識の進歩は、自然の中に中心的な一元性があることを認識する方向に向かうかもしれません。システムの部分間の相互関係によって構成されるような一元性については、疑いの余地がほとんどありませんし、優れた学者たちは、単なるシステムの概念から推測できるよりも深い物理学的力の相関性について考察してきました。

さらに、その一元性の中に、分離と再統合の一般的な原理が含まれている可能性もあります。これらは最高の意志によって生じる自然の関連する多様性の源であるかもしれません。性別や極性の例は、このような見解を支持するために提示されています。私は仮説として申し上げますが、そのような方法で、物事の外側と心の内側の構成が一致している可能性は高くないでしょうか。

しかし、もしこのような一致が証明されるとすれば、それは人間の知識から最後の推論に現れるものであり、科学的な調査の第一原理にはならないでしょう。発見の自然な順序は個別の事実から普遍的な原理へと移行することです。そして、私たちはまだこの道を十分に進んでいないため、すべての特別な差異が最終的にどのような形で融合し、その説明を受け取るのかを決定するのに至っていないという断言が可能です。

%-----------------------File: 221.png----------------------------
私たちはまだこの軌道を十分に進んでいないため、自然のすべての特別な差異が最終的にどのような形で融合し、その説明を受け取るのかを決定するのに至っていないという断言が可能です。

2番目には、仮に思考の形態と実際の自然の構成との間の対応関係が証明されたとしても、その二つのシステム間にどのような結びつきや関係があろうとも、それらの相互独立性という問題には一切影響を与えないでしょう。それは、一方のシステムが他方のシステムの単なる「製品」であるという結果を導くものではありません。

過度な形而上学的な推測への依存は、ある場合においてこのような幻覚を生む傾向を引き起こすことがあります。例えば、多くの悪の存在を説明する試みの中で、その事実に単なる「相対的」な性格があると定義しようとする試みが見られます。それは、善という等価の相対的な要素に対する論理的な反対関係に基づくものであるとします。しかし、この仮定は純粋に任意なものでしかありません。

無限なる智慧と清浄さの眼から見た悪とはどのようなものかについては、私たちには明確な予測ができないでしょう。しかし、私たちは痛みや欠陥、道徳的な反逆、報復的な苦しみなどの形態を通じて、その全ては悲しげで厳格な現実の姿を呈するのみです。そして、それに対する私たちの自然の全防衛力は、純粋な考慮よりも高いレベルでのみ正当化されるでしょう。

先ほど述べた特定の問題についての議論は、哲学における多くの他の争点にも適用されます。例えば、空間と時間の外部的実在性に関する争点などです。私たちはこれらの要素を単なる理解の形式に分解する根拠を持っていませんが、それは私たちの知識の現在の範囲を決定します。

さらに一般化して言えば、多くの現代的な推測には極めて「主観的」な傾向がありますが、その正当性は認められません。知性の視点から見て異なる仮説が観察された事実と同等に整合的である場合、意識の本能的な証言がそれらの相対的な価値についての証拠として権威を持つべきです。

3番目には、思考の法則を研究しても、物事の実際的な構造を決定することができず、その構造に含まれるかつて賢者たちを困惑させ、思索深い人々を悲しませた事実を説明することができるわけではありません。ましてや、それによって現在の存在状態から脱却し、無限と未定義の知識についての直感的な認識が可能であるという教義に正当性を与えることもできるわけではありません。このような知識を求めようと自然の領域内でもその上でもです。

私たちが無限を完全に理解することができないとしても、人間の本質は無限との関係にある可能性があるかもしれません。思考の法則を完全に理解し、最も一般的な意味でそれらに基づく方法論を確立しようとすると、少なくとも「宇宙」や「永遠」という言葉で表現される要素が含まれていると言えるでしょう。幾何学の純粋な抽象においても同様ですが、論理の領域においても、真実の帝国は想像力のそれよりも広い範囲にあると見ることができます。多くの知識の分野があり、これらを完全に把握するには外部から観察する必要がある場合がありますように、知的プロセスの理論が有限な対象のみに適用されるとしても、思考の領域においてすべての限界が撤廃される領域が認識されることが含まれているようです。したがって、一方では思考の法則とその類似性から得られる証拠は、あまり大胆な神秘主義の結論を十分な基礎とするには不十分であるとしても、それらを全く説明的でないものとは見なすことは誤りです。

これらの法則は私たちの知的な本質の一部であり、単に形式的な推論の条件を規定するだけでなく、その存在を他の形で示す可能性があるかもしれません。イオニアとイタリアの哲学者たちによる物理的考察の特有な傾向と結びついている理由と同じ理由により、それらはより一般的かつ正当な思考の傾向とも関連付けられるべきであると考えられます。少なくとも偶然の影響によってそのような深遠な精神が生まれたわけではありません。自然の外的な壮麗さと最も強く接しているときに、永遠の真実性に意識させられ、夜空から秩序の明るい示唆を読み取り、山々の中の荒涼とした場面から抽象的な永遠を超える暗示を感じ取ったときです。

\footnote{聖書詩篇90章2節}

9. さらに、一部の人々にとっては過度に推測的な内容のように思われる思考の流れを更に追跡するのを控えめにするならば、私たちは導き出された肯定的な結果について一時的に振り返ってみましょう。
私たちの性質には、具体的な経験事実から一般的な科学の基礎となる命題へと昇華する能力があることが明らかになりました。また、一般に受け入れられた真実の命題から特定の結論を導き出す能力もあります。
これらの能力は、自然法則とは異なる、独自性を持ち、直接的ではない正確な科学的な表現が可能な法則によって制御されていることが示されました。
さらに、この知的な手続きが我々を取り巻くシステムの条件と一致していることも明らかになりました。つまり、一般的な類似性を持つ種の存在、一般の法則に従って関連付けられた事実、そして獲得した知識の恒久性と秩序との統合性は、既得の知識を安定させると共に、無限の進歩への希望の基礎を築きます。
人間の性質は、観察されたまたは現れた傾向とは独立して、真実に対する一定の関係を持っていることが明らかになりました。この関係は、推測的な知識の対象として詳細に研究することができ、また物理学の各分野と同様にその価値が高く、それらを同じ観点から検討することも可能です。
特に注意すべきは、知性の構成についての見解で、それが決定的な性格を持つ法則によって制御されているものの、必然性による力ではなく作用していることを示すものです。これは、運命の支配から救われたまま、偶然性に委ねられることなく表現されます。この視点を支持することは、類推に基づいて提供される暗示について、我々の性質の別の高次の側面についても可能性があると受け取る少なくとも確かな見解を含んでいます。
つまり、知識の領域だけでなく義務の領域でも固定された法則に従って制約されているという点です。その構成は理想的な基準や最終的な目的に関連付けられています。
確かに、科学的な研究は、道徳的世界と物質的世界の具体的な差異を無視する傾向があると考えられたり、前者が正確な知識の対象とは適切に言えないものであるという見方を助長することがあります。すべての専門的な追求は偏った視点を生む可能性がありますし、物理的必然性の支配が絶対的かつ至上である場面で長い間深く考え込む心は、別の秩序がある可能性を認めることを難しくするかもしれません。
しかし、この専門性の排他的さが知識の特定の領域への献身による偏見を生む理由です。知的な現象に対する科学的方法の適用は、不偏な探求精神を持って行われ、誤りや混乱の要素を無視せずに受け入れるという前提のもとで行われることで、より公正な類推材料が提供されることが示唆されます。

10. 以上の問いがどのような実用的な目的に結びついているのかと問われれば、教育という対象物についての探求は、人間の性質に対する推測的な見解によって主に決定される諸要素の中の一つであることを指摘することができます。教育を最も広い意味で考えるとすれば、その本質と方法に関する全ての問いへの最終的な基礎は、人間とは何か、それぞれの能力が設計された目的は何であるか、どのような動機が持続的な行動を引き出し、そして最善かつ最も安定した結果を引き出すことができるかという前段階の理論に置かれるべきです。これらの質問がここ示される関係において完全かつ公正に考慮されたかどうかについては疑問が残ります。

古代の純粋なモデルによる Taste の最高の陶冶、現代物理学の事実と理論に対する最大限の知識を、私たちの性質の広い視点から見ることで、これらは完璧な知的訓練の一端に過ぎないと思われます。同じ視点から手段を考えるとき、現在の時代における競争心や欲求の精神へのほぼ独占的な依頼が、我々の性質において直ちなる目的のために植え付けられたより永久的な動機を弱める傾向があるかどうかという問いに導かれるかもしれません。

これらの問題と、多くの他の問題について、正当な権限の範囲、思想の自由と感情、習慣、風習に対する訓練との調和、そして全体としての道徳的観点において何等の不確実性、何等の多様な実践を我々は遭遇するのでしょうか！しかし、理性の冷静な視点から見れば、人間の努力の範囲内でより重大かつ重要な対象は、私が述べたように最も広い意味で考えたときです。

何らかの形で人間性に対するより正確かつ明確な理解を促すものは、その実在するあらゆる側面においてこれらの問いをより狭い範囲に収め、そして可能な解決の限界を制限します。すなわち、推測的な問いが行動原理として最も重要な原則を結実させる可能性があることを示すことができます。

11. このような推測が最も正当に及ぼす影響は、数学が人間の知識体系における位置と役割、そして数学的研究が知的鍛錬としての性質や目的についての問いに対する答えである点にあるかもしれません。最近の議論を注意深く追った者は、この問いが重要なものではないと考える人はいないでしょう。数学の地位が両面において基礎的なものであるという主張者たちは、実際の事象の構造から最も強い証拠を得ていると主張します。物質的世界は、その全体にわたって数的关系によって支配されています。全ての動力学的、化学的、電気的、熱学的な作用も、自己同一性を超越して相互変換されるという程度まで、数的关系によって結びついているように見えます。

しかし、この見解は単なる事象の構造に基づくものではなく、より深い基盤に立っています。思考の法則、概念形成と論理推理の全ての過程において、言語がその表現または道具である全ての操作についても、これらは数学の公認された過程の法則と同じ種類です。我々が共通の意味で適切に考えたり、論じたりするためには、これらの法則を理解する必要があるとは主張されていません。人々は意識的にそれらの要素に基づいて推論を行いますが、それでもなお数学的な形式に至る人間の思考を究極まで追及すると、我々の性質によって数学的科学が人間の知識体系において基礎的な位置を占めているという前提が浮かび上がります。そして、これらを完全に無視する精神的文化システムは根本的に不完全であると推測できます。

%-----------------------File: 226.png----------------------------
この証拠に基づいて、我々は数学の科学が人間の知識体系において構成上基礎的な位置を占めているという前提を持てます。そして、これらを全く無視する精神的文化システムは完全でなく、根本的とも言えません。

しかし、同じような考慮は、数学の研究とその応用が知識や教養の基礎であるとする見解にも等しく誤りを示します。もし物質的な枠組みが数学的に構成されているなら、それは単にそうなるだけではありません。もし意識的または無意識的に形式的な論理で思考する心が、その他の感情や行動の能力、美と道徳的な適応性に対する感覚、深い情動と愛情の源を通じて、異なる秩序の事物との関係を主張するなら、それは数学的な法則に従うだけでなく、それらの他の側面にも依存します。さらに、広範な知性的視野と、真理がすべての形態や表現において示す同情力は、論理的論証の強さや精巧さで測定されるものではありません。無限に広大で、普遍的な秩序を保ち、法則の恒常性を持つ物質宇宙の揭示も、最も精密な証明の過程を通じて完全に理解されるとは限りません。そして、生活の利益と義務を含む視野においては、単なる論理的推論だけでは、それらが提示する重い問題をどれほど理解できるでしょうか！したがって、数学的または推論的な能力の養成は知性の教養の一端であることは事実ですが、それはただの一端に過ぎません。知識や心のどの一分野も完全に排除したり、その上位に押し上げたりするという偏見は、判断上の誤りだけでなく、純粋な真理への devotion から不可分に離れる知性的謙虚性の欠如を示すものです。それは、人間が制定し、取り消すことのできないものであると考えられている現実の構造を批判的評価する役割を主張します。

%-----------------------File: 227.png----------------------------
古代で公正に理解された真理の一元性でありながら多様性を持つという概念を破壊します。この誤りの多くは、科学教育の特殊かつ孤立した性格により存在しているように思われます——その性格はまたそれを促進します。哲学の研究は、例外的な数々の例を除いて、それぞれの知識分野の発展に伴ってそれらが互いに関連する関係を決定することの任務を果たすために遅れています。しかし、特定の悪状態をより大きなシステムの一端として見ること、そしてそれが知識が単なる世俗的なものであるという過度な観念と結びついていること、それに加えて、その世俗的利益に基づく正当な動機が既に述べたように過度に優位にあることを考慮することにより、それは不可能ではありません。したがって、この問題を特定するためには、それ自体の目的に対する哲学的研究の考察から自然と生じるか、またはその達成のために必要な読書や反省を通じて生じる必要があります。

極端な場合においても、これらの動機がより高い行動原理によって制御されず、優れた心の個人的な影響で補正されない限り、知識の対象に対する思考基準を低下させる傾向があることは明白です。そして、これらの条件が現実化するほど、同じ結果が生じます。おそらくこれが、しばしば変幻自在な人類の姿よりも、光の父への知性的な忠誠を持つ者たちよりも、真理の異なる部分の統一性、生命的な関連性、そして道徳的目的に対する下部に置かれた地位をより正しく理解している者がいる理由かもしれません。しかし、これらの問題はここではさらに追及することができません。一部の人々には、この作品の明示的な目的から外れているように思われるかもしれませんが、それらの考慮はその目的に関連する推測から自然と生じたか、またはその達成のために必要な読書や反省を通じて生じたものです。

\centerline{\textsc{終わり。}}

\begin{text}
*** PROJECT GUTENBERG の電子書籍『論理の法則に関する調査』終了 ***

***** このファイルは 12345-pdf.pdf や 15114-pdf.zip と名付けられ、以下の場所に各種形式で保管されます：
    https://www.gutenberg.org/1/5/1/1/15114/

更新版が既存の版を置き換える場合、古い版は再命名されます。

印刷物から著作権保護されていない作品を作成することにより、アメリカ合衆国ではこれらの作品に対する著作権は存在せず、財団（そしてあなたも）が米国の著作権者の許可なく且つ著作権料を支払うことなく複製と配布を行うことができます。著作権法の一般条項部分に記載されている特別な規則により、Project Gutenberg-tm の概念と商標を保護します。Project Gutenbergは登録商標であり、電子書籍を有偿で販売する場合、その使用には著作権許諾料を支払う必要があります。電子書籍を無料で配布する場合は、著作権許諾料を支払うことなく商標を使用することは非常に簡単です。この電子書籍は、創作物の作成、報告、公演、研究など、アメリカ合衆国でのほとんどの目的に使用することができます。米国の著作権法で保護されていない電子書籍については、 Redistribution は商標許諾に従って行われます。

START: 全体的なライセンス

プロジェクト・ゲーテンベルグの完全な著作権免許
この作品を配布または使用する前に、以下をご覧ください。

プロジェクト・ゲーテンベルグの使命を推進するために電子書籍の無料配布を促進することを保護するため、`Project Gutenberg' と関連付けられた言葉を使用してこの作品（またはその他の関連作品）を使用または配布する場合、以下の全著作権免許に従うことに同意します。ファイルと共に提供されるか、www.gutenberg.org/license でオンラインで利用可能です。

第1条. 著作権免許の一般条項と Project Gutenberg-tm 電子書籍の再配布
\end{text}

1.A. このProject Gutenberg-tm電子作品の一部を読んだり使用したりすることにより、あなたはその著作権および知的財産（商標・著作権）に関する契約のすべての条項に同意し、理解し、受け入れたことを示します。このライセンスと知的財産（商標・著作権）に関する契約のすべての条項に同意しない場合、あなたはProject Gutenberg-tm電子作品の所有コピーを全て返却または破棄する必要があります。もし料金を支払い、Project Gutenberg-tm電子作品のコピーを取得したりアクセスしたが、この契約の条項に従うことに同意しない場合は、第1.E.8条に定められたように、その料金を支払った人や機関から返金を求めることができます。

1.B. 「Project Gutenberg」は登録商標です。この協定のすべての条項に同意することを約束した人々によってのみ、電子作品に関連付けることができます。多くのProject Gutenberg-tm電子作品に対しては、この協定の全部の条項を遵守せずにでもいくつかの行動が可能であることに注意してください。第1.C条をご覧ください。この協定の条項に従ってProject Gutenberg-tm電子作品への自由なアクセスを保護する手助けをする場合、多くのProject Gutenberg-tm電子作品に対してはさらに多くの行動が可能です。第1.E条をご覧ください。

1.C. ジェンキンズ文学アーカイブ財団（以下、「財団」または「PGLAF」という）は、プロジェクト・ゲッutenberg-tm電子作品の集まりに所有する編纂著作権を有しています。この集められた個々の作品のほとんどは、アメリカ合衆国において公の領域に入っています。もし個々の作品がアメリカ合衆国における著作権法によって保護されていない場合で、かつあなたがアメリカ合衆国に所在している場合、私たちはあなたをプロジェクト・ゲッutenberg-tmの著作を複製、配布、上演、表示または派生作品を作成することから防ぐ権利を主張しません。ただし、すべてのプロジェクト・ゲッutenberg-tmに関する参照が除去されている限りです。もちろん、私たちはプロジェクト・ゲッutenberg-tmの使命を促進するための無料で電子的作品へのアクセスを共有することを支援したいと考えています。この契約に基づき、プロジェクト・ゲッutenberg-tm著作権許諾証明書とともに無償で他の人に共有する際にこの作品を同じ形式で保持することで、これらの条件に従うことができます。

1.D. あなたが所在する場所の著作権法も、この作品に対するあなたの行為を制限します。多くの国における著作権法は常に変化しています。アメリカ合衆国以外に所在している場合、この契約と同様にダウンロード、複製、表示、上演、配布またはこの作品やプロジェクト・ゲッutenberg-tmの他の作品に基づいて派生作品を作成する前に、あなたの国の著作権法を確認してください。財団は、アメリカ合衆国以外のどの国においても著作権状況についての表明を行っていません。

1.E. プロジェクト・ゲッutenberg-tmに関する参照がすべて除去されていない限り：

1.E.1. 以下の文は、`プロジェクト・ゲッutenberg'という言葉が現れるか、またはその言葉と関連付けられているどのプロジェクト・ゲッutenberg-tm作品のコピー（`プロジェクト・ゲッutenberg'という言葉が現れるか、またはその言葉と関連付けられているすべての作品）をアクセス、表示、上演、視聴、複製または配布する際には、目立つ位置に表示される必要があります。これは、全著作権許諾証明書（$$\copyright \, \text{The Project Gutenberg Literary Archive Foundation}$$）へのアクティブなリンクや、またはその著作権許諾証明書への直ちなアクセスが含まれている場合です。

この電子書籍は、アメリカ合衆国およびその大部分の地域で誰でも無料で利用でき、制限がほとんどありません。これをコピーしたり、無償で配布したり再利用したりする際には、この電子書籍に添付されているプロジェクト・ゲッutenbergライセンスまたはwww.gutenberg.orgでオンラインで見つけることができるライセンスの条項に従ってください。あなたがアメリカ合衆国以外の地域に所在している場合、この電子書籍を使用する前に、所在する国の法規を確認してください。

1.E.2. 个别なプロジェクト・ゲッutenberg-tm電子書籍が米国著作権法によって保護されていないテキスト（著作権保持者の許可を得て掲載されている通知が含まれていない）から派生した場合、その作品はアメリカ合衆国の誰でも無償でコピーしたり配布したりすることができます。もしこの電子書籍を再配布するか、`プロジェクト・ゲッutenberg'という言葉と共にまたはそれに付随して掲載されている場合、あなたは1.E.1から1.E.7までの条項に従うか、または1.E.8または1.E.9の条項に基づいて著作権者の使用許可を得る必要があります。

1.E.3. 个别なプロジェクト・ゲッutenberg-tm電子書籍が著作権保持者の許可を得て掲載された場合、その利用と配布は1.E.1から1.E.7までの条項に従うとともに、著作権保持者が付与した追加の条件にも従う必要があります。著作権保持者の許可を得て掲載されたすべての作品のプロジェクト・ゲッutenberg-tmライセンスが最初に記述されている場合、これらの追加の条件はこの作品の開始部分で見つけることができます。

1.E.4. この電子書籍またはその一部を含むファイルから全プロジェクト・ゲッutenberg-tmライセンス条項を削除したり、取り外したりしないでください。

1.E.5. この電子書籍またはその一部の任意の部分をコピーしたり表示したり実行したり配布したり再配布したりする際には、段落1.E.1に記載された文を目立つ位置で表示し、プロジェクト・ゲッutenberg-tmライセンスの全文へのリンクや即時アクセスを提供する必要があります。

1.E.6. この作品をバイナリ形式、圧縮形式、マークアップされた形式、非所有権的または所有権的な形式（包括的な単語処理形式やハイパーテキスト形式を含む）に変換し、配布することは可能です。ただし、Project Gutenberg-tmの作品へのアクセスやコピーを、` Plain Vanilla ASCII ' 形式以外の形式で提供する場合、または Project Gutenberg-tm公式ウェブサイト（www.gutenberg.org）に掲載されている公式版とは異なる形式で複製品を配布する場合は、ユーザーに対して追加料金を請求せず、そのようなコピーを提供する方法、または要求に基づいてコピーを取得できる方法を提供する必要があります。代替形式は、第1.E.1条に規定されるProject Gutenberg-tmライセンスの全文を含む必要があります。

1.E.7. Project Gutenberg-tm作品へのアクセス、閲覧、表示、演奏、複製、配布に対して料金を徴収しないでください。ただし、第1.E.8条または第1.E.9条に準拠している場合を除きます。

1.E.8. Project Gutenberg-tm電子作品のコピーを提供するか、またはアクセスを提供し、配布する料金を徴収することは可能です。ただし、以下の条件を満たす場合に限ります：

* あなたが既に適用税額を計算するために使用している方法を使用して、Project Gutenberg-tm作品から得られる総利益の20%の版税料を支払います。この費は、Project Gutenberg-tm商標の所有者に対して支払われるべきですが、彼はこの条項に基づいて版税をプロジェクト・ゲッutenberg文学アーカイブ財団に寄付することに同意しています。版税は、あなたが（または法律により準備する必要がある場合を含む）定期的な税金申告書を作成した日から60日以内に支払う必要があります。版税は明確にそのようなものとしてマークされ、プロジェクト・ゲッutenberg文学アーカイブ財団の指定された住所に送付されるべきです。

* 30日以内に、著作権の全条項に同意しない旨を書面（またはメール）で通知したユーザーに対して、支払った金銭全額を返還する。あなたはそのようなユーザーに対し、物理的媒体に保持されている作品のすべてのコピーを返却または破棄させ、Project Gutenberg-tm作品への使用および他のコピーへのアクセスを停止させるよう要求しなければならない。

* 90日以内に電子作品に缺陷があることを発見し、あなたに通知した場合、第1.F.3条に基づき、支払った金銭全額または代替コピーの返還を行う。

* この契約の無料配布条件を遵守する。

1.E.9. 異なる条件でProject Gutenberg-tm電子作品またはそのグループを料金を徴収して提供したい場合や、配布したい場合は、米国著作権法によって保護されていない作品を作成するためにProject Gutenberg-tmコレクションを構築するためのProject Gutenberg文学アーカイブ財団、Project Gutenberg-tm商標の管理者から書面による許可が必要である。第3条に従って財団に連絡すること。

1.F.

1.F.1. Project Gutenbergのボランティアおよび従業員は、米国著作権法によって保護されていない作品を特定し、著作権調査を行い、転記し、校正することでProject Gutenberg-tmコレクションを作成します。これらの努力にもかかわらず、Project Gutenberg-tm電子作品やその保存媒体には、「欠陥」が含まれる可能性があり、これらは但し書に定めるもの以外の不完全なデータ、転記誤り、著作権または他の知的財産権侵害、破損または損傷したディスクその他の媒体、コンピュータウイルス、またはあなたの機器で読み取れないコンピュータコードを含むことができる。

1.F.2 保証の制限、損失に関する放棄 - 第1.F.3段落で説明される「交換または返金権」以外については、プロジェクト・ゲッutenberg文学アーカイブ財団は、プロジェクト・ゲッutenberg-tm商標所有者であり、本契約に基づき電子書籍を配布する他のいかなる第三者も、あなたに対して損失、費用、包括的な弁護士報酬等に対する一切の責任を放棄します。あなたは、過失、厳格な責任、保証違反または契約違反による損害に対する補償の remedies がないことに同意します。本財団、商標所有者、および本契約に基づき配布する第三者は、このような損失が発生した場合でも、実質的、直接的、間接的、結果的、懲罰的または偶然的な損害についてあなたに対して責任を負わないことに同意します。

1.F.3 交換または返金の制限的权利 - この電子書籍を受け取って90日以内に欠陥を見つけた場合、あなたはその作品から受け取った人物に対して、説明書と共に支払った金銭（あれば）を返還するための書面による説明を送ることができます。物理媒体で受け取った場合は、説明書と共にその媒体も返却してください。欠陥のある作品を提供した人物または法人は、代替版を提供することで返還を放棄することを選択することができます。電子的に受け取った場合、その作品を提供した人物または法人は、代替の電子配信を提供することで返還を放棄することを選択することができるでしょう。第二のコピーも欠陥がある場合は、問題を修正する追加の機会を与えることなく、書面による返還要求を行うことができます。

1.F.4 本契約第1.F.3段落で定められた交換または返金の制限以外については、この作品は「現状」で提供され、「明示的または黙示的ないかなる種類の保証もなし」と提供されます。商品性または特定目的への適合性に関する保証を含むがこれらに限定されないすべての保証も含まれます。

1.F.5. 一部の州では、特定の明示的保証を否認したり、特定の種類の損害の除外または制限を許可しない場合があります。この契約において定められた否認または制限が適用される州の法に違反する場合、この契約は、該当する州の法律によって可能な最大限の否認または制限と解釈されます。この契約のどの規定も無効であったり執行不能であったりしても、その他の規定は依然として有効です。

1.F.6. 被告人化 - あなたは、以下のいずれかが発生した場合に、被告、著作権所有者、被告の代理人または従業員、Project Gutenberg-tm電子作品をこの契約に基づいて提供する人々、そしてProject Gutenberg-tm電子作品の制作、推奨および配布に関連するボランティアから、すべての責任、費用、および経費（法的費用を含む）を免責することに同意します：(a) この著作物またはProject Gutenberg-tmの他の著作物の配布、(b) Project Gutenberg-tmの著作物への変更、修正、追加または削除、および(c) あなたが引き起こした欠陥。

2. Project Gutenberg-tmの使命についての情報

Project Gutenberg-tmは、最も多様な種類のコンピュータで読み取れる電子作品を自由に配布することと同義です。そのコンピュータには古めかしいものから新しいものまでが含まれます。それは数百人のボランティアの努力と、様々な分野の人々からの寄付によって存在しています。

ボランティアの支援と資金提供は、これらのボランティアが必要とする援助を提供することにより、Project Gutenberg-tmの目標を達成し、そしてProject Gutenberg-tmのコレクションが今後も自由に利用可能であることを確保するための不可欠な要素です。2001年には、Project Gutenberg文学アーカイブ財団が設立され、Project Gutenberg-tmとその将来世代のために安定した永久的な未来を提供することを目指しています。あなたの努力と寄付がどのようにしてそれが可能になるかについて学ぶための情報は、www.gutenberg.orgの財団情報ページをご覧ください。

第3節 プロジェクト・ゲンジューントの文学アーカイブ財団に関する情報

プロジェクト・ゲンジューントの文学アーカイブ財団は、501(c)(3)条項に基づく非営利教育法人で、ミシシッピ州の法律に従って組織され、国税庁から税制上の優遇措置が与えられています。財団のEIN（米国内収入局の税番号）は64-6221541です。プロジェクト・ゲンジューントの文学アーカイブ財団への寄付は、米国連邦税法と州税法により許容される最大限まで税金が控除されます。

財団の事務所は、ユタ州サルト・ラックシティの北1500西809番地にあり、電話番号は(801) 596-1887です。電子メールでの連絡先リンクと最新の連絡先情報は、財団のウェブサイトおよび公式ページwww.gutenberg.org/contactで見つけることができます。

第4節 プロジェクト・ゲンジューントの文学アーカイブ財団への寄付に関する情報

プロジェクト・ゲンジューント-tmは広範な公衆の支援と寄付なしでは存続することができません。機械読取可能な形式で自由に配布できる公的著作権作品やライセンスされた作品の数を増やすという使命を遂行するために、古式的な機器も含む最も広範な設備によってアクセスできる形での配布が可能となるようにするためには、多くの小額寄付（$1から$5,000まで）が特に重要です。

財団は、全50州の米国の慈善法と慈善寄付に関する法律を遵守することにコミットしています。これらの要件は一様ではなく、これらを満たし維持するためには相当な努力と多くの書類作成が必要となります。寄付を求める要件が満たされていない州では寄付を募集していませんが、そのような州から無償で寄付を申し出る donors からの寄付は禁止されていないことを知っています。

国際的な寄付は心から歓迎しますが、私たちは国外からの寄付についての税制上の扱いに関していかなる声明も発行できません。アメリカ合衆国の法規だけでも、私たちの小さなスタッフにとっては洪水のように押し寄せます。

現在の寄付方法と連絡先については、Project Gutenbergのウェブページをご確認ください。寄付は切手、オンライン支払い、クレジットカードによる寄付など、さまざまな方法で受け入れられます。寄付をするには、以下のリンクを訪問してください：www.gutenberg.org/donate

5. Project Gutenberg-tm電子作品に関する一般情報

Project Gutenberg-tmの概念を提唱したのは、Michael S. Hart教授です。彼は40年間にわたり、自由に共有できる電子版作品のライブラリを作り出し、配布しました。その活動には、志願者による松散なネットワークが支えになっていました。

Project Gutenberg-tmの電子書籍は、しばしば複数の印刷版から作成されますが、これらはすべてアメリカ合衆国では著作権保護対象外と確認されています。したがって、私たちは特定の紙媒体版に準拠する必要はありません。

多くの人々は私たちのウェブサイトから始めます。そこには主要なPG検索機能があります：www.gutenberg.org

このウェブサイトでは、Project Gutenberg-tmに関する情報も提供しています。寄付方法や、新しい電子書籍を制作する手助けをする方法、そしてメールニュースレターに登録して最新の電子書籍情報を入手する方法などについても説明されています。


\backmatter

\chapter{訳者・数理校閲後記}

本書は、近代情報科学・デジタル文明の精神的始祖であるジョージ・ブールの主著『思考の法則』（\textit{An Investigation of the Laws of Thought}, 1854）の全22章を、Project Gutenberg（\#15114）の底本に基づき完全翻訳・数理校閲したものである。

ブールが提示した「論理とは計算である」という洞察は、170年の時を経て、現代の半導体チップ、プログラミング言語、検索エンジン、そして巨大なAIモデルの根底を貫く普遍的原理となった。

本書の全訳版が、論理と確率の数学的源流に触れ、人間精神の思考の深遠さに思いを馳せるすべての探求者にとってかけがえのない道標となることを願ってやまない。

\begin{flushright}
2026年10月\\
教育情報分析研究会
\end{flushright}

\end{document}
